\documentclass[11pt,oneside]{article}

\usepackage[T1]{fontenc}
\usepackage{lmodern}
\usepackage{microtype}
\usepackage{geometry}
\geometry{margin=1in,headheight=14pt}
\usepackage{amsmath,amssymb,amsthm,mathtools}
\usepackage{booktabs,longtable,tabularx,array,multirow}
\usepackage{enumitem}
\usepackage{xcolor}
\usepackage{xurl}
\usepackage{hyperref}
\usepackage{listings}
\usepackage{fancyhdr}
\usepackage{titlesec}
\usepackage{needspace}

\definecolor{leanblue}{RGB}{24,91,145}
\definecolor{newgreen}{RGB}{15,105,68}
\definecolor{computeviolet}{RGB}{98,55,145}
\definecolor{classicalgray}{RGB}{85,92,99}
\definecolor{codegray}{RGB}{248,248,248}
\definecolor{rulegray}{RGB}{92,101,111}

\hypersetup{
  colorlinks=true,
  linkcolor=leanblue,
  citecolor=newgreen,
  urlcolor=leanblue,
  pdftitle={The Alaoglu--Erdos Integer-Exponent Conjecture: Unified Report},
  pdfauthor={The ProveIt project},
  pdfsubject={Machine-checked partial results, exact conditional structure, determinant
methods, and research directions}
}
\urlstyle{same}

\pagestyle{fancy}
\fancyhf{}
\lhead{Alaoglu--Erd\H{o}s integer-exponent conjecture}
\rhead{Unified report}
\cfoot{\thepage}
\renewcommand{\headrulewidth}{0.4pt}

\titleformat{\section}{\Large\bfseries}{\thesection}{0.75em}{}
\titleformat{\subsection}{\large\bfseries}{\thesubsection}{0.75em}{}
\titleformat{\subsubsection}{\normalsize\bfseries}{\thesubsubsection}{0.75em}{}

\setlist[itemize]{leftmargin=2em,itemsep=0.28em,topsep=0.4em}
\setlist[enumerate]{leftmargin=2.2em,itemsep=0.28em,topsep=0.4em}
\setlist[description]{leftmargin=3.2cm,labelsep=0.6cm,itemsep=0.35em,topsep=0.4em}
\setlength{\parindent}{0pt}
\setlength{\parskip}{0.55em}
\setlength{\emergencystretch}{3.5em}

\newtheorem{theorem}{Theorem}[section]
\newtheorem{proposition}[theorem]{Proposition}
\newtheorem{lemma}[theorem]{Lemma}
\newtheorem{corollary}[theorem]{Corollary}
\newtheorem{conjecture}[theorem]{Conjecture}
\newtheorem{question}[theorem]{Question}
\theoremstyle{definition}
\newtheorem{definition}[theorem]{Definition}
\newtheorem{observation}[theorem]{Observation}
\newtheorem{target}[theorem]{Research target}
\newtheorem{warning}[theorem]{Caution}
\newtheorem{experiment}[theorem]{Computational experiment}
\theoremstyle{remark}
\newtheorem{remark}[theorem]{Remark}

\newcommand{\N}{\mathbb N}
\newcommand{\Nzero}{\mathbb N_0}
\newcommand{\Z}{\mathbb Z}
\newcommand{\Q}{\mathbb Q}
\newcommand{\R}{\mathbb R}
\newcommand{\C}{\mathbb C}
\newcommand{\Abar}{\overline{\mathbb Q}}
\newcommand{\thetaq}{\vartheta}
\newcommand{\Sol}{\mathcal S}
\newcommand{\Cand}{\mathcal C}
\newcommand{\Pair}{\mathcal P}
\newcommand{\LambdaTheta}{\Lambda_{\thetaq}}
\newcommand{\Zthree}{\Z[1/3]}
\newcommand{\Ztwo}{\Z[1/2]}
\newcommand{\fract}[1]{\left\{#1\right\}}
\newcommand{\distZ}[1]{\left\lVert #1\right\rVert_{\Z}}
\newcommand{\vpp}[2]{v_{#1}\!\left(#2\right)}
\newcommand{\fall}[2]{\left(#1\right)_{\!\underline{#2}}}
\newcommand{\clog}{\operatorname{clog}}
\newcommand{\Apos}{\Abar_{>0}}
\newcommand{\LL}{\mathcal L}
\newcommand{\GammaSat}{\Gamma}
\newcommand{\PGL}{\operatorname{PGL}}
\newcommand{\Qbar}{\Abar}
\newcommand{\sha}[1]{\texttt{\detokenize{#1}}}
\newcommand{\lean}[1]{\texttt{\detokenize{#1}}}
\newcommand{\commit}[1]{\texttt{#1}}
\newcommand{\statuslean}{\textcolor{leanblue}{\textbf{[kernel-verified Lean]}}}
\newcommand{\statusnew}{\textcolor{newgreen}{\textbf{[paper proof in this report]}}}
\newcommand{\statuscomp}{\textcolor{computeviolet}{\textbf{[software computation]}}}
\newcommand{\statusexactcomp}{\textcolor{computeviolet}{\textbf{[exact computation]}}}
\newcommand{\statusknown}{\textcolor{classicalgray}{\textbf{[classical/open background]}}}
\newcommand{\statusdraft}{\textcolor{computeviolet}{\textbf{[Lean draft, not yet accepted]}}}
\newcommand{\statusrepo}{\textcolor{leanblue}{\textbf{[repository build reported]}}}
\newcommand{\statusopen}{\textcolor{classicalgray}{\textbf{[proposed target]}}}
\newcommand{\DeltaTri}{\mathcal T}
\newcommand{\supp}{\operatorname{supp}}
\newcommand{\Span}{\operatorname{span}}
\newcommand{\rank}{\operatorname{rank}}
\newcommand{\vp}{v_p}
\newcommand{\ee}{\mathrm e}
\newcommand{\ord}{\operatorname{ord}}
\newcommand{\num}{\operatorname{num}}
\newcommand{\den}{\operatorname{den}}
\newcommand{\rad}{\operatorname{rad}}
\newcommand{\Law}{\operatorname{Law}}
\newcommand{\vol}{\operatorname{vol}}
\newcommand{\cO}{\mathcal O}
\newcommand{\E}{\mathbb E}
\newcommand{\Var}{\operatorname{Var}}
\newcommand{\Cov}{\operatorname{Cov}}
\newcommand{\dist}{\operatorname{dist}}
\newcommand{\AEname}{Alaoglu--Erd\H{o}s}
\newcommand{\DeltaV}{\mathop{\Delta}\nolimits}

\newenvironment{statusbox}[1]
{\begin{center}\begin{minipage}{0.94\textwidth}\raggedright\hbadness=10000\hrule
\vspace{0.45em}\textbf{Status:} #1\par\vspace{0.35em}}
{\vspace{0.35em}\hrule\end{minipage}\end{center}}

\newenvironment{resultbox}
{\begin{center}\begin{minipage}{0.94\textwidth}\hrule\vspace{0.45em}}
{\vspace{0.35em}\hrule\end{minipage}\end{center}}

\lstdefinestyle{leanstyle}{
  basicstyle=\ttfamily\footnotesize,
  backgroundcolor=\color{codegray},
  frame=single,
  rulecolor=\color{rulegray},
  showstringspaces=false,
  breaklines=true,
  columns=fullflexible,
  keepspaces=true,
  tabsize=2,
  captionpos=b
}
\lstdefinestyle{pythonstyle}{
  language=Python,
  basicstyle=\ttfamily\footnotesize,
  backgroundcolor=\color{codegray},
  frame=single,
  rulecolor=\color{rulegray},
  numbers=left,
  numberstyle=\tiny\color{rulegray},
  stepnumber=1,
  numbersep=8pt,
  showstringspaces=false,
  breaklines=true,
  columns=fullflexible,
  keepspaces=true,
  tabsize=4,
  captionpos=b
}
\lstdefinestyle{proofpython}{
  language=Python,
  basicstyle=\ttfamily\scriptsize,
  columns=fullflexible,
  keepspaces=true,
  breaklines=true,
  breakatwhitespace=false,
  frame=single,
  showstringspaces=false,
  tabsize=4
}

\begin{document}
\pagenumbering{gobble}

\begin{titlepage}
\centering
\vspace*{0.75cm}
{\Huge\bfseries The Alaoglu--Erd\H{o}s\par}
\vspace{0.2cm}
{\Huge\bfseries Integer-Exponent Conjecture\par}
\vspace{0.65cm}
{\LARGE Machine-Checked Partial Results, Exact Conditional Structure,\par}
\vspace{0.12cm}
{\LARGE Determinant Methods, and the Remaining Barrier\par}
\vspace{1.0cm}
{\large A unified report on the corpus in\par}
\vspace{0.12cm}
{\large \url{https://github.com/VladimirReshetnikov/ProveIt}\par}
\vfill
\begin{minipage}{0.9\textwidth}
\small
\textbf{Bottom line.}
The conjecture --- that $2^x,3^x\in\Z$ forces $x\in\Z$ --- remains open, and the reduction to
the four exponentials conjecture explains why no routine combination of present methods can
close it.  What the corpus does establish is unusually complete.  Kernel-verified in Lean: a
zero-free region $2^x<4096$; transcendence of every hypothetical counterexample; multiplicative
rank at most two with a canonical primitive odd core; and --- conditional on a single
counterexample --- the \emph{exact} classification
\[
  \Sol=\{n+k\beta:n,k\in\Nzero\},
  \qquad
  \Cand=\{2^nM^k:n,k\in\Nzero\},
\]
with unique coordinates, exact dyadic block counts, a sharp quadratic cumulative asymptotic,
and vanishing block Fourier modes.  These reduce the problem to one narrow statement: no
primitive odd $w>1$ has a positive power of $w^{\thetaq}$ in $\Zthree$.

This report adds, at paper level, two determinant methods --- an arbitrary-node Vandermonde
repulsion inequality and a generalized-Vandermonde hierarchy over the mixed exponent semigroup
$\Nzero+\Nzero\thetaq$, the latter giving an independent explicit $O((\log X)^2)$ dyadic
bound --- together with a diagnosis of exactly why they stop short: archimedean interpolation
needs $\asymp(\log X)^2$ nodes at scale $X$, while a counterexample supplies only
$\asymp\log X$.  It also records a new sufficient condition of a different kind: if the
products $\log p\log q$ of two prime logarithms are $\Q$-linearly independent, the conjecture
follows.  The independent continuations reconciled from \lean{Article/new-reports/} add a
kernel-verified
restricted matrix-coefficient equivalence and an exact sparse-interpolation obstruction:
$R$ positive points on a hypothetical irrational power curve require exactly $R+1$
monomials.  Its Newton-polygon and boundary-defect refinements isolate a concrete
one-monomial saving problem, but do not solve it.  The other independent continuations add
exact rational-orbit collision counts, a universal superfactorial determinant divisor, a
signed-progression denominator-tax identity, and a random-phase/Fej\'er control experiment.
The seventh and eighth independent reports add, respectively, compact/profinite limit laws
with exact adelic drift, and an arithmetic audit of polynomial and Pad\'e approximation:
denominator clearing, an all-degree normalization-tax obstruction, and rational contact
rigidity.  Reports 9--13 add de-integralized common divisors and natural boundaries, a fixed
integral-matrix and Loewner formulation, finite-place/Fermat-quotient transversality, and a
microscopic prime-switch analysis, followed by an exterior-power and boundary-displacement
matrix audit.  Reports 14--20 add, independently, a structural-prime boundary-resultant
package, an information-theoretic audit of the synchronized semigroup, a special-function
transfer and mixed-Mahler audit, and a weighted Smith-profile/Loewner hierarchy.  They include
an exact ordinary/power-sum decoder, algebraic-decoder degree growth, a generic algebraic-base
spectrum, a finite mixed-Mahler identity, multilevel modular divisors, and exact
geometric-Vandermonde Smith factors, and a Mellin--spectral/approximation audit of the exact
conditional monoid.  The last two reports add an arithmetic approximation/minimax audit with
an exact geometric Newton shadow, and an operator-theoretic classification of which local,
finite-lattice, affine, first-jet, and finite-realization extensions would actually add new
information.  Each sharpens a finite interface while leaving the same
four-exponentials barrier intact.
\end{minipage}
\vfill
{\large August 21, 2026\par}
\end{titlepage}

\pagenumbering{roman}

\begin{abstract}
Let $\thetaq=\log_2 3$.  The Alaoglu--Erd\H{o}s conjecture asserts that a real $x$ with
$2^x,3^x\in\Z$ must be an integer; equivalently, the only positive integers $m$ with
$m^{\thetaq}\in\Z$ are the powers of two.  It is the first unresolved $2\times2$ instance of
the four exponentials problem, while its three-base analogue follows from six exponentials and
is formalized here by an interpolation determinant.

This unified report has five purposes.  First, it catalogues the machine-checked corpus with
exact Lean identifiers: zero-free regions, the Gelfond--Schneider dichotomy, multiplicative
rank two, the canonical primitive odd core, the rational-power index and radical degree,
localized-radical equivalences, arbitrary-order arithmetic-progression obstructions,
shrinking-target Diophantine bounds, fractional-part gap equivalences, and the exact
conditional structure theory --- the free monoid $\Sol=\Nzero+\Nzero\beta$ with a unique least
generator, exact block counts, and the sharp quadratic asymptotic.  Second, it develops two new
determinant methods.  Arbitrary candidate tuples on the lattice curve $y=x^{\thetaq}$ satisfy
\[
 \prod_{0\le i<j\le r}(m_j-m_i)\ \ge\ \frac{r!}{|\fall{\thetaq}{r}|}\,m_0^{\,r-\thetaq},
\]
with no arithmetic-progression hypothesis; and the mixed exponent semigroup supplies positive
integer generalized Vandermonde determinants whose repulsion exponent tends to one.  The
paper proof yields the explicit unconditional bound
$\#(\Cand\cap[X,2X])\le10000\max\{1,\log X\}^2$ by a route independent of the
prime-valuation rank theorem; the corresponding accepted Lean endpoint carries the explicit
hypothesis \lean{NewtonFactorization}.  On the nonarchimedean side, the accepted corpus now
identifies the exact Beatty minimizing fiber and gap, groups every tropical layer by row
partitions, truncates the normalized determinant exactly modulo $p^K$, bounds the number of
surviving row partitions by a stars-and-bars coefficient, and extracts finite
first-fiber divisors.  In the block-scaled consecutive-power model, the all-layer unit gain
adds directly to the tropical minimum, giving $p^{\tau+U}\mid\det$ uniformly across the
assignment cascade.  These gains are lower order; exact synthetic
experiments expose genuine cross-layer carrying but do not exhibit a leading-order cascade at
the tested sizes.  Third, it reconciles the genuinely new material independently developed
in all twenty reports under \lean{Article/new-reports/}.  The reports share this unified
document as their common base; they are audited independently and reconciled by mathematical
content rather than treated as a sequential chain.  Reports 1--2 show that the
Alaoglu--Erd\H{o}s conjecture is equivalent to a
restricted $2\times2$ Matrix Coefficient principle, while an irrational power curve through
$R$ positive nodes has exact sparse interpolation threshold $R+1$.  The corresponding formal
prime-symbol determinant, Newton-polygon calibration, and boundary-dilation identities expose
an exact factor-four auxiliary-polynomial gap and a finite one-monomial target.  Reports 3--5
add the exact rational-orbit collision spectrum, universal factorial and Smith divisors,
adelic determinant accounting, and a signed-progression denominator tax.  Report 6 adds a
uniform metric estimate, exact random-phase extinction, and a positive-kernel Fourier target
for the exceptional phase zero.  Report 7 adds triangle--Haar and profinite Beatty laws,
Sturmian covariance, exact denominator drift and $p$-adic escape, and a fixed-compact-factor
no-go theorem.  Report 8 adds exact denominator-clearing identities for rational
approximation, a finite-difference normalization-tax obstruction valid in every degree, and
the rational contact-equals-multiplicity principle for irrational power curves.  Report 9
adds an exact de-integralized common-zero lattice, a P\'olya--Carlson boundary, and a compact
$\Theta(\log H)$ height dichotomy.  Report 10 gives a fixed $2\times2$ integral-power
reformulation and denominator-sensitive Loewner certificates, together with their optimized
one-half barrier.  Report 11 proves structural-prime $p$-adic transversality and isolates the
Fermat-quotient first layer and adelic error wedge.  Report 12 returns to the historical
colossally abundant prime-switch problem and shows that simultaneous switches require a
microscopic prescribed-prime resonance far below modern short-interval scales.
Report 13 develops the matrix direction independently: projective arithmetic density,
determinantal-divisor and exterior-power ceilings, and an exact rank-$K$ boundary compression
of a $K^2$-point spectrum.  Report 14 develops the finite-place direction independently:
minimal interpolation denominators, structural-prime residual scaling, Newton polygons, and
odd-degree dilation resultants in two favorable residue branches.  Report 15 develops an
independent information-theoretic direction: common digit-length and Sturmian channels,
fixed-rank $S$-unit sum entropy, exact synchronized-sum decoding, algebraic-decoder degree
growth, and a finite Fourier formulation of tied-term cancellation.
Report 16 independently tests gamma, Hurwitz, polylogarithmic, Mahler, Barnes, double-sine,
and $q$-digamma transfer mechanisms.  It promotes the algebraic-base spectrum to arbitrary
independent positive algebraic pairs and isolates a corrected double-pole resonance without
confusing resonance with integrality.  Report 17 independently replaces the determinant
endpoint by a weighted Smith profile at every exterior degree, derives a universal
multilevel modular divisor and an exact geometric-Vandermonde Smith form, and couples this
arithmetic profile to rational cross-Loewner singular scales.  Its entropy analysis proves
that a hypothetical semigroup is already extremal for ordinary anti-concentration and
Shannon inequalities.
Report 18 independently reads the exact monoid as a diagonal spectrum: rank one and rank two
become simple and double spectral-zeta poles, linear and quadratic logarithmic Weyl laws, and
one or two Mellin frequency lattices.  It also proves a sharp fixed-window polynomial
node-budget barrier and isolates an exact two-scale support equation, while showing that none
of these formal rank signatures uses the decisive extra integrality $M^{\thetaq}\in\N$.
Report 19 independently develops the corresponding arithmetic approximation lattice: an
exact unit-locking principle, the same optimized fixed-window capacity, finite dyadic
interpolation certificates, a cofactor formula for discrete minimax error, and a geometric
Newton expansion whose first omitted value is explicit.  Report 20 independently audits the
operator interface: two-mode shifts, heat data, and semisimple functional calculus are blind,
whereas a global or finite valuation lattice, one affine coherence, one rational first jet,
or one exact finite trace/Hankel realization would force integrality.
Fourth, it
quantifies the remaining gap: the determinant hierarchy misses the conditional population by
exactly one factor of $\log X$.
Fifth, it records new equivalent and sufficient conditions --- block-growth dichotomies
(including the observation that the conjecture is equivalent to the dyadic block count being
$1$ for \emph{infinitely many} exponents), a compact synchronized $S$-integral orbit
formulation, and the prime-log-product independence criterion --- and closes with a
prioritized research program and a formalization blueprint.
\end{abstract}

\tableofcontents
\clearpage
\pagenumbering{arabic}
\section{Executive summary}
\label{sec:executive}

\subsection{What is proved, and at what level of trust}
\label{ex:trust}

The full conjecture remains open.  Several routes were pushed to their exact failure points:
the three-base determinant method does not shrink to two bases without proving a case of the
four exponentials conjecture; equidistribution of multiples does not defeat exponentially
shrinking integrality gaps; local curvature obstructions cannot reach the sparse global
structure; rational approximation to a hypothetical least generator stalls at the same
transcendence input; and every successful reduction terminates at the same
transcendence-theoretic wall, isolated in Section~\ref{sec:radical}.  Two lines of attack
opened after that survey -- interpolation determinants on arbitrary candidate nodes, and
generalized Vandermonde determinants over the full semigroup of integral mixed monomials --
do produce new unconditional partial results, and a third, the prime-log-product criterion of
Section~\ref{sec:logproduct}, isolates a clean sufficient condition of a different flavor.
None of them closes the final gap.

To keep unlike kinds of evidence separate, every important statement in this report carries
one of five labels.

\begin{description}[leftmargin=3.9cm,style=nextline]
\item[\statuslean]
A proof source is present in the repository, lies on the audited
\lean{ExponentialIdentities.lean} import surface, and has been accepted by the Lean kernel in
direct or full builds of the cited modules (toolchain 4.31/4.32 series, Mathlib pinned by the
repository manifest).  No \lean{sorry}, no extra axioms, no \lean{native_decide}.
\item[\statusnew]
A conventional mathematical proof is supplied in this document.  It is intended to be
suitable for formalization, but no claim of kernel checking is made.
\item[\statusdraft]
A Lean source for the statement exists in the repository tree, but it has \emph{not} yet been
accepted by the kernel at the pinned state.  A draft is weaker evidence than \statusnew{}
prose backed by a written proof, not stronger: it records intent to formalize, nothing more.
\item[\statuscomp]
A finite statement was checked by an explicitly identified software computation, usually
with outward rounding.  The local section records the script, parameters, margins, artifacts,
and relevant trust base: this includes CPython and \lean{mpmath} for the interval checks and
C, MPFR, and GMP for the directed-rounding scans.  Such evidence is not interchangeable with
a Lean certificate.
\item[\statusknown]
A classical theorem quoted from the literature.
\end{description}

An earlier audit, performed against \commit{423ac24}, correctly observed that three
modules cited by the survey --- \lean{FiniteCheck4096}, \lean{OddCore}
and \lean{FractionalPartDensity} --- were absent then.  That discrepancy is
resolved: the three modules were subsequently committed in \commit{98646b455} and are
kernel-verified; the survey's temporary pending-verification remark was removed in
\commit{716617d72}.  Statements from these modules are therefore labeled \statuslean{} below.
The kernel-checked feature wave also includes
\begin{center}
\lean{PrimitiveOddCore}, \lean{LeastSolution}, \lean{ThreeDenominatorNormalization},\\
\lean{RationalPowerIndex}, \lean{PrimitiveGenerator}, \lean{ExactCounting},\\
\lean{QuadraticAsymptotic}, \lean{BlockWeyl}, \lean{BlockEquidistribution},\\
\lean{OutputNormalization}, \lean{RationalThirdBase}, \lean{RadicalReformulation},\\
\lean{PrimitiveRadical}, \lean{RadicalDegree}, \lean{SharperProgressions},\\
\lean{ThirdDifference}, \lean{HigherDifference}.
\end{center}
The determinant and reformulation modules that earlier editions of this report listed as
drafts have since been accepted, and a further round added
\begin{center}
\lean{CandidateLattice}, \lean{KernelDichotomy}, \lean{SecondIterateKernel},\\
\lean{TowerRigidity}, \lean{SymmetricTowerConstants}, \lean{PrimeLogSpaceOperator},\\
\lean{ArbitraryNodeRepulsion}, \lean{SemigroupGapBound}.
\end{center}
The current algebraic-locus and nonarchimedean wave further includes
\begin{center}
\lean{AlgebraicExponentLocus}, \lean{AlgebraicOutputLocus},\\
\lean{RationalFunctionAlgebraicLocus}, \lean{AlgebraicTowerTransition},\\
\lean{CrossBaseTowerTransition}, \lean{BlockVandermondeValuationBudget},
\lean{AllLayerUnitBasis}, \lean{UnitOrderingGain},\\
\lean{SimultaneousRadicalDegrees}, \lean{CanonicalRadicalMixedOutput},\\
\lean{TropicalInitialForm}, \lean{BlockFiberFactorization},
\lean{TropicalBlockApplication}, \lean{BeattyTropicalFiber},\\
\lean{RowBlockLaplaceExpansion}, \lean{AssignmentExcessGeometry},
\lean{RowBlockUnitDivisibility}, \lean{MixedDeterminantUnitBridge},\\
\lean{RestrictedMatrixCoefficient}, \lean{FormalPrimeDeterminant},
\lean{SparsePowerCurve}, \lean{RationalOrbitGrid}, \lean{BoundaryDilationDefect},
\lean{SuperfactorialDeterminant}, \lean{SignedProgressionDeterminant},
\lean{FejerKernelCertificate}, \lean{GlobalMixedPrefixCascade},\\
\lean{DilationDifferenceQuotient}, \lean{SynchronizedSumDecoder},
\lean{RationalSumDecoder}, \lean{AlgebraicBaseSpectrum},
\lean{MixedMahlerFinite}, \lean{GeometricVandermondeSmith},
\lean{TwoScaleSupportFinite}, \lean{GeometricNewtonInterpolation},
\lean{JordanFirstJet}.
\end{center}
Public headline theorems throughout were audited with \lean{#print axioms}; only the standard
Mathlib axioms \lean{propext}, \lean{Classical.choice}, and \lean{Quot.sound} occur, and the
audit for the most recent round is itself a module, \lean{AuditNewModules}.  The two analytic
inputs still missing from the determinant chain are carried as named hypotheses of the
theorems that use them, not as axioms.  The per-statement ledger is Appendix~\ref{app:status}
and the module inventory is Appendix~\ref{app:inventory}.

\subsection{Headline results}
\label{ex:headline}

\begin{enumerate}
\item \textbf{Verified zero-free region.}  If $2^x,3^x\in\Z$ and $2^x<4096$, then $x\in\Z$;
equivalently every nonintegral solution has $x\ge12$.  \statuslean{}  A directed-rounding
scan extends the finite barrier to $2^x<2^{27}$, i.e.\ $x>27$.  \statuscomp

\item \textbf{Transcendence.}  Every nonintegral solution is transcendental; for $x$ with
$2^x\in\Z$, algebraicity of $x$ is \emph{equivalent} to integrality.  $\thetaq=\log_23$ is
transcendental.  \statuslean

\item \textbf{Exact conditional structure.}  Conditional on a counterexample there is a unique
gcd-normalized odd core $w>1$, every candidate has unique coordinates $2^iw^j$, and the
solution set has the exact free-monoid classification $\Sol=\Nzero+\Nzero\beta$ with a unique
least nonintegral generator $\beta$, which is affinely primitive.  Candidate counts below $N$
are at most $\clog_2N\cdot\clog_3N$, each dyadic block $[2^T,2^{T+1})$ holds at most
$\clog_3(2^{T+1})$ candidates, the exact number of \emph{nonintegral solutions} in the
exponent block $[N,N+1)$ is $\lfloor(N+1)/\beta\rfloor$, while the full solution count ---
equivalently the candidate count in $[2^N,2^{N+1})$ --- is
$1+\lfloor(N+1)/\beta\rfloor$, the
cumulative count satisfies $C(T)/T^2\to1/(2\beta)$, and the conjecture is equivalent to
subquadratic growth along $2^T$.  Exact block averages of every finite trigonometric
polynomial converge to its constant Fourier coefficient.  \statuslean{}  The extension to all
continuous tests and interval indicators remains paper-level.  \statusnew

\item \textbf{Unconditional gap equivalences.}  The conjecture is equivalent to the existence
of any nonempty open subinterval of $(0,1)$ free of fractional parts $\fract{x}$ of
nonintegral solutions, and to those fractional parts being bounded away from $0$, or from
$1$; the convergent denominators of $\thetaq$ obey $q_{n+1}<2m^{q_n}$.  \statuslean

\item \textbf{Sharper local combinatorics.}  For every $r\ge2$, the sharp criterion
$|\fall{\thetaq}{r}|\,h^r<n^{r-\thetaq}$ forbids $r+1$ candidates in arithmetic progression,
uniformly in the progression length.  The underlying arbitrary-order forward-difference
mean-value identity is formalized as well, and for the three-term case the rational effective
criterion $h^{400}\le n^{83}$ (step exponent $83/400=0.2075$) is verified independently.
\statuslean

\item \textbf{New: arbitrary-node repulsion and the semigroup hierarchy.}  Any $r+1$ distinct
natural candidates -- with no arithmetic-progression hypothesis -- obey a quantitative
product-of-gaps lower bound, giving explicit local packing bounds and a best ordinary-power
exponent $0.241503\ldots$ at determinant order $4$ (Section~\ref{sec:arbnode}).  Enlarging the
exponent set to the mixed semigroup $\Nzero+\Nzero\thetaq$ yields arbitrarily large positive
integer generalized Vandermonde determinants, a near-linear repulsion exponent
$1-O(r^{-1/2})$, and the paper-level unconditional dyadic bound
$\#(\Cand\cap[X,2X])\le10000\max\{1,\log X\}^2$ (Section~\ref{sec:semigroup}).  This route is
independent of the repository's prime-valuation rank theorem.  The paper theorems are
unconditional.  Their current Lean forms isolate the two missing analytic ingredients as the
explicit hypotheses \lean{RpowDividedDifferenceMeanValue} and \lean{NewtonFactorization};
all determinant identities, integrality, nonvanishing, positivity, and the deductions from
those hypotheses are kernel-accepted.  \statusnew{} \statuslean{}  Lean modules
\lean{ArbitraryNodeDeterminant}, \lean{MixedExponentSemigroup}, \lean{SemigroupDeterminant},
\lean{FallingFactorialBound}, \lean{SemigroupWeightBound},
\lean{ArbitraryNodeRepulsion}, and \lean{SemigroupGapBound} are on the aggregate import
surface.

\item \textbf{New: a sufficient condition from products of prime logarithms.}  Suppose
$2^x=M\in\Z$ and $3^x=A\in\Z$.  Eliminating $x$ gives the exact identity
\[
  \log M\,\log3=\log A\,\log2 .
\]
Expanding $A$ and $M$ into prime factorizations turns this into an integer-coefficient linear
form in the pairwise products $\log p\,\log q$ of prime logarithms, and the form is nontrivial
precisely when $x\notin\Z$.  Hence: \emph{if the products $\log p\,\log q$, over unordered
pairs of primes, are linearly independent over $\Q$, the conjecture holds}
(Section~\ref{sec:logproduct}).  That independence is a consequence of Schanuel's conjecture.
It is different in kind from the four exponentials conjecture -- it constrains a bilinear
lattice of logarithm products rather than a $2\times2$ logarithm determinant -- so it supplies
a genuinely separate route to the same conclusion, not a reformulation of the same wall.
\statusnew{} \statuslean{}  The exact conditional implication is accepted in
\lean{LogProductIndependence}; the asserted prime-log-product independence itself remains a
classical hypothesis, not a theorem of the corpus.

\item \textbf{Exact localization of the $p$-adic cascade.}  For the mixed Beatty determinant,
the minimum assignments, their floor-sized gap, and the first block-determinant quotient are
now exact.  More generally, after removing $p^\tau$, reduction modulo $p^K$ can receive
contributions only from row partitions of assignment excess below $K$, with every coefficient a signed
product of square unit minors.  The first-fiber valuation has an elementary finite cube budget,
and monic unit-adapted bases extract additional whole-Vandermonde divisors.  For the
block-scaled consecutive-power model, their common exponent $U$ is independent of the row
partition and the exact additive divisor $p^{\tau+U}\mid\det$ follows.  These results
also bound the number of terms below any fixed excess cutoff by an explicit binomial
coefficient.  They localize every possible leading excess to the valuations of, and cancellation among, the
surviving low-excess minor products; they bound neither phenomenon.  \statuslean

\item \textbf{New: restricted matrix coefficients and exact sparse interpolation.}
For positive integral outputs $M,A$, the singular logarithm matrix
$\bigl[\begin{smallmatrix}\log2&\log M\\\log3&\log A\end{smallmatrix}\bigr]$
has a zero coefficient between nonzero rational vectors exactly as required by the restricted
Matrix Coefficient principle if and only if the Alaoglu--Erd\H{o}s conjecture holds.
Moreover, on an irrational curve $Y=X^\beta$, vanishing at $R$ distinct positive nodes
requires at least $R+1$ monomials, and a positive-node interpolant attains this bound with
every coefficient nonzero.  The lower bound, integral evaluation determinant, matching
interpolant, and matrix-coefficient equivalence are kernel verified.  The induced BK-degree
threshold $R-1$ and boundary compression leave a precise one-monomial saving target.
\statuslean{} \statusnew

\item \textbf{The irreducible remainder.}  The conjecture is equivalent to: no odd integer
$u>1$ satisfies $u^{\thetaq}\in\Zthree$, with a symmetric base-swapped formulation and an
equivalent reduction to canonical power-primitive odd bases.  The least radical index gives an
irreducible binomial and the exact algebraic degree.  \statuslean{}  The exclusion itself
remains open even for $u=3$.
\end{enumerate}

\subsection{The strongest new conclusions}
\label{ex:newconclusions}

\begin{resultbox}
\textbf{Determinant hierarchy (Sections~\ref{sec:arbnode} and~\ref{sec:semigroup}).}
For every $r\ge2$, any $r+1$ distinct natural candidates obey a quantitative product-of-gaps
lower bound, with no arithmetic-progression hypothesis.  Enlarging the exponent set from
ordinary powers to the mixed semigroup $\Nzero+\Nzero\thetaq$ supplies arbitrarily large
positive integer generalized Vandermonde determinants; their archimedean upper bounds give a
near-linear repulsion exponent $1-O(r^{-1/2})$ and, unconditionally,
\[
  \#(\Cand\cap[X,2X])\ll(\log X)^2,
\]
with the explicit constant $10000$ proved below.  Paper proof in this report; a Lean
formalization is accepted with its single Newton-factorization input visible as a theorem
hypothesis.  \statusnew{} \statuslean
\end{resultbox}

\begin{resultbox}
\textbf{The one-logarithm deficit (Section~\ref{sec:deficit}).}
To make the semigroup determinant nearly force a gap comparable with $X$, one needs
$r\asymp(\log X)^2$ candidate nodes.  Conditional on a counterexample, the exact ProveIt
classification supplies only $\asymp\log X$ candidates in $[X,2X]$.  Purely archimedean
interpolation, at its present strength, therefore misses the conditional population by exactly
one power of $\log X$.  A successful refinement must either reduce the determinant's node
requirement to $O(\!\log X)$ or extract additional nonarchimedean size from the same nodes.
\statusnew
\end{resultbox}

\begin{resultbox}
\textbf{Matrix coefficients and sparse interpolation
(Section~\ref{sec:matrix-continuation}).}
The conjecture is exactly a restricted $2\times2$ Matrix Coefficient principle.  Independently,
every nonzero relation supported on $m$ monomials restricts to an exponential polynomial with
at most $m-1$ positive-curve zeros; the matching positive-node interpolant has all $R+1$
coefficients nonzero.  Thus $R$ counterexample points have exact sparse interpolation
complexity $R+1$, so saving even one monomial over ordinary interpolation would rule them out.
The matrix equivalence, sparse lower bound, nonzero integral determinant, and matching
interpolant are kernel checked; the BK-area calibration and resulting factor-four comparison
with the general Matrix Coefficient target remain paper-level.  \statuslean{} \statusnew
\end{resultbox}

\subsection{Why this is not a proof}
\label{ex:notproof}

The conjecture is already a special integral instance of the real four exponentials
conjecture.  In the natural $2\times2$ matrix
\[
  \begin{pmatrix}
   1\cdot\log2 & 1\cdot\log3\\
   x\log2 & x\log3
  \end{pmatrix},
\]
the four exponentials are $2,3,2^x,3^x$.  For a nonintegral solution, both the row pair and
the column pair are rationally linearly independent, yet all four exponentials are algebraic.
The conclusion needed here is therefore exactly what four exponentials predicts and what the
known six exponentials theorem does not supply in a $2\times2$ configuration.

The additional integrality hypotheses are stronger than mere algebraicity, and they do yield
the new determinant inequalities of Sections~\ref{sec:arbnode} and~\ref{sec:semigroup} and the
bilinear criterion of Section~\ref{sec:logproduct}.  They do not, by themselves, produce
enough nodes, enough lower-bound divisibility, or the required independence of logarithm
products to cross the four-exponentials threshold.  Conjecture~\ref{conj:AE} therefore remains
marked open throughout this report.
\section{The problem and its equivalent forms}
\label{sec:problem}

\subsection{Original statement}

\begin{conjecture}[Alaoglu--Erd\H{o}s, two-base integer form]
\label{conj:AE}
For every $x\in\R$,
\[
   2^x\in\Z\quad\text{and}\quad 3^x\in\Z
   \qquad\Longrightarrow\qquad x\in\Z.
\]
\end{conjecture}

\begin{statusbox}{recorded in Lean as the proposition \lean{AlaogluErdosConjecture}, without
being asserted.}
\end{statusbox}

Since $2^x$ and $3^x$ are positive, the hypotheses actually put them in $\N$; and $2^x\in\N$
forces $x\ge0$, since a negative exponent gives $0<2^x<1$.  All solutions therefore live in
$[0,\infty)$ (\lean{nonneg_of_two_rpow_integer}), and ``integer'' may be replaced throughout
by ``positive natural number.''

\paragraph{Historical placement.}
Questions of this type go back to the 1944 study of colossally abundant numbers by Alaoglu
and Erd\H{o}s \cite{AE1944}, where the arithmetic nature of simultaneous rational powers
$p^x,q^x$ of distinct primes controls the quotients of consecutive colossally abundant
numbers --- a theme reaching back to Ramanujan's superior highly composite numbers
\cite{Ram1915}.  Alaoglu and Erd\H{o}s isolated the two-prime assertion explicitly, regarded
it as very likely, could not prove it, and reported a communicated three-prime result of
Siegel.  In modern language Siegel's case is supplied by the six exponentials theorem,
whereas the two-base statement sits at the unresolved $2\times2$ threshold of the four
exponentials conjecture (Schneider's first problem \cite{Schneider1957}; see Waldschmidt
\cite{Waldschmidt2000,Waldschmidt2023}).  The elementary-looking integrality question is thus
one of the canonical concrete tests of the gap between six and four exponentials, not an
isolated recreational puzzle.  As of the dated literature check for this report, the four
exponentials conjecture remains open.

\subsection{Candidate-base normalization}

Put
\[
  \thetaq:=\log_2 3=\log_2 3
  =1.584962500721156\ldots
\]
and define the \emph{candidate set}
\[
 \Cand:=\{m\in\N:m>0,\ m^{\thetaq}\in\N\}.
\]

\begin{lemma}[Exponent--candidate bijection]
\label{pr:bijection}
The maps $x\mapsto 2^x$ and $m\mapsto\log_2m$ are mutually inverse bijections between
$\Sol:=\{x\in\R:2^x,3^x\in\Z\}$ and $\Cand$.  Under the bijection, $3^x=m^{\thetaq}$.
\end{lemma}

\begin{proof}
If $x\in\Sol$, set $m=2^x\in\N$.  Then
$m^{\thetaq}=\exp(\log m\cdot\log_2 3)=\exp(x\log3)=3^x\in\N$.  Conversely, if
$m\in\Cand$ and $x=\log_2m$, then $2^x=m$ and the same identity gives
$3^x=m^{\thetaq}\in\N$.
\end{proof}

\begin{statusbox}{\statuslean{} ---
\lean{exists_twoBaseNaturalCandidate_of_two_three_rpow_integer},
\lean{TwoBaseNaturalCandidate.integer_powers}.}
\end{statusbox}

Integral solutions correspond exactly to powers of two, by the identifier
\begin{center}
\lean{twoBaseCandidateExponent_isInteger_iff_isPowerOfTwo},
\end{center}
so
\begin{equation}
 \boxed{\quad\text{Conjecture~\ref{conj:AE}}
 \iff \Cand=\{2^n:n\in\Nzero\}.\quad}
 \label{pr:candidate-equivalence}
\end{equation}
This equivalence is itself formalized, as
\begin{center}
\lean{alaogluErdosConjecture_iff_candidates_are_powers_of_two}.
\end{center}
The candidate formulation
is well suited to unique factorization, finite certification, finite differences of
$t^{\thetaq}$, multiplicative-rank arguments, and the determinant constructions of this
report, and it is the formulation in which most of the corpus is stated.

Two closure properties of $\Sol$ are elementary but essential: $\Sol$ is closed under
addition and under multiplication by positive naturals, so it is an additive submonoid of
$[0,\infty)$ containing $\N$ (\lean{rpow_integer_add}, \lean{rpow_integer_nat_mul}).

\subsection{The logarithmic determinant}

Writing
\[
  M=2^x,\qquad A=3^x,
\]
a solution is equivalently a positive integer pair satisfying
\begin{equation}
  \log M\,\log3-\log A\,\log2=0,
  \label{pr:log-det}
\end{equation}
and the conjecture says that all positive integer solutions of \eqref{pr:log-det} are
\[
  (M,A)=(2^n,3^n),\qquad n\in\Nzero.
\]
Thus the problem is the vanishing of a $2\times2$ determinant of logarithms of algebraic
numbers,
\[
  \det\begin{pmatrix}\log M&\log 2\\ \log A&\log 3\end{pmatrix}=0 ,
\]
with the additional information that two of the four entries are logarithms of \emph{unknown
integers}.  This is not merely an analogy with the four exponentials conjecture; it is its
exact local geometry.

\subsection{One-parameter subgroup form}

Let
\[
  \Gamma=\{(2^t,3^t):t\in\R\}\subset\R_{>0}^2 .
\]
Then Conjecture~\ref{conj:AE} is the statement
\[
  \Gamma\cap\N^2=\{(2^n,3^n):n\in\Nzero\}.
\]
The curve $\Gamma$ is a real one-parameter subgroup of the multiplicative torus
$\mathbb G_m^2(\R)$, and the assertion is that it meets the integer lattice points of the
torus only at the obvious places.  The difficulty is that its direction vector
$(\log2,\log3)$ is transcendental rather than algebraic, so standard algebraic-subgroup
intersection theorems do not apply directly.

\subsection{The four-exponentials reduction}

Suppose $x$ is nonintegral and $2^x,3^x$ are algebraic.  The rational-exponent argument
(Lemma~\ref{pr:rational}) shows that $x$ cannot be rational, so $1,x$ are $\Q$-linearly
independent; and $\log2,\log3$ are $\Q$-linearly independent since $2^a3^b=1$ forces
$a=b=0$.  The four exponentials built from these two pairs are
\[
 e^{1\cdot\log2}=2,\quad e^{1\cdot\log3}=3,\quad e^{x\log2}=2^x,\quad e^{x\log3}=3^x,
\]
all four algebraic --- contrary to the four exponentials conjecture.  This reduction is
formalized: \lean{alaogluErdosConjecture_of_realFourExponentialsConjecture}.  \statuslean

The reduction is a diagnostic, not just context: any proposed elementary argument must
exploit integrality, positivity, or arithmetic structure beyond bare algebraicity of those
four values, or else it would prove a substantial case of the four exponentials conjecture
by another name.

\subsection{Rational and algebraic exponents}

\begin{lemma}[Rational exponent]
\label{pr:rational}
If $x\in\Q$ and $2^x\in\Z$, then $x\in\Z$.  Consequently every nonintegral solution is
irrational.
\end{lemma}

\begin{proof}
Write $x=a/b$ in lowest terms with $b>0$.  If $2^{a/b}=m\in\N$, then $2^a=m^b$; comparing
the exponent of $2$ gives $a=b\,\vpp{2}{m}$, so $b\mid a$ and $b=1$.
\end{proof}

\begin{statusbox}{\statuslean{} ---
\lean{IntegerExponent.integer_of_rational_of_two_rpow_integer},
\lean{irrational_of_not_integer_of_two_rpow_integer}.}
\end{statusbox}

Gelfond--Schneider upgrades ``irrational'' to ``transcendental'' and removes the algebraic
irrational case entirely.

\begin{lemma}[Transcendence dichotomy]
\label{pr:dichotomy}
If $2^x\in\Z$, then either $x\in\Z$ or $x$ is transcendental.  In particular every
counterexample to Conjecture~\ref{conj:AE} is a transcendental number, and
$\thetaq=\log_2 3$ is itself transcendental.
\end{lemma}

\begin{statusbox}{\statuslean{} ---
\lean{transcendental_of_not_integer_of_two_three_rpow_integer},
\lean{transcendental_logThreeDivLogTwo}; the supporting Gelfond--Schneider material is
discussed in Section~\ref{sec:corpus}.}
\end{statusbox}

Together with the finite verification these lemmas already remove several whole regimes.  In
summary form, the pinned corpus proves the following unconditional restrictions on a
hypothetical counterexample $x$.

\begin{itemize}
\item $x$ is irrational, and in fact transcendental.  \statuslean
\item $\thetaq$ is transcendental.  \statuslean
\item $x\ge12$, by the kernel-checked zero-free region $2^x<4096$.  \statuslean
\item $2^x\ge2^{27}$, hence $x\ge27$, by the directed-rounding scan; this raises the
barrier only under the stated software trust boundary and is not a kernel certificate.
\statuscomp
\item The real four exponentials conjecture implies Conjecture~\ref{conj:AE}.  \statuslean
\end{itemize}

These facts eliminate the rational, algebraic-irrational, small, and routine
three-exponentials cases.  They also explain why an elementary manipulation of logarithms is
unlikely to finish the proof: the unknown lies exactly where present transcendence theory
changes from theorem to conjecture.

\subsection{The three-base theorem and its determinant architecture}

By contrast with the two-base problem, three multiplicatively independent bases fall to the
six exponentials mechanism, and the special case needed here is fully formalized.

\begin{theorem}[Three-base theorem]
\label{pr:threebase}
If $2^x,3^x,5^x\in\Z$ for a real $x$, then $x\in\Z$.
\end{theorem}

\begin{statusbox}{\statuslean{} --- modules \lean{IntegerExponent},
\lean{IntegerExponent.DeterminantBound}, \lean{IntegerExponent.ExponentialKernel},
\lean{IntegerExponent.Grid}, following Waldschmidt's square-determinant proof
\cite{WaldschmidtSix}.}
\end{statusbox}

In outline, for irrational $x$ one forms a large square matrix
$\mathcal A_{ij}=\exp(a_ib_j)$, where the row arguments $a_i$ encode monomials in $2,3,5$ and
the column arguments $b_j$ encode affine forms in $1,x$.  Irrationality makes the column
parameters pairwise distinct and multiplicative independence makes the row parameters
pairwise distinct; a total-positivity argument gives $\det\mathcal A\ne0$; integrality of all
entries forces $|\det\mathcal A|\ge1$; and an interpolation estimate with explicit size
choices gives $|\det\mathcal A|<1$, a contradiction.

The dimension count is exactly where the argument stops short of Conjecture~\ref{conj:AE}.
With three independent row directions and two column directions the analytic upper bound
decays faster than the arithmetic lower bound $1$ permits, and the contradiction closes.
Shrinking to two bases removes one row direction, and the same balance then lands precisely
on the classical boundary where the decay below one fails.  Nothing in the proof is wasteful
enough to be recovered by sharper constants: a successful two-base determinant of this shape
would be a new four-exponentials argument, which is why Section~\ref{sec:arbnode} and
Section~\ref{sec:semigroup} enlarge the \emph{node set} rather than tune the estimate.

\paragraph{A denominator-controlled rational third output.}
The rational-output extension now formalized is deliberately conditional.  Let $a>1$ and
assume the monomial map
\[
 (i,j,k)\longmapsto 2^i3^ja^k
\]
is injective.  If $2^x=m_2$ and $3^x=m_3$ are natural numbers, $a^x=q\in\Q$, and for some
$u,v\in\Nzero$ the reduced denominator satisfies
\[
 \operatorname{den}(q)\mid m_2^u m_3^v,
\]
then $x\in\Q$, and in fact $x\in\Z$ because $2^x$ is integral.  A concrete corollary obtains
the monomial injectivity when $a$ has a prime divisor different from $2$ and $3$.

On the determinant side, a nonzero $m\times m$ real determinant whose entries become integers
after multiplication by a common positive denominator $Q$ satisfies
$|\det A|\ge Q^{-m}$.  The explicit analytic estimate also retains a denominator margin:
for $m=2r$, $r>2$, $192C\le r$, and $bs+2r<r^2$, the verified upper-bound expression remains
strictly below $1$ even after multiplication by $b^s$.

\begin{statusbox}{\statuslean{} ---
\lean{one_div_pow_le_abs_det_of_common_denominator},
\lean{exponential_det_numeric_bound_with_denominator},
\lean{rational_of_two_three_a_rpow_rational_of_den_dvd_outputs}, and
\lean{integer_of_two_three_a_rpow_rational_of_prime_factor_of_den_dvd_outputs} in
\lean{RationalThirdBase}.  These are denominator-divisibility and explicit-margin results;
the unrestricted determinant grid for an arbitrary rational third-base output has
\emph{not} been completed.}
\end{statusbox}

\section{Repository timeline and verification status}
\label{sec:timeline}

The formal development advanced in three waves of Lean work, interleaved with four
documentation milestones.  All commits are in the repository \cite{ProveIt}; the
mathematical role of each is summarized below.

\begin{longtable}{@{}p{0.13\textwidth}p{0.80\textwidth}@{}}
\toprule
Commit & Role \\
\midrule
\endfirsthead
\toprule
Commit & Role \\
\midrule
\endhead
\commit{f492c18f1} & Three-base theorem via interpolation determinants
(\lean{IntegerExponent} and submodules). \\
\addlinespace
\commit{5ed4e6185} & Merge of the three-base branch. \\
\addlinespace
\commit{a37d97f9d} & Basic two-base module: the conjecture as a proposition, closure under
natural multiples, irrationality of nonintegral solutions, endpoint gaps, numerical bounds
for $\thetaq$, the $2^x<10$ exclusion, second-difference obstruction, aligned-block $2/3$
counting. \\
\addlinespace
\commit{4a9b03d12} & Major corpus: finite checks through $2^x<256$; arithmetic descent;
multiplicative rank; candidate structure; zero density; shrinking targets; the $3$-smooth
theorem; the four-exponentials reduction; transcendence consequences; and a complete
Gelfond--Schneider port. \\
\addlinespace
\commit{18a85e6b8} & Merge of the two-base branch. \\
\addlinespace
\commit{423ac24f0} & First survey document (\lean{two-base-partial-results.tex}) with
rendered PDF. \\
\addlinespace
\commit{98646b455} & Three further modules, kernel-verified:
\lean{FractionalPartDensity} (density dichotomy and gap equivalences),
\lean{OddCore} (common odd core, improved counting, dyadic blocks),
\lean{FiniteCheck4096} (zero-free region to $2^x<4096$). \\
\addlinespace
\commit{716617d72} & Survey finalized after kernel acceptance of the $4096$ region. \\
\addlinespace
\commit{11ac3093f} & First independent research report (audit at \commit{423ac24f0}, exact
conditional structure, interval scan to $2^{20}$); one of the sources merged into the present
document. \\
\addlinespace
\commit{83ef82ab1} & Merge of the two earlier documents into a single expanded report. \\
\addlinespace
\commit{46f5b67e3} & Localization, sharp arithmetic progressions and an exact generator:
\lean{Localization}, \lean{SharpProgression}, and the first form of
\lean{PrimitiveGenerator}. \\
\addlinespace
\commit{5d3982ceb} & Merge of the parallel conditional-structure and asymptotic
formalizations: canonical primitive generator and exact solution monoid, least solution,
exact dyadic counting, the quadratic cumulative asymptotic, block Weyl sums and block
equidistribution, simultaneous output normalization and power-of-three denominator
normalization, the rational-power index, the radical modules (radical reformulation,
primitive radical, exact radical degree), the higher-difference obstruction, and
denominator-controlled rational third-base outputs. \\
\addlinespace
\commit{8c9c59224} & Denominator-growth constraint
$q_{n+1}<2m^{\,q_n}$ on the continued fraction of a hypothetical solution
(\lean{DenominatorGrowth}), together with the report analysis naming the obstruction as a
vanishing integer-coefficient linear form in products of two prime logarithms. \\
\addlinespace
\commit{2ece309cc} & Second research report: arbitrary-node Vandermonde repulsion, the
semigroup generalized-Vandermonde hierarchy and its dyadic $O((\log X)^2)$ bound, the
one-logarithm deficit analysis, and the compact-orbit and block-growth reformulations; the
other source merged into the present document. \\
\bottomrule
\end{longtable}

The audit performed in \commit{11ac3093f} treated results of the then-uncommitted modules as
article-level claims; with \commit{98646b455} in the tree and kernel-verified, those results
are now cited as committed Lean.  The interval scan to $2^{20}$ remains at the
\statuscomp{} level by design.  The feature modules merged in \commit{5d3982ceb} prove the
cumulative quadratic asymptotic and the finite-trigonometric-polynomial block averages; they
do \emph{not} yet formalize the stronger continuous-test-function or interval-counting
equidistribution theorem, and that distinction is maintained throughout this report.

\subsection{Evidence convention}
\label{pr:evidence-convention}

The evidence convention of this report is source-tree plus build-level, not
rebuild-on-demand.  Every claim marked \statuslean{} was checked against the cited Lean
source and theorem statement at the repository state above, and the corresponding module is
on the aggregate import surface \lean{ExponentialIdentities.lean}, where it is recorded as
accepted by the Lean kernel.  The label does not assert that a separate fresh build was run
while preparing this document; it asserts that the theorem exists, in the stated form, on an
import surface the kernel has accepted.

Three weaker labels are used deliberately.  \statuscomp{} marks a finite interval
computation with an explicit software trust boundary and no kernel certificate;
\statusnew{} marks a conventional proof given in this report, written with formalization in
mind but not yet translated; and \statusdraft{} marks a Lean module that exists in the tree
as a draft but has \emph{not} been accepted by the kernel, and therefore carries no more
weight than \statusnew.  The complete label-by-label inventory is in
Appendix~\ref{app:status}, and the per-module inventory in Appendix~\ref{app:inventory}.
\section{The kernel-verified corpus}
\label{sec:corpus}

Throughout this section $x$ denotes a hypothetical nonintegral solution, $m=2^x\in\N$,
$n=3^x\in\N$.  Every statement here is \statuslean{} unless labeled otherwise.

\subsection{Fractional-part gaps and shrinking targets}
\label{co:gaps}

An integral power caught strictly between consecutive integral powers of its base must clear
both by at least one full unit; dividing through normalizes this to the fractional part.

\begin{theorem}[Fractional-part gaps]
\label{co:thm-gaps}
Let $x\in(N,N+1)$ with $N\in\N$, $t=x-N$, and suppose $2^x,3^x\in\Z$.  Then
\[
1+2^{-N}\le 2^{t}\le 2-2^{-N},
\qquad
1+3^{-N}\le 3^{t}\le 3-3^{-N}.
\]
\end{theorem}

\begin{statusbox}{\statuslean{} --- \lean{two_three_fractional_part_gaps},
\lean{two_three_rpow_gaps_near_nat}, from
\lean{rpow_integer_gap_between_consecutive_powers}.}
\end{statusbox}

Near the endpoints the excluded zones have width $\asymp2^{-N}$: they \emph{shrink} as the
integer part grows.  Applied to all multiples $kx$ (which are again solutions), this yields
explicit Diophantine lower bounds.

\begin{theorem}[Effective irrationality bound]
\label{co:thm-shrinking}
Let $x\notin\Z$ satisfy $2^x\in\Z$ and write $m=2^x$.  Then for all integers $p$ and all
$k\ge1$,
\[
\Bigl|x-\frac pk\Bigr|\ \ge\ \frac1{2\,k\,m^{k}} .
\]
Exact (slightly stronger) radii, and two-base versions taking the better of the base-$2$ and
base-$3$ radii, are formalized as well.
\end{theorem}

\begin{statusbox}{\statuslean{} --- \lean{rational_approximation_exponential_lower_bound},
\lean{two_three_multiple_distance_to_integer_lower_bound},
\lean{exists_natural_base_rational_approximation_lower_bound} in \lean{ShrinkingTarget}.}
\end{statusbox}

In shorthand, $\distZ{kx}\gtrsim2^{-kx}$.  This bound is compatible with ordinary Diophantine
approximation: Dirichlet guarantees infinitely many $k$ with $\distZ{kx}<1/k$, and the
shrinking radius $2^{-kx}$ is incomparably smaller.  The interplay between this exponential
gap and equidistribution is discussed in Section~\ref{sec:attempts}.

The bound dualizes into a growth constraint on the continued fraction of a solution: if
$|x-p/q|<1/(qQ)$ with $q,Q\ge1$, then $Q<2m^{\,q}$.  Applied to consecutive convergent
denominators, which satisfy $|x-p_n/q_n|<1/(q_nq_{n+1})$, this gives
\[
  q_{n+1}\;<\;2\,m^{\,q_n}:
\]
the partial quotients of a hypothetical nonintegral solution are capped by a single
exponential in the preceding denominator, with base its own output $m=2^x$.
\begin{statusbox}{\statuslean{} ---
\lean{approximation_quality_lt_of_noninteger_solution} in \lean{DenominatorGrowth}.}
\end{statusbox}

\subsection{Verified zero-free regions and certificate architecture}
\label{co:finite}

\begin{theorem}[Zero-free region]
\label{co:thm-finite4096}
If $2^x,3^x\in\Z$ and $2^x<4096$, then $x\in\Z$.  Equivalently, every nonintegral solution
satisfies $x\ge12$, and every candidate that is not a power of two is at least $4096$.
\end{theorem}

\begin{statusbox}{\statuslean{} ---
\lean{integer_of_two_three_rpow_integer_of_two_rpow_lt_4096},
\lean{twelve_le_of_not_integer_of_two_three_rpow_integer},
\lean{le_of_twoBaseNonintegerCandidate_4096} in \lean{FiniteCheck4096}, extending
\lean{FiniteCheck} ($2^x<100$) and \lean{FiniteCheck256} ($2^x<256$, hence $x\ge8$).}
\end{statusbox}

The certificate architecture is worth recording, both because it is the repository's model
of a trustworthy large computation and because scaling it is a named research direction.
For each non-power-of-two $m$ the certificate is a pair of exact integer power comparisons
against rational enclosures $L_p/L_q<\thetaq<U_p/U_q$ drawn from continued-fraction
convergents of $\thetaq$: from
\[
a^{L_q}<m^{L_p}
\quad\text{and}\quad
m^{U_p}<(a+1)^{U_q}
\]
one traps $a<m^{\thetaq}<a+1$, so $m^{\thetaq}\notin\Z$.  The $2^x<256$ stage verifies
per-value certificates by \lean{norm_num} with three enclosure tiers (exponent sizes
${\sim}500$, ${\sim}1500$, ${\sim}25000$).  The extension to $4096$ verifies $3836$
certificates through a single kernel-replayed \lean{decide} over a table generated by
PARI/GP and revalidated independently with exact big-integer arithmetic in Python.  Five
tiers suffice, with tier populations $78$, $488$, $3241$, $28$, $1$:

\begin{center}
\small
\begin{tabular}{@{}crr@{}}
\toprule
tier & lower convergent $L_p/L_q$ & upper convergent $U_p/U_q$ \\
\midrule
0 & $569/359$ & $485/306$ \\
1 & $1054/665$ & $1539/971$ \\
2 & $25781/16266$ & $24727/15601$ \\
3 & $50508/31867$ & $125743/79335$ \\
4 & $176251/111202$ & $301994/190537$ \\
\bottomrule
\end{tabular}
\end{center}

Only the value $m=2762$, whose $\thetaq$-power is within $3\cdot10^{-9}$ of an integer in
logarithmic scale, needs tier $4$; its certificate compares integers of roughly $90{,}000$
digits, which the kernel's exact arithmetic replays without difficulty.  No axioms beyond
the Lean kernel are involved; in particular no \lean{native_decide}.

\subsection{Arithmetic descent}
\label{co:descent}

\begin{theorem}[Affine descent]
\label{co:thm-descent}
If the two integral outputs factor simultaneously as $2^x=2^ru^k$ and $3^x=3^rv^k$ with
$r\ge0$, $k\ge1$, then $(x-r)/k$ is again a solution.  Combined with the $x\ge8$ region this
gives the unconditional constraint
\[
r+8k\le x
\]
for nonintegral $x$; in particular common perfect-power shapes of the two outputs are
strongly restricted.
\end{theorem}

\begin{statusbox}{\statuslean{} --- \lean{affine_descent_integral_powers},
\lean{base_depth_add_eight_mul_degree_le},
\lean{min_output_base_factorization_add_eight_le} in \lean{ArithmeticRigidity}.  With the
$4096$ region the constant $8$ improves to $12$ by bookkeeping; with the interval scan of
Section~\ref{sec:finitecheck} it would improve to $20$ at the \statuscomp{} level.}
\end{statusbox}

For the \emph{least} counterexample the same descent later becomes an exact primitivity
statement in Section~\ref{sec:conditional}: no nontrivial matched decomposition exists at
all.

\subsection{Local progression obstructions and zero density}
\label{co:ap}

The function $t\mapsto t^{\thetaq}$ is strictly convex with second differences at moderate
steps lying strictly between $0$ and $1$ --- hence never integers.

\begin{theorem}[Committed AP obstruction]
\label{co:thm-ap}
Let $n,h\ge1$ with $h^5\le n$.  Then $n,n+h,n+2h$ are not all candidates.  In particular no
three consecutive naturals are candidates, and among the positive integers up to $3N$ at
most $2N$ are candidates.
\end{theorem}

\begin{statusbox}{\statuslean{} ---
\lean{not_three_arithmetic_progression_twoBaseNaturalCandidates},
\lean{not_three_consecutive_logThreeDivLogTwo_rpow_integers},
\lean{twoBaseNaturalCandidate_Icc_card_le}.  The sharper criterion
$\thetaq(\thetaq-1)h^2<n^{2-\thetaq}$ is kernel-verified in
\lean{SharperProgressions}, and the analogous four-term criterion is kernel-verified in
\lean{ThirdDifference}.  The arbitrary-order statement for every $r\ge2$ is
\lean{not_arithmetic_progression_twoBaseNaturalCandidates_of_curvature} in
\lean{HigherDifference}, which formalizes a generic iterated forward-difference mean-value
identity and applies uniformly in the order $r$.}
\end{statusbox}

For finite verification the same curvature mechanism is also available in a purely rational
effective form: if $h^{400}\le n^{83}$ --- step exponent $83/400=0.2075$, against the $1/5$
of Theorem~\ref{co:thm-ap} --- then $n,n+h,n+2h$ are not all candidates.

\begin{statusbox}{\statuslean{} ---
\lean{not_three_arithmetic_progression_twoBaseNaturalCandidates_sharp} in
\lean{SharpProgression}, with the kernel-replayed enclosure $\thetaq<317/200$ obtained from
the exact integer comparison $3^{200}<2^{317}$.}
\end{statusbox}

The hypothesis bounding $h$ is essential: for steps comparable to $n$ the second difference
is large and the obstruction vanishes.  This locality is a genuine barrier --- a
\emph{global} progression obstruction would, for example, forbid candidates with odd part
$2^s\pm1$ via progressions such as $(2^i,\,2^{i+t-1},\,2^i(2^t-1))$, but no such global
statement is available; under a counterexample the exact monoid of
Section~\ref{sec:conditional} shows sparse candidates whose additive relations the curvature
argument cannot see.

Feeding the local obstruction into the theory of Roth numbers gives quantitative sparsity.

\begin{theorem}[Zero density]
\label{co:thm-density}
For every $H\ge1$ the number of candidates below $N$ is at most
$H^5+\lceil(N-H^5)/H\rceil\,r_3(H)$, where $r_3(H)$ is the maximal size of a $3$-AP-free
subset of $\{0,\dots,H-1\}$.  Consequently the candidate counting function is $o(N)$: the
candidates, and a fortiori the values $2^x$ of nonintegral solutions, have asymptotic
density zero.
\end{theorem}

\begin{statusbox}{\statuslean{} --- \lean{twoBaseNaturalCandidatesBelow_card_le},
\lean{twoBaseNaturalCandidatesBelow_isLittleO_id},
\lean{tendsto_twoBaseNonintegerCandidatesBelow_density_zero} in \lean{ZeroDensity} and
\lean{NonintegerDensity}.}
\end{statusbox}

\subsection{Multiplicative rank two and the common odd core}
\label{co:rank}

\begin{theorem}[Rank two]
\label{co:thm-rank2}
The prime-exponent vectors $\mathbf v(a)=(v_p(a))_p$ of all candidates span a rational space
of dimension at most two.  Consequently at most $(\clog_2N)^2$ candidates lie below $N$.
\end{theorem}

\begin{statusbox}{\statuslean{} --- \lean{finrank_span_primeExponentVectors_le_two},
\lean{twoBaseNaturalCandidatesBelow_card_le_clog_sq} in \lean{MultiplicativeRank}.  The
mechanism abstracts the determinant method: three candidates with independent exponent
vectors would give three multiplicatively independent bases, contradicting the three-base
theorem of Section~\ref{sec:problem} through the irrationality of $\thetaq$
(\lean{irrational_logThreeDivLogTwo}).}
\end{statusbox}

Since $2$ is itself a candidate, all odd parts lie on at most one further rational
direction.  Pairwise this yields the cross-valuation identity, and globally a common odd
core.

\begin{lemma}[Cross-valuation identity]
\label{co:lem-cross-input}
Let $a,b\in\Cand$ and let $p$ be an odd prime with $\vpp{p}{a}>0$.  Then for every odd prime
$q$,
\[
\vpp{p}{a}\,\vpp{q}{b}=\vpp{p}{b}\,\vpp{q}{a}.
\]
\end{lemma}

\begin{statusbox}{\statuslean{} --- \lean{odd_factorization_cross_mul} in
\lean{CandidateStructure}; the fixed-power relations
\lean{fixed_power_relation_of_odd_prime_factor} follow.}
\end{statusbox}

\begin{theorem}[Common odd core]
\label{co:thm-oddcore-lean}
There is a single odd $w>1$ such that every candidate equals $2^iw^j$ for some
$i,j\in\Nzero$; candidates that are not powers of two have $j\ge1$.  In exponent
coordinates: every solution has the form $x=i+j\log_2w$ with $i,j\in\Nzero$.
\end{theorem}

\begin{statusbox}{\statuslean{} --- \lean{exists_common_odd_core},
\lean{exists_common_odd_core_split}, \lean{exists_common_odd_core_exponent} in
\lean{OddCore}.  The canonical (gcd-normalized) version with uniqueness of $w$ and of each
coordinate pair is \lean{existsUnique_primitive_common_odd_core} in
\lean{PrimitiveOddCore}.}
\end{statusbox}

\begin{corollary}[Improved counting]
\label{co:cor-counting}
The number of candidates below $N$ is at most $\clog_2N\cdot\clog_3N$, and each dyadic
block $[2^T,2^{T+1})$ contains at most $\clog_3(2^{T+1})\approx0.631\,(T+1)$ candidates.
In exponent coordinates, the number of solutions with $x\in[T,T+1)$ grows at most linearly
in $T$, against the quadratic global bound.
\end{corollary}

\begin{statusbox}{\statuslean{} --- \lean{twoBaseNaturalCandidatesBelow_card_le_clog_mul_clog},
\lean{twoBaseNonintegerCandidatesBelow_card_le_clog_mul_clog},
\lean{twoBaseNaturalCandidatesInDyadicBlock_card_le} in \lean{OddCore}.}
\end{statusbox}

The logarithmic-square order resists improvement by counting alone: if a candidate with an
odd prime factor exists, multiplicative closure generates injective boxes $2^ia^j$, so the
count below $2^ka^l$ is at least $kl$
(\lean{twoBaseNaturalCandidatesBelow_card_ge_box},
\lean{twoBaseNonintegerCandidatesBelow_card_ge_box}) --- a conditional lower bound of the
same quadratic-in-logarithm shape, made exact in Section~\ref{sec:conditional}.

\subsection{Transcendence}
\label{co:transcendence}

The corpus contains a complete Lean proof of the Gelfond--Schneider theorem (Hilbert's
seventh problem), following Hua's exposition \cite{Hua1982}, ported from Michail
Karatarakis's Mathlib formalization (Mathlib PR \#42911, Apache~2.0).

\begin{theorem}[Gelfond--Schneider]
\label{co:thm-gs}
If $\alpha,\beta$ are algebraic, $\alpha\ne0,1$, and $\beta$ is irrational, then
$\alpha^\beta$ is transcendental.
\end{theorem}

\begin{statusbox}{\statuslean{} ---
\lean{GelfondSchneider.transcendental_cpow_of_isAlgebraic_of_irrational}.}
\end{statusbox}

\begin{corollary}[Algebraic--integral dichotomy]
\label{co:cor-transcendental}
If $2^x\in\Z$ then $x$ is either an integer or transcendental; algebraicity of such $x$ is
equivalent to integrality.  Every nonintegral two-base solution is transcendental, and
$\thetaq=\log_23$ is transcendental.
\end{corollary}

\begin{statusbox}{\statuslean{} --- \lean{transcendental_of_not_integer_of_two_rpow_integer},
\lean{isAlgebraic_iff_integer_of_two_rpow_integer},
\lean{transcendental_of_not_integer_of_two_three_rpow_integer},
\lean{transcendental_logThreeDivLogTwo}.}
\end{statusbox}

The dichotomy does not resolve the conjecture --- the unknown exponent is allowed to be
transcendental --- but it constrains every construction below: the primitive generator
$\beta$ of Section~\ref{sec:conditional} is transcendental, as is $\log_2w$ for the odd
core.

\subsection{Classification of auxiliary bases: the \texorpdfstring{$3$}{3}-smooth theorem}
\label{co:threesmooth}

\begin{theorem}[$3$-smooth classification]
\label{co:thm-threesmooth}
Suppose $x\notin\Z$ and $2^x,3^x\in\Z$.  Then for a positive natural $a$, the power $a^x$ is
an integer if and only if $a=2^u3^v$ is $3$-smooth.  In particular a third base with any
prime factor other than $2$ and $3$ forces $x\in\Z$.
\end{theorem}

\begin{statusbox}{\statuslean{} ---
\lean{rpow_integer_iff_eq_two_pow_mul_three_pow_of_not_integer},
\lean{integer_of_two_three_a_rpow_integer_of_prime_factor} in \lean{ThreeSmooth}.}
\end{statusbox}

This is an exact classification of bases for a \emph{fixed} hypothetical counterexample.  It
should not be confused with the classification of candidate values $2^x$ as $x$ ranges over
all solutions, which is the subject of the next section.  Note also the scope limitation
explained in Section~\ref{sec:attempts}: the theorem controls the quantity $b^x$, not
divisibility of the integer $2^x$ by $b$.

\subsection{Unconditional gap equivalences}
\label{co:dichotomy}

Because natural multiples of a solution are solutions, one nonintegral solution generates
infinitely many, with irrational generator; Dirichlet approximation and a rotation-stepping
argument spread their fractional parts densely.

\begin{theorem}[Density dichotomy]
\label{co:thm-dichotomy}
If some $x\notin\Z$ satisfies $2^x,3^x\in\Z$, then for every $t\in[0,1]$ and
$\varepsilon>0$ there is a nonintegral solution $y$ with $|\fract{y}-t|<\varepsilon$; and
the set of nonintegral solutions is infinite.
\end{theorem}

\begin{theorem}[Gap equivalences]
\label{co:thm-equiv}
Each of the following is equivalent to Conjecture~\ref{conj:AE}:
\begin{enumerate}[label=\textup{(\roman*)}]
\item some nonempty open interval $(u,v)\subseteq(0,1)$ contains no fractional part of a
nonintegral solution;
\item the fractional parts of nonintegral solutions are bounded away from $1$;
\item the fractional parts of nonintegral solutions are bounded away from $0$.
\end{enumerate}
\end{theorem}

\begin{statusbox}{\statuslean{} --- \lean{dense_fract_of_noninteger_solution},
\lean{infinite_noninteger_solutions_of_exists},
\lean{alaogluErdosConjecture_iff_fract_avoids_interval},
\lean{alaogluErdosConjecture_iff_fract_bounded_away_one},
\lean{alaogluErdosConjecture_iff_fract_bounded_away_zero} in
\lean{FractionalPartDensity}, built on the reusable
\lean{exists_pos_nat_fract_close_of_irrational}.}
\end{statusbox}

\begin{remark}[Consistency pincer]
\label{co:rem-pincer}
Theorem~\ref{co:thm-gaps} bounds $\fract{x}$ away from $0$ and $1$ by margins
$\asymp2^{-\lfloor x\rfloor}$ that decay with the height of the solution, while
Theorem~\ref{co:thm-dichotomy} produces fractional parts arbitrarily close to $0$ and $1$
along multiples of unbounded height.  The two results meet edge-to-edge without
contradiction, and Theorem~\ref{co:thm-equiv} shows the decay is essential: \emph{any}
height-independent gap, however small and wherever located, is exactly as strong as the full
conjecture.  The elementary gap bounds are therefore of optimal shape, and the
fractional-part axis measures the entire remaining distance to Conjecture~\ref{conj:AE}.
\end{remark}

These equivalences are \emph{unconditional} biconditionals.  Conditional on failure, the
exact structure of the next section gives kernel-verified convergence for every finite
trigonometric polynomial and the paper-level strengthening to full blockwise Weyl
equidistribution (Section~\ref{sec:conditional}).
\section{Exact conditional structure of all solutions}
\label{sec:conditional}

The most consequential fact in the current corpus is not another exclusion range.  It is the
complete description of \emph{all} solutions under the hypothesis that even one nonintegral
solution exists.  Conditional on failure of Conjecture~\ref{conj:AE}, the solution set is not
merely sparse or contained in some rank-two set: it is a free rank-two additive monoid with an
explicit canonical generator, exactly countable in every dyadic block, and equidistributed
blockwise on the circle.

Throughout this section we write
\[
 \Sol:=\{x\in\R:2^x\in\N\text{ and }3^x\in\N\},\qquad
 \Cand:=\{2^x:x\in\Sol\},\qquad
 \Pair:=\{(2^x,3^x):x\in\Sol\},
\]
and we say a counterexample exists when $\Sol\not\subseteq\Nzero$.  The canonical primitive odd
core, the rational-power index, the power-of-three denominator normalization, the unique least
solution, the simultaneous primitive output normalization, the exact free-monoid
classification, the exact unit-block and dyadic counts, the cumulative quadratic asymptotic,
and the finite-Fourier block averages are all \statuslean{}.  Only the extension of the block
averages from finite trigonometric polynomials to arbitrary continuous tests or interval
indicators (Theorem~\ref{cd:block-equidist}) remains \statusnew{}.

\subsection{The primitive common odd core}

\begin{theorem}[Primitive common odd core]
\label{cd:odd-core}
Assume $\Cand$ contains a non-power of two.  Then there is a unique odd integer $w>1$ that is
power-free in the sense
\[
\gcd\{\vpp{p}{w}:p\mid w\}=1,
\]
and every candidate $b\in\Cand$ has a unique representation
\[
b=2^iw^j,\qquad i,j\in\Nzero.
\]
\end{theorem}

\begin{proof}
Choose a non-power-of-two candidate $a$.  Its odd valuation vector $A=(A_q)_{q\ne2}$,
$A_q=\vpp{q}{a}$, is nonzero and finitely supported.  Let
\[
g=\gcd\{A_q:A_q>0\},\qquad c_q=A_q/g,\qquad w=\prod_{q\ne2}q^{c_q}.
\]
Then $w$ is odd, greater than one, and the positive $c_q$ have gcd one.

Let $b\in\Cand$ and set $B_q=\vpp{q}{b}$.  Fix an odd prime $p$ with $A_p>0$.  By the
cross-valuation identity of Section~\ref{sec:corpus}, $A_pB_q=B_pA_q$ for all odd $q$.  If
$A_q=0$, this forces $B_q=0$.  On the support of $A$ it says
\[
B_q=\frac{B_p}{A_p}A_q=\Bigl(\frac{gB_p}{A_p}\Bigr)c_q .
\]
Put $t=gB_p/A_p\in\Q_{\ge0}$.  Every $tc_q$ is an integer.  Since the $c_q$ have gcd one, there
are integers $r_q$, finitely many nonzero, with $\sum r_qc_q=1$; therefore
$t=\sum_qr_q(tc_q)\in\Z$.  Writing $j=t\in\Nzero$, the odd part of $b$ is exactly $w^j$, so
$b=2^iw^j$ with $i=\vpp{2}{b}$.

Uniqueness of $(i,j)$ for fixed $w$ follows from the odd valuations and the $2$-adic valuation.
If $w'$ is another primitive common core, the odd valuation vectors of $w$ and $w'$ generate the
same rational ray and are both primitive integer vectors, hence equal.
\end{proof}

\begin{statusbox}{\statuslean{}}
\lean{existsUnique_primitive_common_odd_core} in \lean{PrimitiveOddCore}.  The module defines
\lean{oddPrimeValuationGCD}, proves its equivalence with power-theoretic primitivity for $w>1$,
and kernel-checks both uniqueness of the normalized core and uniqueness of every pair $(i,j)$.
The earlier non-normalized existence statement \lean{exists_common_odd_core} recorded in
Section~\ref{sec:corpus} is subsumed by it.
\end{statusbox}

The gcd normalization is exactly what makes $w$ canonical: it prevents a hidden perfect power
from being absorbed into the exponent $j$, so that both the core and the coordinates are
determined by $\Cand$ alone.

Set
\[
\alpha:=\log_2w,\qquad E:=3^{\alpha}=w^{\thetaq}.
\]
Every candidate exponent then has the form $x=i+j\alpha$.  The number $\alpha$ is irrational
(otherwise the rational-exponent lemma applied to $2^\alpha=w$ would make it integral, forcing
the odd $w>1$ to be a power of two), and in fact transcendental by the algebraic--integral
dichotomy of Section~\ref{sec:corpus}.  The candidate condition becomes
\[
3^{i+j\alpha}=3^iE^j\in\N,
\]
and at this stage the admissible pairs $(i,j)$ still form an unknown subset of $\Nzero^2$.  The
next subsection shows that this subset is completely rigid.

\subsection{The rational-power index and the denominator normalization}

Because a non-power-of-two candidate exists, some $j>0$ has $E^j\in\Q$.  Let
\[
d:=\min\{j\in\N_{>0}:E^j\in\Q\}.
\]

\begin{lemma}[All rational powers are multiples of the least one]
\label{cd:rational-power-index}
For every $j\in\Z$:\quad $E^j\in\Q\iff d\mid j$.
\end{lemma}

\begin{proof}
The set $J:=\{j\in\Z:E^j\in\Q\}$ is an additive subgroup of $\Z$ (it contains $0$, products add
exponents, and reciprocals negate them), and every nonzero subgroup of $\Z$ equals $d\Z$ for its
least positive element $d$.
\end{proof}

\begin{statusbox}{\statuslean{}}
\lean{rationalPowerIndices}, \lean{AddSubgroup.Int.exists_least_positive_generator}, and
\lean{exists_rationalPowerIndex} in \lean{RationalPowerIndex}.
\end{statusbox}

Write the positive rational $E^d$ in lowest terms.  Its denominator has no prime other than $3$:
taking a candidate with odd-core exponent $j=dk>0$ and $2$-exponent $i$, integrality of
$3^i(E^d)^k$ shows that any other prime in the reduced denominator would survive in the $k$-th
power.  Hence there are unique integers $a\ge0$ and $A\ge1$ with
\begin{equation}
E^d=\frac{A}{3^a},\qquad\text{and}\qquad 3\nmid A\ \text{ when }a>0,
\label{cd:eq-Ed}
\end{equation}
so that $a=0$ or $3\nmid A$ in every case.  Define
\begin{equation}
\beta:=a+d\alpha=a+d\log_2w,\qquad
M:=2^{\beta}=2^aw^d,\qquad
A=3^{\beta}=3^aE^d .
\label{cd:eq-betaMA}
\end{equation}
Thus $\beta$ is itself a solution, and it is irrational because $\alpha$ is.  (Where the
denominator normalization is written with a named integer numerator $c$, the identity
$E^d=c/3^a$ with $a=0$ or $3\nmid c$ is exactly \eqref{cd:eq-Ed} with $c=A=3^\beta$.)

\begin{statusbox}{\statuslean{}}
The denominator-support step is \lean{exists_three_pow_denominator_of_rational_power} in
\lean{ThreeDenominatorNormalization}; the normalized integrality criterion is
\lean{three_pow_mul_normalized_pow_integer_iff} in \lean{RationalPowerIndex}.
\end{statusbox}

Equation~\eqref{cd:eq-Ed} also has a field-theoretic reading: $E$ is a root of the binomial
$X^d-A/3^a$, which is the minimal polynomial of $E$ over $\Q$, so $[\Q(E):\Q]=d$.  That
computation, together with the equivalent one-number and localized-radical formulations of the
conjecture that follow from it, is developed in Section~\ref{sec:radical}; we do not repeat it
here.

\subsection{The exact primitive-generator theorem}

\begin{theorem}[Exact primitive generator and solution monoid]
\label{cd:generator}
Assume Conjecture~\ref{conj:AE} is false.  With $w,d,a,\beta,M,A$ as above:
\begin{enumerate}[label=\textup{(\roman*)}]
\item every solution has a unique representation $x=n+k\beta$ with $n,k\in\Nzero$, that is,
\[
 \Sol=\{n+k\beta:n,k\in\Nzero\};
\]
\item every candidate has a unique representation $m=2^nM^k$ with $n,k\in\Nzero$, that is,
\[
 \Cand=\{2^nM^k:n,k\in\Nzero\};
\]
\item every output pair has a unique representation $(2^x,3^x)=(2^nM^k,3^nA^k)$;
\item $\beta$ is the unique least nonintegral solution.
\end{enumerate}
\end{theorem}

\begin{proof}
Let $x\in\Sol$.  By Theorem~\ref{cd:odd-core}, $2^x=2^iw^j$ for unique $i,j\in\Nzero$, so
$x=i+j\alpha$.  From $3^x=3^iE^j\in\N$ we get $E^j\in\Q$, so Lemma~\ref{cd:rational-power-index}
gives $j=dk$ for some $k\in\Nzero$.  Using \eqref{cd:eq-Ed},
\[
3^x=3^i(E^d)^k=3^{\,i-ak}A^k .
\]
Because $A/3^a$ is reduced, this is an integer exactly when $i\ge ak$.  Put $n=i-ak\in\Nzero$;
then
\[
x=i+dk\alpha=n+k(a+d\alpha)=n+k\beta .
\]
Conversely, for any $n,k\in\Nzero$,
\[
2^{\,n+k\beta}=2^nM^k\in\N,\qquad 3^{\,n+k\beta}=3^nA^k\in\N,
\]
so $n+k\beta\in\Sol$.  If $n+k\beta=n'+k'\beta$ with $k\ne k'$, then $\beta=(n'-n)/(k-k')\in\Q$,
a contradiction; hence the representation is unique.  This proves (i), and (ii)--(iii) follow by
exponentiation.  A nonintegral solution has $k\ge1$; the least such value is attained at
$(n,k)=(0,1)$ and equals $\beta$, uniquely so by uniqueness of the representation.
\end{proof}

\begin{statusbox}{\statuslean{}}
\lean{exists_canonical_primitiveGenerator_of_not_alaogluErdosConjecture} in
\lean{PrimitiveGenerator}, built from
\lean{exists_exact_solution_monoid_of_fixed_common_oddCore} in \lean{RationalPowerIndex}.  The
Lean statement supplies the canonical odd core $w$, the rational-power index $d$, the depth $a$,
the numerator $A$, and the unique-coordinate clauses in one package.  Conditional on failure,
$\Sol$ is not merely contained in a rank-two additive set: every solution has a unique pair of
coordinates in the free additive monoid on $1$ and the irrational least generator $\beta$.  The
set of occurring odd-core exponents is exactly $d\,\Nzero$, with the power of two governed by
one linear threshold.
\end{statusbox}

Packaged algebraically, Theorem~\ref{cd:generator} says that
\[
 (\Nzero^2,+)\xrightarrow{\ \sim\ }(\Sol,+),\quad(n,k)\mapsto n+k\beta,
 \qquad
 (\Nzero^2,+)\xrightarrow{\ \sim\ }(\Cand,\cdot),\quad(n,k)\mapsto2^nM^k,
\]
are monoid isomorphisms, and that the output-pair monoid $\Pair$ is freely generated by $(2,3)$
and $(M,A)$.  Every multiplicative relation among candidates is the obvious coordinate relation;
there is no room for an unexpected coincidence.  This exact conditional model is the benchmark
against which any future density, progression, or descent argument must be tested.

The phrase ``least nonintegral'' is materially stronger than choosing an arbitrary
counterexample.  It forces affine primitivity: no nontrivial common root can be extracted from
the two output integers while preserving the same base-adic offset.  It is the arithmetic
analogue of choosing a primitive lattice direction, and the next subsection turns it into exact
exclusions.

\subsection{Primitivity of the least output pair}

\begin{corollary}[No matched base depth]
\label{cd:no-depth}
For the primitive pair $(M,A)=(2^{\beta},3^{\beta})$ we have
$\min\{\vpp{2}{M},\vpp{3}{A}\}=0$; that is, $M$ is odd or $3\nmid A$ (possibly both).
\end{corollary}

\begin{proof}
If both valuations were positive, then $2^{\beta-1}=M/2$ and $3^{\beta-1}=A/3$ would be
integers, so $\beta-1$ would be a smaller nonintegral solution.
\end{proof}

\begin{corollary}[No common perfect-power degree]
\label{cd:no-common-power}
There is no $k\ge2$ for which both $M$ and $A$ are perfect $k$-th powers.
\end{corollary}

\begin{proof}
If $M=U^k$ and $A=V^k$, then $2^{\beta/k}=U$ and $3^{\beta/k}=V$, so $\beta/k$ is a smaller
nonintegral solution.
\end{proof}

\begin{corollary}[Exact affine primitivity]
\label{cd:affine-primitive}
If $M=2^rU^k$ and $A=3^rV^k$ with $r\in\Nzero$, $k\in\N_{>0}$, and $U,V\in\N$, then $r=0$ and
$k=1$.
\end{corollary}

\begin{proof}
The affine descent theorem of Section~\ref{sec:corpus} produces the solution $(\beta-r)/k$,
which is nonintegral since $\beta$ is irrational; if $r>0$ or $k>1$ it is strictly smaller than
$\beta$, contradicting leastness.
\end{proof}

\begin{statusbox}{\statuslean{}}
\lean{exists_least_twoBaseNonintegerSolution},
\lean{IsLeastTwoBaseNonintegerSolution.unique},
\lean{IsLeastTwoBaseNonintegerSolution.not_both_base_dvd_outputs},
\lean{IsLeastTwoBaseNonintegerSolution.no_common_perfect_power}, and
\lean{IsLeastTwoBaseNonintegerSolution.affine_primitive} in \lean{LeastSolution}.
\end{statusbox}

The odd-core normalization of $M$ has a symmetric counterpart on the other output, and the two
can be imposed simultaneously.

\begin{theorem}[Simultaneous primitive output normalization]
\label{cd:simultaneous-output}
If the conjecture fails, its least nonintegral solution $\beta$ admits simultaneous
factorizations
\[
 M=2^\beta=2^aw^d,\qquad A=3^\beta=3^bv^e,
\]
where $w>1$ is odd and power-primitive, $v>1$ is power-primitive with $3\nmid v$, and $d,e>0$.
Moreover $a=0$ or $b=0$; the exponents $a$ and $b$ are the exact base-adic depths; $d$ and $e$
are the gcds of the prime valuations of the corresponding base-free parts; and
\[
 \gcd\!\left(\gcd(d,e),\,|b-a|\right)=1.
\]
Since every counterexample satisfies $\beta\ge12$, these normalizations also obey the size
dichotomy
\[
 w^d\ge4096\qquad\text{or}\qquad v^e\ge3^{12}.
\]
\end{theorem}

\begin{statusbox}{\statuslean{}}
\lean{exists_simultaneous_primitive_output_normalization},
\lean{IsLeastTwoBaseNonintegerSolution.simultaneousCoordinateGCD_eq_one}, and
\lean{exists_large_primitiveCore_of_not_alaogluErdosConjecture} in
\lean{OutputNormalization}.  The gcd-one statement simultaneously controls the two outer power
degrees and the difference of the two base-adic depths; it is strictly stronger than merely
excluding a common perfect-power degree (Corollary~\ref{cd:no-common-power}).
\end{statusbox}

These restrictions are strong, but they still admit infinitely many formal integer parameter
patterns.  The unresolved analytic content is untouched: the pair $(M,A)$ must satisfy
\[
 \frac{\log M}{\log2}=\frac{\log A}{\log3},
\]
and no known local valuation argument turns the depth dichotomy $a=0$ or $b=0$, or the
gcd-one condition, into a contradiction.  Every constraint proved so far is compatible with the
existence of such a pair; what fails to exist must fail for a transcendence reason, not a
valuation-theoretic one.

\subsection{Exact counting}
\label{cd:counting}

The free-monoid description makes the population of a hypothetical counterexample completely
explicit, block by block.

\begin{theorem}[Exact unit-block and dyadic counts]
\label{cd:block-count}
Assume a counterexample exists and let $\beta$ be the primitive generator.  For every
$N\in\Nzero$, the nonintegral solutions in $[N,N+1)$ are precisely
\[
x_{N,k}=\lceil N-k\beta\rceil+k\beta,
\qquad 1\le k\le\Bigl\lfloor\frac{N+1}{\beta}\Bigr\rfloor,
\]
so their number is exactly $\lfloor(N+1)/\beta\rfloor$.  Consequently
\[
\#\bigl(\Sol\cap[N,N+1)\bigr)
=\#\bigl(\Cand\cap[2^N,2^{N+1})\bigr)
=1+\Bigl\lfloor\frac{N+1}{\beta}\Bigr\rfloor .
\]
\end{theorem}

\begin{proof}
For fixed $k\ge1$ there is at most one integer $n$ with $N\le n+k\beta<N+1$; such an $n\ge0$
exists if and only if $k\beta<N+1$, and then $n=\lceil N-k\beta\rceil$.  Irrationality of
$\beta$ makes $(N+1)/\beta$ nonintegral, which gives the stated range.  The unique integral
solution in the block is $N$ itself, whence the additional $1$; passing through $m=2^x$ turns
the exponent block $[N,N+1)$ into the magnitude block $[2^N,2^{N+1})$ bijectively.
\end{proof}

\begin{statusbox}{\statuslean{}}
\lean{ExactCounting} proves the explicit ceiling-image parametrization,
\lean{twoBaseNonintegerSolutionsInUnitBlock_ncard},
\lean{twoBaseNaturalCandidatesInDyadicBlock_card}, and the exact finite-sum identity
\lean{cumulativeTwoBaseCandidateCount_eq}.  \lean{QuadraticAsymptotic} identifies that
cumulative count with the candidate count below $2^T$ and proves the sharp bounds and the limit
in \eqref{cd:eq-cumulative}.
\end{statusbox}

In particular a failed conjecture would produce only linearly many candidates in the logarithmic
block index,
\[
 \#\bigl(\Cand\cap[X,2X]\bigr)\asymp\log X ,
\]
and the kernel-verified zero-free region gives $\beta\ge12$, hence $\beta>12$ because $\beta$ is
irrational.  The number of \emph{nonintegral} solutions in the block is therefore at most
$\lfloor N/12\rfloor$, while the full solution count --- equivalently the dyadic candidate
count --- is at most $1+\lfloor N/12\rfloor$.  Under the directed-rounding scan these bounds
improve respectively to $\lfloor N/27\rfloor$ and $1+\lfloor N/27\rfloor$, at the weaker
\statuscomp{} trust level.  Summing over blocks,
the cumulative candidate count below $2^T$ is
\begin{equation}
C(T)=T+\sum_{r=1}^{T}\Bigl\lfloor\frac r\beta\Bigr\rfloor
=\frac{T^2}{2\beta}+O(T).
\label{cd:eq-cumulative}
\end{equation}
More precisely, Lean proves the two-sided bound and the limit
\[
 \frac{T(T+1)}{2\beta}\le C(T)\le\frac{T(T+1)}{2\beta}+T,
 \qquad
 \frac{C(T)}{T^2}\longrightarrow\frac1{2\beta}.
\]
This is conditionally sharper than the unconditional verified bounds
($(\clog_2N)^2$ and $\clog_2N\cdot\clog_3N$ of Section~\ref{sec:corpus}) and matches them in
leading order: with $\beta>12$ the leading coefficient $1/(2\beta)$ is below $1/24$.

There is also an unconditional growth reformulation, obtained by comparing the two cases.  Under
the conjecture the count below $2^T$ is exactly $T$; under failure the preceding limit is the
strictly positive constant $1/(2\beta)$.  Hence
\begin{equation}
 \text{Conjecture~\ref{conj:AE} holds}
 \iff
 \frac{\#\bigl(\Cand\cap[1,2^T)\bigr)}{T^2}\longrightarrow0 ,
 \label{cd:eq-growth}
\end{equation}
which is \lean{alaogluErdosConjecture_iff_candidate_count_below_two_pow_subquadratic}.
\statuslean{}

This exact population is the correct benchmark for every proposed density or determinant
argument.  A method that proves only $O((\log X)^2)$ candidates in $[X,2X]$, for instance, is a
valid unconditional theorem but cannot contradict Theorem~\ref{cd:block-count}: it is weaker
than the conditional truth it would need to refute.

\subsection{Blockwise irrational rotation and equidistribution}

By Theorem~\ref{cd:block-count} the fractional part of $x_{N,k}$ is $\fract{k\beta}$, so a high
unit block contains exactly an initial segment
\[
 \fract{\beta},\ \fract{2\beta},\ \ldots,\ \fract{K_N\beta},
 \qquad K_N=\Bigl\lfloor\frac{N+1}{\beta}\Bigr\rfloor,
\]
of one irrational rotation orbit, and $K_N\to\infty$.

\begin{theorem}[Finite Fourier tests on exact blocks]
\label{cd:block-fourier}
Assume a counterexample exists and let $\beta$ be its primitive generator.  For every finitely
supported coefficient function $c:\Z\to\C$, put
\[
 P_c(t)=\sum_{h\in\Z}c(h)e^{2\pi iht}.
\]
Then, on the exact ceiling-parametrized points $x_{N,k}=\lceil N-k\beta\rceil+k\beta$,
\[
 \frac1{K_N}\sum_{k=1}^{K_N}P_c(x_{N,k})\longrightarrow c(0)
 \qquad(N\to\infty).
\]
\end{theorem}

\begin{statusbox}{\statuslean{}}
\lean{BlockWeyl} proves vanishing of every nonzero normalized Fourier mode on the exact block
points, and \lean{BlockEquidistribution} proves
\lean{tendsto_primitiveUnitBlockTrigonometricAverage} for every finitely supported coefficient
function, with limit \lean{c 0}.  This is convergence against finite trigonometric polynomials;
it is not yet a theorem for arbitrary continuous tests.
\end{statusbox}

\begin{theorem}[Blockwise Weyl equidistribution]
\label{cd:block-equidist}
Assume a counterexample exists.  For every interval $I\subseteq[0,1]$ whose boundary has measure
zero,
\[
\frac{\#\{x\in(\Sol\setminus\Nzero)\cap[N,N+1):\fract{x}\in I\}}
{\#\bigl((\Sol\setminus\Nzero)\cap[N,N+1)\bigr)}\longrightarrow|I|
\qquad(N\to\infty).
\]
\end{theorem}

\begin{proof}
By Theorem~\ref{cd:block-count} the relevant fractional parts are
$\fract{\beta},\fract{2\beta},\dots,\fract{K_N\beta}$ with $K_N=\lfloor(N+1)/\beta\rfloor\to
\infty$.  For every nonzero integer $h$ the Weyl sum is a finite geometric series,
\[
\frac1K\sum_{k=1}^{K}e^{2\pi ihk\beta}
=\frac{e^{2\pi ih\beta}\bigl(1-e^{2\pi ihK\beta}\bigr)}{K\bigl(1-e^{2\pi ih\beta}\bigr)}
\longrightarrow0,
\]
since $\beta$ is irrational and hence $e^{2\pi ih\beta}\ne1$.  Weyl's criterion concludes.
\end{proof}

\begin{statusbox}{\statusnew{}}
The finite-trigonometric-polynomial input is kernel-verified by
Theorem~\ref{cd:block-fourier}, but the Stone--Weierstrass and Weyl-criterion extension to all
continuous tests, and then to interval indicators, has not been formalized.  No effective
discrepancy rate follows without further Diophantine input on $\beta$: the geometric-sum proof
controls each fixed Fourier mode but not uniformly over many modes.
\end{statusbox}

This regularity is attractive, and it is worth stating plainly why it does not by itself decide
anything.  The integrality condition being tested is not a fixed positive-measure target in the
circle; it is an exponentially shrinking target, since the relevant approximation quality must
improve like the reciprocal of the integer sizes involved.  Equidistribution at a fixed scale is
simply the wrong instrument for a target whose measure tends to zero geometrically, and that
mismatch recurs in the analytic approaches examined later in this report.
\section{Equivalent one-number and radical formulations}
\label{sec:radical}

The exact conditional structure of Section~\ref{sec:conditional} compresses a statement
quantified over all real $x$ into a statement about a single real number at a time.  This
section records that compression in its two kernel-verified forms --- a localization at the
prime $3$ and a radical-algebraicity normalization --- and then names the exact arithmetic
obstruction that remains.

\subsection{Localization at the prime three}

\begin{theorem}[Localization reformulation]
\label{ra:thm-localization}
Conjecture~\ref{conj:AE} is false if and only if there exists an odd integer $u>1$ with
\[
u^{\thetaq}\in\Zthree,
\]
equivalently, an odd $u>1$ and $a\in\Nzero$ with $3^au^{\thetaq}\in\N$.
\end{theorem}

\begin{proof}
If $x$ is a counterexample, write $2^x=2^au$ with $u$ odd; since $2^x$ is not a power of
two, $u>1$.  Using $(2^a)^{\thetaq}=3^a$,
\[
u^{\thetaq}=\frac{(2^x)^{\thetaq}}{(2^a)^{\thetaq}}=\frac{3^x}{3^a}\in\Zthree .
\]
Conversely, if $u^{\thetaq}=B/3^a$ with $B\in\N$, set $x=a+\log_2u$; then $2^x=2^au\in\N$
and $3^x=3^au^{\thetaq}=B\in\N$, and $x\notin\Z$ since otherwise the odd $u>1$ would be a
power of two.
\end{proof}

\begin{statusbox}{\statuslean{} ---
\lean{alaogluErdosConjecture_iff_odd_rpow_not_localized} in \lean{Localization} and
\lean{alaogluErdosConjecture_iff_oddLocalizedRadicalExclusion} in
\lean{RadicalReformulation}; the candidate-level variant
\lean{alaogluErdosConjecture_iff_no_odd_multiple_candidate}, also in \lean{Localization},
is proved independently and phrases the same content in terms of admissible candidate
bases.  The equivalence is kernel-verified; the localized-radical exclusion on its
right-hand side remains the open conjecture.  The allowance of $3$-power denominators is
exactly the freedom created by shifting the exponent by an integer.}
\end{statusbox}

The denominator restriction to powers of $3$ is not cosmetic.  Only the prime $3$ can occur
in the denominator, because the exponent shift $x\mapsto x+1$ multiplies $2^x$ by $2$ and
$3^x$ by $3$; the odd part $u$ absorbs the $2$-adic freedom completely and $\Zthree$
absorbs the remaining $3$-adic freedom.  Everything else has been eliminated.

\subsection{Reduction to power-primitive odd bases}

An odd $u>1$ that is a proper power, $u=w^d$ with $w$ odd and power-primitive, satisfies
$u^{\thetaq}=(w^{\thetaq})^d$.  Unique primitive-power decomposition therefore lets the
quantifier in Theorem~\ref{ra:thm-localization} range over power-primitive bases only, at
the cost of an extra exponent $d\ge1$.

\begin{statusbox}{\statuslean{} --- \lean{PrimitiveRadical} proves unique primitive-power
decomposition of every integer $u>1$ together with the exact equivalence
\lean{alaogluErdosConjecture_iff_primitiveOddLocalizedRadicalExclusion}.  Proper-power
choices of $u$ introduce no genuinely new cases.}
\end{statusbox}

\subsection{Canonical radical algebraicity}

In the notation of the least generator of Section~\ref{sec:conditional}, a counterexample
produces natural numbers $w,d,a,A$ with $w>1$ odd and power-primitive, $d>0$, and
\[
\beta=a+d\log_2w,\qquad 2^\beta=2^aw^d,\qquad 3^\beta=A,
\]
normalized so that $a=0$ or $3\nmid A$.  Setting
\[
E:=w^{\thetaq}=3^{\log_2w},
\]
one has $E^d=A/3^a\in\Q$, and $d$ is the least positive exponent for which $E^d$ is
rational.  Minimality of $d$ has an exact algebraic consequence.

\begin{proposition}[Exact radical degree]
\label{ra:prop-radical-degree}
The binomial $X^d-A/3^a$ is irreducible over $\Q$, it is the minimal polynomial of $E$, and
$E$ is algebraic of exact degree $d$.
\end{proposition}

\begin{proof}
For a positive rational $q$, the binomial $X^d-q$ can be reducible only if $q$ is a $p$-th
power in $\Q$ for some prime $p\mid d$ (the additional $4\mid d$ exceptional case concerns
$-4$ times a fourth power and is irrelevant for $q>0$).  If $q=r^p$ with $p\mid d$, the
positive real $E^{d/p}$ is the positive $p$-th root of $q$, hence rational, contradicting
minimality of $d$.
\end{proof}

\begin{statusbox}{\statuslean{} ---
\lean{irreducible_X_pow_sub_C_of_least_rational_power} in \lean{RadicalDegree} proves the
general norm argument, and \lean{oddCoreRpow_exact_radical_degree} specializes it to the
least normalized radical index.  It identifies the minimal polynomial with $X^d-A/3^a$ and
proves that the degree is exactly $d$.}
\end{statusbox}

A counterexample therefore produces the very special algebraic number $E=3^{\log_2w}$,
reached from a transcendental exponent.  Taking logarithms of $E^d=A/3^a$:
\begin{equation}
d\,\log w\,\log3=(\log A-a\log3)\log2 .
\label{ra:eq-quadratic-logs}
\end{equation}
This is a \emph{bilinear} relation among logarithms of integers: each side is a product of
two logarithms, not a linear form in logarithms with algebraic coefficients.  Baker's
theorem \cite{Baker1975} bounds linear forms $\sum\alpha_i\log\gamma_i$ with algebraic
$\alpha_i$ away from zero, and no substitution turns \eqref{ra:eq-quadratic-logs} into such
a form, because the coefficient of $\log w$ on the left is $d\log3$, itself a logarithm
rather than an algebraic number.  The four exponentials conjecture does forbid the relation.
Equation~\eqref{ra:eq-quadratic-logs} is thus a precise description of the
transcendence-theoretic wall, not a disguised elementary factorization problem.

\subsection{Where the wall stands in the smallest case}

The smallest structural question already meets the open core.  A $3$-smooth candidate
$m=2^a3^b$ with $b\ge1$ would require $3^aE_3^{\,b}\in\N$ for
\[
E_3=3^{\log_23}=e^{(\log3)^2/\log2}=5.7045224946911176\ldots,
\]
and even the irrationality of $E_3$ is open: it is a ``three-logarithms'' quantity outside
the reach of Gelfond--Schneider, Baker's method, and the six exponentials theorem.  No
currently known technique excludes the odd core $w=3$; this is why the corpus emphasizes
counting, structure, and verified finite regions rather than infinite families of
exclusions, and why Conjecture~\ref{ra:conj-radical-target} below is the honest name of the
obstruction.

\begin{conjecture}[Radical four-exponentials target]
\label{ra:conj-radical-target}
For every odd integer $u>1$, $u^{\log_23}\notin\Zthree$.  Equivalently, for every
power-primitive odd $w>1$ and every $d\ge1$, $(3^{\log_2w})^d\notin\Zthree$.
\end{conjecture}

By Theorem~\ref{ra:thm-localization} and the primitive reduction above, both formulations
are exactly equivalent to Conjecture~\ref{conj:AE}; neither is proved.  A proof of the
special case $u=3$ --- even of the weaker assertion that $3^{\log_23}$ is irrational --- by
any genuinely new mechanism would already be a meaningful breakthrough.

\section{The exact four-exponentials barrier}
\label{sec:barrier}
\label{sec:fourexp}

\subsection{The two-by-two exponential matrix}

Suppose $x$ is a nonintegral solution.  Consider the row vector of parameters $(1,x)$ and
the column vector of parameters $(\log2,\log3)$.  The four exponentials formed from the
products are
\[
\exp(1\cdot\log2)=2,\quad
\exp(1\cdot\log3)=3,\quad
\exp(x\log2)=2^x,\quad
\exp(x\log3)=3^x,
\]
all four algebraic --- indeed positive integers.  The two row parameters $1,x$ are
$\Q$-linearly independent because $x$ is irrational: a rational nonintegral exponent is
excluded by the rational-exponent argument of Section~\ref{sec:problem}.  The two column
parameters $\log2,\log3$ are $\Q$-linearly independent because $2^a3^b=1$ forces $a=b=0$,
that is, no nonzero powers of $2$ and $3$ agree.  The real four exponentials conjecture
predicts that at least one of the four displayed values must be transcendental, a
contradiction \cite{Waldschmidt2000}.

\begin{statusbox}{\statuslean{} ---
\lean{alaogluErdosConjecture_of_realFourExponentialsConjecture} in
\lean{FourExponentialsReduction} formalizes exactly this implication.  It does not assume
four exponentials as an axiom: the conjectural statement is defined as a proposition and the
implication is proved conditionally on it.  Four exponentials itself is open, so this is
context and diagnosis, not a proof of Conjecture~\ref{conj:AE}.}
\end{statusbox}

The reduction is a diagnostic and not merely background.  Any proposed elementary argument
must exploit integrality, positivity, or arithmetic structure beyond bare algebraicity of
those four values, or else it would prove a substantial case of the four exponentials
conjecture by another name.

\subsection{Why six exponentials does not specialize}

The six exponentials theorem applies to two independent row parameters together with three
independent column parameters, or to three rows and two columns.  Here only two canonical
logarithmic directions are available.  Adding a third base $5$ supplies the missing
direction and yields the kernel-verified three-base theorem
\[
2^x,3^x,5^x\in\Z\quad\Longrightarrow\quad x\in\Z,
\]
proved by an interpolation determinant.  In the two-base problem, any additional base built
from powers of $2$ and $3$ has logarithm in the rational span of $\log2$ and $\log3$ and
therefore does not create an independent column.

The dimension count is what decides the outcome.  With three independent row directions and
two column directions the analytic size estimate beats the arithmetic lower bound
$|\det|\ge1$; in the $2\times2$ configuration the same balance lands exactly on the
classical boundary, where the required decay below one fails.  A successful two-base
determinant argument would, by construction, be a new four-exponentials argument.

The classification of auxiliary bases makes the obstruction precise: \lean{ThreeSmooth} and
the denominator-controlled results in \lean{RationalThirdBase} show that an auxiliary base
with a prime factor outside $\{2,3\}$ would supply a genuinely new column, but its
corresponding exponential value is not forced to be an integer --- or even algebraic --- by
the original hypotheses.  Bases supported only on $2$ and $3$ do not increase the
transcendence rank.

\subsection{Integrality is stronger than algebraicity, but not by enough}

Four exponentials asks only for algebraicity.  The values here are integers, so there is
extra discrete structure available, and the determinant machinery of
Section~\ref{sec:arbnode} and Section~\ref{sec:semigroup} extracts some of it: the relevant
determinants have integer entries, so a nonzero determinant is a nonzero integer and
therefore satisfies $|\det|\ge1$.

That inequality is the whole of the extra archimedean information.  A nonzero integer has
absolute value at least one, and nothing in the integrality hypotheses improves the constant
$1$ by itself.  To cross the four-exponentials threshold one would need a strictly stronger
arithmetic input, for instance:
\begin{itemize}
\item a proof that the determinants are divisible by a large integer, replacing $|\det|\ge1$
      by $|\det|\ge D$ with $D$ growing in the size parameters; or
\item a proof that a suitable common denominator is anomalously small, improving the
      arithmetic side rather than the analytic side.
\end{itemize}
Both routes ask for divisibility rather than size, which is precisely why the $p$-adic
direction is given priority in the research program of Section~\ref{sec:program}.

\begin{statusbox}{\statusnew{} for the unconditional paper proofs of the determinant hierarchy.
\statuslean{} for the accepted algebraic, integrality, nonvanishing, and positivity layers in
\lean{ArbitraryNodeDeterminant}, \lean{SemigroupDeterminant},
\lean{MixedExponentSemigroup}, and \lean{OrderedChamberPositivity}.  The quantitative Lean
endpoints carry their still-missing calculus inputs as explicit theorem hypotheses rather than
axioms.  The reduction of Conjecture~\ref{conj:AE} to four exponentials is \statuslean{}, but
four exponentials is open and the conjecture remains open with it.}
\end{statusbox}
\section{Rational-function rigidity and closure properties}\label{sec:rigidity}

The exact primitive-generator theorem of Section~\ref{sec:conditional} determines the
additive structure of $\Sol$ completely under failure of Conjecture~\ref{conj:AE}.  It has a
consequence that deserves to be isolated: a hypothetical solution monoid is not merely
additively free of rank two, it is \emph{rigid} against every nonlinear rational operation
over the algebraic numbers.  Addition and multiplication by naturals are the only operations
that stay inside $\Sol$, and every apparent exception collapses to an identity in one
indeterminate.  The results below are exact equivalences, several of them unconditional, and
they translate directly into closure reformulations of the conjecture itself.

\subsection{Evaluation injectivity at the generator}

Throughout this section, $\beta$ denotes the canonical least nonintegral solution supplied by
the generator theorem of Section~\ref{sec:conditional}, so that
\[
 \Sol=\{n+k\beta:n,k\in\Nzero\}
\]
with unique coordinates $n,k$.  The transcendence dichotomy of Section~\ref{sec:corpus} makes
$\beta$ transcendental over $\Q$.  The rational-function statements need the slightly stronger
form over the algebraic closure.

\begin{lemma}[Evaluation is injective over $\Abar$]
\label{rr:lem-eval-inj}
Let $\gamma\in\C$ be transcendental over $\Q$.  Then no nonzero $P\in\Abar[T]$ satisfies
$P(\gamma)=0$; equivalently, the evaluation homomorphism
$\Abar[T]\to\C$, $P\mapsto P(\gamma)$, is injective.
\end{lemma}

\begin{proof}
If $P(\gamma)=0$ for some nonzero $P\in\Abar[T]$, then $\gamma$ is algebraic over $\Abar$.
Since $\Abar/\Q$ is an algebraic extension, the tower property makes $\gamma$ algebraic over
$\Q$, contradicting the hypothesis.  Injectivity follows by applying this to differences.
\end{proof}

\begin{statusbox}{\statusnew{} --- the transcendence of $\beta$ itself is \statuslean{}
(Section~\ref{sec:corpus}, \lean{Transcendence}); the only new content here is the standard
tower argument that upgrades transcendence over $\Q$ to transcendence over $\Abar$.}
\end{statusbox}

\subsection{The rational-function intersection theorem}

\begin{theorem}[Rational-function rigidity at the primitive generator]
\label{rr:thm-rational-rigidity}
Assume Conjecture~\ref{conj:AE} is false and let $\beta$ be the canonical generator.  Let
$R(T)\in\Abar(T)$ have no pole at $\beta$.  Then
\[
 \boxed{\quad
 R(\beta)\in\Sol
 \quad\Longleftrightarrow\quad
 R(T)=n+kT\ \text{in}\ \Abar(T)
 \ \text{for some}\ n,k\in\Nzero.
 \quad}
\]
\end{theorem}

\begin{proof}
The reverse implication is additive closure of $\Sol$ together with closure under
multiplication by naturals: $R(\beta)=n+k\beta\in\Sol$.

For the forward implication write $R=P/Q$ with $P,Q\in\Abar[T]$ and $Q(\beta)\ne0$.  If
$R(\beta)\in\Sol$, the exact coordinates give $R(\beta)=n+k\beta$ for unique
$n,k\in\Nzero$, hence
\[
 P(\beta)-(n+k\beta)\,Q(\beta)=0 .
\]
The left-hand side is the value at $\beta$ of the polynomial
$P(T)-(n+kT)Q(T)\in\Abar[T]$.  By Lemma~\ref{rr:lem-eval-inj} that polynomial is identically
zero, so $P(T)=(n+kT)Q(T)$ and $R(T)=n+kT$ in $\Abar(T)$.
\end{proof}

\begin{statusbox}{\statusnew}
\end{statusbox}

The theorem allows poles away from $\beta$, permits arbitrary cancellation, and takes the full
algebraic closure as coefficient field.  It can be read as the exact intersection
\begin{equation}
 \Sol\cap\Abar(\beta)=\Nzero+\Nzero\beta ,
 \label{rr:eq-intersection}
\end{equation}
with the extra information that every rational expression realizing an element of that
intersection is \emph{identically} the corresponding affine expression --- there is no
accidental representation.

\begin{corollary}[Rational self-maps of the solution set]
\label{rr:cor-selfmaps}
Under failure, the rational functions over $\Abar$ that are regular on $\Sol$ and map $\Sol$
into itself are exactly the maps
\[
 T\longmapsto n+kT,
 \qquad n,k\in\Nzero .
\]
\end{corollary}

\begin{proof}
Such a function in particular sends $\beta$ into $\Sol$, so
Theorem~\ref{rr:thm-rational-rigidity} identifies it as $n+kT$ with $n,k\in\Nzero$.
Conversely, translation by a natural and dilation by a natural preserve
$\Nzero+\Nzero\beta$.
\end{proof}

\begin{statusbox}{\statusnew}
\end{statusbox}

\begin{proposition}[Multivariate reduction to one indeterminate]
\label{rr:prop-multivariate}
Assume failure, let $x_i=n_i+k_i\beta\in\Sol$ for $i=1,\ldots,r$, and let
$F\in\Abar(X_1,\ldots,X_r)$ be regular at $(x_1,\ldots,x_r)$.  Then
\[
 F(x_1,\ldots,x_r)\in\Sol
\]
if and only if the one-variable specialization
\[
 G(T):=F(n_1+k_1T,\ldots,n_r+k_rT)
\]
is identically $p+qT$ in $\Abar(T)$ for some $p,q\in\Nzero$.
\end{proposition}

\begin{proof}
Write $F=P/Q$ with $P,Q\in\Abar[X_1,\ldots,X_r]$ and $Q(x_1,\ldots,x_r)\ne0$.  The
specialized denominator $Q(n_1+k_1T,\ldots,n_r+k_rT)$ is a polynomial in $T$ that does not
vanish at $T=\beta$, hence is not the zero polynomial.  Therefore $G$ is a well-defined
element of $\Abar(T)$ with no pole at $\beta$, and $G(\beta)=F(x_1,\ldots,x_r)$.  Now apply
Theorem~\ref{rr:thm-rational-rigidity}.
\end{proof}

\begin{statusbox}{\statusnew}
\end{statusbox}

Every rational operation on finitely many hypothetical solutions is thus reduced to an
elementary identity between polynomials in a single indeterminate, with coordinates in
$\Nzero$.  No analytic input is required beyond the transcendence already verified by the
kernel.

\subsection{Exact multiplication criterion}

The next statement needs no assumption at all: the branch in which the conjecture holds is
trivial, and the branch in which it fails is governed by the generator.

\begin{theorem}[Products of two solutions]
\label{rr:thm-product}
For all $x,y\in\Sol$, unconditionally,
\[
 \boxed{\quad
 xy\in\Sol
 \quad\Longleftrightarrow\quad
 x\in\Nzero\ \text{or}\ y\in\Nzero .
 \quad}
\]
\end{theorem}

\begin{proof}
If one factor, say $x$, lies in $\Nzero$, then $xy$ is a natural multiple of a solution and
therefore a solution; this direction is unconditional.

For the converse, split on Conjecture~\ref{conj:AE}.  If the conjecture holds then
$\Sol=\Nzero$ and the right-hand side is automatically true, so there is nothing to prove.
Under failure, write the unique coordinates
\[
 x=n+k\beta,
 \qquad
 y=m+\ell\beta ,
 \qquad n,k,m,\ell\in\Nzero .
\]
If neither factor is integral then $k,\ell>0$ and
\[
 xy=nm+(n\ell+mk)\beta+k\ell\,\beta^2 ,
\]
with $k\ell>0$.  Membership $xy\in\Sol$ would force $xy=r+s\beta$ for some $r,s\in\Nzero$,
producing a nonzero polynomial of degree exactly two over $\Q$ vanishing at the
transcendental number $\beta$.  Equivalently, apply Theorem~\ref{rr:thm-rational-rigidity} to
$R(T)=(n+kT)(m+\ell T)$, which is not affine.  Both routes are impossible.
\end{proof}

\begin{statusbox}{\statusnew{} --- note that only transcendence over $\Q$ is used, which is
\statuslean{}; the appeal to $\Abar$ in Theorem~\ref{rr:thm-rational-rigidity} is a
convenience, not a necessity.}
\end{statusbox}

\begin{corollary}[Finite products]
\label{rr:cor-finite-products}
Let $x_1,\ldots,x_r\in\Sol$ be positive.  Then
\[
 \prod_{i=1}^r x_i\in\Sol
 \quad\Longleftrightarrow\quad
 \text{at most one }x_i\text{ is nonintegral} .
\]
If a zero factor is allowed, the product is $0\in\Sol$ and the criterion is vacuous.
\end{corollary}

\begin{proof}
If the conjecture holds, every $x_i$ is integral and both sides are true.  Under failure,
the positive integral factors multiply to a positive natural, so after removing them the
product is a natural multiple of a product of nonintegral solutions.  Writing each remaining
factor as $n_i+k_i\beta$ with $k_i>0$ and expanding, the result is a polynomial in $\beta$
whose degree equals the number of nonintegral factors and whose leading coefficient is
positive.  Degree at least two is excluded by Theorem~\ref{rr:thm-rational-rigidity}, while
degree zero or one is preserved by natural dilation and translation.
\end{proof}

\begin{statusbox}{\statusnew}
\end{statusbox}

\begin{corollary}[Internal power test]
\label{rr:cor-power-test}
For every $x\in\Sol$ and every integer $r\ge2$,
\[
 x^r\in\Sol
 \quad\Longleftrightarrow\quad
 x\in\Nzero .
\]
In particular $x\in\Nzero$ holds exactly when $x^2\in\Sol$: a single squaring test detects
integrality.
\end{corollary}

\begin{proof}
Apply Theorem~\ref{rr:thm-product} with $y=x$ and induct, or apply
Corollary~\ref{rr:cor-finite-products} to $r$ equal factors, noting that $x=0$ satisfies both
sides.
\end{proof}

\begin{statusbox}{\statusnew}
\end{statusbox}

\subsection{Exact quotient criterion}

The same evaluation argument classifies division.  Division is more delicate than
multiplication because even the true solution set $\Nzero$ is not closed under arbitrary
quotients, so the criterion cannot be a clean dichotomy; it is instead an exact divisibility
condition on coordinates.

\begin{theorem}[Quotients of solutions]
\label{rr:thm-quotient}
Assume Conjecture~\ref{conj:AE} is false and write the unique coordinates
\[
 x=n+k\beta,
 \qquad
 y=m+\ell\beta>0 .
\]
Exactly one of the two cases $\ell=0$ and $\ell>0$ applies, and in that case
$x/y\in\Sol$ if and only if the corresponding condition holds:
\begin{enumerate}[label=\textup{(\alph*)}]
\item $\ell=0$, so that $y=m$ is a positive natural, and $m\mid n$ and $m\mid k$;
\item $\ell>0$, and $x=qy$ for some $q\in\Nzero$, equivalently $n=qm$ and $k=q\ell$.
\end{enumerate}
In particular, if $x$ and $y$ are both nonintegral, then
\[
 \frac{x}{y}\in\Sol
 \quad\Longleftrightarrow\quad
 \frac{x}{y}\in\Nzero .
\]
\end{theorem}

\begin{proof}
Suppose $x/y=r+s\beta$ with $r,s\in\Nzero$.  Clearing the denominator gives
\[
 n+k\beta
 =mr+(ms+\ell r)\beta+\ell s\,\beta^2 .
\]
The difference of the two sides is a polynomial in $\beta$ with rational coefficients, so
transcendence of $\beta$ forces every coefficient to vanish:
\[
 \ell s=0,
 \qquad
 k=ms+\ell r,
 \qquad
 n=mr .
\]
If $\ell=0$ then $y=m$ and $y>0$ gives $m\ge1$; the equations read $n=mr$ and $k=ms$, which
is precisely $m\mid n$ and $m\mid k$.  If $\ell>0$ then $s=0$, so $n=mr$ and $k=\ell r$, that
is $x=ry$ with $r\in\Nzero$.  The converses follow by direct substitution.

For the final claim, both $x$ and $y$ nonintegral means $k,\ell>0$, so case~(b) applies and
$x/y=r\in\Nzero$.
\end{proof}

\begin{statusbox}{\statusnew}
\end{statusbox}

\begin{corollary}[Reciprocals]
\label{rr:cor-reciprocal}
Unconditionally, for $x\in\Sol$ with $x>0$,
\[
 \frac1x\in\Sol
 \quad\Longleftrightarrow\quad
 x=1 .
\]
\end{corollary}

\begin{proof}
If the conjecture holds then $\Sol=\Nzero$ and the claim is the elementary fact that a
positive natural with natural reciprocal equals $1$.  Under failure apply
Theorem~\ref{rr:thm-quotient} with numerator $1=1+0\cdot\beta$ and denominator $y=x$.  If
$x=m$ is integral, case~(a) gives $m\mid1$, so $m=1$.  If $x$ is nonintegral, case~(b) gives
$1=qx$ for some $q\in\Nzero$; here $q=0$ is impossible, and $q\ge1$ would make $x=1/q$
rational, contradicting the irrationality of every nonintegral solution.
\end{proof}

\begin{statusbox}{\statusnew}
\end{statusbox}

A hypothetical solution monoid is therefore additively rich --- a free additive monoid of
rank two --- and almost maximally hostile to every multiplicative operation: products,
powers, reciprocals and quotients all escape $\Sol$ except in the degenerate cases forced by
the additive structure itself.

\subsection{Closure formulations of the conjecture}

Because the multiplication criterion is unconditional, it converts directly into equivalent
statements of the conjecture in terms of closure.

\begin{theorem}[Multiplicative closure equivalences]
\label{rr:thm-closure-equiv}
The following are equivalent.
\begin{enumerate}[label=\textup{(\roman*)}]
\item Conjecture~\ref{conj:AE} is true.
\item $\Sol$ is closed under multiplication.
\item $\Sol$ is closed under squaring.
\item For every real $x$,
\[
 2^x,3^x\in\Z
 \quad\Longrightarrow\quad
 2^{x^2},3^{x^2}\in\Z .
\]
\end{enumerate}
\end{theorem}

\begin{proof}
If the conjecture holds then $\Sol=\Nzero$, which is closed under multiplication, so
(i) implies (ii); and (ii) trivially implies (iii).  Statements (iii) and (iv) are the same
assertion, written once for $\Sol$ and once for the defining integrality conditions.
Finally, assume (iii).  For any $x\in\Sol$ we get $x^2\in\Sol$, so
Corollary~\ref{rr:cor-power-test} forces $x\in\Nzero$.  Hence $\Sol=\Nzero$, which is
Conjecture~\ref{conj:AE}.
\end{proof}

\begin{statusbox}{\statusnew}
\end{statusbox}

The candidate formulation of Section~\ref{sec:problem} carries the same content in
multiplicative language.  For positive reals define the commutative operation
\begin{equation}
 m\star n
 :=2^{(\log_2m)(\log_2n)}
 =m^{\log_2n}
 =n^{\log_2m} .
 \label{rr:eq-star}
\end{equation}
If $m=2^x$ and $n=2^y$, then $m\star n=2^{xy}$, so $\star$ is exactly the transport of
multiplication of exponents through the exponent--candidate bijection of
Section~\ref{sec:problem}.

\begin{corollary}[Candidate star-product]
\label{rr:cor-star}
For $m,n\in\Cand$,
\[
 m\star n\in\Cand
 \quad\Longleftrightarrow\quad
 m\ \text{is a power of }2
 \ \text{or}\
 n\ \text{is a power of }2 .
\]
Consequently Conjecture~\ref{conj:AE} is equivalent both to closure of $\Cand$ under $\star$
and to the apparently weaker diagonal condition
\[
 m\in\Cand
 \quad\Longrightarrow\quad
 m\star m\in\Cand .
\]
\end{corollary}

\begin{proof}
Put $x=\log_2m$ and $y=\log_2n$, so that $x,y\in\Sol$ and $m\star n=2^{xy}$.  Membership
$m\star n\in\Cand$ is equivalent to $xy\in\Sol$, and $m$ is a power of two exactly when
$x\in\Nzero$; now apply Theorem~\ref{rr:thm-product}.  The two closure equivalences are the
image of Theorem~\ref{rr:thm-closure-equiv}(ii) and (iii) under the same bijection.
\end{proof}

\begin{statusbox}{\statusnew}
\end{statusbox}

These are genuine equivalences and not a proof.  There is at present no independent mechanism
forcing nonlinear closure: nothing in the archimedean, valuation-theoretic, or determinant
toolkits of Sections~\ref{sec:corpus}, \ref{sec:arbnode} and \ref{sec:semigroup} supplies a
reason for $\Cand$ to be closed under $\star$.  Their value is diagnostic --- they identify
the missing property as a single crisp failure of semiring-like structure, rather than as a
quantitative deficit.

\subsection{Formalization cost}

The multiplication, power and quotient criteria are inexpensive Lean targets, and a route
exists that avoids building a general rational-function library.  For a new module ---
\lean{RationalRigidity} would be a natural name --- the plan is:
\begin{enumerate}
\item case split on \lean{AlaogluErdosConjecture}, discharging the true branch by
$\Sol=\Nzero$;
\item in the failure branch obtain the canonical generator and the unique coordinates from
the packaged statement
\begin{center}
\lean{exists_canonical_primitiveGenerator_of_not_alaogluErdosConjecture}
\end{center}
in the module \lean{PrimitiveGenerator};
\item expand the relevant equality as a polynomial of degree at most two in $\beta$ with
rational coefficients;
\item invoke the kernel-verified transcendence in \lean{Transcendence} to force every
coefficient to vanish;
\item discharge the remaining natural-number divisibility and coordinate conditions with
\lean{omega}.
\end{enumerate}
Only the full rational-function theorem, Theorem~\ref{rr:thm-rational-rigidity}, needs
evaluation injectivity over $\Abar$ and hence Lemma~\ref{rr:lem-eval-inj}; the multiplication
theorem, the power test, the quotient criterion and the closure equivalences need nothing
beyond transcendence over $\Q$, which is already accepted by the kernel.

\begin{statusbox}{Planned; no draft yet.  The mathematics is \statusnew{} throughout this
section, and the prerequisites are \statuslean{}.}
\end{statusbox}

\subsection{What the rigidity statements buy}

The practical payoff is that they sharply restrict how a counterexample could be manufactured
from a smaller one.  A natural strategy for producing new solutions from old ones is to apply
some algebraic operation --- square a solution, multiply two of them, take a ratio, invert,
or substitute solutions into a rational function --- and hope to land back in $\Sol$, either
to run a descent or to build an infinite family with better arithmetic properties than the
canonical generator.  Equation~\eqref{rr:eq-intersection} closes that door completely: the
only rational functions over $\Abar$ that map solutions to solutions are the affine maps
$T\mapsto n+kT$ with natural coefficients, which reproduce nothing that the additive monoid
$\Nzero+\Nzero\beta$ did not already contain.  Any construction of a new counterexample from
an old one must therefore be transcendental in an essential way, or must leave the solution
set and return by a route that is not rational over the algebraic numbers.  In the other
direction, the closure equivalences of Theorem~\ref{rr:thm-closure-equiv} and
Corollary~\ref{rr:cor-star} convert the conjecture into a purely structural question --- is
$\Cand$ closed under $\star$? --- which is attractive precisely because it has no
quantitative parameter to optimize, and which is a warning that no counting or
equidistribution bound alone will produce the answer.  Finally, all of these statements are
short, elementary consequences of results already accepted by the kernel, which makes them
the cheapest genuine additions to the verified corpus that this report proposes.

\section{Three-smooth outputs and transcendental product exponents}\label{sec:smooth}

Call a positive integer \emph{three-smooth} when it has no prime factor outside
$\{2,3\}$, that is, when it equals $2^u3^v$ for some $u,v\in\Nzero$.  The kernel-verified
classification of auxiliary bases recorded in Section~\ref{sec:corpus} settles which bases
$B$ satisfy $B^x\in\Z$ for a \emph{fixed} nonintegral solution $x$: exactly the three-smooth
ones.  This section turns that classification onto the outputs $2^x,3^x$ themselves, and
then onto second-level exponents $xy$ formed from two solutions.

Two distinct things come out, and they have different provenance.  The first is an
unconditional integrality test: a solution is integral precisely when its own two outputs
are three-smooth.  It uses only the kernel-verified transcendence of $\thetaq$ and unique
factorization, so it is cheap to add to the Lean development.  The second is an upgrade of
``not an integer'' to ``transcendental'', obtained from the six exponentials theorem.  That
upgrade converts several apparently weaker \emph{algebraicity} statements into exact
reformulations of Conjecture~\ref{conj:AE}, at the cost of a classical input that the
repository has not formalized in general form.

\subsection{The two outputs cannot both be three-smooth}

\begin{theorem}[Three-smooth output test]
\label{sm:output-smooth}
For every $x\in\Sol$ the following are equivalent:
\begin{enumerate}[label=\textup{(\roman*)}]
\item $x\in\Nzero$;
\item both $2^x$ and $3^x$ are three-smooth;
\item $6^x$ is three-smooth.
\end{enumerate}
\end{theorem}

\begin{proof}
Condition (i) gives (ii) immediately, since $2^x$ and $3^x$ are then powers of $2$ and $3$
respectively.  Conditions (ii) and (iii) are equivalent because $2^x$ and $3^x$ are positive
integers with product $6^x$, and a product of positive integers is three-smooth exactly when
both factors are.

For (ii)$\Rightarrow$(i), write
\[
 2^x=2^p3^q,\qquad 3^x=2^r3^s
\]
with $p,q,r,s\in\Nzero$.  Taking logarithms and dividing by $\log2$ and $\log3$ respectively
gives
\[
 x=p+q\thetaq,\qquad x=s+\frac{r}{\thetaq},
\]
where $\thetaq=\log_23$.  Eliminating $x$ and clearing the denominator yields the quadratic
relation
\begin{equation}
 q\,\thetaq^{2}+(p-s)\,\thetaq-r=0
 \label{sm:eq-quadratic}
\end{equation}
with rational --- indeed integer --- coefficients.  The repository proves $\thetaq$
transcendental, so every coefficient in \eqref{sm:eq-quadratic} vanishes: $q=0$, $p=s$, and
$r=0$.  Hence $2^x=2^p$ and $x=p\in\Nzero$.
\end{proof}

\begin{statusbox}{\statusnew{}}
The statement is unconditional and does not use the conditional generator, the exact solution
monoid, or any six-exponentials input.  Its only nonelementary ingredient is the
kernel-verified transcendence of $\thetaq$,
\begin{center}
\lean{transcendental_logThreeDivLogTwo}
\end{center}
in \lean{Transcendence}, together with unique factorization.  A Lean proof is expected to be
short; see Section~\ref{sec:blueprint} and the skeleton at the end of this section.
\end{statusbox}

The test is sharp in a useful direction: it rules out one of the four a priori smoothness
patterns of the least output pair.

\begin{corollary}[Primitive-core trichotomy]
\label{sm:core-trichotomy}
Assume Conjecture~\ref{conj:AE} fails, and use the simultaneous least-output normalization of
Section~\ref{sec:conditional},
\[
 M=2^{\beta}=2^aw^d,\qquad A=3^{\beta}=3^bv^e,
\]
with $w>1$ odd and power-primitive, $v>1$ power-primitive and coprime to $3$, and $d,e>0$.
Then at most one of $w=3$ and $v=2$ holds, and exactly one of the following three patterns
occurs:
\begin{center}
\begin{tabular}{@{}p{0.24\textwidth}p{0.24\textwidth}p{0.32\textwidth}@{}}
\toprule
$M$ & $A$ & primitive cores\\
\midrule
three-smooth & not three-smooth & $w=3$, $v\ne2$\\
not three-smooth & three-smooth & $w\ne3$, $v=2$\\
not three-smooth & not three-smooth & $w\ne3$, $v\ne2$\\
\bottomrule
\end{tabular}
\end{center}
The missing fourth pattern, both outputs three-smooth, is impossible.
\end{corollary}

\begin{proof}
An odd, three-smooth, power-primitive integer greater than one is a power of $3$, and
primitivity forces the exponent to be $1$; hence $M$ is three-smooth exactly when $w=3$.
Symmetrically a three-smooth power-primitive integer greater than one that is coprime to $3$
is exactly $2$, so $A$ is three-smooth exactly when $v=2$.  The least nonintegral solution
$\beta$ lies in $\Sol\setminus\Nzero$, so Theorem~\ref{sm:output-smooth} forbids both
outputs from being three-smooth, which is the exclusion $\lnot(w=3\wedge v=2)$.
\end{proof}

\begin{statusbox}{\statusnew{}}
The normalization itself, including power-primitivity of $w$ and $v$, the depth dichotomy
$a=0$ or $b=0$, and the gcd-one condition, is \statuslean{} in \lean{OutputNormalization}
(Section~\ref{sec:conditional}).  Only the exclusion of the fourth pattern is new here.
\end{statusbox}

The two exceptional rows are precisely the cases the corpus already identifies as hardest:
$w=3$ makes $3^{\thetaq}$ the sharp radical target of Section~\ref{sec:radical}, and $v=2$ is
its base-swapped mirror with $2^{1/\thetaq}=2^{\log_32}$.  The third row is generic and
carries two independent outside-prime supports.

\subsection{Third-base transcendence from six exponentials}

The remaining statements of this section rest on one published transcendence theorem, which
we state in the form we use.

\begin{theorem}[Six exponentials theorem; Siegel, Lang, Ramachandra]
\label{sm:six-exp}
Let $u_1,u_2$ be complex numbers linearly independent over $\Q$, and let $v_1,v_2,v_3$ be
complex numbers linearly independent over $\Q$.  Then at least one of the six numbers
$e^{u_iv_j}$, $1\le i\le2$, $1\le j\le3$, is transcendental.
\end{theorem}

\begin{statusbox}{\statusknown{} --- classical, not formalized in this repository in general
form.}
See Waldschmidt \cite{Waldschmidt2000,Waldschmidt2005}.  Every consequence below is tagged
\statusknown{} for the same reason: the repository formalizes only the integer-entry
specialization of the underlying interpolation determinant, which yields the $3$-smooth
classification of Section~\ref{sec:corpus} but not the transcendence conclusion.
\end{statusbox}

\begin{lemma}[Third-base transcendence]
\label{sm:third-base-trans}
Let $y\in\Sol\setminus\Nzero$ and let $B\in\N_{>0}$ be not three-smooth.  Then $B^y$ is
transcendental.
\end{lemma}

\begin{proof}
The numbers $\log2,\log3,\log B$ are linearly independent over $\Q$: a nontrivial rational
relation among them would give a nontrivial multiplicative relation among $2$, $3$, and $B$,
which unique factorization forbids because $B$ has a prime factor outside $\{2,3\}$.  Also
$1$ and $y$ are linearly independent over $\Q$, since $y$ is irrational --- a rational
solution is integral by the kernel-verified rational-exponent lemma of
Section~\ref{sec:corpus}, and $y\notin\Nzero$.

Apply Theorem~\ref{sm:six-exp} with $u_1=1$, $u_2=y$ and $v_1=\log2$, $v_2=\log3$,
$v_3=\log B$.  The six values are
\[
 2,\quad 3,\quad B,\quad 2^y,\quad 3^y,\quad B^y,
\]
and the first five are integers, hence algebraic: $2,3,B$ by hypothesis and $2^y,3^y$
because $y\in\Sol$.  Therefore the sixth, $B^y$, is transcendental.
\end{proof}

\begin{statusbox}{\statusknown{} --- consequence of the six exponentials theorem
(Theorem~\ref{sm:six-exp}).}
\end{statusbox}

Combining Lemma~\ref{sm:third-base-trans} with the kernel-verified $3$-smooth classification
of Section~\ref{sec:corpus} gives an exact algebraicity trichotomy for integer bases: for
every $y\in\Sol\setminus\Nzero$ and every $B\in\N_{>0}$,
\begin{equation}
 B^y\in\Abar
 \quad\Longleftrightarrow\quad
 B^y\in\Z
 \quad\Longleftrightarrow\quad
 B=2^u3^v\ \text{ for some }u,v\in\Nzero .
 \label{sm:eq-trichotomy}
\end{equation}
There is no algebraic-irrational middle case.  Any base that is not three-smooth produces a
transcendental value outright.

\subsection{Products of noninteger solutions}

The set $\Sol$ is closed under addition and under multiplication by natural numbers, but it
is not known to be closed under multiplication of two of its own nonintegral elements.  The
next theorem says exactly what happens at the product exponent $xy$.

\begin{theorem}[Product-exponent dichotomy]
\label{sm:product-trans}
Let $x,y\in\Sol$.
\begin{enumerate}[label=\textup{(\alph*)}]
\item If $x\in\Nzero$ or $y\in\Nzero$, then $2^{xy}\in\N$ and $3^{xy}\in\N$.
\item If $x\notin\Nzero$ and $y\notin\Nzero$, then each of $2^{xy}$ and $3^{xy}$ is either an
integer or transcendental, at least one of the two is transcendental, and
\[
 6^{xy}\ \text{is transcendental}.
\]
\end{enumerate}
\end{theorem}

\begin{proof}
Part (a) is natural-number dilation: if $y=n\in\Nzero$ then $xy=nx$ is a natural multiple of
a solution and therefore a solution, and symmetrically in $x$.

For part (b) put $M_x=2^x$ and $A_x=3^x$, both positive integers, and note
\[
 2^{xy}=M_x^{\,y},\qquad 3^{xy}=A_x^{\,y}.
\]
If $M_x$ is three-smooth, the $3$-smooth classification of Section~\ref{sec:corpus} makes
$M_x^{\,y}$ an integer; if $M_x$ is not three-smooth, Lemma~\ref{sm:third-base-trans} makes
it transcendental.  The same alternative holds for $A_x^{\,y}$.  By
Theorem~\ref{sm:output-smooth} applied to the nonintegral solution $x$, the two outputs
$M_x$ and $A_x$ are not both three-smooth, so at least one of $2^{xy},3^{xy}$ is
transcendental.

Finally $6^x=M_xA_x$ is a positive integer, and it is not three-smooth, since a prime outside
$\{2,3\}$ dividing one factor divides the product.  Applying
Lemma~\ref{sm:third-base-trans} with base $6^x$ and exponent $y$ gives
$(6^x)^y=6^{xy}$ transcendental.
\end{proof}

\begin{statusbox}{\statusknown{} for part (b) --- via Theorem~\ref{sm:six-exp}.}
Part (a) is elementary and its underlying closure property, that natural multiples of a
solution are solutions, is \statuslean{}.
\end{statusbox}

\begin{corollary}[Diagonal transcendence]
\label{sm:diagonal-trans}
Let $x\in\Sol\setminus\Nzero$.  Then $6^{x^2}$ is transcendental and at least one of
$2^{x^2}$, $3^{x^2}$ is transcendental.  More precisely,
\[
\begin{aligned}
 2^{x^2}\in\N &\iff 2^x\ \text{is three-smooth},\\
 3^{x^2}\in\N &\iff 3^x\ \text{is three-smooth},
\end{aligned}
\]
and failure of either equivalence's right-hand side makes the corresponding value
transcendental.
\end{corollary}

\begin{proof}
Take $y=x$ in Theorem~\ref{sm:product-trans}(b).  The two displayed equivalences are the
alternative used in that proof, applied to the bases $M_x=2^x$ and $A_x=3^x$: a three-smooth
base gives an integer by the $3$-smooth classification of Section~\ref{sec:corpus}, and a
base that is not three-smooth gives a transcendental value by
Lemma~\ref{sm:third-base-trans}.
\end{proof}

\begin{corollary}[Mixed three-smooth bases]
\label{sm:mixed-base-trans}
Let $x,y\in\Sol\setminus\Nzero$ and let $u,v\ge1$ be integers.  Then $(2^u3^v)^{xy}$ is
transcendental.
\end{corollary}

\begin{proof}
The positive integer
\[
 (2^u3^v)^x=(2^x)^u(3^x)^v
\]
is divisible by every prime occurring in either output, because $u,v\ge1$.  By
Theorem~\ref{sm:output-smooth} at least one output has a prime factor outside $\{2,3\}$, so
this base is not three-smooth.  Apply Lemma~\ref{sm:third-base-trans} with exponent $y$.
\end{proof}

\begin{statusbox}{\statusknown{} for both corollaries --- via Theorem~\ref{sm:six-exp}.}
\end{statusbox}

\subsection{A global second-level trichotomy}

The exact generator theorem of Section~\ref{sec:conditional} makes the preceding dichotomy
uniform: which of the three patterns occurs is a property of the counterexample as a whole,
not of the particular pair $(x,y)$.

\begin{theorem}[Uniform second-level pattern]
\label{sm:uniform-second-level}
Assume Conjecture~\ref{conj:AE} fails and let $M=2^{\beta}$, $A=3^{\beta}$ for the primitive
generator $\beta$.  For every $x,y\in\Sol\setminus\Nzero$, exactly one of the following three
patterns holds, and the pattern does not depend on $x$ or $y$:
\begin{center}
\begin{tabular}{@{}p{0.36\textwidth}p{0.20\textwidth}p{0.24\textwidth}@{}}
\toprule
condition on $(M,A)$ & $2^{xy}$ & $3^{xy}$\\
\midrule
$M$ three-smooth, $A$ not & integer & transcendental\\
$A$ three-smooth, $M$ not & transcendental & integer\\
neither three-smooth & transcendental & transcendental\\
\bottomrule
\end{tabular}
\end{center}
\end{theorem}

\begin{proof}
By the exact primitive-generator theorem of Section~\ref{sec:conditional}, write
$x=n+k\beta$ with $n,k\in\Nzero$ and $k\ge1$, so that
\[
 2^x=2^nM^k,\qquad 3^x=3^nA^k .
\]
Multiplying by a power of the distinguished smooth prime cannot remove a prime factor
outside $\{2,3\}$, and $k\ge1$ keeps every such prime of $M$ present in $M^k$; hence $2^x$ is
three-smooth exactly when $M$ is, and likewise $3^x$ is three-smooth exactly when $A$ is.
Now apply Theorem~\ref{sm:product-trans}(b) together with the alternative used in its proof.
The fourth pattern, both $M$ and $A$ three-smooth, is excluded by
Theorem~\ref{sm:output-smooth} applied to $\beta$.
\end{proof}

\begin{statusbox}{\statusknown{} --- via Theorem~\ref{sm:six-exp}.}
The structural input, namely $\Sol=\{n+k\beta:n,k\in\Nzero\}$ with unique coordinates, is
\statuslean{} in \lean{PrimitiveGenerator} (Section~\ref{sec:conditional}).
\end{statusbox}

This suggests a clean three-case research split, matching
Corollary~\ref{sm:core-trichotomy}.  The two one-smooth-output cases reduce to the
exceptional primitive cores $w=3$ and $v=2$, where the target radicals are $3^{\thetaq}$ and
$2^{1/\thetaq}$ respectively and where Section~\ref{sec:radical} already locates the wall.
In the generic case both second iterates are transcendental and two independent
outside-prime supports are available, which is the configuration a mixed archimedean and
non-archimedean determinant would exploit; see Section~\ref{sec:program}.

\subsection{Algebraicity formulations equivalent to the conjecture}

\begin{theorem}[Algebraicity tests]
\label{sm:algebraicity}
Each of the following is equivalent to Conjecture~\ref{conj:AE}.
\begin{enumerate}[label=\textup{(\roman*)}]
\item For every $x\in\Sol$, $6^{x^2}\in\Abar$.
\item For every $x\in\Sol$, both $2^{x^2}\in\Abar$ and $3^{x^2}\in\Abar$.
\item For every $x,y\in\Sol$, $6^{xy}\in\Abar$.
\item For one fixed pair of integers $u,v\ge1$, every $x\in\Sol$ satisfies
$(2^u3^v)^{x^2}\in\Abar$.
\item For every $x\in\Sol$, the integer $6^x$ is three-smooth.
\end{enumerate}
\end{theorem}

\begin{proof}
If Conjecture~\ref{conj:AE} holds, then $\Sol=\Nzero$, every exponent appearing in
(i)--(iv) is a nonnegative integer, and all the listed values are positive integers, hence
algebraic; condition (v) is immediate because $6^x=2^x3^x$ is then a power of $6$.

Conversely, suppose the conjecture fails and fix $x\in\Sol\setminus\Nzero$.  Then (i) and
(ii) fail by Corollary~\ref{sm:diagonal-trans}, (iii) fails by taking $y=x$ in
Theorem~\ref{sm:product-trans}(b), (iv) fails by Corollary~\ref{sm:mixed-base-trans} with
$y=x$, and (v) fails by Theorem~\ref{sm:output-smooth}.
\end{proof}

\begin{statusbox}{\statusknown{} for (i)--(iv) --- via Theorem~\ref{sm:six-exp}.
\statusnew{} for (v).}
Equivalence (v) needs only Theorem~\ref{sm:output-smooth}, hence only the kernel-verified
transcendence of $\thetaq$; it is therefore the one item in this list that is within
immediate reach of the Lean development.
\end{statusbox}

Items (i)--(iv) replace integrality by the apparently much weaker target of algebraicity, and
the six exponentials theorem makes the replacement lossless.  That sharpens the statement of
what is missing: one would need an independent reason why some nonlinear iterate such as
$6^{x^2}$ is algebraic, derived from the two original integer values $2^x$ and $3^x$ alone.
No mechanism in the present corpus produces algebraicity at the second level, which is
consistent with the four-exponentials diagnosis of Section~\ref{sec:fourexp}: the two-base
configuration supplies only two independent logarithmic directions, and squaring the exponent
does not create a third.

\subsection{Formalization boundary}

The division of labour in this section is unusually clean, and it marks exactly where the
repository can advance without new transcendence theory.

Theorem~\ref{sm:output-smooth}, Corollary~\ref{sm:core-trichotomy}, part (a) of
Theorem~\ref{sm:product-trans}, and equivalence (v) of Theorem~\ref{sm:algebraicity} use only
kernel-verified ingredients: transcendence of $\thetaq$, unique factorization, the closure of
$\Sol$ under natural multiples, the simultaneous output normalization, and the $3$-smooth
classification in \lean{ThreeSmooth}.  A proposed module \lean{OutputSmoothness} would host
them.  With the predicate
\[
 \operatorname{ThreeSmooth}(m):\Longleftrightarrow\exists\,u,v\in\Nzero,\ m=2^u3^v,
\]
the headline statement is a biconditional between integrality of the exponent and joint
smoothness of the two outputs:

\begin{lstlisting}[style=leanstyle,caption={Proposed output-smoothness signature.}]
theorem integer_iff_outputs_threeSmooth
    {x : R} (hx : TwoBaseIntegralSolution x) :
    x in Set.range ((^) : Z -> R) <->
      ThreeSmooth (natOutputTwo hx) /\
      ThreeSmooth (natOutputThree hx)
\end{lstlisting}

The reverse direction is the interesting one, and it follows the proof of
Theorem~\ref{sm:output-smooth}: take logarithms of the two smooth factorizations, assemble
\eqref{sm:eq-quadratic} as a polynomial of degree at most two in $\thetaq$ with integer
coefficients, apply
\begin{center}
\lean{transcendental_logThreeDivLogTwo}
\end{center}
to force every coefficient to vanish, and discharge the remaining natural-number arithmetic
with \lean{omega}.  Mathlib's polynomial evaluation interface is the natural way to build
$qX^2+(p-s)X-r$ and invoke the definition of transcendence.

Everything else in this section --- Lemma~\ref{sm:third-base-trans}, part (b) of
Theorem~\ref{sm:product-trans}, Corollaries~\ref{sm:diagonal-trans} and
\ref{sm:mixed-base-trans}, Theorem~\ref{sm:uniform-second-level}, and items (i)--(iv) of
Theorem~\ref{sm:algebraicity} --- needs the six exponentials theorem in general form.  The
repository's interpolation determinant is currently specialized to integer entries: a
candidate matrix has integer entries, a nonzero integer determinant satisfies $|D|\ge1$, and
that inequality is what drives the $3$-smooth classification.  The natural formal route to
Lemma~\ref{sm:third-base-trans} is to replace the bound $|D|\ge1$ by a number-field norm
bound for a nonzero \emph{algebraic} determinant, retaining the same analytic upper bound
from the interpolation estimate.  Packaging that as an ``algebraic entries of bounded
height'' interface would recover the transcendence conclusion while reusing the existing
determinant infrastructure, and would place the general six exponentials theorem within the
same framework.  This is a substantial addition, not a refactor, and it is listed among the
transcendence-theoretic priorities of Section~\ref{sec:program}.

\section{Perfect-power content, affine descent, and primitive-output statistics}\label{sec:powers}

Section~\ref{sec:conditional} closes the conditional structure question: if
Conjecture~\ref{conj:AE} fails, then $\Sol$ is the free rank-two additive monoid generated by
$1$ and the canonical least nonintegral solution $\beta$, and the least output pair
$(M,A)=(2^{\beta},3^{\beta})$ is affinely primitive.  Both facts are \statuslean{}.  What that
corpus does not yet record is the \emph{exact} perfect-power content of the output pair at an
arbitrary point of the monoid: the kernel-verified corollaries exclude matched decompositions
at $\beta$, but say nothing about which matched decompositions actually occur elsewhere.

This section supplies that classification.  It determines exactly which simultaneous
factorizations
\[
 2^x=2^rU^d,\qquad 3^x=3^rV^d
\]
can occur for $x\in\Sol$, identifies the maximal common perfect-power degree of the output
pair with a gcd of coordinates, and then converts the classification into exact limiting
statistics for primitive output pairs.  Everything here is conditional on failure and rests on
the kernel-verified monoid classification; the classification and the statistics themselves are
\statusnew{}.

Throughout, assume failure and fix the canonical generator $\beta$ of
Section~\ref{sec:conditional}, with
\[
 M=2^{\beta},\qquad A=3^{\beta}.
\]
For $x=n+k\beta$ with $n,k\in\Nzero$, exact coordinates give
\begin{equation}
 2^x=2^nM^k,
 \qquad
 3^x=3^nA^k,
\label{pp:eq-output-coordinates}
\end{equation}
and the pair $(n,k)$ is uniquely determined by $x$.  The coordinate gcd $\gcd(n,k)$ turns out
to be an intrinsic arithmetic invariant of the output pair, not merely of its coordinates.

\subsection{Simultaneous perfect powers}

\begin{theorem}[Exact common-power criterion]
\label{pp:common-power}
Let $x=n+k\beta\in\Sol$ and let $r\ge1$.  The following are equivalent:
\begin{enumerate}[label=\textup{(\roman*)}]
\item both $2^x$ and $3^x$ are perfect $r$-th powers in $\N$;
\item $x/r\in\Sol$;
\item $r\mid n$ and $r\mid k$.
\end{enumerate}
\end{theorem}

\begin{proof}
If $2^x=U^r$ and $3^x=V^r$ with $U,V>0$, uniqueness of positive real $r$-th roots gives
\[
 2^{x/r}=U,\qquad 3^{x/r}=V,
\]
so $x/r\in\Sol$.  The converse follows by raising the two integer values at $x/r$ to the
$r$-th power.

For the coordinate form, write $x/r=p+q\beta$ in the unique coordinates supplied by the exact
primitive-generator theorem of Section~\ref{sec:conditional}.  Then $n+k\beta=rp+rq\beta$, and
uniqueness of coordinates forces $n=rp$ and $k=rq$, which is exactly the pair of divisibility
conditions.  Conversely, if $r\mid n$ and $r\mid k$, then $(n/r)+(k/r)\beta\in\Sol$ and equals
$x/r$.
\end{proof}

For $x>0$ define the \emph{common output power degree}
\[
 g(x):=\max\{r\ge1:\text{$2^x$ and $3^x$ are both perfect $r$-th powers}\}.
\]

\begin{corollary}[Coordinate gcd]
\label{pp:gcd-degree}
If $x=n+k\beta>0$, then
\[
 g(x)=\gcd(n,k),
\]
with the convention $\gcd(n,0)=n$.  The output pair is simultaneously power-primitive exactly
when $\gcd(n,k)=1$.
\end{corollary}

\begin{proof}
By Theorem~\ref{pp:common-power} the admissible degrees $r$ are precisely the common divisors
of $n$ and $k$, and the set of common divisors of a pair that is not $(0,0)$ has a largest
element, namely $\gcd(n,k)$.
\end{proof}

Consequently every nonzero output pair has a unique decomposition
\[
 (2^x,3^x)=(U^{g},V^{g}),\qquad g=\gcd(n,k),
\]
where $(U,V)=(2^{x/g},3^{x/g})$ is a simultaneously power-primitive output pair belonging to
the solution $x/g$.  The repository proves power-primitivity for the least pair only; this
corollary computes the intrinsic power degree at every point of the solution monoid.  For
integral $x=n$ it returns $g(n)=n$, which is the familiar statement that $2^n$ and $3^n$ are
perfect $n$-th powers and no better; for $x=\beta$ it returns $g(\beta)=\gcd(0,1)=1$, which is
exactly the kernel-verified exclusion of a common perfect-power degree at the least pair.

\begin{statusbox}{\statusnew{} --- Theorem~\ref{pp:common-power} and
Corollary~\ref{pp:gcd-degree} are paper results of this report.  Their inputs are
kernel-verified: the unique-coordinate clause of the exact primitive-generator theorem
(\lean{PrimitiveGenerator}) and the least-solution corollaries of \lean{LeastSolution}.  The
implication (i)$\Rightarrow$(ii) is the special case $r=0$ of the Lean affine descent theorem
\lean{affine_descent_integral_powers} in \lean{ArithmeticRigidity}; the converse and the
coordinate criterion (iii) are new here.}
\end{statusbox}

\subsection{Matched affine decompositions}

A descent that strips the same distinguished-base depth and the same outer power from both
outputs is the natural common generalization of the two operations above.  The repository's
affine descent theorem shows that such a stripping produces a solution; the following criterion
says exactly when one exists.

\begin{theorem}[Matched affine decomposition criterion]
\label{pp:affine-decomp}
Let $x=n+k\beta\in\Sol$, let $d\ge1$, and let $r\in\Nzero$.  The following are equivalent:
\begin{enumerate}[label=\textup{(\roman*)}]
\item there exist $U,V\in\N_{>0}$ with
\[
 2^x=2^rU^d,\qquad 3^x=3^rV^d;
\]
\item $(x-r)/d\in\Sol$;
\item $0\le r\le n$, $d\mid k$, and $d\mid(n-r)$.
\end{enumerate}
When these conditions hold,
\[
 \frac{x-r}{d}=\frac{n-r}{d}+\frac kd\,\beta,
 \qquad
 U=2^{(x-r)/d},\qquad V=3^{(x-r)/d}.
\]
\end{theorem}

\begin{proof}
Conditions (i) and (ii) are equivalent by uniqueness of positive real roots, applied after
removing the common factor $2^r$ from the first output and $3^r$ from the second.  Expanding
(ii) in unique coordinates gives
\[
 n+k\beta=r+d(p+q\beta)
\]
for some $p,q\in\Nzero$, so $n=r+dp$ and $k=dq$; this is exactly (iii), and it also yields the
displayed formulas.  Conversely, (iii) lets one define $p=(n-r)/d\in\Nzero$ and
$q=k/d\in\Nzero$, and then $p+q\beta\in\Sol$ equals $(x-r)/d$.
\end{proof}

At the least pair the criterion collapses.  For $x=\beta$ we have $n=0$ and $k=1$, so (iii)
forces $r=0$ and $d=1$: no nontrivial matched decomposition of $(M,A)$ exists.  That is the
kernel-verified exact affine primitivity statement of Section~\ref{sec:conditional}, recovered
here as one instance of a criterion valid at every point.

For a nonintegral solution $x=n+k\beta$ with $k>0$, Euclidean division $n=qk+r$ with
$0\le r<k$ gives a canonical reduced affine presentation
\begin{equation}
 x=r+k(q+\beta),
\label{pp:eq-euclidean-affine}
\end{equation}
and hence
\[
 2^x=2^r\bigl(2^qM\bigr)^k,
 \qquad
 3^x=3^r\bigl(3^qA\bigr)^k .
\]
The normalization $0\le r<k$ makes the presentation unique.  It separates the outer affine
degree $k$ --- an invariant of the point, being its $\beta$-coordinate --- from the reduced
common base depth $r$.

\begin{statusbox}{\statusnew{} --- Theorem~\ref{pp:affine-decomp} and the reduced presentation
\eqref{pp:eq-euclidean-affine} are paper results.  The implication (i)$\Rightarrow$(ii) is the
kernel-verified affine descent theorem; the equivalence with the coordinate condition (iii),
and hence the classification of \emph{all} matched decompositions, is new.}
\end{statusbox}

\subsection{What is already formal, and what is new here}

The table below separates the two layers.  The middle column is the kernel-verified statement
about the least pair $(M,A)$; the right column is what this section adds for a general point
$x=n+k\beta$ of the monoid.

\begin{center}
\begin{tabular}{@{}>{\raggedright\arraybackslash}p{0.26\textwidth}
                   >{\raggedright\arraybackslash}p{0.30\textwidth}
                   >{\raggedright\arraybackslash}p{0.34\textwidth}@{}}
\toprule
Statement & Kernel-verified at $\beta$ & Extension in this section \\
\midrule
Matched base depth &
$\min\{\vpp{2}{M},\vpp{3}{A}\}=0$, in \lean{LeastSolution} &
Admissible depths at $x$ are exactly $0\le r\le n$ with $d\mid(n-r)$ \\
\addlinespace
Common perfect-power degree &
no $k\ge2$ with $M,A$ both $k$-th powers, in \lean{LeastSolution} &
the exact degree is $g(x)=\gcd(n,k)$
(Corollary~\ref{pp:gcd-degree}) \\
\addlinespace
Affine primitivity &
$M=2^rU^k$, $A=3^rV^k$ force $r=0$, $k=1$, in \lean{LeastSolution} &
full criterion for every $x$, with the least pair as the case $(n,k)=(0,1)$
(Theorem~\ref{pp:affine-decomp}) \\
\addlinespace
Affine descent &
$(x-r)/k\in\Sol$ from a matched factorization, in \lean{ArithmeticRigidity} &
the descent condition is an equivalence, not merely an implication \\
\addlinespace
Output normalization &
$M=2^aw^d$, $A=3^bv^e$ with $a=0$ or $b=0$ and
$\gcd(\gcd(d,e),|b-a|)=1$, in \lean{OutputNormalization} &
unchanged; that invariant constrains the two \emph{separate} degrees, whereas $g$ measures
only the common one \\
\bottomrule
\end{tabular}
\end{center}

The kernel-verified identifiers behind the middle column are

\begin{center}
\lean{IsLeastTwoBaseNonintegerSolution.not_both_base_dvd_outputs}
\end{center}
\begin{center}
\lean{IsLeastTwoBaseNonintegerSolution.no_common_perfect_power}
\end{center}
\begin{center}
\lean{IsLeastTwoBaseNonintegerSolution.affine_primitive}
\end{center}

in \lean{LeastSolution}, together with
\lean{exists_simultaneous_primitive_output_normalization} in \lean{OutputNormalization} and
\lean{affine_descent_integral_powers} in \lean{ArithmeticRigidity}.  None of the right-hand
column is machine-checked.  A formalization would be light: the coordinate criteria follow from
the already accepted unique-representation clause, so the missing work is bookkeeping over
$\Nzero^2$ rather than new arithmetic.

\subsection{Why this does not yet descend below \texorpdfstring{$\beta$}{beta}}

If $n$ and $k$ have a common divisor, Theorem~\ref{pp:common-power} descends to
$x/\gcd(n,k)$.  If a matched decomposition exists, Theorem~\ref{pp:affine-decomp} descends to
$(x-r)/d$.  But exact coordinates guarantee that the descended point is again $p+q\beta$ with
$p,q\in\Nzero$, and it lies below $\beta$ only when $q=0$, in which case it is integral and
carries no information.  The least nonintegral point $\beta$ is already primitive for every
descent of this shape, which is precisely why the kernel-verified corollaries at $\beta$ are
exclusions rather than contradictions.

A successful descent argument therefore needs a transformation not captured by common perfect
powers and common distinguished-base depths.  Both operations act inside the monoid
$\Nzero^2$; anything that could close the problem must leave it.

\subsection{Total and primitive counts}

Corollary~\ref{pp:gcd-degree} converts ordinary coprimality statistics of lattice points into
exact statistics of hypothetical output pairs.  For $T>0$ put
\[
 N_{\beta}(T)=\#\bigl\{(n,k)\in\Nzero^2\setminus\{(0,0)\}:n+k\beta<T\bigr\}
\]
and
\[
 P_{\beta}(T)=\#\bigl\{(n,k)\in\Nzero^2:n+k\beta<T,\ \gcd(n,k)=1\bigr\}.
\]
The first counts nonzero solutions below $T$; by Corollary~\ref{pp:gcd-degree} the second
counts exactly those output pairs with no nontrivial common perfect-power degree.

Counting each column separately, for every real $U>0$,
\begin{equation}
 N_{\beta}(U)=\sum_{0\le k<U/\beta}\bigl\lceil U-k\beta\bigr\rceil-1
 =\frac{U^2}{2\beta}+O_{\beta}(U+1),
\label{pp:eq-lattice-count}
\end{equation}
uniformly in $U$.  For integer $T$ this is the kernel-verified cumulative count $C(T)$ of
Section~\ref{sec:conditional} minus the point $x=0$, so \eqref{pp:eq-lattice-count} agrees in
leading order with \lean{QuadraticAsymptotic}; the uniform real-variable form used below is
elementary.

\begin{theorem}[Primitive output pairs]
\label{pp:primitive-asymptotic}
As $T\to\infty$,
\[
 P_{\beta}(T)
 =\frac{T^2}{2\beta\zeta(2)}+O_{\beta}(T\log T)
 =\frac{3T^2}{\pi^2\beta}+O_{\beta}(T\log T).
\]
Consequently
\[
 \frac{P_{\beta}(T)}{N_{\beta}(T)}\longrightarrow\frac1{\zeta(2)}=\frac6{\pi^2}.
\]
\end{theorem}

\begin{proof}
Every nonzero coordinate pair has the unique factorization
\[
 (n,k)=d\,(n_0,k_0),\qquad d=\gcd(n,k),\quad\gcd(n_0,k_0)=1,
\]
and $n+k\beta<T$ holds if and only if $n_0+k_0\beta<T/d$.  Hence
\[
 N_{\beta}(T)=\sum_{d\ge1}P_{\beta}(T/d),
\]
a finite sum for each $T$, and M\"obius inversion gives
\[
 P_{\beta}(T)=\sum_{d\le T}\mu(d)\,N_{\beta}(T/d).
\]
Inserting the uniform estimate \eqref{pp:eq-lattice-count},
\[
\begin{aligned}
 P_{\beta}(T)
 &=\frac{T^2}{2\beta}\sum_{d\le T}\frac{\mu(d)}{d^2}
   +O_{\beta}\!\left(\sum_{d\le T}\Bigl(\frac Td+1\Bigr)\right)\\
 &=\frac{T^2}{2\beta\zeta(2)}+O_{\beta}(T\log T),
\end{aligned}
\]
where the tail $\sum_{d>T}\mu(d)/d^2=O(1/T)$ contributes $O_{\beta}(T)$.  Dividing by
\eqref{pp:eq-lattice-count} gives the stated limit.
\end{proof}

The constant $6/\pi^2$ is universal: it does not depend on the hypothetical generator at all.
The geometry of the monoid contributes the factor $1/(2\beta)$, and the primitive-lattice
density contributes $1/\zeta(2)$.  In particular no choice of counterexample can make the
primitive proportion small.

\subsection{Distribution of the maximal common power degree}

For fixed $g\ge1$ let
\[
 P_{\beta,g}(T):=\#\bigl\{x=n+k\beta:0<x<T,\ \gcd(n,k)=g\bigr\}
\]
be the number of solutions below $T$ whose output pair has common power degree exactly $g$.

\begin{corollary}[Exact-degree distribution]
\label{pp:degree-distribution}
For fixed $g\ge1$,
\[
 P_{\beta,g}(T)=P_{\beta}(T/g)
 =\frac{3T^2}{\pi^2\beta g^2}
  +O_{\beta}\!\left(\frac Tg\log\frac{2T}{g}\right).
\]
The limiting probability that an output pair below $T$ has maximal common power degree exactly
$g$ is
\[
 \frac6{\pi^2g^2}.
\]
\end{corollary}

\begin{proof}
The map $(n_0,k_0)\mapsto g\,(n_0,k_0)$ is a bijection from the primitive pairs below $T/g$
onto the pairs of gcd exactly $g$ below $T$, so $P_{\beta,g}(T)=P_{\beta}(T/g)$; now apply
Theorem~\ref{pp:primitive-asymptotic} and divide by \eqref{pp:eq-lattice-count}.
\end{proof}

The probabilities sum to one, since $\sum_{g\ge1}6/(\pi^2g^2)=1$.  Likewise, by
Theorem~\ref{pp:common-power} both outputs are perfect $r$-th powers exactly when $r\mid n$ and
$r\mid k$, so
\begin{equation}
 \#\bigl\{0<x<T:\text{$2^x$ and $3^x$ are both perfect $r$-th powers}\bigr\}
 =N_{\beta}(T/r),
\label{pp:eq-rth-power-frequency}
\end{equation}
and the limiting proportion is $1/r^2$.

\subsection{Exact generating series}

The coordinate monoid has an equally explicit Laplace series.  For $s>0$,
\begin{equation}
\begin{aligned}
 Z_{\beta}(s)
 &:=\sum_{x\in\Sol\setminus\{0\}}e^{-sx}
  =\sum_{n,k\ge0}e^{-s(n+k\beta)}-1\\
 &=\frac1{(1-e^{-s})(1-e^{-\beta s})}-1 ,
\end{aligned}
\label{pp:eq-laplace-series}
\end{equation}
the two factors being the two free generators $1$ and $\beta$.  M\"obius inversion in the degree
gives the primitive-output series
\[
 Z_{\beta}^{\mathrm{prim}}(s)=\sum_{d\ge1}\mu(d)\,Z_{\beta}(ds).
\]
At the candidate level the same factorization appears multiplicatively: since
$\Cand=\{2^nM^k:n,k\in\Nzero\}$ with unique coordinates,
\begin{equation}
 \sum_{m\in\Cand}m^{-s}
 =\frac1{(1-2^{-s})(1-M^{-s})},
 \qquad \Re s>0.
\label{pp:eq-candidate-dirichlet}
\end{equation}
These identities make the two free generators visible in analytic form and give clean starting
points for weighted Tauberian refinements of Theorem~\ref{pp:primitive-asymptotic}.

\begin{statusbox}{\statusnew{} --- Theorem~\ref{pp:primitive-asymptotic},
Corollary~\ref{pp:degree-distribution}, and the series
\eqref{pp:eq-laplace-series}--\eqref{pp:eq-candidate-dirichlet} are paper results of this
report.  Their only nonelementary input is the kernel-verified exact monoid classification of
Section~\ref{sec:conditional}, whose leading-order count is itself \statuslean{} through
\lean{ExactCounting} and \lean{QuadraticAsymptotic}.  The finite scans of
Section~\ref{sec:finitecheck} play no role here; the exponent bound in force remains the
kernel-verified $x\ge12$.}
\end{statusbox}

\subsection{Interpretation}

Roughly $60.79\%$ of hypothetical output pairs would be simultaneously power-primitive, and
the remaining ones split across degrees $g\ge2$ with the exact frequencies
$6/(\pi^2g^2)$.  This is not pressure toward a contradiction.  It is precisely the primitive
proportion expected of a free rank-two monoid, and it would be the same for any generator: the
statistic is diagnostic, not restrictive.

Its diagnostic value is a warning about method.  The kernel-verified primitivity corollaries at
$\beta$, and the classification of this section at a general point, both say that
perfect-power and matched-depth structure is the exception among hypothetical solutions.  Any
descent argument built on extracting a common root or a common base depth therefore acts on a
minority of the conditional population --- asymptotically the proportion
$1-6/\pi^2\approx0.3921$ of pairs that admit any nontrivial common degree at all --- and
leaves the primitive majority untouched.  A method that closes the problem must engage the
power-primitive pairs directly, which is exactly where the transcendence content of
Section~\ref{sec:fourexp} sits.

\section{Candidate divisibility and ambient root saturation}\label{sec:lattice}

The exact classification of Section~\ref{sec:conditional} describes the conditional candidate
set by its \emph{generators}: under failure, $\Cand=\{2^nM^k:n,k\in\Nzero\}$ with unique
coordinates.  Generation alone says nothing about how candidates sit inside $\N$ under the two
operations that descent arguments and searches actually use --- divisibility, and extraction of
roots.  This section answers both exactly.  The candidate monoid has a simple valuation
geometry: divisibility is coordinatewise order on a two-dimensional lattice cone, $\Cand$ is
closed under $\gcd$ and $\operatorname{lcm}$, and the set of integers some power of which is a
candidate is again an explicit cone --- one that contains $\Cand$ with an exact finite-index
defect measured by the rational-power index $d$ of Section~\ref{sec:conditional}, which is also
the exact radical degree of Section~\ref{sec:radical}.

Every statement below is \statusnew{}.  Its inputs --- the canonical primitive odd core, the
rational-power index, and the exact primitive-generator theorem --- are \statuslean{}; what is
added here is elementary lattice combinatorics on top of them.  None of it is a step toward a
contradiction: as the closing discussion makes explicit, the closure properties proved here are
compatible with failure by construction, and the kernel-verified exponent bound remains
$x\ge12$.

\subsection{A two-dimensional valuation cone}

\begin{lemma}[Valuation coordinates and divisibility]
\label{lt:lem-cone}
Assume Conjecture~\ref{conj:AE} is false and let $w,d,a,\beta,M$ be the canonical data of
Section~\ref{sec:conditional}, so that $w>1$ is odd and power-primitive, $d>0$, and
\[
 M=2^{\beta}=2^aw^d .
\]
Put $U:=w^d$, the odd part of $M$, and for $n,k\in\Nzero$ write
\[
 m(n,k):=2^nM^k=2^{\,n+ak}U^k,
 \qquad s:=\vpp{2}{m(n,k)}=n+ak .
\]
Then:
\begin{enumerate}[label=\textup{(\roman*)}]
\item $U>1$;
\item the map $m(n,k)\mapsto(s,k)$ is a bijection from $\Cand$ onto the lattice cone
\[
 \Lambda_a:=\{(s,k)\in\Nzero^2:s\ge ak\};
\]
\item writing $m_i=m(n_i,k_i)$ and $s_i=n_i+ak_i$,
\begin{equation}
 m_1\mid m_2
 \quad\Longleftrightarrow\quad
 s_1\le s_2\ \text{ and }\ k_1\le k_2 .
 \label{lt:eq-divisibility-coordinates}
\end{equation}
\end{enumerate}
\end{lemma}

\begin{proof}
(i) If $U=1$ then $M=2^a$ and $\beta=a\in\Z$, contradicting the irrationality of the least
generator.

(ii) By the exact primitive-generator theorem of Section~\ref{sec:conditional}, every candidate
is $m(n,k)$ for a unique pair $(n,k)\in\Nzero^2$.  The reparameterization $(n,k)\mapsto(n+ak,k)$
is affine with inverse $(s,k)\mapsto(s-ak,k)$, and $n\ge0$ holds exactly when $s\ge ak$.

(iii) Since $U$ is odd, $\vpp{2}{m(n,k)}=s$, and for an odd prime $p$ we have
$\vpp{p}{m(n,k)}=k\,\vpp{p}{U}$, which vanishes unless $p\mid U$.  Divisibility in $\N$ is
comparison of all prime valuations.  If $m_1\mid m_2$ then comparison at $2$ gives $s_1\le s_2$,
and comparison at any prime $p\mid U$ --- one exists by (i) --- gives $k_1\vpp{p}{U}\le
k_2\vpp{p}{U}$ with $\vpp{p}{U}>0$, hence $k_1\le k_2$.  Conversely those two inequalities
force every prime valuation of $m_1$ to be at most the corresponding valuation of $m_2$.
\end{proof}

Thus candidate divisibility is nothing but the coordinatewise order on $\Lambda_a$.  The
inequality $s\ge ak$ is the only constraint, and it records the base-adic depth $a$ of the least
generator: a candidate cannot carry the odd part $U^k$ without carrying at least $ak$ factors of
two along with it.

\subsection{The candidate gcd/lcm lattice}

\begin{theorem}[Candidate gcd/lcm lattice]
\label{lt:thm-gcd-lcm}
Unconditionally, $\Cand$ is closed under $\gcd$ and $\operatorname{lcm}$.  If
Conjecture~\ref{conj:AE} fails, then in the notation of Lemma~\ref{lt:lem-cone},
\[
\begin{aligned}
 \gcd(m_1,m_2)
 &=m\!\left(\min(s_1,s_2)-a\min(k_1,k_2),\ \min(k_1,k_2)\right),\\
 \operatorname{lcm}(m_1,m_2)
 &=m\!\left(\max(s_1,s_2)-a\max(k_1,k_2),\ \max(k_1,k_2)\right);
\end{aligned}
\]
equivalently, in cone coordinates $\gcd$ and $\operatorname{lcm}$ are coordinatewise minimum
and maximum on $\Lambda_a$.  Candidate divisibility is therefore a distributive lattice, and
$\Lambda_a$ is a sublattice of $\Nzero^2$.
\end{theorem}

\begin{proof}
If the conjecture holds then $\Cand=\{2^n:n\in\Nzero\}$, whose gcds and lcms are again powers
of two; this gives the unconditional half.

Assume failure.  By Lemma~\ref{lt:lem-cone} the candidate $m(n,k)$ has valuation vector
$s$ at the prime $2$ and $k\,\vpp{p}{U}$ at each odd $p$.  Gcd and lcm take coordinatewise
minima and maxima of valuation vectors, so
\[
 \gcd(m_1,m_2)=2^{\min(s_1,s_2)}U^{\min(k_1,k_2)},
 \qquad
 \operatorname{lcm}(m_1,m_2)=2^{\max(s_1,s_2)}U^{\max(k_1,k_2)} .
\]
It remains to check that both exponent pairs lie in $\Lambda_a$.  Suppose $s_1\le s_2$; then
$\min(s_1,s_2)=s_1\ge ak_1\ge a\min(k_1,k_2)$ because $a\ge0$.  For the maximum, suppose
$k_1\le k_2$; then $\max(s_1,s_2)\ge s_2\ge ak_2=a\max(k_1,k_2)$.  Hence $\Lambda_a$ is closed
under coordinatewise $\min$ and $\max$, that is, it is a sublattice of $\Nzero^2$, and the
translation back to $(n,k)$ coordinates is the stated substitution $n=s-ak$.  A sublattice of a
product of chains is distributive.
\end{proof}

\begin{statusbox}{\statuslean{} --- \lean{twoBaseNaturalCandidate_gcd} and
\lean{twoBaseNaturalCandidate_lcm} in \lean{CandidateLattice}, resting on the coordinatewise
arithmetic lemmas \lean{gcd_two_pow_mul_odd_pow} and \lean{lcm_two_pow_mul_odd_pow} of the
same module and on the \statuslean{} exact primitive-generator theorem of
Section~\ref{sec:conditional}.  The formal statements are unconditional: the Lean proof splits
on the conjecture and disposes of the affirmative branch through the powers of two.  No
transcendence input and no analytic estimate enters.}
\end{statusbox}

\begin{corollary}[Counting candidate divisors]
\label{lt:cor-divisors}
Under failure, the number of candidates dividing $m(n,k)$, including $1$ and $m(n,k)$ itself, is
\[
 (k+1)(n+1)+\frac{ak(k+1)}2 .
\]
\end{corollary}

\begin{proof}
By \eqref{lt:eq-divisibility-coordinates} a candidate divisor corresponds to a point
$(s',j)\in\Lambda_a$ with $0\le j\le k$ and $aj\le s'\le n+ak$.  For fixed $j$ there are
$n+ak-aj+1$ admissible values of $s'$.  Summing the arithmetic progression over
$j=0,1,\ldots,k$ gives $(k+1)(n+1)+a\sum_{i=0}^{k}i$, which is the stated value.
\end{proof}

\begin{remark}[Candidates are not divisor-closed]
\label{lt:rem-not-divisor-closed}
The full divisor count of $m(n,k)=2^sU^k$ is
$(s+1)\prod_{p\mid U}\left(k\,\vpp{p}{U}+1\right)$, which for $k\ge1$ and $U$ with more than one
prime factor grows much faster than the count of Corollary~\ref{lt:cor-divisors}.  So $\Cand$ is
emphatically not closed under taking divisors; it is closed only under the two lattice
operations.  The gap between the two counts is a concrete measure of how thin the candidate
sublattice is inside the divisor lattice of a single candidate.
\end{remark}

\subsection{Fixed prime support}

The lattice picture has an immediate consequence for the shape of a hypothetical counterexample
branch: the supports do not vary at all.

\begin{corollary}[Prime-support rigidity]
\label{lt:cor-prime-support}
Assume failure and use the simultaneous primitive output normalization of
Section~\ref{sec:conditional},
\[
 M=2^{\beta}=2^aw^d,
 \qquad
 A=3^{\beta}=3^bv^e .
\]
Then for $x=n+k\beta$,
\[
 2^x=2^{\,n+ak}w^{dk},
 \qquad
 3^x=3^{\,n+bk}v^{ek},
\]
so every nonintegral solution --- that is, every solution with $k\ge1$ --- has exactly the same
odd prime support $\operatorname{supp}(w)$ in its base-$2$ output, and exactly the same prime
support away from $3$, namely $\operatorname{supp}(v)$, in its base-$3$ output.  The only
variation across the entire nonintegral branch is whether the distinguished primes $2$ and $3$
occur at all.
\end{corollary}

\begin{proof}
Substituting $M$ and $A$ into $(2^x,3^x)=(2^nM^k,3^nA^k)$ gives the displayed factorizations.
For $k\ge1$ the odd part of $2^x$ is $w^{dk}$, and
$\operatorname{supp}(w^{dk})=\operatorname{supp}(w)$ because $dk\ge1$; similarly for $v^{ek}$.
The exponents $n+ak$ and $n+bk$ may vanish, which is exactly the stated freedom at $2$ and $3$.
\end{proof}

A counterexample therefore locks the whole nonintegral branch onto a single fixed pair of finite
prime sets.  Computational reconnaissance should exploit this directly, enumerating primitive
cores and supports rather than raw integers scanned independently; the finite verifications
reported in Section~\ref{sec:finitecheck} are \statuscomp{} and this observation does not change
their status, only the shape of the search that would extend them.

\subsection{Ambient root saturation and the index \texorpdfstring{$d$}{d}}

Reparameterizing Lemma~\ref{lt:lem-cone} by the exponent of the primitive core $w$ rather than
of $U=w^d$ gives the description that governs roots.  With $s=\vpp{2}{m}$ and $t$ the exponent
of $w$ in the odd part of $m$,
\begin{equation}
 \Cand=\{2^sw^t:s,t\in\Nzero,\ d\mid t,\ ds\ge at\}.
 \label{lt:eq-candidate-cone}
\end{equation}
Indeed $m=2^nM^k$ has $s=n+ak$ and $t=dk$, so $d\mid t$, and $n\ge0$ is equivalent to
$s\ge at/d$, that is, to $ds\ge at$.

\begin{definition}[Ambient root saturation]
\label{lt:def-saturation}
For a multiplicative submonoid $H\subseteq\N_{>0}$ put
\[
 \operatorname{Sat}(H):=\{u\in\N_{>0}:u^r\in H\text{ for some }r\ge1\},
\]
the set of positive integers admitting some positive power inside $H$.
\end{definition}

\begin{theorem}[Exact ambient root saturation]
\label{lt:thm-root-saturation}
Assume Conjecture~\ref{conj:AE} is false.  Then
\[
 \operatorname{Sat}(\Cand)=\{2^sw^t:s,t\in\Nzero,\ ds\ge at\},
\]
and for $u=2^sw^t$ in this set the least positive $r$ with $u^r\in\Cand$ is
\[
 r_{\min}(u)=\frac d{\gcd(d,t)} .
\]
\end{theorem}

\begin{proof}
First let $u=2^sw^t$ with $s,t\in\Nzero$ and $ds\ge at$, and put $r=d/\gcd(d,t)$.  Then
$rt=d\cdot t/\gcd(d,t)$ is divisible by $d$, and multiplying $ds\ge at$ by $r>0$ gives
$d(rs)\ge a(rt)$.  By \eqref{lt:eq-candidate-cone}, $u^r=2^{rs}w^{rt}\in\Cand$, so
$u\in\operatorname{Sat}(\Cand)$.  Conversely, if $u^{r'}\in\Cand$ for some $r'\ge1$ then
\eqref{lt:eq-candidate-cone} forces $d\mid r't$, since $w$ is power-primitive and hence the
exponent of $w$ in the odd part of $u^{r'}$ is exactly $r't$.  Now $d\mid r't$ is equivalent to
$\bigl(d/\gcd(d,t)\bigr)\mid r'$, so $r'\ge r$ and $r$ is least.

Now suppose merely that $u\in\N_{>0}$ satisfies $u^r\in\Cand$ for some $r\ge1$, and write
$u^r=2^Sw^{dk}$ with $dS\ge a\,dk$.  A prime divides $u$ if and only if it divides $u^r$, so no
prime outside $\{2\}\cup\operatorname{supp}(w)$ divides $u$.  Put $c_p:=\vpp{p}{w}$ for
$p\mid w$; comparing valuations in $u^r=2^Sw^{dk}$ gives
\[
 r\,\vpp{p}{u}=dk\,c_p
 \qquad\text{for every }p\mid w .
\]
Because $w$ is power-primitive, $\gcd\{c_p:p\mid w\}=1$, so there are integers $r_p$, almost all
zero, with $\sum_pr_pc_p=1$; then
\[
 \frac{dk}r=\sum_pr_p\cdot\frac{dk\,c_p}r=\sum_pr_p\,\vpp{p}{u}\in\Z .
\]
Call this integer $t\ge0$.  Then $\vpp{p}{u}=tc_p$ for every $p\mid w$, so the odd part of $u$
is exactly $w^t$, while $S=rs$ for $s=\vpp{2}{u}$.  Substituting $S=rs$ and $dk=rt$ into
$dS\ge a\,dk$ and dividing by $r>0$ yields $ds\ge at$.  Hence $u=2^sw^t$ lies in the stated
cone.
\end{proof}

\begin{statusbox}{\statusnew{} --- the only nonformal ingredient is the B\'ezout step, which is
the same primitivity mechanism already kernel-verified in \lean{PrimitiveOddCore} through
\lean{oddPrimeValuationGCD}.  Everything else is unique factorization and integer arithmetic.}
\end{statusbox}

\begin{corollary}[Ambient saturation index]
\label{lt:cor-ambient-index}
Under failure, $\Cand=\operatorname{Sat}(\Cand)$ if and only if $d=1$, and in general
\[
 d=\min\{r\ge1:\exists\,s\in\Nzero,\ (2^sw)^r\in\Cand\}.
\]
\end{corollary}

\begin{proof}
Take $t=1$ and any $s\ge a$, so that $ds\ge a=at$; by Theorem~\ref{lt:thm-root-saturation} the
integer $2^sw$ lies in $\operatorname{Sat}(\Cand)$ with $r_{\min}=d/\gcd(d,1)=d$.  This proves
the displayed formula, and it shows that $2^sw\in\Cand$ forces $d=1$.  Conversely, if $d=1$ then
the divisibility condition $d\mid t$ in \eqref{lt:eq-candidate-cone} is vacuous and the two sets
coincide.
\end{proof}

\begin{corollary}[Which integer roots of a candidate are again candidates]
\label{lt:cor-roots}
Assume failure and let $m=2^nM^k\in\Cand$ and $r\ge1$.
\begin{enumerate}[label=\textup{(\alph*)}]
\item There is an integer $u$ with $u^r=m$ if and only if $r\mid n+ak$ and $r\mid dk$; such a
$u$ always lies in $\operatorname{Sat}(\Cand)$.
\item That integer root is itself a candidate if and only if $r\mid n$ and $r\mid k$, in which
case $u=2^{n/r}M^{k/r}$.
\end{enumerate}
\end{corollary}

\begin{proof}
(a) Write $m=2^sw^t$ with $s=n+ak$ and $t=dk$.  If $u^r=m$ then, exactly as in the proof of
Theorem~\ref{lt:thm-root-saturation}, $u$ is supported in $\{2\}\cup\operatorname{supp}(w)$,
$r\mid s$, and $t/r\in\Z$ by the B\'ezout step; conversely $u=2^{s/r}w^{t/r}$ is an integer
$r$-th root when $r\mid s$ and $r\mid t$.  Membership in $\operatorname{Sat}(\Cand)$ is
immediate from Definition~\ref{lt:def-saturation}.

(b) If $u\in\Cand$ then $u=2^{n'}M^{k'}$ for unique $n',k'\in\Nzero$, and $u^r=2^{rn'}M^{rk'}$
together with uniqueness of coordinates gives $n=rn'$ and $k=rk'$.  The converse is the direct
computation $\bigl(2^{n/r}M^{k/r}\bigr)^r=m$.
\end{proof}

\begin{observation}[Normal intrinsically, saturated only up to index $d$]
\label{lt:obs-index}
Let $L:=\{2^sw^t:s,t\in\Z\}$ be the ambient multiplicative group generated by $2$ and $w$ inside
$\Q_{>0}^{\times}$, and let $L_{\Cand}$ be the group generated by $\Cand$.  Then
\[
 L_{\Cand}=\{2^sw^t:s\in\Z,\ d\mid t\},
 \qquad
 [L:L_{\Cand}]=d,
 \qquad
 \Cand=\operatorname{Sat}(\Cand)\cap L_{\Cand}.
\]
In particular $\operatorname{Sat}(\Cand)$ is itself a submonoid of $\N_{>0}$, namely the full
lattice points of the rational cone $ds\ge at$, $s,t\ge0$.
\end{observation}

\begin{proof}
$L_{\Cand}$ is generated by $2$ and $M=2^aw^d$, so it consists of the elements
$2^{i+aj}w^{dj}$ with $i,j\in\Z$, which is the stated set; the quotient $L/L_{\Cand}$ is
isomorphic to $\Z/d\Z$ through the exponent of $w$.  Intersecting the cone description of
Theorem~\ref{lt:thm-root-saturation} with the congruence $d\mid t$ returns
\eqref{lt:eq-candidate-cone}.  Closure of the cone under multiplication is additivity of the
inequality $ds\ge at$ in $(s,t)$.
\end{proof}

A terminology caution is worth stating explicitly, because the word ``normal'' is used
differently in affine-semigroup theory.  There, normality of a semigroup is measured relative to
the group that the semigroup itself generates.  By Observation~\ref{lt:obs-index} that group is
$L_{\Cand}$, in which the congruence $t\equiv0\pmod d$ already holds identically, and $\Cand$ is
the full set of cone points of $L_{\Cand}$ for every value of $d$: intrinsically, $\Cand$ is
normal.  Theorem~\ref{lt:thm-root-saturation} measures something different, namely saturation
inside the larger ambient lattice $L$ containing every $2^sw^t$.  So $d$ is an over-lattice
saturation defect of finite index, not a failure of intrinsic normality.  This is precisely
parallel to its two other exact meanings established elsewhere in this report: $d$ is the least
rational-power index of Section~\ref{sec:conditional} and the exact radical degree
$[\Q(w^{\thetaq}):\Q]$ of Section~\ref{sec:radical}.  One integer invariant, three
characterizations --- algebraic degree, least rational power, and lattice saturation defect.

\subsection{Why lattice saturation is a good Lean target}

Two points must be kept apart.  The gcd/lcm closure of Theorem~\ref{lt:thm-gcd-lcm} and the
saturation formula of Theorem~\ref{lt:thm-root-saturation} are compatible with failure by
construction: they describe the counterexample world rather than obstruct it, and no
contradiction can be extracted from them alone.  They are structural information and search
normalization.  With that said, they are unusually attractive as the next formalization
increment.

The mathematical inputs are already kernel-verified.  The canonical primitive odd core, the
rational-power index, the three-power denominator normalization, and the exact primitive
generator with unique coordinates are all \statuslean{} and live in
\lean{ExponentialIdentities/TwoBaseIntegerExponent/}.  Everything added in this section is
finite combinatorics over those results: unique factorization in $\N$, one B\'ezout step that
already exists in \lean{PrimitiveOddCore} as \lean{oddPrimeValuationGCD}, and elementary
$\min$/$\max$ arithmetic on $\Nzero^2$.  There is no analysis, no transcendence input, and no
numerical certificate anywhere in the chain --- a rare profile in this corpus, and the reason the
work should be cheap relative to its yield.  A single module, sitting beside
\lean{PrimitiveGenerator} and \lean{RationalPowerIndex}, would carry the coordinate bijection of
Lemma~\ref{lt:lem-cone}, the divisibility criterion
\eqref{lt:eq-divisibility-coordinates}, the lattice operations, the divisor count, the cone form
\eqref{lt:eq-candidate-cone}, the saturation set, the exact value of $r_{\min}$, and the index
characterization of $d$.  The natural name is
\begin{center}
\lean{CandidateLattice}
\end{center}
and the blueprint conventions of Section~\ref{sec:blueprint} apply to it unchanged.

The contribution to the search is threefold.  First, a kernel-checked normal form: candidate
enumeration can be phrased over lattice coordinates rather than over raw integers, so rejection
tests become integer inequalities of the form $ds\ge at$ and $d\mid t$ instead of factorizations
recomputed for every scanned value.  This is exactly the discipline recommended in the
computational priority of Section~\ref{sec:program}, and it applies equally to the certificate
architecture that any extension of the \statuscomp{} finite verifications of
Section~\ref{sec:finitecheck} would need; the kernel-verified region itself would still stop at
$x\ge12$ until new certificates are replayed.  Second, a single invariant with three
independently proved meanings: once Corollary~\ref{lt:cor-ambient-index} is formal alongside the
radical-degree proposition of Section~\ref{sec:radical}, any future argument that bounds $d$ ---
from Kummer theory, from ramification, or from a determinant estimate --- bounds all three at
once, and the Lean statements make that transfer mechanical rather than a matter of prose.
Third, and most concretely, root extraction is the step where descent arguments live.  Every
descent so far --- no matched base depth, no common perfect-power degree, exact affine
primitivity in Section~\ref{sec:conditional} --- proceeds by extracting a root of the output pair
and contradicting leastness of $\beta$.  Corollary~\ref{lt:cor-roots} says exactly which integer
roots of a candidate survive as candidates and which merely land in the saturation, so a
formalized version supplies the reusable side condition that such arguments currently re-prove
by hand.  That is a modest but genuine improvement in the leverage of the conditional structure,
and it costs no new mathematics.

\section{A sufficient condition: independence of prime-log products}
\label{sec:logproduct}

The reductions of Section~\ref{sec:radical} isolate the obstruction as a statement about one
radical at a time.  This section isolates it differently, as a single statement about
$\Q$-linear independence of a natural family of real numbers.  The resulting criterion is
weaker in appearance than the four exponentials conjecture and is of a different kind: it does
not mention exponentials at all.

\subsection{Failure as a vanishing linear form in log-products}

Suppose the conjecture fails, and let $x$ be a nonintegral solution with outputs
\[
  M=2^x\in\N,
  \qquad
  A=3^x\in\N .
\]
Taking logarithms of $M=2^x$ and $A=3^x$ and eliminating $x$ gives the defining relation
already recorded in \eqref{pr:log-det}:
\begin{equation}
  \log M\,\log3-\log A\,\log2=0 .
  \label{eq:lp-relation}
\end{equation}
Now expand both integers over their prime factorizations,
\[
  M=\prod_p p^{\alpha_p},
  \qquad
  A=\prod_p p^{\gamma_p},
  \qquad
  \alpha_p,\gamma_p\in\Nzero,
\]
so that $\log M=\sum_p\alpha_p\log p$ and $\log A=\sum_p\gamma_p\log p$.  Substituting into
\eqref{eq:lp-relation},
\begin{equation}
  \sum_p\alpha_p\bigl(\log p\,\log3\bigr)
  -\sum_p\gamma_p\bigl(\log p\,\log2\bigr)=0 .
  \label{eq:lp-expanded}
\end{equation}
Every term of \eqref{eq:lp-expanded} is an integer multiple of a product of exactly two
logarithms of primes.  Collecting coefficients in the family
\[
  \bigl\{\log p\,\log q\ :\ p\le q\ \text{prime}\bigr\},
\]
the unordered pair $\{p,3\}$ with $p\notin\{2,3\}$ receives the coefficient $\alpha_p$; the
pair $\{p,2\}$ with $p\notin\{2,3\}$ receives $-\gamma_p$; the diagonal pair $\{3,3\}$
receives $\alpha_3$; the diagonal pair $\{2,2\}$ receives $-\gamma_2$; and the mixed pair
$\{2,3\}$ receives $\alpha_2-\gamma_3$.

Thus \emph{failure of the conjecture is exactly the existence of a nontrivial vanishing
integer-coefficient linear form in products of two prime logarithms}, with the additional
positivity constraint that the coefficients arise from prime factorizations.

\subsection{The independence criterion}

\begin{definition}[Prime-log-product independence]
\label{lp:def-independence}
Say that \emph{prime-log products are independent} if the real numbers
$\log p\,\log q$, indexed by pairs of primes $p\le q$, are linearly independent over $\Q$.
\end{definition}

Equivalently, in the integer-coefficient form convenient for formalization: for every finite
set $S$ of primes and every $c:S\times S\to\Z$,
\[
  \sum_{p\in S}\sum_{q\in S}c_{pq}\,\log p\,\log q=0
  \qquad\Longrightarrow\qquad
  c_{pq}+c_{qp}=0\ \text{ for all }p,q\in S .
\]
(The diagonal case $p=q$ reads $2c_{pp}=0$, hence $c_{pp}=0$.)  Clearing denominators shows
the two formulations agree.

\begin{theorem}[Prime-log-product independence implies the conjecture]
\label{lp:thm-main}
If prime-log products are independent, then Conjecture~\ref{conj:AE} holds.
\end{theorem}

\begin{proof}
Let $x$ satisfy $2^x,3^x\in\Z$.  Both powers are positive, so $M:=2^x$ and $A:=3^x$ are
positive integers, and $x\ge0$.  Relation \eqref{eq:lp-relation} holds, hence so does
\eqref{eq:lp-expanded}.

Apply the independence hypothesis to the finite set
$S=\operatorname{primes}(M)\cup\operatorname{primes}(A)\cup\{2,3\}$ and to the coefficient
family
\[
 c_{pq}
 =\begin{cases}
   \alpha_p, & q=3,\\
   -\gamma_p, & q=2,\\
   0, & \text{otherwise,}
  \end{cases}
\]
which reproduces the left-hand side of \eqref{eq:lp-expanded}.  The conclusion
$c_{pq}+c_{qp}=0$ yields, coefficient by coefficient:
\begin{itemize}
\item taking $q=3$ and $p\in S\setminus\{2,3\}$: $\alpha_p+0=0$, so $\alpha_p=0$;
\item taking $p=q=3$: $2\alpha_3=0$, so $\alpha_3=0$;
\item taking $q=2$ and $p\in S\setminus\{2,3\}$: $-\gamma_p+0=0$, so $\gamma_p=0$;
\item taking $p=q=2$: $-2\gamma_2=0$, so $\gamma_2=0$;
\item taking $p=2$, $q=3$: $\alpha_2-\gamma_3=0$.
\end{itemize}
Primes outside $S$ divide neither $M$ nor $A$, so their exponents vanish as well.  Hence
$\alpha_p=0$ for every $p\neq2$, that is
\[
  M=2^{\alpha_2},
\]
and symmetrically $A=3^{\gamma_3}=3^{\alpha_2}$.  Writing $n=\alpha_2\in\Nzero$ we get
$2^x=M=2^n$, and strict monotonicity of $t\mapsto2^t$ gives $x=n\in\Z$.
\end{proof}

\begin{statusbox}{\statuslean{} --- \lean{LogProductIndependence} implements this argument,
with the independence hypothesis recorded as an explicit \lean{Prop} and never asserted; the
accepted conditional theorem is
\lean{alaogluErdosConjecture_of_primeLogProductIndependence}.}
\end{statusbox}

\subsection{What the criterion is, and is not}

Three remarks place Theorem~\ref{lp:thm-main} correctly.

\emph{It is a consequence of Schanuel's conjecture.}  Schanuel's conjecture implies that
logarithms of multiplicatively independent positive rationals are algebraically independent
over $\Q$; in particular the numbers $\log p$, over all primes $p$, would be algebraically
independent, and then the products $\log p\log q$ are $\Q$-linearly independent because
distinct monomials of degree two in algebraically independent elements are linearly
independent.  So Definition~\ref{lp:def-independence} is strictly weaker than Schanuel and
lies well inside the standard expectations.

\emph{Its strength beyond the conjecture lies in pairs that never occur.}  This calibration
matters and is easy to miss.  Collecting the relation \eqref{eq:lp-expanded} back up,
\[
  \sum_p\alpha_p\bigl(\log p\log3\bigr)-\sum_p\gamma_p\bigl(\log p\log2\bigr)
  =\log3\,\log M-\log2\,\log A,
\]
so the only products ever exercised are those of the form $\log2\log p$ and $\log3\log p$.
Write $F$ for that sub-family.  A vanishing rational relation over $F$ factors as
$\log2\,\log u+\log3\,\log v=0$ for positive rationals $u,v$, i.e.\ says precisely that some
positive rational $\ne1$ has a rational $\thetaq$-th power.  Equivalently, $\Q$-linear
independence of $F$ is the \emph{rational-output} form of the problem:
\begin{equation}
  2^x\in\Q\ \text{ and }\ 3^x\in\Q
  \qquad\Longrightarrow\qquad
  x\in\Z .
  \label{eq:lp-rational-form}
\end{equation}
This is worth separating from Conjecture~\ref{conj:AE}.  It implies the conjecture, since
integers are rational, and it is \emph{a priori strictly stronger}: a rational counterexample
such as $2^x=3/2$ is not an integer counterexample, and nothing in the corpus reduces one to
the other.  Both follow from four exponentials, because rationals are algebraic.

The proof of Theorem~\ref{lp:thm-main} establishes \eqref{eq:lp-rational-form} verbatim: for
rational outputs the prime exponents $\alpha_p,\gamma_p$ are integers of either sign, and the
argument never used their positivity.  So the criterion is not a restatement of the
conjecture; it delivers the stronger rational-output form.  What the calibration does show is
where its surplus lies: independence of pairs $\log p\log q$ with $p,q\ge5$ is never
exercised, so one should not expect the full criterion to be easier than
\eqref{eq:lp-rational-form}.  Its unconditional dividend is the two-prime case of
Section~\ref{lp:twoprime}.

\begin{statusbox}{\statuslean{} for the integer-output form
(\lean{alaogluErdosConjecture_of_primeLogProductIndependence}); the rational-output
strengthening \eqref{eq:lp-rational-form} is \statusnew{} and awaits a Lean version stated
over \lean{Rat} valuations.}
\end{statusbox}

\emph{It is nevertheless of a different kind from four exponentials.}  The four exponentials conjecture is a
statement about $2\times2$ matrices of logarithms and their exponentials; the present criterion
mentions no exponentials and is purely a linear-independence assertion about a countable
family of real numbers.  Neither is known, and neither is known to imply the other in general;
what Theorem~\ref{lp:thm-main} shows is that \emph{either} suffices for the Alaoglu--Erd\H{o}s
conjecture.  This is worth recording because the two hypotheses invite different techniques:
four exponentials is attacked by auxiliary functions and zero estimates, while linear
independence of log-products is a question about the arithmetic of the multiplicative group of
$\Q_{>0}$ and its second symmetric power.

\emph{The smallest instances are already open.}  Independence of the family includes, as its
very first nontrivial case, the assertion that $(\log3)^2$, $(\log2)(\log3)$ and $(\log2)^2$
are $\Q$-linearly independent, equivalently that $\thetaq=\log_23$ is not a root of a rational
quadratic --- a statement about the irrationality of $\thetaq^2$ modulo $\Q(\thetaq)$ that is
not decided by present transcendence theory.  Indeed the relation exhibited in
Section~\ref{sec:radical} for a $3$-smooth counterexample would already contradict the single
independence
\[
  \alpha_2(\log2)(\log3)+\alpha_3(\log3)^2-\gamma_2(\log2)^2-\gamma_3(\log2)(\log3)\ne0
\]
for the relevant nonzero integer vector.  Restricted to the two primes $2$ and $3$, the
criterion needed is thus exactly: \emph{$1$, $\thetaq$ and $\thetaq^2$ are linearly independent
over $\Q$}, which is the $3$-smooth shadow of Conjecture~\ref{ra:conj-radical-target}.

\subsection{A finite form of the criterion}

Because a counterexample has a fixed finite prime support, the full strength of
Definition~\ref{lp:def-independence} is never needed at once.

\begin{corollary}[Finite prime-support form]
\label{lp:cor-finite}
Let $S$ be a finite set of primes containing $2$ and $3$.  If the products $\log p\log q$ with
$p\le q$ in $S$ are $\Q$-linearly independent, then no counterexample to
Conjecture~\ref{conj:AE} has both outputs supported on $S$.
\end{corollary}

\begin{proof}
The proof of Theorem~\ref{lp:thm-main} used only the pairs drawn from the prime support of
$MA$ together with $2$ and $3$.
\end{proof}

\subsection{The two-prime case is a theorem, and it is sharp}
\label{lp:twoprime}

The finite form is more than a bookkeeping device, because its smallest instance is
\emph{provable}.

\begin{proposition}[Two-prime log-product independence]
\label{lp:prop-twoprime}
Let $p\neq q$ be primes.  Then $(\log p)^2$, $\log p\,\log q$ and $(\log q)^2$ are linearly
independent over $\Q$.
\end{proposition}

\begin{proof}
A vanishing relation $a(\log p)^2+b\log p\log q+c(\log q)^2=0$ with rational $a,b,c$, divided
by $(\log p)^2$, reads $a+bt+ct^2=0$ where $t=\log q/\log p$.  If $(a,b,c)\neq0$ this makes
$t$ a root of a nonzero rational polynomial, hence algebraic.  But $t$ is transcendental:
it is irrational by unique factorization, and were it algebraic irrational then
Gelfond--Schneider applied to $p^{\,t}=q$ would make the integer $q$ transcendental.
\end{proof}

Consequently the hypothesis of Corollary~\ref{lp:cor-finite} is \emph{unconditionally true}
for $S=\{2,3\}$, and the corollary becomes an unconditional exclusion.  Stated directly, and
proved in Lean without any appeal to Definition~\ref{lp:def-independence}:

\begin{theorem}[Three-smooth outputs force an integral exponent]
\label{lp:thm-smooth}
If $2^x\in\Z$ and $3^x\in\Z$ are both $3$-smooth --- divisible by no prime outside
$\{2,3\}$ --- then $x\in\Z$.
\end{theorem}

\begin{proof}
Write $2^x=2^a3^b$ and $3^x=2^c3^d$.  Taking base-two logarithms gives the two relations
$x=a+b\thetaq$ and $x\thetaq=c+d\thetaq$.  Substituting the first into the second eliminates
$x$ and leaves
\[
  b\,\thetaq^{2}+(a-d)\,\thetaq-c=0 .
\]
If $b\neq0$ this makes $\thetaq$ a root of a nonzero rational quadratic, contradicting its
transcendence.  Hence $b=0$ and $x=a$.
\end{proof}

\begin{corollary}[A smooth candidate has a rough partner]
\label{lp:cor-rough}
If $x$ is a nonintegral solution and $2^x$ is $3$-smooth, then $3^x$ has a prime factor at
least $5$.  The two outputs of a counterexample can never both be built from $\{2,3\}$.
\end{corollary}

\begin{statusbox}{\statuslean{} --- \lean{integer_of_threeSmooth_outputs} and
\lean{exists_prime_five_le_dvd_of_threeSmooth_candidate} in \lean{SmoothOutputs}.  The proofs
use only transcendence of $\thetaq$, itself kernel-verified in this corpus.}
\end{statusbox}

Two remarks fix the position of this result precisely.

\emph{The two-sided hypothesis is essential.}  Assuming only that $2^x$ is $3$-smooth gives
$x=a+b\thetaq$ and then requires $\bigl(3^{\thetaq}\bigr)^{b}\in\Q$, which is exactly the
open localized-radical condition of Section~\ref{sec:radical}.  What the second smoothness
hypothesis supplies is a \emph{second} linear relation between $x$ and $\thetaq$; two such
relations eliminate $x$ and leave an algebraic equation in $\thetaq$ alone.  This degree drop
is the whole mechanism, and it is worth isolating as a template: \emph{each independent
multiplicative constraint on an output buys one linear relation, and two relations force
algebraicity of $\thetaq$.}  The difficulty of the general problem is precisely that one
never gets a second constraint for free.

\emph{Three primes is where the criterion becomes conjectural.}
Proposition~\ref{lp:prop-twoprime} exhausts what transcendence of a single logarithmic ratio
can give.  For $S=\{2,3,5\}$ a vanishing relation among the six products, divided by
$(\log2)^2$, reads
\[
  a+b\thetaq+c\thetaq_5+d\thetaq^{2}+e\thetaq\thetaq_5+f\thetaq_5^{2}=0,
  \qquad \thetaq_5=\log_25,
\]
that is, the point $(\thetaq,\thetaq_5)$ lies on a rational conic.  Both coordinates are
transcendental and $1,\thetaq,\thetaq_5$ are $\Q$-independent by unique factorization, but no
present method excludes a quadratic relation between them.  So
Definition~\ref{lp:def-independence} is a theorem for two primes, open from three onwards,
and the Alaoglu--Erd\H{o}s conjecture sits exactly on that boundary.
\section{The algebraic auxiliary-base spectrum}\label{sec:spectrum}

Section~\ref{sec:smooth} classified the \emph{natural} auxiliary bases whose $x$-th powers
remain algebraic, and had to tag the transcendence half of that classification as classical
background: the repository's interpolation determinant carried integer entries, so it could
prove $B^x\notin\Q$ but not $B^x\notin\Qbar$.  That boundary has moved.  The same square
determinant now accepts algebraic entries of bounded height inside a number field, and the
resulting classification is exact over \emph{all} positive real algebraic bases.  A
hypothetical nonintegral solution turns out to have one fixed, explicitly described set of
good bases --- a dense but very thin subgroup of $\R_{>0}$ --- and every positive algebraic
base outside it produces a transcendental value.

\subsection{The saturated group of positive two-three units}

The multiplicative group $\Apos$ of positive real algebraic numbers is a vector space over
$\Q$ once rational scalars act through the unique positive real root:
\[
  q\cdot\alpha=\alpha^{q},
  \qquad q\in\Q,\quad \alpha\in\Apos.
\]
For multiplicatively independent $a,b\in\Apos$ define the \emph{saturated group}
\[
  \GammaSat(a,b)=\{a^{r}b^{s}:r,s\in\Q\}.
\]
The map $\Q^{2}\to\GammaSat(a,b)$, $(r,s)\mapsto a^{r}b^{s}$, is an isomorphism of
$\Q$-vector spaces; injectivity is exactly multiplicative independence after clearing
denominators.  For the original problem write
\[
  \GammaSat=\GammaSat(2,3)=\{2^{r}3^{s}:r,s\in\Q\}.
\]

This group has a second description that is the one the Lean development actually uses.  Let
$\langle2,3\rangle_{\Z}=\{2^{i}3^{j}/(2^{k}3^{\ell}):i,j,k,\ell\in\Nzero\}$ be the group of
positive rational $2,3$-units, and let
\[
  \sqrt[\infty]{\langle2,3\rangle_{\Z}}
  :=\{\gamma>0:\exists\,n\ge1,\ \gamma^{n}\in\langle2,3\rangle_{\Z}\}
\]
be its divisible hull inside $\R_{>0}$.

\begin{observation}\label{sp:hull-equals-gamma}
$\sqrt[\infty]{\langle2,3\rangle_{\Z}}=\GammaSat$.  In particular every element of the
divisible hull is automatically algebraic.
\end{observation}

Indeed, if $\gamma^{n}=2^{i}3^{j}/(2^{k}3^{\ell})$ then $\gamma$ is the unique positive real
$n$-th root of a positive rational, hence $\gamma=2^{(i-k)/n}3^{(j-\ell)/n}\in\GammaSat$;
conversely $2^{r}3^{s}$ raised to a common denominator of $r,s$ lands in
$\langle2,3\rangle_{\Z}$.  The Lean predicate \lean{HasPositiveTwoThreeUnitPower} is the
literal transcription of membership in the hull, so it and membership in $\GammaSat$ may be
used interchangeably below.

\subsection{The exact spectrum theorem}

The kernel-verified corpus contains the determinant in fully symmetric form: three positive
algebraic bases with injective nonnegative monomials cannot all have algebraic $x$-th powers
when $x$ is irrational.  The next theorem is the coordinate-free synthesis for an arbitrary
independent pair; its proof is a direct application of the same classical six-exponentials
input, and it is stated here at paper level because the Lean development fixes the pair
$(2,3)$.

\begin{theorem}[Exact algebraic-base spectrum]\label{sp:general-spectrum}
Let $a,b\in\Apos$ be multiplicatively independent, let $x\in\R\setminus\Q$, and assume
$a^{x},b^{x}\in\Qbar$.  Then, for every $\alpha\in\Apos$,
\[
  \alpha^{x}\in\Qbar
  \quad\Longleftrightarrow\quad
  \alpha\in\GammaSat(a,b).
\]
\statusnew
\end{theorem}

\begin{proof}
If $\alpha=a^{r}b^{s}$ with $r,s\in\Q$, positivity gives $\alpha^{x}=(a^{x})^{r}(b^{x})^{s}$,
which is algebraic.

Conversely, suppose $\alpha\notin\GammaSat(a,b)$.  Then $\log a$, $\log b$, $\log\alpha$ are
linearly independent over $\Q$: a rational relation clears to an integer multiplicative
relation, and if the coefficient of $\log\alpha$ vanishes it contradicts the multiplicative
independence of $a,b$, while otherwise it places $\alpha$ in the saturated group.  The pair
$1,x$ is linearly independent over $\Q$ because $x$ is irrational.  Apply the six
exponentials theorem (Theorem~\ref{sm:six-exp}) to the three rows $\log a,\log b,\log\alpha$
and the two columns $1,x$.  Five of the six exponential values,
\[
  a,\quad b,\quad \alpha,\quad a^{x},\quad b^{x},
\]
are algebraic.  Therefore the sixth, $\alpha^{x}$, is transcendental.
\end{proof}

Conceptually this says more than any list of good auxiliary bases: two algebraic outputs have
already exhausted the maximum algebraic multiplicative rank that six exponentials permits.

Specializing to $a=2$, $b=3$ gives the statement the repository proves outright.

\begin{corollary}[Complete algebraic spectrum of a counterexample]\label{sp:AE-spectrum}
Let $x\in\Sol\setminus\Z$.  For every $\alpha\in\Apos$,
\[
  \alpha^{x}\in\Qbar
  \quad\Longleftrightarrow\quad
  \alpha\in\GammaSat.
\]
In particular the spectrum does not depend on which hypothetical nonintegral solution is
chosen.  \statuslean
\end{corollary}

A rational solution is integral, so $x$ is irrational and Theorem~\ref{sp:general-spectrum}
applies; the formal route is the divisible-hull equivalence of the next subsection together
with Observation~\ref{sp:hull-equals-gamma}.

\begin{statusbox}{\statuslean{} --- the divisible-hull equivalence lives in
\lean{AlgebraicBaseUnitPower}.}
Its two forms are
\begin{center}
\lean{real_rpow_isAlgebraic_iff_integer_or_hasPositiveTwoThreeUnitPower}
\end{center}
together with its counterexample-relative form
\begin{center}
\lean{real_rpow_isAlgebraic_iff_hasPositiveTwoThreeUnitPower}
\end{center}
in the \lean{TwoBaseNonintegerSolution} namespace, both resting on the injectivity criterion
\lean{real_two_three_gamma_monomial_injective_iff}.
The natural-base and rational-base specializations are
\lean{rpow_isAlgebraic_iff_integer_or_eq_two_pow_mul_three_pow} in \lean{AlgebraicThirdBase}
and \lean{rat_rpow_isAlgebraic_iff_isTwoThreeUnit_of_not_integer} in
\lean{AlgebraicRationalBase}; the half-plane, mixed-iterate and rank-one ratio consequences
are in \lean{AlgebraicOutputLocus}; and the simultaneous exclusion of both canonical radicals
from the entire hull is in \lean{CanonicalRadicalDivisibleHull}.
\end{statusbox}

The group $\GammaSat$ is countable, divisible, of rational rank two, and dense in $\R_{>0}$
--- density already follows from the subgroup $\{2^{r}:r\in\Q\}$.  A counterexample would
therefore possess a dense set of algebraic bases with algebraic output, while every algebraic
base outside that exact dense subgroup would give a transcendental one.

\begin{corollary}[Rank bound for any family of algebraic outputs]\label{sp:rank-bound}
Let $x\in\Sol\setminus\Z$ and let $\alpha_{1},\ldots,\alpha_{t}\in\Apos$ satisfy
$\alpha_{i}^{x}\in\Qbar$.  Then every $\alpha_{i}$ lies in $\GammaSat$ and the multiplicative
rank of $\alpha_{1},\ldots,\alpha_{t}$ is at most two.  \statusnew
\end{corollary}

Both ingredients are kernel-verified --- membership in the hull by
\lean{AlgebraicBaseUnitPower} and the rank ceiling by \lean{AlgebraicRankThreeEquivalence} ---
but the combined family statement is recorded here at paper level.

\subsection{The divisible-hull formulation}

Stated without reference to a fixed counterexample, the kernel-verified form is a single
biconditional valid for every positive real algebraic base.  For $\gamma\in\Apos$, under
$2^{x},3^{x}\in\Z$,
\begin{equation}
  \gamma^{x}\in\Qbar
  \quad\Longleftrightarrow\quad
  x\in\Z
  \quad\text{or}\quad
  \gamma\in\sqrt[\infty]{\langle2,3\rangle_{\Z}}.
\label{sp:eq-hull}
\end{equation}
The arithmetic core is an exact combinatorial equivalence: $\gamma$ has \emph{no} positive
natural power in $\langle2,3\rangle_{\Z}$ precisely when the monomial map
$(i,j,k)\mapsto2^{i}3^{j}\gamma^{k}$ is injective on $\Nzero^{3}$.  The forward implication of
\eqref{sp:eq-hull} then feeds that injectivity into the generic algebraic determinant, and the
reverse implication descends algebraicity from a positive integral power.  Specializing
$\gamma$ to a natural number recovers the $3$-smooth classification of
Section~\ref{sec:smooth} with ``rational'' upgraded to ``algebraic'', so for a hypothetical
nonintegral solution the three properties $a^{x}\in\Z$, $a^{x}\in\Q$ and $a^{x}\in\Qbar$
coincide for $a\in\N_{>0}$, and hold exactly at $2,3$-smooth $a$.  Specializing to a positive
rational $q$ gives $q^{x}\in\Qbar\iff q\in\langle2,3\rangle_{\Z}$; there is no
algebraic-irrational middle case anywhere in the picture.

\subsection{The determinant engines behind the spectrum}

Two generic determinant theorems, neither of which mentions the bases $2,3$, carry the whole
classification.  The rational one states that if $a,b,c\in\Q_{>0}$ have injective nonnegative
monomials $(i,j,k)\mapsto a^{i}b^{j}c^{k}$ and $a^{x},b^{x},c^{x}\in\Q$, then $x\in\Q$.  The
algebraic one removes both restrictions: if $a,b,c$ are positive real algebraic numbers with
injective nonnegative monomials and $a^{x},b^{x},c^{x}$ are all algebraic, then $x\in\Q$.
Both allow bases below $1$ and negative exponents, and both use $1+|x|$ in the size estimate.

\begin{statusbox}{\statuslean{} ---
\lean{rational_of_three_rat_rpows_rational_of_monomial_injective} in
\lean{RationalSixExponentials}, and
\lean{rational_of_three_real_rpows_isAlgebraic_of_monomial_injective} in
\lean{AlgebraicSixExponentials}.}
\end{statusbox}

The step from rational to algebraic outputs is exactly the replacement of the trivial bound
$|D|\ge1$ for a nonzero integer determinant by a field-norm bound for a nonzero algebraic
integer determinant, which is the extension anticipated in Section~\ref{sec:smooth}.  The six
input and output values are placed in one number field $K$ with one common nonzero natural
scale $s$ making every entry lie in $\mathcal O_{K}$.  For side parameter $n$ the square
matrix has size $n^{6}$ and each entry acquires the factor $s^{n^{5}}$, so the determinant
acquires $s^{n^{11}}$; at every embedding other than the distinguished one the six factor
groups in an entry are bounded by $T^{6n^{5}}$, and the norm of the nonzero algebraic-integer
determinant gives
\[
  \frac{1}{\bigl((n^{6})!\,T^{6n^{11}}\bigr)^{[K:\Q]-1}}
  \ \le\ \bigl|\sigma(\det A)\bigr|.
\]
The verified estimates $(n^{6})!\le2^{n^{11}}$ and
$\bigl((n^{6})!\,T^{6n^{11}}\bigr)^{[K:\Q]-1}\le\bigl((2T^{6})^{[K:\Q]-1}\bigr)^{n^{11}}$ put
the entire conjugate and degree loss on the same $n^{11}$ scale as the old rational
denominator loss, where the existing explicit analytic margin already absorbs it.  No new
analytic estimate was required --- only honest bookkeeping.

\begin{statusbox}{\statuslean{} ---
\lean{one_div_conjugate_entry_bound_le_norm_det} and
\lean{one_div_entry_bound_le_scaled_norm_det} in \lean{AlgebraicIntegerDeterminant}, with the
number-field realizations \lean{AlgebraicOutputBridge.ofRealIsAlgebraic} and
\lean{FiniteAlgebraicOutputBridge.ofFinSixIsAlgebraic}.}
\end{statusbox}

\subsection{The strong spectrum via Baker and Roy}

Let $\LL$ denote the $\Qbar$-vector space generated by $1$ and by all logarithms of nonzero
algebraic numbers, with arbitrary complex branches allowed; for positive real algebraic
numbers we use the real logarithm, which is one of those branches.  Roy's strong six
exponentials theorem \cite{Roy1992,Waldschmidt2005} states that if three complex numbers are
linearly independent over $\Qbar$ and two more are linearly independent over $\Qbar$, then at
least one of the six pairwise products lies outside $\LL$.

\begin{theorem}[Strong algebraic-base spectrum]\label{sp:strong-spectrum}
Under the hypotheses of Theorem~\ref{sp:general-spectrum}, if
$\alpha\notin\GammaSat(a,b)$, then
\[
  x\log\alpha\notin\LL.
\]
In particular $\alpha^{x}$ is transcendental.  \statusnew
\end{theorem}

\begin{proof}
As in Theorem~\ref{sp:general-spectrum}, the logarithms $\log a,\log b,\log\alpha$ are
$\Q$-linearly independent, and Baker's homogeneous theorem \cite{Baker1975} upgrades that to
linear independence over $\Qbar$.  The number $x$ is transcendental: were it algebraic, its
irrationality together with $a\ne0,1$ would make $a^{x}$ transcendental by
Gelfond--Schneider.  Hence $1,x$ are linearly independent over $\Qbar$.  Apply Roy's theorem
to the rows $\log a,\log b,\log\alpha$ and the columns $1,x$.  Five products lie in $\LL$,
namely $\log a$, $\log b$, $\log\alpha$, $x\log a=\log(a^{x})$ and $x\log b=\log(b^{x})$.
Therefore the sixth product $x\log\alpha$ lies outside $\LL$.
\end{proof}

\begin{definition}\label{sp:log-strong}
Call a positive exponential value $e^{z}$ \emph{log-strongly transcendental} when
$z\notin\LL$.  This implies ordinary transcendence of $e^{z}$, since the logarithm of a
positive algebraic number belongs to $\LL$.
\end{definition}

So outside $\GammaSat$ the failure is not merely that $\alpha^{x}$ is transcendental: the
number $x\log\alpha$ is not even a $\Qbar$-linear combination of $1$ and logarithms of
algebraic numbers.  This is a genuinely stronger exclusion than
Corollary~\ref{sp:AE-spectrum}, and it is the form in which the spectrum interacts with
linear forms in logarithms.

\subsection{One base is enough: detectors and rank equivalences}

Because the spectrum is the same for every hypothetical nonintegral solution, a
\emph{single} auxiliary base outside $\GammaSat$ already decides the whole conjecture.

\begin{theorem}[Algebraic detector equivalence]\label{sp:detector}
Fix $\alpha\in\Apos\setminus\GammaSat$.  Then the Alaoglu--Erd\H{o}s conjecture is equivalent
to
\[
  \forall x\in\R,\quad 2^{x},3^{x}\in\Z\Longrightarrow \alpha^{x}\in\Qbar,
\]
and equally to the assertion that every such $x$ satisfies $x\log\alpha\in\LL$.  \statusnew
\end{theorem}

\begin{proof}
If the conjecture holds, every solution is a nonnegative integer and $\alpha^{x}$ is
algebraic, with $x\log\alpha\in\LL$.  Conversely a nonintegral solution would make
$\alpha^{x}$ transcendental by Corollary~\ref{sp:AE-spectrum} and would put $x\log\alpha$
outside $\LL$ by Theorem~\ref{sp:strong-spectrum}.
\end{proof}

This is a paper-level universalization of a pointwise locus that the kernel already knows.
Valid detectors include, for instance,
\[
  5,\qquad \sqrt5,\qquad 1+\sqrt2,\qquad \frac{7+3\sqrt5}{2},
\]
provided the chosen number does not accidentally lie in $\GammaSat$; for an algebraic integer
with a prime divisor outside $\{2,3\}$ that exclusion is immediate, and in general it is a
multiplicative independence check.  Nothing arithmetically special about $5$ is used.

Four instances of this phenomenon are themselves kernel-verified, and they are the most
quotable consequences of the whole algebraic-output layer.

\begin{resultbox}
\textbf{One-base equivalents of the conjecture.}  Each of the following is equivalent to the
Alaoglu--Erd\H{o}s conjecture, and each equivalence is accepted by the Lean kernel.
\begin{enumerate}[label=\textup{(\roman*)}]
\item Every $x$ with $2^{x},3^{x}\in\Z$ has $5^{x}$ algebraic.
\item Every $x$ with $2^{x},3^{x}\in\Z$ has $6^{x^{2}}$ algebraic.
\item For every $x$ with $2^{x},3^{x}\in\Z$, the positive-algebraic output locus
\[
  \{\alpha\in\Apos:\alpha^{x}\in\Qbar\}
\]
has multiplicative rank three, i.e.\ contains $a,b,c$ with injective nonnegative monomials.
\item For every $x$ with $2^{x},3^{x}\in\Z$, the rational squared-output ratio locus has rank
two, i.e.\ contains algebraic members $(2^{i}/3^{j})^{x^{2}}$ and $(2^{k}/3^{\ell})^{x^{2}}$
with $i\ell\ne kj$.
\end{enumerate}
A counterexample would collapse (iii) to rank at most two and confine the entire locus in
(iv) to a single rational ray in the exponent lattice.  \statuslean
\end{resultbox}

\begin{statusbox}{\statuslean{} --- all four equivalences are accepted by the kernel; the
Lean names follow.}
Items (i) and (ii) are \lean{alaogluErdosConjecture_iff_five_rpow_isAlgebraic} and
\lean{alaogluErdosConjecture_iff_six_squared_exponent_isAlgebraic} in
\lean{AlgebraicConsequences}, which also carries the two-iterate variant
\lean{alaogluErdosConjecture_iff_both_iterated_outputs_isAlgebraic}.  Item (iii) is
\begin{center}
\lean{alaogluErdosConjecture_iff_algebraicRpowOutputLocusHasRankThree}
\end{center}
in \lean{AlgebraicRankThreeEquivalence}, and item (iv) is
\begin{center}
\lean{alaogluErdosConjecture_iff_ratioSquaredAlgebraicLocusHasRankTwo}
\end{center}
in \lean{AlgebraicOutputLocus}.
The same modules also record the pointwise forms
\lean{ratioSquaredAlgebraicLocusHasRankTwo_iff_integer} and
\lean{algebraicRpowOutputLocusHasRankThree_iff_integer}, and the explicit $p$-adic line
\lean{exists_external_prime_algebraic_output_ratio_padic_line} on which a counterexample would
force every algebraic ratio output to sit.
\end{statusbox}

\subsection{Exact exponent, rational-function, and tower loci}
\label{sec:exact-exponent-loci}

Fix a hypothetical nonintegral solution $x$ and put
\[
 \mathcal E_x=\{y\in\R:2^y,3^y\in\Qbar\}.
\]
The transposed six-exponentials determinant determines this set exactly:
\begin{equation}
 \mathcal E_x=\Q+\Q x.
 \label{sp:eq-exponent-plane}
\end{equation}
Thus every exponent outside the affine plane makes at least one of $2^y$ and $3^y$
transcendental.  In particular $y=x^n$ for every $n\ge2$, and also $y=x^{-1}$, lie outside
the plane because $x$ is transcendental.  For every fixed $n\ge2$, requiring both
$2^{x^n}$ and $3^{x^n}$ to be algebraic at every original two-base solution is therefore
equivalent to Conjecture~\ref{conj:AE}.  The theorem says ``at least one'': it neither selects
which base is transcendental nor supplies the missing algebraicity needed to prove the
conjecture.

The exact plane makes every genuinely nonaffine rational function an equally sharp detector.
For $P,Q\in\Q[X]$ with $Q(x)\ne0$,
\begin{equation}
 2^{P(x)/Q(x)},3^{P(x)/Q(x)}\in\Qbar
 \quad\Longleftrightarrow\quad
 P=(q+rX)Q\quad\text{for some }q,r\in\Q.
 \label{sp:eq-rational-function-locus}
\end{equation}
The reverse direction is literal affine membership in \eqref{sp:eq-exponent-plane}; the
forward direction uses transcendence of $x$ to turn equality after evaluation into a
polynomial identity.  Consequently any fixed nonaffine rational function represented as
$P/Q$ with the displayed denominator $Q$ having no real zero gives a global formulation
equivalent to Conjecture~\ref{conj:AE}, and every
$P\in\Q[X]$ of degree at least two does so with $Q=1$.  These are detector equivalences, not
new algebraicity results: they relocate the open conclusion to another exponent without
proving it.

Under failure, the same module chooses the canonical generator $\beta$ so that the two loci
can be compared on one rational-function chart.  For every $P/Q$ defined at $\beta$, both
$2^{P(\beta)/Q(\beta)}$ and $3^{P(\beta)/Q(\beta)}$ are algebraic exactly when
$P=(q+rX)Q$ with $q,r\in\Q$, whereas $P(\beta)/Q(\beta)$ is itself an original integral
solution exactly when $P=(n+kX)Q$ with $n,k\in\Nzero$.  Thus the exact solution monoid is the
nonnegative integral cone inside the rational affine algebraic-output plane; the comparison
is exact, but it does not collapse the larger plane to that cone.

There is also an exact two-dimensional picture for the squared monomials.  Under a
counterexample, the set
\[
 \mathcal F_x=\{(i,j)\in\Nzero^2:
      (2^i3^j)^{x^2}\in\Qbar\}
\]
is exactly one of
\[
 \Nzero\times\{0\},\qquad \{0\}\times\Nzero,\qquad \{(0,0)\}.
\]
Writing $M=2^x$ and $B=3^x$, the three cases say respectively that $M$ is $2,3$-smooth and
$B$ is not, that $B$ is $2,3$-smooth and $M$ is not, or that neither is; both cannot be
$2,3$-smooth for a nonintegral solution.  Under the original two integrality hypotheses this
coordinate-face property is pointwise equivalent to nonintegrality, and the conjecture is
equivalent to excluding it for every solution.  All three faces remain compatible with the
present theory, so the trichotomy does not itself yield a contradiction.

A related transition theorem has a different, solution-independent plane.  With
$\thetaq=\log_2 3$, define
\[
 \mathcal T_x=\{y\in\R:2^y\in\Qbar\text{ and }2^{xy}\in\Qbar\}.
\]
For every hypothetical nonintegral solution,
\begin{equation}
 \mathcal T_x=\Q+\Q\thetaq.
 \label{sp:eq-tower-transition}
\end{equation}
Two consecutive algebraic base-two tower levels therefore correspond exactly to one exponent
in this fixed plane.  More strongly, at most one positive depth $n$ can satisfy
$x^n\in\mathcal T_x$.  Indeed, transcendence of $x$ makes the $\thetaq$-coefficients in the
two affine expressions nonzero; raising the two expressions to cross-powers and using
transcendence of $\thetaq$ promotes equality at $\thetaq$ to a polynomial identity, whose
degree comparison forces the two depths to coincide.  Hence every three consecutive
positive-depth levels
\[
 2^{x^n},\qquad 2^{x^{n+1}},\qquad 2^{x^{n+2}}
\]
contain a transcendental value, and over any finite set the collection of depths starting an
adjacent algebraic pair has cardinality at most one.  Conversely, for any fixed distinct
positive depths $n\ne m$, the assertion that every original two-base solution has algebraic
adjacent pairs at both depths $n$ and $m$ is equivalent to Conjecture~\ref{conj:AE}.  The
restriction still permits isolated algebraic levels and a single adjacent algebraic pair; it
does not prove the full conjecture.

The symmetric transition picture is stronger than merely repeating the same argument with
base three.  Put
\[
 \mathcal S_x=\{y\in\R:3^y\in\Qbar\text{ and }3^{xy}\in\Qbar\}.
\]
Under a nonintegral solution,
\begin{equation}
 y\in\mathcal S_x
 \quad\Longleftrightarrow\quad
 \thetaq y\in\Q+\Q\thetaq.
 \label{sp:eq-three-tower-transition}
\end{equation}
No positive depth in the base-two tower can start an algebraic adjacent pair simultaneously
with any positive depth in the base-three tower, even when the two depths are unrelated.
The base-three tower itself has at most one such positive depth.  Consequently, across the
two principal towers together there is at most one positive transition depth, and at that
depth at most one of the two bases can occur.  There is also a direction-uniform version in
the multiplicatively dependent-output branch: once its kernel relation produces a nonzero
rational $2,3$-unit direction with a depth-one transition, every nonzero rational
$2,3$-unit direction is excluded at every depth $n\ge2$.  These are strict structural
obstructions, but they still allow the single depth-one transition forced by the kernel
relation and therefore do not close that branch.

\begin{statusbox}{\statuslean{} --- all statements in this subsection are accepted by the
kernel.}
The exact exponent plane and its natural-power and reciprocal consequences are
\lean{TwoBaseNonintegerSolution.twoThreeRpowAlgebraic_iff_affine},
\lean{transcendental_two_or_three_rpow_nat_pow},
\lean{transcendental_two_or_three_rpow_inv}, and
\lean{alaogluErdosConjecture_iff_twoThreeRpowAlgebraic_nat_pow} in
\lean{AlgebraicExponentLocus}.  Equation~\eqref{sp:eq-rational-function-locus} and its global
and polynomial specializations are
\lean{TwoBaseNonintegerSolution.twoThreeRpowAlgebraic_rationalFunction_iff},
\lean{alaogluErdosConjecture_iff_twoThreeRpowAlgebraic_rationalFunction}, and
{\small\lean{alaogluErdosConjecture_iff_twoThreeRpowAlgebraic_polynomial_of_two_le_natDegree}}
in
\lean{RationalFunctionAlgebraicLocus}.  Its simultaneous two-locus package is
\lean{exists_generator_with_exact_rationalFunction_algebraic_and_integral_loci}.

The exact squared-coordinate-face statements are
\lean{TwoBaseNonintegerSolution.exists_exact_algebraic_output_monomial_locus},
\lean{algebraic_two_three_squared_monomial_locus_trichotomy},
\lean{squaredTwoThreeMonomialAlgebraicLocusIsCoordinateFace_iff_not_integer}, and
\lean{alaogluErdosConjecture_iff_no_squaredTwoThreeMonomial_coordinateFace} in
\lean{AlgebraicOutputLocus}.  Finally, \eqref{sp:eq-tower-transition}, global uniqueness of a
transition depth, and the fixed-distinct-depth equivalence are
\lean{TwoBaseNonintegerSolution.twoRpowAlgebraicTransition_iff_affine_logRatio},
\lean{transcendental_one_of_three_consecutive_two_rpow_tower},
\lean{not_two_distinct_algebraic_tower_transitions},
\lean{card_algebraicTowerTransitionDepths_le_one}, and
\lean{alaogluErdosConjecture_iff_two_distinct_tower_transitions_algebraic} in
\lean{AlgebraicTowerTransition}.  The symmetric locus
\eqref{sp:eq-three-tower-transition}, the cross-base exclusion, the combined uniqueness
theorem, and the rational-unit-direction exclusions are respectively
\lean{TwoBaseNonintegerSolution.threeRpowAlgebraicTransition_iff_scaled_affine_logRatio},
\lean{not_baseTwo_and_baseThree_tower_transitions},
\lean{eq_of_twoOrThree_tower_transitions}, and
\lean{not_higher_unitDirectionTowerTransition_of_kernelRelation} in
\lean{CrossBaseTowerTransition}.  Their axiom audit uses only \lean{propext},
\lean{Classical.choice}, and \lean{Quot.sound}.  None asserts Conjecture~\ref{conj:AE}.
\end{statusbox}

There is a corresponding zero-one law: assuming the conjecture fails and fixing
$\alpha\in\Apos$, exactly one of two things happens.  Either $\alpha\in\GammaSat$ and
$\alpha^{x}$ is algebraic for \emph{every} nonintegral solution $x$, or
$\alpha\notin\GammaSat$ and $\alpha^{x}$ is transcendental for every nonintegral solution.  No
base can behave one way for one counterexample and the other way for another.

\subsection{Why the exact spectrum does not settle the conjecture}

It is worth being precise about what has and has not been gained.  The spectrum is an exact
answer to the question ``which algebraic bases stay algebraic?'', and its answer is
$\GammaSat=\{2^{r}3^{s}:r,s\in\Q\}$ --- but that is precisely the set of bases whose
logarithms lie in the rational span of $\log2$ and $\log3$.  Every base the spectrum
certifies is, logarithmically, a rational recombination of the two columns already present in
the four-exponentials matrix of Section~\ref{sec:fourexp}.  Six exponentials needs a third
row genuinely independent of $\log2,\log3$ together with algebraicity of its output; the
spectrum theorem shows that a counterexample supplies algebraic outputs at exactly the bases
that fail this requirement, and transcendental outputs at exactly the bases that would meet
it.  The classification is therefore self-consistent with a counterexample rather than
obstructive to one: it converts every attempt to manufacture a third independent base into
the very transcendence statement one was trying to prove.  The kernel-verified radical work
makes the same wall concrete: the simultaneous normalization of the two canonical radicals
produces an explicit multiplicative relation among $2$, the primitive odd core, and the
swapped radical, so that the most tempting three-base application of the algebraic
six-exponentials theorem fails its monomial-injectivity hypothesis for a structural,
kernel-verified reason (\lean{swappedRadical_relation_blocks_monomial_injectivity} in
\lean{CanonicalRadicalFieldNorm}).  Closing the gap needs a new
source of size --- nonarchimedean, or from a sharper five/six-exponentials interface --- not
a wider supply of algebraic bases.  The conjecture remains open.

\section{Rational and integral outputs inside the spectrum}\label{sec:valuationlattice}

The algebraic spectrum identifies the divisible group
\[
\GammaSat=\{2^r3^s:r,s\in\Q\}
\]
as the exact set of positive algebraic bases whose powers at a hypothetical nonintegral
solution remain algebraic: outside $\GammaSat$ every such power is transcendental, and inside
$\GammaSat$ every such power is algebraic.  A finer arithmetic question survives inside the
spectrum.  If $\alpha=2^r3^s\in\GammaSat$, when is $\alpha^x$ rational, and when is it a
positive integer?

For a fixed solution the answer is completely encoded by prime valuations of the output pair.
Rationality becomes a congruence condition, integrality adds a cone condition, and the gap
between the obvious rational bases $2^u3^v$ with $u,v\in\Z$ and the full rational locus is one
explicit finite abelian group whose order is a single gcd of $2\times2$ minors.  This is a
strict refinement of the divisible-hull description of the algebraic locus, and it is exact
rather than asymptotic.

\subsection{The valuation matrix of an output pair}

Let $x\in\Sol\setminus\Z$ be a nonintegral solution and write
\[
  M=2^x\in\N,\qquad A=3^x\in\N .
\]
Let $P$ be the finite union of the prime supports of $M$ and $A$, and form the column vectors
\[
  \mathbf m=(\vpp{p}{M})_{p\in P},
  \qquad
  \mathbf a=(\vpp{p}{A})_{p\in P},
  \qquad
  V_x=\begin{pmatrix}\mathbf m&\mathbf a\end{pmatrix}\in M_{P\times2}(\Z).
\]
For a column $q=(r,s)^{T}\in\Q^2$ put $\alpha(q)=2^r3^s$, so that $\GammaSat$ is exactly the
image of $q\mapsto\alpha(q)$.  Taking the unique positive real interpretation of the rational
powers,
\begin{equation}
  \alpha(q)^x=\bigl(2^x\bigr)^r\bigl(3^x\bigr)^s=M^rA^s
  =\prod_{p\in P}p^{\,r\vpp{p}{M}+s\vpp{p}{A}} .
  \label{vl:eq-output}
\end{equation}
The exponent vector appearing in \eqref{vl:eq-output} is precisely $V_xq$.

\begin{lemma}[The valuation matrix has rank two]
\label{vl:lem-rank-two}
The two columns of $V_x$ are linearly independent over $\Q$.
\end{lemma}

\begin{proof}
Suppose $r\mathbf m+s\mathbf a=0$ with $r,s\in\Q$.  Then every prime exponent in
\eqref{vl:eq-output} vanishes, so $M^rA^s=1$.  Taking real logarithms and using
$\log M=x\log2$, $\log A=x\log3$ gives $x(r\log2+s\log3)=0$.  Since a solution with
$M=2^x\in\N$, $M>1$ has $x>0$, and since $\thetaq=\log_2 3$ is irrational, $r=s=0$.
\end{proof}

\begin{proposition}[Exact rational and integral auxiliary-output spectrum]
\label{vl:prop-spectrum}
Let $x\in\Sol\setminus\Z$, let $V_x$ be as above, and let $q\in\Q^2$.  Then
\[
 \alpha(q)^x\in\Q_{>0}
 \iff V_xq\in\Z^{P},
 \qquad\qquad
 \alpha(q)^x\in\N
 \iff V_xq\in\Nzero^{P}.
\]
Together with the algebraic spectrum, this classifies algebraic, rational and integral outputs
for every positive algebraic auxiliary base.  \statusnew
\end{proposition}

\begin{proof}
Equation \eqref{vl:eq-output} exhibits $\alpha(q)^x$ as a finite product of primes with
rational exponents.  Choose a common denominator $d$ for the entries of $V_xq$.  If
$\alpha(q)^x$ is rational, then raising to the $d$-th power gives a rational number whose
ordinary $p$-adic valuation equals $d$ times the $p$-th entry of $V_xq$; comparing valuations
of numerator and denominator shows every entry of $V_xq$ is an integer.  Conversely, integer
exponents produce a positive rational number.  Such a rational number is a positive integer
exactly when all of its prime valuations are nonnegative, which is the second equivalence.
\end{proof}

Three nested spectra therefore occur, all determined by $V_x$:
\[
 \begin{array}{rcl}
 \mathcal B_{\mathrm{alg}}(x)
  &=&\{\alpha(q):q\in\Q^2\}=\GammaSat,\\[2mm]
 \mathcal B_{\mathrm{rat}}(x)
  &=&\{\alpha(q):q\in\Lambda_x\},\\[1mm]
 \mathcal B_{\mathrm{int}}(x)
  &=&\{\alpha(q):q\in\Lambda_x^{+}\},
 \end{array}
 \qquad
 \begin{aligned}
 \Lambda_x&=\{q\in\Q^2:V_xq\in\Z^{P}\},\\[1mm]
 \Lambda_x^{+}&=\{q\in\Lambda_x:V_xq\ge0\}.
 \end{aligned}
\]
The outer spectrum does not depend on $x$; the two inner ones may.

\subsection{Smith normal form and the finite rational-output defect}

By Lemma~\ref{vl:lem-rank-two} the matrix $V_x$ has rank two, so $\Lambda_x$ is a rank-two
lattice in $\Q^2$ containing $\Z^2$ with finite index.  Write
\[
 \delta_1(V_x)=\gcd\{\text{all entries of }V_x\},
 \qquad
 \delta_2(V_x)=\gcd\left\{
  \left|\det\begin{pmatrix}
   \vpp{p}{M}&\vpp{p}{A}\\
   \vpp{p'}{M}&\vpp{p'}{A}
  \end{pmatrix}\right|
  :p,p'\in P\right\}
\]
for the two determinantal divisors.  The second is positive precisely because the rank is two.

\begin{proposition}[Smith classification of the rational locus]
\label{vl:prop-smith}
Let $d_1\mid d_2$ be the Smith invariant factors of $V_x$, so that $d_1=\delta_1(V_x)$ and
$d_1d_2=\delta_2(V_x)$.  After a unimodular change of coordinates in $\Q^2$,
\[
  \Lambda_x=\tfrac1{d_1}\Z\oplus\tfrac1{d_2}\Z ,
\]
and consequently
\[
  \Lambda_x/\Z^2\cong\Z/d_1\Z\oplus\Z/d_2\Z,
  \qquad
  \bigl|\Lambda_x/\Z^2\bigr|=\delta_2(V_x).
  \qquad\statusnew
\]
\end{proposition}

\begin{proof}
Choose unimodular matrices $U\in\mathrm{GL}_P(\Z)$ and $W\in\mathrm{GL}_2(\Z)$ with
\[
 UV_xW=\begin{pmatrix}d_1&0\\0&d_2\\0&0\\\vdots&\vdots\end{pmatrix}.
\]
Since $U$ preserves $\Z^{P}$, the condition $V_xq\in\Z^{P}$ is equivalent to
$UV_xq\in\Z^{P}$, and substituting $q=Wq'$ turns it into $d_1q'_1\in\Z$ and $d_2q'_2\in\Z$.
Since $W$ preserves $\Z^2$, it identifies $\Lambda_x/\Z^2$ with the quotient computed in the
primed coordinates, which is $\Z/d_1\Z\oplus\Z/d_2\Z$.  The identifications
$d_1=\delta_1(V_x)$ and $d_1d_2=\delta_2(V_x)$ are the standard determinantal-divisor
formulae for the Smith normal form.
\end{proof}

The first invariant factor has a direct output interpretation: $d_1$ is the largest integer
$d$ such that $M$ and $A$ are both perfect $d$-th powers, because $d\mid\delta_1(V_x)$ says
exactly that $d$ divides every prime valuation of $M$ and of $A$.  For the canonical least
nonintegral solution $\beta$, primitivity of the least output pair (Section~\ref{sec:conditional})
gives $d_1=1$, and then
\[
  \Lambda_\beta/\Z^2\cong\Z/\delta_2(V_\beta)\Z .
\]
Thus a single integer --- the gcd of all $2\times2$ valuation minors of the least output pair
--- measures the entire finite defect between the rational $\beta$-outputs at algebraic
$2,3$-radical bases and the obvious rational $2,3$-unit bases.  We call $\delta_2(V_\beta)$
the \emph{rational-output defect}.  It is a finite, in principle computable, invariant of a
hypothetical counterexample, and it equals $1$ exactly when no auxiliary base outside
$\{2^u3^v:u,v\in\Z\}$ produces a rational output.

The integral spectrum is the intersection of this lattice with the rational polyhedral cone
$V_xq\ge0$.  By Gordan's lemma $\Lambda_x^{+}$ is a finitely generated normal affine
semigroup, so $\mathcal B_{\mathrm{int}}(x)$ has a finite generating set of auxiliary bases.
This is an exact structural replacement for attempting to enumerate algebraic auxiliary bases
one at a time.

\begin{statusbox}{\statuslean{} and \statusnew{}}
The lattice statements themselves --- Lemma~\ref{vl:lem-rank-two},
Proposition~\ref{vl:prop-spectrum}, Proposition~\ref{vl:prop-smith} and the Gordan
consequence --- are \statusnew{}: they are proved on paper in this report and are not yet
formalized in that form.  What \emph{is} kernel-verified is the arithmetic core on which they
rest, in three modules.

\lean{CanonicalRadicalValuation} proves the divisibility dictionary between perfect powers and
valuations that underlies the meaning of $d_1$: a divisor of the odd-prime valuation gcd of a
positive integer is an honest exponent of an exact power
(\lean{exists_pow_eq_of_dvd_oddPrimeValuationGCD}), the least rational-power index is coprime
to the joint gcd of the denominator exponent and the numerator's valuation gcd
(\lean{rationalPowerIndex_gcd_denominator_valuationGCD_eq_one}), and every prime dividing that
index leaves some prime valuation of the normalized numerator undivided, giving the exact
radical degree together with the valuation restriction
(\lean{exists_exact_radical_degree_and_valuation_restriction}).

\lean{SimultaneousRadicalDegrees} kernel-verifies the exactness statements that the $d_1=1$
normalization buys: under the simultaneous primitive normalization the displayed exponents are
the least rational-power indices and hence the exact degrees over $\Q$
(\lean{oddCoreRpow_exact_rationalPowerIndex_of_simultaneous_normalization} and
\lean{threeFreeCoreRpow_exact_radical_degree_of_simultaneous_normalization}), and when both
base-adic depths vanish the two indices are coprime and the mixed product has exact degree
equal to their product (\lean{mixed_coreRpow_exact_radical_degree_of_zero_baseDepths}).  The
descent used there is precisely the one behind $d_1=1$: a proper index divisor would make both
outputs of the least nonintegral solution common perfect powers, which is excluded by
\lean{IsLeastTwoBaseNonintegerSolution.no_common_perfect_power} in \lean{LeastSolution}
(Section~\ref{sec:conditional}).  That exclusion is exactly the assertion $d_1=1$ for
$V_\beta$.

\lean{CanonicalRadicalMixedOutput} kernel-verifies the transcendence side of the mixed case: a
genuinely mixed positive monomial in the two canonical radicals retains an external prime in
every rational power, so all of its outputs at the least nonintegral solution are
transcendental, and an algebraic mixed output is therefore forced onto a coordinate face.  The
theorem is
\begin{center}
\lean{TwoBaseNonintegerSolution.transcendental_mixed_canonical_radical_output},
\end{center}
packaged under failure of the conjecture as
\begin{center}
\lean{exists_simultaneous_canonical_radical_mixed_output_transcendence}.
\end{center}
These are statements about the canonical radicals rather than about $\Lambda_x$ itself; they
confirm the geometry of the lattice picture without yet formalizing it.
\end{statusbox}

\begin{resultbox}
\textbf{Exact simultaneous radical degrees and mixed-output exclusion.}
Under failure of the conjecture, let $\beta$ be the least nonintegral solution and choose its
simultaneous primitive normalization
\[
 2^\beta=2^a w^d,\qquad 3^\beta=3^b v^e,
 \qquad a=0\ \text{or}\ b=0.
\]
For
\[
 E=w^{\thetaq},\qquad F=v^{1/\thetaq},
\]
the simultaneous power identities are
\[
 E^d=\frac{3^b v^e}{3^a},
 \qquad
 F^e=\frac{2^a w^d}{2^b}.
\]
The least positive rational-power indices of $E$ and $F$ are exactly $d$ and $e$.
Consequently their minimal polynomials are the displayed pure binomials
\[
 X^d-\frac{3^b v^e}{3^a},
 \qquad
 X^e-\frac{2^a w^d}{2^b},
\]
and their degrees over $\Q$ are exactly $d$ and $e$.  In the zero-depth case $a=b=0$ one has
$\gcd(d,e)=1$, and for $G=EF$,
\[
 \operatorname{minpoly}_{\Q}(G)=X^{de}-v^{e^2}w^{d^2},
 \qquad [\Q(G):\Q]=de.
\]
Moreover, the canonical data supplied by failure contain an external prime, and for every
$i,j>0$ the mixed output $(E^iF^j)^\beta$ is transcendental; if a nonnegative mixed output is
algebraic, its exponent pair lies on one of the two coordinate axes.  \statuslean
\end{resultbox}

These conclusions are mutually compatible, not contradictory.  The degree statements concern
the algebraic numbers $E$, $F$, and $EF$, whereas the real irrational-power outputs at
$\beta$ are transcendental.  A field embedding of $\Q(E,F)$ need not commute with the chosen
real operation $z\mapsto z^\beta$, so taking conjugates or norms does not turn the mixed-output
transcendence into a forbidden algebraic relation.  In the zero-depth case the product degree
$de$ is the maximal degree allowed by the two coprime pure-radical degrees, not an impossible
dependence.  Thus the simultaneous theory sharpens the exact algebraic data and excludes the
interior of the mixed-output quadrant, but leaves both coordinate axes --- and the full
Alaoglu--Erd\H{o}s conjecture --- open.

\subsection{Two translated solutions remove the fractional defect}

The algebraic spectrum is solution-independent, but the rational locus $\Lambda_x$ need not
be.  Translation of the exponent by one gives a sharp intersection identity, and it shows that
the defect cannot persist across the whole solution monoid.

\begin{proposition}[Translation intersection]
\label{vl:prop-translation}
Let $x\in\Sol\setminus\Z$.  Then
\[
  \Lambda_x\cap\Lambda_{x+1}=\Z^2 .
\]
Equivalently, for $\alpha\in\Apos$, if $\alpha^x$ and $\alpha^{x+1}$ are both rational, then
$\alpha=2^u3^v$ with $u,v\in\Z$.  \statusnew
\end{proposition}

\begin{proof}
One inclusion is immediate: for $q\in\Z^2$ the base $\alpha(q)=2^u3^v$ has
$\alpha(q)^y=M^uA^v\in\Q_{>0}$ for every solution $y$, in particular for $y=x$ and $y=x+1$.
Conversely, suppose $\alpha^x$ and $\alpha^{x+1}$ are rational.  Both are algebraic, so the
algebraic spectrum gives $\alpha=2^r3^s$ with $r,s\in\Q$, and
\[
  \alpha=\frac{\alpha^{x+1}}{\alpha^x}\in\Q_{>0}.
\]
A positive rational of the form $2^r3^s$ with $r,s\in\Q$ has $r,s\in\Z$, by comparing ordinary
$2$-adic and $3$-adic valuations after clearing a common denominator.
\end{proof}

\begin{corollary}[Universal rational and integral spectra under failure]
\label{vl:cor-universal}
Assume Conjecture~\ref{conj:AE} fails and let $\alpha\in\Apos$.  \statusnew
\begin{enumerate}
\item $\alpha^y$ is rational for every nonintegral solution $y$ if and only if
      $\alpha=2^u3^v$ with $u,v\in\Z$.
\item $\alpha^y$ is a positive integer for every solution $y$ if and only if
      $\alpha=2^u3^v$ with $u,v\in\Nzero$.
\end{enumerate}
\end{corollary}

\begin{proof}
For the first part, apply Proposition~\ref{vl:prop-translation} to a nonintegral solution $x$
and to $x+1$, which is again a nonintegral solution.  Conversely, rational $2,3$-units have
rational outputs $M^uA^v$ at every solution.

For the second part, assume $\alpha^y\in\N$ for every $y\in\Sol$.  Taking $y=x$ and $y=x+1$
makes $\alpha=\alpha^{x+1}/\alpha^x$ a positive rational, say $a/b$ in lowest terms.  Put
$z=\alpha^x\in\N$.  Every $x+n$ with $n\ge0$ is again a solution, so the number
$z(a/b)^n=\alpha^{x+n}$ is an integer for every $n$, whence
$b^n$ divides the fixed integer $z$ for every $n$, forcing $b=1$ and $\alpha\in\N$.  Since
$\alpha^x$ is an integer, the algebraic spectrum gives $\alpha\in\GammaSat$, and an integer in
$\GammaSat$ is exactly $2^u3^v$ with $u,v\in\Nzero$.  The converse is immediate.
\end{proof}

Corollary~\ref{vl:cor-universal} says that the rational-output defect is a genuinely
per-solution phenomenon.  A base can enjoy a rational output at one nonintegral solution only
by failing at the next one, unless it is a rational $2,3$-unit already.

\subsection{The two smooth branches are mutually exclusive}

The number of rows of $V_x$ is another exact invariant of a counterexample, and the smallest
conceivable case is excluded outright.  Call a positive integer $\{2,3\}$-smooth (equivalently
$3$-smooth) when it is divisible by no prime outside $\{2,3\}$.  In the notation above, $M$ is
$\{2,3\}$-smooth exactly when the column $\mathbf m$ is supported on $\{2,3\}$, and both
outputs are $\{2,3\}$-smooth exactly when $P\subseteq\{2,3\}$, that is, when $V_x$ is a
$2\times2$ integer matrix.

\begin{proposition}[The smooth cases are mutually exclusive]
\label{vl:prop-not-both-smooth}
For a nonintegral solution $x$, the outputs $M=2^x$ and $A=3^x$ cannot both be
$\{2,3\}$-smooth.  Equivalently, the valuation matrix $V_x$ of a counterexample always has a
row indexed by a prime $p\ge5$.  \statuslean
\end{proposition}

\begin{proof}
Suppose $M=2^a3^b$ and $A=2^c3^e$ with $a,b,c,e\in\Nzero$.  Taking base-two logarithms of
$2^x=M$ and of $3^x=A$ gives the two relations
\[
  x=a+b\thetaq,\qquad x\thetaq=c+e\thetaq ,\qquad \thetaq=\log_23 .
\]
Substituting the first into the second eliminates $x$ and leaves
\[
  b\,\thetaq^{2}+(a-e)\,\thetaq-c=0 .
\]
Were any of $b$, $a-e$, $c$ nonzero, this would exhibit $\thetaq$ as a root of a nonzero
rational polynomial of degree at most two, contradicting its transcendence.  Hence $b=c=0$ and
$a=e$, so $x=a\in\Z$, contrary to hypothesis.
\end{proof}

\begin{corollary}[A smooth output has a rough partner]
\label{vl:cor-rough-partner}
Let $x$ be a nonintegral solution with $M=2^x$ $\{2,3\}$-smooth.  Then $A=3^x$ has a prime
divisor $p\ge5$.  \statuslean
\end{corollary}

\begin{proof}
If every prime divisor of $A$ were $2$ or $3$, then $A$ would be $\{2,3\}$-smooth, and
Proposition~\ref{vl:prop-not-both-smooth} would force $x\in\Z$.
\end{proof}

\begin{statusbox}{\statuslean{}}
\lean{integer_of_threeSmooth_outputs} and
\lean{exists_prime_five_le_dvd_of_threeSmooth_candidate} in \lean{SmoothOutputs}, with the
integral restatement \lean{integer_of_two_three_rpow_integer_of_threeSmooth_outputs}.  The
proofs use only the transcendence of $\thetaq$, itself kernel-verified in this corpus from
Gelfond--Schneider.  The same statement is also derived independently at paper level in the
research report by the identical degree-drop elimination, so the formalized and unformalized
routes confirm each other.
\end{statusbox}

Two remarks fix the position of this exclusion.

\emph{The two-sided hypothesis is essential.}  Assuming only that $M$ is $\{2,3\}$-smooth
gives $x=a+b\thetaq$ and then requires $\bigl(3^{\thetaq}\bigr)^{b}$ to be rational up to a
power of $3$, which is exactly the open localized-radical condition of
Section~\ref{sec:radical}.  What the second smoothness hypothesis supplies is a \emph{second}
linear relation between $x$ and $\thetaq$; two such relations eliminate $x$ and leave an
algebraic equation in $\thetaq$ alone.  This degree drop is the entire mechanism, and it is
worth isolating as a template: each independent multiplicative constraint on an output buys
one linear relation, and two relations force algebraicity of $\thetaq$.  The difficulty of the
general problem is precisely that a second constraint is never free.

\emph{The exclusion feeds back into the lattice.}  Because $P\not\subseteq\{2,3\}$, the
valuation matrix of a counterexample is never square, so $\delta_2(V_x)$ is a genuine gcd over
at least three rows rather than a single determinant, and the rational locus $\Lambda_x$ is cut
out by at least three congruences.  In particular the rational-output defect is bounded by the
absolute value of any one $2\times2$ minor of $V_x$, and any two minors that are coprime force
$\delta_2(V_x)=1$, in which case no auxiliary base outside $\{2^u3^v:u,v\in\Z\}$ has a rational
output at $x$ at all.

\section{The second-iterate kernel and multiplicative dependence}
\label{sec:kernel}

The algebraic spectrum of Section~\ref{sec:spectrum} controls one application of the exponent
$x$.  Asking which bases of $\GammaSat$ \emph{return} to $\GammaSat$ after that application
turns second iterates into finite-dimensional linear algebra, and --- as this section records
--- into an entirely elementary statement about four integers.

\subsection{The kernel and its invariance}

For a nonintegral solution $x$ with outputs $M=2^x$ and $A=3^x$, set
\begin{equation}
  K_x=\bigl\{(r,s)\in\Q^2\ :\ M^rA^s\in\GammaSat\bigr\},
  \qquad \GammaSat=\{2^u3^v:u,v\in\Q\}.
  \label{eq:kn-def}
\end{equation}
Writing $\mathbf m_{\mathrm{ext}}=(\vpp pM)_{p\ne2,3}$ and
$\mathbf a_{\mathrm{ext}}=(\vpp pA)_{p\ne2,3}$ for the external valuation vectors, membership
of $M^rA^s$ in $\GammaSat$ is exactly the vanishing of every prime valuation outside
$\{2,3\}$, so
\begin{equation}
  K_x=\ker\bigl((r,s)\mapsto r\,\mathbf m_{\mathrm{ext}}+s\,\mathbf a_{\mathrm{ext}}\bigr).
  \label{eq:kn-valuations}
\end{equation}

\begin{proposition}[Kernel dimension]
\label{kn:prop-dim}
$\dim_\Q K_x\le1$; the kernel contains no vector with $rs>0$; it equals $\Q(1,0)$ exactly
when $M$ is $3$-smooth and $\Q(0,1)$ exactly when $A$ is $3$-smooth.
\end{proposition}

\begin{proof}
At least one external vector is nonzero: otherwise both $M$ and $A$ would be $3$-smooth,
which Theorem~\ref{lp:thm-smooth} excludes.  Hence the map in \eqref{eq:kn-valuations} has
rank at least one and its kernel has dimension at most one.  If $r,s>0$ then any external
prime dividing $M$ or $A$ has positive valuation in $M^rA^s$, and if $r,s<0$ it has negative
valuation, so no same-sign vector lies in the kernel.  The two smooth cases are immediate
from \eqref{eq:kn-valuations} by testing $(1,0)$ and $(0,1)$.
\end{proof}

\begin{statusbox}{\statuslean{} --- \lean{SecondIterateKernel} carries the kernel as an
integral subgroup \lean{kernelAddSubgroup} of $\Z^2$, with \lean{kernelPairs_mul_nonpos} for
the absence of same-sign vectors, \lean{kernelPairs_det_eq_zero} and
\lean{kernelPairs_prop_dim} for the rank bound, \lean{mem_kernelPairs_one_zero_iff} for the
two smooth cases, and \lean{saturatedKernel_invariant} for $K_x=K_y$.  The input that $M$ and
$A$ cannot both be $3$-smooth is \lean{integer_of_threeSmooth_outputs} in \lean{SmoothOutputs}.
Clearing denominators lets the formal statement use integer rather than rational pairs, which
is equivalent and lighter.}
\end{statusbox}

The kernel is not an artifact of the chosen counterexample.  If $x,y$ are two nonintegral
solutions and $\alpha=2^r3^s$, then $B=\alpha^{x}$ is algebraic, and applying the spectrum
theorem to $y$ gives $\alpha^{xy}=B^{y}\in\Qbar$ if and only if $B\in\GammaSat$, that is, if
and only if $(r,s)\in K_x$.  The left-hand condition is symmetric in $x$ and $y$, so
$K_x=K_y$: the kernel is an invariant of the problem, not of the solution.

\subsection{Multiplicative dependence of four integers}

The kernel condition has a formulation with no logarithms, no projective linear group, and no
transcendence vocabulary at all.  It is the observation of this section.

\begin{theorem}[Multiplicative form of the kernel dichotomy]
\label{kn:thm-multiplicative}
Let $x$ be a nonintegral solution with outputs $M=2^x$ and $A=3^x$.  The following are
equivalent:
\begin{enumerate}[label=\textup{(\roman*)}]
\item $K_x\ne\{0\}$;
\item the four integers $M$, $A$, $2$, $3$ are multiplicatively dependent, that is,
      \[
        M^{a}A^{b}2^{c}3^{d}=1
        \qquad\text{for some }(a,b,c,d)\in\Z^4\setminus\{0\};
      \]
\item $1$, $\thetaq$, $x$, $x\thetaq$ are linearly dependent over $\Q$;
\item $x$ lies in the $\PGL_2(\Q)$-orbit of $\thetaq$, that is,
      $x=(u+v\thetaq)/(r+s\thetaq)$ with rational $u,v,r,s$ and $us-vr\ne0$.
\end{enumerate}
\end{theorem}

\begin{proof}
(i) implies (ii).  If $(r,s)\in K_x$ is nonzero then $M^rA^s=2^u3^v$ with $u,v\in\Q$; clearing
a common denominator of $r,s,u,v$ produces integers with $M^aA^b2^{c}3^{d}=1$ and
$(a,b)\ne(0,0)$.

(ii) implies (i).  Given such a relation, note first that $(a,b)\ne(0,0)$: otherwise
$2^c3^d=1$, and unique factorization forces $c=d=0$, contradicting nontriviality.  Then
$M^aA^b=2^{-c}3^{-d}\in\GammaSat$, so $(a,b)\in K_x\setminus\{0\}$.

(ii) is equivalent to (iii).  Taking logarithms of $M^aA^b2^c3^d=1$ and using
$\log M=x\log2$ and $\log A=x\log3$ gives
\[
  (ax+c)\log2+(bx+d)\log3=0 ,
\]
and division by $\log2$ turns this into
\begin{equation}
  x\,(a+b\thetaq)+(c+d\thetaq)=0 ,
  \label{eq:kn-linear}
\end{equation}
which is precisely a nontrivial rational linear relation among $1,\thetaq,x,x\thetaq$.  Each
step reverses.

(iii) is equivalent to (iv).  In \eqref{eq:kn-linear} the coefficient $a+b\thetaq$ cannot
vanish, since $\thetaq$ is irrational and $(a,b)\ne(0,0)$; hence
$x=-(c+d\thetaq)/(a+b\thetaq)$.  If the determinant $ad-bc$ vanished, numerator and
denominator would be rationally proportional and $x$ would be rational, hence integral by
Lemma~\ref{pr:rational}, contrary to hypothesis.  Conversely, any such expression
cross-multiplies to \eqref{eq:kn-linear}.
\end{proof}

\begin{statusbox}{\statuslean{} --- \lean{KernelDichotomy} formalizes the equivalence of
(ii), (iii) and (iv), as
\lean{multiplicativelyDependentOutputs_iff_secondIterateKernelRelation} together with
\lean{exists_moebius_of_multiplicativelyDependentOutputs} and
\lean{multiplicativelyDependentOutputs_of_moebius}; the packaged branching is
\lean{twoBaseNonintegerSolution_kernel_dichotomy}.  Clause (i) is
\lean{SecondIterateKernel}.  The formalization is slightly stronger than the statement above:
(ii) $\iff$ (iii) needs only irrationality of $\thetaq$, and nonintegrality of $x$ enters
just for the nonvanishing of the M\"obius determinant.}
\end{statusbox}

\subsection{Why the multiplicative form is worth having}

Two consequences follow from stating the dichotomy multiplicatively.

\emph{It removes the transcendence vocabulary from the branching.}  Every counterexample
falls into exactly one of two cases, and the case distinction is a question about four
ordinary integers: either they satisfy a multiplicative relation or they do not.  In the
dependent branch Proposition~\ref{kn:prop-dim} says the relation is essentially unique up to
scaling, and its exponents on $M$ and $A$ have opposite signs unless one of the outputs is
$3$-smooth.

\emph{It makes Baker's theorem applicable to the independent branch.}  This is the point that
the Diophantine formulations obscure.  In the independent branch $M,A,2,3$ are
multiplicatively independent, so $\log M,\log A,\log2,\log3$ are $\Q$-linearly independent
\emph{logarithms of algebraic numbers} --- all four arguments are positive integers.  Baker's
homogeneous theorem then gives more than $\Q$-linear independence: these four logarithms are
linearly independent over $\Qbar$.  The corpus has not previously used this input.

\begin{corollary}[Algebraic-coefficient rigidity in the independent branch]
\label{kn:cor-baker}
Suppose $x$ is a nonintegral solution for which $M,A,2,3$ are multiplicatively independent.
Then no relation
\[
  \gamma_1\log M+\gamma_2\log A+\gamma_3\log2+\gamma_4\log3=0
\]
holds with algebraic $\gamma_1,\gamma_2,\gamma_3,\gamma_4$ not all zero.
\end{corollary}

\begin{statusbox}{\statusknown{} --- immediate from Baker's homogeneous theorem on
$\Qbar$-linear independence of logarithms of algebraic numbers, applied to the four positive
integers $M,A,2,3$.}
\end{statusbox}

It must be said plainly what this does and does not do.  It does \emph{not} close the
independent branch: the defining relation $\log M\log3=\log A\log2$ has coefficients $\log3$
and $-\log2$, which are transcendental by Lindemann, so it is not a $\Qbar$-linear relation
and Corollary~\ref{kn:cor-baker} does not apply to it.  That is exactly the ``Baker is linear,
not bilinear'' obstruction recorded in Appendix~\ref{app:falseshortcuts}.  What changes is
that the independent branch is no longer structureless: it carries a genuine unconditional
rigidity statement, and any future argument that manages to produce an algebraic-coefficient
relation among these four logarithms --- for instance by clearing the transcendental
coefficients against a second, independent relation --- would close it immediately.

\section{Depth-three tower rigidity}
\label{sec:tower}

A hypothetical counterexample makes the first level of the tower $2\mapsto2^x$ integral by
assumption.  This section answers the exact question that the corpus keeps circling: starting
from an arbitrary positive algebraic base, how many times can one apply the same nonintegral
exponent $x$ before transcendence is forced?  The answer is at most three, the bound is
attained, and the depth at which the first transcendental value appears is determined by a
single rational-linear invariant.

For positive real bases there is no branch ambiguity between iterated powers and powers of
powers, so the tower is unambiguous:
\[
 \bigl((\gamma^x)^x\bigr)^x=\gamma^{x^3}.
\]
Throughout this section $x$ denotes a hypothetical nonintegral solution, that is
$x\in\Sol\setminus\Nzero$, and we write
\[
 m=2^x\in\N,\qquad n=3^x\in\N .
\]
The letter $\gamma$ is reserved for a positive algebraic base, matching the Lean naming in
\lean{AlgebraicBaseUnitPower}.

\subsection{The saturated spectrum and its strong form}

Give $\Apos$ its structure of $\Q$-vector space, with rational scalars acting through the
unique positive real root, $q\cdot\gamma=\gamma^q$, and set
\[
 \GammaSat=\{2^r3^s:r,s\in\Q\}\subset\R_{>0}.
\]
This is the divisible hull of the group generated by $2$ and $3$: a countable, divisible
subgroup of rational rank two, dense in $\R_{>0}$ already through $\{2^r:r\in\Q\}$, and
isomorphic to $\Q^2$ as a $\Q$-vector space via $(r,s)\mapsto2^r3^s$.  An \emph{integer} lies
in $\GammaSat$ exactly when it is three-smooth, since clearing denominators in $N=2^r3^s$
gives $N^k=2^p3^q$ with $p,q\in\Z$ and unique factorization then confines the primes of $N$
to $\{2,3\}$.

\begin{theorem}[Exact algebraic-base spectrum]
\label{tw:spectrum}
Let $x\in\Sol\setminus\Nzero$ and let $\gamma\in\Apos$.  Then
\[
 \gamma^x\in\Abar
 \quad\Longleftrightarrow\quad
 \gamma\in\GammaSat .
\]
In particular the algebraic spectrum of a counterexample does not depend on which nonintegral
solution is used.
\end{theorem}

\begin{statusbox}{\statuslean{} --- in \lean{AlgebraicBaseUnitPower}, the pair
\begin{center}
{\scriptsize\lean{TwoBaseNonintegerSolution.real_rpow_isAlgebraic_iff_hasPositiveTwoThreeUnitPower}}
\end{center}
and
\begin{center}
{\scriptsize\lean{TwoBaseNonintegerSolution.transcendental_real_rpow_of_no_positiveTwoThreeUnitPower}}
\end{center}
}
The Lean predicate is \lean{HasPositiveTwoThreeUnitPower}: some positive natural power of
$\gamma$ is a positive rational $2,3$-unit.  For positive algebraic $\gamma$ this is exactly
membership in $\GammaSat$ --- one direction raises $2^r3^s$ to a common denominator, the other
takes the unique positive real root --- so the kernel-verified biconditional is the displayed
spectrum verbatim.
\end{statusbox}

The spectrum is already a strong statement: two algebraic outputs exhaust the maximal
algebraic multiplicative rank that the two-base configuration can support, and every positive
algebraic base outside one explicit dense subgroup produces a transcendental value.  For the
tower we need the sharper conclusion that the offending \emph{logarithm} is not merely
irrational but escapes the algebraic span of all logarithms.

Let $\LL$ denote the $\Abar$-vector space spanned by $1$ and by all logarithms of nonzero
algebraic numbers, arbitrary complex branches allowed; for a positive real algebraic number
the real logarithm is one of those branches.

\begin{definition}
\label{tw:def-logstrong}
A positive real value $e^z$ is \emph{log-strongly transcendental} when $z\notin\LL$.  This
implies ordinary transcendence, because the logarithm of a positive algebraic number lies in
$\LL$ by definition.
\end{definition}

\begin{theorem}[Strong algebraic-base spectrum]
\label{tw:strong-spectrum}
Let $x\in\Sol\setminus\Nzero$ and let $\gamma\in\Apos\setminus\GammaSat$.  Then
\[
 x\log\gamma\notin\LL,
\]
so $\gamma^x$ is log-strongly transcendental.  The same conclusion holds with $2,3$ replaced
by any multiplicatively independent pair $a,b\in\Apos$ with $a^x,b^x\in\Abar$, and
$\GammaSat$ replaced by $\GammaSat(a,b)=\{a^rb^s:r,s\in\Q\}$.
\end{theorem}

\begin{proof}
The logarithms $\log2,\log3,\log\gamma$ are linearly independent over $\Q$: a rational
relation clears to an integral multiplicative relation, which either contradicts the
multiplicative independence of $2$ and $3$ or places $\gamma$ in $\GammaSat$.  Baker's
homogeneous theorem \cite{Baker1975} upgrades this to linear independence over $\Abar$.

The exponent $x$ is transcendental: it is irrational by the kernel-verified rational-exponent
lemma of Section~\ref{sec:corpus}, and an algebraic irrational $x$ would make $2^x$
transcendental by Gelfond--Schneider \cite{Schneider1957}.  Hence $1$ and $x$ are linearly
independent over $\Abar$.

Apply Roy's strong six exponentials theorem \cite{Roy1992,Waldschmidt2005} to the three rows
$\log2,\log3,\log\gamma$ and the two columns $1,x$.  Five of the six pairwise products lie in
$\LL$, namely
\[
 \log2,\quad\log3,\quad\log\gamma,\quad x\log2=\log m,\quad x\log3=\log n .
\]
Therefore the sixth product $x\log\gamma$ lies outside $\LL$.
\end{proof}

\begin{statusbox}{\statusnew{} --- paper proof in this report.  It rests on Baker's theorem
and on Roy's strong form of the six exponentials theorem, neither of which is formalized in
the repository; the repository's own interpolation determinant is specialized to integer
entries.  Theorem~\ref{tw:spectrum}, the ordinary-transcendence shadow of this statement, is
\statuslean{}.}
\end{statusbox}

\subsection{The return kernel}

The spectrum controls one application of $x$.  The tower needs to know which bases
\emph{inside} $\GammaSat$ come back to $\GammaSat$ after that application.  On prime
valuations this is finite-dimensional linear algebra.

\begin{proposition}[The two outputs are not both three-smooth]
\label{tw:not-both-smooth}
For $x\in\Sol\setminus\Nzero$, at most one of $m=2^x$ and $n=3^x$ is three-smooth.
\end{proposition}

\begin{proof}
Suppose $m=2^a3^b$ and $n=2^c3^d$ with $a,b,c,d\in\Nzero$.  Taking logarithms in base $2$ and
in base $3$ gives $x=a+b\thetaq$ and $x=d+c/\thetaq$; eliminating $x$ and clearing the
denominator yields
\begin{equation}
 b\,\thetaq^2+(a-d)\,\thetaq-c=0 .
 \label{tw:eq-quadratic}
\end{equation}
The kernel-verified transcendence of $\thetaq$ forces $b=c=0$ and $a=d$, hence $x=a\in\Nzero$,
contrary to hypothesis.
\end{proof}

\begin{statusbox}{\statusnew{} --- the only nonelementary input is the kernel-verified
transcendence of $\thetaq$,
\begin{center}
\lean{transcendental_logThreeDivLogTwo}
\end{center}
in \lean{Transcendence}.  This is the same computation as the three-smooth output test of
Section~\ref{sec:smooth}, restated here in the form the kernel below consumes.}
\end{statusbox}

For $x\in\Sol\setminus\Nzero$ define the \emph{return kernel}
\begin{equation}
 K_x=\{\gamma\in\GammaSat:\gamma^x\in\GammaSat\}
 =\{2^r3^s:(r,s)\in\Q^2,\ m^rn^s\in\GammaSat\} .
 \label{tw:eq-kernel}
\end{equation}
Under the coordinate isomorphism $\GammaSat\cong\Q^2$ this is a $\Q$-subspace, because
$(2^r3^s)^x=m^rn^s$ is multiplicative in the coordinates.  Write
\[
 \mathbf m_{\mathrm{ext}}=\bigl(\vpp{p}{m}\bigr)_{p\ne2,3},
 \qquad
 \mathbf n_{\mathrm{ext}}=\bigl(\vpp{p}{n}\bigr)_{p\ne2,3},
\]
two finitely supported nonnegative integer vectors.

\begin{proposition}[Kernel formula and dimension]
\label{tw:kernel-dim}
In the coordinates $(r,s)$,
\begin{equation}
 K_x=\ker\Bigl(\Q^2\longrightarrow\Q^{(\mathbb P\setminus\{2,3\})},\quad
 (r,s)\longmapsto r\,\mathbf m_{\mathrm{ext}}+s\,\mathbf n_{\mathrm{ext}}\Bigr).
 \label{tw:eq-kernel-valuations}
\end{equation}
Moreover:
\begin{enumerate}[label=\textup{(\roman*)}]
\item $\dim_\Q K_x\le1$;
\item $K_x$ contains no $2^r3^s$ with $rs>0$;
\item $K_x=\{2^r:r\in\Q\}$ exactly when $m$ is three-smooth;
\item $K_x=\{3^s:s\in\Q\}$ exactly when $n$ is three-smooth.
\end{enumerate}
\end{proposition}

\begin{proof}
Membership of $m^rn^s$ in $\GammaSat$ is precisely the vanishing of every prime valuation
outside $\{2,3\}$, which is \eqref{tw:eq-kernel-valuations}.  By
Proposition~\ref{tw:not-both-smooth} at least one of the two external vectors is nonzero, so
the displayed linear map has rank at least one and its kernel has dimension at most one.

If $r,s>0$, then any external prime $p$ occurring in $m$ or in $n$ has
$r\vpp{p}{m}+s\vpp{p}{n}>0$, since both valuations are nonnegative and at least one is
positive; the case $r,s<0$ is symmetric.  Hence no same-sign direction survives.

Finally, if $m$ is three-smooth then $\mathbf m_{\mathrm{ext}}=0$ while
$\mathbf n_{\mathrm{ext}}\ne0$, so the kernel is exactly the first coordinate axis; the
converse follows by testing the direction $(1,0)$, which lies in $K_x$ exactly when
$m\in\GammaSat$, that is, exactly when the integer $m$ is three-smooth.  The statement for
$n$ is the mirror image.
\end{proof}

\begin{statusbox}{\statusnew{} --- elementary given Proposition~\ref{tw:not-both-smooth}.  No
transcendence input beyond the transcendence of $\thetaq$ enters the proof; the six
exponentials theorem is used only to interpret $K_x$, not to compute it.}
\end{statusbox}

The dynamical content of the kernel is that it admits no nontrivial internal orbit of length
two.

\begin{lemma}[No two-step internal orbit]
\label{tw:no-two-step}
Let $x\in\Sol\setminus\Nzero$ and let $\gamma\in K_x\setminus\{1\}$.  Then
$\gamma^x\notin K_x$; equivalently,
\[
 \gamma\in\GammaSat\ \text{and}\ \gamma^x\in\GammaSat
 \quad\Longrightarrow\quad
 \gamma^{x^2}\notin\GammaSat .
\]
\end{lemma}

\begin{proof}
Suppose $\gamma$ and $\gamma_1=\gamma^x$ both lie in $K_x\setminus\{1\}$.  By
Proposition~\ref{tw:kernel-dim}(i) the space $K_x$ is at most one-dimensional over $\Q$, so
two of its nonzero elements are rational scalar multiples of one another in the additive
notation: $\gamma_1=\gamma^q$ for some $q\in\Q$.  Since also $\gamma_1=\gamma^x$ and
$\gamma\ne1$, comparing logarithms gives $x=q$, contradicting the irrationality of $x$.
Hence $\gamma_1\notin K_x$, which by \eqref{tw:eq-kernel} says exactly that
$\gamma_1^x=\gamma^{x^2}\notin\GammaSat$.
\end{proof}

There are therefore no cycles, no fixed rays, and no length-two paths lying entirely inside
$\GammaSat$: a rational direction may survive one return to the algebraic spectrum, but the
returned direction cannot itself survive.

\subsection{Exact survival depth}

\begin{theorem}[Depth-three algebraic-tower theorem]
\label{tw:depth-three}
Let $x\in\Sol\setminus\Nzero$ and let $\gamma\in\Apos\setminus\{1\}$.  Then at least one of
\[
 \gamma^x,\qquad\gamma^{x^2},\qquad\gamma^{x^3}
\]
is transcendental.  More precisely, the first transcendental term occurs at exactly one of
the following three depths.
\begin{enumerate}[label=\textup{(\arabic*)}]
\item \textbf{Depth one.}  If $\gamma\notin\GammaSat$, then $\gamma^x$ is log-strongly
transcendental.
\item \textbf{Depth two.}  If $\gamma\in\GammaSat\setminus K_x$, then
$\gamma^x\in\Abar\setminus\GammaSat$ and $\gamma^{x^2}$ is log-strongly transcendental.
\item \textbf{Depth three.}  If $\gamma\in K_x\setminus\{1\}$, then $\gamma^x\in\GammaSat$,
$\gamma^{x^2}\in\Abar\setminus\GammaSat$, and $\gamma^{x^3}$ is log-strongly transcendental.
\end{enumerate}
Equivalently, for some $j\in\{1,2,3\}$ one has $x^j\log\gamma\notin\LL$.
\end{theorem}

\begin{proof}
Depth one is Theorem~\ref{tw:strong-spectrum}.

Assume $\gamma\in\GammaSat$, so that $\gamma_1=\gamma^x$ is algebraic by
Theorem~\ref{tw:spectrum}.  If $\gamma\notin K_x$, then $\gamma_1\notin\GammaSat$ by
\eqref{tw:eq-kernel}, and $\gamma_1$ is a positive algebraic base outside $\GammaSat$; the
strong spectrum applied to $\gamma_1$ and the same exponent $x$ makes
$\gamma_1^x=\gamma^{x^2}$ log-strongly transcendental.

If instead $\gamma\in K_x\setminus\{1\}$, then $\gamma_1\in\GammaSat$, so
$\gamma_2=\gamma_1^x=\gamma^{x^2}$ is algebraic.  By Lemma~\ref{tw:no-two-step} we have
$\gamma_1\notin K_x$, hence $\gamma_2\notin\GammaSat$.  One last application of the strong
spectrum, now to the base $\gamma_2$, gives $\gamma_2^x=\gamma^{x^3}$ log-strongly
transcendental.

The three cases are mutually exclusive and exhaustive, and in each case all earlier terms of
the tower are algebraic, so the stated depth is exact rather than an upper bound.
\end{proof}

The classification is compact enough to tabulate.  For a nonintegral solution $x$ and a
positive algebraic base $\gamma\ne1$, the first level of the tower whose logarithm escapes
$\LL$ is:

\begin{center}
\begin{tabular}{@{}p{0.28\textwidth}p{0.20\textwidth}p{0.40\textwidth}@{}}
\toprule
Location of $\gamma$ & First log outside $\LL$ & Structural reason\\
\midrule
$\gamma\notin\GammaSat$ & $x\log\gamma$ & exact algebraic spectrum\\
$\gamma\in\GammaSat\setminus K_x$ & $x^2\log\gamma$ & the first image leaves $\GammaSat$\\
$\gamma\in K_x\setminus\{1\}$ & $x^3\log\gamma$ & no two-step internal orbit\\
\bottomrule
\end{tabular}
\end{center}

\begin{statusbox}{\statusnew{} --- paper proof in this report, resting on
Theorem~\ref{tw:strong-spectrum} and hence on Baker's theorem and Roy's strong six
exponentials theorem.  The ordinary-transcendence version of depth one is \statuslean{}; the
ordinary-transcendence version of depth two for the mixed integer bases $2^i3^j$ with
$i,j>0$ is also \statuslean{}, as recorded below.}
\end{statusbox}

The theorem places a strict ceiling on how far an algebraic power tower can imitate the
integral first level of a counterexample.  A counterexample buys two algebraic outputs; it
never buys a third independent one, and the tower makes the deficit visible after at most
three steps regardless of the base chosen.

\subsection{Exact behaviour of the bases two, three, and six}

\begin{corollary}[Base-two tower dichotomy]
\label{tw:base2-tower}
Let $x\in\Sol\setminus\Nzero$ and $m=2^x$.
\begin{enumerate}[label=\textup{(\alph*)}]
\item If $m$ is not three-smooth, then $2^{x^2}$ is log-strongly transcendental.
\item If $m$ is three-smooth, say $m=2^a3^b$ with $a,b\in\Nzero$, then
\begin{equation}
 2^{x^2}=m^x=(2^x)^a(3^x)^b=m^an^b\in\N,
 \label{tw:eq-base2-integer}
\end{equation}
and $2^{x^3}$ is log-strongly transcendental.
\end{enumerate}
The mirror statement holds with $2$ replaced by $3$ and $m$ by $n$.
\end{corollary}

\begin{proof}
The base $2$ lies in $\GammaSat$, and its coordinate vector is $(1,0)$, which lies in $K_x$
exactly when $m\in\GammaSat$, that is, by Proposition~\ref{tw:kernel-dim}(iii), exactly when
the integer $m$ is three-smooth.  Now apply Theorem~\ref{tw:depth-three}: the non-smooth case
is depth two and the smooth case is depth three.  In the smooth case the intermediate value
is computed by \eqref{tw:eq-base2-integer}, where the middle equality is the identity
$(2^a3^b)^x=(2^x)^a(3^x)^b$ for positive real bases and the last membership holds because
$m,n\in\N$.
\end{proof}

Equation \eqref{tw:eq-base2-integer} is worth pausing on.  In the smooth branch the second
level of the tower is not merely algebraic but an explicit positive integer, expressed in the
two original outputs with the same exponents that describe the first output.  The tower does
not become transcendental because the arithmetic runs out immediately; it becomes
transcendental because the third level has nowhere left to land.

\begin{corollary}[Base six always fails at depth two]
\label{tw:base6}
For every $x\in\Sol\setminus\Nzero$ and all $r,s\in\Q$ of the same nonzero sign,
\[
 x^2\bigl(r\log2+s\log3\bigr)\notin\LL .
\]
In particular $x^2\log6\notin\LL$, so $6^{x^2}$ is log-strongly transcendental.
\end{corollary}

\begin{proof}
The base $\gamma=2^r3^s$ lies in $\GammaSat$, and its coordinate vector $(r,s)$ has positive
coordinate product, so $\gamma\notin K_x$ by Proposition~\ref{tw:kernel-dim}(ii).  Its exact
survival depth is therefore two, by Theorem~\ref{tw:depth-three}(2).  The case $r=s=1$ is
$\gamma=6$.
\end{proof}

This recovers and strengthens a kernel-verified family.  The Lean theorem
\begin{center}
\lean{TwoBaseNonintegerSolution.transcendental_mixed_two_three_squared_exponents}
\end{center}
in \lean{AlgebraicOutputLocus} proves that $\bigl(2^i3^j\bigr)^{x^2}$ is transcendental for
all natural $i,j>0$, which is exactly the statement that no genuinely mixed three-smooth
integer base survives to the second level.  Corollary~\ref{tw:base6} extends that to all
rational same-sign directions, including negative ones, and upgrades the conclusion from
transcendence to escape from $\LL$; more importantly it explains \emph{why} the mixed bases
fail, namely that the exceptional directions of $K_x$ are confined to the two coordinate
axes, which the mixed monomials avoid by construction.  The kernel-verified case is the
positive-integer corner of a statement that is really about a line in $\Q^2$.

\begin{statusbox}{\statusnew{} for Corollaries~\ref{tw:base2-tower} and~\ref{tw:base6} as
stated.  \statuslean{} for the ordinary transcendence of $\bigl(2^i3^j\bigr)^{x^2}$ with
integer $i,j>0$, in \lean{AlgebraicOutputLocus}.  The gap between the two is exactly the
general six exponentials input discussed in Section~\ref{sec:smooth}.}
\end{statusbox}

\subsection{One-base tower equivalences}

The depth theorem converts into a reformulation of the conjecture that is striking in its
indifference to the base.

\begin{resultbox}
\textbf{Universal three-level detector (Theorem~\ref{tw:universal-detector}).}
Fix \emph{any} positive algebraic $\gamma\ne1$ --- $\gamma=5$, $\gamma=\sqrt2$, $\gamma=2$,
$\gamma=3$, $\gamma=\tfrac12(7+3\sqrt5)$, anything at all.  Conjecture~\ref{conj:AE} is
equivalent to
\[
 \forall x\in\R,\quad 2^x,3^x\in\Z
 \ \Longrightarrow\
 \gamma^x,\ \gamma^{x^2},\ \gamma^{x^3}\in\Abar,
\]
and equally to the assertion that all three numbers $x\log\gamma$, $x^2\log\gamma$,
$x^3\log\gamma$ lie in $\LL$.  No arithmetic feature of the base is used: three levels of the
tower detect a counterexample for every positive algebraic base simultaneously. \statusnew
\end{resultbox}

\begin{theorem}[Universal three-level detector]
\label{tw:universal-detector}
Fix $\gamma\in\Apos\setminus\{1\}$.  Then Conjecture~\ref{conj:AE} is equivalent to
\begin{equation}
 \forall x\in\R,\quad
 2^x,3^x\in\Z
 \ \Longrightarrow\
 \gamma^x,\gamma^{x^2},\gamma^{x^3}\in\Abar,
 \label{tw:eq-three-level}
\end{equation}
and also to the same implication with $\Abar$-algebraicity replaced by the three membership
conditions $x^j\log\gamma\in\LL$, $j=1,2,3$.
\end{theorem}

\begin{proof}
If the conjecture holds, then $\Sol=\Nzero$, every exponent $x,x^2,x^3$ appearing in
\eqref{tw:eq-three-level} is a nonnegative integer, and all three powers of the algebraic
number $\gamma$ are algebraic, with logarithms in $\LL$.  If the conjecture fails, apply
Theorem~\ref{tw:depth-three} to a nonintegral solution: one of the three values is log-strongly
transcendental, so \eqref{tw:eq-three-level} fails and so does its $\LL$-version.
\end{proof}

Bases already inside $\GammaSat$ have their first level algebraic for free, which removes one
of the three conditions and yields the shortest form of the equivalence.

\begin{corollary}[Base-two cubic-tower equivalence]
\label{tw:base2-equivalence}
Conjecture~\ref{conj:AE} is equivalent to
\begin{equation}
 \forall x\in\R,\quad
 2^x,3^x\in\Z
 \ \Longrightarrow\
 2^{x^2}\in\Abar\ \text{ and }\ 2^{x^3}\in\Abar .
 \label{tw:eq-base2-cubic}
\end{equation}
The same statement holds with the base $3$ in place of the base $2$.
\end{corollary}

\begin{proof}
Take $\gamma=2$ in Theorem~\ref{tw:universal-detector} and drop the first condition, which is
automatic because $2^x$ is an integer under the hypothesis.
\end{proof}

This is a genuinely different equivalence from the second-level algebraicity tests of
Section~\ref{sec:smooth}, which require \emph{both} $2^{x^2}$ and $3^{x^2}$ to be algebraic.
One original base suffices if one is willing to climb a single extra level.  The trade is
informative: a second base and a third tower level carry exactly the same amount of
information, which is another way of seeing that the two-base configuration supplies only two
independent logarithmic directions no matter how the exponent is iterated.

\begin{statusbox}{\statusnew{} --- paper proof in this report.  Both equivalences inherit the
classical inputs of Theorem~\ref{tw:strong-spectrum}; the ordinary-algebraicity forms
\eqref{tw:eq-three-level} and \eqref{tw:eq-base2-cubic} need only the six exponentials
theorem, since Theorem~\ref{tw:spectrum} is \statuslean{} and only the depth-two and
depth-three steps require the general form.  Neither is a proof of
Conjecture~\ref{conj:AE}, which remains open: the tower reformulations relocate the
difficulty to the second and third levels without supplying a mechanism that produces
algebraicity there.}
\end{statusbox}

\section{Symmetric tower constants}
\label{sec:symmetric}

The exponent $\thetaq=\log_2 3$ and its reciprocal
\[
 \eta=\frac1{\thetaq}=\frac{\log2}{\log3}
\]
generate two symmetric ``next outputs'',
\[
 E_3=3^{\thetaq}=5.704522494691117635\ldots,
 \qquad
 E_2=2^{\eta}=1.548562652630242907\ldots
\]
The first is the canonical radical of Section~\ref{sec:radical} in the special branch where
the primitive odd core equals $3$; the second is its base-swapped mirror.  Neither constant
is known to be transcendental, and for $E_3$ even irrationality is open.  What is available
unconditionally is a statement about the \emph{pair}, and what a counterexample would add is
an exact classification of which of the two becomes algebraic.

\subsection{An unconditional strong dichotomy}

\begin{theorem}[At least one symmetric tower constant is strongly transcendental]
\label{tw:symmetric-unconditional}
At least one of the two numbers
\[
 \thetaq\log3=\log E_3,
 \qquad
 \eta\log2=\log E_2
\]
lies outside $\LL$.  In particular, at least one of $E_3$ and $E_2$ is transcendental.
\end{theorem}

\begin{proof}
The three numbers $1,\thetaq,\eta$ are linearly independent over $\Abar$: a relation
$a+b\thetaq+c\eta=0$ with algebraic coefficients, multiplied by $\thetaq$, becomes
$b\thetaq^2+a\thetaq+c=0$, which is impossible for a transcendental $\thetaq$ unless
$a=b=c=0$.  The two numbers $\log2,\log3$ are linearly independent over $\Q$ by unique
factorization, hence over $\Abar$ by Baker's theorem \cite{Baker1975}.

Apply Roy's strong six exponentials theorem \cite{Roy1992,Waldschmidt2005} to the three rows
$1,\thetaq,\eta$ and the two columns $\log2,\log3$.  Four of the six pairwise products are
already logarithms of algebraic numbers,
\[
 \log2,\qquad\log3,\qquad\thetaq\log2=\log3,\qquad\eta\log3=\log2,
\]
and therefore lie in $\LL$.  Hence at least one of the two remaining products,
$\thetaq\log3$ and $\eta\log2$, lies outside $\LL$.
\end{proof}

\begin{statusbox}{\statusnew{} --- paper proof in this report, using Baker's theorem and
Roy's strong six exponentials theorem.  The ordinary six exponentials theorem gives the
weaker dichotomy that one of $E_3,E_2$ is transcendental.  The strong form is worth the extra
input because it survives verbatim into the conditional branches below.}
\end{statusbox}

Note what the theorem does not say.  It does not identify which constant is transcendental,
and no known method identifies either one individually; that is precisely the wall described
in Section~\ref{sec:radical}.  The value of the dichotomy is that a counterexample would
resolve it, and would resolve it in a completely explicit way.

\subsection{Exact conditional classification}

Throughout this subsection $x\in\Sol\setminus\Nzero$, $m=2^x$, and $n=3^x$.

\begin{theorem}[The $E_3$ equivalences]
\label{tw:E3}
The following are equivalent.
\begin{enumerate}[label=\textup{(\roman*)}]
\item $E_3=3^{\thetaq}$ is algebraic;
\item $\thetaq\log3\in\LL$;
\item $m$ is three-smooth;
\item there are $a,b\in\Nzero$ with $b>0$ such that $m=2^a3^b$ and $x=a+b\thetaq$;
\item $x\in\Q+\Q\thetaq$.
\end{enumerate}
If these conditions fail, then $\thetaq\log3\notin\LL$ and $E_3$ is log-strongly
transcendental.  If they hold, then
\begin{equation}
 E_3^{\,b}=\frac{n}{3^a}\in\Q_{>0},
 \label{tw:eq-E3-radical}
\end{equation}
so $E_3$ is an explicit algebraic radical.
\end{theorem}

\begin{proof}
(iv)$\Rightarrow$(iii)$\Rightarrow$(v): the first is trivial, and taking logarithms of
$m=2^a3^b$ gives $x=a+b\thetaq$, which lies in $\Q+\Q\thetaq$.

(v)$\Rightarrow$(iv): write $x=u+v\thetaq$ with $u,v\in\Q$ and clear a common positive
denominator $d$, so that $m^d=2^p3^q$ with $p,q\in\Z$.  The left side is a positive integer,
so prime valuations force $p,q\ge0$ and $d\mid p$, $d\mid q$.  Hence $u=a$ and $v=b$ with
$a,b\in\Nzero$, and $b>0$ because $x\notin\Nzero$.

(i)$\Rightarrow$(iv): apply the six exponentials theorem to the rows $1,x,\thetaq$ and the
columns $\log2,\log3$.  The six exponential values are
\[
 2,\quad3,\quad m,\quad n,\quad3,\quad E_3,
\]
all algebraic by hypothesis.  The theorem therefore forbids both independence conditions
simultaneously; the columns $\log2,\log3$ are $\Q$-independent, so the rows are $\Q$-dependent
and $x=u+v\thetaq$ for some $u,v\in\Q$.  Now (v)$\Rightarrow$(iv) applies.

(iv)$\Rightarrow$(i): from $x=a+b\thetaq$ we get $n=3^x=3^aE_3^{\,b}$, which is
\eqref{tw:eq-E3-radical} and exhibits $E_3$ as a positive real radical of a positive rational,
hence algebraic.

Finally the strong statement.  Suppose $m$ is not three-smooth.  Then $\log2,\log m,\log3$ are
linearly independent over $\Q$: a rational relation with nonzero coefficient on $\log m$ would
place $m$ in $\GammaSat$, and an integer in $\GammaSat$ is three-smooth, while a relation with
zero coefficient on $\log m$ contradicts the independence of $\log2,\log3$.  Baker's theorem
upgrades this to independence over $\Abar$, so after division by $\log2$ the three numbers
$1,x,\thetaq$ are $\Abar$-linearly independent.  Apply Roy's theorem to the rows $1,x,\thetaq$
and the columns $\log2,\log3$: five of the six products, namely
$\log2,\log3,\log m,\log n$, and $\thetaq\log2=\log3$, lie in $\LL$, so
$\thetaq\log3\notin\LL$.  Thus failure of (iii) implies failure of (ii); and (i) implies (ii)
because $\log E_3=\thetaq\log3$ is then the logarithm of an algebraic number.
\end{proof}

\begin{theorem}[The $E_2$ equivalences]
\label{tw:E2}
The following are equivalent.
\begin{enumerate}[label=\textup{(\roman*)}]
\item $E_2=2^{\eta}$ is algebraic;
\item $\eta\log2\in\LL$;
\item $n$ is three-smooth;
\item there are $c,d\in\Nzero$ with $c>0$ such that $n=2^c3^d$ and $x=d+c\eta$;
\item $x\in\Q+\Q\eta$.
\end{enumerate}
If these conditions fail, then $\eta\log2\notin\LL$ and $E_2$ is log-strongly transcendental.
If they hold, then
\begin{equation}
 E_2^{\,c}=\frac{m}{2^d}\in\Q_{>0}.
 \label{tw:eq-E2-radical}
\end{equation}
\end{theorem}

\begin{proof}
Interchange the roles of $2$ and $3$, of $m$ and $n$, and of $\thetaq$ and $\eta$ in the proof
of Theorem~\ref{tw:E3}.  For the strong step, use the $\Q$-linear independence of
$\log3,\log n,\log2$ when $n$ is not three-smooth.
\end{proof}

\begin{corollary}[The two branches are mutually exclusive]
\label{tw:not-both-algebraic}
For a nonintegral solution, $E_3$ and $E_2$ cannot both be algebraic.  Equivalently, at least
one of $\thetaq\log3$ and $\eta\log2$ lies outside $\LL$ under a counterexample as well as
unconditionally.
\end{corollary}

\begin{proof}
By Theorems~\ref{tw:E3} and~\ref{tw:E2}, algebraicity of $E_3$ is equivalent to
three-smoothness of $m$ and algebraicity of $E_2$ to three-smoothness of $n$.
Proposition~\ref{tw:not-both-smooth} forbids both, since eliminating $x$ from
$x=a+b\thetaq$ and $x=d+c\eta$ produces the quadratic \eqref{tw:eq-quadratic} for the
transcendental $\thetaq$.
\end{proof}

So under a counterexample the unconditional dichotomy of
Theorem~\ref{tw:symmetric-unconditional} becomes a genuine classification rather than an
alternative: at most one of the two constants is algebraic, and each is algebraic exactly in
its own smooth-output branch, with an explicit radical description.

\begin{center}
\begin{tabular}{@{}p{0.24\textwidth}p{0.20\textwidth}p{0.44\textwidth}@{}}
\toprule
Smooth branch & Return kernel & Symmetric constants\\
\midrule
$m=2^a3^b$, $b>0$ & $\{2^r:r\in\Q\}$ &
$E_3^{\,b}=n/3^a\in\Q_{>0}$; $E_2$ log-strongly transcendental\\
$n=2^c3^d$, $c>0$ & $\{3^s:s\in\Q\}$ &
$E_2^{\,c}=m/2^d\in\Q_{>0}$; $E_3$ log-strongly transcendental\\
neither output smooth & no coordinate axis &
both $E_3$ and $E_2$ log-strongly transcendental\\
\bottomrule
\end{tabular}
\end{center}

The three rows of the table are exactly the three surviving rows of the primitive-core
trichotomy of Section~\ref{sec:smooth}: the first is the case $w=3$, in which the canonical
radical $E=w^{\thetaq}$ of Section~\ref{sec:radical} degenerates to $E_3$ itself with index
$d=b$; the second is its mirror $v=2$; and the third is the generic case with two independent
external prime supports.  The missing fourth pattern, both outputs three-smooth, is the one
excluded by Proposition~\ref{tw:not-both-smooth}.

Two remarks are worth making about the direction of the implications.  First, nothing here
proves that $E_3$ or $E_2$ is transcendental: a counterexample would make one of them an
explicit algebraic radical, and Corollary~\ref{tw:not-both-algebraic} only says that it cannot
make both algebraic.  The conjecture remains open, and the smallest instance --- irrationality
of $E_3=3^{\log_23}$ --- is exactly the wall named in Section~\ref{sec:radical}.  Second, the
implication that is usable in the other direction is the strong one: proving
$\thetaq\log3\notin\LL$ would kill the entire $w=3$ branch, and proving $\eta\log2\notin\LL$
would kill the $v=2$ branch, leaving only the generic case, in which both outputs carry prime
support outside $\{2,3\}$ and the local determinant methods have two independent external
primes to work with.

\begin{statusbox}{\statusnew{} --- paper proofs in this report.
Theorems~\ref{tw:E3} and~\ref{tw:E2} use the six exponentials theorem for the algebraicity
equivalences and Baker's theorem together with Roy's strong six exponentials theorem for the
$\LL$-statements; none of those classical inputs is formalized in the repository.
Corollary~\ref{tw:not-both-algebraic} adds only Proposition~\ref{tw:not-both-smooth}, whose
sole nonelementary ingredient is the kernel-verified transcendence of $\thetaq$, so the
mutual exclusion of the two smooth branches --- as opposed to the transcendence conclusions
attached to them --- is within immediate reach of the Lean development.}
\end{statusbox}

\section{Finite-difference obstructions}
\label{sec:differences}

The first local obstruction recorded in the corpus forbids three candidates in arithmetic
progression under the convenient hypothesis $h^{5}\le n$.  That hypothesis is an artifact of
the estimate, not of the mechanism: the sharp arbitrary-order curvature theorem is now
kernel-verified for every $r\ge2$, and it contains the three- and four-term cases as
specializations.

For a function $f$ and a step $h>0$ write $\Delta_hf(t)=f(t+h)-f(t)$ and
$\Delta_h^{r}=\Delta_h\circ\cdots\circ\Delta_h$ ($r$ factors).  Iterating the fundamental
theorem of calculus gives the mean-value identity in integral form,
\begin{equation}
\Delta_h^{r}f(n)=\int_{[0,h]^{r}}f^{(r)}(n+t_1+\cdots+t_r)\,dt_1\cdots dt_r .
\label{df:finite-diff-integral}
\end{equation}
For $f(t)=t^{\thetaq}$ and $r\ge2$ the derivative is
\[
 f^{(r)}(t)=\fall{\thetaq}{r}\,t^{\thetaq-r},
 \qquad
 \fall{\thetaq}{r}:=\prod_{j=0}^{r-1}(\thetaq-j)
 =\thetaq(\thetaq-1)\cdots(\thetaq-r+1).
\]
Since $1<\thetaq<2$ the falling factorial never vanishes, so $f^{(r)}$ has a fixed nonzero
sign on $(0,\infty)$, and $|f^{(r)}|$ decreases in $t$.  These two facts --- nonvanishing and
decay --- are the whole content of the obstruction.

\begin{theorem}[Higher-difference progression obstruction]
\label{df:thm-higher-difference}
Let $r\ge2$, let $n,h\in\N$ be positive, and set $c_r=|\fall{\thetaq}{r}|$.  If all of
$n,n+h,\dots,n+rh$ are candidates, then
\[
 1\le c_r\,h^{r}n^{\thetaq-r}.
\]
Such a progression is therefore impossible whenever
\begin{equation}
 c_r\,h^{r}<n^{\,r-\thetaq}.
\label{df:progression-condition}
\end{equation}
\end{theorem}

\begin{proof}
If all $n+jh$ are candidates, every $(n+jh)^{\thetaq}$ is an integer, so the integer
combination $\Delta_h^{r}f(n)=\sum_{j=0}^{r}(-1)^{r-j}\binom{r}{j}f(n+jh)$ is an integer.  By
\eqref{df:finite-diff-integral} and the fixed sign of $f^{(r)}$ it is nonzero, hence of
absolute value at least one.  Bounding $|f^{(r)}|$ by its value at the left endpoint $n$,
\[
 1\le|\Delta_h^{r}f(n)|\le c_r\,h^{r}n^{\thetaq-r}. \qedhere
\]
\end{proof}

\begin{statusbox}{\statuslean{} --- \lean{HigherDifference} proves the generic iterated
mean-value identity \lean{exists_iterated_fwdDiff_eq_pow_mul}, its real-power specialization
\lean{exists_iterated_fwdDiff_rpow_point}, integrality of the binomial forward difference, and
the obstruction \lean{not_arithmetic_progression_twoBaseNaturalCandidates_of_curvature}
uniformly for every $r\ge2$.  The displayed iterated-integral identity
\eqref{df:finite-diff-integral} is not used by the Lean proof; the formal argument iterates
the mean-value theorem instead, which avoids any integration theory.}
\end{statusbox}

\subsection{The sharp three-term criterion}

For $r=2$ the criterion \eqref{df:progression-condition} reads
\begin{equation}
 \thetaq(\thetaq-1)h^{2}<n^{\,2-\thetaq},
 \qquad\text{i.e.}\qquad
 h<1.038547802\ldots\cdot n^{0.2075187\ldots},
\label{df:sharp-3ap}
\end{equation}
strictly weaker as a hypothesis than $h\le n^{1/5}$, whose exponent is only $0.2$.  The real
form is \lean{second_difference_logThreeDivLogTwo_ap_of_curvature} in
\lean{SharperProgressions}, built on the curvature estimate
\lean{second_difference_logThreeDivLogTwo_ap_bound}.

Almost all of the gain survives in a purely rational hypothesis, which is what a kernel can
check without any real-number decision procedure.  The step exponent $83/400=0.2075$ is
admissible:

\begin{theorem}[Sharp rational three-term criterion]
\label{df:thm-sharp-rational}
Let $n,h\ge1$ be integers with $h^{400}\le n^{83}$.  Then $n$, $n+h$, $n+2h$ are not all
candidates.
\end{theorem}

\begin{statusbox}{\statuslean{} --- \lean{SharpProgression} proves
\lean{second_difference_logThreeDivLogTwo_ap_sharp} and
\lean{not_three_arithmetic_progression_twoBaseNaturalCandidates_sharp}.  The two numerical
ingredients are the enclosure $\thetaq<317/200$, replayed by the kernel from the exact power
inequality $3^{200}<2^{317}$ through \lean{logThreeDivLogTwo_lt_317_div_200}, and the
consequence $\thetaq(\thetaq-1)<37089/40000<1$.  For $r=3$, \lean{ThirdDifference}
kernel-checks the exact third-difference identity and the corresponding four-candidate
obstruction.}
\end{statusbox}

\subsection{Thresholds for higher order}

Writing \eqref{df:progression-condition} as a bound on the step, a progression of length
$r+1$ starting at $n$ is impossible once
\[
 h<C_r\,n^{\,1-\thetaq/r},
 \qquad C_r=c_r^{-1/r}.
\]
Table~\ref{df:tab-difference-thresholds} lists the first thresholds.

\begin{table}[ht]
\centering
\small
\begin{tabular}{@{}rrrr@{}}
\toprule
$r$ & length $r+1$ & exponent $1-\thetaq/r$ & $C_r$ \\
\midrule
2 & 3 & 0.207518750 & 1.038547802 \\
3 & 4 & 0.471679166 & 1.374849673 \\
4 & 5 & 0.603759375 & 1.164124485 \\
5 & 6 & 0.683007500 & 0.946705202 \\
6 & 7 & 0.735839583 & 0.778537835 \\
7 & 8 & 0.773576786 & 0.652646038 \\
8 & 9 & 0.801879687 & 0.557368999 \\
\bottomrule
\end{tabular}
\caption{Step thresholds for the higher-difference obstruction.  A progression of length
$r+1$ with step $h<C_rn^{1-\thetaq/r}$ cannot consist of candidates.}
\label{df:tab-difference-thresholds}
\end{table}

\subsection{Why longer progressions do not close the problem}

As $r$ grows the admissible step exponent $1-\thetaq/r$ approaches one, which looks
tantalizing: for large $r$ the theorem forbids progressions with steps almost as large as the
starting point itself.  It does not contradict the exact candidate monoid, and it cannot be
bootstrapped into a proof, for a purely combinatorial reason.

\begin{observation}[Combinatorial limitation]
\label{df:obs-combinatorial}
Conditionally on the structure described in Section~\ref{sec:conditional}, a dyadic block near
$X$ contains $O(\log X)$ candidates, so the average additive gap inside the block is of order
$X/\log X$.  For each fixed $r$ this exceeds $X^{\,1-\thetaq/r}$ once $X$ is large.  A
rank-two multiplicative monoid carries no reason whatsoever to contain long \emph{additive}
progressions; the hypothesis of Theorem~\ref{df:thm-higher-difference} is therefore vacuous
for the configurations that actually occur.
\end{observation}

Extracting leverage from long progressions would require $r$ growing with the height, control
of the rapidly growing constant $c_r=|\fall{\thetaq}{r}|$, and --- decisively --- a
combinatorial source guaranteeing such progressions exist.  None of the three is available.
The productive response is to drop the additive-structure hypothesis altogether, which is
what Section~\ref{sec:arbnode} does.

\subsection{The optimized variable-order scale}\label{od:optimized}

Theorem~\ref{df:thm-higher-difference} is uniform in $r$, so nothing forces the difference
order to stay fixed while the starting point grows.
Table~\ref{df:tab-difference-thresholds} records the first fixed-order thresholds, but each
row is only a snapshot of a one-parameter family that may be re-optimized at every height.
Carrying out that optimization strengthens every fixed-order corollary at once and, more
valuably, replaces the vague complaint that the candidate density ``feels too small'' by a
single explicit constant.  Throughout this subsection $\log$ denotes the natural logarithm.

\paragraph{The exact derivative coefficient.}
Fix $r\ge2$ and recall the constant $c_r=|\fall{\thetaq}{r}|$ of
Theorem~\ref{df:thm-higher-difference}.  Since $1<\thetaq<2$, the first two factors
$\thetaq$ and $\thetaq-1$ are positive while every later factor $\thetaq-j$ with $j\ge2$ is
negative, so taking absolute values termwise gives
\[
 c_r=\thetaq(\thetaq-1)\prod_{j=2}^{r-1}(j-\thetaq),
\]
and the recursion $\Gamma(z+1)=z\Gamma(z)$ telescopes the product into closed form:
\begin{equation}
 c_r=\frac{\thetaq(\thetaq-1)}{\Gamma(2-\thetaq)}\,\Gamma(r-\thetaq).
 \label{od:eq-cr-gamma}
\end{equation}
So $c_r$ grows at factorial speed, which is exactly why no fixed order can be optimal for all
$n$: the constant eventually defeats the improving exponent.  Stirling's formula applied to
\eqref{od:eq-cr-gamma} gives the root asymptotic that the optimization needs,
\begin{equation}
 c_r^{1/r}=\frac re\left(1+O\!\left(\frac{\log r}{r}\right)\right).
 \label{od:eq-cr-root}
\end{equation}

\paragraph{Optimizing the step threshold over the order.}
The formalized condition \eqref{df:progression-condition} forbids a progression of length
$r+1$ as soon as $h<c_r^{-1/r}n^{\,1-\thetaq/r}$, and substituting \eqref{od:eq-cr-root}
rewrites that threshold as
\begin{equation}
 c_r^{-1/r}n^{\,1-\thetaq/r}
 =\frac{en}{r}\exp\!\left(-\frac{\thetaq\log n}{r}\right)
  \left(1+O\!\left(\frac{\log r}{r}\right)\right).
 \label{od:eq-threshold-approx}
\end{equation}
The two $r$-dependent factors pull in opposite directions: $c_r^{-1/r}\sim e/r$ decays in
$r$, while $n^{-\thetaq/r}$ increases to $1$.  Discarding the lower-order factor and
differentiating the logarithm of the leading term,
\[
 \frac{d}{dr}\left(1+\log n-\log r-\frac{\thetaq\log n}{r}\right)
 =-\frac1r+\frac{\thetaq\log n}{r^{2}},
\]
which is positive for $r<\thetaq\log n$, vanishes at $r=\thetaq\log n$, and is negative
afterwards.  The leading term is therefore unimodal in real $r>0$ with a unique maximum at
$r=\thetaq\log n$, and at that point the explicit factor $e$ cancels against
$\exp(-\thetaq\log n/r)=e^{-1}$, leaving the clean scale $n/(\thetaq\log n)$.

\begin{theorem}[Optimized variable-order obstruction]
\label{od:thm-optimized}
Let $0<\varepsilon<1$.  For all sufficiently large $n$ choose an integer
\[
 r_n=\thetaq\log n+O(1),\qquad r_n\ge2 .
\]
If
\[
 1\le h\le(1-\varepsilon)\,\frac{n}{\thetaq\log n},
\]
then $n,n+h,\dots,n+r_nh$ are not all candidates.  Moreover, the largest leading-order step
scale obtainable by optimizing the curvature condition \eqref{df:progression-condition} over
$r$ is
\[
 h_{\max}(n)\sim\frac{n}{\thetaq\log n}
 =\frac{\log2}{\log3}\cdot\frac{n}{\log n}
 =0.6309297535\ldots\cdot\frac{n}{\log n}.
\]
\end{theorem}

\begin{proof}
Take $r_n$ to be either nearest integer to $\thetaq\log n$.  Then $r_n\to\infty$, so
\eqref{od:eq-cr-root} and hence \eqref{od:eq-threshold-approx} apply, and the quotient of
$c_{r_n}^{-1/r_n}n^{\,1-\thetaq/r_n}$ by $n/(\thetaq\log n)$ tends to $1$.  The fixed margin
$1-\varepsilon$ absorbs both the Stirling error $1+O(\log r/r)$ and the $O(1)$ rounding of
$\thetaq\log n$ to an integer, so \eqref{df:progression-condition} holds for every $h$ in the
stated range, and Theorem~\ref{df:thm-higher-difference} excludes the progression.

For the optimality assertion, the leading approximation in \eqref{od:eq-threshold-approx} is
unimodal in real $r>0$ with its unique maximum at $r=\thetaq\log n$, where its value is
$n/(\thetaq\log n)$.  Any regime in which $r$ stays bounded, or grows outside a constant
multiple of $\log n$, returns a strictly smaller order or a strictly smaller leading constant.
\end{proof}

\begin{statusbox}{\statusnew{} --- the obstruction being optimized is \statuslean{} and holds
uniformly for every $r\ge2$ with no upper bound on the order, so substituting $r=r_n$ inside
the formalized statement is legitimate; the relevant module is \lean{HigherDifference}.  The
optimization itself, the Gamma evaluation \eqref{od:eq-cr-gamma}, and the Stirling asymptotic
\eqref{od:eq-cr-root} are paper arguments and have not been submitted to the kernel.}
\end{statusbox}

\paragraph{Comparison with the conditional candidate spacing.}
A dyadic block $[X,2X)$ is an interval of length one in exponent space, so the exact block
count of Section~\ref{sec:conditional} applies verbatim: conditional on a counterexample with
least nonintegral generator $\beta$,
\begin{equation}
 \#(\Cand\cap[X,2X))
 =\frac{\log_2X}{\beta}+O(1)
 =\frac{\log X}{\beta\log2}+O_\beta(1).
 \label{od:eq-dyadic-count}
\end{equation}
The average additive spacing between consecutive candidates inside that block is therefore
\[
 \sim\beta\log2\cdot\frac{X}{\log X},
\]
whereas Theorem~\ref{od:thm-optimized} taken at $n\asymp X$ reaches only
\[
 h_{\max}\sim\frac1\thetaq\cdot\frac{X}{\log X}.
\]
Both are of order $X/\log X$, so the comparison is a pure constant, and in it the
transcendental $\thetaq$ cancels against the $\log2$:
\begin{equation}
 \frac{\beta\log2}{1/\thetaq}
 =\beta\,\thetaq\log2
 =\beta\log3 .
 \label{od:eq-spacing-ratio}
\end{equation}
The kernel-verified zero-free region $2^x<4096$ forces $\beta>12$ (Section~\ref{sec:conditional}),
which already makes the ratio exceed $12\log3=13.1833\ldots$; the interval scan of
Section~\ref{sec:finitecheck} raises it past $20\log3=21.9722\ldots$, and the extended
directed-rounding scan through $2^{27}$ raises it past $27\log3=29.6625\ldots$.  The last two
improvements are \statuscomp{} and do not enter any kernel-verified statement; the exponent
bound the kernel actually certifies remains $x\ge12$.

The same constant governs a second, purely cardinal comparison, which is the more telling of
the two.  The configuration forbidden by Theorem~\ref{od:thm-optimized} consists of
\[
 r_n+1\sim\thetaq\log X
\]
candidates arranged in one structured pattern, while the whole dyadic block contains only
$\sim\log X/(\beta\log2)$ candidates by \eqref{od:eq-dyadic-count}.  The ratio is again
exactly $\beta\log3$: the optimized theorem asks a block to supply about $\beta\log3$ times
as many candidates as it can possibly hold.  Concretely, at $X=2^{100}$ the block $[X,2X)$
holds $100/\beta+O(1)$ candidates --- about eight if $\beta$ were as small as its
kernel-certified floor, and about three under the $2^{27}$ scan --- while the shortest
progression the theorem forbids at that height already needs $r_n+1\approx111$ terms.

\begin{remark}[Quantified barrier]
\label{od:rem-barrier}
Optimizing over $r$ materially strengthens the local theorem: the admissible step scale rises
from $n^{\,1-\thetaq/r}$ at fixed $r$ to $n^{\,1-o(1)}$, missing $n$ itself only by the single
logarithmic factor $\thetaq\log n$.  It still cannot contradict the conditional monoid,
because the conditional candidate density in a dyadic block is too small by the fixed factor
$\beta\log3>13.18$.  No argument using only the number of candidates in a dyadic block can
force the progression that the current finite-difference functional requires.  A successful
continuation must exploit additional arithmetic correlation --- for example valuation
congruences that force a progression of a special shape --- or replace $\Delta_h^{r}$ by a
different integer-valued functional with a better curvature-to-spacing constant.
\end{remark}

This is a limitation of the method, not evidence against Conjecture~\ref{conj:AE}.  Its value
is to replace ``the density seems too weak'' by a sharp asymptotic constant and a concrete
numerical target that any improved functional must beat.

\paragraph{Agreement with the one-logarithm deficit.}
Section~\ref{sec:deficit} audits the semigroup determinant in the same currency and identifies
the same bottleneck.  In both analyses the binding quantity is the conditional dyadic
population $\asymp\log X$ of \eqref{od:eq-dyadic-count}, and in both the archimedean
inequality turns out to be comfortably satisfied by the exact monoid rather than violated by
it.  To that extent the two diagnoses agree.  They differ, instructively, on the size of the
shortfall.  The semigroup determinant needs $r\asymp(\log X)^2$ nodes and is supplied
$\asymp\log X$: a deficit of a full power of $\log X$.  The optimized finite difference needs
$\thetaq\log X$ nodes and is supplied $\log X/(\beta\log2)$: a deficit of the bounded factor
$\beta\log3$.  The finite difference is thus quantitatively much closer --- but only because
it buys that closeness with a hypothesis it cannot discharge, namely that the candidates lie
in exact arithmetic progression, which Observation~\ref{df:obs-combinatorial} says a rank-two
multiplicative monoid has no reason to provide.  Discarding that hypothesis is exactly what
the determinants of Sections~\ref{sec:arbnode} and~\ref{sec:semigroup} do, and the
one-logarithm deficit is the precise price charged for it.  Read together, the two sections
say that the present archimedean toolkit is short by one factor of $\log X$ in generality or
by a factor of $\beta\log3$ in strength, and that tuning constants converts neither shortfall
into the other.

\section{New result I: arbitrary-node lattice repulsion}
\label{sec:arbnode}

The higher-difference theorem of Section~\ref{sec:differences} uses equally spaced nodes,
because the forward-difference operator is built from an arithmetic progression.  The same
integer-versus-small-real principle extends to \emph{arbitrary} nodes through interpolation
determinants.  The resulting statement constrains every configuration of candidates, not only
the additively structured ones, which is exactly the gap identified in
Observation~\ref{df:obs-combinatorial}.

\subsection{The determinant identity}

For distinct real numbers $x_0<\cdots<x_r$ and a function $f$, write $f[x_0,\ldots,x_r]$ for
the divided difference.  The standard determinant identity is
\begin{equation}
 \det
 \begin{pmatrix}
  1&x_0&\cdots&x_0^{r-1}&f(x_0)\\
  1&x_1&\cdots&x_1^{r-1}&f(x_1)\\
  \vdots&\vdots&&\vdots&\vdots\\
  1&x_r&\cdots&x_r^{r-1}&f(x_r)
 \end{pmatrix}
 =V(x_0,\ldots,x_r)\,f[x_0,\ldots,x_r],
\label{df:det-divdiff}
\end{equation}
where
\[
 V(x_0,\ldots,x_r)=\prod_{0\le i<j\le r}(x_j-x_i).
\]
If $f\in C^{r}$ on the interval spanned by the nodes, the generalized mean-value theorem for
divided differences supplies some $\xi\in(x_0,x_r)$ with
\begin{equation}
 f[x_0,\ldots,x_r]=\frac{f^{(r)}(\xi)}{r!}.
\label{df:divdiff-mvt}
\end{equation}
For $f(t)=t^{\thetaq}$ this is again $f^{(r)}(t)=\fall{\thetaq}{r}t^{\thetaq-r}$, and since
$\thetaq$ is irrational --- in particular not an integer --- the coefficient
$\fall{\thetaq}{r}$ never vanishes.

\subsection{The arbitrary-node theorem}

\begin{theorem}[Vandermonde repulsion for candidates]
\label{df:thm-arbitrary-node}
Let $r\ge2$ and let
\[
 1\le m_0<m_1<\cdots<m_r
\]
be natural candidates, so that $m_i^{\thetaq}\in\N$ for every $i$.  Then
\begin{equation}
 \prod_{0\le i<j\le r}(m_j-m_i)
 \ge
 \frac{r!}{|\fall{\thetaq}{r}|}\,m_0^{\,r-\thetaq}.
\label{df:arbitrary-node-bound}
\end{equation}
\end{theorem}

\begin{proof}
Apply \eqref{df:det-divdiff} to $f(t)=t^{\thetaq}$ with the nodes $m_i$.  Every entry of the
matrix is an integer: the ordinary powers $m_i^{k}$ are integral by construction, and the last
column is integral precisely because the $m_i$ are candidates.  By \eqref{df:divdiff-mvt} the
determinant equals
\[
 D=V(m_0,\ldots,m_r)\,\frac{\fall{\thetaq}{r}}{r!}\,\xi^{\thetaq-r}
 \qquad(m_0<\xi<m_r).
\]
Each factor is nonzero, so $D\in\Z\setminus\{0\}$ and therefore $|D|\ge1$.  Because
$r\ge2>\thetaq$ the exponent $\thetaq-r$ is negative, so $\xi^{\thetaq-r}\le m_0^{\thetaq-r}$.
Therefore
\[
 1\le|D|
 \le V(m_0,\ldots,m_r)\,\frac{|\fall{\thetaq}{r}|}{r!}\,m_0^{\thetaq-r},
\]
which rearranges to \eqref{df:arbitrary-node-bound}.
\end{proof}

\begin{statusbox}{\statusnew{} for the unconditional paper theorem; \statuslean{} for the
formal determinant and repulsion chain.  \lean{ArbitraryNodeDeterminant} proves integrality,
the Vandermonde identities, and the hypothesis-carrying repulsion step, while
\lean{ArbitraryNodeRepulsion} packages the final bound under the explicit hypothesis
\lean{RpowDividedDifferenceMeanValue}.  The corpus does not assert that analytic hypothesis.}
\end{statusbox}

\begin{remark}[Formalization state]
\label{df:rem-lean-draft}
The accepted module \lean{ArbitraryNodeDeterminant} supplies the
algebraic half deliberately free of analysis --- the node matrix, an explicit integral model,
integrality of the determinant with the resulting bound $1\le|\det|$, the Vandermonde product
and its positivity on strictly increasing nodes.  The later endpoint theorem is also accepted,
but its divided-difference mean-value statement for $t\mapsto t^{\thetaq}$ remains a visible
hypothesis rather than a proved ingredient.  Thus the conditional Lean implication may be
cited as machine-checked; the unconditional analytic discharge is supplied only by the paper
proof here.
\end{remark}

\subsection{Three points: a lattice-triangle inequality}

For $r=2$ the determinant is
\[
 D=\det\begin{pmatrix}
 1&m_0&m_0^{\thetaq}\\
 1&m_1&m_1^{\thetaq}\\
 1&m_2&m_2^{\thetaq}
 \end{pmatrix},
\]
and $|D|$ is twice the area of the lattice triangle with vertices $(m_i,m_i^{\thetaq})$.
Strict convexity of $t\mapsto t^{\thetaq}$ makes the area positive; lattice integrality makes
it at least $1/2$.  The bound below is the quantitative form of that picture.

\begin{corollary}[Three-candidate span]
\label{df:cor-three-span}
If three candidates lie in $[N,N+H]$ with $N\ge1$, then
\begin{equation}
 H\ge\left(\frac{8}{\thetaq(\thetaq-1)}\right)^{1/3}N^{\,(2-\thetaq)/3}.
\label{df:three-span}
\end{equation}
Rounded to ten decimal places the constant is $2.0510723957$ and the exponent
$0.1383458331$; in particular $H\ge2.0510723956\,N^{0.1383458330}$.
\end{corollary}

\begin{proof}
Write $a=m_1-m_0$ and $b=m_2-m_1$.  Then $a,b>0$ and $a+b\le H$, so the Vandermonde product
satisfies
\[
 V=ab(a+b)\le\frac{H^{3}}{4},
\]
the continuous maximum being attained at $a=b=H/2$.  Theorem~\ref{df:thm-arbitrary-node} at
$r=2$ gives
\[
 V\ge\frac{2}{\thetaq(\thetaq-1)}N^{\,2-\thetaq}.
\]
Combining the two inequalities yields \eqref{df:three-span}.
\end{proof}

\subsection{General local packing bound}

Put
\[
 R_r:=\binom{r+1}{2}=\frac{r(r+1)}{2},
 \qquad
 C_r:=\left(\frac{r!}{|\fall{\thetaq}{r}|}\right)^{1/R_r},
 \qquad
 \alpha_r:=\frac{r-\thetaq}{R_r}=\frac{2(r-\thetaq)}{r(r+1)}.
\]
Every pairwise gap among points of $[N,N+H]$ is at most $H$, so the Vandermonde product is at
most $H^{R_r}$, and Theorem~\ref{df:thm-arbitrary-node} immediately gives the following.

\begin{corollary}[At most $r$ points below the repulsion scale]
\label{df:cor-r-packing}
If $r+1$ candidates lie in $[N,N+H]$, then
\begin{equation}
 H\ge C_rN^{\alpha_r}.
\label{df:r-span}
\end{equation}
Equivalently, every interval of length strictly less than $C_rN^{\alpha_r}$ beginning at or
after $N$ contains at most $r$ candidates.
\end{corollary}

The exponent can be optimized in $r$.  As a function of a real variable $r>\thetaq$,
\[
 \alpha(r)=\frac{2(r-\thetaq)}{r(r+1)}
\]
attains its maximum at
\[
 r_{\ast}=\thetaq+\sqrt{\thetaq^{2}+\thetaq}=3.6090841941\ldots,
\]
so the best integer order in this ordinary-power family is $r=4$ --- not the three-point case
that geometry makes most vivid.  Table~\ref{df:tab-ordinary-constants} lists the constants.

\begin{table}[ht]
\centering
\small
\begin{tabular}{rrrrr}
\toprule
$r$ & $R_r$ & $|\fall{\thetaq}{r}|$ & $\alpha_r$ & $C_r$ \\
\midrule
2 & 3  & 0.9271436280 & 0.1383458331 & 1.2920946431 \\
3 & 6  & 0.3847993728 & 0.2358395832 & 1.5805909191 \\
4 & 10 & 0.5445055422 & 0.2415037499 & 1.4602287461 \\
5 & 15 & 1.3150013031 & 0.2276691666 & 1.3510881062 \\
6 & 21 & 4.4907787617 & 0.2102398809 & 1.2735045942 \\
7 & 28 & 19.8269566337 & 0.1933941964 & 1.2187063909 \\
8 & 36 & 107.3637136680 & 0.1781954861 & 1.1790125192 \\
\bottomrule
\end{tabular}
\caption{Ordinary-power arbitrary-node repulsion constants.  For $r=2$ the sharper
three-point constant of Corollary~\ref{df:cor-three-span} uses the exact maximum of the
three-gap product rather than the crude bound $H^{R_2}$.}
\label{df:tab-ordinary-constants}
\end{table}

\begin{corollary}[Best ordinary-power local statement]
\label{df:cor-four-span}
Every interval
\[
 [N,N+H],
 \qquad H<1.4602287460\,N^{0.2415037499},
\]
contains at most four natural candidates.  The exact constant is
$C_4=(4!/|\fall{\thetaq}{4}|)^{1/10}=1.4602287461$ to ten decimal places, and the exact
exponent is $\alpha_4=(4-\thetaq)/10=0.2415037499\ldots$
\end{corollary}

\subsection{A bound for any longer interval}

The span estimate converts into a counting inequality.  Let $m_1<\cdots<m_s$ be all candidates
in $[N,N+H]$.  For each $i\le s-r$ the consecutive window $m_i,\ldots,m_{i+r}$ has span at
least $C_rN^{\alpha_r}$ by Corollary~\ref{df:cor-r-packing}.  Summing these spans, each
adjacent gap $m_{j+1}-m_j$ is counted at most $r$ times, so
\[
 (s-r)C_rN^{\alpha_r}\le r(m_s-m_1)\le rH.
\]

\begin{corollary}[Arbitrary-node counting inequality]
\label{df:cor-ordinary-count}
For every $r\ge2$,
\begin{equation}
 \#(\Cand\cap[N,N+H])\le r+\frac{r}{C_r}\,HN^{-\alpha_r}.
\label{df:ordinary-count}
\end{equation}
\end{corollary}

This is a genuine strengthening of progression avoidance: it constrains every possible
configuration of candidates in an interval, with no additive hypothesis at all.  It remains
far from the conditional benchmark of Section~\ref{sec:conditional}.  At the optimal $r=4$ and
$H=N$, the right-hand side of \eqref{df:ordinary-count} is $O(N^{0.7585\ldots})$, whereas the
exact population of a failed conjecture would be only $O(\log N)$.  The gap is enormous, and
closing it by ordinary powers alone is hopeless --- which motivates enlarging the column space,
the subject of Section~\ref{sec:semigroup}.

\subsection{Relation to higher differences}

For equally spaced nodes $m_i=n+ih$ the Vandermonde factor is explicit:
\[
 V=h^{R_r}\prod_{j=1}^{r}j! .
\]
Substituting in \eqref{df:det-divdiff} recovers an integer interpolation remainder closely
related to the $r$th forward difference of Section~\ref{sec:differences}.  The committed
repository theorem is sharper for this special configuration, because it uses the exact
binomial combination rather than bounding all pairwise gaps by a common span.  Conversely,
Theorem~\ref{df:thm-arbitrary-node} applies when the nodes carry no additive structure
whatsoever.  The two results are complementary, and only the arbitrary-node form escapes the
vacuity diagnosed in Observation~\ref{df:obs-combinatorial}.
\section{New result II: semigroup generalized Vandermonde determinants}
\label{sec:semigroup}

The determinant of Section~\ref{sec:arbnode} used only the integral columns
$1,m,m^2,\ldots$ and one extra integral column $m^{\thetaq}$.  But a candidate carries much
more integral data.  If $m^{\thetaq}=n\in\N$, then for every $a,b\in\Nzero$,
\begin{equation}
 m^{a+b\thetaq}=m^a n^b\in\N.
 \label{sg:eq-mixed-integrality}
\end{equation}
The full set of available exponents is therefore
\[
 \LambdaTheta:=\{a+b\thetaq:a,b\in\Nzero\}.
\]
Since $\thetaq$ is irrational, all pairs $(a,b)$ give distinct exponents.

\subsection{Positivity of generalized Vandermonde determinants}

Let
\[
 0\le\lambda_0<\lambda_1<\cdots<\lambda_r
\qquad\text{and}\quad
 0<x_0<x_1<\cdots<x_r.
\]
The functions $x^{\lambda_j}$ form an extended complete Chebyshev system on $(0,\infty)$.
Their Wronskian is classical in the Karlin--Studden theory of Tchebycheff systems:
\begin{equation}
 W(x)
 =x^{\sum_j\lambda_j-R_r}
   \prod_{0\le i<j\le r}(\lambda_j-\lambda_i)>0.
 \label{sg:eq-wronskian}
\end{equation}
Indeed, after factoring $x^{\lambda_j}$ from column $j$ and $x^{-i}$ from derivative row
$i$, the remaining determinant is the Vandermonde determinant of the $\lambda_j$, expressed
in the falling-factorial basis.  Positivity of all initial Wronskians implies strict positivity
of every ordered evaluation determinant.

For completeness, there is also an elementary zero-count proof.  Put $x=e^t$.  A nonzero
linear combination of $x^{\lambda_j}$ becomes an exponential polynomial
$\sum c_je^{\lambda_jt}$.  After factoring the lowest exponential, differentiation removes
the constant term and leaves one fewer exponential.  Induction with Rolle's theorem shows
that a combination of $r+1$ such functions has at most $r$ real zeros counting multiplicity.
Consequently an evaluation determinant at $r+1$ ordered nodes cannot vanish.  Its sign is
constant on the connected chamber $x_0<\cdots<x_r$ and is positive in the coalescing-node
limit by \eqref{sg:eq-wronskian}.

\begin{lemma}[Generalized Vandermonde positivity]
\label{sg:lem-generalized-vandermonde}
For strictly increasing nonnegative exponents $\lambda_j$ and strictly increasing positive
nodes $x_i$,
\begin{equation}
 \det\bigl(x_i^{\lambda_j}\bigr)_{0\le i,j\le r}>0.
 \label{sg:eq-gv-positive}
\end{equation}
\end{lemma}

The two arguments just sketched are carried out in full, with all induction bookkeeping made
explicit, in Appendix~\ref{app:gvproof}.

\subsection{Integer determinants from candidates}

\begin{theorem}[Semigroup integer determinant]
\label{sg:thm-semigroup-integer-det}
Let $m_0<\cdots<m_r$ be natural candidates.  Choose distinct pairs
$(a_j,b_j)\in\Nzero^2$ and order
\[
 \lambda_j=a_j+b_j\thetaq
\]
so that $\lambda_0<\cdots<\lambda_r$.  Then
\begin{equation}
 D_{\boldsymbol\lambda}(\boldsymbol m)
 :=\det\bigl(m_i^{\lambda_j}\bigr)_{0\le i,j\le r}
 \label{sg:eq-semigroup-det}
\end{equation}
is a positive integer.  In particular, $D_{\boldsymbol\lambda}(\boldsymbol m)\ge1$.
\end{theorem}

\begin{proof}
By \eqref{sg:eq-mixed-integrality}, each entry equals
\[
 m_i^{a_j}(m_i^{\thetaq})^{b_j}\in\N,
\]
so the determinant is an integer.  It is strictly positive by
Lemma~\ref{sg:lem-generalized-vandermonde}.
\end{proof}

\begin{statusbox}{\statusnew}
This theorem packages all mixed integrality consequences into one reusable object.  It is
not the same as the repository's exponential interpolation determinant in
\lean{MultiplicativeRank.lean}: here the rows are arbitrary natural candidates and the
columns are chosen from the two-dimensional exponent semigroup.
\end{statusbox}

\subsection{Divided-difference factorization}

Let $f_j(x)=x^{\lambda_j}$ and let
\[
 V(\boldsymbol m)=\prod_{i<j}(m_j-m_i).
\]
Applying the Newton divided-difference row transformation gives
\begin{equation}
 D_{\boldsymbol\lambda}(\boldsymbol m)
 =V(\boldsymbol m)
  \det\bigl(f_j[m_0,\ldots,m_i]\bigr)_{0\le i,j\le r}.
 \label{sg:eq-newton-factorization}
\end{equation}
The transformation is exact; its determinant is $V(\boldsymbol m)^{-1}$.

Suppose
\[
 N\le m_0<\cdots<m_r\le N+H\le2N.
\]
For each divided difference, the mean-value theorem gives a point in $[N,2N]$ and
\begin{equation}
 |f_j[m_0,\ldots,m_i]|
 \le
 \frac{|\fall{\lambda_j}{i}|}{i!}
 \sup_{N\le x\le2N}x^{\lambda_j-i}.
 \label{sg:eq-dd-first-bound}
\end{equation}
We need a uniform elementary coefficient estimate.

\begin{lemma}[Falling-factorial coefficient bound]
\label{sg:lem-falling-bound}
For $\lambda\ge0$ and $i\in\Nzero$,
\begin{equation}
 \frac{|\fall{\lambda}{i}|}{i!}\le2^{\lambda+i}.
 \label{sg:eq-falling-bound}
\end{equation}
\end{lemma}

\begin{proof}
If $0\le\lambda\le1$, each factor satisfies
$|\lambda-k|\le k+1$, so the quotient is at most $1$.
If $\lambda\ge1$, let $n=\lceil\lambda\rceil$.  Then
\[
 \frac{|\fall{\lambda}{i}|}{i!}
 \le\prod_{k=0}^{i-1}\frac{\lambda+k}{k+1}
 \le\prod_{k=0}^{i-1}\frac{n+k}{k+1}
 =\binom{n+i-1}{i}.
\]
The last binomial coefficient is at most the sum of its row,
$2^{n+i-1}$, and $n-1\le\lambda$.  Hence it is at most $2^{\lambda+i}$.
\end{proof}

For $x\in[N,2N]$,
\begin{equation}
 x^{\lambda-i}\le2^\lambda N^{\lambda-i}.
 \label{sg:eq-power-strip-bound}
\end{equation}
If $\lambda-i\ge0$, use $x\le2N$; if it is negative, use $x\ge N$.
Combining \eqref{sg:eq-dd-first-bound}, \eqref{sg:eq-falling-bound}, and
\eqref{sg:eq-power-strip-bound},
\begin{equation}
 |f_j[m_0,\ldots,m_i]|
 \le2^{2\lambda_j+i}N^{\lambda_j-i}.
 \label{sg:eq-dd-final-bound}
\end{equation}

\subsection{Quantitative semigroup gap theorem}

Set
\[
 R:=R_r=\frac{r(r+1)}2,
 \qquad
 \Sigma:=\sum_{j=0}^{r}\lambda_j.
\]
By the Leibniz formula and \eqref{sg:eq-dd-final-bound}, every permutation term in the small
determinant in \eqref{sg:eq-newton-factorization} is at most
\[
 2^{2\Sigma+R}N^{\Sigma-R},
\]
and there are $(r+1)!$ terms.  Since every pairwise node difference is at most $H$,
$V(\boldsymbol m)\le H^R$.  The positive integer lower bound from
Theorem~\ref{sg:thm-semigroup-integer-det} now gives the following.

\begin{theorem}[Semigroup determinant gap]
\label{sg:thm-semigroup-gap}
Let $r\ge1$.  Let $\lambda_0<\cdots<\lambda_r$ be any $r+1$ exponents in
$\LambdaTheta$, and put $R=r(r+1)/2$ and $\Sigma=\sum_j\lambda_j$.  If $r+1$ natural
candidates lie in
\[
 [N,N+H]\subseteq[N,2N],
 \qquad N\ge1,
\]
then
\begin{equation}
 H\ge
 ((r+1)!)^{-1/R}
 2^{-1-2\Sigma/R}
 N^{1-\Sigma/R}.
 \label{sg:eq-semigroup-gap}
\end{equation}
\end{theorem}

\begin{proof}
Equations \eqref{sg:eq-newton-factorization} and \eqref{sg:eq-dd-final-bound} imply
\[
 1\le D_{\boldsymbol\lambda}(\boldsymbol m)
 \le
 H^R(r+1)!2^{R+2\Sigma}N^{\Sigma-R}.
\]
Taking $R$th roots and rearranging proves \eqref{sg:eq-semigroup-gap}.
\end{proof}

The exponent $1-\Sigma/R$ is the key.  We should choose the $r+1$ smallest available
semigroup exponents to minimize $\Sigma$.

\subsection{The ordered exponent semigroup}

Write the elements of $\LambdaTheta$ in increasing order:
\[
 0=\lambda_0<\lambda_1<\lambda_2<\cdots.
\]
The beginning is
\begin{align*}
 0,\ 1,\ \thetaq,\ 2,\ 1+\thetaq,\ 3,\ 2\thetaq,\ 2+\thetaq,\ 4,
 \ 1+2\thetaq,\ 3+\thetaq,\ 3\thetaq,\ 5,\ldots
\end{align*}
The counting function is elementary:
\begin{align}
 L(T)
 &:=\#\{\lambda\in\LambdaTheta:\lambda\le T\}\notag\\
 &=\sum_{0\le b\le T/\thetaq}
   \left(\left\lfloor T-b\thetaq\right\rfloor+1\right)
 =\frac{T^2}{2\thetaq}+O(T).
 \label{sg:eq-semigroup-counting}
\end{align}
Irrationality of $\thetaq$ ensures that different pairs $(a,b)$ are counted distinctly.
Inverting \eqref{sg:eq-semigroup-counting} and summing the order statistics gives
\begin{align}
 \lambda_r
 &=\sqrt{2\thetaq r}+O(1),
 \label{sg:eq-lambda-asymp}\\
 \Sigma_r:=\sum_{j=0}^{r}\lambda_j
 &=\frac23\sqrt{2\thetaq}\,r^{3/2}+O(r),
 \label{sg:eq-sigma-asymp}\\
 \frac{\Sigma_r}{R_r}
 &=\frac43\sqrt{\frac{2\thetaq}{r}}+O(r^{-1}).
 \label{sg:eq-sigma-ratio-asymp}
\end{align}
Consequently the power of $N$ in \eqref{sg:eq-semigroup-gap} tends to $1$:
\begin{equation}
 1-\frac{\Sigma_r}{R_r}
 =1-\frac43\sqrt{\frac{2\thetaq}{r}}+O(r^{-1}).
 \label{sg:eq-delta-asymp}
\end{equation}

\begin{table}[ht]
\centering
\small
\begin{tabular}{rcc}
\toprule
$r$ & $\Sigma_r$ & $1-\Sigma_r/R_r$ \\
\midrule
2  & $2.5849625007$  & $0.1383458331$ \\
3  & $4.5849625007$  & $0.2358395832$ \\
4  & $7.1699250014$  & $0.2830074999$ \\
5  & $10.1699250014$ & $0.3220049999$ \\
6  & $13.3398500029$ & $0.3647690475$ \\
7  & $16.9248125036$ & $0.3955424106$ \\
8  & $20.9248125036$ & $0.4187552082$ \\
9  & $25.0947375050$ & $0.4423391666$ \\
10 & $29.6797000058$ & $0.4603690908$ \\
12 & $39.4345875079$ & $0.4944283653$ \\
20 & $87.7939875195$ & $0.5819333928$ \\
\bottomrule
\end{tabular}
\caption{Exponent gain from using the first $r+1$ mixed exponents in
$\Nzero+\Nzero\thetaq$.  The ordinary-power family peaks and then decays; the semigroup
family improves monotonically in these data and tends to exponent $1$.}
\label{sg:tab-semigroup-delta}
\end{table}

\subsection{An explicit coarse form}

For the global counting corollary we need no delicate lattice-point error term.  A simple
bound suffices.

\begin{lemma}[Crude semigroup weight bound]
\label{sg:lem-crude-semigroup}
For $r\ge2$, the first $r+1$ semigroup exponents satisfy
\begin{equation}
 \lambda_r\le6\sqrt r,
 \qquad
 \frac{\Sigma_r}{R_r}\le\frac{12}{\sqrt r}.
 \label{sg:eq-crude-semigroup}
\end{equation}
\end{lemma}

\begin{proof}
Let $q=\lceil\sqrt{r+1}\rceil$.  The $(q+1)^2$ pairs
$0\le a,b\le q$ give distinct exponents and each is at most
$q(1+\thetaq)<3q$.  Since $(q+1)^2>r+1$, the $r$th ordered exponent is below $3q$.
For $r\ge2$,
$q\le\sqrt{r+1}+1\le2\sqrt r$, hence $\lambda_r\le6\sqrt r$.
Then
\[
 \Sigma_r\le(r+1)\lambda_r\le6(r+1)\sqrt r,
\]
and division by $R_r=r(r+1)/2$ gives the second inequality.
\end{proof}

For $r\ge144$, \eqref{sg:eq-crude-semigroup} gives $\Sigma_r/R_r\le1$.  Also
\[
 ((r+1)!)^{1/R_r}
 \le(r+1)^{(r+1)/R_r}
 =(r+1)^{2/r}\le3
\]
for $r\ge2$.  Therefore Theorem~\ref{sg:thm-semigroup-gap} yields a particularly clean form.

\begin{corollary}[Explicit near-linear span]
\label{sg:cor-explicit-semigroup-span}
If $r\ge144$ and $r+1$ natural candidates lie in
$[N,N+H]\subseteq[N,2N]$, then
\begin{equation}
 H\ge\frac1{24}N^{1-12/\sqrt r}.
 \label{sg:eq-explicit-semigroup-span}
\end{equation}
\end{corollary}

\begin{proof}
In \eqref{sg:eq-semigroup-gap}, use
$((r+1)!)^{-1/R}\ge1/3$,
$2^{-1-2\Sigma/R}\ge2^{-3}=1/8$, and, since $N\ge1$,
$N^{1-\Sigma/R}\ge N^{1-12/\sqrt r}$.
\end{proof}

\subsection{Counting candidates in a general interval}

As in the ordinary-power count of Section~\ref{sec:arbnode}, sum the spans of consecutive
windows of $r+1$ candidates.  Every adjacent gap is counted at most $r$ times.

\begin{corollary}[Semigroup interval count]
\label{sg:cor-semigroup-interval-count}
For $r\ge144$ and $1\le H\le N$,
\begin{equation}
 \#(\Cand\cap[N,N+H])
 \le r+24rH N^{-1+12/\sqrt r}.
 \label{sg:eq-semigroup-interval-count}
\end{equation}
\end{corollary}

\begin{proof}
If $s$ is the number of candidates and $s\le r$, there is nothing to prove.  Otherwise,
Corollary~\ref{sg:cor-explicit-semigroup-span} gives a lower bound
$L=N^{1-12/\sqrt r}/24$ for each span $m_{i+r}-m_i$.  Thus
\[
 (s-r)L\le rH,
\]
which is equivalent to \eqref{sg:eq-semigroup-interval-count}.
\end{proof}

\subsection{An independent dyadic logarithmic-square bound}

\begin{corollary}[Explicit $O((\log X)^2)$ dyadic bound]
\label{sg:cor-semigroup-log-square}
For every real $X\ge1$,
\begin{equation}
 \#(\Cand\cap[X,2X])
 \le10000\max\{1,\log X\}^2.
 \label{sg:eq-semigroup-log-square}
\end{equation}
\end{corollary}

\begin{proof}
Put $L=\max\{1,\log X\}$ and
\[
 r=\left\lceil144L^2\right\rceil.
\]
Then $r\ge144$, $\sqrt r\ge12L$, and
\[
 X^{12/\sqrt r}
 =\exp\!\left(\frac{12\log X}{\sqrt r}\right)
 \le e.
\]
Apply \eqref{sg:eq-semigroup-interval-count} with $N=H=X$:
\[
 \#(\Cand\cap[X,2X])
 \le r+24rX^{12/\sqrt r}
 \le(1+24e)r.
\]
Since $L\ge1$, $r\le145L^2$, and
$(1+24e)145<10000$.
\end{proof}

\begin{statusbox}{\statusnew}
The bound is numerically coarse by design.  Its value is conceptual independence: it comes
from positivity and integrality of generalized Vandermonde determinants, not from the
prime-factor valuation rank of candidates.  Refining constants is much less important than
finding an additional source of determinant lower bound.
\end{statusbox}

\begin{statusbox}{\statuslean{} and \statusnew{}}
The hierarchy is now kernel-accepted in \lean{MixedExponentSemigroup},
\lean{SemigroupDeterminant}, \lean{FallingFactorialBound}, and
\lean{SemigroupWeightBound}.  The final span and $10^4(\log X)^2$ deductions are accepted in
\lean{SemigroupGapBound} conditional on its explicit \lean{NewtonFactorization} argument.
That predicate is a visible input, not an axiom; the unconditional analytic proof remains the
paper proof in this section.
\end{statusbox}

\section{The min-plus ceiling: why term-wise \texorpdfstring{$p$}{p}-adic amplification cannot
close a determinant argument}
\label{sec:minplus}

Every interpolation-determinant argument in this report squeezes a nonzero integer determinant
$D$ between a lower and an upper bound.  A recurring proposal --- registered in the deficit
analysis of Section~\ref{sec:deficit} as the most natural way to supply the missing logarithm
--- is to manufacture the lower bound from $p$-adic content: choose the nodes and the column
exponents so that every monomial in the Leibniz expansion of $D$ is divisible by a high power
of $2$ and of $3$, and conclude
\begin{equation}
  v_p(D)\ \ge\ V_p^{*}:=\min_{\sigma}\sum_i v_p\!\left(E_{i\sigma(i)}\right),
  \qquad |D|\ \ge\ 2^{V_2^{*}}3^{V_3^{*}} .
  \label{eq:mp-minplus}
\end{equation}
This section records why the proposal cannot work, and --- as importantly --- corrects the
quantitative form in which the obstruction has been stated.

\subsection{The threshold}

Fix a column set $S=\{(a_j,b_j)\}_{j=0}^{r}$ of distinct mixed exponents, put
$\lambda_j=a_j+b_j\thetaq$ and $W=\sum_j\lambda_j$, and let $D$ be the semigroup determinant
$\det\bigl(c_i^{a_j}(c_i^{\thetaq})^{b_j}\bigr)$ built on the candidates $c_0<\cdots<c_r$ of a
single dyadic block $[2^T,2^{T+1})$.  Every entry satisfies
$\log_2E_{ij}\le(T+1)\lambda_j$, so the Leibniz bound gives
\begin{equation}
  \log_2|D|\ \le\ (T+1)W+\log_2(r+1)!\ =\ TW\,(1+o(1)).
  \label{eq:mp-bar}
\end{equation}
An amplification scheme yielding a guaranteed lower bound $|D|\ge2^{V}$ therefore needs
$V>(T+1)W+\log_2(r+1)!$ for a contradiction.  We call the right-hand side of
\eqref{eq:mp-bar} the \emph{bar}.

\subsection{The ceiling}

The obstruction is not quantitative at all.  It is the observation that the $\{2,3\}$-part of
an integer divides the integer.

\begin{theorem}[Min-plus ceiling]
\label{mp:thm-ceiling}
Let $n\ge1$ be a natural number, let $V_2,V_3\ge0$ satisfy $2^{V_2}\mid n$ and $3^{V_3}\mid n$,
and suppose $n\le2^{B}$.  Then
\[
  V_2+V_3\,\thetaq\ \le\ B .
\]
More generally, for any finite set $S$ of primes,
$\sum_{p\in S}\vpp pn\log p\le\log n$.
\end{theorem}

\begin{proof}
$2^{V_2}$ and $3^{V_3}$ are coprime and both divide $n$, so $2^{V_2}3^{V_3}\mid n$ and hence
$2^{V_2}3^{V_3}\le n\le2^{B}$.  Taking $\log_2$ gives the claim.  For general $S$ the same
argument applies to $\prod_{p\in S}p^{\vpp pn}$, which divides $n$ by pairwise coprimality.
\end{proof}

\begin{statusbox}{\statuslean{} --- \lean{minPlusContent_le_of_le} and
\lean{minPlusContent_le_of_natAbs_le} in \lean{MinPlusCeiling}, with the general prime-set form
\lean{primeContent_le_log} and the divisibility input \lean{prod_primePow_dvd}.}
\end{statusbox}

\begin{corollary}[Term-wise amplification is inert]
\label{mp:cor-inert}
Let $D\ne0$ be the determinant of an interpolation system, let $V_2^{*},V_3^{*}$ be any
certified min-plus lower bounds as in \eqref{eq:mp-minplus}, and let $B$ be any correct
archimedean upper bound $\log_2|D|\le B$.  Then $V_2^{*}+V_3^{*}\thetaq\le B$.

Consequently no choice of node system, no choice of column set, and no enlargement of the prime
set can make a term-wise divisibility bound contradict a correct archimedean bound.
\end{corollary}

The corollary is unconditional and needs no equidistribution input, no averaging, and no
property of the candidates beyond $D\in\Z\setminus\{0\}$.  Its content is that term-wise
divisibility contributes nothing beyond what $|D|\ge1$ already gives: it is a lower bound for
$|D|$ derived from $|D|$ itself.

\subsection{A correction to the averaging form of the obstruction}

The obstruction has previously been argued by averaging: the minimum over permutations is at
most the mean over permutations, the mean of $\sum_iv_p(E_{i\sigma(i)})$ is
$\sum_j\bigl(a_j\bar\nu+b_j\bar\mu\bigr)$ with $\bar\nu,\bar\mu$ the row-averaged valuations,
and the depths $n_k$ equidistribute over $[0,T]$, so $\bar\nu=T/2$.  That computation is
correct, and its conclusion --- that the ceiling is governed by the \emph{same} block
equidistribution that defeats the shrinking-target route of Section~\ref{sec:reform} --- is a
genuine structural insight worth keeping: at the structural places the two walls of this
problem are one wall.

The numerical conclusion drawn from it, however, is too strong.  Bounding the $2$-adic place
and the $\thetaq$-weighted $3$-adic place separately gives $TW/2$ \emph{each}; their sum is
$TW$, which is the bar, not half of it.  With $M=2^{\beta}$ and $A=3^{\beta}$ and the
normalization $s_M=2^{\vpp2M}3^{\vpp3M}$, $s_A=2^{\vpp2A}3^{\vpp3A}$ for the $3$-smooth parts,
the mean bound reads
\begin{equation}
  V^{*}\ \le\ \frac{T}{2}\left[\,W+\frac1\beta\sum_j\Bigl(a_j\log_2s_M+b_j\log_2s_A\Bigr)\right],
  \label{eq:mp-mean}
\end{equation}
and the bracket is at most $2W$, with equality exactly when $M$ and $A$ are both $3$-smooth.
Theorem~\ref{lp:thm-smooth} excludes that case, so the mean bound is strictly below the bar;
but the deficiency it provides is proportional to the roughness of the outputs and degenerates
as the outputs approach smoothness.

The factor $1/2$ requires $3\nmid M$ \emph{and} $2\nmid A$.  The matched-depth theorem of
Section~\ref{sec:conditional} excludes only the pair $2\mid M$ and $3\mid A$; the crossed pair
$3\mid M$ and $2\mid A$ is not excluded, and it is exactly the configuration in which the
ceiling climbs.

\subsection{Calibration: the exact ceilings}

The mean bound is not the true ceiling, because the minimum is genuinely lower than the mean.
The exact constants were obtained in the continuation report by an optimal-transport
computation, and they replace every heuristic figure that has circulated for this obstruction.

Write $\mathfrak G$ for the Gini norm of a normalized prime-valuation slope vector, and $c_p$
for the slope vector at $p$.  The termwise all-prime divisor fraction of the balanced
mixed-semigroup determinant is exactly $1-\tfrac38\sum_p\mathfrak G(c_p)+o(1)$, and the
primitive condition $\min\bigl(\vpp2M,\vpp3A\bigr)=0$ forces $\sum_p\mathfrak G(c_p)\ge8/15$.

\begin{center}
\begin{tabular}{lc}
\toprule
System & Termwise-divisor ceiling \\
\midrule
Ordinary input/output Vandermonde product & $<0.8710491$ \\
Mixed single block, arbitrary branch & $4/5$ \\
Mixed single block, cross-coprime & $7/10$ \\
Mixed single block, places $2$ and $3$ only, cross-coprime & $3/10$ \\
Complete triangular row grid & $\pi/4=0.7853982\ldots$ \\
First tied fibers at $2$ and $3$ & $o(1)$ of the block budget \\
\bottomrule
\end{tabular}
\end{center}

Two figures that have been proposed for this obstruction are not valid in the stated
generality.  A structural ceiling of $1/2$ requires dropping the depths $k\,u\,\vpp2M$ and
$k\,v\,\vpp3A$ simultaneously, which assumes $\vpp2M=\vpp3A=0$; the primitive theorem of
Section~\ref{sec:conditional} supplies only $\min\bigl(\vpp2M,\vpp3A\bigr)=0$.  An all-prime
ceiling of $2/3$ does not follow from ``minimum $\le$ average'': for an external positive weight
proportional to $X_1+X_2$ on the standard triangle one has $\mathbb E Y=2/3$ and
$\Phi(1,1)=4/15$, so the minimum assignment is $4/5$ of half the mean, not $2/3$.

An independent numerical check, computing exact assignment minima for the block node system
over the places $2$ and $3$ only, in the cross-coprime branch and with a square-box column
family, gives a ratio settling near $0.33$ across block heights $T=400,1200,3600,10000$,
consistent with the exact continuum value $3/10$ for those places.  The same computation shows
what happens outside the balanced primitive normalization: with a degenerate column family such
as the pure-$b$ columns $(0,j)$, and with $A$ within $\log_23$ of being a power of three, the
ratio rises to $0.9964$.  That configuration violates none of the kernel-verified constraints,
but it saturates \emph{tautologically} --- the entries are then nearly powers of three, their
$3$-adic content is nearly their size, and \eqref{eq:mp-minplus} degenerates.  The same
degeneration occurs on the $a$-side at $b=0$, where $D=\prod_{i<i'}(c_{i'}-c_i)$ and
$v_2(D)=\sum_{i<i'}\min(n_i,n_{i'})$ exactly.  This is why the balanced and primitive
normalizations are not cosmetic: they are what make the $4/5$ and $7/10$ ceilings meaningful.

\begin{statusbox}{\statusnew{} for the exact ceilings, from the continuation report's transport
computation --- they follow from the closed-form Gini identity, not from experiment.
\statuscomp{} for the independent $3/10$ cross-check and the degenerate-family observation:
exact rational node data, exact Jonker--Volgenant assignment minima, double-precision
$\thetaq$.  Neither is needed for Corollary~\ref{mp:cor-inert}, which is \statuslean{} and
unconditional.}
\end{statusbox}

\subsection{Where cancellation can live}

The ceilings bound only the \emph{tropical} divisor.  They say nothing about the true valuation
$v_p(D)$, which can be larger when several permutations attain the minimum.  The following
observation localizes that possibility exactly.

\begin{proposition}[A unique minimizer forbids cancellation]
\label{mp:prop-unique}
Let $D=\det(E_{ij})$ be a nonzero integer determinant and let $p$ be a prime.  If a unique
permutation $\sigma_0$ minimizes $\sum_iv_p\bigl(E_{i\sigma(i)}\bigr)$, then
$v_p(D)=\sum_iv_p\bigl(E_{i\sigma_0(i)}\bigr)$.
\end{proposition}

\begin{proof}
Divide the Leibniz expansion by $p^{\tau_p}$ with $\tau_p$ the minimum.  The term indexed by
$\sigma_0$ becomes a $p$-adic unit and every other term stays divisible by $p$, so the sum is a
unit modulo $p$.
\end{proof}

\begin{statusbox}{\statuslean{} --- \lean{det_dvd_and_not_dvd_of_unique_min} in
\lean{MinPlusCeiling}, stated in the divisibility form `$q^{\tau}\mid\det$ and
$q^{\tau+1}\nmid\det$' so that no $p$-adic valuation machinery is needed.}
\end{statusbox}

The corresponding tied-minimum statement is now also exact at the algebraic level.  Suppose
each Leibniz product has been written
\[
 \prod_i E_{i,\sigma(i)}=p^\tau q_\sigma
\]
with $q_\sigma\in\Z$.  Then the determinant factors as
\begin{equation}
 \det E=p^\tau\sum_\sigma\operatorname{sgn}(\sigma)q_\sigma,
 \label{mp:eq-tropical-quotient}
\end{equation}
and, for every $k\ge0$,
\begin{equation}
 p^{\tau+k}\mid\det E
 \quad\Longleftrightarrow\quad
 p^k\mid\sum_\sigma\operatorname{sgn}(\sigma)q_\sigma.
 \label{mp:eq-tropical-all-depth}
\end{equation}
If $p$ is prime and $S$ is the tied minimum fiber, while the quotients outside $S$ are
divisible by $p$, the first extra factor is characterized exactly by the initial form
\begin{equation}
 p^{\tau+1}\mid\det E
 \quad\Longleftrightarrow\quad
 \sum_{\sigma\in S}\operatorname{sgn}(\sigma)q_\sigma=0
 \quad\text{in }\Z/p\Z.
 \label{mp:eq-tied-fiber}
\end{equation}
Thus higher-cost terms disappear modulo $p$ and all first-layer cancellation is concentrated
in the tied fiber.

\begin{statusbox}{\statuslean{} --- \lean{TropicalInitialForm} proves the factorization and
all-depth divisibility criterion as
\lean{det_eq_pow_mul_signedLeibnizQuotientSum} and
\lean{det_pow_add_dvd_iff_signedLeibnizQuotientSum_dvd}; its prime first-layer forms are
\lean{det_pow_succ_dvd_iff_signedLeibnizQuotientSum_eq_zero},
\lean{signedLeibnizQuotientSum_mod_eq_tiedFiberSum}, and
\lean{det_pow_succ_dvd_iff_tiedFiberSum_eq_zero}.}
This module is a purely algebraic prerequisite and diagnostic.  It proves neither that the
tied sum vanishes nor that it has a large higher $p$-adic valuation; in particular it supplies
no cancellation lower bound and does not close Conjecture~\ref{conj:AE}.
\end{statusbox}

The abstract block factorization needed for the report's first-fiber calculation is also now
kernel-verified.  If a block map $b$ describes the surviving permutations and $A_k$ is the
square submatrix on the fiber $b^{-1}(k)$, then
\begin{equation}
 \sum_{\sigma\text{ preserving }b}
   \operatorname{sgn}(\sigma)\prod_i A_{\sigma(i),i}
 =\prod_k\det A_k.
 \label{mp:eq-block-factorization}
\end{equation}
The formula allows varying block sizes and a fixed preliminary row permutation, whose sign is
the only extra factor.  When, after reindexing, each $A_k$ consists of consecutive powers of
nodes $x_{k,0},\ldots,x_{k,s_k-1}$, it becomes the exact Vandermonde product
\begin{equation}
 \prod_k\prod_{0\le i<j<s_k}(x_{k,j}-x_{k,i}).
 \label{mp:eq-block-vandermonde}
\end{equation}
Composing \eqref{mp:eq-block-factorization} with the tropical quotient criterion gives an
exact first-extra-factor test: under the explicit hypotheses that every quotient outside the
block-preserving set vanishes modulo $p$ and that each block-preserving quotient equals the
corresponding full Leibniz product of $A$,
$p^{\tau+1}\mid\det E$ if and only if $\prod_k\det A_k$ vanishes modulo $p$; in the
consecutive-power specialization the right side is precisely
\eqref{mp:eq-block-vandermonde} modulo $p$.

\begin{statusbox}{\statuslean{} --- \lean{BlockFiberFactorization} proves
\lean{sum_blockPreserving_eq_prod_blockDet},
\lean{sum_basePerm_mul_blockPreserving_eq_sign_mul_prod_blockDet},
\lean{sum_blockPreserving_eq_prod_vandermonde_of_reindex}, and
\lean{sum_blockPreserving_consecutivePowers_eq_prod_vandermonde}.
\lean{TropicalBlockApplication} packages the composition as
\lean{det_pow_succ_dvd_iff_prod_blockDet_eq_zero} and
\lean{det_pow_succ_dvd_iff_prod_vandermonde_eq_zero}.}
The rank-one Beatty scaffolding is now kernel-verified separately.  For matrices
\[
 E_{ij}=p^{\,n_i c(b_j)}A_{ij},
 \qquad n_i=T-\lfloor i\beta\rfloor,
\]
with strictly decreasing Beatty depths, sorted block labels $b_j$, and strictly increasing
block values $c$, \lean{BeattyTropicalFiber} proves that the minimum assignments are exactly
the block-preserving permutations.  Every other assignment gains at least
$\lfloor\beta\rfloor$, its normalized quotient is divisible by
$p^{\lfloor\beta\rfloor}$, and the determinant has the exact expansion
\[
 \det E=p^\tau\left(\prod_k\det A_k+p^{\lfloor\beta\rfloor}Z\right).
\]
The headline declarations are
\lean{rankOneAssignmentCost_eq_iff_blockPreserving},
\lean{beattyAssignmentCost_add_floor_le_of_not_blockPreserving},
\lean{beattyPower_pow_floor_dvd_quotient_of_not_blockPreserving}, and
\lean{det_beattyPowerMatrix_eq_pow_mul_blockDet_add_pow_floor_mul}.  Thus the minimizing fiber,
the tropical gap, and the full-product quotient identity no longer remain paper-only
hypotheses.  The generic scaled consecutive-power wiring is now formalized below, and
\lean{GlobalMixedPrefixCascade} closes the formerly paper-level finite specialization: its
sorted box prefix is the genuine global first-$N$ mixed-semigroup prefix, its coordinate
fibers are lower-closed, and it constructs the lexicographic and colexicographic fiber
equivalences needed for the dyadic and triadic determinant identities.  What remains globally
is the valuation and cancellation of the residual low-excess sums.  No leading-order upper
bound follows from the exact expansion.
\end{statusbox}

The complete assignment cascade can now be grouped more economically than the raw Leibniz
sum.  A \emph{row-block assignment} records only which row set is sent to each repeated column
weight, with the prescribed fiber cardinalities.  For the rank-one power matrix, the whole
Leibniz fiber over such an assignment $a$ factors as one sign times a product $m_a$ of square
unit minors.  If $c(a)$ is its structural assignment cost and $\tau$ is the minimum, then
\begin{equation}\label{mp:eq-row-block-all-layer}
 \det E=p^\tau\sum_a p^{c(a)-\tau}m_a.
\end{equation}
More sharply, for every cutoff $K$ this is the exact finite-depth decomposition
\begin{equation}\label{mp:eq-row-block-truncation}
 \det E=p^\tau\left(
   \sum_{c(a)-\tau<K}p^{c(a)-\tau}m_a
   +p^K\sum_{c(a)-\tau\ge K}p^{c(a)-\tau-K}m_a
 \right).
\end{equation}
Thus modulo $p^K$ every high-excess row partition disappears, and the surviving coefficients
are explicit products of unit minors rather than sums over individual permutations.  In the
Beatty specialization, every noncanonical row partition has excess at least
$\lfloor\beta\rfloor$.

\begin{statusbox}{\statuslean{} --- \lean{RowBlockLaplaceExpansion} proves the fiber
factorization \lean{sum_permutationRowBlockAssignmentFiber_eq_sign_mul_prod_blockDet}, the
normalized all-layer identity
\lean{det_rankOnePowerMatrix_eq_minPow_mul_sum_assignmentExcess}, the exact cutoff formula
\lean{det_rankOnePowerMatrix_eq_minPow_mul_truncatedAssignmentSum}, and the Beatty gap
\lean{beattyRowBlockAssignment_excess_ge_floor}.}
This is an exact support theorem for every finite cascade depth.  It gives no upper bound for
the number of surviving row partitions by itself, no upper bound for the valuations of the residual
minor-product quotients, and no noncancellation statement for their signed sum.  It therefore reorganizes the remaining
wall but does not bound $v_p(\det E)-\tau$.
\end{statusbox}

There is, however, now an exact entropy bound for the first of those quantities.  For a row
partition $a$, let $P_k(a)$ be the excess of the total block value assigned to the first
$k+1$ rows over the oppositely sorted assignment, and define its crossing energy by
\[
 \mathcal E(a)=\sum_{k=0}^{N-1}P_k(a).
\]
Every $P_k(a)$ is nonnegative, the vector $(P_k(a))_k$ determines $a$ when the block values
are distinct, and summation by parts identifies $\mathcal E(a)$ with the excess for the unit
depth staircase.  For Beatty depths,
\begin{equation}\label{mp:eq-beatty-crossing-energy}
 c(a)-\tau\ge \lfloor\beta\rfloor\,\mathcal E(a).
\end{equation}
Consequently, for $N,K\ge1$, if $q=\lfloor(K-1)/\lfloor\beta\rfloor\rfloor$, stars and bars gives the explicit
finite support upper bound
\begin{equation}\label{mp:eq-low-excess-support}
 \#\{a:c(a)-\tau<K\}\le \binom{N+q-1}{q}.
\end{equation}
For fixed normalized depth $q$ this is polynomial in $N$; for growing $q$ it is an explicit
entropy bound, not a cancellation estimate.

\begin{statusbox}{\statuslean{} --- \lean{AssignmentExcessGeometry} proves the exact identity
\lean{rowBlockCrossingEnergy_eq_reverseFinRank_excess}, the Beatty lower bound
\lean{beattyFloor_mul_rowBlockCrossingEnergy_le_excess}, the stars-and-bars injection
\lean{card_rowBlockAssignment_crossingEnergy_lt_choose_of_pos}, and the specialized bound
\lean{card_beattyRowBlockAssignment_excess_lt_of_pos}.}
This controls how many row partitions can occur in the cutoff sum.  It supplies no upper
bound for their minor-product valuations and no noncancellation theorem for their signed sum,
so \eqref{mp:eq-low-excess-support} is cascade-support control rather than a proof of the
conjecture.
\end{statusbox}

The fixed-divisor contribution can be threaded through this entire expansion, rather than
applied only to the canonical block.  For a block-scaled consecutive-power unit matrix, let
$s_k$ be the fixed size of block $k$, define
\[
 U_2(s)=\sum_{0\le j<s}\bigl(j+v_2(j!)\bigr),\qquad
 U_3(s)=\sum_{0\le j<s}\left(\left\lfloor\frac j2\right\rfloor+
 v_3\!\left(\left\lfloor\frac j2\right\rfloor!\right)\right),
\]
and put
\[
 \mathcal U_2=\sum_k U_2(s_k),\qquad \mathcal U_3=\sum_k U_3(s_k).
\]
Every row-block assignment has the same column-fiber sizes, so its minor product is divisible
by $2^{\mathcal U_2}$, respectively $3^{\mathcal U_3}$.  Factoring this common divisor from the exact
row-partition expansion gives
\begin{equation}\label{mp:eq-additive-unit-gain}
 2^{\tau_2+\mathcal U_2}\mid\det E_2,
 \qquad
 3^{\tau_3+\mathcal U_3}\mid\det E_3.
\end{equation}
The row scale is allowed to depend on both the target block and the row; this includes the
mixed factors $M^{ku}$ in the dyadic blocks and $A^{kv}$ in the triadic blocks.  Thus the
gain is genuinely uniform across every tropical layer, not a statement only about the first
initial form.

\begin{statusbox}{\statuslean{} --- \lean{RowBlockUnitDivisibility} defines
\lean{fiberScaledConsecutivePowerMatrix}, proves the common-minor divisors
\lean{two_pow_rowBlockOrderingExponent_dvd_minorProduct_fiberScaledConsecutivePowers} and
\lean{three_pow_rowBlockOrderingExponent_dvd_minorProduct_fiberScaledConsecutivePowers}, and
packages \eqref{mp:eq-additive-unit-gain} as
\lean{two_pow_minCost_add_rowBlockOrderingExponent_dvd_det_rankOnePowerMatrix} and
\lean{three_pow_minCost_add_rowBlockOrderingExponent_dvd_det_rankOnePowerMatrix}.}
The exponent is exactly the sum of the block gains from \lean{UnitOrderingGain} and is
independent of assignment cost.  Consequently it factors every cutoff sum in
\eqref{mp:eq-row-block-truncation}, but it cannot distinguish layers or upper-bound the
valuation of the residual signed sum.
\end{statusbox}

There is also a cascade-independent fixed divisor of the unit determinants.  A monic
triangular change of columns from $1,X,X^2,\ldots$ to an integer polynomial basis preserves
the full evaluation determinant.  On positive odd nodes with consecutive local degrees, the dyadic $p$-ordering basis
\[
 \prod_{0\le t<j}\bigl(X-(2t+1)\bigr)
\]
extracts the column weight
\[
 j+v_2(j!),
\]
while on positive nodes prime to three with consecutive local degrees, the ordering $1,2,4,5,7,8,\ldots$ extracts
\[
 \left\lfloor\frac j2\right\rfloor+
 v_3\!\left(\left\lfloor\frac j2\right\rfloor!\right).
\]
The resulting prime powers divide the \emph{whole} univariate Vandermonde determinant,
independently of which tropical layer survives.  At the univariate block level this extracts a
universal lower-order divisor rather than estimating successive initial forms.

\begin{statusbox}{\statuslean{} --- \lean{AllLayerUnitBasis} proves the generic column-content
and determinant-one column-change lemmas \lean{prod_columnDivisor_dvd_det},
\lean{pow_sum_columnWeight_dvd_det_of_unimodular_mul}, and
\lean{pow_sum_dvd_det_vandermonde_of_monicBasis}.  Its strengthened unit specializations are
\lean{two_pow_sum_degree_add_factorialVal_dvd_det_vandermonde} and
\lean{three_pow_sum_halfDegree_add_factorialVal_dvd_det_vandermonde}.}
The companion module \lean{UnitOrderingGain} proves exact digit-sum identities for the block
weights and the finite aggregate bounds
\[
 \sum_k U_2(s_k)\le S^3,
 \qquad
 2\sum_k U_3(s_k)\le S^3
\]
whenever there are at most $S$ blocks and every block has size at most $S$
(\lean{sum_twoUnitBlockGain_le_cube} and
\lean{two_mul_sum_threeUnitBlockGain_le_cube}).  For the report's $O(L)$ fibers of size
$O(L)$, the remaining paper-level block count therefore gives $O(L^3)$, equivalently
$O(N^{3/2})$, lower order than $H_T\asymp T^{5/2}$.  In fact, for the triangular cutoff
$u+\thetaq v\le L$ and $N\sim L^2/(2\thetaq)$, elementary summation gives the paper-level
leading terms for the total gains
\[
 \mathcal U_2\sim\frac{2^{3/2}}{3\sqrt{\thetaq}}N^{3/2}
   =0.7488834673\ldots N^{3/2},
 \qquad
 \mathcal U_3\sim\sqrt{\frac{\thetaq}{8}}N^{3/2}
   =0.4451070799\ldots N^{3/2}.
\]
The finite divisors and cubic aggregation are kernel-verified; the triangular fiber
asymptotics and displayed constants are not Lean theorems.  The determinant-one fiberwise
change and its additive combination with the tropical minimum are kernel-verified for the
generic scaled block model by \lean{RowBlockUnitDivisibility}; the concrete triangular index
bridge, including column permutation and transfer back to the raw mixed integer matrix, is
kernel-verified under explicit finite fiber certificates in
\lean{MixedDeterminantUnitBridge}.  The module \lean{GlobalMixedPrefixCascade} now constructs
those certificates for the genuine global first-$N$ mixed-semigroup enumeration and proves
the exact all-layer and cutoff-truncated dyadic and triadic formulas directly.  The resulting
all-layer gain neither supplies nor upper-bounds the remaining leading-order cancellation
excess.
\end{statusbox}

Consequently the missing fifth, or more, can arise only at primes where the assignment problem
ties, and the tie geometry is explicit: for $p\ne2,3$ columns tie along the rational lines
orthogonal to the external valuation vector, and in the cross-coprime branch the $2$-adic cost
depends only on $u$ and the $3$-adic cost only on $v$, so every vertical or horizontal column
fiber creates equal-cost permutations.  The continuation report evaluates the first such tied
layer.  The kernel-verified finite cube budget, combined with the elementary paper-level count
of the Beatty fiber sizes, gives $O(T^2)$ against $H_T\asymp T^{5/2}$, hence $o(1)$ of the
budget; the sharper quoted two-logarithm estimate improves this to
$O\bigl(T^{3/2}(\log T)^2\bigr)$.

\subsection{What this leaves}

Three statements survive every normalization.

First, term-wise amplification is retired unconditionally against the crude entry bound.
Corollary~\ref{mp:cor-inert} is a structural identity, not a numerical near-miss, and it holds
for every prime set, every node system and every column set.

Second, the quantitative ceilings above say how much of the archimedean budget automatic
divisibility can supply in each concrete construction, and none reaches $1$.  Adding every
prime in the fixed support does not close any determinant; it improves the diagnosis from a
heuristic ``about one half'' to a proved $4/5$.

Third --- and this is the redirection --- a successful $p$-adic proof has to be a
\emph{cancellation} theorem, not a divisibility-by-every-term theorem.
Proposition~\ref{mp:prop-unique} says exactly where such a theorem could live: inside the
equal-cost fibers, and it would have to repeat through many tropical layers or be combined with
a sharper archimedean free energy than the crude bar.  Quantifying the gain of the refined
divided-difference bounds of Sections~\ref{sec:arbnode} and~\ref{sec:semigroup} over the crude
bar is the concrete companion question, and it concerns the archimedean side alone.

\begin{statusbox}{\statusnew{} for the reading of the obstruction; \statuslean{} for
Theorem~\ref{mp:thm-ceiling}, Corollary~\ref{mp:cor-inert} and
Proposition~\ref{mp:prop-unique}.  The structural observation that the ceiling and the
shrinking-target obstruction are governed by the same block equidistribution is recorded as an
interpretive claim, not a theorem.}
\end{statusbox}

\section{Why the new hierarchy still stops short}
\label{sec:deficit}

The semigroup determinant is much stronger locally than the ordinary-power determinant: its
repulsion exponent approaches $1$.  It is therefore important to identify exactly why it does
not prove the conjecture.

\subsection{How many nodes are needed to approach scale $X$?}

The asymptotic exponent in \eqref{sg:eq-delta-asymp} is
\[
 \delta_r:=1-\frac{\Sigma_r}{R_r}
 =1-\frac{c_{\thetaq}}{\sqrt r}+O(r^{-1}),
 \qquad
 c_{\thetaq}=\frac43\sqrt{2\thetaq}=2.3739044262\ldots.
\]
For nodes near $X$, the lower span has the shape
\begin{equation}
 X^{\delta_r}
 =X\exp\!\left(-\frac{c_{\thetaq}\log X}{\sqrt r}+O\!\left(\frac{\log X}{r}\right)\right).
 \label{sg:eq-scale-loss}
\end{equation}
To make the exponential loss in \eqref{sg:eq-scale-loss} bounded independently of $X$, one
needs
\begin{equation}
 \sqrt r\asymp\log X,
 \qquad\text{that is,}\qquad
 r\asymp(\log X)^2.
 \label{sg:eq-node-threshold}
\end{equation}
This is precisely why Corollary~\ref{sg:cor-semigroup-log-square} has a logarithmic square.

\subsection{How many nodes a counterexample supplies}

Conditional on a counterexample, the exact dyadic block count established in
Section~\ref{sec:conditional} gives
\begin{equation}
 \#(\Cand\cap[X,2X])
 =\frac{\log_2X}{\beta}+O(1)
 \asymp\log X.
 \label{sg:eq-conditional-log-population}
\end{equation}
Thus the actual conditional population is smaller than the determinant threshold
\eqref{sg:eq-node-threshold} by one factor of $\log X$.

Taking $r\asymp\log X$ in \eqref{sg:eq-scale-loss} yields only
\[
 X\exp(-C\sqrt{\log X}),
\]
which is much smaller than the average spacing $X/\log X$ among the
$\asymp\log X$ conditional candidates.  The inequality is therefore compatible with the
exact monoid even at the level of expected spacing.

\subsection{Broader multiplicative intervals do not help}

The cumulative conditional count below $X$ is $\asymp(\log X)^2$, so one might try to use all
those points at once.  But their nodes range from $1$ to $X$, and the gap theorem depends on
the smallest node.  Restricting to $[X^\gamma,X]$ with fixed $0<\gamma<1$ does retain
$\asymp(\log X)^2$ points, but its additive length is $\asymp X$, while the lower scale is at
most $X^{\gamma+o(1)}$.  Again there is no contradiction.

This is not a mere defect of coarse constants.  It reflects a mismatch between additive node
geometry and the logarithmic triangular lattice of exponents $(n,k)$ in
$m=2^nM^k$.

\subsection{The determinant needs one more source of size}

The archimedean argument uses only
\[
 D\in\Z_{>0}\quad\Longrightarrow\quad D\ge1.
\]
If the exact monoid coordinates forced
\[
 D\equiv0\pmod{Q}
\qquad\text{with}\quad
 \log Q\gg R\log X/\sqrt r,
\]
then the lower bound $D\ge Q$ could remove the one-logarithm deficit.  Alternatively, one
could seek a determinant with total exponent weight $\Sigma=O(r)$ rather than
$\Sigma\asymp r^{3/2}$, but the two-dimensional semigroup counting law
$L(T)\asymp T^2$ makes that impossible using only distinct nonnegative mixed monomials.

This gives a concrete design specification:

\begin{quote}
A successful determinant proof must add nonarchimedean divisibility, introduce lower-weight
functions outside the mixed-monomial semigroup, or exploit the exact two-dimensional row grid
so efficiently that $O(\!\log X)$ nodes suffice.
\end{quote}

\begin{statusbox}{\statusnew}
The deficit analysis of this section is a paper argument about the limits of the method, not
a formalized theorem.  It does not weaken any statement proved above; it explains why the
explicit $O((\log X)^2)$ bound of Corollary~\ref{sg:cor-semigroup-log-square} cannot be
pushed to the conditional $\asymp\log X$ population by tuning constants alone.  The
conjecture remains open.
\end{statusbox}
\section{Further equivalent reformulations}
\label{sec:reform}

Several equivalent forms are already formalized in the repository and were catalogued in
Section~\ref{sec:corpus}.  The exact conditional structure of Section~\ref{sec:conditional}
also yields useful paper-level reformulations that expose different possible proof
technologies.  Each restatement below is an exact equivalence, not a heuristic analogy: a
proof of any one of them is a proof of Conjecture~\ref{conj:AE}.

\subsection{Integer points on a transcendental lattice curve}

The candidate formulation of Section~\ref{sec:problem} says that
\begin{equation}
 \{(m,n)\in\N^2:n=m^{\thetaq}\}
 =\{(2^k,3^k):k\in\Nzero\}.
 \label{rf:eq:lattice-curve}
\end{equation}
The curve $n=m^{\thetaq}$ is strictly convex and transcendental, and \eqref{rf:eq:lattice-curve}
asserts that its only lattice points are the trivial exponential ones.  The arbitrary-node
repulsion theorem of Section~\ref{sec:arbnode} can be read as a lower bound for the volumes of
lattice simplices cut from this curve; the deficit analysis of Section~\ref{sec:deficit}
measures exactly how far the current bound is from excluding the nontrivial points.

\subsection{Vanishing determinant of four logarithms}

Writing $M=2^x$ and $A=3^x$, a solution is a positive integer pair with
\begin{equation}
 \log M\,\log3-\log A\,\log2=0,
 \quad M,A\in\N,
 \quad\Longrightarrow\quad
 (M,A)=(2^k,3^k).
 \label{rf:eq:four-log-form}
\end{equation}
The hypothesis of \eqref{rf:eq:four-log-form} is a vanishing $2\times2$ determinant of
logarithms of positive integers.  This is the most direct bridge to the four-exponentials
circle of ideas and to algebraic independence of logarithms, and it is the form in which the
kernel-verified four-exponentials reduction is stated.

\subsection{Primitive localized-radical exclusion}

The kernel-verified radical form of Section~\ref{sec:radical} is
\begin{equation}
 \forall\,w>1\text{ odd and power-primitive},\quad
 \forall\,d>0,\ a\ge0,\qquad
 3^a\bigl(w^{\thetaq}\bigr)^d\notin\N.
 \label{rf:eq:primitive-localized}
\end{equation}
If \eqref{rf:eq:primitive-localized} fails with least rational-power index $d$, then
$X^d-c/3^a$ is the exact minimal polynomial of $w^{\thetaq}$ over $\Q$ for the corresponding
rational $c$.

This formulation is well suited to Kummer-theoretic, norm, and ramification investigations,
because all redundant proper powers have been removed by primitivity and the algebraic degree
of the hypothetical radical is known exactly rather than merely bounded.

\subsection{Exact growth dichotomy}

Write
\[
 B(T):=\#\bigl(\Cand\cap[2^T,2^{T+1})\bigr)
\]
for the $T$th dyadic block count.  The repository formalizes the cumulative subquadratic
equivalence of Section~\ref{sec:conditional}; the exact block formula there, namely
\begin{equation}
 B(T)=1+\left\lfloor\frac{T+1}{\beta}\right\rfloor
 \qquad\text{under failure with least generator }\beta,
 \label{rf:eq:dyadic-count}
\end{equation}
gives several immediate variants.

\begin{proposition}[Block-growth reformulations]
\label{rf:prop:block-reforms}
Conjecture~\ref{conj:AE} is equivalent to each of the following:
\begin{enumerate}
\item the dyadic block counts $B(T)$ are bounded;
\item $B(T)=o(T)$;
\item $\displaystyle\liminf_{T\to\infty}B(T)/T=0$;
\item $B(T)=1$ for every $T\in\Nzero$.
\end{enumerate}
\end{proposition}

\begin{proof}
If the conjecture holds, every block contains only $2^T$, so all four properties hold.  If it
fails, \eqref{rf:eq:dyadic-count} gives $B(T)=1+\lfloor(T+1)/\beta\rfloor$, so
$B(T)/T\to1/\beta>0$ and none of the first three properties holds.  The fourth is visibly
equivalent to the candidate statement.
\end{proof}

This is useful because a future analytic theorem need not classify candidates individually.
Any uniform $o(T)$ bound for the number of candidates in the $T$th dyadic block would prove
the conjecture.

The exactness of \eqref{rf:eq:dyadic-count} permits a sharpening of item~(4) that is, at first
sight, dramatically weaker than the conjecture and yet still equivalent to it.

\begin{proposition}[Infinitely many singleton blocks suffice]
\label{rf:prop:infinitely-many-blocks}
Conjecture~\ref{conj:AE} is equivalent to the statement
\[
 B(T)=1\quad\text{for infinitely many }T\in\Nzero.
\]
\end{proposition}

\begin{proof}
If the conjecture holds then $B(T)=1$ for every $T$, so certainly for infinitely many.
Conversely, suppose the conjecture fails, and let $\beta>0$ be the irrational least generator
of the exact solution monoid.  By \eqref{rf:eq:dyadic-count},
\[
 B(T)=1
 \iff
 \left\lfloor\frac{T+1}{\beta}\right\rfloor=0
 \iff
 T+1<\beta ,
\]
which holds for at most the finitely many indices $T<\beta-1$.  Hence failure forces $B(T)=1$
for only finitely many $T$, and the contrapositive gives the equivalence.
\end{proof}

\begin{remark}
\label{rf:rem:not-some-T}
Proposition~\ref{rf:prop:infinitely-many-blocks} cannot be weakened further to ``$B(T)=1$ for
some $T$''.  That statement is already known unconditionally and carries no information: the
kernel-verified zero-free region $2^x<4096$ gives $\beta\ge12$, hence $B(T)=1$ for
$T=0,\dots,10$, and the interval computation of Section~\ref{sec:finitecheck} pushes the same
conclusion to $T\le19$ at the weaker \statuscomp{} trust level.  The gap between ``for some
$T$'' and ``for infinitely many $T$'' is exactly the gap between a finite check and the
conjecture, because none of the present results bounds $\beta$ from above: a counterexample
with a very large least generator would have $B(T)=1$ for every $T<\beta-1$, that is, on any
range a fixed finite computation could ever reach.  Only the infinitude of singleton blocks is
decisive.
\end{remark}

\begin{statusbox}{\statuslean{} for Proposition~\ref{rf:prop:block-reforms}:
\lean{BlockGrowthReformulations} proves the bounded, little-$o$, liminf-zero, and identically
one equivalences.  \statusnew{} for Proposition~\ref{rf:prop:infinitely-many-blocks}, whose
short paper proof rests on the same kernel-verified exact dyadic block formula.}
\end{statusbox}

\subsection{Compact synchronized \texorpdfstring{$S$}{S}-integral orbit}

The $S$-integral orbit of Section~\ref{sec:conditional} yields another equivalence, this time
one that isolates the synchronization of the two denominators.

\begin{proposition}[Compact-orbit reformulation]
\label{rf:prop:compact-orbit}
Conjecture~\ref{conj:AE} fails if and only if there exist integers $M,A>1$ and an irrational
$\beta>0$ such that, for every $k\ge1$, the point
\begin{equation}
 \left(
  \frac{M^k}{2^{\lfloor k\beta\rfloor}},
  \frac{A^k}{3^{\lfloor k\beta\rfloor}}
 \right)
 \label{rf:eq:compact-orbit-points}
\end{equation}
lies on the compact arc
\[
 \mathcal A=\{(2^t,3^t):0\le t<1\}.
\]
Under these conditions the orbit is dense in $\mathcal A$.
\end{proposition}

\begin{proof}
Failure supplies the least generator $\beta$, with $M=2^\beta$ and $A=3^\beta$; then
\eqref{rf:eq:compact-orbit-points} is exactly
$\bigl(2^{\fract{k\beta}},3^{\fract{k\beta}}\bigr)$, and density follows from irrational
rotation.

Conversely, arc membership gives $t_k\in[0,1)$ such that
\[
 k\log_2M-\lfloor k\beta\rfloor=t_k
 =k\log_3A-\lfloor k\beta\rfloor.
\]
Thus $\log_2M=\log_3A=:\gamma$, and $0\le k\gamma-\lfloor k\beta\rfloor<1$ for every $k$.
Hence $\lfloor k\gamma\rfloor=\lfloor k\beta\rfloor$, so $|k(\gamma-\beta)|<1$ for every $k$;
therefore $\gamma=\beta$.  Consequently $M=2^\beta$ and $A=3^\beta$ are integers while $\beta$
is nonintegral, so the conjecture fails.  The same identity now shows that the orbit equals
the irrational-rotation orbit and is dense.
\end{proof}

\begin{statusbox}{\statuslean{} --- \lean{not_alaogluErdosConjecture_iff_exists_compactOrbit}
and \lean{alaogluErdosConjecture_iff_no_compactOrbit} in \lean{CompactOrbit}.}
\end{statusbox}

This formulation isolates the synchronized-denominator feature that a specialized
$S$-integral counting theorem would need to exploit: the two coordinates are $S$-integral for
disjoint sets of primes, yet their denominator exponents are forced to agree exactly at every
height.

\subsection{Two-dimensional exponential grid}

Under failure, every solution coordinate $(n,k)$ gives
\[
 m_{n,k}=2^nM^k,
 \qquad
 m_{n,k}^{\thetaq}=3^nA^k,
\]
and for every mixed exponent $a+b\thetaq$ with $a,b\in\Nzero$,
\begin{equation}
 m_{n,k}^{\,a+b\thetaq}
 =(2^a3^b)^n(M^aA^b)^k.
 \label{rf:eq:two-dimensional-grid}
\end{equation}
Thus the candidate evaluation matrix factors through a two-dimensional row lattice and a
two-dimensional column lattice.  The conjecture can be viewed as the assertion that this
integer exponential grid cannot exist unless $(M,A)$ is a common integral power of $(2,3)$.
This is the most natural starting point for a bespoke $2\times2$ interpolation determinant
that uses more than the arbitrary-node geometry of Section~\ref{sec:arbnode}, and it is the
structure the semigroup hierarchy of Section~\ref{sec:semigroup} exploits only partially.

\subsection{An operator formulation on the prime-logarithm space}
\label{rf:operator}

All the reformulations above are Diophantine.  There is also a purely linear-algebraic one,
which recasts the problem as a statement about multiplication by $\thetaq$ acting on a single
$\Q$-vector space.

Let
\[
  W:=\operatorname{span}_{\Q}\{\log p\ :\ p\ \text{prime}\}\subset\R ,
\]
a $\Q$-vector space of countably infinite dimension, with the $\log p$ as a basis by unique
factorization.  Multiplication by $\thetaq$ does not preserve $W$, so consider the largest
part of $W$ that it does return to $W$:
\[
  U:=\{\,w\in W\ :\ \thetaq\,w\in W\,\}.
\]
This is a $\Q$-subspace, and it is never zero: $\thetaq\log2=\log3\in W$, so
$\Q\log2\subseteq U$.

\begin{proposition}[Operator form]
\label{rf:prop-operator}
$\dim_{\Q}U=\operatorname{rank}K$, where
$K=\{v\in\Q_{>0}:v^{\thetaq}\in\Q\}$ is the multiplicative group of
Section~\ref{sec:logproduct}.  Consequently the rational-output form
\eqref{eq:lp-rational-form} --- and hence Conjecture~\ref{conj:AE} --- is equivalent to
\[
  \boxed{\ \dim_{\Q}\{\,w\in W:\thetaq\,w\in W\,\}=1 .\ }
\]
\end{proposition}

\begin{proof}
An element of $W$ is $w=\sum_pc_p\log p$ with rational $c_p$, almost all zero.  Clearing a
common denominator $N$ gives $Nw=\log v$ for a positive rational $v$.  Now $\thetaq w\in W$
says $\thetaq\log v=\sum_pd_p\log p$ with rational $d_p$; clearing a denominator $M$ turns
this into $\thetaq\log\bigl(v^{M}\bigr)=\log u$ with $u\in\Q_{>0}$, that is
$\bigl(v^{M}\bigr)^{\thetaq}=u$, i.e.\ $v^{M}\in K$.  So $U$ is exactly the $\Q$-span of
$\log K$ inside $W$, whence $\dim_{\Q}U=\operatorname{rank}K$.  The conjecture in its
rational-output form says $K=\langle2\rangle$, of rank one.
\end{proof}

\begin{statusbox}{\statusnew{} --- an elementary translation, but a genuinely different
framing: it replaces a Diophantine assertion by the statement that the operator ``multiply by
$\thetaq$'' carries only a single line of $W$ back into $W$.}
\end{statusbox}

Two remarks on what this framing does and does not give.

\emph{It explains why finite-dimensional arguments cannot apply directly.}  If some nonzero
finite-dimensional subspace $V\subseteq W$ were $\thetaq$-\emph{stable}, then $\thetaq$ would
be an eigenvalue of the rational matrix representing multiplication by $\thetaq$ on $V$, hence
algebraic, contradicting transcendence.  So a proof cannot proceed by exhibiting an invariant
finite-dimensional subspace: none exists, unconditionally.  The subspace $U$ above is not
$\thetaq$-stable --- $\thetaq U\subseteq W$ but generally $\thetaq U\not\subseteq U$ --- and
this failure of stability is precisely what blocks the eigenvalue argument.

\emph{The natural stable candidate is empty for the same reason.}  Setting
$U_\infty=\{w\in W:\thetaq^{k}w\in W\ \text{for all}\ k\ge0\}$ does produce a
$\thetaq$-stable space, but it is either zero or infinite-dimensional, and infinite dimension
is automatic whenever it is nonzero: if $0\ne w\in U_\infty$, the powers $\thetaq^{k}w$ are
$\Q$-linearly independent, since a relation $\sum_kc_k\thetaq^{k}w=0$ with $w\ne0$ would make
$\thetaq$ algebraic.  So the operator picture reproduces the transcendence barrier exactly,
rather than circumventing it.  Its value is diagnostic: it shows the obstruction is not a
missing estimate but the absence of any finite-dimensional invariant structure to work with.

\section{A consolidated catalogue of equivalent formulations}\label{sec:catalogue}

The preceding sections restate Conjecture~\ref{conj:AE} many times over: as a candidate-base
identity, as a localized radical exclusion, as a counting dichotomy, as a statement about
fractional parts, and as a closure property of the solution monoid.  Because these
restatements were introduced where they were needed rather than where they belong together,
this section collects them in one place.

Throughout, $\thetaq=\log_23$,
\[
 \Sol=\{x\in\R:2^x,3^x\in\Z\},
 \qquad
 \Cand=\{m\in\N:m>0,\ m^{\thetaq}\in\N\},
\]
which correspond under $x\mapsto2^x$ by the exponent--candidate bijection of
Section~\ref{sec:problem}, and
\[
 B(T):=\#\bigl(\Cand\cap[2^T,2^{T+1})\bigr)
\]
is the $T$th dyadic block count.

\subsection{How to read the catalogue}

Every item (E1)--(E26) below is an \emph{exact biconditional}: a proof of any one of them is a
proof of Conjecture~\ref{conj:AE}, and a disproof of any one of them is a counterexample.
None of them is proved.  The status tag on each item records the strength of the evidence for
the \emph{equivalence}, never for the conjecture: \statuslean{} means the biconditional itself
has been accepted by the Lean kernel, \statusnew{} means it is a paper proof in this report,
and \statusknown{} means the equivalence is obtained by combining a paper proof here with a
classical transcendence theorem.

Two adjacent notions are deliberately excluded from the catalogue and treated separately in
Section~\ref{eq:sufficient} below: \emph{sufficient conditions}, which are one-way
implications from hypotheses that are themselves unproved, and \emph{unconditional partial
results}, which settle a bounded region rather than the statement.  Conflating the three is
the most common way to overstate what the corpus establishes.

\subsection{Cluster A: normalized restatements}

These three items are the conjecture rewritten in the coordinates in which the rest of the
corpus works.  They carry no arithmetic content beyond that bijection, and their value is that
each suggests a different toolkit.

\begin{enumerate}[label=\textup{(E\arabic*)}]
\item \statuslean{} \emph{Candidate form.} $\Cand=\{2^n:n\in\Nzero\}$; equivalently, the only
positive integers $m$ with $m^{\thetaq}\in\N$ are the powers of two
(Section~\ref{sec:problem}).
\item \statuslean{} \emph{Vanishing four-log determinant.} The only positive integer pairs
$(M,A)$ with $\log M\,\log3-\log A\,\log2=0$ are $(M,A)=(2^n,3^n)$, $n\in\Nzero$
(Section~\ref{sec:problem}); this is the $2\times2$ logarithmic determinant whose entries
include two logarithms of unknown integers, and the form in which the four-exponentials
reduction of Section~\ref{sec:fourexp} is stated.
\item \statuslean{} \emph{One-parameter subgroup intersection.} The real one-parameter
subgroup $\Gamma=\{(2^t,3^t):t\in\R\}$ of the multiplicative torus meets $\N^2$ exactly in
$\{(2^n,3^n):n\in\Nzero\}$ (Section~\ref{sec:problem}); a geometric restatement of (E2), with
the difficulty that the direction vector $(\log2,\log3)$ is transcendental rather than
algebraic.
\end{enumerate}

\begin{statusbox}{\statuslean{} --- the exponent--candidate bijection is
\lean{exists_twoBaseNaturalCandidate_of_two_three_rpow_integer} together with
\lean{TwoBaseNaturalCandidate.integer_powers}, and the equivalence (E1) is
\begin{center}
\lean{alaogluErdosConjecture_iff_candidates_are_powers_of_two}.
\end{center}
Items (E2) and (E3) are the same content transported along that bijection and along
$m\mapsto(m,m^{\thetaq})$; they carry no separate Lean identifier.}
\end{statusbox}

\subsection{Cluster B: localized radicals}

This cluster is the most compressed form of the problem known: a statement quantified over all
real $x$ becomes a statement about one integer radical at a time, with all $2$-adic and all but
one $3$-adic degree of freedom eliminated.  It is also the cluster closest to established
transcendence theory.

\begin{enumerate}[resume,label=\textup{(E\arabic*)}]
\item \statuslean{} \emph{Localized radical exclusion.} There is no odd integer $u>1$ with
$u^{\thetaq}\in\Zthree$ (Section~\ref{sec:radical}).
\item \statuslean{} \emph{Integral localized form.} There is no pair $(u,a)$ with $u>1$ odd,
$a\in\Nzero$ and $3^au^{\thetaq}\in\N$; equivalently, no candidate has odd part greater than
one (Section~\ref{sec:radical}).
\item \statuslean{} \emph{Primitive localized radical exclusion.} There is no power-primitive
odd $w>1$ and $d\ge1$ with $\bigl(w^{\thetaq}\bigr)^d\in\Zthree$
(Section~\ref{sec:radical}); primitivity removes every proper-power redundancy, and under
failure the minimal polynomial of $w^{\thetaq}$ is known exactly.
\item \statuslean{} \emph{Symmetric base-swapped form.} With $\eta=\log_32=1/\thetaq$, there is
no integer $v>1$ with $3\nmid v$ and $v^{\eta}\in\Ztwo$ (Appendix~\ref{app:symmetric}); the two
normalizations produce the same least exponent, so the reduction is not an artifact of
privileging the base $2$.
\end{enumerate}

\begin{statusbox}{\statuslean{} --- (E4) is
\begin{center}
\lean{alaogluErdosConjecture_iff_odd_rpow_not_localized}
\end{center}
in \lean{Localization}, restated as
\begin{center}
\lean{alaogluErdosConjecture_iff_oddLocalizedRadicalExclusion}
\end{center}
in \lean{RadicalReformulation}; (E5) is the independently proved candidate-level variant
\begin{center}
\lean{alaogluErdosConjecture_iff_no_odd_multiple_candidate};
\end{center}
(E6) is
\begin{center}
\lean{alaogluErdosConjecture_iff_primitiveOddLocalizedRadicalExclusion}
\end{center}
in \lean{PrimitiveRadical}; and (E7) is
\begin{center}
\lean{alaogluErdosConjecture_iff_threeFreeLocalizedRadicalExclusion}
\end{center}
in \lean{RadicalReformulation}.  The equivalences are kernel-verified; the exclusions on their
right-hand sides are the open conjecture.}
\end{statusbox}

\subsection{Cluster C: counting and block growth}

The counting cluster replaces an integrality statement about individual candidates by an upper
bound on how many of them there can be.  It is attractive because an analytic theorem need
never identify a candidate, and unattractive because the target quantities are extremely far
from the arithmetic conclusion; Section~\ref{sec:deficit} measures that distance precisely.

\begin{enumerate}[resume,label=\textup{(E\arabic*)}]
\item \statuslean{} \emph{Subquadratic cumulative growth.} $\#\bigl(\Cand\cap[1,2^T]\bigr)$ is
$o(T^2)$ as $T\to\infty$ (Section~\ref{sec:conditional}).
\item \statuslean{} \emph{Bounded blocks.} The dyadic block counts $B(T)$ are bounded
(Section~\ref{sec:reform}).
\item \statuslean{} \emph{Sublinear blocks.} $B(T)=o(T)$ (Section~\ref{sec:reform}).
\item \statuslean{} \emph{Sublinear along a subsequence.} $\liminf_{T\to\infty}B(T)/T=0$
(Section~\ref{sec:reform}).
\item \statuslean{} \emph{Singleton blocks.} $B(T)=1$ for every $T\in\Nzero$
(Section~\ref{sec:reform}); an immediate restatement of (E1) blockwise.
\item \statusnew{} \emph{Infinitely many singleton blocks.} $B(T)=1$ for infinitely many
$T\in\Nzero$ (Section~\ref{sec:reform}).
\end{enumerate}

\begin{remark}
\label{eq:rem-infinitely-many}
Item (E13) is the sharpest illustration of how exact the conditional structure theory is.  It
looks dramatically weaker than (E12), yet the exact block formula
$B(T)=1+\lfloor(T+1)/\beta\rfloor$ under failure with least generator $\beta$ forces $B(T)=1$
to hold only for the finitely many $T<\beta-1$.  It cannot be weakened further to ``$B(T)=1$
for some $T$'': that statement is already known unconditionally and carries no information,
since the kernel-verified zero-free region gives $\beta\ge12$.  The gap between ``for some''
and ``for infinitely many'' is exactly the gap between a finite check and the conjecture.
\end{remark}

\begin{statusbox}{\statuslean{} for (E8) ---
\begin{center}
\lean{alaogluErdosConjecture_iff_candidate_count_below_two_pow_subquadratic}.
\end{center}
\statuslean{} for (E9)--(E12) ---
\lean{alaogluErdosConjecture_iff_dyadicBlock_bounded},
\lean{alaogluErdosConjecture_iff_dyadicBlock_isLittleO},
\lean{alaogluErdosConjecture_iff_dyadicBlock_liminf_zero}, and
\lean{alaogluErdosConjecture_iff_dyadicBlock_eq_one} in
\lean{BlockGrowthReformulations}.  \statusnew{} for (E13), whose paper proof rests on the
kernel-verified exact block formula.}
\end{statusbox}

\subsection{Cluster D: orbits and fractional parts}

Under failure the solution set is an irrational-rotation orbit in disguise.  This cluster
converts the conjecture into the assertion that no such orbit can be arithmetically
synchronized, and into three gap statements on the unit interval.

\begin{enumerate}[resume,label=\textup{(E\arabic*)}]
\item \statuslean{} \emph{Compact synchronized $S$-integral orbit.} There are no integers
$M,A>1$ and irrational $\beta>0$ such that
$\bigl(M^k/2^{\lfloor k\beta\rfloor},A^k/3^{\lfloor k\beta\rfloor}\bigr)$ lies on the compact
arc $\{(2^t,3^t):0\le t<1\}$ for every $k\ge1$ (Section~\ref{sec:reform}).  The two coordinates
are $S$-integral for disjoint prime sets, yet their denominator exponents must agree exactly at
every height.
\item \statuslean{} \emph{Avoided interval.} Some nonempty open interval $(u,v)\subseteq(0,1)$
contains no fractional part $\fract{y}$ of a nonintegral solution $y$
(Section~\ref{sec:corpus}).
\item \statuslean{} \emph{Gap at the top.} The fractional parts of nonintegral solutions are
bounded away from $1$ (Section~\ref{sec:corpus}).
\item \statuslean{} \emph{Gap at the bottom.} The fractional parts of nonintegral solutions are
bounded away from $0$ (Section~\ref{sec:corpus}).
\end{enumerate}

Items (E15)--(E17) are unconditional biconditionals, and they say something quantitative about
the elementary gap bounds of Section~\ref{sec:corpus}: those bounds place $\fract{x}$ away from
$0$ and $1$ by margins that decay like $2^{-\lfloor x\rfloor}$, and the equivalences show that
\emph{any} height-independent margin, however small and wherever located, is already the full
conjecture.  The decay is therefore essential rather than an artifact of the estimate.

\begin{statusbox}{\statuslean{} for (E14) ---
\lean{alaogluErdosConjecture_iff_no_compactOrbit} in \lean{CompactOrbit}.  \statuslean{} for
(E15)--(E17) ---
\begin{center}
\lean{alaogluErdosConjecture_iff_fract_avoids_interval},
\end{center}
\begin{center}
\lean{alaogluErdosConjecture_iff_fract_bounded_away_one},
\end{center}
\begin{center}
\lean{alaogluErdosConjecture_iff_fract_bounded_away_zero}
\end{center}
in \lean{FractionalPartDensity}.}
\end{statusbox}

\subsection{Cluster E: closure, smoothness, and algebraicity}

The last cluster is the least classical and, for that reason, the most interesting as a source
of new attacks.  It says that the conjecture is equivalent to the failure of a single nonlinear
closure property: the solution monoid, which is trivially closed under addition, must not also
be closed under multiplication.

\begin{enumerate}[resume,label=\textup{(E\arabic*)}]
\item \statusnew{} \emph{Multiplicative closure.} $\Sol$ is closed under multiplication.
\item \statusnew{} \emph{Square closure.} $\Sol$ is closed under squaring.
\item \statusnew{} \emph{Star closure.} $\Cand$ is closed under the operation
$m\star n=2^{\log_2m\,\log_2n}$.
\item \statusnew{} \emph{Star idempotence.} Every $m\in\Cand$ satisfies $m\star m\in\Cand$.
\item \statusnew{} \emph{Smooth outputs.} Every $x\in\Sol$ has both $2^x$ and $3^x$
$3$-smooth.
\item \statusnew{} \emph{Smooth product output.} Every $x\in\Sol$ has $6^x$ $3$-smooth.
\item \statusknown{} \emph{Square-exponent algebraicity.} Every $x\in\Sol$ has
$6^{x^2}\in\Abar$.
\item \statusknown{} \emph{Split square-exponent algebraicity.} Every $x\in\Sol$ has both
$2^{x^2}\in\Abar$ and $3^{x^2}\in\Abar$.
\item \statusknown{} \emph{Mixed-product algebraicity.} All $x,y\in\Sol$ satisfy
$6^{xy}\in\Abar$.
\end{enumerate}

\begin{remark}[Manufacturing a third base]
\label{eq:rem-third-base}
Items (E22)--(E26) admit a concrete reading.  By the kernel-verified $3$-smooth classification
of Section~\ref{sec:corpus}, a hypothetical nonintegral solution $\beta$ admits an integer base
$B$ with $B^{\beta}\in\Z$ only when $B$ is $3$-smooth.  It would therefore suffice to
manufacture \emph{one} non-$3$-smooth positive integer $B$ with $B^{\beta}\in\Z$.  The obvious
candidates $B=2^{\beta}$, $B=3^{\beta}$, $B=6^{\beta}$ have $\beta$-th powers
$2^{\beta^2}$, $3^{\beta^2}$, $6^{\beta^2}$, so ``manufacture a third base'' is exactly the
square-closure problem of (E19) in different language.  Whenever the chosen base is not
$3$-smooth its $\beta$-th power is already known to be transcendental by six exponentials, so
any proof of algebraicity would finish the conjecture at once.  What is missing is a single
bridge:
\[
 \boxed{\text{derive one nonlinear algebraic value from the two linear integer values}.}
\]
Neither unique factorization, nor the ordinary exponent laws, nor linear forms in logarithms
supplies it.
\end{remark}

\begin{statusbox}{\statusnew{} for (E18)--(E23) --- paper proofs in this report, from the
closure equivalences and the output-smoothness theorem; no Lean formalization yet.
\statusknown{} for (E24)--(E26) --- the equivalences combine those paper proofs with the
classical six exponentials theorem, which supplies transcendence of a non-$3$-smooth base
raised to the hypothetical exponent.  None of the nine items has been accepted by the kernel.}
\end{statusbox}

\subsection{Which formulations are worth attacking}

The twenty-six items are not equally promising, and it is worth saying so plainly.  Cluster C
is exact but analytically distant: the deficit analysis of Section~\ref{sec:deficit} shows the
determinant hierarchy reaching $O((\log X)^2)$ while a counterexample supplies $\asymp\log X$
candidates, a full logarithm short of the $o(T)$ that (E10) would need.  Cluster A is a change
of coordinates, useful for orientation and for locating the four-exponentials barrier, but not
itself a method.  Cluster B is closest to established transcendence theory and is where a
Kummer-theoretic, norm, or ramification argument would land; its smallest open instance is the
irrationality of $3^{\log_23}$.  Cluster D isolates a synchronization phenomenon that a bespoke
$S$-integral counting theorem might exploit.  Cluster E is the only one whose hypotheses are
elementary and nonlinear at the same time, and it is the natural home of any attempt to
manufacture a third base.

\subsection{Equivalences versus sufficient conditions}
\label{eq:sufficient}

An equivalence and a sufficient condition play very different roles, and the distinction is
easy to lose in a corpus this size.  Items (E1)--(E26) are biconditionals: each is exactly as
strong as Conjecture~\ref{conj:AE}, so proving any of them proves the conjecture and no
weakening of any of them is safe without checking that the reverse implication survives.  The
statements below are strictly one-way.  Each is a hypothesis that is itself open, each implies
the conjecture, and none is known to follow from it --- so none belongs in the catalogue.

\begin{enumerate}[label=\textup{(S\arabic*)}]
\item \emph{The real four exponentials conjecture.}  If for all $\Q$-linearly independent
$x_1,x_2$ and $\Q$-linearly independent $y_1,y_2$ some $e^{x_iy_j}$ is transcendental, then
Conjecture~\ref{conj:AE} holds.  The implication is kernel-verified
(Section~\ref{sec:fourexp}); the hypothesis is open, and is the reason no routine combination
of present methods closes the problem.
\item \emph{Schanuel's conjecture.}  It implies algebraic independence of the logarithms of the
primes, hence (S3), hence Conjecture~\ref{conj:AE} (Section~\ref{sec:logproduct}).
\item \emph{Prime-log-product independence.}  If the real numbers $\log p\,\log q$, indexed by
pairs of primes $p\le q$, are linearly independent over $\Q$, then Conjecture~\ref{conj:AE}
holds (Section~\ref{sec:logproduct}).  This criterion mentions no exponentials at all, is
strictly weaker than Schanuel, and is not known to imply or to follow from (S1).
\end{enumerate}

\begin{statusbox}{\statuslean{} for (S1) ---
\begin{center}
\lean{alaogluErdosConjecture_of_realFourExponentialsConjecture}
\end{center}
in \lean{FourExponentialsReduction}, which formalizes the implication without assuming its
hypothesis.  \statusnew{} for the Schanuel implication in (S2).  \statuslean{} for the
conditional implication in (S3): \lean{LogProductIndependence} records the independence
hypothesis as an explicit \lean{Prop}, never asserted, with headline statement
\begin{center}
\lean{alaogluErdosConjecture_of_primeLogProductIndependence}.
\end{center}}
\end{statusbox}

The finite prime-support form of (S3) is weaker still and correspondingly more approachable: if
the products $\log p\log q$ over a fixed finite prime set $S\ni2,3$ are $\Q$-linearly
independent, then no counterexample has both outputs supported on $S$.  For $S=\{2,3\}$ this
reduces to $\Q$-linear independence of $1$, $\thetaq$ and $\thetaq^2$, which would exclude
every $3$-smooth counterexample.  That is a partial exclusion, not an equivalence.

A third category deserves the same care.  Unconditional partial results settle a bounded region
and are neither equivalences nor sufficient conditions: the kernel-verified zero-free region
gives $2^x\ge4096$, hence $x\ge12$, for every nonintegral solution
(Section~\ref{sec:finitecheck}), \statuslean{}; the directed-rounding scan through
$2^{27}$ extends the same conclusion at the weaker trust level of a software computation,
\statuscomp{}, and is not a kernel proof.  The kernel-verified exponent bound remains $x\ge12$.
No item in this section, and no combination of them, proves Conjecture~\ref{conj:AE}; the
catalogue records where a proof could enter, not that one has.

\section{Independent finite verification}\label{sec:finitecheck}

Exactly one zero-free region in this report is a kernel certificate.  Lean proves that a
solution with $2^x<4096$ forces $x\in\Z$, so that a nonintegral solution satisfies $x\ge12$;
the statement lives in \lean{FiniteCheck4096} as
\begin{center}
\lean{twelve_le_of_not_integer_of_two_three_rpow_integer}
\end{center}
and it is the only finite exclusion that any other statement here may use without adding
trust assumptions.  \statuslean

One further finite computation pushes the excluded range far past $4096$: a
directed-rounding exclusion through $2^{27}$, described in
Section~\ref{sec:finite27}.  It supersedes the earlier interval scans, which are not
restated here.  It is a software result: it does not replay in the Lean kernel, it is not
interchangeable with \lean{FiniteCheck4096}, and it does not raise the kernel-verified
exponent bound above $x\ge12$.  They are recorded here because they strengthen the explicit
conditional picture, test the conditional models of Section~\ref{sec:conditional} against
data, and supply the extremal cases that a future generated certificate must clear.

\subsection{What each scan must certify}
\label{cp:sub-certificate}

Write $\thetaq=\log_2 3$, and let $m$ be a positive integer that is not a power of two.
Excluding $m$ means proving $m^{\thetaq}\notin\N$, and it suffices to exhibit an integer $n$
with $n<m^{\thetaq}<n+1$.  Exponentiation is never performed: the pair of strict inequalities
is equivalent, after taking logarithms and clearing the positive factor $\log 2$, to
\begin{align}
\log m\,\log3-\log n\,\log2&>0,\label{cp:eq-lower-log}\\
\log(n+1)\,\log2-\log m\,\log3&>0.\label{cp:eq-upper-log}
\end{align}
Every logarithm and every product occurring in the scanned ranges is positive, so the
required rounding directions are monotone and unambiguous.

In both scans an ordinary floating-point evaluation merely \emph{proposes}
$n\approx\lfloor m^{\thetaq}\rfloor$.  Once certified lower bounds for
\eqref{cp:eq-lower-log} and \eqref{cp:eq-upper-log} are both strictly positive, the floating
proposal has no remaining logical role: it is a witness generator, not a step of the
argument.  A locator error can therefore create a spurious failure report, but it cannot
create a spurious exclusion.  Powers of two are skipped rather than tested, since $2^q$ has
the integral value $(2^q)^{\thetaq}=3^q$; an exactly integral value at a non-power of two
would satisfy neither strict inequality for any $n$ and would be reported as a failure.

\subsection{Superseded scans}

Two earlier software exclusions --- an interval-arithmetic scan through $2^{20}$ and a
directed-rounding scan through $2^{26}$ --- have been absorbed by the extension through
$2^{27}$ of Section~\ref{sec:finite27}, which uses the same directed-rounding
method, the same margin conventions and the same manifest discipline over the wider range.
Only the strongest exclusion is stated in this report; the superseded ranges and their
chunk tables are not reproduced.

\subsection{The trust boundary, stated plainly}
\label{cp:sub-boundary}

\begin{resultbox}
\textbf{Kernel.}  $2^x<4096$ and $2^x,3^x\in\Z$ imply $x\in\Z$; a nonintegral solution has
$x\ge12$.  \statuslean{} \lean{FiniteCheck4096}.

\textbf{Interval scan.}  Under CPython and \lean{mpmath} interval arithmetic, a nonintegral
solution has $x>20$.  \statuscomp

\textbf{Directed-rounding scan.}  Under the C source, GCC code generation, OpenMP scheduling,
the MPFR and GMP shared libraries, the x86-64 arithmetic environment, and the operating
system, a nonintegral solution has $x>27$.  \statuscomp
\end{resultbox}

MPFR's correct-rounding contract makes the numerical reasoning crisp, and the interval scan
is reproducible from a short script, but neither removes the software assumptions listed
above.  Nothing in Section~\ref{sec:conditional} or Section~\ref{sec:program} that is tagged
\statuslean{} depends on either scan; wherever a numeric threshold is used inside a
kernel-verified statement, that threshold is $12$.

\paragraph{A drafted extension of the kernel region.}
An attempt to raise the kernel-checked region from $2^x<4096$ to $2^x<16384$ exists in draft.
It follows the \lean{FiniteCheck4096} pattern, splitting the enlarged table into four chunks
so that no single \lean{decide} call carries the whole certificate.  The certificate rows
themselves verify: each chunk's Boolean recursion evaluates and each row's exact integer
comparison is accepted.  What has not been accepted is the assembly --- the per-chunk length
lemmas that glue the four tables into a single covering statement exceeded the elaborator's
recursion limit, so the module does not compile end to end and no theorem is available for
use.  Until that is resolved the kernel bound is unchanged at $x\ge12$.  \statusdraft

\paragraph{From scan to kernel certificate.}
The repository's own certificate architecture (Section~\ref{sec:corpus}) shows the stronger
trust model: generate rational certificates externally, revalidate them independently, and
replay only exact integer comparisons in Lean by \lean{decide}.  The next step is certificate
extraction, not a larger blind scan.  For each block one records rational lower and upper
approximants to $\thetaq$ together with integer inequalities sufficient to replay every
exclusion, exploiting monotonicity and convexity so that a single certificate covers a run of
consecutive $m$ rather than storing one record per input.  A scalable formal certificate
should share convergent enclosures across ranges, chunk the data modules to bound elaboration
memory --- the failure mode of the $16384$ draft is precisely that the chunk bookkeeping, not
the arithmetic, is the expensive part --- and rely on verified binary powering rather than
giant literals.  The extremal pairs identified by both scans predict the required tier depth.
A practical intermediate target is a generated, kernel-replayed certificate beyond $4096$
(with the existing $16384$ draft as an engineering test), followed by a hierarchical
certificate that approaches the current software frontier $2^{27}$.
\section{Attempts at a full proof and their exact failure points}
\label{sec:attempts}

This section records the principal avenues that were pushed until they broke, together with
the exact place where each one breaks.  The purpose is not merely negative.  Every failure
point names a quantitative or structural feature that a successful proof must improve, and
several of them turn out to be the same feature seen from different sides; the last
subsection gives that common feature its exact name.

\subsection{Collapsing rank two to rank one}
\label{at:rank}

The multiplicative-rank bound is kernel-verified: the exponent vectors of the candidate set
span a subgroup of rank at most two.  Under failure the rank-two possibility is not vague
but known exactly, $\Cand=\langle2,M\rangle$, with $M$ the primitive generator of
Section~\ref{sec:conditional}.  Nothing in the rank machinery distinguishes this free
rank-two monoid from the allowed monoid of powers of two.

\emph{Failure point.}  A rank bound is an upper bound on dimension; it supplies no mechanism
for forcing the dimension down to one.  Excluding the second generator \emph{is} the
original transcendence problem, restated.

\subsection{Rational approximation to a least generator}
\label{at:approx}

Assume failure and let $\beta$ be the least nonintegral solution, with $M=2^{\beta}$ and
$A=3^{\beta}$.  Let $p/q$ approximate $\beta$ and put $\varepsilon=q\beta-p$.  Then
\begin{align}
  M^{q}-2^{p} &= 2^{p}\bigl(e^{\varepsilon\log2}-1\bigr), \label{at:eq-approx-two}\\
  A^{q}-3^{p} &= 3^{p}\bigl(e^{\varepsilon\log3}-1\bigr). \label{at:eq-approx-three}
\end{align}
Both left-hand sides are integers, so a nonzero value has absolute value at least $1$, which
for small $\varepsilon$ gives bounds of the shape
\[
  |\varepsilon|\gtrsim2^{-p}
  \qquad\text{and}\qquad
  |\varepsilon|\gtrsim3^{-p}.
\]
These are exponentially tiny in $q$.  Continued fractions only guarantee infinitely many
approximants with $|\varepsilon|\ll q^{-1}$, and Baker-type lower bounds for ordinary linear
forms in logarithms are polynomial in the relevant heights
\cite{Baker1975,Waldschmidt2000}.  Neither scale conflicts with a lower bound as weak as
$M^{-q}$ or $A^{-q}$.  Multiplying \eqref{at:eq-approx-two} by \eqref{at:eq-approx-three},
or eliminating $\varepsilon$ between them, only recovers $\log_2 3$ to first order; the
higher terms are far too small to create an integer contradiction.

\emph{Failure point.}  The two integer gaps are correlated rather than independent, so the
two exponential inequalities carry essentially one piece of information, not two.

\begin{observation}
\label{at:obs-approx}
Any rational-approximation proof needs more than ``a nonzero integer has absolute value at
least one.''  It needs a divisibility lower bound growing exponentially with $q$, or an
approximation to $\beta$ exponentially better than continued fractions supply in general.
\end{observation}

\subsection{Equidistribution against exponentially shrinking targets}
\label{at:equi}

At paper level the orbit $\fract{k\beta}$ is equidistributed modulo $1$
(Section~\ref{sec:conditional}); the finite-Fourier convergence needed for Weyl's criterion,
covering every fixed mode and every finite trigonometric polynomial, is kernel-verified.
Meanwhile, integrality of $2^{n+t}$ between adjacent integral powers pushes the fractional
part $t$ away from both endpoints by an amount comparable with $2^{-n}$, and the base-$3$
condition imposes a comparable $3^{-n}$ restriction; at height $k$ the admissible target
width is of exponential scale $M^{-k}$, a bound also kernel-verified in the
shrinking-target module.

Equidistribution controls visits to fixed intervals and, with effective discrepancy, to
polynomially shrinking ones along subsequences.  It does not force visits to a summable
family of targets of length $e^{-ck}$: the Borel--Cantelli heuristic predicts only finitely
many such visits for a typical rotation, and numbers with bounded partial quotients show
that equidistributed multiples can avoid exponentially small targets forever.  The exact
conditional monoid is perfectly compatible with that behaviour.

\emph{Failure point.}  The blockwise Weyl theorems are structurally correct but point the
wrong way: they describe how already existing solutions distribute inside each block, not
how likely arbitrary exponents are to hit the integer lattice.  Closing the route needs
Diophantine information specific to the transcendental $\beta$ with $2^{\beta},3^{\beta}\in\N$,
not generic metric theory.

\subsection{Finite differences and additive progressions}
\label{at:diff}

For $f(t)=t^{\thetaq}$ the repository proves the exact mean-value formula for every forward
difference order $r$,
\[
  \Delta_h^{r}f(n)=\fall{\thetaq}{r}\,h^{r}\xi^{\thetaq-r}
  \qquad(n<\xi<n+rh),
\]
so that if all $r+1$ values $f(n),f(n+h),\dots,f(n+rh)$ are integers, the left side is an
integer; under $|\fall{\thetaq}{r}|h^{r}<n^{r-\thetaq}$ its absolute value lies strictly
between $0$ and $1$.  This is a sharp uniform obstruction to candidate arithmetic
progressions, and its rational effective three-term shadow $h^{400}\le n^{83}$ is
kernel-verified together with the arbitrary-order form.

The difficulty is combinatorial rather than analytic.  Conditional failure produces only
$O(\log X)$ candidates per multiplicative block, far below the positive-density regime of
Roth-type theorems, and a rank-two multiplicative monoid need not contain long additive
progressions at all: many points in a short interval need not be equally spaced.  A
three-term progression of candidates would solve the $S$-unit equation
\[
  2^{n_1}M^{k_1}+2^{n_3}M^{k_3}=2^{n_2+1}M^{k_2},
\]
and general $S$-unit theory yields finiteness statements, not the required nonexistence.

\emph{Failure point.}  Sparse exact structure and sharp local curvature are mutually
consistent.  The arbitrary-node determinant of Section~\ref{sec:arbnode} removes the
equal-spacing hypothesis, but its packing exponent is still too weak to contradict the exact
conditional count.

\subsection{Radical fields and conjugates}
\label{at:radical}

The normalized radical $E=w^{\thetaq}$ is algebraic with exact binomial minimal polynomial
$X^{d}-c/3^{a}$, and its conjugates are $\zeta_d^{j}E$; degree, norm, and denominator
support are all kernel-verified (Section~\ref{sec:radical}).  What this does not supply is a
way to transport the real analytic identity $\log E=\thetaq\log w$ to the other conjugates.

\emph{Failure point.}  The chosen positive real logarithm is not Galois-equivariant: the
remaining roots differ by complex logarithmic periods, and taking norms merely reproduces
$E^{d}=c/3^{a}$.  Ramification is equally inconclusive, since the equation relates the
constrained primes through a quotient of logarithms rather than through a field embedding;
standard Kummer theory sees the radical equation but not the special value
$\thetaq=\log_2 3$.

\begin{observation}
\label{at:obs-radical}
A radical-field proof needs a genuinely mixed invariant: something simultaneously sensitive
to the algebraic conjugates of $E$ and to the logarithmic relation among $2,3,w,E$.  At
present the algebraic and the real-analytic arguments meet only at the four-logarithm
determinant.
\end{observation}

\subsection{The compact $S$-integral orbit}
\label{at:orbit}

From the least generator define
\[
  U_k:=\frac{M^{k}}{2^{\lfloor k\beta\rfloor}}=2^{\fract{k\beta}}\in\Z[1/2],
  \qquad
  V_k:=\frac{A^{k}}{3^{\lfloor k\beta\rfloor}}=3^{\fract{k\beta}}\in\Z[1/3].
\]
The points $(U_k,V_k)$ lie on the compact analytic arc
$\mathcal A=\{(2^{t},3^{t}):0\le t<1\}$ and are dense there because $\beta$ is irrational.
This is the compact-orbit reformulation of Section~\ref{sec:reform}, and it invites
o-minimal and $S$-integral point counting; Butler relates special cases of Wilkie-type
counting to real three- and four-exponentials statements \cite{Butler2017}.

\emph{Failure point.}  The counting is not strong enough by an enormous margin.  The
logarithmic heights of the orbit points grow linearly in $k$, so the arc carries only about
$\log H$ orbit points of height at most $H$, far below the subpolynomial count $H^{\varepsilon}$
that Pila--Wilkie style theorems already tolerate outside algebraic parts.  What would
suffice is a denominator-sensitive refinement counting only points whose coordinate
denominators are supported exactly on $\{2\}$ and on $\{3\}$ with synchronized exponents; no
theorem of that strength appears in the inspected corpus.

\subsection{Forcing an auxiliary base}
\label{at:auxbase}

If one could exhibit an integer $b>1$ multiplicatively independent of $2$ and $3$ with
$b^{x}$ algebraic --- or rational with suitably controlled denominator --- then the six
exponentials theorem, or the repository's rational-third-base theorem, would force
$x\in\Q$ and hence $x\in\Z$.  The five exponentials theorem \cite{Waldschmidt2000,Roy1992}
similarly makes $e^{\gamma x}$ transcendental for every nonzero algebraic $\gamma$, given
the four algebraic corner exponentials, and the kernel-verified $3$-smooth classification is
the formalized integer-level shadow of the six exponentials theorem.

\emph{Failure point.}  The hypotheses provide no such $b$.  The natural auxiliary quantities
$p^{x}=e^{x\log p}$ have transcendental coefficient $\log p$, so the five exponentials
theorem does not convert a divisor of $M$ or $A$ into a contradiction; and the six
exponentials theorem controls $b^{x}$, not the distinct question of whether $b$ divides the
integer $2^{x}$.  The obvious constructions $b=2^{u}3^{v}$, and any construction from $M$ and
$A$, collapse to the same rank-two logarithmic span.  An auxiliary base must come from an
arithmetic operation whose exponential behaviour stays controlled while its logarithm leaves
that span --- an algebraic independence problem in disguise.

\subsection{Exploiting the primitive output pair}
\label{at:primitive}

The primitive pair $(M,A)$ is canonical and rigid: kernel-verified theorems show that it
admits no matched depth, no common perfect power, and no nontrivial affine decomposition,
and that the odd core and its coordinates are unique.  Every structural handle one might
hope to grip has therefore been shown to be absent by design rather than by accident.

\emph{Failure point.}  After all of that normalization, what remains is exactly the single
relation $\log M\,\log3=\log A\,\log2$, a product of logarithms lying outside the scope of
linear-forms machinery.

\subsection{The linear-form-in-log-products formulation of the wall}
\label{at:logproduct}

The obstruction can be given one further exact name.  A counterexample is a pair of integers
$m,A\ge2$ with $\log m\,\log3=\log A\,\log2$.  Expanding both integers over their prime
factorizations $m=\prod_pp^{m_p}$ and $A=\prod_pp^{a_p}$, failure of Conjecture~\ref{conj:AE}
is equivalent to the nontrivial vanishing of
\begin{equation}
  D\;=\;\sum_{p}\bigl(m_p\log3-a_p\log2\bigr)\log p,
  \label{at:eq-logform}
\end{equation}
an \emph{integer-coefficient linear form in the products of two prime logarithms}
$\{\log p\cdot\log2,\ \log p\cdot\log3\}$.  Baker's theory bounds linear forms in logarithms
with \emph{algebraic} coefficients \cite{Baker1975}; here every coefficient of $\log p$ is
itself a transcendental logarithm, and no lower-bound technology exists for such forms.
Indeed even the $\Q$-linear independence of the family
$\{\log p\cdot\log q:p<q\ \text{prime}\}$ is open --- it follows from, and is essentially a
fragment of, the four exponentials circle of conjectures.  Section~\ref{sec:logproduct}
turns exactly that independence statement into a sufficient condition for the conjecture.

\emph{Failure point.}  Every route examined above --- determinants, product formulas in the
radical field, auxiliary exponentials, equidistribution against shrinking targets, orbit
counting --- terminates at a special case of this one independence statement.  A genuine
proof therefore requires either a fundamentally new lower bound for linear forms in
log-products, or a mechanism exploiting the positivity and integrality of the coefficients
$(m_p,a_p)$ in \eqref{at:eq-logform} that has no analogue in the general four-exponentials
setting.

\section{Consolidated research program}
\label{sec:program}

The two source programs overlapped substantially; the exact structure theory of
Section~\ref{sec:conditional} settles some directions outright and re-prioritizes the rest.
The list below is ordered by how directly each item attacks the quantitative deficit of
Section~\ref{sec:deficit} --- the one missing logarithm between what the determinant
hierarchy delivers and what a contradiction requires.

\subsection{Priority A: $p$-adic amplification of semigroup determinants}
\label{at:prog-padic}

The semigroup determinant $D=\det\bigl(m_i^{a_j}(m_i^{\thetaq})^{b_j}\bigr)$ is a positive
integer, and under the exact monoid every entry has explicit valuations obtained from
$m_i=2^{n_i}M^{k_i}$ and $m_i^{\thetaq}=3^{n_i}A^{k_i}$:
\begin{align*}
  \vpp{2}{\text{entry}_{ij}}&=a_j\bigl(n_i+\vpp{2}{M}k_i\bigr)+b_j\vpp{2}{A}k_i,\\
  \vpp{3}{\text{entry}_{ij}}&=a_j\vpp{3}{M}k_i+b_j\bigl(n_i+\vpp{3}{A}k_i\bigr).
\end{align*}
What this would buy is decisive: archimedean reasoning offers only the universal lower bound
$|D|\ge1$, whereas a systematic $p$-adic factor supplies $|D|\ge p^{L}$ and could erase the
missing logarithm outright.

\begin{question}
\label{at:q-padic}
Can one choose $O(\log X)$ candidate rows and mixed-monomial columns so that the resulting
nonzero determinant satisfies $\log|D|_p^{-1}\gg(\log X)^2$ for $p=2$ or $p=3$, while the
archimedean estimate remains $\log|D|_\infty=o\bigl((\log X)^2\bigr)$?
\end{question}

The difficulty is that a naive row or column factor is lost because the exponent set contains
$(0,0)$; only more subtle min-plus determinant estimates, applied after suitable row and
column finite-difference operations, have a chance.  The exact affine primitivity and the
gcd condition on the normalized degrees should be fed into the estimate rather than
discarded.

\subsection{Priority B: exploit the full $(n,k)$ grid, not sorted magnitudes}
\label{at:prog-grid}

The arbitrary-node argument forgets the canonical coordinates and retains only the sorted
integers $m_i$.  Restoring them exposes rectangular, often Kronecker-like, bivariate
exponential matrices, which is exactly the balanced $2\times2$ setting where ordinary six
exponentials estimates stop; integrality, positivity, and the one-depth-zero normalization
may create a margin absent from the general four-exponentials problem.  The right experiment
is not another large generic determinant but an optimized minor selected from triangular
coordinate regions $n+k\beta\le T$ and $a+b\thetaq\le L$, with exact accounting of
archimedean size, zero estimates, and $2$- and $3$-adic valuations.  The hard part is that
the optimization is genuinely two-parameter and the useful shapes are not predicted by any
current heuristic, so a search must precede a theorem.

The continuation in Section~\ref{sec:matrix-continuation} turns this priority into an exact
fewnomial target.  On $R$ distinct counterexample points, every vanishing auxiliary
polynomial has at least $R+1$ monomials; therefore one must either construct a relation with
one fewer monomial than ordinary interpolation would use, or exploit additional arithmetic
structure not measured by support.  On the $2N$ grid, the BK-degree calibration shows a
factor-four gap between the extremal interpolant and the general Matrix Coefficient target.
The boundary-dilation defects are the most concrete finite compatibility equations currently
available for attacking that gap.

\subsection{Priority C: denominator-sensitive counting on the compact arc}
\label{at:prog-counting}

The compact-orbit reformulation asks for infinitely many points on $y=x^{\thetaq}$ in
$[1,2)\times[1,3)$ whose first denominator is a power of $2$, whose second denominator is a
power of $3$, and whose denominator exponents are synchronized by the same integer
$\lfloor k\beta\rfloor$.  A theorem asserting that a compact transcendental arc with nonzero
curvature carries $o(H)$ points of logarithmic height at most $H$ with denominators
supported on two prescribed disjoint prime sets and a common exponent parameter would close
the case, since the failed-conjecture orbit has $\asymp H$ such points.  What makes it hard
is that standard rational-point counting is sensitive to denominator \emph{size} but not to
denominator \emph{support} or synchronization; Butler's connection between Wilkie-type
statements and real exponential conjectures \cite{Butler2017} suggests this direction sits
on the transcendence frontier rather than being a routine application.

\subsection{Priority D: mixed archimedean--radical invariants}
\label{at:prog-radical}

The normalized radical $E=w^{\thetaq}$ has exact degree $d$ and binomial minimal polynomial,
so the Kummer field $\Q(E,\zeta_d)$ is completely explicit.  The questions worth pursuing
are whether the logarithmic relation $\log E\log2-\log w\log3=0$ forces an impossible
regulator or height relation there; whether simultaneous normalization at both outputs
constrains the two normalized degrees through field intersections or discriminants beyond
the known gcd condition; whether all complex branches $\log(\zeta_d^{j}E)=\log E+2\pi ij/d$
can populate a multi-logarithm determinant with algebraic entries of controlled height; and
whether a $p$-adic logarithm branch is compatible with the real identity after raising to
the normalized degree.  A warning is essential, and it is the same one that closed
Section~\ref{at:radical}: merely taking norms or conjugates of $X^{d}-c/3^{a}$ reproduces
known relations, so the new invariant must mix conjugation with the special logarithmic
direction.

\subsection{Priority E: discharge and sharpen the determinant hypotheses}
\label{at:prog-formalize}

The algebraic and combinatorial determinant hierarchy has now been formalized and accepted:
\begin{center}
\lean{ArbitraryNodeDeterminant}, \lean{ExponentialPolynomialZeros},\\
\lean{OrderedChamberPositivity}, \lean{MixedExponentSemigroup},\\
\lean{SemigroupDeterminant}, \lean{FallingFactorialBound},\\
\lean{SemigroupWeightBound}, \lean{ArbitraryNodeRepulsion}, \lean{SemigroupGapBound}.
\end{center}
They are all on the aggregate import surface.  \statuslean{}  The two
quantitative endpoints retain their analytic inputs as explicit hypotheses,
\lean{RpowDividedDifferenceMeanValue} and \lean{NewtonFactorization}; those propositions are
not asserted by the corpus.  The priority is therefore no longer to translate the algebra,
but to prove these named hypotheses in the required uniform form, or to replace them by a
stronger factorization whose archimedean loss is small enough to interact with only
$O(\log X)$ nodes.  The exact tropical initial-form criterion of
\lean{TropicalInitialForm} is a useful prerequisite on the nonarchimedean side, but it does
not supply the cancellation lower bound that this program needs.

\subsection{Priority F: force a third algebraic exponential}
\label{at:prog-thirdbase}

The denominator-controlled theorem in \lean{RationalThirdBase} is a precise gateway: for an
auxiliary integer $q$ containing a prime outside $\{2,3\}$ it suffices, in several forms, to
prove that $q^{x}$ is rational with denominator supported by the known outputs, after which
monomial injectivity and the three-base determinant force $x\in\Q$ and hence $x\in\Z$.  What
this would buy is a complete proof.  The obstacle is the one identified in
Section~\ref{at:auxbase}: straight monomials cannot introduce a new prime direction, so
candidates must come from norms, resultants, or algebraic combinations arising in the
primitive radical field, whose $x$th powers would then have to simplify.

\subsection{Priority G: five and sharp six exponentials}
\label{at:prog-fiveexp}

Formalizing the five exponentials theorem, or Roy's sharp six exponentials theorem
\cite{Roy1992}, would add genuinely new restrictions --- for instance transcendence of
$e^{\gamma x}$ for algebraic $\gamma\ne0$ --- and would extend the rational-third-base
result from a controlled denominator to an unrestricted rational output.  The effort is
comparable to the completed Gelfond--Schneider port, which is encouraging; the risk is that
it be undertaken without a specific downstream lemma in mind, since as
Section~\ref{at:auxbase} shows these theorems constrain $b^{x}$ and not the divisibility
facts the conjecture actually needs.

\subsection{Priority H: larger kernel-replayed regions and effective measures}
\label{at:prog-finite}

The zero-free region $2^{x}<4096$ is kernel-verified, while the current software frontier is
$2^{27}$ (\statuscomp{}) in Section~\ref{sec:finite27}.  A source draft aiming at $2^{14}$
has verified arithmetic rows but no accepted assembled theorem, so the next target is a
generated, kernel-replayed certificate beyond $4096$, with $16384$ a useful engineering
milestone rather than an established bound.  Every descent constant
improves by bookkeeping once the region grows, so the payoff is broad but shallow.  The
bottleneck is a compact kernel-replayable proof object: adaptive convergent enclosures,
chunked certificate modules, and tiers shared along ranges.  In the same vein, effective
irrationality measures for $\thetaq$ from P\'ade or hypergeometric methods would replace
per-value certificate tiers with a uniform polynomial criterion and likely extend verified
regions by orders of magnitude.  Their limitation is sharp: they say nothing directly about
the generator $\beta=\log_2M$ for unknown $M$, which would require an effective measure
derived from the log-product relation \eqref{at:eq-logform} and therefore lies beyond
ordinary linear forms in logarithms.

\subsection{Priority I: computational reconnaissance with exact certificates}
\label{at:prog-compute}

Finite scanning is evidence, not a path to infinity, but computation remains valuable in
three focused roles: searching candidate determinant shapes for large systematic $2$- or
$3$-adic valuations, feeding Priority~A and Priority~B; enumerating small primitive radical
patterns $(w,d,a,c)$ and recording local obstructions with exact rational inequalities
rather than floating-point rejection; and rejecting candidates early using the primitive-pair
conditions --- no matched depth, no common perfect power, no matched affine decomposition ---
while recording near misses keyed by odd core, denominator, and rational-power index.  The
goal must be pattern discovery and theorem generation; the difficulty is the discipline not
to spend the budget on extending a decimal barrier that proves nothing new.

\subsection{Priority J: the colossally abundant interface}
\label{at:prog-ca}

Formalizing the original application --- that the rational-power form of the conjecture
implies that quotients of consecutive colossally abundant numbers are prime \cite{AE1944} ---
would connect the corpus to its historical source and give the whole development a visible
consequence.  It is logically downstream of everything above, and the work is combinatorial
and definitional rather than transcendental, so it should follow the core structure modules
rather than compete with them.
\section{Lean formalization blueprint}
\label{sec:blueprint}

The new arguments of Sections~\ref{sec:arbnode}, \ref{sec:semigroup}, \ref{sec:deficit},
\ref{sec:reform} and \ref{sec:logproduct} were written so that each analytic step is a
reusable lemma rather than a step inside one monolithic proof.  The intended module layout
sits under \lean{ExponentialIdentities/TwoBaseIntegerExponent/} alongside the existing
kernel-verified corpus of Section~\ref{sec:corpus}.

Since the first draft of this blueprint was written, all of the modules described below have
been accepted by the Lean kernel and placed on the aggregate import surface
\lean{ExponentialIdentities.lean}.  The formal determinant chain deliberately exposes its
two remaining analytic inputs --- \lean{RpowDividedDifferenceMeanValue} and
\lean{NewtonFactorization} --- as hypotheses of the quantitative endpoint theorems.  Thus
the algebraic identities and every deduction from those hypotheses are kernel-verified, but
the corpus does not assert the hypotheses themselves.  The descriptions below record the
implemented architecture rather than a pending plan.

\subsection{\lean{ArbitraryNodeDeterminant}}
\label{bp:arbnode-module}

\begin{statusbox}{\statuslean{} --- \lean{ArbitraryNodeDeterminant.lean}, with the analytic
mean-value input exposed in the later module \lean{ArbitraryNodeRepulsion}.}
\end{statusbox}

This module formalizes the ordinary-power determinant and the arbitrary-node repulsion
theorem of Section~\ref{sec:arbnode}.  Its dependencies are:

\begin{enumerate}
\item a determinant/divided-difference identity for the matrix whose columns are
$1,x,\ldots,x^{r-1},f(x)$ evaluated at the nodes, so that the determinant factors as the
Vandermonde product times the $r$th divided difference of $f$;
\item a generalized mean value theorem for divided differences at distinct ordered nodes,
which converts that divided difference into a derivative value at an interior point;
\item the existing derivative tower \lean{rpowDerivativeFamily} and the falling-power
coefficient \lean{fallingRpowCoeff}, both already available in \lean{HigherDifference};
\item the integer-entry determinant lower bound: a matrix all of whose entries are integers
has integer determinant, so a nonzero determinant has absolute value at least one.
\end{enumerate}

The accepted module splits exactly along those four lines.  It defines \lean{nodeMatrix} and its
integer shadow \lean{nodeIntMatrix}, proves that candidate nodes make every entry an
integer, derives \lean{one_le_abs_det_nodeMatrix} from that, and then isolates the algebra
of the final step in a hypothesis-carrying form: the analytic factorization
$\det=V\cdot dd$ and the divided-difference bound $|dd|\le B$ are inputs, and the conclusion
is the repulsion inequality.  A target statement for the fully analytic version is:

\begin{lstlisting}[style=leanstyle,caption={Schematic fully analytic arbitrary-node signature.}]
theorem vandermonde_lower_bound_of_twoBaseNaturalCandidates
    {r : N} (hr : 2 <= r)
    (m : Fin (r + 1) -> N)
    (hstrict : StrictMono m)
    (hcand : forall i, TwoBaseNaturalCandidate (m i)) :
    (r.factorial : R) / |fallingRpowCoeff logThreeDivLogTwo r| *
        (m 0 : R) ^ ((r : R) - logThreeDivLogTwo)
      <= vandermondeProduct r m
\end{lstlisting}

The intermediate form accepted by the kernel is the hypothesis-carrying one, which is
worth stating separately because it is reusable for any $f$ whose divided differences can be
bounded:

\begin{lstlisting}[style=leanstyle,caption={Accepted packaged repulsion step.}]
theorem inv_le_vandermondeProduct_of_dividedDifference (r : N)
    (m : Fin (r + 1) -> N) (hstrict : StrictMono m)
    (hcand : forall i, TwoBaseNaturalCandidate (m i)) {dd B : R}
    (hB : 0 < B)
    (hdet : (nodeMatrix r m).det = vandermondeProduct r m * dd)
    (hdd_bound : |dd| <= B) (hne : (nodeMatrix r m).det <> 0) :
    1 / B <= vandermondeProduct r m
\end{lstlisting}

Working with a \lean{Fin}-indexed strictly monotone family and an explicit pair product is
easier than working with a \lean{Finset} of nodes; the \lean{Finset} form is recovered at the
end through the monotone enumeration \lean{Finset.orderEmbOfFin}.

\subsection{\lean{GeneralizedVandermonde}}
\label{bp:gv-module}

\begin{statusbox}{\statuslean{} --- the zero-count/nonvanishing engine is
accepted in \lean{ExponentialPolynomialZeros}, and ordered-chamber positivity is accepted in
\lean{OrderedChamberPositivity}.  Appendix~\ref{app:gvproof} gives the paper proof and its
Wronskian interpretation.}
\end{statusbox}

The positivity theorem for generalized Vandermonde determinants with arbitrary real
exponents is useful well beyond this conjecture, and two formalization routes were
available.

\begin{description}
\item[Wronskian route.] Prove the Wronskian determinant formula of
Appendix~\ref{app:gvproof} directly, then invoke or develop the criterion that positive
initial Wronskians produce an extended complete Chebyshev system and hence positive ordered
evaluation determinants.  This needs a small Chebyshev-system API that mathlib does not yet
have.

\item[Zero-count route.] After the substitution $x=e^t$, prove by induction that a nonzero
exponential polynomial with $r+1$ distinct real exponents has at most $r$ real zeros.  Deduce
nonvanishing of the evaluation determinant, and if the sign is wanted, obtain it by
continuity and a coalescing-node limit.
\end{description}

The decisive observation for the present application is that only \emph{nonvanishing} is
needed downstream.  The determinant in question has integer entries, and a nonzero integer
already has absolute value at least one; the sign never enters the repulsion inequality.
That removes the coalescing-limit argument from the critical path.  The formalization first
took the zero-count route in \lean{ExponentialPolynomialZeros}; the independent accepted
module \lean{OrderedChamberPositivity} now supplies the strict sign as well.

\subsection{\lean{ExponentialPolynomialZeros}}
\label{bp:expzeros-module}

\begin{statusbox}{\statuslean{} --- \lean{ExponentialPolynomialZeros.lean}.}
\end{statusbox}

This module carries the zero-count induction.  Write
$g(t)=\sum_{j\le n}c_je^{\lambda_jt}$ with $\lambda_0<\cdots<\lambda_n$.  Normalizing by
$e^{-\lambda_0t}$ preserves the zero set; differentiating kills the constant term and lowers
the number of exponents by one; Rolle's theorem between consecutive zeros supplies the zeros
of the derivative.  Induction on $n$ then forces every coefficient to vanish.

Two Lean engineering difficulties dominate the effort, and neither is visible in the paper
proof.

First, the Rolle step must produce $n$ \emph{distinct and ordered} interior points from
$n+1$ ordered zeros.  Mathlib's \lean{exists_hasDerivAt_eq_zero} yields one interior point
per gap, but nothing about how the points from different gaps compare.  The implementation chooses
them with \lean{choose}, then proves strict monotonicity by chaining
$y_i<x_{i+1}\le x_j<y_j$ through the monotonicity of the node family.  When the hypothesis
arrives as a \lean{Finset} of at least $n+1$ reals rather than an indexed family, the
monotone enumeration \lean{Finset.orderEmbOfFin} converts it, after cutting the set down to
exactly $n+1$ elements with \lean{Finset.exists_subset_card_eq}.

Second, the induction reindexes.  The derivative of an $(m+1)$-term exponential polynomial
is an $m$-term one whose exponents are $\lambda_{j+1}-\lambda_0$ and whose coefficients are
$c_{j+1}(\lambda_{j+1}-\lambda_0)$, so every statement has to be transported along
\lean{Fin.succ}, and the induction has to be set up with the order quantified first so that
\lean{induction} may generalize over exponents, coefficients and nodes simultaneously.  In
the accepted module this is done through a private auxiliary theorem of the shape
\lean{forall (n : N) (lam c : Fin (n + 1) -> R), ...}, from which the user-facing statements
are one-line consequences.

\begin{lstlisting}[style=leanstyle,caption={Zero-count theorem and its determinant form.}]
theorem card_le_of_expPoly_eq_zero {n : N} (lam c : Fin (n + 1) -> R)
    (hlam : StrictMono lam) (hc : exists j, c j <> 0) (S : Finset R)
    (hzero : forall t, t in S -> expPoly lam c t = 0) :
    S.card <= n

theorem det_rpowVandermonde_ne_zero {n : N} (x lam : Fin (n + 1) -> R)
    (hx : StrictMono x) (hx0 : 0 < x 0) (hlam : StrictMono lam) :
    (rpowVandermonde x lam).det <> 0
\end{lstlisting}

A vanishing determinant yields a nonzero kernel vector, via
\lean{Matrix.exists_mulVec_eq_zero_iff}, and it is a real combination of
the functions $x\mapsto x^{\lambda_j}$ vanishing at $n+1$ distinct positive nodes.  So the
determinant statement follows from the zero-count statement.

\subsection{\lean{MixedExponentSemigroup} and \lean{SemigroupDeterminant}}
\label{bp:semigroup-module}

\begin{statusbox}{\statuslean{} --- \lean{MixedExponentSemigroup.lean} and
\lean{SemigroupDeterminant.lean}.}
\end{statusbox}

The pair-indexed exponent is the whole of the definitional layer:

\begin{lstlisting}[style=leanstyle]
noncomputable def mixedExponent (ab : N x N) : R :=
  (ab.1 : R) + (ab.2 : R) * logThreeDivLogTwo
\end{lstlisting}

and the module proves, in this order:

\begin{itemize}
\item injectivity of \lean{mixedExponent}, from irrationality of $\thetaq$;
\item integrality of $m^{a+b\thetaq}=m^a\,(m^{\thetaq})^b$ for every candidate $m$, packaged
as a natural-number witness \lean{candidateRpowNat};
\item closure under addition, so that the image is an \lean{AddSubmonoid R};
\item the elementary comparisons $a+b\le a+b\thetaq\le a+2b$ and
$a+b\thetaq\le 3q$ for $a,b\le q$, which are what the counting arguments actually consume.
\end{itemize}

\lean{SemigroupDeterminant} then builds the matrix $m_i^{\lambda_j}$ for candidate nodes and
mixed exponents, shows every entry is a natural number, and concludes integrality and the
lower bound $|\det|\ge1$ whenever the determinant is nonzero --- the nonvanishing being
imported from \lean{ExponentialPolynomialZeros}.

\begin{lstlisting}[style=leanstyle,caption={Semigroup determinant integrality.}]
theorem one_le_abs_det_semigroupMatrix (r : N) (m : Fin (r + 1) -> N)
    (hcand : forall i, TwoBaseNaturalCandidate (m i))
    (ab : Fin (r + 1) -> N x N)
    (hne : (semigroupMatrix r m ab).det <> 0) :
    1 <= |(semigroupMatrix r m ab).det|
\end{lstlisting}

It is deliberately unnecessary to define the globally sorted semigroup here.  The
determinant estimate accepts an arbitrary finite family of pairs together with a proof that
their mixed exponents are strictly increasing; the sorted enumeration is a convenience of
the paper proof, not a requirement of the formal one.

\subsection{\lean{FallingFactorialBound}}
\label{bp:falling-module}

\begin{statusbox}{\statuslean{} --- \lean{FallingFactorialBound.lean}, on the aggregate
import surface.}
\end{statusbox}

This module carries the two elementary estimates behind every Newton-row bound of
Section~\ref{sec:semigroup}: the falling-coefficient bound
$|\fall{\lambda}{i}|/i!\le2^{\lambda+i}$ for $\lambda\ge0$, and the strip bound
$x^{\lambda-i}\le2^{\lambda}N^{\lambda-i}$ for $N\le x\le2N$, $N\ge1$.  The proof of the
first compares each factor $|\lambda-k|$ with $m+1+k$ where $m=\lfloor\lambda\rfloor$,
identifies $\prod_{k<i}(m+1+k)$ with $i!\binom{m+i}{i}$, and bounds the binomial coefficient
by the row sum $2^{m+i}$.  Their product is the single inequality consumed downstream:

\begin{lstlisting}[style=leanstyle,caption={Newton-row entry bound.}]
theorem abs_fallingRpowCoeff_div_factorial_mul_rpow_le
    {N x lam : R} (hN : 1 <= N) (h1 : N <= x) (h2 : x <= 2 * N)
    (hlam : 0 <= lam) (i : N) :
    |fallingRpowCoeff lam i| / (i.factorial : R) * x ^ (lam - (i : R))
      <= (2 : R) ^ (2 * lam + (i : R)) * N ^ (lam - (i : R))
\end{lstlisting}

The determinant estimate then follows from \lean{Matrix.det_apply} or, better, from a
generic Leibniz-bound lemma worth adding once:
\begin{quote}
if every entry in row $i$, column $j$ is bounded by a separable quantity $A_iB_j$, then the
determinant is at most $(r+1)!\prod_iA_i\prod_jB_j$.
\end{quote}

\subsection{\lean{SemigroupWeightBound}}
\label{bp:weight-module}

\begin{statusbox}{\statuslean{} --- \lean{SemigroupWeightBound.lean}.}
\end{statusbox}

For the explicit logarithmic-square corollary the full asymptotic count of the ordered
semigroup is optional.  What is needed is the square-box construction of
Section~\ref{sec:semigroup} in \emph{existence form}, which sidesteps real-valued order
statistics entirely:

\begin{itemize}
\item the $(q+1)^2$ pairs in $[0,q]^2$ have pairwise distinct mixed exponents;
\item every such exponent is at most $3q$, using $1<\thetaq<2$;
\item taking $q=\lfloor\sqrt r\rfloor$ makes $r+1$ pairs fit, with total weight at most
$6(r+1)\sqrt r$;
\item dividing by $R_r=r(r+1)/2$ gives the deficit ratio $\Sigma_r/R_r\le12/\sqrt r$.
\end{itemize}

The module realizes the injective family explicitly by base-$(q+1)$ digits,
$j\mapsto(j\bmod(q+1),\;j\div(q+1))$, so no sorting is performed and no canonical ordered
semigroup is ever constructed.  The statements are:

\begin{lstlisting}[style=leanstyle,caption={Square-box weight bound, existence form.}]
theorem exists_mixedExponent_family (r : N) :
    exists ab : Fin (r + 1) -> N x N, Function.Injective ab and
      (forall j, mixedExponent (ab j) <= 6 * Real.sqrt (r : R)) and
      (sum j, mixedExponent (ab j)) <= 6 * ((r : R) + 1) * Real.sqrt (r : R)
\end{lstlisting}

The module also contains the abstract counting bridge that turns any span lower bound into a
cardinality bound, and which serves the ordinary-power corollary of Section~\ref{sec:arbnode}
equally well: if $S\subseteq[N,N+H]$ and every $r+1$ elements of $S$ contain two points at
distance at least $L>0$, then $\#S\le r+rH/L$.  The proof partitions $[N,N+H]$ into
$\lfloor H/L\rfloor+1$ half-open blocks of length $L$, observes that each block meets $S$ in
at most $r$ points, and applies pigeonhole.

\subsection{\lean{SemigroupDyadicBound}}
\label{bp:dyadic-module}

\begin{statusbox}{\statuslean{} --- implemented in \lean{SemigroupGapBound}; the final
cardinality theorem is conditional on the explicit \lean{NewtonFactorization} hypothesis.}
\end{statusbox}

This layer combines the span theorem with the consecutive-window summation and produces the
explicit $O((\log X)^2)$ dyadic bound.  The accepted endpoint is
\lean{card_le_ten_thousand_max_log_sq}; schematically it states:

\begin{lstlisting}[style=leanstyle,caption={Schematic explicit dyadic bound.}]
theorem card_twoBaseNaturalCandidates_Icc_le_log_sq (X : N) (hX : 1 <= X) :
  ((Finset.Icc X (2 * X)).filter TwoBaseNaturalCandidate).card
    <= Nat.ceil (10000 * max 1 (Real.log X) ^ 2)
\end{lstlisting}

A cleaner all-natural statement may replace the real logarithm by the binary ceiling
logarithm, obtaining a larger but much easier constant times \lean{Nat.clog 2 X ^ 2}.  Since
the repository already has logarithmic-square counting infrastructure in
\lean{MultiplicativeRank}, its finite-set definitions and cardinality lemmas can be reused
while keeping the proof mechanism entirely independent --- which is the point of proving the
same bound twice.

\subsection{\lean{BlockGrowthReformulations}}
\label{bp:blockgrowth-module}

\begin{statusbox}{\statuslean{} --- \lean{BlockGrowthReformulations.lean}.}
\end{statusbox}

The exact conditional block formula proved in the corpus makes the dyadic block count
$B(T)=\#\{\text{candidates in }[2^T,2^{T+1})\}$ either identically $1$ or bounded below by a
positive multiple of $T$.  That dichotomy collapses a whole family of apparently weaker
growth hypotheses onto the conjecture itself, and the module states all four equivalences
together so that a future argument may pick whichever is most convenient:

\begin{lstlisting}[style=leanstyle,caption={Block-growth reformulations.}]
theorem alaogluErdosConjecture_block_growth_reformulations :
    (AlaogluErdosConjecture <-> exists C : N, forall T : N, dyadicBlockCount T <= C) and
    (AlaogluErdosConjecture <->
       (fun T : N => (dyadicBlockCount T : R)) =o[atTop] fun T : N => (T : R)) and
    (AlaogluErdosConjecture <->
       forall eps : R, 0 < eps ->
         frequently T in atTop, (dyadicBlockCount T : R) < eps * (T : R)) and
    (AlaogluErdosConjecture <-> forall T : N, dyadicBlockCount T = 1)
\end{lstlisting}

The practical value is negative and worth stating plainly: proving mere boundedness, or even
vanishing lower density along a subsequence, is exactly as hard as the conjecture.  No
softening of the counting target buys anything.

\subsection{\lean{CompactOrbit}}
\label{bp:compactorbit-module}

\begin{statusbox}{\statuslean{} --- \lean{CompactOrbit.lean}.}
\end{statusbox}

This module states the conjecture as the nonexistence of a synchronized integral orbit on
the compact arc $\{(2^t,3^t):0\le t<1\}$.  The orbit point is
\[
 \bigl(M^k/2^{\lfloor k\beta\rfloor},\;A^k/3^{\lfloor k\beta\rfloor}\bigr),
\]
both coordinates divided by a power of their own base with the \emph{same} integer exponent
$\lfloor k\beta\rfloor$; that shared exponent is what makes the reformulation nontrivial.

\begin{lstlisting}[style=leanstyle,caption={Compact-orbit reformulation.}]
theorem alaogluErdosConjecture_iff_no_compactOrbit :
    AlaogluErdosConjecture <->
      forall M A : N, forall beta : R,
        1 < M -> 1 < A -> Irrational beta -> 0 < beta ->
          exists k : N, 1 <= k and
            synchronizedOrbitPoint M A beta k not_in compactExponentialArc
\end{lstlisting}

The forward direction routes through the canonical primitive generator, so the produced
$\beta$ is the least noninteger solution and $M,A$ are its normalized outputs.  The converse
needs neither $M,A>1$ nor $\beta>0$: arc membership by itself already pins $M=2^\beta$ and
$A=3^\beta$.  Formalizing that asymmetry was the only delicate part.

\subsection{\lean{LogProductIndependence}}
\label{bp:logprod-module}

\begin{statusbox}{\statuslean{} --- \lean{LogProductIndependence.lean}.  The independence
property is an explicit hypothesis, never asserted; the implication from it is accepted.}
\end{statusbox}

This module records the prime-log-product criterion of Section~\ref{sec:logproduct}.  The
open hypothesis is never asserted; it is a named \lean{Prop} that appears only as an explicit
argument, so the exported theorem is unconditional Lean about a conditional mathematical
statement.

\begin{lstlisting}[style=leanstyle,caption={Prime-log-product criterion.}]
def PrimeLogProductIndependence : Prop :=
  forall (S : Finset N) (c : N -> N -> Z),
    (forall p, p in S -> Nat.Prime p) ->
    (sum p in S, sum q in S, (c p q : R) * (Real.log p * Real.log q)) = 0 ->
    forall p, p in S -> forall q, q in S -> c p q + c q p = 0

theorem alaogluErdosConjecture_of_primeLogProductIndependence
    (hind : PrimeLogProductIndependence) : AlaogluErdosConjecture
\end{lstlisting}

The supporting lemma is the prime-factorization expansion
$\log n=\sum_{p\mid n}v_p(n)\log p$, after which the relation $M^{\log3}=A^{\log2}$ becomes a
vanishing integer combination of the products $\log p\log q$; symmetrized coefficientwise
vanishing then forces $M$ to be a power of $2$ and the exponent to be an integer.

\subsection{Proof-audit discipline}
\label{bp:audit}

The discipline that governs every status label in this report is mechanical, and it is worth
restating because several modules above are new.  Each headline theorem must enter the
aggregate import surface \lean{ExponentialIdentities.lean} and be checked with
\texttt{\#print axioms}, so that the trust boundary of each result is visible rather than
inferred.  A result stays \statusdraft{} or \statusnew{} until a full build succeeds; the
existence of a \lean{sorry}-free source file is not verification, and neither is a successful
build of an earlier revision.  Numerical constants appearing in the tables of this report are
explanatory only --- the formal statements use exact expressions or safe rational constants
such as $1/24$ and $10000$, chosen so that the arithmetic side conditions close by
\lean{norm_num} rather than by a tuned estimate.

\section{Comparison of methods}
\label{sec:comparison}

\subsection{Prime-valuation rank versus interpolation geometry}
\label{bp:rank-vs-geometry}

The module \lean{MultiplicativeRank} proves that any three candidates have linearly dependent
prime-exponent vectors, of rank at most two.  It does so by a three-base, six-exponentials
determinant: if three candidate bases were multiplicatively independent, integrality of all
three $\thetaq$th powers would force $\thetaq$ rational, contradicting its irrationality.  A
finite candidate set therefore lies in a rational plane inside the free vector space on
primes.  Choosing two prime coordinates injectively encodes the set and yields an explicit
logarithmic-square count.

The semigroup determinant of Section~\ref{sec:semigroup} proves a logarithmic-square dyadic
bound by an entirely different mechanism.

\begin{center}
\begin{tabular}{p{0.42\textwidth}p{0.42\textwidth}}
\toprule
Prime-valuation method & Semigroup determinant method \\
\midrule
Uses prime factorization of candidate bases & Uses only integer values on $y=x^{\thetaq}$ \\
Global multiplicative rank at most two & Local total positivity and interpolation \\
Invokes a six-exponentials consequence & Uses elementary determinant integrality once
$\thetaq$ is fixed \\
Produces coordinates in two prime directions & Produces repulsion for arbitrary ordered
nodes \\
Naturally global in magnitude & Naturally local in additive intervals \\
Kernel-verified \statuslean{} & Paper theorem \statusnew{}; hypothesis-carrying Lean endpoint
\statuslean{} \\
\bottomrule
\end{tabular}
\end{center}

Their agreement at logarithmic-square scale is suggestive rather than coincidental.  Both
methods see a two-dimensional structure, but they see different two-dimensional structures:
prime-exponent rank two on one side, mixed-exponent lattice dimension two on the other.
Neither refinement alone has closed the gap, and a future proof may well need to combine the
two dual lattices rather than push either one further in isolation.

\subsection{Exact conditional rank one over two generators}
\label{bp:exact-rank}

Conditional on failure of the conjecture, the candidate monoid is not merely of rank at most
two --- it is exactly free on $2$ and $M$.  In prime-valuation space its points form the
nonnegative lattice generated by the prime-exponent vectors of $2$ and $M$.  In
mixed-exponent evaluation space, the values factor across a two-dimensional grid indexed by
the pair $(n,k)$ with value $2^nM^k$.  This exact grid is far stronger than the
unconditional rank theorem, and yet every presently known consequence remains compatible
with it.  That compatibility is the honest measure of how much room is left.

The strongest likely synthesis is a determinant whose rows use primitive-generator
coordinates and whose columns use mixed exponents.  Such a matrix would be simultaneously:

\begin{itemize}
\item integral, hence subject to the $|\det|\ge1$ lower bound;
\item totally positive after ordering by magnitude, hence nonvanishing for free;
\item explicitly factored into powers of $2,3,M,A$;
\item constrained by affine primitivity and by simultaneous output normalization.
\end{itemize}

No current module combines all four properties.  \lean{MultiplicativeRank} has the first and
third; the semigroup modules have the first and second; \lean{PrimitiveGenerator} and
\lean{OutputNormalization} have the fourth.  Building the module that has all four is the
most concrete open formalization task suggested by this report.

\subsection{Finite differences as confluent determinants}
\label{bp:confluent}

Forward differences are the equally spaced specialization of interpolation, and Wronskians
are its confluent limit as nodes coalesce.  The determinant hierarchy of
Sections~\ref{sec:arbnode} and \ref{sec:semigroup} therefore provides a single conceptual
envelope for four things that the corpus currently treats separately:

\begin{itemize}
\item second-, third-, and arbitrary-order finite differences;
\item arbitrary-node ordinary-power determinants;
\item mixed-monomial generalized Vandermonde determinants;
\item confluent determinants using derivatives or repeated coordinate directions.
\end{itemize}

The corresponding formalization strategy is to build a general
interpolation-determinant layer once and recover \lean{HigherDifference} as a specialization
of it, rather than to keep the two developments parallel.  That would remove duplicated real
analysis and, more usefully, make experimental determinant choices cheap to certify: a new
choice of nodes and exponents would need only a divided-difference bound, with integrality
and nonvanishing supplied by the shared layer.

\section{Conclusions}
\label{sec:conclusion}

The Alaoglu--Erd\H{o}s conjecture (Conjecture~\ref{conj:AE}) remains open.  Nothing in this
report proves it, and no combination of the results collected here proves it.  What has
changed is that the reason can now be stated exactly rather than impressionistically.

The kernel-verified core is broad and its themes are complementary.  Interpolation
determinants settle three independent bases; Gelfond--Schneider excludes algebraic irrational
exponents; exact certificates create a growing zero-free region, currently $2^x<4096$;
unique factorization yields a canonical primitive odd core and, conditionally, an exact
solution monoid; arbitrary-order forward differences force sharp local additive sparsity,
with the rational three-term criterion $h^{400}\le n^{83}$ as its effective form; leastness
gives exact affine primitivity; and unconditional equivalences with fractional-part gaps and
localized radicals identify the entire remaining distance to the conjecture.

Conditional on a single counterexample, the corpus determines the whole picture.  A
counterexample would not produce a messy exceptional set that density or semigroup arguments
might squeeze; it would produce the perfectly rigid free monoid
\[
\Sol=\Nzero+\Nzero\beta,\qquad
\Cand=\{2^nM^k:n,k\in\Nzero\},
\]
whose block counts, cumulative asymptotics, fractional parts, rotation statistics,
rational-power index, radical degree, and simultaneous output normalization are all explicit
and internally consistent with every verified gap and density bound.  The surviving
obstruction is correspondingly concentrated:
\[
 \text{exclude }w^{\log_2 3}\in\Zthree
 \text{ for every primitive odd }w>1.
\]
That single exclusion is equivalent to, and exactly as open as, the original conjecture.

The new determinant hierarchy sharpens the picture quantitatively.  The arbitrary-node
theorem of Section~\ref{sec:arbnode} shows that candidate lattice points cannot cluster
freely: every Vandermonde product over candidate nodes obeys a power-law lower bound.  The
mixed-exponent semigroup of Section~\ref{sec:semigroup} upgrades this to arbitrarily
high-order positive integer determinants with a repulsion exponent tending to one, and yields
a second, conceptually independent logarithmic-square counting theorem.

The most valuable outcome, however, is the quantified failure point of
Section~\ref{sec:deficit}.  Pure archimedean interpolation needs about $(\log X)^2$ nodes to
operate at scale $X$, whereas the exact conditional monoid supplies only about $\log X$ of
them.  The method misses by exactly one factor of $\log X$ --- not by a constant, and not by
an epsilon.  This makes future work sharper by exclusion: improving constants, tightening the
falling-factorial bound, or extending finite searches cannot close a gap of that shape.  The
plausible routes are precisely those capable of supplying one additional logarithm of
determinant strength: $2$- or $3$-adic divisibility used simultaneously with the archimedean
estimate, a synchronized-denominator counting theorem on the compact exponential arc of
Section~\ref{sec:reform}, or a mechanism forcing a third algebraically independent
exponential direction.

Section~\ref{sec:logproduct} isolates a different sufficient condition, and it is worth
separating from the rest because it does not live at the four-exponentials boundary at all.
If the products $\log p\log q$ of prime logarithms are linearly independent over $\Q$, the
conjecture follows by an argument that is elementary once the hypothesis is granted.  That
independence is a consequence of Schanuel's conjecture and is not known unconditionally ---
indeed even the irrationality of $(\log3)^2/(\log2)^2$ is open --- but it is a different
classical conjecture, with a different literature and different partial results.  Having two
inequivalent sufficient conditions, one transcendence-theoretic and one arithmetic, is more
informative than having one.

The near-term program is the one consolidated in Section~\ref{sec:program}: finish the exact
conditional consequences, in particular the functional-analytic extension of blockwise
equidistribution to arbitrary continuous tests and interval indicators; exploit the uniform
higher-difference theorem combinatorially, allowing the order to grow with height;
kernel-replay the $2^{20}$ region as a compact certificate; use the explicit radical field for
a genuinely number-field argument; and pursue auxiliary-function methods tailored to
integrality, where all corner values are positive integers and the primitive pair is
canonical.  Alongside those, the immediate formalization work is the module list of
Section~\ref{sec:blueprint}: bring the nine existing drafts through a full kernel build and
onto the aggregate import surface, then attempt the combined determinant identified in
Section~\ref{sec:comparison}.

A complete proof still requires a new idea.  But the target is now narrow, the conditional
structure is exact, the failure mode of the best available elementary method is measured
rather than guessed, and the surrounding formal infrastructure is strong enough to test the
next idea precisely.
\clearpage

\section{Continuation: executive answer, scope, and trust boundaries}
\label{sec:continuation}

The sections that follow were produced as a continuation of the preceding material and
are merged here unchanged in substance.  They supply a symbolic-dynamical model of a
hypothetical counterexample, an effective replacement for the elementary
continued-fraction estimate, exact all-support valuations of the block determinants, and
the sharp constants that govern how much of the archimedean budget term-wise divisibility
can supply.  Section~\ref{sec:minplus} states the unconditional ceiling that
underlies all of them and records the two constants this part corrects.


\subsection{The requested goal and the outcome}

The requested primary goal was to continue the attached research program and prove the
following statement.

\begin{conjecture}[Alaoglu--Erd\H{o}s]
If $x\in\R$ and $2^x,3^x\in\Z$, then $x\in\Z$.
\end{conjecture}

I did not find an unconditional proof.  In particular, none of the arguments below converts
\[
  \log(2^x)\log3-\log(3^x)\log2=0
\]
into a forbidden linear form in logarithms of algebraic numbers.  The equality is bilinear in
logarithms, and normalizing it gives exactly the unresolved $2\times2$ determinant behind the
four exponentials conjecture.  Baker-type estimates control every \emph{nonzero} linear form
that occurs after fixing integer coefficients; they do not prove that this special bilinear
form cannot vanish.

The continuation nevertheless changes the technical map in useful ways.  It obtains an exact
symbolic-dynamical formulation, replaces an elementary continued-fraction estimate by an
effective polynomial one, computes the full leading $p$-adic content of ordinary block
Vandermondes, and gives sharp continuum formulas for the termwise divisibility of two much
larger determinant constructions.  These calculations do more than say that a method ``seems
too weak'': they quantify the leading constant that is available and isolate the precise
kind of extra cancellation a successful proof would have to produce.

\subsection{The principal new results}

\begin{longtable}{@{}>{\raggedright\arraybackslash}p{0.25\textwidth}>{\raggedright\arraybackslash}p{0.55\textwidth}>{\raggedright\arraybackslash}p{0.13\textwidth}@{}}
\toprule
Result & Content & Status\\
\midrule
\endfirsthead
\toprule
Result & Content & Status\\
\midrule
\endhead
Synchronized Sturmian orbit & A counterexample is equivalent to one irrational mechanical word simultaneously controlling the dyadic and triadic denominator jumps of two compact rational recurrences. & \statusnew\\
\addlinespace
Effective irrationality of $\beta$ & Baker--W\"ustholz/Matveev gives $|\beta-p/q|\ge c(M)q^{-\mu(M)}$ and hence polynomial, rather than exponential, recurrence bounds for continued-fraction denominators. & \statusknown{} plus \statusnew\\
\addlinespace
Ordinary Vandermonde audit & The valuations at every prime in the fixed supports of $M$ and $A$ are computed to leading order.  The product construction has universal ceiling $1-\log2/(3\log6)=0.871049\ldots$. & \statusnew\\
\addlinespace
Gini assignment formula & For the balanced mixed-semigroup determinant, the exact termwise all-prime divisor fraction is $1-\frac38\sum_p\mathfrak G(c_p)+o(1)$ for an explicit norm $\mathfrak G$. & \statusnew\\
\addlinespace
Universal $4/5$ ceiling & The primitive condition forces $\sum_p\mathfrak G(c_p)\ge8/15$.  Thus entrywise/min-plus divisibility can supply at most $80\%$ of the leading block determinant height. & \statusnew\\
\addlinespace
Cross-coprime constants & If $\gcd(MA,6)=1$, the all-prime ceiling is $7/10$, and the places $2$ and $3$ alone contribute exactly $3/10$ in the continuum limit. & \statusnew\\
\addlinespace
Full-grid transport ceiling & On the complete triangular $(n,k)$ grid, any termwise all-prime divisor is bounded by the countermonotone energy $\pi/8$, while the largest determinant term has energy $1/2$; the ratio is $\pi/4$. & \statusnew\\
\addlinespace
HCIZ reformulation & The full-grid archimedean problem is expressed by the Harish--Chandra--Itzykson--Zuber integral, identifying a concrete free-energy inequality as the remaining analytic target. & \statusknown{} plus \statusnew\\
\bottomrule
\end{longtable}

\subsection{Evidence labels used in this part}

This part uses the report-wide convention of Section~\ref{pr:evidence-convention},
with two additions.  \statusrepo{} marks a statement for which a repository
commit or build log reports success but which was not independently replayed by the
author of the continuation; it is repository metadata, not a kernel audit.
\statusopen{} marks a statement proposed as a precise target rather than asserted.
Where the continuation wrote \statusknown{} for standard literature the
report-wide label is used unchanged.

\subsection{What the determinant ceilings do and do not prove}

A ceiling such as $4/5$ has a specific meaning.  Given an integer determinant
$D=\det(E_{ij})$, define at each prime $p$ the minimum valuation of a Leibniz term,
\[
  \tau_p=\min_{\sigma}\sum_i v_p(E_{i,\sigma(i)}).
\]
Every term, and therefore $D$, is divisible by $p^{\tau_p}$.  The product
$\prod_pp^{\tau_p}$ is the \emph{termwise} or \emph{tropical} divisor.  It ignores extra
valuation caused by cancellation among several terms attaining the same minimum.  The $4/5$
theorem says that this automatic divisor cannot occupy more than four-fifths of the leading
archimedean height for the balanced block construction.  It does \emph{not} say that the true
valuation $v_p(D)$ is at most the tropical minimum; the opposite inequality holds, and extra
cancellation may be large.  Thus the result retires only the naïve entrywise-divisibility
version of the method.  It leaves a sharply defined cancellation problem.

\section{Baseline inherited from the unified report}

\subsection{Notation}

Set
\[
  \vartheta=\log_2 3.
\]
The solution and candidate sets are
\[
  \Sol=\{x\ge0:2^x,3^x\in\N\},
  \qquad
  \Cand=\{m\in\N:m^{\vartheta}\in\N\}.
\]
The map $x\mapsto2^x$ is a bijection $\Sol\to\Cand$.  Integer exponents give the obvious
solutions $\Nzero\subseteq\Sol$ and candidates $2^{\Nzero}\subseteq\Cand$.

Throughout the conditional sections, suppose the conjecture fails.  The exact structure
proved in the source report then supplies a least nonintegral solution $\beta>0$ and primitive
outputs
\[
  M=2^\beta\in\N,
  \qquad
  A=3^\beta\in\N.
\]
Every solution and candidate has a unique coordinate representation
\begin{equation}\label{eq:exact-structure}
  \Sol=\{n+k\beta:n,k\in\Nzero\},
  \qquad
  \Cand=\{2^nM^k:n,k\in\Nzero\}.
\end{equation}
The associated output is
\[
  (2^nM^k)^\vartheta=3^nA^k.
\]
We write
\[
  a=v_2(M),\qquad b=v_3(A),\qquad c=v_2(A),\qquad d=v_3(M).
\]
The inherited primitive-output theorem includes the crucial condition
\begin{equation}\label{eq:no-matched-depth}
  \min(a,b)=0.
\end{equation}
It is important not to silently strengthen \eqref{eq:no-matched-depth} to $a=b=0$.

\subsection{Inherited results used here}

The following facts are taken from the attached unified report and its referenced Lean
modules.

\begin{itemize}
\item \statuslean{} A nonintegral solution, if it exists, is transcendental.  The constant
      $\vartheta$ is transcendental, by Gelfond--Schneider.
\item \statuslean{} Failure gives the exact free rank-two monoid structure
      \eqref{eq:exact-structure}, including uniqueness of $(n,k)$.
\item \statuslean{} In a dyadic exponent block $[T,T+1)$, the solutions are exactly
      \[
        T
        \quad\text{and}\quad
        T-\lfloor k\beta\rfloor+k\beta,
        \qquad 1\le k\le\left\lfloor\frac{T+1}{\beta}\right\rfloor.
      \]
      Thus the block count is $1+\lfloor(T+1)/\beta\rfloor$.
\item \statuslean{} The primitive pair obeys \eqref{eq:no-matched-depth}; both outputs cannot
      simultaneously be $\{2,3\}$-smooth.
\item \statuslean{} Mixed monomials $m^u(m^\vartheta)^v$ are integers for
      $(u,v)\in\Nzero^2$, and generalized Vandermonde determinants built from distinct
      candidate nodes and distinct exponents $u+\vartheta v$ are nonzero integers.
\item \statuslean{} The kernel-checked finite exclusion in the source report gives
      $x\ge12$ for a nonintegral solution.  The report separately records a directed-rounding
      software exclusion through $x\le26$.
\end{itemize}

The continuation does not reprove this large inherited corpus.  It restates every inherited
fact at the point where it enters a new argument.

\subsection{Repository reconciliation as of August 21, 2026}

The attached source is internally newer than one bibliographic pin printed near its end: its
body and status matrix already discuss the enlarged algebraic-output and determinant module
surface, although a repository item still names commit
\sha{5d3982ceb12644fdd3499569d9db5f62f9a31dc6}.  During this continuation I checked the
current public repository history.  Commit
\sha{78f80c5d3f6f9afd64ecb72ddb154c85e2036bc1} reports a successful aggregate build with
$67$ imported modules and merges the later algebraic-output work.  Subsequent commits add a
literature compass and a short $p$-adic priority note; the latest snapshot consulted while
writing was \sha{ff9a9f5aa3a3f02599de93fc053e60918f269855}.

\begin{statusbox}{\statusrepo}
The build claims in the preceding paragraph are taken from repository commit messages and
files.  I did not independently run all Lean jobs.  The mathematical arguments newly stated
below are therefore labelled \statusnew, not \statuslean.
\end{statusbox}

\subsection{The exact remaining logarithmic wall}

Writing $M=2^\beta$ and $A=3^\beta$ gives
\[
  \beta=\frac{\log M}{\log2}=\frac{\log A}{\log3}
\]
and hence
\begin{equation}\label{eq:bilinear-wall}
  \log M\,\log3-\log A\,\log2=0.
\end{equation}
The four numbers $\log2,\log3,\log M,\log A$ are logarithms of algebraic numbers.  Equation
\eqref{eq:bilinear-wall} is a vanishing $2\times2$ determinant.  Four exponentials would
exclude it because $2,3$ are multiplicatively independent and $1,\beta$ are linearly
independent over $\Q$, but four exponentials remains open.  The new work below therefore
pursues information that is invisible in the bare real logarithmic determinant: exact floor
coding, effective separation, and valuations of large integer determinants.

\section{A synchronized Sturmian model of a counterexample}
\label{sec:sturmian}

\subsection{Fractional-part normalization}

Write
\[
  \beta=B+\alpha,
  \qquad B=\lfloor\beta\rfloor\in\Nzero,
  \qquad 0<\alpha<1.
\]
Since a nonintegral solution is transcendental, $\alpha$ is irrational.  Define the rational
localized multipliers
\begin{equation}\label{eq:UV}
  U=\frac{M}{2^B}=2^\alpha\in\Q\cap(1,2),
  \qquad
  V=\frac{A}{3^B}=3^\alpha\in\Q\cap(1,3).
\end{equation}
For $k\ge0$, put
\[
  e_k=\lfloor k\alpha\rfloor,
  \qquad
  s_k=e_{k+1}-e_k\in\{0,1\}.
\]
The binary sequence $s=s_0s_1s_2\cdots$ is the lower mechanical word of slope $\alpha$.
Define compact normalized orbits
\begin{equation}\label{eq:normalized-orbits}
  R_k=\frac{U^k}{2^{e_k}}=2^{\fract{k\alpha}}\in[1,2),
  \qquad
  S_k=\frac{V^k}{3^{e_k}}=3^{\fract{k\alpha}}\in[1,3).
\end{equation}
Both orbits use the \emph{same} fractional part and hence the same symbolic itinerary.

\begin{theorem}[Synchronized dyadic--triadic Sturmian recurrence]\label{thm:sturmian-recurrence}
Assume a primitive nonintegral solution $\beta$ exists.  Then the rational sequences
$R_k,S_k$ in \eqref{eq:normalized-orbits} satisfy
\begin{equation}\label{eq:sturmian-recurrence}
  R_0=S_0=1,
  \qquad
  R_{k+1}=\frac{U R_k}{2^{s_k}},
  \qquad
  S_{k+1}=\frac{V S_k}{3^{s_k}},
\end{equation}
where
\begin{equation}\label{eq:sturmian-threshold}
  s_k=1
  \quad\Longleftrightarrow\quad
  \fract{k\alpha}\ge1-\alpha
  \quad\Longleftrightarrow\quad
  UR_k\ge2
  \quad\Longleftrightarrow\quad
  VS_k\ge3.
\end{equation}
Moreover,
\begin{equation}\label{eq:prefix-count}
  \sum_{j=0}^{N-1}s_j=\lfloor N\alpha\rfloor
  \qquad(N\ge1).
\end{equation}
\end{theorem}

\begin{proof}
The recurrence follows by subtracting consecutive floor exponents:
\[
  R_{k+1}
  =\frac{U^{k+1}}{2^{e_{k+1}}}
  =\frac{U U^k}{2^{e_k+s_k}}
  =\frac{UR_k}{2^{s_k}},
\]
and the same calculation with base $3$ gives the formula for $S_k$.  Since
$e_{k+1}=e_k+1$ exactly when $\fract{k\alpha}+\alpha\ge1$, the first equivalence in
\eqref{eq:sturmian-threshold} holds.  The identities
$UR_k=2^{\alpha+\fract{k\alpha}}$ and
$VS_k=3^{\alpha+\fract{k\alpha}}$ give the other two.  Finally,
\eqref{eq:prefix-count} telescopes:
$\sum_{j<N}s_j=e_N-e_0=\lfloor N\alpha\rfloor$.
\end{proof}

\subsection{Sturmian consequences}

The following facts are classical consequences of irrational mechanical coding
\cite{MorseHedlund1940,Lothaire2002}.

\begin{corollary}[Aperiodicity, balance, and minimal complexity]\label{cor:sturmian}
The word $s$ in Theorem~\ref{thm:sturmian-recurrence} is aperiodic.  Any two factors of equal
length contain numbers of $1$'s differing by at most one, and the number of distinct factors
of length $n$ is exactly $n+1$.
\end{corollary}

Aperiodicity has an elementary proof in this setting.  If $s$ were eventually periodic, then
its limiting density of $1$'s would be rational.  Equation \eqref{eq:prefix-count} shows that
this density is $\alpha$, contradicting irrationality.  Balance and factor complexity are the
standard Sturmian characterization.

The point is not merely terminological.  A hypothetical counterexample is no longer described
only by two rational numbers $U,V$ satisfying a logarithmic equality.  It determines one
minimal-complexity aperiodic word, and that same word must be realized by two rational
piecewise-multiplicative systems with multiplicatively independent reset bases $2$ and $3$.
This is a much more rigid synchronized object.

\subsection{Exact reduced denominators}

The primitive factor depths determine how much cancellation occurs in
\eqref{eq:normalized-orbits}.  Write
\[
  M=2^aW,\quad W\text{ odd},
  \qquad
  A=3^bZ,\quad 3\nmid Z.
\]
Because $M$ is not a power of $2$ and $A$ is not a power of $3$, one has $W,Z>1$.

\begin{proposition}[Denominator exponents and common jump word]\label{prop:denominators}
In lowest terms,
\begin{align}
  R_k&=\frac{W^k}{2^{q^{(2)}_k}},
  &q^{(2)}_k&=\lfloor k\beta\rfloor-ak
              =k(B-a)+\lfloor k\alpha\rfloor,\label{eq:q2}\\
  S_k&=\frac{Z^k}{3^{q^{(3)}_k}},
  &q^{(3)}_k&=\lfloor k\beta\rfloor-bk
              =k(B-b)+\lfloor k\alpha\rfloor.\label{eq:q3}
\end{align}
Consequently
\[
  q^{(2)}_{k+1}-q^{(2)}_k=B-a+s_k,
  \qquad
  q^{(3)}_{k+1}-q^{(3)}_k=B-b+s_k.
\]
At least one of the two sequences has baseline jump $B$, because $\min(a,b)=0$.
\end{proposition}

\begin{proof}
Using $\lfloor k\beta\rfloor=kB+\lfloor k\alpha\rfloor$,
\[
 R_k=\frac{M^k}{2^{\lfloor k\beta\rfloor}}
     =\frac{2^{ak}W^k}{2^{\lfloor k\beta\rfloor}}
     =\frac{W^k}{2^{\lfloor k\beta\rfloor-ak}}.
\]
The numerator is odd, so this is reduced.  The base-$3$ calculation is identical.  Taking
successive differences gives the jump formulas, and the last claim is the inherited
primitive condition \eqref{eq:no-matched-depth}.
\end{proof}

Thus the same Sturmian bit $s_k$ records whether \emph{both} reduced denominators take their
large or small permitted jump at time $k$.  One denominator can have a shallower constant
baseline, but the aperiodic fluctuation is synchronized exactly.

\subsection{An exact equivalent exclusion principle}

The source report already contains localization reformulations such as
$2^\alpha\in\Z[1/2]$ and $3^\alpha\in\Z[1/3]$.  The recurrence packages them into a symbolic
statement.

\begin{theorem}[Synchronized Sturmian exclusion is equivalent to Alaoglu--Erd\H{o}s]
\label{thm:sturmian-equivalence}
The Alaoglu--Erd\H{o}s conjecture is equivalent to the nonexistence of data
\[
  \alpha\in(0,1)\setminus\Q,
  \qquad
  U\in\Q\cap(1,2),
  \qquad
  V\in\Q\cap(1,3)
\]
such that, for the lower mechanical word
$s_k=\lfloor(k+1)\alpha\rfloor-\lfloor k\alpha\rfloor$, the recurrences
\[
  R_0=S_0=1,
  \qquad R_{k+1}=UR_k/2^{s_k},
  \qquad S_{k+1}=VS_k/3^{s_k}
\]
remain in $[1,2)\cap\Q$ and $[1,3)\cap\Q$, respectively, for every $k$, with
$U=2^\alpha$ and $V=3^\alpha$.
Equivalently, it suffices to rule out irrational $\alpha$ for which
\[
  2^\alpha\in\Z[1/2],
  \qquad
  3^\alpha\in\Z[1/3].
\]
\end{theorem}

\begin{proof}
A counterexample $\beta=B+\alpha$ gives the data by
Theorem~\ref{thm:sturmian-recurrence}.  Conversely, if such $\alpha,U,V$ exist, choose an
integer $B$ at least as large as the powers of $2$ and $3$ occurring in the respective
reduced denominators of $U$ and $V$.  Then
\[
  2^{B+\alpha}=2^BU\in\N,
  \qquad
  3^{B+\alpha}=3^BV\in\N,
\]
while $B+\alpha\notin\Z$.  The recurrence identities are equivalent to
$R_k=2^{\fract{k\alpha}}$ and $S_k=3^{\fract{k\alpha}}$, so they add an exact dynamic coding
but no hidden hypothesis.
\end{proof}

\subsection{Cobham-type and Mahler-type diagnostics}

Cobham's theorem says that a sequence automatic in two multiplicatively independent integer
bases is ultimately periodic \cite{Cobham1972,AlloucheShallit2003}.  Sturmian words of
irrational slope are not automatic in any integer base.  Theorem~\ref{thm:sturmian-equivalence}
therefore suggests a tempting strategy: show that the dyadic rational recurrence makes $s$
$2$-automatic and that the triadic recurrence makes it $3$-automatic.

That implication is \emph{not} presently available.  Rationality of the multiplier does not
turn the compact rational state $R_k$ or $S_k$ into a finite-state object; the exact reduced
denominators grow without bound.  Cobham's theorem becomes applicable only after proving a
finite-kernel or finite-substitution property, and that is essentially a new arithmetic
statement.  The synchronized formulation is useful because it identifies this missing lemma
cleanly rather than treating ``automaticity'' as a metaphor.

\begin{target}[Finite-kernel extraction]
Show that simultaneous rational realizability of the same irrational mechanical word by
\eqref{eq:sturmian-recurrence} forces a finite $2$-kernel and a finite $3$-kernel, or derive a
weaker pair of incompatible substitutive structures.
\end{target}

A Mahler-function variant is also conceivable: encode the jump word by
$F(z)=\sum_{k\ge0}s_kz^k$.  Aperiodicity and the finite coefficient set already imply, by
classical results of Fatou type, that $F$ is transcendental over $\Q(z)$.  A proof would have
to derive incompatible functional equations from the two rational recurrences.  No such
equations were found here; the recurrence is multiplicative in the state, not self-similar in
the index.

\section{Effective logarithmic separation}
\label{sec:baker}

\subsection{From linear forms in two logarithms to an irrationality measure}

The source report proves an elementary denominator-growth statement: a very good rational
approximation to a hypothetical solution $\beta$ would force a correspondingly large output,
and consecutive continued-fraction denominators obey a single-exponential recurrence.  Since
$\beta$ is itself a ratio of logarithms of fixed integers, Baker's theory gives a substantially
stronger conclusion.

\begin{theorem}[Effective finite irrationality measure for the primitive generator]
\label{thm:baker-beta}
Assume a primitive counterexample $\beta$ exists, and put $M=2^\beta$.  There are effectively
computable constants $c_M>0$ and $\mu_M<\infty$ such that, for every $p\in\Z$ and $q\ge2$,
\begin{equation}\label{eq:beta-irrat-measure}
  \left|\beta-\frac pq\right|\ge c_Mq^{-\mu_M}.
\end{equation}
One may choose $\mu_M$ bounded by an effectively computable absolute constant times
$1+\log M$, hence by an absolute constant times $1+\beta$.  The same conclusion follows from
$\beta=\log A/\log3$.
\end{theorem}

\begin{proof}
Consider the linear form
\[
  \Lambda=q\log M-p\log2.
\]
It is nonzero: otherwise $M^q=2^p$, which would make $\beta=p/q$ rational; a rational solution
is integral by the elementary rational-exponent classification already formalized in the
source corpus.  The Baker--W\"ustholz theorem, or Matveev's explicit lower bound for a
homogeneous linear form in logarithms of algebraic numbers, gives
\[
  \log|\Lambda|\ge-C\,(1+\log M)\log H-C'(1+\log M),
\]
where $H\ge\max(3,|p|,q)$ and $C,C'$ are effective absolute constants after the fixed number
$2$ is absorbed.  For approximations with $|p|\le(\beta+1)q$, one has $H\ll_Mq$; for all other
$p$ the desired inequality is trivial after decreasing $c_M$.  Since
\[
  \left|\beta-\frac pq\right|=\frac{|\Lambda|}{q\log2},
\]
absorbing fixed factors and one additional power of $q$ gives
\eqref{eq:beta-irrat-measure}.  Matveev's dependence on the logarithmic heights is linear in
$\log M=\beta\log2$ here, which gives the final scale assertion.
\end{proof}

\begin{statusbox}{\statusknown{} plus \statusnew}
The lower-bound theorem for logarithmic forms is classical
\cite{Baker1975,BakerWustholz1993,Matveev2000}.  Its application to the primitive generator,
the continued-fraction consequences below, and the comparison with the source report are
new deductions in this continuation.  No useful small numerical value of $\mu_M$ is claimed;
standard explicit constants are very large.
\end{statusbox}

\subsection{Polynomial recurrence for convergent denominators}

Let $p_n/q_n$ be the continued-fraction convergents of $\beta$.  The standard upper bound
\[
  \left|\beta-\frac{p_n}{q_n}\right|<\frac1{q_nq_{n+1}}
\]
combined with \eqref{eq:beta-irrat-measure} gives the following.

\begin{corollary}[Polynomial denominator growth]\label{cor:q-growth}
There is an effective constant $C_M>0$ such that
\begin{equation}\label{eq:q-growth}
  q_{n+1}\le C_M q_n^{\mu_M-1}
\end{equation}
for every $n$ with $q_n\ge2$.  Consequently the next partial quotient satisfies
$a_{n+1}=O_M(q_n^{\mu_M-2})$.
\end{corollary}

This is qualitatively stronger than the inherited elementary estimate
$q_{n+1}<2M^{q_n}$.  It rules out Liouville behavior and says that the exceptionally long
near-repetitions of the Sturmian word are polynomially controlled in the preceding standard
word length.  The exponent depends on the unknown integer $M$, however, and may be enormous.

\subsection{Polynomial spacing in conditional blocks}

The irrationality measure also upgrades the separation of the block nodes.  From
\eqref{eq:beta-irrat-measure}, for $d\ge1$ and the nearest integer $r$ to $d\beta$,
\begin{equation}\label{eq:beta-mod-one}
  \distZ{d\beta}=|d\beta-r|
  \ge c_M d^{1-\mu_M}.
\end{equation}

For an integer $T\ge0$, put
\[
  K_T=\left\lfloor\frac{T+1}{\beta}\right\rfloor,
  \qquad
  x_{T,k}=T-\lfloor k\beta\rfloor+k\beta
         =T+\fract{k\beta},
  \quad 0\le k\le K_T.
\]
The associated candidate and output nodes are
\begin{equation}\label{eq:block-pairs}
  m_{T,k}=2^{x_{T,k}}=2^{T-\lfloor k\beta\rfloor}M^k,
  \qquad
  a_{T,k}=3^{x_{T,k}}=3^{T-\lfloor k\beta\rfloor}A^k.
\end{equation}

\begin{corollary}[Effective block separation]\label{cor:block-separation}
For distinct $i,j\in\{0,\ldots,K_T\}$,
\begin{align}
 |x_{T,j}-x_{T,i}|&\ge c_MK_T^{1-\mu_M},\label{eq:x-sep}\\
 |m_{T,j}-m_{T,i}|&\ge (\log2)c_M2^T K_T^{1-\mu_M},\label{eq:m-sep}\\
 |a_{T,j}-a_{T,i}|&\ge (\log3)c_M3^T K_T^{1-\mu_M}.
 \label{eq:a-sep}
\end{align}
\end{corollary}

\begin{proof}
The ordinary distance between two fractional parts is at least their circular distance, so
\[
 |\fract{j\beta}-\fract{i\beta}|
 \ge\distZ{(j-i)\beta}
 \ge c_M|j-i|^{1-\mu_M}
 \ge c_MK_T^{1-\mu_M},
\]
because $1-\mu_M<0$.  This proves \eqref{eq:x-sep}.  The mean value theorem applied to $2^x$
and $3^x$ on $[T,T+1]$ gives the remaining inequalities.
\end{proof}

\subsection{Separation on the full \texorpdfstring{$(n,k)$}{(n,k)} lattice}

The same estimate applies to the two-dimensional conditional exponent grid.

\begin{corollary}[Polynomial separation of conditional exponents]\label{cor:grid-separation}
There are effective $c_M,\mu_M>0$ such that, for distinct integer pairs
$(n,k),(n',k')$ with $0\le k,k'\le H$,
\[
  |(n+k\beta)-(n'+k'\beta)|
  \ge \min\{1,c_MH^{1-\mu_M}\}.
\]
\end{corollary}

\begin{proof}
If $k=k'$, the difference is a nonzero integer.  Otherwise divide the linear form
$(n-n')+(k-k')\beta$ by $|k-k'|$ and apply
\eqref{eq:beta-irrat-measure}.
\end{proof}

The column exponent lattice $u+\vartheta v$ has an analogous effective polynomial spacing
with absolute constants, because $\vartheta=\log_2 3$ is a fixed ratio of two logarithms.
More refined explicit estimates are known for linear forms in $\log2$ and $\log3$; for
example Wu's simultaneous linear-independence measure leads, in the usual normalization, to
an irrationality exponent for $\vartheta$ below approximately $8.62$
\cite{Wu2003}.  The exact convention-dependent numerical conversion is not needed below.

\subsection{Why the Baker upgrade does not close the conjecture}

The effective spacing is useful, but its scale is still compatible with the exact conditional
population.  A dyadic exponent block contains only $K_T+1\asymp T$ nodes.  A polynomial lower
bound on their minimum spacing imposes no contradiction on $T$ points in a fixed interval.
Likewise, the full triangular grid has $\asymp T^2$ points in an interval of exponent length
$T$, and polynomially small gaps are expected.

Baker's theorem therefore strengthens compactness and determinant estimates but does not
supply the missing logarithm by itself.  Its most promising roles are indirect:

\begin{itemize}
\item controlling the smallest row and column gaps in an exact Harish--Chandra--Itzykson--Zuber
      factorization;
\item bounding how often a $p$-adic assignment problem can be nearly degenerate;
\item limiting the lengths of near-periodic Sturmian blocks that a congruence argument must
      handle;
\item making computational searches over continued-fraction branches provably finite once a
      candidate output $M$ is fixed.
\end{itemize}

\section{Exact all-support valuations of ordinary block Vandermondes}

\subsection{The block determinants}

Fix an integer $T$ and abbreviate $K=K_T$.  Let $m_k=m_{T,k}$ and $a_k=a_{T,k}$ be the
$K+1$ input--output pairs in \eqref{eq:block-pairs}.  Their ordinary Vandermonde determinants
are
\begin{align}
  V_T^{(2)}&=\det(m_i^j)_{0\le i,j\le K}
            =\prod_{0\le i<j\le K}(m_j-m_i),\label{eq:Vm}\\
  V_T^{(3)}&=\det(a_i^j)_{0\le i,j\le K}
            =\prod_{0\le i<j\le K}(a_j-a_i).
  \label{eq:Va}
\end{align}
They are nonzero integers.  Since $m_k\in[2^T,2^{T+1})$ and
$a_k\in[3^T,3^{T+1})$,
\begin{align}
  \log|V_T^{(2)}|&\le R_TT\log2,\label{eq:Vm-arch}\\
  \log|V_T^{(3)}|&\le R_T(T\log3+\log2),\label{eq:Va-arch}
\end{align}
where
\[
  R_T=\binom{K+1}{2}=\frac{K(K+1)}2.
\]
The additive $R_T\log2$ in \eqref{eq:Va-arch} is lower order.

\subsection{The input determinant: an exact formula at every support prime}

Factor
\[
  M=2^aW,
  \qquad W\text{ odd},
  \qquad
  \delta_2=\beta-a=\log_2W.
\]
As noted above, $W>1$, so $W\ge3$ and
\begin{equation}\label{eq:delta2-gt1}
  \delta_2\ge\log_2 3>1.
\end{equation}
Because $a$ is an integer,
\begin{equation}\label{eq:m-delta-form}
  m_k=2^{T-\lfloor k\delta_2\rfloor}W^k.
\end{equation}
The $2$-adic valuations in \eqref{eq:m-delta-form} strictly decrease with $k$, while every
odd support-prime valuation strictly increases.

\begin{proposition}[Exact support valuations for $V_T^{(2)}$]
\label{prop:input-vander-valuations}
For $0\le i<j\le K$,
\begin{align}
  v_2(m_j-m_i)&=T-\lfloor j\delta_2\rfloor,
  \label{eq:v2-gap}\\
  v_p(m_j-m_i)&=i\,v_p(W)
  \qquad(p\mid W).
  \label{eq:vp-gap}
\end{align}
Consequently
\begin{align}
 v_2(V_T^{(2)})
 &=\sum_{j=1}^Kj\bigl(T-\lfloor j\delta_2\rfloor\bigr),
 \label{eq:v2-V}\\
 v_p(V_T^{(2)})
 &=v_p(W)\sum_{i=0}^{K-1}i(K-i)
 \qquad(p\mid W).
 \label{eq:vp-V}
\end{align}
\end{proposition}

\begin{proof}
Put $\rho_k=T-\lfloor k\delta_2\rfloor$.  By
\eqref{eq:delta2-gt1}, $\rho_i>\rho_j$ when $i<j$.  From
\eqref{eq:m-delta-form},
\[
 m_j-m_i
 =2^{\rho_j}W^i\left(W^{j-i}-2^{\rho_i-\rho_j}\right).
\]
The parenthesis is odd, proving \eqref{eq:v2-gap}.  For an odd prime $p\mid W$, the two terms
$m_i$ and $m_j$ have distinct valuations $i v_p(W)<jv_p(W)$, so the ultrametric equality gives
\eqref{eq:vp-gap}.  Summing over pairs gives \eqref{eq:v2-V}--\eqref{eq:vp-V}.
\end{proof}

Let $Q_T^{(2)}$ be the part of $V_T^{(2)}$ supported on the primes dividing $2M$.  Then
Proposition~\ref{prop:input-vander-valuations} gives the exact identity
\begin{equation}\label{eq:logQm-exact}
 \log Q_T^{(2)}
 =\log2\sum_{j=1}^Kj\bigl(T-\lfloor j\delta_2\rfloor\bigr)
  +\log W\sum_{i=0}^{K-1}i(K-i).
\end{equation}
The two elementary sums have asymptotics
\begin{align*}
 \sum_{j=1}^Kj\bigl(T-\lfloor j\delta_2\rfloor\bigr)
 &=\frac{TK^2}{2}-\frac{\delta_2K^3}{3}+O_\beta(T^2),\\
 \sum_{i=0}^{K-1}i(K-i)&=\frac{K^3}{6}+O(K^2).
\end{align*}
Using $\log W=\delta_2\log2$ and $K/T\to1/\beta$ yields
\begin{equation}\label{eq:Qm-ratio}
 \frac{\log Q_T^{(2)}}{R_TT\log2}
 \longrightarrow
 1-\frac{\delta_2}{3\beta}.
\end{equation}
This includes every prime appearing in the entries.  Any further prime divisor of the
Vandermonde comes from subtraction and hence from genuine cancellation beyond termwise
entry valuations.

\subsection{The output determinant and bounded tie corrections}

Factor
\[
  A=3^bZ,
  \qquad 3\nmid Z,
  \qquad
  \delta_3=\beta-b=\log_3Z>0.
\]
Then
\begin{equation}\label{eq:a-delta-form}
  a_k=3^{T-\lfloor k\delta_3\rfloor}Z^k.
\end{equation}
Unlike $\delta_2$, the number $\delta_3$ can be smaller than $1$; for example the least
possible base-free factor is $2$.  Equal consecutive $3$-adic valuations can therefore occur.
They contribute only a lower-order correction.

\begin{proposition}[Output support valuations]\label{prop:output-vander-valuations}
For every prime $p\mid Z$ and $0\le i<j\le K$,
\[
  v_p(a_j-a_i)=i\,v_p(Z).
\]
At $p=3$,
\[
 v_3(a_j-a_i)=T-\lfloor j\delta_3\rfloor
\]
after adding a nonnegative correction that is zero whenever
$\lfloor i\delta_3\rfloor<\lfloor j\delta_3\rfloor$.  The total correction over all pairs is
$O_{A,\beta}(K^2)$.  Consequently, if $Q_T^{(3)}$ denotes the part of
$V_T^{(3)}$ supported on primes dividing $3A$, then
\begin{equation}\label{eq:Qa-ratio}
  \frac{\log Q_T^{(3)}}{R_TT\log3}
  \longrightarrow
  1-\frac{\delta_3}{3\beta}.
\end{equation}
\end{proposition}

\begin{proof}
The assertion for $p\mid Z$ is the same distinct-valuation argument as before.  Put
$\rho_k=T-\lfloor k\delta_3\rfloor$.  If $\rho_i>\rho_j$, then
\[
 a_j-a_i=3^{\rho_j}Z^i
 \left(Z^{j-i}-3^{\rho_i-\rho_j}\right)
\]
and the parenthesis is not divisible by $3$, so the valuation is exactly $\rho_j$.  If
$\rho_i=\rho_j$, then
\[
 v_3(a_j-a_i)=\rho_j+v_3(Z^{j-i}-1).
\]
Equality of the floors implies $(j-i)\delta_3<1$, so $j-i$ belongs to a fixed finite set
depending only on $\delta_3$.  Hence the extra valuation is bounded by a constant depending on
$A$ and $\beta$.  There are $O(K^2)$ pairs.  The asymptotic calculation is then identical to
\eqref{eq:logQm-exact}.
\end{proof}

\subsection{The depth-sensitive product ceiling}

Combining \eqref{eq:Vm-arch}, \eqref{eq:Va-arch},
\eqref{eq:Qm-ratio}, and \eqref{eq:Qa-ratio} gives the principal conclusion of the ordinary
Vandermonde audit.

\begin{theorem}[All-support ordinary Vandermonde fraction]\label{thm:ordinary-fraction}
For a hypothetical primitive counterexample,
\begin{equation}\label{eq:ordinary-fraction}
 \frac{\log(Q_T^{(2)}Q_T^{(3)})}{R_TT\log6}
 \longrightarrow
 1-\frac{\delta_2\log2+\delta_3\log3}{3\beta\log6},
 \qquad
 \delta_2=\beta-a,
 \quad
 \delta_3=\beta-b.
\end{equation}
Because $\min(a,b)=0$, the limit is at most
\begin{equation}\label{eq:ordinary-universal-ceiling}
 1-\frac{\log2}{3\log6}
 =0.871049\,064\ldots.
\end{equation}
If $b=0$, the sharper ceiling is
\[
 1-\frac{\log3}{3\log6}=0.795599\,921\ldots.
\]
If $a=b=0$, the limit is exactly $2/3$.
\end{theorem}

\begin{proof}
Only the numerical ceiling remains to be shown.  If $a=0$, then $\delta_2=\beta$, so the
deficit in \eqref{eq:ordinary-fraction} is at least $\log2/(3\log6)$.  If $b=0$, it is at
least $\log3/(3\log6)$.  Taking the larger possible fraction gives
\eqref{eq:ordinary-universal-ceiling}.  If both vanish, the numerator of the deficit is
$\beta\log6$.
\end{proof}

\begin{table}[ht]
\centering
\caption{Leading termwise-divisor ceilings for ordinary block Vandermondes.}
\begin{tabular}{@{}lll@{}}
\toprule
Branch information & Exact/upper fraction & Decimal\\
\midrule
$a=b=0$ & $2/3$ & $0.666666\ldots$\\
$b=0$ & $\le1-\log3/(3\log6)$ & $0.795599\ldots$\\
No matched depth only & $\le1-\log2/(3\log6)$ & $0.871049\ldots$\\
\bottomrule
\end{tabular}
\end{table}

\subsection{Why even \texorpdfstring{$87\%$}{87 percent} is not enough}

The deficit in Theorem~\ref{thm:ordinary-fraction} is a fixed positive fraction once the
hypothetical primitive pair is fixed.  The archimedean height is $\asymp T^3$, whereas all
factorial, spacing, and tie corrections are $O_\beta(T^2\log T)$.  Thus no lower-order
refinement of the ordinary Vandermonde can close the remaining gap.

This conclusion is stronger than the earlier observation that the $2$-adic valuation alone
captures one third in the base-free case.  It includes every prime already present in the
entries and every possible base depth.  A successful ordinary-Vandermonde proof would need a
new leading-order source of divisibility in the prime factors of the \emph{differences}; the
fixed-support valuations are completely accounted for by
Theorem~\ref{thm:ordinary-fraction}.

\section{The mixed-semigroup determinant and its Gini norm}
\label{sec:mixed}

\subsection{Balanced low-weight columns}

Let
\[
  \Lambda=\{u+\vartheta v:(u,v)\in\Nzero^2\}.
\]
Because $\vartheta$ is irrational, every element has a unique pair $(u,v)$.  List the pairs
in increasing order of $u+\vartheta v$:
\[
  (u_0,v_0),(u_1,v_1),\ldots.
\]
For $N\ge1$, let
\[
  L_N=\max_{0\le j<N}(u_j+\vartheta v_j).
\]
The elementary lattice count
\begin{equation}\label{eq:semigroup-count}
 \#\{(u,v)\in\Nzero^2:u+\vartheta v\le L\}
 =\frac{L^2}{2\vartheta}+O_\vartheta(L)
\end{equation}
shows that
\begin{equation}\label{eq:L-asymp}
  L_N\sim\sqrt{2\vartheta N}.
\end{equation}
Moreover, the empirical distribution of
\begin{equation}\label{eq:scaled-cols}
  X_j=\left(\frac{u_j}{L_N},\frac{\vartheta v_j}{L_N}\right)
\end{equation}
converges to uniform area measure on the standard triangle
\begin{equation}\label{eq:triangle}
  \DeltaTri=\{(x,z)\in\R^2:x\ge0,
  z\ge0,
  x+z\le1\}.
\end{equation}
Here and below ``uniform'' means probability measure with density $2$ on this triangle.

For the block at height $T$, set $K=K_T$ and $N=K+1$.  Define
\begin{equation}\label{eq:mixed-D}
  D_T=\det\left(m_{T,k}^{u_j}a_{T,k}^{v_j}\right)_{0\le k,j<N}.
\end{equation}
The generalized-Vandermonde theorem inherited from the source report implies that $D_T$ is a
nonzero integer.  Put
\[
  B_j=2^{u_j}3^{v_j}.
\]
Since $m_{T,k}=2^{x_{T,k}}$ and $a_{T,k}=3^{x_{T,k}}$,
\begin{equation}\label{eq:entry-B}
  m_{T,k}^{u_j}a_{T,k}^{v_j}=B_j^{x_{T,k}}
  =B_j^{T+\fract{k\beta}}.
\end{equation}
Thus the Leibniz expansion gives the clean archimedean bound
\begin{equation}\label{eq:mixed-arch}
  \log|D_T|
  \le (T+1)\sum_{j<N}\log B_j+\log N!.
\end{equation}
The leading height is
\begin{equation}\label{eq:HT}
  H_T=T\sum_{j<N}\log B_j.
\end{equation}
By \eqref{eq:L-asymp} and the mean
$\mathbb E(X_1+X_2)=2/3$ on $\DeltaTri$,
\begin{equation}\label{eq:H-asymp}
 H_T\sim \frac23\,TNL_N\log2
 \asymp_\beta T^{5/2}.
\end{equation}
Both $\sum_j\log B_j$ and $\log N!$ in the error of \eqref{eq:mixed-arch} are
$o(H_T)$.

\subsection{The all-prime tropical divisor}

For a prime $p$, define the assignment minimum
\begin{equation}\label{eq:taup}
 \tau_{p,T}
 =\min_{\sigma\in S_N}
   \sum_{k=0}^{N-1}
   v_p\!\left(m_{T,k}^{u_{\sigma(k)}}a_{T,k}^{v_{\sigma(k)}}\right).
\end{equation}
Every Leibniz term is divisible by $p^{\tau_{p,T}}$, so
\begin{equation}\label{eq:Qtrop}
  Q_T^{\mathrm{trop}}=\prod_p p^{\tau_{p,T}}
  \quad\text{divides}\quad D_T.
\end{equation}
Only primes dividing $6MA$ occur in the product.

The problem is now to compute the leading asymptotic of
$\log Q_T^{\mathrm{trop}}/H_T$.  A crude estimate ``minimum $\le$ average'' loses the precise
branch dependence.  The exact answer is controlled by a simple probability norm.

\subsection{Antimonotone assignment and the Gini norm}

Let $X=(X_1,X_2)$ be uniform on $\DeltaTri$, and let $X'$ be an independent copy.  For
$c=(r,s)\in\R^2$, define
\begin{equation}\label{eq:Gini}
  \mathfrak G(r,s)
  =\mathbb E\left|r(X_1-X_1')+s(X_2-X_2')\right|.
\end{equation}
This is a norm on $\R^2$: homogeneity and the triangle inequality are immediate, and the
support of $X-X'$ spans $\R^2$, so only $(r,s)=(0,0)$ has norm zero.

Let $Y=rX_1+sX_2$, with ascending quantile $Q_Y(t)$.  Pairing an increasing uniform row
parameter with the decreasing ordering of $Y$ gives the continuum assignment functional
\begin{equation}\label{eq:Phi-def}
  \Phi(r,s)=\int_0^1 t\,Q_Y(1-t)\,dt.
\end{equation}

\begin{lemma}[Gini formula for the minimum assignment]\label{lem:gini-assignment}
For every $(r,s)\in\R^2$,
\begin{equation}\label{eq:Phi-G}
  \Phi(r,s)
  =\frac12\,\mathbb E Y-\frac14\,\mathfrak G(r,s)
  =\frac{r+s}{6}-\frac14\,\mathfrak G(r,s).
\end{equation}
\end{lemma}

\begin{proof}
The rearrangement inequality says that the minimum scalar product pairs the increasing row
parameter with the decreasing list of column values.  Changing variables $u=1-t$ in
\eqref{eq:Phi-def} gives
\[
  \Phi=\int_0^1(1-u)Q_Y(u)\,du.
\]
If $Y'$ is an independent copy, the minimum of $Y,Y'$ has quantile density
$2(1-u)$, so
\[
  \Phi=\frac12\mathbb E\min(Y,Y').
\]
Since
$\min(Y,Y')=(Y+Y'-|Y-Y'|)/2$, this becomes
$\frac12\mathbb EY-\frac14\mathbb E|Y-Y'|$.  Finally
$\mathbb EX_1=\mathbb EX_2=1/3$.
\end{proof}

\subsection{Closed form of the Gini norm}

The norm \eqref{eq:Gini} is elementary.  Sort the three vertex values
\[
  0,
  \qquad r,
  \qquad s
\]
as $\xi_0\le\xi_1\le\xi_2$, and put
\[
  A_0=\xi_1-\xi_0,
  \qquad
  B_0=\xi_2-\xi_1.
\]

\begin{proposition}[Explicit triangular Gini norm]\label{prop:G-explicit}
If $A_0+B_0>0$, then
\begin{equation}\label{eq:G-explicit}
 \mathfrak G(r,s)
 =\frac{2\left(2A_0^2+3A_0B_0+2B_0^2\right)}
        {15(A_0+B_0)}.
\end{equation}
At $(r,s)=(0,0)$ the value is $0$.  In particular,
\begin{equation}\label{eq:G-special}
  \mathfrak G(1,0)=\mathfrak G(0,1)
  =\mathfrak G(1,1)=\frac4{15},
  \qquad
  \mathfrak G(1,-1)=\frac7{15}.
\end{equation}
\end{proposition}

\begin{proof}
Translation of the three vertex values does not change a difference $Y-Y'$, so take
$\xi_0=0$, $\xi_1=A_0$, and $\xi_2=A_0+B_0=L$.  Projection of uniform area measure on a
triangle by a linear function has density
\[
 f(y)=
 \begin{cases}
 \displaystyle\frac{2y}{A_0L},&0\le y\le A_0,\\[6pt]
 \displaystyle\frac{2(L-y)}{B_0L},&A_0\le y\le L.
 \end{cases}
\]
Thus its distribution function is
$F(y)=y^2/(A_0L)$ on the first interval and
$F(y)=1-(L-y)^2/(B_0L)$ on the second.  The standard identity
\[
  \mathbb E|Y-Y'|=2\int_{-\infty}^{\infty}F(y)(1-F(y))\,dy
\]
gives \eqref{eq:G-explicit} after two polynomial integrations.  Cases with one gap zero
follow by continuity.
\end{proof}

For later use, if $\gamma\ge0$, then
\begin{equation}\label{eq:G-minus1gamma}
 \mathfrak G(-1,\gamma)
 =\frac{2(2+3\gamma+2\gamma^2)}{15(1+\gamma)}
 =\frac4{15}+\frac{2\gamma(1+2\gamma)}{15(1+\gamma)}
 \ge\frac4{15}.
\end{equation}

\subsection{Normalized prime-valuation slope vectors}

Recall
\[
 a=v_2(M),\quad b=v_3(A),\quad c=v_2(A),\quad d=v_3(M).
\]
For each prime, define a normalized slope vector $c_p\in\R^2$ by
\begin{align}
 c_2&=\left(-1+\frac a\beta,\ \frac{c}{\beta\vartheta}\right),
 \label{eq:c2}\\
 c_3&=\left(\frac{\vartheta d}{\beta},\ -1+\frac b\beta\right),
 \label{eq:c3}\\
 c_p&=\left(
  \frac{v_p(M)\log p}{\beta\log2},
  \frac{v_p(A)\log p}{\beta\log3}
 \right),\qquad p\ne2,3.
 \label{eq:cp-external}
\end{align}
The logarithmic factorizations of $M$ and $A$ give the conservation law
\begin{equation}\label{eq:slope-conservation}
  \sum_p c_p=(0,0).
\end{equation}
Indeed, the first coordinate sum is
\[
 -1+\frac{a\log2+d\log3+
   \sum_{p\ne2,3}v_p(M)\log p}{\beta\log2}=0,
\]
and similarly for the second coordinate.

The meaning of these vectors can be read directly from the valuation matrices.  Put
$s_k=\beta k/T$.  Up to $o(L_N)$ after division by $T$, the logarithmic $p$-adic cost of the
entry indexed by a scaled column $(x,z)$ is
\begin{center}
\begin{tabular}{@{}ll@{}}
\toprule
Prime & Cost divided by $T L_N\log2$\\
\midrule
$2$ & $x+s_k\,c_2\mathbin{\cdot}(x,z)$\\
$3$ & $z+s_k\,c_3\mathbin{\cdot}(x,z)$\\
$p\ne2,3$ & $s_k\,c_p\mathbin{\cdot}(x,z)$\\
\bottomrule
\end{tabular}
\end{center}
For example,
\[
 v_2(m_k^ua_k^v)=u\bigl(T-\lfloor k\beta\rfloor+ak\bigr)+vck,
\]
and substituting $u=L_Nx$, $v=L_Nz/\vartheta$, $k\sim Ts_k/\beta$ gives the first row.
The $O(u)$ errors from the floors sum to $O(NL_N)=o(H_T)$.

\subsection{The exact continuum limit}

\begin{theorem}[All-prime tropical fraction for the mixed block determinant]
\label{thm:mixed-gini}
Under a hypothetical primitive counterexample, the determinant \eqref{eq:mixed-D} satisfies
\begin{equation}\label{eq:mixed-gini-limit}
 \frac{\log Q_T^{\mathrm{trop}}}{H_T}
 \longrightarrow
 1-\frac38\sum_p\mathfrak G(c_p),
\end{equation}
where the finitely many slope vectors $c_p$ are given by
\eqref{eq:c2}--\eqref{eq:cp-external}.
\end{theorem}

\begin{proof}
The row parameters $s_k=\beta k/T$ become uniformly distributed on $[0,1]$, because
$K_T\sim T/\beta$.  The scaled columns become uniformly distributed on $\DeltaTri$ by
\eqref{eq:semigroup-count}.  The finite rearrangement inequality and Riemann-sum convergence
therefore show that the assignment minimum at prime $p$, after normalization by
$TL_NN\log2$, tends to its constant-column mean plus $\Phi(c_p)$.

Only $p=2$ has constant part $x$, and only $p=3$ has constant part $z$.  Their combined mean is
$2/3$.  By Lemma~\ref{lem:gini-assignment}, the sum of the slope contributions is
\[
 \sum_p\Phi(c_p)
 =\frac16\sum_p(c_{p,1}+c_{p,2})
  -\frac14\sum_p\mathfrak G(c_p).
\]
The first sum vanishes by \eqref{eq:slope-conservation}.  Hence
\[
 \frac{\log Q_T^{\mathrm{trop}}}{TL_NN\log2}
 \longrightarrow
 \frac23-\frac14\sum_p\mathfrak G(c_p).
\]
On the other hand, \eqref{eq:H-asymp} gives
$H_T/(TL_NN\log2)\to2/3$.  Dividing proves
\eqref{eq:mixed-gini-limit}.
\end{proof}

This theorem is exact at the level of minimum Leibniz-term valuations.  It replaces a family
of separate $p$-adic estimates by one finite norm sum satisfying a conservation law.

The four base and cross depths already give a useful branch-sensitive bound without knowing
any external prime factorization.

\begin{corollary}[Depth-sensitive three-vector bound]\label{cor:three-vector}
For every hypothetical primitive counterexample,
\begin{equation}\label{eq:three-vector}
 \limsup_{T\to\infty}
 \frac{\log Q_T^{\mathrm{trop}}}{H_T}
 \le
 1-\frac38\Bigl(
   \mathfrak G(c_2)+\mathfrak G(c_3)
   +\mathfrak G(-c_2-c_3)\Bigr).
\end{equation}
Equality in this upper bound can occur only when all nonzero external slope vectors are
positively collinear.
\end{corollary}

\begin{proof}
The external vectors satisfy
$\sum_{p\ne2,3}c_p=-c_2-c_3$.  Hence the triangle inequality gives
\[
 \sum_{p\ne2,3}\mathfrak G(c_p)
 \ge\mathfrak G(-c_2-c_3).
\]
Substitute this in Theorem~\ref{thm:mixed-gini}.  Equality in a norm triangle inequality
requires all summands to lie on one nonnegative ray.
\end{proof}

The equality condition detects exactly the second-iterate kernel branch reviewed in the source
report.  Indeed the external vectors $c_p$ are rescaled versions of
$(v_p(M),v_p(A))$.  They are all collinear if and only if the two external valuation vectors
of $M$ and $A$ are rationally dependent, equivalently if the kernel $K_\beta$ is nonzero.
Thus:

\begin{corollary}[The tropical norm sees the kernel dichotomy]\label{cor:gini-kernel}
In the multiplicatively independent branch $K_\beta=0$, the inequality in
\eqref{eq:three-vector} is strict.  In the dependent branch $K_\beta\ne0$, the external
triangle inequality is an equality (including the axis cases in which one output is
$\{2,3\}$-smooth).  For fixed $M,A$, the strict gap in the independent branch is explicit:
\begin{equation}\label{eq:external-angular-defect}
 \Delta_{\rm ext}
 =\sum_{p\ne2,3}\mathfrak G(c_p)
  -\mathfrak G\left(\sum_{p\ne2,3}c_p\right)>0.
\end{equation}
\end{corollary}

There is no uniform positive lower bound for $\Delta_{\rm ext}$: distinct rational rays can
approach one another as the valuations grow.  Nevertheless, the quantity is a computable
branch invariant and gives a sharper determinant threshold for every fixed hypothetical
primitive pair.

\subsection{The universal \texorpdfstring{$4/5$}{4/5} ceiling}

\begin{theorem}[Sharp universal min-plus ceiling]\label{thm:four-fifths}
For every hypothetical primitive counterexample,
\begin{equation}\label{eq:G-lower}
  \sum_p\mathfrak G(c_p)\ge\frac8{15}.
\end{equation}
Consequently
\begin{equation}\label{eq:four-fifths}
 \limsup_{T\to\infty}
 \frac{\log Q_T^{\mathrm{trop}}}{H_T}
 \le\frac45.
\end{equation}
The constant $4/5$ is the sharp supremum allowed by the normalized valuation constraints,
although equality would require a degenerate boundary incompatible with a genuine
nonintegral solution.
\end{theorem}

\begin{proof}
By the primitive theorem, either $a=0$ or $b=0$.  Suppose first that $a=0$.  Then
$c_2=(-1,\gamma)$ with $\gamma=c/(\beta\vartheta)\ge0$.  By
\eqref{eq:G-minus1gamma}, $\mathfrak G(c_2)\ge4/15$.  Conservation and the triangle inequality
for the norm give
\[
  \sum_{p\ne2}\mathfrak G(c_p)
  \ge\mathfrak G\!\left(\sum_{p\ne2}c_p\right)
  =\mathfrak G(-c_2)
  =\mathfrak G(c_2).
\]
Thus the total is at least $8/15$.  If $b=0$, apply the same argument to
$c_3=(\eta,-1)$, where $\eta=\vartheta d/\beta\ge0$.  Substitution in
\eqref{eq:mixed-gini-limit} gives
$1-(3/8)(8/15)=4/5$.
\end{proof}

The proof is short because the difficult arithmetic has been compressed into two facts: the
slope vectors sum to zero, and one structural vector reaches from $-1$ to a nonnegative
vertex.  No averaging estimate alone reveals this geometry.

\subsection{Cross-coprime constants: \texorpdfstring{$7/10$ and $3/10$}{7/10 and 3/10}}

Assume the stronger branch condition
\begin{equation}\label{eq:cross-coprime}
  \gcd(MA,6)=1.
\end{equation}
Then $a=b=c=d=0$,
\[
  c_2=(-1,0),
  \qquad
  c_3=(0,-1),
  \qquad
  \sum_{p\ne2,3}c_p=(1,1).
\]

\begin{corollary}[All-prime cross-coprime ceiling]\label{cor:seven-tenths}
Under \eqref{eq:cross-coprime},
\begin{equation}\label{eq:seven-tenths}
 \limsup_{T\to\infty}
 \frac{\log Q_T^{\mathrm{trop}}}{H_T}
 \le\frac7{10}.
\end{equation}
If the external normalized valuation vectors are not all positively collinear, the inequality
is strict by a branch-dependent positive amount.
\end{corollary}

\begin{proof}
By \eqref{eq:G-special}, the structural vectors each contribute $4/15$.  The triangle
inequality gives
\[
  \sum_{p\ne2,3}\mathfrak G(c_p)
  \ge\mathfrak G(1,1)=\frac4{15}.
\]
Thus the total norm is at least $4/5$, and
$1-(3/8)(4/5)=7/10$.  Equality in the triangle inequality requires positive collinearity.
\end{proof}

The places $2$ and $3$ alone have an even smaller exact continuum contribution.

\begin{corollary}[The structural-place $3/10$ constant]\label{cor:three-tenths}
Under \eqref{eq:cross-coprime}, let
$Q_T^{(2,3)}=2^{\tau_{2,T}}3^{\tau_{3,T}}$.  Then
\begin{equation}\label{eq:three-tenths}
  \frac{\log Q_T^{(2,3)}}{H_T}\longrightarrow\frac3{10}.
\end{equation}
\end{corollary}

\begin{proof}
For $c=(-1,0)$,
\[
 \Phi(c)=\frac{-1}{6}-\frac14\cdot\frac4{15}=-\frac7{30},
\]
and the same holds for $(0,-1)$.  Adding the constant mean $2/3$ gives
$2/3-7/15=1/5$.  Dividing by the height mean $2/3$ gives $3/10$.
\end{proof}

Equivalently, one can compute this constant from quantiles.  The descending quantile of
$X_1$ on $\DeltaTri$ is $1-\sqrt t$, so
\[
 \int_0^1t(1-\sqrt t)\,dt=\frac1{10},
 \qquad
 \mathbb EX_1=\frac13.
\]
The ratio is $(1/10)/(1/3)=3/10$ in each structural coordinate.

\subsection{Correction and calibration of a later repository note}

A short $p$-adic priority note added to the repository after the attached unified report makes
the correct strategic observation that minimum Leibniz-term valuations leave a constant
fraction of the archimedean height uncovered.  Two details require correction.

First, the displayed structural valuation formulas in that note simultaneously omit the
terms $kuv_2(M)$ and $kvv_3(A)$; they therefore assume both $a=0$ and $b=0$.  The inherited
primitive theorem gives only $\min(a,b)=0$.  Equations
\eqref{eq:c2}--\eqref{eq:c3} retain all four cross/base depths.

Second, the proposed all-prime $2/3$ ceiling is not a valid universal consequence of
``minimum $\le$ average.''  For an external positive weight proportional to $X_1+X_2$, the
minimum assignment is $4/5$ of its average assignment, not at most $2/3$.  The exact Gini
formula shows that the universal normalized ceiling is $4/5$, while the cross-coprime
all-prime ceiling is $7/10$.  The note's qualitative conclusion survives---entrywise
valuation is insufficient---but the constants and branch hypotheses are materially different.

\begin{statusbox}{\statusnew}
The correction is not based on a numerical counterexample to the Alaoglu--Erd\H{o}s
conjecture.  It is an exact audit of the assignment inequality used in the method.  Rational
external valuation directions can approximate the equality configurations arbitrarily well,
so no stronger uniform constant follows merely from integrality of the valuation vectors.
\end{statusbox}

\subsection{Where extra divisibility can occur}

The tropical minimum equals the true determinant valuation whenever its minimizing
permutation is unique.

\begin{proposition}[Unique tropical term forbids cancellation]\label{prop:unique-tropical}
Let $D=\det(E_{ij})$ be a nonzero integer determinant.  Fix a prime $p$, and suppose a unique
permutation $\sigma_0$ minimizes
$\sum_i v_p(E_{i,\sigma(i)})$.  Then
\[
  v_p(D)=\sum_i v_p(E_{i,\sigma_0(i)}).
\]
\end{proposition}

\begin{proof}
Divide the Leibniz expansion by the corresponding power of $p$.  The term indexed by
$\sigma_0$ is a $p$-adic unit, while every other term is divisible by $p$.  The sum is
therefore a unit modulo $p$.
\end{proof}

Consequently, the missing $20\%$ or more can arise only from primes at which the assignment
problem has several minimizers.  The tie geometry is explicit.

\begin{itemize}
\item For $p\ne2,3$, columns tie when
      $v_p(M)(u-u')+v_p(A)(v-v')=0$, so ties lie on rational lines orthogonal to the external
      valuation vector.
\item In the cross-coprime branch, the $2$-adic cost depends only on $u$ and the $3$-adic cost
      only on $v$.  Every vertical or horizontal column fiber creates many equal-cost
      permutations.  Any leading extra valuation must be produced inside these fibers.
\item Cross valuations $v_2(A)$ and $v_3(M)$ tilt the fibers.  The floor term in the row
      profile perturbs the exact affine ordering but does not change the continuum slope
      vector.
\end{itemize}

This turns ``look for $p$-adic cancellation'' into a concrete algebra problem: factor or
control the residual determinants inside equal-slope column fibers, and show that their
valuation supplies at least the missing Gini deficit.


\section{The complete triangular grid and a \texorpdfstring{$\pi/4$}{pi/4} transport ceiling}
\label{sec:fullgrid}

The one-block determinant deliberately uses the exact $O(T)$ population in a unit interval of
exponent space.  The source report also suggests exploiting the full two-dimensional
coordinate grid rather than sorting only by magnitude.  Conditional on a counterexample, that
proposal gives $\asymp T^2$ genuine integer rows below exponent height $T$, exactly the number
that the mixed-semigroup interpolation argument seemed to be missing.  It is therefore
important to determine whether entrywise $p$-adic divisibility becomes strong enough on this
larger node set.

The answer is again negative, but for a different and pleasantly universal reason.  Once all
places are summed, the prime-by-prime assignment minima are bounded by one antimonotone
transport problem for the logarithmic entry matrix.  Both the row heights and the balanced
column weights have the same triangular distribution, and the ratio between antimonotone and
monotone transport is exactly $\pi/4$.

\subsection{Rows and columns}

For $T>0$ define the complete conditional row set
\begin{equation}\label{eq:grid-rows}
 \mathcal R_T=
 \{(n,k)\in\Nzero^2:n+k\beta\le T\},
 \qquad N_T=\#\mathcal R_T.
\end{equation}
For $i=(n_i,k_i)\in\mathcal R_T$ put
\begin{equation}\label{eq:grid-row-values}
 x_i=n_i+k_i\beta,
 \qquad m_i=2^{n_i}M^{k_i}=2^{x_i},
 \qquad a_i=3^{n_i}A^{k_i}=3^{x_i}.
\end{equation}
Because $\beta$ is irrational, the $x_i$ are pairwise distinct.

Let $\mathcal C_T$ be the first $N_T$ points $(u,v)\in\Nzero^2$ in increasing order of
\[
 \lambda=u+\vartheta v,
 \qquad
 y=\log(2^u3^v)=\lambda\log2.
\]
Write $L_T$ for the largest selected $\lambda$.  Standard lattice-point estimates for the two
triangles give
\begin{equation}\label{eq:grid-cardinality}
 N_T=\frac{T^2}{2\beta}+O_\beta(T),
 \qquad
 L_T=\sqrt{2\vartheta N_T}+O_\vartheta(1)
      =T\sqrt{\frac{\vartheta}{\beta}}+O_\beta(1).
\end{equation}

The associated square determinant is
\begin{equation}\label{eq:grid-determinant}
 D_T^{\triangle}
 =\det\!\left(m_i^{u_j}a_i^{v_j}\right)_{i,j=1}^{N_T}
 =\det\!\left(e^{x_i y_j}\right)_{i,j=1}^{N_T}.
\end{equation}
Every entry is an integer.  Ordering the $x_i$ and $y_j$ increasingly, strict total positivity
of the exponential kernel makes $D_T^{\triangle}>0$; equivalently, this is the same
exponential-polynomial zero theorem used by the generalized-Vandermonde part of the source
report.

\begin{statusbox}{\statusnew{} for the asymptotic transport statements below.  Integrality and
nonvanishing are inherited from the mixed-semigroup determinant architecture
\statuslean.}
\end{statusbox}

\subsection{The common triangular distribution}

The normalized row heights $h_i=x_i/T$ and column weights $\ell_j=y_j/(L_T\log2)$ both live in
$[0,1]$.  Their empirical measures have the same limit.

\begin{lemma}[Triangular height law]\label{lem:triangular-law}
For every continuous $f:[0,1]\to\R$,
\begin{align}
 \frac1{N_T}\sum_{i=1}^{N_T}f(h_i)
   &\longrightarrow \int_0^1 2t f(t)\,dt, \label{eq:row-law}\\
 \frac1{N_T}\sum_{j=1}^{N_T}f(\ell_j)
   &\longrightarrow \int_0^1 2t f(t)\,dt. \label{eq:column-law}
\end{align}
Thus the common limiting distribution has cumulative distribution function $F(t)=t^2$ and
quantile function
\begin{equation}\label{eq:triangular-quantile}
 Q(s)=\sqrt s,\qquad 0\le s\le1.
\end{equation}
\end{lemma}

\begin{proof}
After scaling $(n,k)$ to $(n/T,\beta k/T)$, the row lattice becomes equidistributed in the
unit right triangle
\[
 \{(r,s):r,s\ge0,\ r+s\le1\}.
\]
The subtriangle on which $r+s\le t$ has area $t^2/2$, while the whole triangle has area
$1/2$.  Hence the height $r+s$ has distribution function $t^2$.  Boundary discrepancies are
$O(T)$ against $N_T\asymp T^2$, proving \eqref{eq:row-law}.

For the columns, scale $(u,v)$ to $(u/L_T,\vartheta v/L_T)$.  The same unit triangle appears,
and $\ell=(u+\vartheta v)/L_T$ is again the sum of the scaled coordinates.  The boundary shell
created by taking the first exactly $N_T$ points has $O(L_T)$ points against
$N_T\asymp L_T^2$, which proves \eqref{eq:column-law}.
\end{proof}

\subsection{Summing all primewise assignment minima}

For each prime $p$, let
\begin{equation}\label{eq:grid-tau}
 \tau_{p,T}
 =\min_{\sigma\in S_{N_T}}
   \sum_{i=1}^{N_T}
   v_p\!\left(m_i^{u_{\sigma(i)}}a_i^{v_{\sigma(i)}}\right),
 \qquad
 Q_T^{\triangle}=\prod_p p^{\tau_{p,T}}.
\end{equation}
As before, $Q_T^{\triangle}$ divides every Leibniz term and hence divides
$D_T^{\triangle}$.  The key observation is that taking a separate minimizing permutation at
each prime can only decrease the sum relative to using one permutation at all primes.

\begin{lemma}[All-place transport domination]\label{lem:all-place-transport}
For every $T$,
\begin{equation}\label{eq:all-place-min}
 \log Q_T^{\triangle}
 \le
 \min_{\sigma\in S_{N_T}}
 \sum_{i=1}^{N_T}x_i y_{\sigma(i)}.
\end{equation}
\end{lemma}

\begin{proof}
For a permutation $\sigma$, put
\[
 C_p(\sigma)=
 \sum_i v_p\!\left(m_i^{u_{\sigma(i)}}a_i^{v_{\sigma(i)}}\right)\log p.
\]
Then
\[
 \sum_p\min_\sigma C_p(\sigma)
 \le \min_\sigma\sum_p C_p(\sigma).
\]
Unique factorization and \eqref{eq:grid-row-values} give, entry by entry,
\[
 \sum_p v_p(m_i^{u_j}a_i^{v_j})\log p
 =\log(m_i^{u_j}a_i^{v_j})
 =x_i y_j.
\]
Substitution proves \eqref{eq:all-place-min}.
\end{proof}

Notice how little arithmetic remains in the last inequality: it is valid in every valuation
branch, includes all primes, and does not use the primitive-depth normalization.  The price is
that it is only an upper bound for the common termwise divisor.  That is exactly what is needed
for a no-go theorem.

\subsection{Optimal transport and the \texorpdfstring{$\pi/4$}{pi/4} ratio}

Let
\begin{align}
 E_T^-&=\min_{\sigma\in S_{N_T}}\sum_i x_i y_{\sigma(i)},
 \label{eq:grid-Eminus}\\
 E_T^+&=\max_{\sigma\in S_{N_T}}\sum_i x_i y_{\sigma(i)}.
 \label{eq:grid-Eplus}
\end{align}
The rearrangement inequality pairs the two ordered lists oppositely for $E_T^-$ and in the
same order for $E_T^+$.

\begin{theorem}[Triangular-grid transport constants]\label{thm:grid-transport}
As $T\to\infty$,
\begin{align}
 \frac{E_T^-}{N_TTL_T\log2}
   &\longrightarrow
     \int_0^1\sqrt{t(1-t)}\,dt
      =\frac\pi8, \label{eq:grid-min-limit}\\
 \frac{E_T^+}{N_TTL_T\log2}
   &\longrightarrow
     \int_0^1t\,dt
      =\frac12. \label{eq:grid-max-limit}
\end{align}
Consequently
\begin{equation}\label{eq:grid-ratio}
 \frac{E_T^-}{E_T^+}\longrightarrow\frac\pi4
 =0.7853981633\ldots .
\end{equation}
\end{theorem}

\begin{proof}
By Lemma~\ref{lem:triangular-law}, the increasing empirical quantile functions of the two
normalized lists converge in $L^1$ to $Q(t)=\sqrt t$.  The discrete rearrangement sums are
Riemann sums for the monotone and antimonotone couplings.  Thus
\[
 \frac{E_T^-}{N_TTL_T\log2}
 \longrightarrow\int_0^1Q(t)Q(1-t)\,dt,
 \qquad
 \frac{E_T^+}{N_TTL_T\log2}
 \longrightarrow\int_0^1Q(t)^2\,dt.
\]
The first integral is the beta integral
$B(3/2,3/2)=\Gamma(3/2)^2/\Gamma(3)=\pi/8$; the second is $1/2$.
\end{proof}

The determinant has $N_T!$ Leibniz terms, each at most $e^{E_T^+}$.
Since $N_T\asymp T^2$ while $TL_TN_T\asymp T^4$,
\begin{equation}\label{eq:grid-factorial-lower-order}
 \log(N_T!)=O(T^2\log T)=o(N_TTL_T).
\end{equation}
Combining the preceding statements gives the promised ceiling.

\begin{corollary}[Full-grid termwise-divisor ceiling]\label{cor:grid-pi-four}
With the notation above,
\begin{equation}\label{eq:grid-pi-four}
 \limsup_{T\to\infty}
 \frac{\log Q_T^{\triangle}}
      {E_T^++\log(N_T!)}
 \le\frac\pi4.
\end{equation}
Thus the complete $\asymp T^2$ row grid does not make the all-prime common divisor exceed the
sharp largest-Leibniz-term archimedean bound.  Entrywise valuation alone leaves at least the
fraction $1-\pi/4=0.2146018\ldots$ uncovered at leading order.
\end{corollary}

\begin{proof}
Lemma~\ref{lem:all-place-transport} gives
$\log Q_T^{\triangle}\le E_T^-$.  Divide by the standard determinant upper bound
$E_T^++\log(N_T!)$ and apply Theorem~\ref{thm:grid-transport} and
\eqref{eq:grid-factorial-lower-order}.
\end{proof}

\begin{warning}[What the $\pi/4$ theorem does not exclude]\label{warn:grid-actual}
The denominator in \eqref{eq:grid-pi-four} is the largest-term upper bound, not the true size
of $D_T^{\triangle}$.  Generalized Vandermonde determinants exhibit substantial alternating
cancellation.  A sufficiently sharp archimedean estimate could lower the relevant budget
below $E_T^-$, in which case the termwise divisor might become decisive.  The theorem excludes
the naive ``more rows plus the same Hadamard/Leibniz estimate'' route; it does not exclude a
joint $p$-adic/free-energy argument.
\end{warning}

\subsection{The Harish--Chandra--Itzykson--Zuber identity}

The exact analytic object behind that warning is classical.  For diagonal Hermitian matrices
$X=\operatorname{diag}(x_1,\ldots,x_N)$ and
$Y=\operatorname{diag}(y_1,\ldots,y_N)$ with distinct increasing entries, normalized Haar
measure on the unitary group satisfies
\begin{equation}\label{eq:hciz}
 \int_{U(N)}
   \exp\!\bigl(\operatorname{Tr}(XUYU^*)\bigr)\,dU
 =\frac{\prod_{r=1}^{N-1}r!}{\Delta(x)\Delta(y)}
   \det(e^{x_i y_j})_{i,j=1}^{N},
\end{equation}
where
$\Delta(x)=\prod_{i<j}(x_j-x_i)$ and similarly for $\Delta(y)$
\cite{HarishChandra1957,ItzyksonZuber1980,McSwiggen2018}.  Hence
\begin{equation}\label{eq:hciz-det}
 \log D_T^{\triangle}
 =\log\Delta(x)+\log\Delta(y)
  -\sum_{r=1}^{N_T-1}\log(r!)
  +\log\mathcal I_T,
\end{equation}
with $\mathcal I_T$ the unitary integral in \eqref{eq:hciz}.

The large powers of $T$ hidden separately in the two Vandermonde factors and the factorial
product cancel at leading logarithmic scale.  What remains is a genuine free-energy problem
for two triangular empirical measures.  The Baker separation theorem of
Section~\ref{sec:baker} gives a polynomial lower bound for the spacings in $\Delta(x)$, so the
row Vandermonde cannot collapse through anomalously close conditional exponents.  What is not
known is a bound on the full combination \eqref{eq:hciz-det} strong enough to beat the
primewise divisor.

\begin{target}[HCIZ--valuation comparison]\label{target:hciz}
Determine the $T^4$-scale limit, or a sufficiently sharp upper bound, of
\[
 \log D_T^{\triangle}-\log Q_T^{\triangle}
\]
for the conditional triangular row lattice and balanced mixed columns.  It would suffice to
prove that this quantity is negative for all sufficiently large $T$.  Corollary
\ref{cor:grid-pi-four} shows that such a result must use either the HCIZ/Vandermonde
cancellation in \eqref{eq:hciz-det}, extra $p$-adic cancellation beyond the assignment minima,
or both.
\end{target}

This formulation has a useful advantage over a generic interpolation determinant: both
limiting measures and all lattice densities are explicit.  It is therefore suitable for
numerical free-energy reconnaissance before attempting a proof.

\section{The first tied \texorpdfstring{$p$}{p}-adic layer is lower order}
\label{sec:firstlayer}

Theorem~\ref{thm:four-fifths} says that extra divisibility must come from ties among
minimum-valuation Leibniz terms.  Ties are especially abundant in the cross-coprime branch
\eqref{eq:cross-coprime}: at $p=2$ all columns with the same $u$ have the same valuation
profile, and at $p=3$ all columns with the same $v$ do.  It is natural to hope that cancellation
inside those large fibers supplies the missing leading-order fraction.

The present section computes that first cancellation layer exactly.  It factors into ordinary
Vandermondes of the odd row parts.  There are two useful bounds.  Standard upper bounds for
$p$-adic linear forms in two logarithms give the sharper
$O(T^{3/2}(\log T)^C)$ estimate.  More importantly for formalization, the elementary fact
$p^{v_p(z)}\le |z|$ already gives an unconditional $O(T^2)$ bound.  Either is lower order
than the mixed determinant height $\asymp T^{5/2}$.  This does not bound the final determinant
valuation---higher tropical layers may cancel again---but it rules out a one-step fiber
explanation of the missing fraction without requiring deep $p$-adic approximation theory.

\subsection{Tropical initial form at two}

Assume throughout this section that $\gcd(MA,6)=1$.  For the block rows, write
\begin{equation}\label{eq:odd-output}
 q_k=a_{T,k}=3^{n_k}A^k,
 \qquad n_k=T-\lfloor k\beta\rfloor.
\end{equation}
Both $M$ and every $q_k$ are odd, and
\begin{equation}\label{eq:two-entry-factor}
 m_{T,k}^{u}a_{T,k}^{v}
 =2^{u n_k}\,M^{ku}q_k^v.
\end{equation}
For each integer $u\ge0$, let
\[
 V_u=\{v:(u,v)\text{ is one of the selected }N\text{ columns}\},
 \qquad s_u=\#V_u.
\]
Because columns are selected by the inequality $u+\vartheta v\le L_N$ up to one boundary
shell, each nonempty $V_u$ is an initial interval of integers.  Sort the distinct $u$ values
in increasing order.  Since $n_0>n_1>\cdots>n_{N-1}$, every minimizing assignment for the
cost $un_k$ sends one consecutive row block $R_u$ of size $s_u$ to the columns in the fiber
$\{u\}\times V_u$.

Let $\tau_{2,T}$ be the minimum from \eqref{eq:taup}.  The sum of all Leibniz terms having
exactly this valuation is the $2$-adic tropical initial form.  Up to an overall sign and an odd
factor, it is
\begin{equation}\label{eq:initial-two}
 I_{2,T}
 =\prod_u
   \det\!\left(q_k^v\right)_{k\in R_u,\ v\in V_u}
 =\prod_u\prod_{\substack{i<j\\i,j\in R_u}}(q_j-q_i).
\end{equation}
The first equality follows by grouping the optimal permutations by fibers and factoring the
odd number $M^{ku}$ from each row in a fixed $u$ block; the second is the ordinary Vandermonde
identity because $V_u$ is consecutive.

\begin{proposition}[Exact first $2$-adic initial form]\label{prop:initial-two}
There is an integer $Z_{2,T}$ such that
\begin{equation}\label{eq:two-layer-expansion}
 2^{-\tau_{2,T}}D_T
 =\varepsilon_{2,T} I_{2,T}+2^{\lfloor\beta\rfloor}Z_{2,T},
\end{equation}
where $\varepsilon_{2,T}$ is odd.  Consequently
\begin{equation}\label{eq:two-layer-min}
 v_2(D_T)-\tau_{2,T}
 \ge \min\{v_2(I_{2,T}),\lfloor\beta\rfloor\}.
\end{equation}
No equality assertion follows without controlling the valuation of the full outside-fiber
remainder $Z_{2,T}$.
\end{proposition}

\begin{proof}
The rearrangement inequality with repeated column weights identifies the minimizing
permutations exactly as above, and summing their signed unit parts gives
\eqref{eq:initial-two}.  Any nonminimizing assignment must move at least one column across two
adjacent $u$ fibers.  An exchange across the corresponding boundary changes the cost by
\[
 (u'-u)(n_i-n_j).
\]
Here $u'-u\ge1$, while consecutive row depths differ by
$\lfloor(k+1)\beta\rfloor-\lfloor k\beta\rfloor\ge\lfloor\beta\rfloor$.  Hence every
nonminimal Leibniz term has valuation at least $\tau_{2,T}+\lfloor\beta\rfloor$, proving the
expansion.  The valuation statement follows from the ultrametric inequality.
\end{proof}

\subsection{The symmetric initial form at three}

Put
\begin{equation}\label{eq:three-unit-input}
 r_k=m_{T,k}=2^{n_k}M^k.
\end{equation}
At $p=3$ the entry factors as
\[
 m_{T,k}^ua_{T,k}^v
 =3^{vn_k}\,r_k^uA^{kv}.
\]
For a selected value of $v$, let $U_v$ be the consecutive set of selected $u$ values and let
$R'_v$ be the corresponding consecutive block of rows in the minimum assignment.  The same
argument gives
\begin{equation}\label{eq:initial-three}
 I_{3,T}
 =\prod_v\prod_{\substack{i<j\\i,j\in R'_v}}(r_j-r_i)
\end{equation}
and
\begin{equation}\label{eq:three-layer-expansion}
3^{-\tau_{3,T}}D_T
 =\varepsilon_{3,T}I_{3,T}+3^{\lfloor\beta\rfloor}Z_{3,T},
\end{equation}
with $\varepsilon_{3,T}$ a $3$-adic unit.

\subsection{Two-logarithm bounds inside the fibers}

For $i<j$, set $d=j-i$ and
$h=n_i-n_j=\lfloor j\beta\rfloor-\lfloor i\beta\rfloor$.  The common factors in the row
values are units at the relevant place, so
\begin{align}
 v_2(q_j-q_i)&=v_2(A^d-3^h), \label{eq:qdiff-two}\\
 v_3(r_j-r_i)&=v_3(M^d-2^h). \label{eq:rdiff-three}
\end{align}
Neither difference vanishes: $A^d=3^h$ or $M^d=2^h$ would give $d\beta=h$, contrary to the
irrationality of $\beta$.

A standard consequence of $p$-adic linear-form theory is that for fixed multiplicatively
independent positive rationals $\alpha_1,\alpha_2$ and a fixed prime $p$, there are effectively
computable constants $C_0,C_1$ such that
\begin{equation}\label{eq:padic-two-log}
 v_p(\alpha_1^d-\alpha_2^h)
 \le C_0+C_1\bigl(\log(2+\max\{d,h\})\bigr)^2.
\end{equation}
Much sharper versions, with linear dependence on $\log B$, are known in important two-logarithm
regimes; the square is more than sufficient here
\cite{BugeaudLaurent1996,Bugeaud2002,Yu1998}.  Apply this to $(A,3,p=2)$ and
$(M,2,p=3)$.  Since $d,h=O_\beta(N)$, each factor in
\eqref{eq:initial-two}--\eqref{eq:initial-three} has valuation
$O_{M,A}((\log N)^2)$.

For the qualitative scale separation, however, no linear-form theorem is needed.  If $p$ is
prime and $a^d\ne b^h$, the divisor-size inequality gives the explicit finite bound
\begin{equation}
 v_p(a^d-b^h)
 \le 1+d\,\lceil\log_p a\rceil+h\,\lceil\log_p b\rceil.
 \label{eq:elementary-padic-power-difference}
\end{equation}
Indeed, $p^{v_p(a^d-b^h)}$ divides the nonzero difference, while
$|a^d-b^h|\le a^d+b^h$ and each power is bounded by the corresponding power of $p$.
Inside a Beatty row block, $h\le B_{\beta}d$ for a fixed integer $B_\beta$, so the right side
is $O_{M,A}(d)$.  A block of size $s$ then contributes $O_{M,A}(s^3)$, because
$\sum_{0\le i<j<s}(j-i)=O(s^3)$.  There are $O(L_N)$ fibers of size $O(L_N)$, hence the
total elementary budget is
\begin{equation}
 O_{M,A}(L_N^4)=O_{M,A}(N^2)=O_{M,A}(T^2).
 \label{eq:elementary-first-layer-budget}
\end{equation}

The number of pairs within the vertical fibers is
\begin{equation}\label{eq:vertical-pairs}
 \sum_u\binom{s_u}{2}
 =\frac{L_N^3}{6\vartheta^2}+O_\vartheta(L_N^2)
 =O_\vartheta(N^{3/2}),
\end{equation}
while the number within the horizontal fibers is
\begin{equation}\label{eq:horizontal-pairs}
 \sum_v\binom{\#U_v}{2}
 =\frac{L_N^3}{6\vartheta}+O_\vartheta(L_N^2)
 =O_\vartheta(N^{3/2}).
\end{equation}
We obtain the following rigorous scale separation.

\begin{theorem}[First tied layer is lower order]\label{thm:first-layer-small}
In the cross-coprime branch,
\begin{equation}\label{eq:first-layer-small}
 v_2(I_{2,T})+v_3(I_{3,T})
 =O_{M,A}(N^2)=O_{M,A}(T^2).
\end{equation}
If one invokes the quoted two-logarithm estimate \eqref{eq:padic-two-log}, the stronger bound
$O_{M,A}(N^{3/2}(\log N)^2)$ also holds.
In particular,
\begin{equation}\label{eq:first-layer-negligible}
 \frac{v_2(I_{2,T})\log2+v_3(I_{3,T})\log3}{H_T}
 \longrightarrow0.
\end{equation}
\end{theorem}

\begin{proof}
For the elementary assertion, sum \eqref{eq:elementary-padic-power-difference} over the pairs
in each fiber and use \eqref{eq:elementary-first-layer-budget}.  For the sharper assertion,
sum \eqref{eq:padic-two-log} over the pair counts
\eqref{eq:vertical-pairs} and \eqref{eq:horizontal-pairs}, using
\eqref{eq:qdiff-two}--\eqref{eq:rdiff-three}.  Finally use $N\asymp_\beta T$ and
$H_T\asymp_\beta T^{5/2}$.
\end{proof}

\begin{statusbox}{\statuslean{} for the exact finite valuation and cube-budget skeleton;
\statusknown{} only for the optional sharper two-logarithm estimate.}
The module \lean{BlockVandermondeValuationBudget} proves the exact collision-depth identity
\lean{padicValNat_natVandermondeProduct_eq_sum_collisionCount}, the elementary power-difference
bound \lean{padicValNat_natAbs_pow_sub_pow_le}, and the variable-block and power-difference cube bounds
\lean{padicValNat_blockVandermondeNatProduct_le_cubeBudget} and
\lean{padicValNat_blockVandermondeNatProduct_le_powDifferenceCubeBudget}.  The displayed
$O(T^2)$ specialization still uses the elementary paper-level count of Beatty fiber sizes;
the finite inequality on which it rests is kernel-verified.  The sharper
$O(T^{3/2}(\log T)^2)$ statement uses the cited classical estimate and is not claimed as a
Lean theorem.
\end{statusbox}

\subsection{What a leading excess would have to do}

Theorem~\ref{thm:first-layer-small} does not prove
$v_2(D_T)-\tau_{2,T}=o(H_T)$ or its $3$-adic analogue.  If the initial form is divisible by a
power exceeding the first tropical gap, terms from the next valuation layer enter; those can
cancel, exposing another layer, and so on.  A leading-order excess is therefore possible only
through control of the valuations and cancellations in the truncated sums at cutoffs growing
with $T$.  It could arise from unusually divisible low-excess minor products, cancellation
within one layer, or carrying across several layers.  Cancellation confined to the
minimum-permutation fibers is quantitatively negligible.
Formula~\eqref{mp:eq-row-block-truncation} makes this statement finite and exact: at depth $K$
one need only control the signed products of unit minors attached to row partitions of excess
less than $K$.  The unit-adapted bases above extract a universal divisor from the relevant
univariate minors, but its paper-level blockwise scale is only $O(N^{3/2})$.  What is missing is
an upper bound on the valuation of the residual truncated sum, uniform for cutoffs growing
with $T$.

This yields a more precise version of the $p$-adic research target.

\begin{target}[No shallow-layer proof]\label{target:cascade}
Either prove that the tropical cancellation cascade has total depth $o(H_T)$---which would
retire the single-block $p$-adic strategy---or identify a structural identity forcing
$\gg H_T$ cumulative cancellation across successively nonminimal assignments.  Pairwise
congruences such as \eqref{eq:qdiff-two} and \eqref{eq:rdiff-three} cannot by themselves
supply the latter, because $p$-adic logarithm bounds make each congruence only polylogarithmic
in the row separation.
\end{target}

\subsection{Exact finite cascade reconnaissance}
\label{sec:cascade-experiments}

Exact integer experiments help rule out two tempting but false simplifications of
Target~\ref{target:cascade}.  For synthetic cross-coprime parameters $M,A$, put
\[
 \beta=\log_2M,\qquad n_k=T-\lfloor k\beta\rfloor,
\]
and take the first $N=1+\lfloor(T+1)/\beta\rfloor$ exponent pairs $(u,v)$ in increasing order
of $u+\thetaq v$.  The tested determinant is exactly
\[
 D_T=\det\left(
   (2^{n_k}M^k)^u(3^{n_k}A^k)^v
 \right)_{0\le k<N, (u,v)}.
\]
Let $e_p=v_p(D_T)-\tau_p$, let $I_p$ denote the first-fiber Vandermonde product, and let $U_p$
be the common row-block ordering exponent determined by the column-fiber sizes.  The weights
come from \lean{AllLayerUnitBasis}, their finite sums from \lean{UnitOrderingGain}, and their
uniform contribution to every row-block term and the full determinant from
\lean{RowBlockUnitDivisibility}.  For $(M,A)=(5,7)$, exact DVR elimination
gives:
\begin{center}
\small
\begin{tabular}{@{}r r|r r r|r r r|c@{}}
\toprule
$T$ & $N$ & $e_2$ & $U_2$ & $v_2(I_2)$ & $e_3$ & $U_3$ & $v_3(I_3)$ &
$(e_2\log2+e_3\log3)/H_T$ \\
\midrule
30  & 14  & 26  & 17  & 29  & 14  & 9   & 14  & $0.0360124$ \\
100 & 44  & 157 & 136 & 157 & 87  & 74  & 87  & $0.0101461$ \\
300 & 130 & 868 & 796 & 870 & 517 & 471 & 530 & $0.00353027$ \\
\bottomrule
\end{tabular}
\end{center}
The data are compatible with an $N^{3/2}$ excess and incompatible with no proved asymptotic;
their purpose is diagnostic.

Two exact small cases are more informative.  At $(M,A,T,p)=(5,7,30,2)$ one has
\[
 \tau_2=177,\qquad v_2(D_T)=203,
 \qquad e_2=26<29=v_2(I_2).
\]
Writing $2^{-\tau_2}D_T=\sum_\delta2^\delta C_\delta$, the first layer has
$v_2(C_0)=29$, whereas the layer $\delta=6$ has $v_2(C_6)=20$ and therefore smaller effective
valuation $26$.  A nonminimal layer can genuinely beat the initial form.  Conversely, at
$(M,A,T,p)=(11,5,20,2)$ the layers $\delta=0$ and $\delta=4$ both have effective valuation
$9$, but their sum is divisible by $2^{10}$; all other layers have effective valuation at
least $12$, and the final excess is $10$.  Thus even ``take the minimum effective layer'' is
false because different layers can carry into one another.

\begin{statusbox}{\statusexactcomp{} --- all displayed determinant, assignment, layer, and valuation
integers use exact integer arithmetic; the logarithmic ratios are ordinary floating
approximations used only diagnostically.  The no-dependency scripts
\texttt{scripts/cascade\_experiments.py},
\texttt{scripts/cascade\_layer\_report.py}, and
\texttt{scripts/cascade\_verify.py} reproduce the computations; the last directly
Bareiss-checks the same twenty small determinants and agrees on their valuations.}
The parameters are deliberately synthetic: they satisfy $\gcd(6,MA)=1$, but do not assert the
unknown identity $A=3^{\log_2M}$.  The decreasing ratios and the apparent $N^{3/2}$ scale are
evidence only.  The two finite counterexamples to naive layer formulas are exact for the
displayed mixed determinant family.
\end{statusbox}

\section{Synthesis of the determinant audits}
\label{sec:synthesis}

The preceding calculations put several superficially different determinant proposals into the
same currency: the fraction of the leading archimedean budget supplied automatically by
valuations of the entries.  The numerical constants should not be compared as though the
underlying determinants were identical---the ordinary, mixed-block, and complete-grid systems
have different sizes and different upper bounds---but they do show which mechanism is closest
to a contradiction before any genuine cancellation is used.

\begin{table}[htbp]
\centering
\caption{Leading termwise-divisor ceilings obtained in this continuation.}
\label{tab:ceilings}
\begin{tabularx}{\textwidth}{@{}>{\raggedright\arraybackslash}p{0.29\textwidth}
 >{\raggedright\arraybackslash}p{0.27\textwidth}
 >{\raggedright\arraybackslash}X@{}}
\toprule
System & Ceiling or exact limit & Meaning \\
\midrule
Ordinary input/output Vandermonde product &
$1-\dfrac{\delta_2\log2+\delta_3\log3}{3\beta\log6}
 <0.8710491$ &
Exact fixed-support divisor fraction; all primes already occurring in $M$ and $A$ included.\\[0.7em]
Mixed single block, arbitrary branch &
$1-\dfrac38\sum_p\mathfrak G(c_p)\le4/5$ &
Exact all-prime tropical limit and universal primitive-depth ceiling.\\[0.7em]
Mixed single block, cross-coprime &
$\le7/10$ &
Uses $c_2=(-1,0)$, $c_3=(0,-1)$ and the external sum $(1,1)$.\\[0.7em]
Mixed single block, only structural primes in the cross-coprime branch &
$3/10$ &
Exact contribution of $p=2,3$ before external primes or cancellation.\\[0.7em]
Complete triangular row grid &
$\le\pi/4=0.7853982\ldots$ &
All-prime common divisor versus the sharp largest-Leibniz-term transport upper bound.\\[0.7em]
First tied fibers at $2$ and $3$ &
$o(1)$ of the mixed-block budget &
Verified finite budgets plus elementary paper-level Beatty counting give $O(T^2)$ against
$H_T\asymp T^{5/2}$; quoted two-logarithm estimates sharpen this to
$O(T^{3/2}(\log T)^2)$.\\
\bottomrule
\end{tabularx}
\end{table}

Three conclusions survive every normalization.

\begin{enumerate}
\item Merely adding every prime in the fixed support does not close any determinant.  It does,
      however, improve the diagnosis substantially: the universal mixed-block ceiling is
      $4/5$, not a heuristic ``about one half''.
\item The exact primitive theorem is quantitatively useful.  The condition
      $\min(v_2(M),v_3(A))=0$ creates one structural slope vector of norm at least $4/15$; the
      conservation law forces a second norm contribution of at least the same size.
\item A successful $p$-adic proof must be a cancellation theorem, not a divisibility-by-every-
      term theorem.  In the easiest branch, even the first large tied family contributes only
      a lower-order amount.  The required phenomenon would have to repeat through many
      tropical layers or interact with a much sharper archimedean free energy.
\end{enumerate}

\subsection{Correction of the post-attachment ``one half / two thirds'' note}
\label{subsec:report7-correction}

After the attached unified report was produced, the repository acquired a short exploratory
note (merged into Section~\ref{sec:minplus} of this report) at commit
\sha{ff9a9f5aa3a3f02599de93fc053e60918f269855}.  That note correctly identified the distinction
between termwise min-plus divisibility and cancellation beyond the tropical minimum, and its
real/$p$-adic compatibility diagnosis remains valuable.  It proposed, however, a structural
$1/2$ ceiling and an all-prime $2/3$ ceiling.  Neither constant is valid in the stated
generality.

The structural averaging argument omitted the base and cross depths that appear in
\eqref{eq:c2}--\eqref{eq:c3}.  Even under $v_2(M)=0$, the terms
$v_2(A)kv$ and $v_3(M)ku+v_3(A)kv$ do not disappear.  The exact structural contribution is
obtained by retaining $c_2$ and $c_3$ in Theorem~\ref{thm:mixed-gini}; it can exceed one half
of $H_T$.

The all-prime step used the assertion that an oppositely sorted nonnegative triangular weight
has assignment minimum at most two thirds of its mean.  For the weight $Y=X+Z$ on the standard
triangle, the exact values from Lemma~\ref{prop:G-explicit} are
\[
 \mathbb E Y=\frac23,
 \qquad
 \Phi(1,1)=\frac4{15},
 \qquad
 \frac{\Phi(1,1)}{\mathbb EY/2}=\frac45.
\]
Thus $4/5$, not $2/3$, is already attained by the relaxed external rank-one grading.  The
Gini-norm identity explains the general case and leads to the rigorous ceiling in
Theorem~\ref{thm:four-fifths}.  This is a correction of a calculation in an exploratory note,
not a criticism of the source report itself; the attached unified report did not contain the
claim.

\begin{statusbox}{\statusnew{} --- the correction follows from the exact transport formula,
not from numerical experimentation.}
\end{statusbox}


\section{Directed-rounding extension through \texorpdfstring{$2^{27}$}{2\^{}27}}
\label{sec:finite27}

The attached report supplied a complete MPFR scan through $2^{26}$.  Because a finite
extension is independent of the unresolved transcendence barrier and gives a useful stronger
floor for every hypothetical generator, I extracted the literal C listing from its
Appendix~``Directed-rounding verifier'', compiled it against MPFR~4.2.2, replayed an old audit
chunk, and scanned the next full dyadic range.

\subsection{Manifest audit}

The literal extracted source has SHA-256 digest
\begin{center}
 \sha{f888c0d89be9366f01e2d26f339dcbe446556d39d4c5dbcf8f579ec0d0052e1d}.
\end{center}
This does \emph{not} equal the digest
\begin{center}
\sha{ec709623f812478fbb7fb5567aa59fbfe3df0ae639ca366bc7c72e278bcfee2f}
\end{center}
printed in the source of that section.  None of the ordinary line-ending, final-newline, or tab
normalizations explains the discrepancy.  The report therefore contains a stale or
externally generated source-identity hash.

The semantic reproducibility check succeeds.  Replaying the old interval
$[29{,}360{,}128,33{,}554{,}432)$ at $192$ bits gives exactly the report's records
\[
 \begin{split}
 &\text{lower: }1.4439124872349507\times10^{-19}
   \text{ at }(30{,}103{,}832,713{,}403{,}245{,}915),\\
 &\text{upper: }5.4132120119800561\times10^{-21}
   \text{ at }(29{,}411{,}704,687{,}581{,}893{,}443),
 \end{split}
\]
with zero failures.  Thus the listing used here reproduces the old executable's mathematical
behavior on an audited four-million-input overlap even though the printed file digest does not
identify it.

\begin{statusbox}{\statuscomp{} --- this section is a directed-rounding software result.  It
is stronger numerically than the Lean theorem but is not a kernel proof.  The manifest hash
mismatch is explicitly part of the trust boundary.}
\end{statusbox}

\subsection{The new range}

The half-open interval $[2^{26},2^{27})$ was divided into sixteen chunks of width $2^{22}$.
The executable used $192$-bit MPFR endpoints and outward rounding exactly as in the source
report.  There are $2^{26}$ integers in the range; the first, $2^{26}$, is the sole power of
two and is intentionally skipped.  Hence the required nonpower count is
$2^{26}-1=67{,}108{,}863$.

The concatenated canonical extension log has SHA-256 digest
\begin{center}
 \sha{b11843448cd8946442aba13758b5b666e4ac6eae50d2243fbdc32d83274cea8d}.
\end{center}
Every chunk reports zero failures, and the checked counts sum to the required total.

{\small
\begin{longtable}{@{}r p{0.27\textwidth} r r r@{}}
\caption{Directed-rounding extension on $[2^{26},2^{27})$.}
\label{tab:scan27}\\
\toprule
Chunk & Integer range & Checked & Smallest lower gap & Smallest upper gap \\
\midrule
\endfirsthead
\toprule
Chunk & Integer range & Checked & Smallest lower gap & Smallest upper gap \\
\midrule
\endhead
0 & 67{,}108{,}864--71{,}303{,}167 & 4{,}194{,}303 & $4.2837\times10^{-20}$ & $1.7723\times10^{-19}$ \\
1 & 71{,}303{,}168--75{,}497{,}471 & 4{,}194{,}304 & $2.2961\times10^{-20}$ & $1.1433\times10^{-20}$ \\
2 & 75{,}497{,}472--79{,}691{,}775 & 4{,}194{,}304 & $1.2813\times10^{-19}$ & $1.6063\times10^{-20}$ \\
3 & 79{,}691{,}776--83{,}886{,}079 & 4{,}194{,}304 & $2.5420\times10^{-20}$ & $5.2012\times10^{-20}$ \\
4 & 83{,}886{,}080--88{,}080{,}383 & 4{,}194{,}304 & $2.6166\times10^{-20}$ & $3.8496\times10^{-20}$ \\
5 & 88{,}080{,}384--92{,}274{,}687 & 4{,}194{,}304 & $1.0923\times10^{-19}$ & $5.7144\times10^{-20}$ \\
6 & 92{,}274{,}688--96{,}468{,}991 & 4{,}194{,}304 & $1.1017\times10^{-20}$ & $2.3227\times10^{-20}$ \\
7 & 96{,}468{,}992--100{,}663{,}295 & 4{,}194{,}304 & $3.0759\times10^{-20}$ & $7.4810\times10^{-20}$ \\
8 & 100{,}663{,}296--104{,}857{,}599 & 4{,}194{,}304 & $5.6465\times10^{-20}$ & $6.4744\times10^{-20}$ \\
9 & 104{,}857{,}600--109{,}051{,}903 & 4{,}194{,}304 & $8.8144\times10^{-20}$ & $8.3662\times10^{-21}$ \\
10 & 109{,}051{,}904--113{,}246{,}207 & 4{,}194{,}304 & $3.8892\times10^{-20}$ & $7.0327\times10^{-21}$ \\
11 & 113{,}246{,}208--117{,}440{,}511 & 4{,}194{,}304 & $9.9365\times10^{-21}$ & $2.1814\times10^{-20}$ \\
12 & 117{,}440{,}512--121{,}634{,}815 & 4{,}194{,}304 & $4.4964\times10^{-20}$ & $5.4132\times10^{-21}$ \\
13 & 121{,}634{,}816--125{,}829{,}119 & 4{,}194{,}304 & $1.0774\times10^{-20}$ & $3.4093\times10^{-20}$ \\
14 & 125{,}829{,}120--130{,}023{,}423 & 4{,}194{,}304 & $4.3609\times10^{-20}$ & $8.3073\times10^{-22}$ \\
15 & 130{,}023{,}424--134{,}217{,}727 & 4{,}194{,}304 & $5.9082\times10^{-21}$ & $5.1326\times10^{-20}$ \\
\midrule
\multicolumn{2}{@{}r}{\textbf{total}} & \textbf{67{,}108{,}863} & & \\
\bottomrule
\end{longtable}
}

\subsection{New extremal records and high-precision replay}

The smallest lower logarithmic margin in the extension is
\begin{equation}\label{eq:new-lower-record}
 5.9082259992879571\times10^{-21}
 \quad\text{at}\quad
 (m,n)=(132{,}833{,}765,7{,}501{,}348{,}219{,}569).
\end{equation}
A $120$-decimal-digit reevaluation gives
\[
 m^{\vartheta}-n
 =6.3939754533470409837680125510\ldots\times10^{-8}.
\]

The smallest upper logarithmic margin is the new global record
\begin{equation}\label{eq:new-upper-record}
 8.3073231820333117\times10^{-22}
 \quad\text{at}\quad
 (m,n)=(129{,}799{,}345,7{,}231{,}571{,}292{,}851),
\end{equation}
where
\[
 n+1-m^{\vartheta}
 =8.6669904355818956853297952288\ldots\times10^{-9}.
\]
Both one-input cases were rerun independently at $320$ and $512$ bits.  The directed margins
and nearest integer were unchanged in every displayed digit, and both replays reported zero
failures.

\begin{theorem}[Software-level finite exclusion through $2^{27}$]
\label{thm:finite27}
At the stated MPFR software-trust level, for every integer $1\le m<2^{27}$,
\[
 m^{\vartheta}\in\N
 \quad\Longrightarrow\quad
 m\text{ is a power of two}.
\]
Consequently every nonintegral real $x$ satisfying $2^x,3^x\in\Z$ obeys
\begin{equation}\label{eq:x-gt-27}
 x>27.
\end{equation}
\end{theorem}

\begin{proof}
The attached report's complete scan covers $1\le m<2^{26}$, and the sixteen new chunks cover
$2^{26}\le m<2^{27}$.  In each case the directed logarithmic interval for
$\vartheta\log m$ lies strictly between the intervals for the logarithms of two consecutive
integers, except when $m$ is a power of two.  If $m=2^x$ came from a nonintegral solution and
$m<2^{27}$, then $m^\vartheta=3^x$ would be an integer at a non-power input, contradicting the
scan.  The endpoint $m=2^{27}$ corresponds to the integral exponent $27$, so a nonintegral
solution must have $m>2^{27}$ and hence $x>27$.
\end{proof}

This raises the software-level spacing ratio in the optimized finite-difference comparison
from $26\log3$ to $27\log3=29.662531\ldots$.  It does not change any kernel-verified theorem;
the formal floor remains $x\ge12$ until the finite certificates are replayed in Lean.

\section{Additional proof routes pushed to their failure points}
\label{ct:attempts}

The determinant calculations were the main new line of attack, but they were not the only one.
This section records four other routes in enough detail to prevent the same hidden loss from
being rediscovered under different notation.

\subsection{Transferring convergents of \texorpdfstring{$\beta$}{beta} to
\texorpdfstring{$\vartheta$}{theta}}
\label{subsec:transfer}

Let $p/q$ be a continued-fraction convergent of a hypothetical primitive generator $\beta$ and
put
\[
 \varepsilon=q\beta-p.
\]
For all sufficiently large $q$, $|\varepsilon|<1/2$.  Define the two nonzero integer gaps
\begin{equation}\label{eq:D2D3}
 D_2=M^q-2^p=2^p(2^\varepsilon-1),
 \qquad
 D_3=A^q-3^p=3^p(3^\varepsilon-1).
\end{equation}
They manufacture an exact rational approximation to $\vartheta$:
\begin{equation}\label{eq:theta-transfer}
 R_{p,q}:=
 \frac{2^pD_3}{3^pD_2}
 =\frac{3^\varepsilon-1}{2^\varepsilon-1}
 \in\Q.
\end{equation}
Taylor expansion at zero gives
\begin{equation}\label{eq:theta-transfer-error}
 R_{p,q}=\vartheta+
 \frac{\log3\,(\log3-\log2)}{2\log2}\,\varepsilon
 +O(\varepsilon^2),
\end{equation}
so $|R_{p,q}-\vartheta|\ll|\varepsilon|$.

This looks promising because $\vartheta$ has an explicit irrationality measure.  Wu's
simultaneous linear-independence estimate for $1,\log2,\log3,\log5$ gives, after setting the
unused coefficients to zero, a conservative exponent
\begin{equation}\label{eq:wu-measure}
 \left|\vartheta-\frac rs\right|
 \ge c_\vartheta s^{-\mu_\vartheta},
 \qquad \mu_\vartheta=8.616
\end{equation}
for all sufficiently large denominators $s$ \cite{Wu2003}.  The obstacle is the denominator
created by \eqref{eq:theta-transfer}.  Before reduction it is at most
\begin{equation}\label{eq:transfer-denominator}
 3^p|D_2|
 \ll 3^p2^p|\varepsilon|
 =6^p|\varepsilon|.
\end{equation}
Combining \eqref{eq:theta-transfer-error}--\eqref{eq:transfer-denominator} with
\eqref{eq:wu-measure} gives only
\begin{equation}\label{eq:transfer-lower}
 |\varepsilon|
 \gg 6^{-p\mu_\vartheta/(1+\mu_\vartheta)}.
\end{equation}
Since
\[
 6^{-\mu_\vartheta/(1+\mu_\vartheta)}
 =2^{-2.316\ldots},
\]
this is substantially weaker than the elementary bound
$|\varepsilon|\gg2^{-p}$ obtained directly from $D_2\ne0$.  Sharper real irrationality
measures cannot repair the architecture: every irrationality exponent is at least two, and
the unreduced denominator already contains both $2^p$ and $3^p$.

\begin{proposition}[Exact denominator-loss diagnosis]\label{prop:transfer-loss}
The convergent-transfer construction can improve the elementary gap only if one proves an
exponential lower bound for
\begin{equation}\label{eq:transfer-gcd}
 G_{p,q}=\gcd(2^pD_3,3^pD_2)
\end{equation}
that survives reduction of $R_{p,q}$.  In the absence of such a bound, every known
irrationality measure for $\vartheta$ is defeated by the factor $6^p$ in
\eqref{eq:transfer-denominator}.
\end{proposition}

No mechanism forcing $G_{p,q}$ to be exponentially large was found.  The two gaps share the
same small real parameter $\varepsilon$, but real closeness does not impose a common prime
divisor.

\subsection{Two-row determinants and fixed-support \texorpdfstring{$S$}{S}-units}

For two block rows $i<j$, put $d=j-i$ and
$h=n_i-n_j=\lfloor j\beta\rfloor-\lfloor i\beta\rfloor$.  Their $2\times2$ determinant has
the exact factorization
\begin{equation}\label{eq:two-row-factor}
 m_{T,i}a_{T,j}-m_{T,j}a_{T,i}
 =6^{n_j}(MA)^i\bigl(2^hA^d-3^hM^d\bigr).
\end{equation}
Moreover
\begin{equation}\label{eq:two-row-ratio}
 \frac{2^hA^d}{3^hM^d}
 =\left(\frac32\right)^{\beta d-h}.
\end{equation}
Thus a good rational approximation $h/d$ to $\beta$ makes the bracket in
\eqref{eq:two-row-factor} a small \emph{relative} difference of two integers supported on the
fixed prime set dividing $6MA$.

Let
\[
 U=2^hA^d,
 \qquad V=3^hM^d,
 \qquad G=\gcd(U,V).
\]
The reduced nonzero integer $(U-V)/G$ gives
\begin{equation}\label{eq:sunit-gap}
 |\beta d-h|
 \gg_{M,A}\frac{G}{\max\{U,V\}}.
\end{equation}
The exponent of the right-hand side can be written explicitly as the positive part of a
valuation-vector difference:
\begin{equation}\label{eq:sunit-rate}
 \log\frac{U}{G}
 =\sum_p
 \max\!\left
 \{h\mathbf1_{p=2}+d v_p(A)
   -h\mathbf1_{p=3}-d v_p(M),0\right\}\log p.
\end{equation}
As $h/d\to\beta$, this is $\kappa_{M,A}d+o(d)$ for a positive constant
$\kappa_{M,A}$.  Hence \eqref{eq:sunit-gap} is merely exponential in $d$, much weaker than the
Baker--Matveev polynomial separation of Section~\ref{sec:baker}.

The $abc$ conjecture applied to the reduced equation $U/G-V/G=(U-V)/G$ would make the
difference nearly as large as a fixed positive power of $U/G$, but it would still allow a
relative gap $e^{-\epsilon d}$.  Continued-fraction gaps of order $1/d$ are much larger.
Thus this direct $S$-unit equation does not become a proof even after inserting $abc$; one
would need an additional theorem controlling the prime support or radical of the difference.

\subsection{Automaticity and the synchronized Sturmian word}

The coding $s_k=\lfloor(k+1)\alpha\rfloor-\lfloor k\alpha\rfloor$ from
Section~\ref{sec:sturmian} is aperiodic Sturmian.  A standard theorem in automatic-sequence
theory says that no aperiodic Sturmian word is automatic in any integer base
\cite{Cobham1972,AlloucheShallit2003}.  One quick obstruction is frequency: the frequency of
$1$ in $s$ is the irrational number $\alpha$, whereas an automatic sequence having a letter
frequency has rational frequency.

This observation blocks a tempting but invalid shortcut.  The rationality of
$U=2^\alpha$ does \emph{not} make the threshold itinerary of
$R\mapsto UR/2^{\mathbf1_{UR\ge2}}$ finite-state; its exact rational denominators grow with
time.  The same is true for $V$ in base three.  Therefore Cobham's theorem cannot be invoked
merely because the same word is produced by a dyadic and a triadic rational system.  A viable
automata proof would first have to manufacture genuine finite kernels in two multiplicatively
independent bases, and doing so is already a new arithmetic theorem about $U$ and $V$.

\subsection{Real and \texorpdfstring{$p$}{p}-adic logarithms}

At the real place the counterexample is exactly
\begin{equation}\label{eq:real-log-wall}
 \log M\,\log3-\log A\,\log2=0.
\end{equation}
At $p=2$ or $3$, after stripping valuation and torsion parts, the units among
$M,A,2,3$ admit $p$-adic logarithms.  The determinant entries then have unit parts whose
$p$-adic logarithms are bilinear forms in the row and column coordinates.  This is the natural
language for the cancellation cascade isolated in Section~\ref{sec:firstlayer}.

What is missing is a transfer principle from \eqref{eq:real-log-wall} to a congruence or rank
drop among those $p$-adic logarithms.  Real exponentiation by the transcendental number
$\beta$ does not define a compatible $p$-adic exponent, and the equality $M=2^\beta$ has no
$p$-adic meaning.  Known $p$-adic linear-form theorems control a nonzero form once its integer
coefficients and algebraic bases are specified; they do not create a $p$-adic form from a real
one.

\begin{target}[Real/$p$-adic compatibility]\label{target:real-padic}
Find an algebraic construction, valid for the integer pair $(M,A)$ alone, that converts
\eqref{eq:real-log-wall} into either
\begin{enumerate}
\item a nontrivial congruence among $p$-adic unit logarithms with precision $\gg T$, repeated
      over $\gg T^{3/2}$ fiber pairs; or
\item a rank drop in the tropical initial matrices persisting through $\gg T$ valuation
      layers.
\end{enumerate}
Without such a construction, ordinary $p$-adic logarithm bounds point in the opposite
direction: they show that individual congruences are only polylogarithmically deep.
\end{target}

\subsection{Transcendental-curve point counting}

The graph $y=x^\vartheta$ is definable in the real exponential field, so Pila--Wilkie type
results~\cite{PilaWilkie2006} give subpolynomial bounds for rational points outside algebraic pieces.  Those generic
bounds are far too weak here.  The curve already contains the infinite sequence
$(2^n,3^n)$, giving $\asymp\log H$ integer points of height at most $H$, while a hypothetical
second generator produces $\asymp(\log H)^2$ points globally.  A proof by point counting would
therefore need a sharp $O(\log H)$ theorem with a structural classification of the extremal
sequence, not merely $O(H^\epsilon)$.  Such a theorem would essentially be another form of the
rank-one conclusion sought by the conjecture.

\section{Exact sufficient conditions extracted from the new analysis}
\label{sec:sufficient}

A useful continuation should not stop at saying that the present determinant fails.  The
calculations above isolate numerical thresholds which, if crossed by a future theorem, would
settle the conjecture.  This section records those thresholds as precise implications.

\subsection{The mixed-block cancellation-excess criterion}

Let
\[
 \mathcal P=\{p:p\mid6MA\}
\]
be the fixed support of all determinant entries, and retain the mixed block determinant $D_T$,
its tropical minima $\tau_{p,T}$, and the archimedean scale $H_T$.  Define the support-prime
cancellation excess
\begin{equation}\label{eq:cancellation-excess}
 \mathcal E_T
 =\sum_{p\in\mathcal P}
   \bigl(v_p(D_T)-\tau_{p,T}\bigr)\log p
 \ge0
\end{equation}
and the exact tropical deficit
\begin{equation}\label{eq:tropical-deficit}
 \delta(M,A)=\frac38\sum_{p\in\mathcal P}\mathfrak G(c_p).
\end{equation}
Theorem~\ref{thm:mixed-gini} says
$\log Q_T^{\rm trop}=(1-\delta(M,A)+o(1))H_T$, while the elementary Leibniz upper bound is
$\log D_T\le H_T+o(H_T)$.

\begin{theorem}[Cancellation-excess criterion]\label{thm:excess-criterion}
If a hypothetical primitive counterexample satisfies
\begin{equation}\label{eq:excess-threshold}
 \liminf_{T\to\infty}\frac{\mathcal E_T}{H_T}
 >\delta(M,A),
\end{equation}
then no such counterexample exists.  In particular it is enough to prove an excess greater
than $H_T/5$ uniformly in all primitive branches; in the cross-coprime branch the exact
threshold is at least $3H_T/10$.
\end{theorem}

\begin{proof}
The selected support primes alone give
\[
 \log D_T
 \ge\sum_{p\in\mathcal P}v_p(D_T)\log p
 =\log Q_T^{\rm trop}+\mathcal E_T.
\]
Under \eqref{eq:excess-threshold}, the right-hand side exceeds $H_T$ by a fixed positive
multiple for all large $T$, contradicting the Leibniz upper bound.  The numerical thresholds
follow from Theorem~\ref{thm:four-fifths} and Corollary~\ref{cor:seven-tenths}.
\end{proof}

This criterion is not tautological: it involves only valuations at the fixed finite set of
primes already present in $6MA$, whereas $D_T$ can acquire arbitrary new prime factors.  It
states exactly how much same-support cancellation is required.

\subsection{A denominator-cancellation criterion}

For the transferred approximant $R_{p,q}$ of Section~\ref{subsec:transfer}, let $s_{p,q}$ be
its reduced denominator.  If one could prove along infinitely many convergents that
\begin{equation}\label{eq:poly-reduced-denominator}
 s_{p,q}\ll q^\eta
 \qquad\text{for some }\eta<\frac{1+\mu_\vartheta}{\mu_\vartheta},
\end{equation}
then Wu's measure would give
$|R_{p,q}-\vartheta|\gg q^{-\eta\mu_\vartheta}$, whereas
\eqref{eq:theta-transfer-error} and the convergent bound give
$|R_{p,q}-\vartheta|\ll q^{-1}$.  Since
$\eta\mu_\vartheta<1+\mu_\vartheta$ is not by itself enough, one tracks
$|\varepsilon|<1/q_{+}$ more precisely; a contradiction follows whenever
$q_{+}\gg q^{\eta\mu_\vartheta+\epsilon}$ along the same subsequence.  A simpler sufficient
condition is
\begin{equation}\label{eq:very-small-reduced-denominator}
 s_{p,q}=q^{o(1)}
\end{equation}
for infinitely many convergents with $q_{+}\ge q^{1+\epsilon}$.

Equivalently, the gcd $G_{p,q}$ in \eqref{eq:transfer-gcd} would have to cancel essentially the
entire exponential factor $6^p|\varepsilon|$, leaving a subpolynomial denominator.  This is a
sharp formulation of the otherwise vague request for ``correlation between the two integer
gaps.''

\subsection{An automaticity criterion}

The synchronized coding gives another exact, if presently inaccessible, implication.

\begin{proposition}[Finite-kernel criterion]\label{prop:automatic-criterion}
If under a hypothetical counterexample the common word
$s_k=\lfloor(k+1)\alpha\rfloor-\lfloor k\alpha\rfloor$ were automatic in any integer base
$q\ge2$, then the Alaoglu--Erd\H{o}s conjecture would follow.
\end{proposition}

\begin{proof}
An automatic Sturmian word is ultimately periodic.  Its slope, equal to the frequency of the
letter $1$, is therefore rational.  But $\alpha=\beta-\lfloor\beta\rfloor$ would then be
rational, and the rational-exponent theorem would make $\beta$ integral.
\end{proof}

The missing step is not symbolic dynamics; it is proving automaticity from the two rational
multipliers $U$ and $V$.  The growing denominators in \eqref{eq:normalized-orbits} show exactly why the
obvious state space is infinite.

\section{Branch-sensitive research program}
\label{ct:program}

The source report already classifies a hypothetical pair much more sharply than ``two
integers on a curve.''  The next attack should exploit that classification rather than seek a
single estimate that treats every branch identically.

\subsection{Branch I: independent external valuation directions}

Here the second-iterate kernel $K_\beta$ is zero, equivalently the external prime-valuation
vectors of $M$ and $A$ are independent.  Corollary~\ref{cor:gini-kernel} gives a strict
angular defect $\Delta_{\rm ext}>0$ in the all-prime tropical fraction.  That makes the
termwise divisor farther from the archimedean upper bound, but the same independence also
supplies at least two genuinely different external primes at which cancellation can be
studied.

The concrete target is a simultaneous $p$-adic rank theorem for the tropical initial matrices:
show that if cancellation persists through $L$ layers at two non-collinear external primes,
then an integer relation appears among the two external valuation vectors.  Such a theorem
would contradict $K_\beta=0$.  Existing one-prime linear-form bounds do not compare the
residue kernels at different primes; this is the genuinely new ingredient required.

\subsection{Branch II: a common external valuation ray}

When $K_\beta\ne0$, all external vectors are collinear and the source report gives the exact
M\"obius-orbit condition
\begin{equation}\label{eq:mobius-branch}
 \beta=\frac{u+v\vartheta}{r+s\vartheta}
 \qquad(u,v,r,s\in\Q,\\ us-vr\ne0).
\end{equation}
This is the branch in which the external triangle inequality can be an equality.  It is also
arithmetically lower-dimensional: after clearing denominators, the external parts of $M$ and
$A$ are powers of one common primitive core.

The best next step is therefore not a generic all-prime determinant.  It is a mixed
real/radical invariant built from the canonical core and the two binomial fields already
formalized in the repository.  Useful intermediate targets are:
\begin{enumerate}
\item compare the field norms of the canonical radicals on the two sides of
      \eqref{eq:mobius-branch};
\item show that simultaneous radical degrees force an additional common perfect-power degree,
      contradicting affine primitivity; or
\item prove that the M\"obius coefficients in \eqref{eq:mobius-branch} cannot have the sign
      pattern imposed by nonnegative prime valuations unless one output is $\{2,3\}$-smooth.
\end{enumerate}
The repository modules
\begin{center}
\lean{CanonicalRadicalFieldNorm}, \lean{CanonicalRadicalValuation},\\
\lean{CanonicalRadicalSharpness}, \lean{CanonicalRadicalDivisibleHull},\\
\lean{SimultaneousRadicalDegrees}
\end{center}
provide much of the algebraic plumbing for this branch.

\subsection{Branch III: one smooth output}

If $M$ is $\{2,3\}$-smooth, then $\beta=a+d\vartheta$ with $d>0$; if $A$ is smooth, there is a
symmetric affine formula.  The source corpus proves that both outputs cannot be smooth and
that the other output contains a prime at least five.  In the extreme case $M=3^d$, the
problem is the localized radical assertion
\[
 (3^d)^\vartheta\notin\Z[1/3].
\]

This branch is the smallest transcendence-theoretic wall and deserves separate treatment.
Potentially useful objects are the irreducible binomial for the least localized radical, its
field norm, and reduction at primes over two and three.  A successful local argument must
remain sensitive to the real identity defining $\vartheta$; ordinary ideal factorization
alone is already exhausted by the canonical-radical modules.

\subsection{Branch IV: cross-coprime primitive outputs}

The cross-coprime branch has the strongest determinant constants and the cleanest residue
matrices.  Here the exact first tropical forms are the fiber Vandermondes of
Section~\ref{sec:firstlayer}.  This is the best laboratory for determining whether the
cancellation cascade is real or illusory.

The next computational experiment should not search for counterexamples.  For synthetic odd
pairs $(M,A)$ and Beatty depths generated by $\beta=\log_2M$, compute
\[
 v_2(D_T)-\tau_{2,T},
 \qquad
 v_3(D_T)-\tau_{3,T}
\]
for increasing $T$, together with the histogram of tropical layers contributing at the final
valuation.  Small exact experiments performed during this work show an excess compatible with
$N^{3/2}$ rather than $N^{5/2}$, matching Theorem~\ref{thm:first-layer-small}; they are
reconnaissance only, since a synthetic $A$ does not satisfy $A=3^\beta$.  A proof should seek a
uniform $o(H_T)$ bound, probably through $p$-adic exponential-polynomial zero estimates.

\subsection{Archimedean priority: solve the triangular HCIZ limit}

The complete grid finally supplies the missing $T^2$ rows, and its empirical measures are
explicit.  This makes Target~\ref{target:hciz} the highest-priority analytic problem.  The
required calculation is a large-$N$ asymptotic for the HCIZ integral with both spectra having
density $2t$ on $[0,1]$, under the coupled scaling $x_i y_j\asymp T^2$ and
$N\asymp T^2$.  Even a nonsharp upper bound improving the largest-term energy by more than
$1-\pi/4$ would materially change the $p$-adic comparison.

A practical sequence is:
\begin{enumerate}
\item numerically evaluate $\log\det(e^{x_i y_j})$ using high-precision total-positive
      algorithms for moderate $T$;
\item subtract the exact Vandermonde/factorial terms in \eqref{eq:hciz-det} to estimate the
      free-energy limit;
\item derive the corresponding variational problem using known HCIZ large-deviation methods;
\item only then reinsert the primewise valuation lower bounds.
\end{enumerate}
This route is attractive because the analytic limit does not depend on the unknown prime
factorization of $M$ and $A$; only $\beta$ enters the overall scaling.

\subsection{Finite verification priority}

The software floor is now $x>27$.  The immediate formal target is not another raw doubling of
the MPFR scan but a replayable certificate architecture.  A realistic hierarchy is:
\begin{enumerate}
\item a fast untrusted generator records only near-boundary pairs and rational interval data;
\item a small verified checker proves each interval enclosure and the monotone coverage between
      records;
\item chunks are combined by a theorem whose statement mentions only integer endpoints;
\item the final theorem is imported into the aggregate surface without
      \lean{native\_decide} or an opaque external axiom.
\end{enumerate}
The source-hash discrepancy found in Section~\ref{sec:finite27} makes this especially
worthwhile: mathematical replay is more robust than a manifest that identifies an external
binary only by prose.

\section{Lean formalization blueprint}
\label{sec:formalization}

The source corpus is already unusually strong on exact conditional structure.  The most useful
continuation is therefore not a second informal reimplementation of those results, but a
small collection of sharply separated modules whose statements expose the genuinely new
invariants.  This section gives a dependency order designed to keep arithmetic, asymptotics,
and transcendence inputs from leaking into one another.

\subsection{Proposed module graph}

{\footnotesize
\begin{longtable}{@{}>{\raggedright\arraybackslash}p{0.22\textwidth}>{\raggedright\arraybackslash}p{0.37\textwidth}>{\raggedright\arraybackslash}p{0.31\textwidth}@{}}
\toprule
Proposed module & Principal statement & Main dependencies \\
\midrule
\endfirsthead
\toprule
Proposed module & Principal statement & Main dependencies \\
\midrule
\endhead
\lean{SturmianNormalization} & Exact recurrences for $R_k,S_k$, common jump word $s_k$, and the exponent/output correspondence. & Existing primitive-generator and denominator-normalization modules; elementary floor arithmetic. \\
\addlinespace
\lean{ReducedOrbitDenominators} & Lowest-term denominators $q_k^{(2)},q_k^{(3)}$ and the full-depth alternative from $\min(v_2(M),v_3(A))=0$. & \lean{SturmianNormalization}; unique factorization and valuations. \\
\addlinespace
\lean{EffectiveGeneratorSeparation} & An abstract interface turning a two-logarithm lower bound into $|\beta-p/q|\ge c q^{-\mu}$ and a polynomial recurrence for continued-fraction denominators. & A stated external Baker--Matveev hypothesis, plus existing continued-fraction lemmas. \\
\addlinespace
\lean{BlockVandermondeValuation} & Exact $v_2$ and fixed-support $v_p$ formulas for the conditional block Vandermonde; output analogue. & Existing exact block parametrization; valuation of a difference with unequal valuations. \\
\addlinespace
\lean{OrdinaryAdelicFraction} & The asymptotic fraction \eqref{eq:ordinary-fraction} and the universal $1-\log2/(3\log6)$ ceiling. & \lean{BlockVandermondeValuation}; polynomial floor sums; primitive-output depth theorem. \\
\addlinespace
\lean{TriangleGini} & $\mathfrak G$ is a norm and satisfies the explicit formula \eqref{eq:G-explicit}. & Elementary measure integration, or a finite polygonal integration library. \\
\addlinespace
\lean{FiniteTropicalAssignment} & Exact rearrangement formula for a finite row list and a finite column-weight list. & Sorting/permutations; finite rearrangement inequality. \\
\addlinespace
\lean{MixedColumnRiemannSums} & Lattice count \eqref{eq:semigroup-count} and convergence of scaled low-weight columns to uniform measure on $\DeltaTri$. & Two-dimensional lattice-point estimates and Riemann sums. \\
\addlinespace
\lean{MixedAdelicGini} & The limit \eqref{eq:mixed-gini-limit}, conservation $\sum_pc_p=0$, and branch-sensitive three-vector bound. & Previous three modules; finite prime support. \\
\addlinespace
\lean{MixedTropicalCeilings} & Universal $4/5$, cross-coprime $7/10$, and structural $3/10$ constants. & \lean{MixedAdelicGini}; primitive-output depth theorem; explicit values from \lean{TriangleGini}. \\
\addlinespace
\lean{FullGridTransport} & Row/column empirical distributions, countermonotone energy $\pi/8$, maximum energy $1/2$, and the ratio $\pi/4$. & Triangular lattice sums, quantile integration, rearrangement. \\
\addlinespace
\lean{BeattyTropicalFiber} [accepted] & Exact minimizing fibers, floor-sized gap, normalized quotient divisibility, and block-determinant expansion for Beatty rank-one matrices. & \lean{TropicalInitialForm}, \lean{BlockFiberFactorization}; rearrangement and floor arithmetic. \\
\addlinespace
\lean{RowBlockLaplaceExpansion} [accepted] & Exact grouping of every assignment layer by row partitions; normalized and depth-$K$-truncated expansions with coefficients equal to signed products of unit minors. & \lean{BeattyTropicalFiber}, \lean{BlockFiberFactorization}; finite fiberwise sums and determinant algebra. \\
\addlinespace
\lean{AssignmentExcessGeometry} [accepted] & Prefix-surplus crossing energy, the Beatty lower bound $c(a)-\tau\ge\lfloor\beta\rfloor\mathcal E(a)$, and for $N>0$ the cutoff support bound $\#\{c(a)-\tau<K\}\le\binom{N+q-1}{q}$. & \lean{RowBlockLaplaceExpansion}; finite rearrangement, summation by parts, and stars and bars. \\
\addlinespace
\lean{BlockVandermondeValuationBudget} [accepted] & Exact collision-energy identity, elementary power-difference bound, and variable-block cube budget; the finite input to the paper-level $O(T^2)=o(H_T)$ first-layer specialization. & \lean{BlockFiberFactorization}; elementary $p$-adic valuation and finite sums. \\
\addlinespace
\lean{AllLayerUnitBasis} [accepted] & Monic unitriangular and determinant-one column changes extract complete-determinant unit divisors on positive dyadic/triadic unit nodes, with $p$-ordering weights $j+v_2(j!)$ and $\lfloor j/2\rfloor+v_3(\lfloor j/2\rfloor!)$. & \lean{BlockVandermondeValuationBudget}; integer polynomial evaluation and Vandermonde determinants. \\
\addlinespace
\lean{UnitOrderingGain} [accepted] & Exact block-weight decompositions and finite quadratic/cubic aggregate bounds; at most $S$ blocks of size at most $S$ have dyadic gain at most $S^3$ and twice the triadic gain at most $S^3$. & \lean{AllLayerUnitBasis}; Legendre digit sums and finite-sum bounds. \\
\addlinespace
\lean{RowBlockUnitDivisibility} [accepted] & Uniform dyadic/triadic unit divisors for every scaled row-block minor product and the additive full-determinant conclusion $p^{\tau+U}\mid\det$. & \lean{RowBlockLaplaceExpansion}, \lean{UnitOrderingGain}; row-scaled Vandermonde reindexing and a determinant-one fiberwise basis change. \\
\addlinespace
\lean{FiniteCheck2Pow27} & A kernel-replayable certificate excluding non-powers of two below $2^{27}$. & Existing finite-check certificate architecture; generated chunk certificates. \\
\bottomrule
\end{longtable}
}

The first two modules are almost entirely combinatorial.  The accepted Beatty, block-budget,
and unit-basis modules likewise import no transcendence theory; in particular, the qualitative
first-layer scale separation no longer needs an abstract $p$-adic logarithm interface.  Any
future use of the sharper estimate \eqref{eq:padic-two-log} should still be carried as a named
classical input unless and until Mathlib contains a suitable theorem.  This keeps the trust
boundary explicit.

\subsection{Finite statements before asymptotics}

The continuum notation in Sections~\ref{sec:mixed} and~\ref{sec:fullgrid} is compact, but Lean
should first see finite inequalities.  For rows $0\le k<N$ and columns $z_j\in\R$, the finite
rearrangement lemma should assert
\begin{equation}\label{eq:finite-rearrangement}
  \min_{\sigma\in S_N}\sum_{k=0}^{N-1} k z_{\sigma(k)}
  =\sum_{k=0}^{N-1} k z_{N-1-k}^{\uparrow},
\end{equation}
where $z^{\uparrow}$ is the nondecreasing rearrangement.  Every tropical assignment in this
report is an affine rescaling of \eqref{eq:finite-rearrangement}.  The limiting Gini identity
can then be obtained by three auditable steps:
\begin{enumerate}
\item prove convergence of empirical distributions for the triangular lattice;
\item prove convergence of the first moments and pairwise absolute-difference moments;
\item invoke the finite identity
\[
 \frac1{N^2}\sum_{i,j}|z_i-z_j|
 =\frac2{N^2}\sum_{j=0}^{N-1}(2j-N+1)z_j^{\uparrow}.
\]
\end{enumerate}
This route avoids formalizing optimal transport as a general theory.

\subsection{Exact theorem sketches}

The intended public statements can be phrased close to the mathematics.

\begin{lstlisting}[style=leanstyle]
/-- A primitive counterexample determines one common lower mechanical word
    driving both normalized rational orbits. -/
theorem exists_synchronized_sturmian_data
    (hfail : ExistsNonintegralSolution) :
    Exists fun B : Nat => Exists fun alpha U V : Real =>
      0 < alpha /\ alpha < 1 /\ Irrational alpha /\
      U = 2 ^ alpha /\ V = 3 ^ alpha /\
      (forall k : Nat,
        let e := Int.floor (k * alpha)
        normalizedTwo k = U ^ k / 2 ^ e /\
        normalizedThree k = V ^ k / 3 ^ e) := by
  ...

/-- The automatic all-prime divisor of the balanced mixed determinant
    cannot fill more than four fifths of its leading height. -/
theorem mixed_tropical_divisor_limsup_le_four_fifths
    (hprimitive : min (vTwo M) (vThree A) = 0) :
    limsup (fun T => log (tropicalDivisor T) / leadingHeight T)
      <= (4 : Real) / 5 := by
  ...

/-- The complete conditional grid has countermonotone/maximal transport
    ratio pi/4. -/
theorem full_grid_transport_ratio :
    countermonotoneEnergy / maximalEnergy = Real.pi / 4 := by
  ...
\end{lstlisting}

The actual APIs will need natural-valued denominator exponents rather than casts of negative
integers, and the asymptotic statements will likely use \lean{Tendsto} and \lean{IsLittleO}
rather than a custom \lean{limsup}.  The sketches are included to fix the mathematical
interfaces, not Lean syntax.

\subsection{The formalized finite core of the first tied layer}

The reusable initial form of a determinant at a prime is now implemented.  For a matrix over
$\Z$ and an assignment minimum $\tau_p$, reduce
\[
 p^{-\tau_p}\det(E_{ij})\pmod p.
\]
Only minimizing permutations survive.  The modules \lean{TropicalInitialForm},
\lean{BlockFiberFactorization}, and \lean{BeattyTropicalFiber} now prove, as purely finite
identities, that the relevant minimizers preserve the equal-$u$ or equal-$v$ fibers, that the
surviving sum is a product of smaller determinants, and that the next assignment layer begins
at least $\lfloor\beta\rfloor$ later.  The module \lean{RowBlockLaplaceExpansion} extends this
from the first fiber to every finite depth: after extracting the minimum power, its cutoff-$K$
identity discards all row partitions of excess at least $K$ modulo $p^K$ and expresses every
surviving coefficient as a signed product of square unit minors.  This is an exact support
decomposition, not a bound on the surviving signed sum.  The companion
\lean{AssignmentExcessGeometry} proves that Beatty excess dominates
$\lfloor\beta\rfloor$ times prefix-crossing energy and, for $N>0$, bounds the number of partitions below
cutoff $K$ by $\binom{N+q-1}{q}$, with $q=\lfloor(K-1)/\lfloor\beta\rfloor\rfloor$.
The finite valuation estimate is
separated into
\lean{BlockVandermondeValuationBudget}.  Combining its kernel-verified cube inequality with
the elementary paper-level count and size bounds for the Beatty fibers gives
\begin{equation}\label{eq:initial-form-bound-formal}
  v_2(I_{2,T})+v_3(I_{3,T})
  =O(N^2)=o(T^{5/2}).
\end{equation}
The optional classical two-logarithm input sharpens $O(N^2)$ to
$O(N^{3/2}(\log N)^2)$, but the sharper exponent is not needed for the qualitative
conclusion.  Finally, \lean{RowBlockUnitDivisibility} applies the unit-ordering divisor to
every row-partition minor product and proves the additive scaled-block conclusion
$p^{\tau+U}\mid\det$.  Thus the reusable finite algebra and valuation inequality are formalized, while
the displayed asymptotic specialization remains a short paper argument.  The remaining
mathematical target is no longer the first tied fiber; it is the valuation and cancellation
of the low-excess truncated sums as their cutoff grows.

\subsection{Certificate architecture for the new finite bound}

The existing verified bound through $4096$ should remain the trusted baseline until a new
certificate is accepted.  A practical route to $2^{27}$ is hierarchical:
\begin{enumerate}
\item partition $[1,2^{27})$ into fixed-width chunks;
\item for each non-power-of-two $m$, store an integer nearest-value locator $n_m$ and rational
      lower/upper enclosures proving either $m^\vartheta<n_m$ or $m^\vartheta>n_m$;
\item hash each chunk and prove a small checker correct in Lean;
\item combine the chunk theorems into one interval theorem.
\end{enumerate}
The MPFR run in Section~\ref{sec:finite27} is reconnaissance and a source of extremal cases;
it is not the certificate itself.

\section{Numerical sanity checks and reproducibility}
\label{sec:numerics}

None of the limiting constants in the determinant theorems depends on numerical evidence.
The computations below were used to catch normalization errors, check convergence rates, and
stress-test the closed forms before writing the proofs.

\subsection{Balanced-block assignment constants}

For each $N$, take the first $N$ pairs $(u,v)$ ordered by $u+\vartheta v$, pair them against
$N$ equally spaced row parameters, and solve the one-dimensional assignment problems by
sorting.  In the cross-coprime model, the structural places should tend to $3/10$ and the
minimal all-prime three-vector model to $7/10$.

\begin{table}[htbp]
\centering
\small
\begin{tabular}{@{}rcc@{}}
\toprule
$N$ & structural $2,3$ fraction & all-prime three-vector fraction \\
\midrule
25   & 0.3224571184 & 0.7258906037 \\
50   & 0.3135103379 & 0.7141968856 \\
100  & 0.3050363664 & 0.7053583986 \\
250  & 0.3024422068 & 0.7025311250 \\
500  & 0.3012614628 & 0.7013136826 \\
1000 & 0.3006780555 & 0.7006854844 \\
2500 & 0.3002789586 & 0.7002894348 \\
5000 & 0.3001354539 & 0.7001384826 \\
\midrule
Limit & $3/10$ & $7/10$ \\
\bottomrule
\end{tabular}
\caption{Finite assignment sums converging to the exact mixed-block constants.}
\label{tab:mixed-numerics}
\end{table}

The convergence from above is consistent with the $O(L^{-1})=O(N^{-1/2})$ boundary error in
triangular lattice counting.  The table also exposes the error in the proposed $2/3$ ceiling:
the finite all-prime ratio is already near $0.70$, and the exact limit is $7/10$.

\subsection{Complete-grid transport ratio}

For a representative hypothetical height parameter $\beta=27.314159$, form the row lattice
$n+\beta k\le T$ and the matching low-weight column lattice.  Normalize the
countermonotone assignment by the largest-term assignment.  The finite ratios are:

\begin{table}[htbp]
\centering
\small
\begin{tabular}{@{}rrc@{}}
\toprule
$T$ & number of rows & transport ratio \\
\midrule
50  & 73    & 0.7415113 \\
100 & 237   & 0.7559915 \\
200 & 841   & 0.7671191 \\
400 & 3157  & 0.7750371 \\
800 & 12210 & 0.7796781 \\
\midrule
Limit & --- & $\pi/4=0.7853981634\ldots$ \\
\bottomrule
\end{tabular}
\caption{Finite full-grid countermonotone/maximal transport ratios.}
\label{tab:grid-numerics}
\end{table}

The slow boundary convergence is expected: both empirical triangles have $O(T)$ boundary
points against $O(T^2)$ bulk points, and the matching weight scale also grows with $T$.

\subsection{Directed-rounding scan manifest}

The scan in Section~\ref{sec:finite27} used the literal C listing extracted from the attached
LaTeX source.  Its SHA-256 digest is
\begin{center}
\sha{f888c0d89be9366f01e2d26f339dcbe446556d39d4c5dbcf8f579ec0d0052e1d}.
\end{center}
This differs from the digest printed in the old report.  To ensure the discrepancy was only a
manifest error, the extracted source was first rerun on the old overlap interval
$[29{,}360{,}128,33{,}554{,}432)$; it reproduced the old extremal records and reported zero
failures.

The extension interval $[2^{26},2^{27})$ was split into sixteen chunks of width $2^{22}$ and
run at $192$ bits.  The concatenated chunk log has SHA-256
\begin{center}
\sha{b11843448cd8946442aba13758b5b666e4ac6eae50d2243fbdc32d83274cea8d}.
\end{center}
The two new extremal records were replayed at $320$ and $512$ bits.  A representative build
and invocation sequence was:
\begin{lstlisting}[style=pythonstyle,language=bash]
cc -O3 -std=c17 mpfr_scan.c \
  /lib/x86_64-linux-gnu/libmpfr.so.6 -lgmp -lm -o mpfr_scan

./mpfr_scan 67108864 71303168 192
# repeat consecutive width-2^22 chunks through 134217728
sha256sum extension-through-2pow27.log
\end{lstlisting}
The scanner excludes powers of two from the tested population because they are the known
solutions.  Across the extension it checked $67{,}108{,}863$ non-powers of two and reported
zero failures.

\subsection{Extremal replay records}

The narrowest new lower-directed separation was
\[
 m=132{,}833{,}765,
 \qquad
 n=7{,}501{,}348{,}219{,}569,
 \qquad
 \text{directed margin }5.9082259992879571\times10^{-21}.
\]
The corresponding high-precision ordinary difference is approximately
\[
 m^\vartheta-n=6.39397545334704\times10^{-8}.
\]
The narrowest new upper-directed separation was
\[
 m=129{,}799{,}345,
 \qquad
 n=7{,}231{,}571{,}292{,}851,
 \qquad
 \text{directed margin }8.3073231820333117\times10^{-22},
\]
with
\[
 n+1-m^\vartheta=8.666990435581896\times10^{-9}.
\]
These decimal differences are explanatory only.  The certification decision uses the MPFR
outward-rounded interval endpoints and the scanner's locator checks.

\section{Status matrix and dependency audit}
\label{sec:status}

\begin{longtable}{@{}>{\raggedright\arraybackslash}p{0.37\textwidth}>{\raggedright\arraybackslash}p{0.18\textwidth}>{\raggedright\arraybackslash}p{0.35\textwidth}@{}}
\toprule
Statement & Status & Dependency or caveat \\
\midrule
\endfirsthead
\toprule
Statement & Status & Dependency or caveat \\
\midrule
\endhead
Exact primitive monoid $\Sol=\Nzero+\Nzero\beta$ and candidate monoid & \statuslean & Inherited from the source report. \\
\addlinespace
Primitive-depth condition $\min(v_2(M),v_3(A))=0$ & \statuslean & Inherited; central to both universal ceilings. \\
\addlinespace
Synchronized dyadic--triadic Sturmian recurrence & \statusnew & Elementary from the exact monoid and floor normalization. \\
\addlinespace
Lowest-term denominator formulas and full-depth alternative & \statusnew & Unique factorization plus primitive depth. \\
\addlinespace
Effective polynomial irrationality measure for fixed $M$ & \statusknown{} plus \statusnew & Baker--W"ustholz/Matveev; constants depend effectively on $M$. \\
\addlinespace
Polynomial recurrence of continued-fraction denominators & \statusnew & Standard convergent inequality applied to the previous row. \\
\addlinespace
Exact ordinary block valuations at all input-support primes & \statusnew & Unequal valuation of the two monomials in every difference. \\
\addlinespace
Ordinary product fraction \eqref{eq:ordinary-fraction} & \statusnew & Exact floor sums; output endpoint errors are $O(T^2)$. \\
\addlinespace
Universal ordinary ceiling $0.871049\ldots$ & \statusnew & Primitive depth only; not a bound on cancellation from new primes in differences. \\
\addlinespace
Explicit triangular Gini norm & \statusnew & Direct integration. \\
\addlinespace
Mixed all-prime tropical limit \eqref{eq:mixed-gini-limit} & \statusnew & Rearrangement plus triangular Riemann sums. \\
\addlinespace
Universal $4/5$ mixed ceiling & \statusnew & Conservation, Gini triangle inequality, primitive depth. \\
\addlinespace
Cross-coprime $7/10$ and structural $3/10$ & \statusnew & Specialization $a=b=c=d=0$. \\
\addlinespace
Full-grid $\pi/4$ ceiling & \statusnew & Countermonotone and comonotone transport for the common quantile $\sqrt t$. \\
\addlinespace
HCIZ determinant identity & \statusknown & Exact finite identity; the needed large-$N$ inequality remains open. \\
\addlinespace
First tied-layer fiber factorization & \statusnew & Exact initial-form determinant algebra in the cross-coprime branch. \\
\addlinespace
First tied layer is $o(H_T)$ & \statusknown{} plus \statusnew & Standard $p$-adic two-logarithm estimates; no claim about higher valuation layers. \\
\addlinespace
Mixed cancellation-excess criterion & \statusnew & Immediate sufficient condition from integrality and the Gini deficit. \\
\addlinespace
Software exclusion through $2^{27}$ & \statuscomp & Complete MPFR directed-rounding scan; not a Lean proof. \\
\addlinespace
Kernel exclusion through $4096$ & \statuslean & Remains the strongest formal finite theorem audited in the source report. \\
\addlinespace
The Alaoglu--Erd\H{o}s conjecture itself & \statusopen & Equivalent to the unresolved special four-exponentials determinant. \\
\bottomrule
\end{longtable}

\subsection{Logical dependency graph}

The new deductions can be summarized as
\[
\begin{array}{c}
\text{exact primitive monoid}\ +\ \text{primitive depth}
\\[2pt]\Downarrow\\[-2pt]
\text{synchronized Sturmian system}
\qquad\text{and}\qquad
\text{exact block valuation profiles}
\\[2pt]\Downarrow\\[-2pt]
\begin{array}{c}
\text{ordinary all-prime fraction}\\[2pt]
\text{mixed Gini conservation law}\\[2pt]
\text{full-grid transport problem}
\end{array}
\\[2pt]\Downarrow\\[-2pt]
\begin{array}{c}
0.871049\ldots\text{ ordinary ceiling},\quad
4/5\text{ mixed ceiling},\quad
\pi/4\text{ full-grid ceiling}
\end{array}
\\[2pt]\Downarrow\\[-2pt]
\text{only cancellation beyond tropical minima or a sharper HCIZ bound can finish these routes.}
\end{array}
\]
Baker--Matveev separation and the finite scan are parallel inputs: they sharpen the possible
counterexample but do not enter the determinant ceiling proofs.

\section{Conclusions}
\label{ct:conclusion}

The Alaoglu--Erd\H{o}s conjecture remains unproved.  The continuation nevertheless changes the
shape of the search in several concrete ways.

First, the exact conditional monoid has a compact symbolic normal form.  After removing the
integer part of the primitive generator, one irrational Sturmian word controls both the
binary and ternary denominator jumps of two rational recurrences.  This exposes a possible
route through finite-state rigidity, but it also prevents a false shortcut: rational
multipliers do not by themselves make the word automatic.

Second, ordinary transcendence theory gives more Diophantine separation than the source
report's elementary argument records.  A fixed hypothetical generator is a ratio of two
logarithms of integers, so Baker--W"ustholz/Matveev yields an effective polynomial
irrationality measure.  Continued-fraction denominators therefore grow at most polynomially
from one stage to the next.  This is useful quantitative structure, but it does not conflict
with the $O(T)$ exact population of a dyadic block.

Third, the most direct $p$-adic determinant plan can now be audited with exact constants.
For ordinary block Vandermondes, every valuation at every prime already supporting the inputs
is accounted for, and primitive depth leaves at least $12.8950\ldots\%$ of the leading
archimedean budget uncovered.  For the balanced mixed determinant, all primes collapse into a
finite family of slope vectors summing to zero.  The triangular Gini norm then proves that
entrywise divisibility misses at least $20\%$ universally, at least $30\%$ in the
cross-coprime branch, and at least $70\%$ if one uses only the structural places $2,3$ there.
The complete two-dimensional grid improves the geometry but still leaves a termwise ceiling
of $\pi/4$.

These are no-go theorems for particular mechanisms, not evidence that the conjecture is false.
They identify the exact remaining mechanisms:
\begin{enumerate}
\item a leading-order cascade of cancellations among many equal-valuation Leibniz terms;
\item an archimedean determinant estimate substantially smaller than its largest-term bound,
      most naturally through the HCIZ free energy;
\item a genuinely mixed real/$p$-adic construction that turns
      $\log M\log3=\log A\log2$ into deep congruences; or
\item a rigidity theorem for the synchronized rational Sturmian system.
\end{enumerate}
The first cancellation layer is provably too small, so any successful $p$-adic proof must
propagate through many valuation layers rather than extracting one convenient Vandermonde
factor.

Finally, the directed-rounding computation was extended by one full binary exponent.  At the
software-trust level, no non-power-of-two candidate occurs below $2^{27}$, so a counterexample
would satisfy $x>27$.  The formal Lean boundary remains $x\ge12$ until a generated certificate
is checked by the kernel.  The old scanner-source hash printed in the source report is
incorrect; the literal source, overlap replay, new log hash, and high-precision extremal
replays are recorded explicitly in this document.

\begin{resultbox}
\textbf{Bottom line.}
The conjecture is still open, but the continuation produces new exact mathematics rather than
only a list of ideas: a synchronized Sturmian equivalence; effective polynomial separation;
closed all-prime Vandermonde asymptotics; a triangular Gini norm governing mixed tropical
divisibility; universal $4/5$, $7/10$, $3/10$, and $\pi/4$ ceilings; a first-layer
cancellation theorem; exact sufficient thresholds for future work; and a reproducible
software exclusion through $2^{27}$.  The next decisive calculation is the triangular HCIZ
free energy or, on the arithmetic side, a proof of a genuinely leading tropical cancellation
cascade.
\end{resultbox}

\clearpage

\section{Second continuation: scope and claim ledger}
\label{sec:continuation2}

The problem is


This document is a sequel rather than a replacement for the preceding parts of this report.  It restates every inherited fact used below, because several later
arguments depend sensitively on the exact quantifiers and on whether a
quantity is integral or merely rational.  In particular, the slogan ``one
counterexample generates a grid'' is proved here as an exact equality of
monoids, not used as a heuristic.

No complete proof of Conjecture~\ref{conj:AE} is claimed.  The strongest new
unconditional outputs of this continuation are:
\begin{enumerate}
\item a streamlined proof of the exact free-monoid structure, including the
      denominator-elimination step that is easy to understate;
\item a universal $N^2$-scale prime-divisibility lower bound for exact-grid
      interpolation determinants;
\item the triangular determinant inequality of
      Theorem~\ref{thm:triangle-inequality};
\item a rational-basis, product-formula extension that quantifies exactly how
      denominators defeat naive lattice reduction;
\item additional shared-prime, congruence, counting, and formalization
      consequences.
\end{enumerate}

\subsection{Claim ledger}

\begin{center}
\begin{tabular}{p{0.43\textwidth}p{0.16\textwidth}p{0.31\textwidth}}
\toprule
Claim & Status & Main input \\
\midrule
A counterexample is positive and transcendental & proved & elementary reductions; Gelfond--Schneider \\
All solutions lie in a rational affine line through $1$ and one nonintegral solution & proved & Six Exponentials Theorem \\
The complete solution set is $\Nzero+\Nzero\beta$ & proved & valuations and minimality \\
A counterexample has finite effective irrationality measure & proved & linear forms in two logarithms \\
Universal square-grid divisor $6^{(1+\beta)L_m}$ & proved & tropical determinant valuation \\
Triangular inequality $g(\sqrt{\beta\log2\log3})\ge0$ & proved & HCIZ bound, assignment asymptotics, Baker uniformity \\
Rational-basis denominator inequality & proved & product formula and signed assignment bounds \\
A complete contradiction for all $\beta$ & not obtained & requires further $N^2$-scale arithmetic gain \\
\bottomrule
\end{tabular}
\end{center}

The elementary reductions, the affine rank bound, the primitive least
counterexample and the exact free-monoid theorem restated at this point in the
source document are already established earlier in this report, and in
kernel-verified form: see Section~\ref{sec:conditional} and the modules
\lean{LeastSolution}, \lean{PrimitiveGenerator} and \lean{RationalPowerIndex}.
They are not repeated here.  What follows is the new material of the second
continuation, which takes the exact grid as given and asks how much arithmetic a
determinant built on it can extract.

\section{Effective Diophantine and counting consequences}

The two logarithmic descriptions of $\beta$ imply more than transcendence.
Because $A$ is not a power of $2$, the linear form
$q\log A-p\log2$ is nonzero whenever $(p,q)\ne(0,0)$.  Effective lower bounds
for linear forms in logarithms therefore give constants $c>0$ and $\tau>0$,
effectively computable from $A$, such that
\begin{equation}\label{eq:finite-type}
  \|q\beta\|\ge c q^{-\tau}\qquad(q\ge1).
\end{equation}
The analogous representation $\beta=\log B/\log3$ gives another such bound.
No useful small absolute exponent is asserted here; standard explicit bounds
are enormous.  The qualitative conclusion is nevertheless important.

\begin{corollary}[Finite irrationality measure]\label{cor:finite-mu}
A counterexample $\beta$ is not a Liouville number.  Its irrationality measure
is finite and effectively bounded in terms of either $A$ or $B$.
\end{corollary}

The exact monoid also gives an exact counting formula.  Let
\[
  N_{\mathcal S}(T):=\#(\mathcal S\cap[0,T]).
\]
Then, with $K=\lfloor T/\beta\rfloor$,
\begin{equation}\label{eq:count-exact}
 N_{\mathcal S}(T)=\sum_{k=0}^{K}\bigl(\lfloor T-k\beta\rfloor+1\bigr).
\end{equation}

\begin{proposition}[Quadratic counting with a power saving]\label{prop:count}
There is an effective $\delta>0$ such that
\[
 N_{\mathcal S}(T)=\frac{T^2}{2\beta}
 +\frac{\beta+1}{2\beta}T+O_\beta(T^{1-\delta}).
\]
\end{proposition}
\begin{proof}[Proof sketch]
Equation~\eqref{eq:finite-type}, the Erd\H{o}s--Tur\'an inequality, and the
geometric-series estimate for $\sum_{k\le K}e^{2\pi i h k\beta}$ give a power
saving for the discrepancy of the rotation sequence $\{k\beta\}$.  In
particular,
\[
 \sum_{k=0}^{K}\{T-k\beta\}=\frac{K+1}{2}+O(K^{1-\delta})
\]
for some effective $\delta>0$.  Insert this in
Equation~\eqref{eq:count-exact}.  Writing $K=T/\beta-\{T/\beta\}$, the
apparently oscillatory terms of order $T$ cancel, leaving the displayed linear
coefficient.
\end{proof}

The same argument gives quantitative equidistribution of the dense orbit
\[
 \left(2^{\{k\beta\}},3^{\{k\beta\}}\right)
 =\left(\frac{A^k}{2^{\lfloor k\beta\rfloor}},
         \frac{B^k}{3^{\lfloor k\beta\rfloor}}\right).
\]
Thus a counterexample would generate a dense family of rational points on the
arc $t\mapsto(2^t,3^t)$, with pure base-power denominators.  Their heights grow
exponentially in $k$, so this observation alone is far below the density needed
for Pila--Wilkie-type counting to force algebraicity; it is nevertheless a
useful reformulation for $S$-unit methods.


\section{Valuation-lattice saturation and an exact duality}
\label{sec:saturation-duality}

The denominator argument has a useful invariant formulation.  It also has a
symmetric consequence that is easy to miss if one works only with the least
counterexample $\beta$.

Write
\[
 a_p=v_p(A),\qquad b_p=v_p(B),
\]
and form the two-column integer matrix whose rows are
\[
  (\mathbf 1_{p=2},a_p)
  \quad\hbox{and}\quad
  (\mathbf 1_{p=3},b_p)
  \qquad (p\text{ prime}).
\]
Only finitely many entries in the second column are nonzero.  Let
$\mathcal L$ be the rank-two lattice generated by the columns in the direct
sum of all row copies of $\mathbb Z$.

\begin{theorem}[Saturated valuation lattice]
\label{thm:saturated-valuation-lattice}
The lattice $\mathcal L$ is primitive (saturated):
\[
  (\mathcal L\otimes_{\mathbb Z}\mathbb Q)
  \cap\bigoplus_p(\mathbb Z\oplus\mathbb Z)=\mathcal L.
\]
Equivalently, the greatest common divisor of all maximal minors of the matrix
is $1$.  In the present sparse matrix this is exactly
\[
 \gcd\!\left(
   \{a_p:p\ne2\}\cup
   \{b_p:p\ne3\}\cup
   \{a_2-b_3\}
 \right)=1.
\]
\end{theorem}

\begin{proof}
The nonzero maximal minors are, up to sign, exactly the displayed $a_p$,
$b_p$, and $a_2-b_3$: the first column is nonzero only in the row belonging to
$p=2$ in the first valuation block and the row belonging to $p=3$ in the
second.  It remains to prove that their gcd is $1$.

If a number $d>1$ divided all these minors, choose
$t\in\{0,\ldots,d-1\}$ with $t\equiv-a_2\pmod d$.  The congruence
$a_2\equiv b_3\pmod d$ then gives $t\equiv-b_3\pmod d$, while every other
valuation of $A$ or $B$ is already divisible by $d$.  Consequently both
$2^tA$ and $3^tB$ are perfect $d$th powers in $\mathbb N$.  Hence
\[
  \frac{\beta+t}{d}\in S.
\]
This number is nonintegral, and it is strictly between $0$ and $\beta$ because
$\beta>1$ and $0\le t\le d-1$, contradicting the choice of $\beta$.  Thus the
gcd of the maximal minors is $1$.  The rank-two Smith normal form criterion
then says exactly that the column lattice is saturated.
\end{proof}

A particularly concrete form of the theorem records exactly how extraction of
common roots behaves throughout the solution monoid.

\begin{corollary}[Perfect-power content is preserved]
\label{cor:perfect-power-content}
For $u,v\in\mathbb N_0$, not both zero, define
\[
 X_{u,v}=2^uA^v,\qquad Y_{u,v}=3^uB^v.
\]
Then
\[
 \gcd\!\left(
   \{v_p(X_{u,v}):p\text{ prime}\}
   \cup
   \{v_p(Y_{u,v}):p\text{ prime}\}
 \right)=\gcd(u,v).
\]
Equivalently, for every $d\ge1$, the two integers $X_{u,v}$ and $Y_{u,v}$
are simultaneously perfect $d$th powers if and only if
$d\mid u$ and $d\mid v$.
\end{corollary}

\begin{proof}
Every divisor of $\gcd(u,v)$ plainly divides all displayed valuations.  In the
other direction, if $d$ divides all of them, then
$X_{u,v}=C^d$ and $Y_{u,v}=D^d$ for positive integers $C,D$.  Thus
$(u+v\beta)/d\in S$.  Uniqueness of the representation in the irrational
basis $(1,\beta)$ gives $u/d,v/d\in\mathbb N_0$.
\end{proof}

For example, $A$ and $B$ cannot themselves be simultaneous proper powers; and
for every $t\ge0$, neither pair
\[
  (2^tA,3^tB),\qquad (2A^t,3B^t)
\]
can consist of simultaneous proper powers of the same degree.  These are not
independent ad hoc restrictions: all of them are specializations of one
saturated rank-two valuation lattice.

There is also an exact reciprocal description.  Put
\[
 S_{A,B}=\{y\in\mathbb R:A^y\in\mathbb Z,
                              \ B^y\in\mathbb Z\}.
\]

\begin{theorem}[Reciprocal solution monoid]
\label{thm:reciprocal-monoid}
Under the counterexample hypothesis,
\[
  S_{A,B}=\mathbb N_0+\frac1\beta\mathbb N_0.
\]
In particular $1/\beta$ is the least positive nonintegral element of
$S_{A,B}$.
\end{theorem}

\begin{proof}
If $y\in S_{A,B}$, then $x=\beta y$ satisfies
\[
  2^x=A^y\in\mathbb Z,
  \qquad
  3^x=B^y\in\mathbb Z.
\]
Hence $x=n+k\beta$ with $n,k\in\mathbb N_0$, and
$y=k+n/\beta$.  Conversely,
\[
 A^{k+n/\beta}=A^k2^n,
 \qquad
 B^{k+n/\beta}=B^k3^n
\]
are integers for all $n,k\ge0$.
\end{proof}

This reciprocal monoid shows that the two columns of the valuation matrix may
be interchanged without losing arithmetic information: the original and dual
root-extraction problems are the same saturated lattice viewed in opposite
bases.


\section{Exact-grid interpolation determinants}

The free monoid of the exact primitive-generator theorem
(Section~\ref{sec:conditional}) supplies two exact
rank-two grids: one in the logarithms of the bases and one in the solution
exponents.  Retaining both coordinates is the central step of this
continuation.

For finite sets $\mathcal R,\mathcal C\subset\Nzero^2$ with
$|\mathcal R|=|\mathcal C|=N$, define
\begin{align*}
 u_{a,b}&:=a\alpha+b\gamma,\qquad (a,b)\in\mathcal R,\\
 v_{n,k}&:=n+k\beta,\qquad (n,k)\in\mathcal C,
\end{align*}
and the matrix
\begin{equation}\label{eq:grid-matrix}
 M_{(a,b),(n,k)}
 :=\exp(u_{a,b}v_{n,k})
 =2^{an}3^{bn}A^{ak}B^{bk}.
\end{equation}
Every entry is a positive integer.

\subsection{Nonvanishing and total positivity}

\begin{lemma}[Strict positivity]\label{lem:tp}
Order the $u$-values and the $v$-values increasingly.  Then
\[
 D(\mathcal R,\mathcal C):=\det\bigl(\exp(u_i v_j)\bigr)_{i,j=1}^N>0.
\]
In particular it is a positive integer and hence at least $1$.
\end{lemma}
\begin{proof}
The $u$-values are distinct because $2$ and $3$ are multiplicatively
independent, and the $v$-values are distinct because $\beta$ is irrational.
The kernel $K(u,v)=e^{uv}$ is strictly totally positive on $\R\times\R$.
Equivalently, the functions $e^{v_1u},\ldots,e^{v_Nu}$ with
$v_1<\cdots<v_N$ form an extended complete Chebyshev system: their Wronskian is
\[
 e^{u(v_1+\cdots+v_N)}\prod_{i<j}(v_j-v_i)>0.
\]
The generalized mean-value theorem for Chebyshev systems gives the desired
sign.
\end{proof}

\subsection{Archimedean upper bound}

Write $\DeltaV(u)=\prod_{i<j}(u_j-u_i)$ and similarly for $v$, and let
$G(N+1)=\prod_{j=1}^{N-1}j!$ be the relevant Barnes product.

\begin{lemma}[HCIZ determinant inequality]\label{lem:HCIZ}
For increasingly ordered real vectors $u,v$,
\begin{equation}\label{eq:HCIZ-bound}
 0<\det(e^{u_i v_j})
 \le
 \frac{\DeltaV(u)\DeltaV(v)}{G(N+1)}
 \exp\!\left(\sum_{i=1}^N u_i v_i\right).
\end{equation}
\end{lemma}
\begin{proof}
The Harish--Chandra--Itzykson--Zuber formula with normalized Haar measure is
\[
 \int_{U(N)}e^{\operatorname{tr}(\operatorname{diag}(u)
 U\operatorname{diag}(v)U^*)}\,dU
 =G(N+1)\frac{\det(e^{u_i v_j})}{\DeltaV(u)\DeltaV(v)}.
\]
Von Neumann's trace inequality bounds the exponent in the integral above by
$\sum_i u_i v_i$.
\end{proof}

The factor $G(N+1)$ is crucial.  When both node sets come from two-dimensional
lattice regions, the leading $N^2\log N$ contributions of the two Vandermonde
products cancel the $N^2\log N$ contribution of $G(N+1)$ exactly.  The problem
therefore lives at the next, $N^2$-constant, scale.  This is the determinant
manifestation of the critical nature of four exponentials.

\subsection{Finite-prime lower bound as an assignment problem}

For a prime $p$, the valuation of an entry in~\eqref{eq:grid-matrix} is
\begin{equation}\label{eq:valuation-cost}
 w_p((a,b),(n,k))
 =\mathbf1_{p=2}an+\mathbf1_{p=3}bn+a_p ak+b_p bk.
\end{equation}

\begin{lemma}[Tropical determinant bound]\label{lem:tropical}
For every prime $p$,
\[
 v_p(D(\mathcal R,\mathcal C))
 \ge \min_{\sigma\in S_N}\sum_{i=1}^N w_p(i,\sigma(i)).
\]
\end{lemma}
\begin{proof}
Each Leibniz term has the displayed summed valuation.  The nonarchimedean
triangle inequality says that the valuation of their signed sum is at least
the minimum of their valuations.
\end{proof}

This converts divisibility into a finite optimal-transport problem.  Even a
coarse componentwise use of it already produces an $N^2$-scale divisor.

\section{Square grids: a universal divisor and the critical cancellation}

Take $\mathcal R=\mathcal C=I_m:=\{0,\ldots,m-1\}^2$ and $N=m^2$.  Denote the
corresponding determinant by $D_m$.

Let $L_m$ be the minimum of $\sum_i a_i n_{\sigma(i)}$ over bijections between
the two copies of $I_m$.  The multiset of first coordinates consists of each
integer $0,\ldots,m-1$ repeated $m$ times.  The rearrangement inequality gives
\begin{equation}\label{eq:Lm}
 L_m=m\sum_{j=0}^{m-1}j(m-1-j)
 =\frac{m^2(m-1)(m-2)}6.
\end{equation}
The same value applies to each of the four pairings $an,bn,ak,bk$.

\begin{theorem}[Universal square-grid divisor]\label{thm:square-divisor}
For every $m\ge1$,
\begin{equation}\label{eq:square-lower}
 \log D_m\ge L_m(\log2+\log3+\log A+\log B)
 =L_m(1+\beta)\log6.
\end{equation}
Equivalently, the integer $D_m$ is at least
$6^{(1+\beta)L_m}$ in logarithmic size.
\end{theorem}
\begin{proof}
For a fixed prime, the minimum of a sum of nonnegative cost components is at
least the sum of their separate minima.  Apply
Lemma~\ref{lem:tropical} to~\eqref{eq:valuation-cost}; each present component
has minimum $L_m$.  Sum the resulting lower bounds with weights $\log p$ and
use
$\sum_pa_p\log p=\log A$ and
$\sum_pb_p\log p=\log B$.
\end{proof}

The estimate is much stronger than the bare fact $D_m\ge1$.  Nevertheless, a
balanced rectangular asymptotic does not yet contradict it.  Let
$z=\sqrt{\alpha\gamma\beta}$.  After choosing rectangle aspect ratios so that
the two contributions to each projected logarithm are equal, the normalized
node distribution is that of
$Z=(X+Y)/2$ for independent $X,Y\sim\mathrm{Unif}[0,1]$.  One has
\[
 \E Z^2=\frac7{24},\qquad
 \iint\log|s-t|\,d\Law(Z)(s)d\Law(Z)(t)
 =\frac13\log2-\frac{25}{12}.
\]
Combining these with Lemma~\ref{lem:HCIZ} and
Theorem~\ref{thm:square-divisor} yields the necessary rectangular inequality
\begin{equation}\label{eq:rectangle-gap}
 h(z):=\frac12\log(4z)+\frac13\log2-\frac43+\frac z2\ge0.
\end{equation}
For the range allowed even by $A\ge3$, this inequality is already positive.
The square-grid divisor is therefore real progress in the arithmetic estimate,
but rectangular geometry spends too much archimedean size.

\section{Triangular grids and a new exclusion inequality}

The natural way to reduce the archimedean cost is to take the first $N$ nodes
rather than a rectangular block.

Let $\mathcal R_N$ consist of the $N$ pairs $(a,b)\in\Nzero^2$ with the
smallest values of $u=a\alpha+b\gamma$, and let $\mathcal C_N$ consist of the
$N$ pairs $(n,k)\in\Nzero^2$ with the smallest values of $v=n+k\beta$.
Ties do not occur.  Let $D_N$ be the positive integer determinant from these
sets.

Define
\[
 U:=\sqrt{2\alpha\gamma},\qquad
 V:=\sqrt{2\beta},\qquad
 z:=\sqrt{\alpha\gamma\beta},
\]
so $UV=2z$.

\subsection{Limiting distributions}

Elementary lattice-point counting gives
\[
 \#\{(a,b)\in\Nzero^2:a\alpha+b\gamma\le T\}
 =\frac{T^2}{2\alpha\gamma}+O(T),
\]
and similarly for $n+k\beta$.  Thus the $N$th thresholds are asymptotic to
$U\sqrt N$ and $V\sqrt N$.  After division by $\sqrt N$, the projected nodes
converge to the distributions
\[
 d\mu_U(t)=\frac{2t}{U^2}\,dt\quad(0\le t\le U),\qquad
 d\nu_V(t)=\frac{2t}{V^2}\,dt\quad(0\le t\le V).
\]
Their quantile functions are
\[
 Q_{\mu_U}(s)=U\sqrt s,\qquad Q_{\nu_V}(s)=V\sqrt s.
\]
The logarithmic energies are
\begin{equation}\label{eq:beta2-energy}
 \mathcal E(\mu_U)=\log U-\frac74,\qquad
 \mathcal E(\nu_V)=\log V-\frac74.
\end{equation}
The finite irrationality type supplied by linear forms in logarithms ensures
that exceptionally close projected lattice points do not spoil convergence of
the logarithmic energies; details are given in
Appendix~\ref{app:energy-convergence}.

\subsection{The archimedean constant}

The Barnes asymptotic is
\[
 \log G(N+1)=\frac{N^2}{2}\log N-\frac{3N^2}{4}+o(N^2).
\]
Using~\eqref{eq:beta2-energy}, the two Vandermonde terms in
Lemma~\ref{lem:HCIZ}, and comonotone coupling for the trace term, gives
\begin{equation}\label{eq:triangle-upper}
 \limsup_{N\to\infty}\frac{\log D_N}{N^2}
 \le \frac12\log(UV)-1+\frac{UV}{2}.
\end{equation}
Indeed,
\[
 \int_0^1Q_{\mu_U}(s)Q_{\nu_V}(s)\,ds
 =UV\int_0^1s\,ds=\frac{UV}{2}.
\]

\subsection{The finite-prime constant}

Uniform measure on the row triangle can be written as
\[
 a=\frac U\alpha X,\qquad b=\frac U\gamma Y,
 \qquad X,Y\ge0,\quad X+Y\le1,
\]
where $(X,Y)$ is uniform on the standard simplex.  Each coordinate has density
$2(1-x)$ and quantile
\[
 q(s)=1-\sqrt{1-s}.
\]
The column coordinates have the corresponding scales
$n=V X'$ and $k=(V/\beta)Y'$.
The countermonotone assignment constant is
\begin{equation}\label{eq:J}
 J:=\int_0^1q(s)q(1-s)\,ds
 =\frac\pi8-\frac13.
\end{equation}
Consequently each of the four logarithmically weighted components
$an\log2$, $bn\log3$, $ak\log A$, and $bk\log B$ contributes asymptotically
$UVJ N^2$ to the componentwise valuation lower bound.  Therefore
\begin{equation}\label{eq:triangle-lower}
 \liminf_{N\to\infty}\frac{\log D_N}{N^2}\ge4UVJ.
\end{equation}

\begin{theorem}[Triangular determinant inequality]\label{thm:triangle-inequality}
Every nonintegral counterexample $\beta$ satisfies
\begin{equation}\label{eq:g-ineq}
 \boxed{
 g\!\left(\sqrt{\beta\log2\log3}\right)\ge0,
 \qquad
 g(z)=\frac12\log(2z)-1+\left(\frac{11}{3}-\pi\right)z.}
\end{equation}
\end{theorem}
\begin{proof}
Subtract the lower bound~\eqref{eq:triangle-lower} from the upper
bound~\eqref{eq:triangle-upper}.  Since $UV=2z$ and
$4J=\pi/2-4/3$, the resulting necessary inequality is exactly
\eqref{eq:g-ineq}.
\end{proof}

The derivative $g'(z)=1/(2z)+11/3-\pi$ is positive for $z>0$, so $g$ has a
unique positive zero
\[
 z_0=1.12895010749713,
 \qquad
 \beta_\triangle:=\frac{z_0^2}{\log2\log3}=1.67370758736671.
\]
At the first arithmetically possible nonintegral value
$\beta=\log_2 3$, one has $z=\log3$ and
\[
 g(\log3)=-0.029549732685145<0.
\]

The numerical consequence of Theorem~\ref{thm:triangle-inequality}
is weaker than what is already known: it gives $\beta\ge\beta_\triangle=1.6737\ldots$, whereas the
kernel-verified finite check of Section~\ref{sec:finitecheck} gives
$\beta\ge12$ and the directed-rounding scan reaches $2^{27}$.
The inequality is recorded for its method, not for its bound: it is the first
exclusion obtained from a uniform asymptotic determinant estimate rather than from
a per-value certificate, and it exposes the arithmetic constant that would have to
improve in order to reach all $A$ at once.


\subsection{Numerical anatomy of the two grid geometries}

The table compares the triangular gap $g$ with the balanced-rectangle gap $h$
from~\eqref{eq:rectangle-gap}.  Negative values would contradict a
counterexample at that $\beta$.

\begin{center}
\begin{tabular}{lrrr}
\toprule
$\beta$ & $z=\sqrt{\alpha\gamma\beta}$ & $g(z)$ & $h(z)$ \\
\midrule
$1$ & 0.8726396796 & -0.2633422644 & -0.04093352603 \\
$\log_2 3$ & 1.098612289 & -0.02954973269 & 0.1871929656 \\
$\beta_{\triangle}$ & 1.128950107 & -1.054219794e-81 & 0.2159820074 \\
$2$ & 1.23409887 & 0.09973735836 & 0.3130828643 \\
$\log_2 5$ & 1.329717364 & 0.1872568287 & 0.3982047952 \\
$3$ & 1.511456262 & 0.3467367942 & 0.5531278372 \\
$5$ & 1.951281643 & 0.7053840783 & 0.900746934 \\
$10$ & 2.759528964 & 1.303060538 & 1.47815739 \\
\bottomrule
\end{tabular}
\end{center}

The triangular region improves the constant substantially, but $g$ is
increasing and becomes comfortably positive in the range $\beta\ge\log_2 5$.
To continue, one must either improve the finite-prime lower bound, find a
better variational region, or enrich the determinant.

\section{Rational changes of basis and the product formula}

A natural reaction to the positive gap is to replace $3$ by a ratio such as
$3/2$, whose logarithm is smaller.  Ignoring denominators makes this look like
an immediate proof.  The product formula shows exactly why that argument is
false and, more usefully, turns it into quantitative denominator constraints.

\subsection{A signed arithmetic functional}

For a rational $q>1$, write it in lowest terms as
$q=\num(q)/\den(q)$ and define
\begin{equation}\label{eq:rho}
 \rho(q):=
 \frac{J\log\num(q)-K\log\den(q)}{\log q},
 \qquad
 J=\frac\pi8-\frac13,\quad K=\frac16.
\end{equation}
The constant $J$ is the countermonotone coordinate cost in a simplex, while
$K$ is the comonotone cost, since a simplex coordinate has second moment
$1/6$.  If $q$ is an integer, then $\rho(q)=J$.  A denominator contributes with
$K$, not $J$, because a negative valuation reverses the relevant assignment
bound.

\begin{theorem}[Adelic triangular inequality]\label{thm:adelic-triangle}
Let $r_1,r_2>1$ be multiplicatively independent positive rationals and suppose
$R_i=r_i^\beta\in\Q_{>1}$ for $i=1,2$.  Put
\[
 z_r:=\sqrt{\beta\log r_1\log r_2}.
\]
Then the existence of these four rational exponentials implies
\begin{equation}\label{eq:adelic-rho}
 \frac12\log(2z_r)-1+z_r
 -2z_r\bigl(\rho(r_1)+\rho(r_2)+\rho(R_1)+\rho(R_2)\bigr)\ge0.
\end{equation}
Equivalently, if
\[
 \Delta_{\rm den}:=
 \sum_{q\in\{r_1,r_2,R_1,R_2\}}
 \frac{\log\den(q)}{\log q},
\]
then
\begin{equation}\label{eq:adelic-den}
 g(z_r)+2z_r\left(\frac12-\frac\pi8\right)\Delta_{\rm den}\ge0.
\end{equation}
\end{theorem}
\begin{proof}
Use the same first-$N$ triangular sets for the positive logarithms
$\log r_1,\log r_2$ and for $1,\beta$.  The determinant entries are now
rational.  For a fixed prime, split each signed valuation coefficient into its
positive and negative parts.  In the tropical determinant bound, a positive
coefficient is multiplied by the countermonotone minimum $J$, while a negative
coefficient is multiplied by the comonotone maximum $K$.  Summing over primes
and using the product formula for the nonzero rational determinant gives the
lower constant
\[
 2z_r\bigl(\rho(r_1)+\rho(r_2)+\rho(R_1)+\rho(R_2)\bigr).
\]
The HCIZ upper constant remains $\tfrac12\log(2z_r)-1+z_r$.  This proves
\eqref{eq:adelic-rho}.  Finally,
\[
 \rho(q)=J-(K-J)\frac{\log\den(q)}{\log q},
\]
and substitution gives~\eqref{eq:adelic-den}.
\end{proof}

\subsection{All unimodular bases of $\langle2,3\rangle$}

Let
\[
 M=\begin{pmatrix}a&b\\c&d\end{pmatrix}\in\operatorname{GL}_2(\Z)
\]
and suppose
\[
 \lambda_1=a\alpha+b\gamma>0,\qquad
 \lambda_2=c\alpha+d\gamma>0.
\]
Set
\[
 r_1=2^a3^b,\quad r_2=2^c3^d,\qquad
 R_1=A^aB^b,\quad R_2=A^cB^d.
\]
Then Theorem~\ref{thm:adelic-triangle} applies.  This supplies an infinite
family of necessary inequalities indexed by unimodular changes of basis.
Continued-fraction vectors make one $\lambda_i$ small, but the corresponding
rational base has a large denominator; Equation~\eqref{eq:adelic-den} measures
the precise price.

\begin{proposition}[Optimality of the integral basis]\label{prop:integral-basis-opt}
Among unimodular bases for which both transformed bases are integers, the
product $\lambda_1\lambda_2$ is at least $\alpha\gamma$.  Equality occurs only
for the original basis $(2,3)$ up to permutation.  Likewise, among unimodular
bases of the solution lattice whose two transformed exponents remain in
$\mathcal S$, the product of the positive exponent generators is at least
$\beta$.
\end{proposition}
\begin{proof}
If $2^a3^b$ is an integer, then $a,b\ge0$.  Thus an integral unimodular basis
has nonnegative entries.  Expanding,
\[
 (a\alpha+b\gamma)(c\alpha+d\gamma)
 \ge(ad+bc)\alpha\gamma\ge\alpha\gamma,
\]
because $|ad-bc|=1$.  Equality forces the standard basis up to permutation.
For the exponent lattice, the exact primitive-generator theorem says that a transformed
exponent is simultaneously integral at bases $2$ and $3$ exactly when its
coefficients in $(1,\beta)$ are nonnegative.  The same argument applies.
\end{proof}

This proposition explains why a denominator-free lattice reduction cannot
improve the triangular inequality: the exact monoid has already selected the
smallest integral basis.

\subsection{The basis $(2,3/2)$ and a gcd constraint}

Take
\[
 r_1=2,\qquad r_2=\frac32,\qquad
 R_1=A,\qquad R_2=\frac BA.
\]
Let $d=A/\gcd(A,B)$ be the denominator of $B/A$ in lowest terms and put
$\delta=\log(3/2)$.  Then
\begin{equation}\label{eq:delta-32}
 \Delta_{\rm den}=\frac{\alpha}{\delta}
 +\frac{\log d}{\beta\delta}.
\end{equation}
Theorem~\ref{thm:adelic-triangle} therefore gives a rigorous lower bound on
$d$, whenever the right side obtained by solving~\eqref{eq:adelic-den} is
positive.  In words: if $3/2$ substantially improves the archimedean
configuration, then $B/A$ must carry enough denominator to pay for that
improvement.  A counterexample cannot have excessive divisibility
$A\mid B$ in the range where the reduced-basis gap is negative.

For orientation, the next table substitutes the hypothetical value
$\beta=\log_2 A$.  The column $d_{\min}$ is the least integer denominator not
excluded by the inequality alone; it is not a claim that the corresponding
$B$ exists.

\begin{center}
\begin{longtable}{rrrrr}
\toprule
$A$ & $\log_2A$ & $z_r$ & gap before $\log d$ & $d_{\min}$ \\
\midrule
\endfirsthead
\toprule
$A$ & $\log_2A$ & $z_r$ & gap before $\log d$ & $d_{\min}$ \\
\midrule
\endhead
3 & 1.5849625 & 0.667419621 & -0.260297461 & 4 \\
5 & 2.32192809 & 0.807818616 & -0.0396109308 & 2 \\
6 & 2.5849625 & 0.852347316 & 0.0269341919 & 1 \\
7 & 2.80735492 & 0.88825597 & 0.0795954278 & 1 \\
9 & 3.169925 & 0.94387388 & 0.159569469 & 1 \\
10 & 3.32192809 & 0.966239056 & 0.191227186 & 1 \\
11 & 3.45943162 & 0.986033907 & 0.219022705 & 1 \\
12 & 3.5849625 & 1.00376439 & 0.243748126 & 1 \\
13 & 3.70043972 & 1.01980266 & 0.265979188 & 1 \\
14 & 3.80735492 & 1.03443012 & 0.286146736 & 1 \\
15 & 3.9068906 & 1.04786443 & 0.304581094 & 1 \\
17 & 4.08746284 & 1.07180649 & 0.337231592 & 1 \\
18 & 4.169925 & 1.08256404 & 0.351820061 & 1 \\
19 & 4.24792751 & 1.09264233 & 0.36544257 & 1 \\
20 & 4.32192809 & 1.10211837 & 0.378212214 & 1 \\
21 & 4.39231742 & 1.11105698 & 0.39022374 & 1 \\
22 & 4.45943162 & 1.11951323 & 0.401557272 & 1 \\
23 & 4.52356196 & 1.12753426 & 0.412281139 & 1 \\
\bottomrule
\end{longtable}
\end{center}

A second useful basis is
\[
 r_1=\frac32,\qquad r_2=\frac43,
 \qquad R_1=\frac BA,\qquad R_2=\frac{A^2}{B}.
\]
Its two logarithms are smaller, but Equation~\eqref{eq:adelic-den} now charges
the denominators of all four displayed ratios.  It therefore constrains the
possibility that $B/A$ and $A^2/B$ are simultaneously close to integers.  More
generally, continued-fraction bases produce a network of constraints on
\[
 \frac{B^q}{A^p}
 \quad\text{when}\quad q\log3-p\log2\text{ is small}.
\]
This is a concrete future route: prove that the denominator requirements for a
carefully chosen finite family of bases are incompatible with the prime
valuation vectors of $A$ and $B$.

\section{Joint prime transport and shared-factor gains}

The componentwise lower bounds used above are intentionally conservative.  At
a fixed prime, a single permutation must handle all valuation components at
once.  If a prime occurs in more than one generator, the separate
countermonotone assignments can conflict, making the true minimum strictly
larger.

For the original triangular determinant, define
\[
 \mathcal T_p(N):=\frac1{N^2}
 \min_{\sigma\in S_N}\sum_i w_p(i,\sigma(i)).
\]
Then the strongest direct product-formula lower bound is
\[
 \liminf_{N\to\infty}\frac{\log D_N}{N^2}
 \ge\sum_p(\log p)\liminf_{N\to\infty}\mathcal T_p(N).
\]
The baseline $4UVJ$ replaces each $\mathcal T_p$ by the sum of four independent
one-coordinate minima.

A simple rigorous gain estimate is available for a prime $p\notin\{2,3\}$ that
divides both $A$ and $B$.  Put $r=v_p(A)$, $s=v_p(B)$ and
\[
 R=\frac r\alpha,\qquad S=\frac s\gamma.
\]
On the normalized row simplex, the relevant scalar is $RX+SY$, while the
column coordinate $k$ has the marginal law of a simplex coordinate.  The
countermonotone cost is at least the product of the means minus the product of
the standard deviations.  Since
\[
 \E X=\E Y=\frac13,\quad
 \Var X=\Var Y=\frac1{18},\quad
 \Cov(X,Y)=-\frac1{36},
\]
one obtains
\begin{equation}\label{eq:overlap-bound}
 \liminf\mathcal T_p(N)
 \ge\frac{UV}{\beta}
 \left(
  \frac{R+S}{9}
  -\frac{\sqrt{R^2+S^2-RS}}{18}
 \right).
\end{equation}
The componentwise baseline for the same two terms is
\[
 \frac{UV}{\beta}J(R+S).
\]
Taking the larger of these two lower bounds gives a fully explicit shared-prime
improvement.  Summed over common prime divisors of $A$ and $B$, it yields a
sufficient contradiction criterion whenever the gain exceeds the positive
gap $g(z)$.

The estimate is strongest when $r/\alpha$ and $s/\gamma$ are comparable and
vanishes continuously as one valuation goes to zero.  It therefore does not
supply a universal proof: the worst arithmetic configuration may have nearly
disjoint prime support.  But it turns ``a large gcd should help'' into a
quantitative theorem and identifies exact valuation patterns worth targeting.

\section{Universal congruence divisibility is lower order}

There is also divisibility at primes not dividing $6AB$.  Modulo such a prime
$p$, a row of the matrix depends on $(a,b)$ only modulo $p-1$, by Fermat's
little theorem.  Hence there are at most $(p-1)^2$ distinct rows modulo $p$.
For an $N\times N$ determinant,
\[
 v_p(D_N)\ge N-\operatorname{rank}_{\mathbb F_p}(M)
 \ge N-(p-1)^2
\]
whenever the right side is positive.  Summing over primes $p\le\sqrt N$ and
using the prime number theorem gives the universal extra estimate
\begin{equation}\label{eq:congruence-gain}
 \sum_{\substack{p\le\sqrt N\\p\nmid6AB}}
  \bigl(N-(p-1)^2\bigr)\log p
 =\left(\frac23+o(1)\right)N^{3/2}.
\end{equation}
Higher prime powers and column congruences can improve the lower-order term.
They do not alter the $N^2$ constant in
Theorem~\ref{thm:triangle-inequality}.  This provides a useful barrier result:
congruence collisions that depend only on finite-group periodicity are too
small by a factor of $\sqrt N$.  A complete determinant proof needs
counterexample-specific valuation structure or a different determinant
geometry.

\section{Why several tempting routes still stop short}

\subsection{Continued fractions and integer gaps}

For integers $p,q>0$,
\[
 \frac{A^q}{2^p}=2^{q\beta-p},\qquad
 \frac{B^q}{3^p}=3^{q\beta-p}.
\]
When $p/q$ is a convergent to $\beta$, both ratios are close to $1$.  Since the
cross-multiplied differences are nonzero integers, one obtains exponential
lower bounds such as
\[
 |q\beta-p|\gg 2^{-p}.
\]
These are far weaker than the polynomial lower bounds already supplied by
linear forms in logarithms.  Conversely, Baker's polynomial lower bounds do
not contradict the $q^{-2}$ approximations supplied by continued fractions.
The two representations of $\beta$ give proportional linear forms, not two
independent small quantities.

\subsection{Catalan, $abc$, and small differences}

The equations $A^q-2^p=\pm1$ and $B^q-3^p=\pm1$ would be highly constrained,
but continued fractions only make the relative difference small, not the
absolute difference.  The absolute scale is exponential in $q$.  Standard
$abc$ heuristics do not bridge that exponential gap without an additional
source of unusually good approximation.

\subsection{Counting rational points}

The orbit of fractional parts gives $O(\log H)$ rational points of height at
most $H$ on a fixed transcendental arc.  This is compatible with every known
subpolynomial upper bound for rational points on definable transcendental
curves.  A counting theorem would need to exploit the pure $2$-power and
$3$-power denominator restrictions, not merely rationality.

\subsection{The four-exponentials boundary}

If $x$ is irrational, the pairs $(\log2,\log3)$ and $(1,x)$ are both
$\Q$-linearly independent, while all four exponentials
\[
 2,\quad3,\quad2^x,\quad3^x
\]
are algebraic.  Thus the Four Exponentials Conjecture would prove
Conjecture~\ref{conj:AE} immediately.  The exact cancellation of the
$N^2\log N$ terms in our determinant estimates is the quantitative shadow of
being exactly at this boundary.  Six Exponentials is sufficient to collapse
all possible exceptions to one affine direction, but not to remove that last
direction.

\section{A variational reformulation of the remaining determinant problem}

The preceding calculations extend to scaled lattice regions
$R,C\subset\R_{\ge0}^2$ of equal area.  Push forward uniform measure on $R$ by
$(a,b)\mapsto a\alpha+b\gamma$, and uniform measure on $C$ by
$(n,k)\mapsto n+k\beta$.  The HCIZ side is determined by their logarithmic
energies and comonotone transport.  The finite-prime side is a sum of
countermonotone or, more accurately, joint optimal-transport costs for the
valuation bilinear forms~\eqref{eq:valuation-cost}.

This suggests the following concrete variational program.
\begin{enumerate}
\item Choose regions $R,C$ and compute the archimedean functional
\[
 \mathcal A(R,C)=\frac12\mathcal E(\mu_R)
 +\frac12\mathcal E(\nu_C)+\frac34
 +\int_0^1Q_{\mu_R}(t)Q_{\nu_C}(t)\,dt,
\]
with the scale normalization determined by $\operatorname{area}(R)
=\operatorname{area}(C)$.
\item Compute or bound the exact prime transport functional
\[
 \mathcal P(R,C;A,B)=
 \sum_p\log p\inf_{\pi}
 \int w_p(r,\pi(r))\,dr.
\]
\item Any counterexample must satisfy
$\mathcal A(R,C)\ge\mathcal P(R,C;A,B)$ for every admissible pair of regions.
A single strict reverse inequality proves a contradiction.
\end{enumerate}

Rectangles and origin triangles are analytically soluble test cases.  The next
step is not merely to scan aspect ratios: it is to optimize jointly over region
shape and valuation configurations allowed by
the two-three unit corollary and the primitive-gcd lemma of
Section~\ref{sec:conditional}.  The latter
constraints are discrete and may prevent the valuation directions that make
the transport cost artificially small.

\section{Confluent and integer-valued determinant extensions}

Ordinary evaluations use one row per lattice point.  A potentially stronger
construction uses integer jets.  View~\eqref{eq:grid-matrix} as evaluating the
monomials $X^nY^k$ at the integer points
\[
 (X,Y)=(2^a3^b,A^aB^b).
\]
The Euler operators
\[
 D_X=X\frac{\partial}{\partial X},\qquad
 D_Y=Y\frac{\partial}{\partial Y}
\]
act by
\[
 D_X^rD_Y^s(X^nY^k)=n^rk^sX^nY^k,
\]
which remains integral at the lattice points.  Hasse derivatives give the
alternative integer values
\[
 \binom nr\binom ksX^{n-r}Y^{k-s}.
\]
Thus one can build confluent evaluation determinants without introducing the
irrational factors $n+k\beta$ that ordinary derivatives along the real curve
would produce.

The unresolved issues are nonvanishing and optimal scaling.  A useful target
is an integer-jet determinant whose archimedean limit behaves like a confluent
HCIZ determinant while its $p$-adic assignment problem acquires an additional
$N^2$ contribution.  This is more promising than generic factorial
normalization: the elementary divisibility obtained from binomial bases alone
is only of order $N^{3/2}\log N$ for balanced two-dimensional grids.

\section{Conditional resolutions and exact reformulations}

\subsection{Four Exponentials and Schanuel}

The Four Exponentials Conjecture immediately implies
Conjecture~\ref{conj:AE}.  Schanuel's conjecture implies Four Exponentials and
therefore also suffices.  The present problem is thus a particularly concrete
integral special case of a central algebraic-independence conjecture.

\subsection{Logarithm determinant formulation}

With $A=2^x$ and $B=3^x$, the condition is
\[
 \det\begin{pmatrix}
 \log2&\log A\\
 \log3&\log B
 \end{pmatrix}=0.
\]
The conjecture says that if $A,B$ are positive integers and this determinant
vanishes, then the two columns are related by a nonnegative integer factor.
Baker's theorem controls nonzero \emph{linear} forms in logarithms, whereas
this is a bilinear relation among four logarithms; that distinction is exactly
where standard transcendence machinery stops.

\subsection{Order-preserving monoid formulation}

A counterexample would give an order-preserving monoid isomorphism
\[
 \langle2,A\rangle_{\Nzero^2}\longrightarrow
 \langle3,B\rangle_{\Nzero^2},
 \qquad2^nA^k\mapsto3^nB^k,
\]
because both sides have logarithm $n+k\beta$ up to a fixed positive scale.
The exact primitive-generator theorem says this monoid is the entire domain on which
$m\mapsto m^{\log_2 3}$ is integer-valued.  A purely arithmetic classification
of such order-preserving power isomorphisms would prove the conjecture without
explicit transcendence estimates.

\section{Formalization roadmap}

A substantial portion of the new work is suitable for formalization in Lean,
with the deep analytic inputs isolated behind explicit interfaces.

The first three items of the source roadmap are already discharged in the
repository and are therefore not repeated as targets: the range, rationality,
additive-monoid and local-finiteness lemmas are \lean{IntegerExponent} and
\lean{LeastSolution}; the affine rank-two
proposition is \lean{MultiplicativeRank} and \lean{AlgebraicExponentLocus}; and the
primitive-gcd and denominator-elimination lemmas, together with the exact monoid
identity they yield, are \lean{PrimitiveOddCore},
\lean{ThreeDenominatorNormalization}, \lean{RationalPowerIndex} and
\lean{PrimitiveGenerator}.  What remains is the following.

\begin{enumerate}
\item Define finite matrices~\eqref{eq:grid-matrix}.  The tropical valuation
      lemma follows directly from the Leibniz formula and valuation of a sum.
      The square-grid assignment value~\eqref{eq:Lm} is a finite rearrangement
      theorem.
\item Treat strict total positivity, HCIZ, logarithmic-energy convergence, and
      asymptotic transport as separate analytic interfaces.  Finite versions
      of all determinant inequalities can still be machine checked once
      supplied with rational interval certificates.
\item For computational searches over region shapes or rational bases, emit
      proof certificates consisting of finite node lists, exact valuation
      assignment dual variables, and rational upper bounds for the real
      determinant.
\end{enumerate}

This separation would let the ProveIt development verify every new arithmetic
step even before the transcendence and asymptotic libraries are available.

\section{Prioritized research directions}

\begin{enumerate}
\item \textbf{Joint valuation optimization.}  Replace the componentwise
assignment lower bound by the exact transport cost for each prime, and optimize
over valuation vectors satisfying the primitive-gcd and base-prime-unit
conditions.  This is the most direct way to seek the missing $N^2$ constant.

\item \textbf{Rational-basis denominator network.}  Apply
Theorem~\ref{thm:adelic-triangle} to a finite set of continued-fraction bases.
Translate the resulting inequalities into lower bounds on denominators of
$B^q/A^p$.  Try to show that unique factorization makes the bounds mutually
incompatible.

\item \textbf{Variational region search.}  Numerically optimize convex and
nonconvex lattice regions using the exact archimedean energy and assignment
functionals.  Any promising negative gap should then be converted into a
rigorous interval proof.

\item \textbf{Integer jets and confluent total positivity.}  Establish a
nonvanishing theorem for Euler/Hasse-jet determinants and determine whether
their prime transport grows faster than their archimedean capacity.

\item \textbf{Exploit minimal-generator valuations.}  The conditions
$\min(v_2(A),v_3(B))=0$ and $g_0=1$ have not yet been fully inserted into the
variational optimization.  They should be imposed from the beginning rather
than used only as post-processing constraints.

\item \textbf{Effective finite certificates.}  For candidate ranges of $A$,
combine rigorous interval arithmetic with determinant and gcd constraints.
Although finite verification cannot settle the conjecture alone, it can reveal
which valuation patterns come closest to saturating the theoretical bounds.
\end{enumerate}


\section{A proof-producing finite verification}
\label{sec:finite-verification}

The structural and determinant arguments above are uniform, but they leave an
infinite range of possible values of $A=2^\beta$.  A complementary finite
calculation can remove an initial segment without assuming any unproved
transcendence statement.  The following computation is deliberately based on
outward-rounded intervals and an explicit positive-series remainder, rather
than on a test such as ``the floating-point value is not visually close to an
integer.''

\begin{theorem}[Certified finite range]
\label{thm:finite-range}
For every integer $A$ with
\[
  2\le A\le 100,000,
\]
the number $A^{\log_2 3}$ is an integer if and only if $A$ is a power
of $2$.  Consequently, a nonintegral Alaoglu--Erd\H{o}s counterexample must
satisfy
\[
  2^\beta>100,000,
  \qquad
  \beta>\log_2(100,000)\approx 16.609640474437.
\]
\end{theorem}

\begin{proof}[Certificate method]
For fixed $A$, put
\[
  F_A(b)=(\log b)(\log 2)-(\log A)(\log 3),\qquad b>0.
\]
The function $F_A$ is strictly increasing, and
$A^{\log_2 3}=b$ is equivalent to $F_A(b)=0$.
Every logarithm is enclosed without calling a transcendental floating-point
routine in the decisive comparison.  For $0\le z\le 1/3$ we use
\[
  \log\frac{1+z}{1-z}
  =2\sum_{j=0}^{M-1}\frac{z^{2j+1}}{2j+1}+R_M(z),
  \qquad
  0<R_M(z)<
  \frac{2z^{2M+1}}{(2M+1)(1-z^2)}.
\]
For an integer $n\ge1$, choose $k=\lfloor\log_2 n\rfloor$ and write
\[
  \log n=k\log2+
  \log\frac{1+z_n}{1-z_n},
  \qquad
  z_n=\frac{n-2^k}{n+2^k}\in[0,1/3).
\]
Thus all terms are positive rational operations, and directed rounding gives
certified lower and upper endpoints.  The run used 70 decimal digits and
$M=28$ terms.  A binary64 value was used only as an initial guess for the
integer immediately below the root; the program then moved monotonically
until it had certified, for consecutive integers $b,b+1$,
\[
  F_A(b)<0<F_A(b+1).
\]
Hence no integer can be the root.  Powers of $2$ are separated at the start,
and for $A=2^m$ the corresponding value is exactly $3^m$.

The completed run checked 99,983 nonpowers of $2$.  The largest number of
integer steps needed to repair an initial binary64 guess was 0.  The
smallest absolute interval endpoint encountered was
$\texttt{6.263278510736485942331545732E-13}$, far outside the interval uncertainty.  Runtime on the
working container was approximately 15.238 seconds.  The exact verifier
used for the run is reproduced in Appendix~\ref{app:finite-verifier}; its
SHA--256 digest is
\begin{center}
\texttt{acc099fb776f830b243a2421bc7ce27d5bedc63a38c3a7b190accb3c621cbadd}.
\end{center}
\end{proof}

\begin{remark}
The numerical bound is not presented as evidence that near misses cannot occur
later.  Its role is narrower and rigorous: it turns all $A$ in the displayed
finite interval into closed cases, while the determinant inequalities address
whole infinite parameter ranges and arithmetic configurations at once.
\end{remark}

\section{Conclusion}

The continuation does not close the final four-exponentials gap, but it
substantially narrows and clarifies the problem.

First, the hypothetical exceptional set is completely rigid: one least
counterexample $\beta$ would generate every solution, and only solutions, as
$n+k\beta$ with $n,k\ge0$.  Second, that exact grid supports determinants that
are simultaneously strictly totally positive over the reals and highly
divisible at finite primes.  The square-grid divisor and the triangular
inequality quantify both sides at the correct $N^2$ scale.  Third, rational
basis reduction is no longer a vague temptation: the adelic inequality shows
exactly how denominators pay for every smaller logarithm.  Finally, the
remaining obstacle has a precise form---one needs an additional $N^2$-scale
arithmetic gain, not merely more lattice points or a sharper polynomial error
term.

The most promising immediate attack is therefore a joint one: optimize the
triangular or variational determinant using the full prime valuation vectors,
while applying several rational bases to force incompatible denominator and
gcd requirements.  That program is concrete, computationally testable, and
amenable to partial formal verification.

\section{Third continuation: matrix coefficients, sparse interpolation, and boundary defects}
\label{sec:matrix-continuation}

Reports 1--2 in \lean{Article/new-reports/} are independent continuations of the same unified
base and therefore share most of their inherited body with this document.  They are reconciled
here by mathematical content rather
than by concatenation: stale module statuses, the old candidate-count off-by-one, and the
incorrect first-layer valuation equality are not reintroduced.  Their genuinely new core is a
matrix-coefficient reformulation, an exact sparse-interpolation threshold, a Newton-polygon
degree calibration, and a boundary-defect identity for the extremal interpolant.  The matrix
coefficient, formal prime-symbol, and sparse-interpolation components now have dedicated
kernel-checked Lean modules.

\begin{statusbox}{The full conjecture remains open.  The restricted matrix-coefficient
equivalence, formal prime-symbol zero classification, Structural Rank equivalence, finite
sparse power-curve lemmas, and boundary-dilation algebra below are \statuslean.
Newton-polygon geometry, the packaged BK-degree threshold, the counting reinterpretations,
and the complex monodromy consequence remain conventional mathematics or consequences of
named classical inputs.}
\end{statusbox}

\subsection{The restricted matrix-coefficient principle}
\label{mc:restricted}

For positive integers $M,A$, put
\[
  \mathsf L(M,A)=
  \begin{pmatrix}
    \log2&\log M\\
    \log3&\log A
  \end{pmatrix}.
\]
The determinant vanishes exactly when there is a real $\beta$ with
\[
  M=2^\beta,\qquad A=3^\beta.
\]
In that case
\[
  \mathsf L(M,A)=
  \binom{\log2}{\log3}\begin{pmatrix}1&\beta\end{pmatrix},
\]
and for rational vectors $w=(r,s)^T$ and $v=(u,z)^T$ one has the exact factorization
\begin{equation}
  w^T\mathsf L(M,A)v
   =(r\log2+s\log3)(u+z\beta).
  \label{mc:eq-factor}
\end{equation}

Define $\mathrm{MCC}_{2,3}^{\mathbb Z}$ to assert that every singular matrix of this
restricted form has nonzero rational $w,v$ with $w^T\mathsf L(M,A)v=0$.

\begin{theorem}[Restricted Matrix Coefficient equivalence]
\label{mc:thm-equivalence}
The Alaoglu--Erd\H{o}s conjecture is equivalent to
$\mathrm{MCC}_{2,3}^{\mathbb Z}$.
\end{theorem}

\begin{proof}
If the conjecture holds, singularity supplies a common exponent $\beta$, which must be an
integer $m$; the rational right vector $(-m,1)^T$ lies in the kernel.  Conversely, suppose
$\beta$ is a nonintegral common exponent.  Integrality of $2^\beta$ implies that $\beta$ is
irrational.  In \eqref{mc:eq-factor} the left factor cannot vanish for nonzero rational
$w$, by multiplicative independence of $2$ and $3$, and the right factor cannot vanish for
nonzero rational $v$, by irrationality of $\beta$.  Thus the required zero matrix
coefficient does not exist.
\end{proof}

The new module \lean{RestrictedMatrixCoefficient} formalizes the concrete $2\times2$ matrix,
its determinant, the common-exponent equivalence, the rank-one identity
\lean{restrictedLogMatrix_eq_rankOne}, the coefficient factorization and anisotropy, and the
exact endpoint
\[
 \lean{alaogluErdosConjecture_iff_restrictedMatrixCoefficientPrinciple}.
\]
This is \statuslean.  It is a strict specialization of the Matrix Coefficient Conjecture of
Dasgupta--Kakde~\cite{DasguptaKakde2026}; in the real $2\times2$ case their fewnomial
condition implies the desired matrix coefficient, while constructing that fewnomial from
singularity remains the open direction.

\subsection{The formal prime-symbol determinant}
\label{mc:formal-prime}

Let $a_p=v_p(M)$ and $b_p=v_p(A)$, and introduce one formal variable $X_p$ for each prime in
$\{2,3\}\cup\operatorname{supp}(MA)$.  The formal determinant is
\begin{equation}
  \Phi_{M,A}(\mathbf X)
   =X_2\sum_p b_pX_p-X_3\sum_p a_pX_p.
  \label{mc:eq-phi}
\end{equation}
Substitution $X_p=\log p$ gives $\det\mathsf L(M,A)$.  Coefficient comparison gives
\begin{equation}
  \Phi_{M,A}\equiv0
  \quad\Longleftrightarrow\quad
  \exists m\in\Nzero,\ M=2^m\ \text{and}\ A=3^m.
  \label{mc:eq-formal-zero}
\end{equation}
Indeed, all $b_p$ away from $3$ and all $a_p$ away from $2$ must vanish, while the
$X_2X_3$ coefficient forces $b_3=a_2$.  Hence the conjecture is equivalently the assertion
that specialization of this signed sparse quadratic at the prime-logarithm point cannot
vanish unless the polynomial is formally zero.

For a hypothetical counterexample $\Phi_{M,A}$ is nonzero and irreducible over $\Q$.
A factorization of a homogeneous quadratic would yield two rational linear factors; one of
them would vanish after $X_p=\log p$, contradicting rational linear independence of logarithms
of distinct primes.  The new \lean{FormalPrimeDeterminant} module defines the genuine
\lean{MvPolynomial} \lean{formalPrimeDeterminant}, proves its specialization
\lean{formalPrimeDeterminant_eval₂_log}, the unconditional classification
\lean{formalPrimeDeterminant_eq_zero_iff}, and the exact equivalence
\lean{alaogluErdosConjecture_iff_sparsePrimeLogStructuralRankPrinciple}.  These claims are
\statuslean.  The homogeneous-factor irreducibility bookkeeping is intentionally kept
separate and remains paper-level.

\subsection{Exact sparse interpolation on an irrational power curve}
\label{mc:sparse}

For irrational $\beta$, a monomial $X^rY^s$ restricted to
$(X,Y)=(e^t,e^{\beta t})$ becomes
\[
  e^{(r+s\beta)t}.
\]
Distinct pairs $(r,s)\in\Nzero^2$ give distinct frequencies.  The classical Chebyshev
zero theorem for real exponential polynomials, already formalized in
\lean{ExponentialPolynomialZeros}, therefore gives the following sharp lower bound.

\begin{theorem}[Sparse power-curve zero bound]
\label{mc:thm-sparse-lower}
An $m$-term nonzero polynomial restricted to $Y=X^\beta$, with $\beta$ irrational, has at
most $m-1$ distinct positive zeros.  Equivalently, vanishing at $R$ distinct positive curve
points requires at least $R+1$ nonzero monomials.
\end{theorem}

The new \lean{SparsePowerCurve} module proves the arbitrary-order indexed form
\lean{sparsePowerCurve_coeff_eq_zero_of_zeros}, its finite-zero-set form
\lean{card_le_of_sparsePowerCurve_eq_zero}, and the nonsingularity of the square evaluation
matrix.  It also proves the integral form
\lean{det_natPowerCurveEvaluationMatrix_ne_zero}: if
$y_i=x_i^\beta$ are natural numbers, then
\[
  \det\bigl(x_i^{r_j}y_i^{s_j}\bigr)\ne0
\]
for distinct monomial pairs and ordered positive $x_i$.  Thus the determinant is a nonzero
integer and has absolute value at least one.  These statements are \statuslean.

The matching upper construction is the node polynomial
\[
  P_0(T)=\prod_{i=1}^{R}(T-x_i).
\]
Because all $x_i$ are positive, every elementary symmetric coefficient is nonzero, so $P_0$
has exactly $R+1$ monomials.  This coefficient assertion and the vanishing at every node are
also formalized as \lean{positiveNodeInterpolant_support_card} and
\lean{positiveNodeInterpolant_eval}.  Consequently the following packaged minimum is a
paper-level corollary whose two substantive halves are kernel checked:
\begin{equation}
 \min\{\#\operatorname{supp}P:0\ne P\in\C[X,Y],\
       \ P(x_i,x_i^\beta)=0\ (1\le i\le R)\}=R+1.
 \label{mc:eq-exact-support}
\end{equation}
Here the paper-level minimum allows complex coefficients; the kernel-checked zero-bound and
matching positive-node construction are stated over the real core needed in the application.

In the restricted integer-output setting this gives an exact dichotomy.  If the common
exponent is an integer $m$, the minimum support on any nonempty family is $2$, attained by
$Y-X^m$; a single nonzero monomial cannot vanish at a positive point.  If the exponent is
nonintegral, it is irrational and the minimum on $R$ distinct points is $R+1$ by
\eqref{mc:eq-exact-support}.  The irrational branch is the kernel-checked sparse theorem;
the integral branch is the displayed elementary two-term witness and positivity argument.

For a counterexample $(M,A,\beta)$ and a finite
$\Omega\subset\Nzero^2$, the points
\[
  (2^a3^b,M^aA^b)
   =(e^{z_{a,b}},e^{\beta z_{a,b}}),
  \qquad z_{a,b}=a\log2+b\log3,
\]
have distinct parameters.  Hence $R=\#\Omega$ gives the exact minimum $R+1$.  On the
$2N\times2N$ box, any polynomial with at most $4N^2$ monomials would therefore save one term
over trivial interpolation and rule out a counterexample.  Dasgupta--Kakde ask for the
stronger support bound $<N^2$ on the same $4N^2$ points; for Alaoglu--Erd\H{o}s, one saved
monomial already suffices.

\subsection{BK degree and the exact factor-four threshold}
\label{mc:bk}

For $0\ne P\in\Q[X,Y]$, first divide out the greatest common monomial of its terms, as in the
Dasgupta--Kakde translation-invariant normalization.  Let $P^{\rm red}$ denote the resulting
polynomial, let $\Delta(P)$ be its Newton polygon, and set
\[
  \operatorname{bkd}(P)
   =\max\{\deg P^{\rm red},\,2\operatorname{Area}\Delta(P)\}.
\]
Pick's theorem gives
\begin{equation}
  \#\operatorname{supp}P\le \operatorname{bkd}(P)+2,
  \label{mc:eq-support-bk}
\end{equation}
with $+1$ when the Newton polygon is one-dimensional.  For irrational $\beta$ and $R\ge2$
distinct positive rational nodes whose ordinates $x_i^\beta$ are rational, let
$L\in\Q[X]$ be the Lagrange interpolation
polynomial of degree at most $R-1$ with $L(x_i)=x_i^\beta$.  Combining
\eqref{mc:eq-exact-support} with the thin interpolant
$Y-L(X)$ forces all $R+1$ available coefficients to occur and yields the exact paper-level
threshold
\begin{equation}
 \min\{\operatorname{bkd}(P):0\ne P\in\Q[X,Y],\
       \ P(x_i,x_i^\beta)=0\}=R-1.
 \label{mc:eq-bk-threshold}
\end{equation}
For the $K\times K$ semigroup grid with $K\ge2$ this is $K^2-1$; on the
Dasgupta--Kakde $2N$ grid it is
$4N^2-1$.  Their target $<N^2$ must therefore cross a genuine asymptotic factor-four gap.
No Pick-theorem or normalized lattice-area library is currently present in the pinned
Mathlib, so \eqref{mc:eq-support-bk}--\eqref{mc:eq-bk-threshold} remain \statusnew.

\subsection{Boundary defects of the extremal interpolant}
\label{mc:boundary}

Assume $K\ge2$ and a counterexample normalization
$M=2^\beta$, $A=3^\beta$ with $\beta$ irrational.  Let $Q_K\in\Q[X]$ be the
degree-$K^2-1$ polynomial interpolating the second coordinate on the $K\times K$ grid, and let
\[
  \Pi_{r,s}(X)=\prod_{\substack{0\le i<r\\0\le j<s}}(X-2^i3^j).
\]
The interior grid is stable under the two dilations.  Therefore
\begin{align}
 Q_K(2X)-M Q_K(X)&=\Pi_{K-1,K}(X)U_K(X),\label{mc:eq-defect-two}\\
 Q_K(3X)-A Q_K(X)&=\Pi_{K,K-1}(X)V_K(X),\label{mc:eq-defect-three}
\end{align}
where $U_K,V_K$ are unique of degree $K-1$; their leading coefficients are respectively
$q_K(2^{K^2-1}-M)$ and $q_K(3^{K^2-1}-A)$ when $q_K$ is the leading coefficient of $Q_K$.
The irrationality of $\beta$ makes both displayed factors nonzero.  More generally, the same
degree conclusion requires the explicit noncancellation hypotheses
$2^{K^2-1}\ne M$ and $3^{K^2-1}\ne A$.

The dilation operators $D_2f=f(2X)-Mf(X)$ and $D_3f=f(3X)-Af(X)$ commute.  Substituting
\eqref{mc:eq-defect-two}--\eqref{mc:eq-defect-three} and cancelling the common interior
factor $\Pi_{K-1,K-1}$ produces an identity of degree at most $2K-2$ involving the top and
right boundary pieces together with their bottom and left dilation preimages.  This boundary
compression is a concrete finite target:
it turns $K^2$ interpolation constraints into a low-degree compatibility equation.  It does
not by itself produce the missing fewnomial, and no such claim is made.

The module \lean{BoundaryDilationDefect} proves the residual polynomial algebra over a general
field and the commutator identity over an infinite field.  Its theorem
\lean{existsUnique_squareGridResiduals} gives both unique residuals, their exact
degree and leading coefficients; \lean{boundaryDilation_commutator} proves the displayed
commuting boundary identity; and \lean{squareBoundaryExpressions_natDegree_le} proves the
$2K-2$ bounds.  The interpolation values, node injectivity, and noncancellation of the top
coefficient remain visible hypotheses.  These statements are \statuslean.

\subsection{Counting thresholds and rational monodromy}
\label{mc:counting-monodromy}

A counterexample gives the rational points
\[
  P_{n,k}=\bigl(2^nM^k,3^nA^k\bigr)
  =\bigl(2^{n+k\beta},3^{n+k\beta}\bigr),\qquad n,k\in\Nzero,
\]
on the fixed curve $Y=X^{\thetaq}$.  Since the second coordinate controls the height, a
triangular lattice count gives
\[
  \#\{(n,k)\in\Nzero^2:n+k\beta\le\log_3H\}
  =\frac{(\log H)^2}{2\beta(\log3)^2}+O_\beta(\log H).
\]
Thus any uniform $o((\log H)^2)$ upper bound for rational points on the unbounded power curve
would prove the conjecture.  The synchronized compact orbit from \lean{CompactOrbit} gives a
sharper arithmetic target: its first $k$ points have height at most $A^k$, so an
$o(\log H)$ bound for points retaining the common floor level would also suffice.  These are
height reinterpretations of the accepted exact and quadratic conditional counts, not a new
point-counting theorem.

Finally, classical six exponentials forces a complex consequence which the current real
determinant formalization does not cover.  If $\beta$ were a nonintegral solution and
$q\in\Q^\times$, then
\begin{equation}
  e^{2\pi i q\beta}
  \quad\text{is transcendental},
  \label{mc:eq-monodromy}
\end{equation}
and so are its sine and cosine.  Apply six exponentials to the independent rows
$1,\beta$ and the independent columns $\log2,\log3,2\pi iq$; five of the six exponentials are
algebraic, forcing the sixth to be transcendental.  This is a valid consequence of the named
classical theorem, not a new kernel theorem.

\subsection{The reported \texorpdfstring{$2^{29}$}{2\^{}29} scan and its trust boundary}
\label{mc:finite29}

The first and fifth new reports record an MPFR directed-rounding extension through $2^{29}$,
including chunk tables, SHA--256 identifiers, extremal records, and high-precision replays;
report 5 also embeds the literal C source and a full manifest in its TeX.  The logical
consequence would be that a nonintegral solution has $x>29$.  However, the merged
\lean{new-reports/} directories contain no standalone scanner, canonical chunk logs, or
replay artifacts named by those hashes, and the long run was not independently repeated as
part of this merge.  It is therefore retained as a \emph{reported} \statuscomp{} result,
while the repository-reproducible software frontier in Section~\ref{sec:finite27} remains
$2^{27}$.

\subsection{What changed, and what did not}

The continuation produces exact finite reformulations, but no unconditional proof of the
conjecture.  Its most useful hierarchy is now:
\begin{enumerate}
\item the conjecture is exactly the restricted matrix-coefficient principle
  (\statuslean);
\item the formal prime-symbol determinant specializes to the logarithmic determinant, and
  its sparse Structural Rank principle is again exactly the conjecture (\statuslean);
\item a counterexample has maximal sparse interpolation complexity $R+1$, with the lower
  bound, matching positive interpolant, and evaluation determinant kernel checked;
\item the corresponding BK-degree threshold is exactly $R-1$ (\statusnew), exposing a
  factor-four gap to the current Matrix Coefficient auxiliary-polynomial target;
\item dilation defects compress the extremal grid interpolant to boundary polynomials of
  degree $K-1$ (\statuslean), but an arithmetic construction saving one monomial remains open;
\item generic polylogarithmic point counting is insufficient unless its exponent is below
  two globally, or below one on the synchronized compact orbit.
\end{enumerate}
The Massold v1 parameter substitution audited in report 2~\cite{Massold2025} does not cover the
$2\times2$ case: its displayed choice gives $t=1$, $\mu=\nu=1$, making the required strict
inequality $1/2>1$.  This is a narrow provenance warning about that version, not a claim
about every possible correction.

\section{Further continuations: support spectra and arithmetic determinant audits}
\label{sec:further-continuations}

Reports 3--5 are likewise independent continuations of that common base.  This section records
their new mathematics after reconciling overlaps rather than treating the reports as a
sequence.  Three themes survive audit: the exact support spectrum of a finite orbit,
integer-valued-polynomial divisors and primitive Smith normalization, and an adelic audit of
exact-grid determinants including a signed-progression no-go theorem.  We do not import a
stated ``quadratic divisor criterion'' from report 5: an inequality containing an unspecified
$o(N^2)$ term is not a sufficient eventual strict inequality.  Likewise, the proposed cusp
lattice realization is kept conditional on the uniform logarithmic-energy convergence that
its current proof does not establish.

\subsection{The exact orbit-collision support spectrum}
\label{fc:support-spectrum}

For a finite nonempty $\Omega\subset\Z^2$, define
\[
 T_x(\Omega)=\{i+jx:(i,j)\in\Omega\},\qquad D_x(\Omega)=\#T_x(\Omega),
\]
and let $\mu_x(\Omega)$ be the least number of monomials in a nonzero Laurent polynomial
in $\C[X^{\pm1},Y^{\pm1}]$ vanishing at all points
\[
 X_x(\Omega)=\{(2^{i+jx},3^{i+jx}):(i,j)\in\Omega\}.
\]
Multiplication by a monomial shows that the same minimum is obtained with ordinary
polynomials.  Restriction to the orbit turns every monomial into an exponential with
frequency $a\log2+b\log3$; these frequencies are distinct by unique factorization.
Consequently the exponential-polynomial zero bound and the positive-node polynomial give
the exact formula
\begin{equation}
  \boxed{\mu_x(\Omega)=D_x(\Omega)+1.}
  \label{fc:eq-support-spectrum}
\end{equation}
More generally, every evaluation matrix over $\C$ has rank
$\min(D_x(\Omega),\#S)$ for any finite monomial support $S$.  Repeated zeros counted with
multiplicity obey the same rule: prescribing total multiplicity $R$ costs at least $R+1$
monomials.

For $L,K\ge1$, on the rectangle
$\Omega_{L,K}=\{0,\ldots,L-1\}\times\{0,\ldots,K-1\}$ the collision count is explicit.
If $x\notin\Q$, then $D_x(L,K)=LK$.  If $x=r/s\in\Q_{\ge0}$ is in lowest terms, then
\begin{equation}
 D_{r/s}(L,K)=LK-(L-r)_+(K-s)_+,
 \qquad
 \mu_{r/s}(L,K)=D_{r/s}(L,K)+1.
 \label{fc:eq-rectangle-collision}
\end{equation}
Indeed two vertices collide exactly along the primitive direction $(r,-s)$; the collision
classes are paths, and vertices minus edges gives \eqref{fc:eq-rectangle-collision}.  Thus
rationality is equivalent to saving one monomial on some square grid, while an irrational
exponent has $\mu_x(L,L)=L^2+1$ for every $L$.  For integer outputs, a complex dependence
may be replaced by a rational one and denominators cleared, giving an integer certificate.
This is the exact finite-support version of the restricted Matrix Coefficient target.

The module \lean{RationalOrbitGrid} kernel-checks the exact rectangle formula, including the
zero-numerator edge case, the irrational full-cardinality statement, the irrational square-grid
sparse lower bound, and the canonical one-monomial-saving rational-orbit certificate once the grid exceeds
the reduced numerator and denominator.  Together with \lean{SparsePowerCurve}, this certifies
the real, nonnegative-exponent core of the spectrum above.  The arbitrary finite Laurent
support/rank formulation and the confluent multiplicity packaging remain paper-level.
\statuslean

\subsection{Coefficient height and the generic-modulus no-go}
\label{fc:sparse-height-no-go}

The exact support count does not control coefficient height.  For $M\ge1$, under an irrational
exponent $x$ with $A=2^x\in\Z_{>0}$, the canonical annihilator of the $M^2$
first-coordinate nodes is
\[
 Q_M(T)=\prod_{0\le i,j<M}(T-2^iA^j)\in\Z[T].
\]
It has exactly $M^2+1$ nonzero coefficients.  Write
$H(Q)=\max_k|q_k|$ for $Q(T)=\sum_kq_kT^k$.  Writing
$C_M=(2A)^{M^2(M-1)/2}$, positivity of the roots gives the exact finite bounds
\begin{equation}
 C_M\le\lVert Q_M\rVert_1\le2^{M^2}C_M,
 \qquad
 \frac{C_M}{M^2+1}\le H(Q_M)\le2^{M^2}C_M,
 \label{fc:eq-canonical-height}
\end{equation}
and hence
\[
 \log H(Q_M)=\frac{(1+x)\log2}{2}M^2(M-1)+O(M^2+\log M).
\]
Thus the canonical support is sharp but carries an $\exp(\Theta_x(M^3))$ archimedean cost.
Adding confluent multiplicities does not lower the support cost: the exponential Chebyshev
bound charges one extra monomial for each unit of prescribed total multiplicity.

The corresponding adelic criterion is elementary and exact.  If $P\in\Z[X,Y]$ and, at every
integer grid point $z$, a finite set of primes $S$ satisfies
\[
 |P(z)|_\infty\prod_{p\in S}|P(z)|_p<1,
\]
then the product formula forces $P(z)=0$.  Equivalently, it is enough to find a divisor
$q_z\mid P(z)$ with $|P(z)|<q_z$ at every row.  This identifies the right target but does not
construct the coefficients.

A cardinality-only modular pigeonhole does not suffice.  Let $E\in\Z^{R\times K}$ with
$1\le K\le R$, let $Q\ge2$ and $H\ge0$, and assume
$|E_{rj}|\le M_r$ with $M_r\ge1$.  Using coefficient vectors in
$\{0,\ldots,H\}^K$ and one modulus $Q$ requires
$(H+1)^K>Q^R$, which forces $H\ge Q$; the generic bound
$|(Ed)_r|\le KHM_r$ can then never prove $|(Ed)_r|<Q$.  Even row-dependent moduli do not
repair this architecture: for $H\ge1$ and $q_r\ge2$, the two sufficient conditions
\[
 (H+1)^K>\prod_{r=1}^Rq_r,
 \qquad q_r>KHM_r\quad(1\le r\le R)
\]
are incompatible.  Finally, if one uses the full adelic product---or a finite prime set
containing every prime divisor of the row content $g_r$---then dividing that row by $g_r$
changes neither the archimedean--$p$-adic product nor the meaningful cancellation: the product
formula for $g_r$ exactly cancels the apparent common divisibility.  The useful arithmetic
object is therefore the Smith profile of the row-primitive evaluation matrix, not a raw common
modulus.

These coefficient-height, adelic, pigeonhole, and row-content statements are paper-level.
The same support proof is base-independent for every multiplicatively independent positive
base tuple: distinct monomials restrict to exponentials with distinct real frequencies.

\subsection{Universal factorial content and primitive interpolation}
\label{fc:factorial-smith}

For $r\ge1$, write
\[
 \mathfrak F_r=\prod_{j=0}^{r-1}j!.
\]
For arbitrary integers $m_0,\ldots,m_r,y_0,\ldots,y_r$, the interpolation determinant with
columns $1,m,\ldots,m^{r-1},y$ is divisible by $\mathfrak F_r$.  Replacing the power columns
by $\binom{m}{0},\ldots,\binom{m}{r-1}$ proves the claim by a triangular Stirling change of basis,
and the nodes $0,\ldots,r$ with ordinate vector $(0,\ldots,0,1)$ show that the universal
divisor is best possible.  Applied to candidate ordinates, this strengthens the arbitrary-node
repulsion constant by the factor $\mathfrak F_r^{1/\binom{r+1}{2}}$ but leaves its exponent
unchanged.

The focused module \lean{SuperfactorialDeterminant} kernel-checks the nonnegative-integer-node
instance of this universal divisor, its exact consecutive-node sharpness witness, the resulting lower bound for a nonzero
candidate interpolation determinant, and the strengthened divided-difference repulsion and
packing inequalities.  As in \lean{ArbitraryNodeRepulsion}, the final two endpoints retain
\lean{RpowDividedDifferenceMeanValue} as an explicit hypothesis; the superfactorial
divisibility itself is unconditional in that formalized scope.  The arbitrary signed-integer
node formulation above remains a paper consequence.  \statuslean

The bivariate version is equally exact.  If $\mathcal S\subset\Nzero^2$ is a finite coordinate
lower set, then for exactly $|\mathcal S|$ integer points $(X_i,Y_i)$,
\begin{equation}
 \prod_{(a,b)\in\mathcal S}a!b!
 \ \bigm|\
 \det\bigl(X_i^aY_i^b\bigr)_{i,(a,b)\in\mathcal S}.
 \label{fc:eq-lower-set-factorial}
\end{equation}
The proof uses the integer-valued basis $\binom{X}{a}\binom{Y}{b}$.  On a weighted triangular lower
set of cardinality $N$, the logarithm of this divisor is
$O(N^{3/2}\log N)=o(N^2)$.  It is a genuine source of primes absent from the matrix entries,
but it cannot by itself repair a positive quadratic-scale determinant screen.

There is a sharper invariant for the unique $(M+1)$-term interpolating relation.  Let $M\ge2$
and let $u_1<\cdots<u_M$ be distinct integers, with $v_1,\ldots,v_M\in\Z$.  For the
$M\times(M+1)$ matrix whose $i$th row
is $(v_i,1,u_i,\ldots,u_i^{M-1})$, let $\Delta_M$ be the gcd of all signed maximal minors and
let
\[
 G_M=\gcd_i\prod_{a<b,\ a,b\ne i}(u_b-u_a),\qquad
 V_M=\prod_{a<b}(u_b-u_a),\qquad
 R_i=\prod_{j\ne i}|u_i-u_j|.
\]
Then a primitive cofactor generator, determined only up to global sign, has
$Y$-coefficient satisfying $|c_Y|=V_M/\Delta_M$, and
\begin{equation}
 \mathfrak F_{M-1}\mid G_M\mid\Delta_M,
 \qquad
 G_M=\frac{V_M}{\operatorname{lcm}_i R_i}.
 \label{fc:eq-smith-node-divisor}
\end{equation}
Moreover $G_M$ is the greatest divisor forced uniformly by the nodes as the integer ordinate
vector varies.  The remaining quotient $\Delta_M/G_M$ is therefore the precise
output-dependent Smith mass that an arithmetic proof must enlarge.  These Smith statements
and their $2^i3^j$-grid asymptotics remain paper-level.

\subsection{The adelic--HCIZ accounting identity}
\label{fc:adelic-hciz}

Let $N\ge1$.  For an exact-grid matrix $E_{ij}=\exp(u_iv_j)\in\N$ with strictly increasing
$u_1<\cdots<u_N$ and $v_1<\cdots<v_N$, let $D_N>0$ be its determinant.  Put
\[
 C_N^- =\sum_i u_iv_{N+1-i},\quad
 M_N=\frac{(\sum_i u_i)(\sum_jv_j)}N,\quad
 C_N^+=\sum_i u_iv_i.
\]
For each prime let $\tau_p$ be the minimum valuation of a Leibniz term, and set
$T_N=\sum_p\tau_p\log p$ and $Q_N=\prod_pp^{\tau_p}$.  The HCIZ identity writes
\[
 \log D_N=V_N+\log I_N,\qquad
 V_N=\log\DeltaV(u)+\log\DeltaV(v)-\log\!\prod_{j=1}^{N-1}j!.
\]
Here $I_N$ is the Harish--Chandra--Itzykson--Zuber integral
$\int_{U(N)}\exp(\operatorname{tr}(\operatorname{diag}(u)U
\operatorname{diag}(v)U^*))\,dU$ for normalized Haar probability measure.
Rearrangement, averaging, Jensen, and von Neumann's trace inequality give the exact sandwich
\begin{equation}
 T_N\le C_N^-\le M_N\le\log I_N\le C_N^+.
 \label{fc:eq-hciz-sandwich}
\end{equation}
Defining
\[
 \Xi_N=\log(D_N/Q_N),\quad S_N=V_N+M_N-C_N^-,\quad
 H_N=\log I_N-M_N,\quad A_N=C_N^--T_N
\]
therefore gives
\begin{equation}
 \boxed{\Xi_N=S_N+H_N+A_N},\qquad H_N,A_N,\Xi_N\ge0.
 \label{fc:eq-adelic-decomposition}
\end{equation}
This separates geometric screen, random-matrix entropy, and incompatible primewise
assignments instead of hiding them in one determinant estimate.

For the inherited triangular spectra, the best possible joint screen has normalized limit
\[
 f_{\rm joint}(c)=\tfrac12\log c-1+\left(\tfrac49-\tfrac\pi8\right)c,\qquad
 c=2\sqrt{\beta\log2\log3}.
\]
Its unique zero corresponds to $\beta=6.9269442924\ldots$.  Hence throughout the
kernel-verified range $\beta\ge12$, no combination of an exact HCIZ evaluation and termwise
prime valuations can close the inherited triangular determinant family.  This is a no-go
theorem for that geometry, not for all determinant methods.

Report 5 proposes power-cusp spectra whose formal continuum screen tends to $-\infty$.
The continuum formula itself is valid under convergence of logarithmic energies, but weak
lattice convergence does not supply that hypothesis near the diagonal.  A complete lattice
realization needs an explicit spacing/truncation estimate.  Even after such a repair, the
separated-prime-support branch makes the adelic alignment defect grow linearly in the cusp
parameter, so cusp concentration alone cannot be a universal proof.

\subsection{Signed progressions and the denominator tax}
\label{fc:signed-progressions}

Assume $N\ge1$, $p,q,r,s,M,A\in\N_{>0}$, and the counterexample relations
$M=2^\beta$, $A=3^\beta$ with $\beta$ irrational.  For $0\le i,j<N$, use the nonnegative exponent vectors
\[
 (a_i,b_i)=(iq,(N-1-i)p),\qquad
 (n_j,k_j)=(jr,(N-1-j)s).
\]
Set
\[
 X=2^{qr},\quad Y=3^{pr},\quad Z=M^{qs},\quad W=A^{ps},
 \qquad P=XW,\quad Q=YZ.
\]
Define the original integer matrix and its determinant by
\[
 E_{ij}=2^{a_in_j}3^{b_in_j}M^{a_ik_j}A^{b_ik_j},
 \qquad D_N=\det(E_{ij})_{0\le i,j<N},
\]
and let $Q_N^{\rm trop}$ be the termwise tropical divisor of this matrix.
Then
\[
\log(P/Q)=(q\log2-p\log3)(r-s\beta),
\]
so $P\ne Q$: unique factorization makes the first factor nonzero, while irrationality makes
the second nonzero.  A direct reduction to a geometric Vandermonde gives
\begin{equation}
 |D_N|=(PQ)^{\binom N3}
 \prod_{d=1}^{N-1}|P^d-Q^d|^{N-d}.
 \label{fc:eq-signed-progression}
\end{equation}
After writing $g=\gcd(P,Q)$, $\widehat P=P/g$, $\widehat Q=Q/g$, every row, column, and common
gauge factor cancels from the tropical quotient:
\begin{equation}
 \frac{|D_N|}{Q_N^{\rm trop}}
 =\prod_{d=1}^{N-1}|\widehat P^{\,d}-\widehat Q^{\,d}|^{N-d}.
 \label{fc:eq-progression-quotient}
\end{equation}
Let $B=\min(\widehat P,\widehat Q)$ and
$\eta=|\log(\widehat P/\widehat Q)|$.  With
$G(N+1)=\prod_{d=1}^{N-1}d^{N-d}$, the exact expansion is
\begin{align}
 \log\frac{|D_N|}{Q_N^{\rm trop}}
 &=\binom{N+1}{3}\log B+\binom N2\log\eta+\log G(N+1)+R_N,\label{fc:eq-tax}\\
 0&\le R_N\le\binom{N+1}{3}\eta.
\end{align}
The integer gap $|\widehat P-\widehat Q|\ge1$ forces
$\log B+\log\eta\ge-\eta$.  Thus the product of two continued-fraction errors is not free:
the denominator height repays the apparent quadratic smallness at cubic scale.  This exact
algebra rules out the naive double-continued-fraction progression, while leaving open more
structured cyclotomic cancellation inside $\widehat P^d-\widehat Q^d$.

The module \lean{SignedProgressionDeterminant} kernel-checks the geometric Vandermonde
factorization, the primitive invariant quotient and its coprimality with both primitive
bases, the exact finite logarithmic decomposition, the nonnegative remainder bound, and the
integer-gap inequality.  The original four-parameter matrix-to-tropical-quotient reduction,
the continued-fraction asymptotics, and any general impossibility theorem for nontropical
determinant cancellation remain paper-level.  \statuslean

\begin{statusbox}{The universal univariate superfactorial divisor and its strengthened
repulsion/packing consequences are \statuslean{} in \lean{SuperfactorialDeterminant}; the
finite signed-progression and denominator-tax algebra is \statuslean{} in
\lean{SignedProgressionDeterminant}.
The general support-spectrum, bivariate factorial/Smith, and HCIZ statements are paper proofs
extracted from reports 3--5.  For signed progressions, the original four-parameter/tropical
identification and continued-fraction conclusion remain paper-level.  No cusp lattice limit
or full conjecture is asserted.}
\end{statusbox}

\section{Sixth report: phase randomization and positive-kernel certificates}
\label{sec:phase-certificates}

Report 6 is another independent continuation of the same common base.  Its matrix-coefficient
and sparse-grid arguments overlap Sections~\ref{mc:restricted}--\ref{mc:sparse}, so they are
reconciled rather than duplicated.
Its genuinely distinct contribution is a probability-theoretic control experiment that
separates a generic random phase from the arithmetically distinguished phase zero.  The
experiment does not declare $\thetaq$ generic and does not prove the conjecture; it identifies
one exact finite Fourier estimate that would suffice.

\subsection{A uniform metric estimate and random-phase extinction}
\label{pc:metric-phase}

Fix $0<a<b$ and, for $m\ge2$ and $0<\varepsilon\le1/4$, set
\[
 E_m(\varepsilon)=
 \{\alpha\in[a,b]:\distZ{m^\alpha}<\varepsilon\}.
\]
The monotone substitution $y=m^\alpha$ and a harmonic-sum estimate give
\begin{equation}
 \operatorname{Leb}(E_m(\varepsilon))
 \le \frac83\left(b+\frac{1+\log3}{\log2}\right)\varepsilon.
 \label{pc:eq-metric-bound}
\end{equation}
The constant is independent of $m$.  Consequently every nonnegative summable sequence
$\sum_m\varepsilon_m<\infty$ satisfies
\[
 \distZ{m^\alpha}<\varepsilon_m
 \quad\hbox{for only finitely many $m$}
\]
for almost every $\alpha>0$, by the first Borel--Cantelli lemma.  This is a metric theorem,
not a specialization theorem for the named exponent $\thetaq$.

For the synchronized model, fix $N\ge1$ and put
\[
 D_N=\{a\in\Z:2^N<a<2^{N+1}\},\qquad
 Z_N(u;\eta)=\#\{a\in D_N:\distZ{a^{\thetaq}-u}<\eta\},
\]
where $u$ is a phase on $\R/\Z$.  A uniform phase makes each radius-$\eta$ arc have Haar
measure $2\eta$, so at the natural resolution $\eta=3^{-N}$,
\begin{equation}
 \mathbb E_u Z_N(u;3^{-N})
 =2(2^N-1)3^{-N}<2(2/3)^N.
 \label{pc:eq-random-phase}
\end{equation}
Summability and Borel--Cantelli show that almost every single phase has no sufficiently late
near hits.  No independence between blocks is used.

Phase zero is qualitatively different.  Let
\[
 \mathcal H_N=\#\{a\in D_N:a^{\thetaq}\in\Z\}.
\]
If the conjecture holds then $\mathcal H_N=0$.  Under failure, the least nonintegral generator
$\beta$ and \lean{ExactCounting} give the exact nonintegral count
\begin{equation}
 \mathcal H_N=\left\lfloor\frac{N+1}{\beta}\right\rfloor,
 \qquad \frac{\mathcal H_N}{N}\longrightarrow\frac1\beta,
 \qquad Z_N(0;3^{-N})\ge\mathcal H_N.
 \label{pc:eq-phase-zero}
\end{equation}
This is deliberately the nonintegral count: the full candidate count in the half-open dyadic
block convention used elsewhere in the report has the additional integral point and equals
$1+\lfloor(N+1)/\beta\rfloor$.

\subsection{The Fej\'er certificate}
\label{pc:fejer}

Write $e(t)=e^{2\pi i t}$ and, for an integer $H\ge2$, define
\[
 F_H(t)=\frac1H\left|\sum_{j=0}^{H-1}e(jt)\right|^2,
 \qquad \Phi_H(t)=\frac{F_H(t)}H.
\]
Then $\Phi_H\ge0$, $\Phi_H(0)=1$, and
\[
 \Phi_H(t)=\frac1H\sum_{|h|<H}
   \left(1-\frac{|h|}{H}\right)e(ht).
\]
With
\[
 S_N(h)=\sum_{a\in D_N}e(ha^{\thetaq}),\qquad
 Q_N(H)=\sum_{a\in D_N}\Phi_H(a^{\thetaq}),
\]
positivity and the finite Fourier expansion give the exact certificate
\begin{align}
 \mathcal H_N&\le Q_N(H),\label{pc:eq-hit-score}\\
 Q_N(H)&=\frac{|D_N|}{H}+
 \frac2H\sum_{h=1}^{H-1}\left(1-\frac hH\right)\Re S_N(h).
 \label{pc:eq-fejer-score}
\end{align}
Thus the conjecture follows if, along an unbounded sequence, one can choose $H=H_N^*$ with
$|D_N|/H_N^*=o(N)$ and make the weighted Fourier term in
\eqref{pc:eq-fejer-score} equal to $o(N)$.  The natural bandwidth $H=3^N$ has exponentially
small random-phase mean, but it demands information at exponentially many frequencies;
fixed-mode equidistribution from \lean{BlockWeyl} is therefore far too coarse.

The same diagnosis appears in the exact second moment.  If
$d_{a,a'}=\distZ{\{a^{\thetaq}\}-\{a'^{\thetaq}\}}$ and $0<\eta\le1/4$, then
\[
 \mathbb E_u Z_N(u;\eta)^2
 =2\eta|D_N|+2\sum_{a<a'}(2\eta-d_{a,a'})_+.
\]
The second term is the microscopic pair-correlation energy that a counterexample must force
to be large at phase zero.

\subsection{Why naive random sparse coefficients do not bridge the gap}
\label{pc:random-signs}

The integrality amplifier behind the sparse certificate is exact: if
$D\in\N_{>0}$, $0\ne P\in\Q[X,Y]$, $DP\in\Z[X,Y]$, and $|P(g)|<D^{-1}$ at every conditional integer-grid
point $g$, then $P(g)=0$ everywhere.  Sections~\ref{mc:sparse} and
\ref{fc:support-spectrum} then turn a subcritical support into rationality of the common
exponent.

Unconditioned random signs do not naturally reach this threshold.  For distinct monomials
$M_\nu$, real $c_\nu$, and independent signs $\epsilon_\nu$, put
$P_\epsilon=\sum_\nu\epsilon_\nu c_\nu M_\nu$.  On every finite grid
$G\subset\R_{>0}^2$,
\begin{equation}
 \mathbb E\sum_{g\in G}|P_\epsilon(g)|^2
 =\sum_\nu |c_\nu|^2\sum_{g\in G}|M_\nu(g)|^2.
 \label{pc:eq-random-sign-energy}
\end{equation}
At a fixed $g$, if $\sigma_g^2=\sum_\nu|c_\nu M_\nu(g)|^2$, the fourth-moment bound and
Paley--Zygmund give
\[
 \Pr\{|P_\epsilon(g)|\ge\sigma_g/\sqrt2\}\ge\frac1{12}.
\]
When $\sigma_g=0$ this is immediate; otherwise it is the usual Paley--Zygmund application to
$|P_\epsilon(g)|^2$.  This gives an RMS-scale baseline at each fixed point.  It does not
exclude rare simultaneous small-ball sign assignments, but any improvement must exploit more
structure than this second-moment calculation records.

\begin{statusbox}{The metric, Haar-phase, second-moment, and random-sign statements above are
self-contained paper proofs from report 6.  The phase-zero count is a repackaging of
\lean{ExactCounting}.  The finite positive-kernel bridge, exact weighted Fourier score, and
integer-phase hit bound are \statuslean{} in \lean{FejerKernelCertificate}; Haar integration
and Borel--Cantelli are not claimed there.  Neither almost-everywhere extinction nor fixed-mode
equidistribution is specialized to $\thetaq$.}
\end{statusbox}

\section{Seventh and eighth reports: compact laws and approximation rigidity}
\label{sec:reports78}

Reports 7 and 8 were produced independently from the same inherited unified report.  They are
therefore not successive revisions of one another.  This section audits each continuation on
its own terms and then reconciles the nonoverlapping conclusions.  Their directions are
complementary.  Report 7 asks what fixed measure-theoretic, profinite, and adelic observables
can see under a hypothetical counterexample.  Report 8 asks whether polynomial, rational, or
Hermite--Pad\'e approximation can exploit the exact integer lattice.  Neither closes the
conjecture: the first proves that every fixed compact factor is already maximally distributed,
while the second proves that the most natural approximation scales lose their gain when
denominators are cleared.

\subsection{Difference groups, triangle--Haar laws, and profinite Beatty factors}
\label{r78:compact-laws}

Assume in this subsection that the conjecture fails and let $\beta=B+\alpha$ be its least
nonintegral solution, where $B=\lfloor\beta\rfloor$ and $0<\alpha<1$.  Write
\[
 M=2^\beta=2^aW,\qquad A=3^\beta=3^bZ,
 \qquad W>1\text{ odd},\quad Z>1,\quad 3\nmid Z.
\]
The inherited structure theorem gives
$\Sol=\Nzero+\Nzero\beta$ and unique coordinates.  Taking differences produces an exact
measure-theoretic dichotomy:
\begin{equation}
 \Sol-\Sol=
 \begin{cases}
   \Z,&\text{if the conjecture holds},\\
   \Z+\Z\beta,&\text{if it fails}.
 \end{cases}
 \label{r78:eq-difference-group}
\end{equation}
Modulo $\Z$, the failure branch is the countable dense Haar-null subgroup
$\{k\alpha:k\in\Z\}\subset\mathbb T$.  Thus closure, Haar measure, and the arithmetic set itself
give three sharply different pictures of the same orbit.  This is the first reason that a
fixed weak-limit argument cannot settle the problem.

The macroscopic and microscopic coordinates can nevertheless be combined exactly.  Put
\[
 L_T=\{(n,k)\in\Nzero^2:n+\beta k\le T\},
 \qquad N_T=\#L_T=\frac{T^2}{2\beta}+O_\beta(T),
\]
and $\Delta=\{(u,v)\in\R_{\ge0}^2:u+v\le1\}$.  Report 7 proves that for every continuous
$F:\Delta\times\mathbb T\to\C$,
\begin{equation}
 \frac1{N_T}\sum_{(n,k)\in L_T}
 F\!\left(\frac nT,\frac{\beta k}{T},\fract{k\beta}\right)
 \longrightarrow
 2\int_\Delta\int_{\mathbb T} F(u,v,t)\,dt\,du\,dv.
 \label{r78:eq-triangle-haar}
\end{equation}
For a nonzero circle Fourier mode, Abel summation reduces the assertion to bounded geometric
partial sums of $e(hk\beta)$; the zero mode is an anisotropic Riemann sum.  In the variables
$H=u+v$ and $C=v/(u+v)$, the limiting density is $2H\,dH\,dC$, independently of the Haar
phase.  This is the precise independence statement; the separate Dirichlet description of
$(u,v,1-u-v)$ is merely another description of uniform triangle area, not a replacement for
the $(H,C)$ factorization.

The finite congruence refinement is stronger.  For every irrational $\alpha$, every $d\ge1$,
residues $r,j\pmod d$, and interval $I\subset\mathbb T$,
\begin{equation}
 \frac1K\#\{0\le k<K:k\equiv r\pmod d,
   \ \lfloor k\alpha\rfloor\equiv j\pmod d,
   \ \fract{k\alpha}\in I\}
 \longrightarrow\frac{|I|}{d^2}.
 \label{r78:eq-profinite-beatty}
\end{equation}
Indeed, after writing $k=d\ell+r$, the map
$x\mapsto(\lfloor dx\rfloor,\fract{dx})$ sends Lebesgue measure to uniform measure on
$\Z/d\Z$ times Lebesgue measure.  Equivalently,
\[
 k\longmapsto(k,\lfloor k\alpha\rfloor,\fract{k\alpha})
\]
is Haar distributed in $\widehat\Z^2\times\mathbb T$.  Every fixed family of congruence
conditions is therefore asymptotically independent of the real phase.  With the inherited
finite irrationality measure
$\distZ{h\beta}\ge c_Mh^{1-\mu_M}$, where we harmlessly enlarge $\mu_M$ so that $\mu_M>1$,
the same argument controls growing moduli only in the range
$d^{\mu_M+1}=o(K)$.  The true denominator moduli grow exponentially in $k$, so this effective
range remains far too small.

\subsection{Exact denominator drift and noncompact escape}
\label{r78:adelic-drift}

The compact arc coordinates
\[
 R_k=\frac{M^k}{2^{\lfloor k\beta\rfloor}},
 \qquad S_k=\frac{A^k}{3^{\lfloor k\beta\rfloor}}
\]
separate into primitive numerators and denominator exponents
\[
 q^{(2)}_k=\lfloor k\beta\rfloor-ka,
 \qquad q^{(3)}_k=\lfloor k\beta\rfloor-kb.
\]
Writing $t_k=\fract{k\beta}$ gives the exact identities
\begin{align}
 q^{(2)}_k&=k(\beta-a)-t_k,&
 q^{(3)}_k&=k(\beta-b)-t_k,\label{r78:eq-centered-drift}\\
 q^{(2)}_k-q^{(3)}_k&=k(b-a).\label{r78:eq-drift-difference}
\end{align}
Both centered valuation defects are therefore the same bounded quantity $-t_k$, rather than
independent noise.  Since $\beta-a=\log_2W>0$ and $\beta-b=\log_3Z>0$,
\begin{equation}
 |R_k|_2=2^{q^{(2)}_k}\longrightarrow\infty,
 \qquad |S_k|_3=3^{q^{(3)}_k}\longrightarrow\infty.
 \label{r78:eq-padic-escape}
\end{equation}
Thus the empirical measures escape every compact subset of $\Q_2\times\Q_3$.  Removing the
denominator powers leaves $(W^k,Z^k)$ in the compact group
$\mathcal H=\overline{\langle(W,Z)\rangle}\subset\Z_2^\times\times\Z_3^\times$; the pair
$(t_k,W^k,Z^k)$ is distributed according to Lebesgue measure times Haar measure on
$\mathcal H$.  In a moving frame the limiting measure is the pushforward of
$dt\,dm_\mathcal H$ under
\[
 (t,u,v)\longmapsto(t,-t,-t,u,v).
\]
This retains exactly the phase information that one-point compactification would erase.

The integer definition, the cast identity, the common centered defect, the cross-base
difference, the natural-exponent layer, and the resulting $2$- and $3$-adic escape are
\statuslean{} in \lean{AdelicDrift}.  The module deliberately defines the raw exponent in
$\Z$ before proving nonnegativity, so no hidden natural-subtraction assumption enters the
statement.

\subsection{Sturmian rigidity and fixed-observable blindness}
\label{r78:sturmian}

Choose a representative $t\in[0,1)$ of an initial circle point and set
\[
 \xi_j(t)=\mathbf1_{[1-\alpha,1)}(\fract{t+j\alpha}).
\]
Each bit is the carry in addition by $\alpha$, so the exact telescoping identity is
\begin{equation}
 \sum_{j=0}^{N-1}\xi_j(t)=\lfloor t+N\alpha\rfloor-\lfloor t\rfloor.
 \label{r78:eq-sturmian-prefix}
\end{equation}
For uniform $t$, this sum takes only $\lfloor N\alpha\rfloor$ and
$\lceil N\alpha\rceil$, with mean $N\alpha$ and variance
$\fract{N\alpha}(1-\fract{N\alpha})\le1/4$.  If
$\delta_\ell=\fract{\ell\alpha}$, then
\begin{equation}
 \operatorname{Cov}(\xi_0,\xi_\ell)
 =\max(0,\alpha-\delta_\ell)+\max(0,\alpha+\delta_\ell-1)-\alpha^2.
 \label{r78:eq-sturmian-covariance}
\end{equation}
The correlations do not decay.  The process is a zero-entropy, pure-point factor of an
irrational rotation; Bernoulli, Poisson, and central-limit heuristics are therefore the wrong
model for its carries.

More generally, let $G$ be a profinite group and let
$\pi:\widehat\Z^2\to G$ be continuous.  Combining
\eqref{r78:eq-triangle-haar} and \eqref{r78:eq-profinite-beatty} shows that for continuous
observables the scaled triangle coordinate, $\pi(k,\lfloor k\alpha\rfloor)$, and the circle
phase have the product limit.  The same holds for indicators whose discontinuity set has
limiting measure zero; this qualification is needed for Sturmian bits, which are not
continuous functions.  Consequently every fixed compact congruence or unit factor is blind
to the counterexample.  A successful observable must move with height, retain an unbounded
valuation, or resolve an exact shrinking target.

Two further paper-level calibrations reinforce the point.  First,
\[
 \frac1{\log X}\int_1^X
 \delta_{(\fract u,\fract{u^{\thetaq}})}\,\frac{du}{u}
 \stackrel{*}{\longrightarrow}m_{\mathbb T^2};
\]
nonzero Fourier modes follow from a first-derivative estimate.  Second, for
$E_m(\psi)=\{c\in[a,b]:\distZ{m^c}<\psi_m\}$ one has
$\lambda(E_m(\psi))\ll_{a,b}\psi_m$, and for $\psi_m=m^{-\tau}$,
\begin{equation}
 \dim_{\mathrm H}\limsup_m E_m(m^{-\tau})
 \le\min\!\left(1,\frac{b+1}{b+\tau}\right).
 \label{r78:eq-metric-dimension}
\end{equation}
This can place $\thetaq$ in a zero-dimensional exceptional set without excluding that named
number.  The local mass estimate near exact hits in report 7 is only a scale diagnostic unless
overlaps and uniform Taylor remainders are controlled, and is not promoted here to a theorem.
The near-Sidon statement for conditional truncations is likewise retained only as a
conditional consequence of the \statusknown{} Evertse--Schlickewei--Schmidt finite-rank unit-equation
theorem.

\subsection{Arithmetic resolution of normalized polynomial approximation}
\label{r78:approximation-tax}

Report 8 starts from a candidate $m\in[2^T,2^{T+1})$, puts
$u=m/2^T\in[1,2)$ and $n=m^{\thetaq}=3^Tu^{\thetaq}\in\N$, and keeps careful track of the
denominators introduced by normalization.  If $p\in\Q[t]$, $\deg p\le d$, and
$Hp\in\Z[t]$, then
\begin{equation}
 H2^{Td}\bigl(n-3^Tp(u)\bigr)\in\Z.
 \label{r78:eq-normalized-clearing}
\end{equation}
Hence an error smaller than $(H3^T2^{Td})^{-1}$ forces exact equality at every candidate in
the block.  For integral $A,B$ of degrees at most $d$, the correct rational-approximation
object is the numerator residual, not the quotient error:
\begin{equation}
 2^{Td}\bigl(B(u)n-3^TA(u)\bigr)
 =3^T2^{Td}\bigl(B(u)u^{\thetaq}-A(u)\bigr)\in\Z.
 \label{r78:eq-rational-residual-clearing}
\end{equation}
A small quotient can therefore be arithmetically useless when $B$ is large or has costly
coefficients.

The decisive analytic comparison is elementary.  For every $d\ge1$ and every real polynomial
$p$ of degree at most $d$,
\begin{equation}
 \|t^{\thetaq}-p(t)\|_{[1,2]}
 \ge c_d:=
 \frac{|\fall{\thetaq}{d+1}|\,2^{\thetaq-2d-2}}{(d+1)^{d+1}}.
 \label{r78:eq-finite-difference-lower}
\end{equation}
Indeed, with $h=(d+1)^{-1}$, the $(d+1)$-st forward difference of $p$ vanishes, the triangle
inequality costs $2^{d+1}$, and the iterated finite-difference mean-value theorem evaluates the
same difference of $t^{\thetaq}$ at an intermediate point.  For $d=0$, endpoint separation
gives the lower bound $1$.

Combining \eqref{r78:eq-normalized-clearing} and
\eqref{r78:eq-finite-difference-lower} yields the all-degree obstruction:
\begin{theorem}[Normalization-tax obstruction]
\label{r78:thm-normalization-tax}
If $T\ge4$, $d\ge0$, $H\ge1$, $\deg p\le d$, and $Hp\in\Z[t]$, then
\begin{equation}
 H3^T2^{Td}\|t^{\thetaq}-p(t)\|_{[1,2]}\ge1.
 \label{r78:eq-normalization-tax}
\end{equation}
Thus no normalized rational polynomial, in any degree, crosses the arithmetic forcing
threshold.
\end{theorem}

The proof reduces to $T=4,H=1$.  Writing $L_d=3^4 2^{4d}c_d$, the rational bounds
$3/2<\thetaq<8/5$ show $L_1>\cdots>L_5>4$, $L_6>1$, and
$L_{d+1}/L_d>1$ for $d\ge6$; in particular the exact calculation includes
$L_5>13328/3125$.  No finite search bound is used.  The forward-difference estimate,
polynomial degree floor, exact denominator clearing, and the all-degree tax theorem are
\statuslean{} in \lean{PolynomialApproximationTax}; the sharp Bernstein-ellipse rate and its
degree constant remain classical paper inputs.

Working in the physical variable avoids the factor $2^{Td}$, but then the same lower bound
requires degree
\[
 d\ge\frac{\log3}{\log(4e)}T-O(\log T)
 =0.460384\ldots T-O(\log T).
\]
Classical Bernstein approximation improves the slope to
$\log3/\log(3+2\sqrt2)=0.623238\ldots$.  Either slope exceeds the conditional block
population slope by a wide margin; this is a structural mismatch, not a numerical failure at
the beginning of the search range.

\subsection{Rational contact rigidity and the area-to-diameter collapse}
\label{r78:contact}

The local approximation audit has a sharper exact conclusion.  For a polynomial $P$ and a
rational point $(a,b)$, let $\operatorname{mult}_{(a,b)}P$ be the least total degree in the
Taylor expansion of $P(a+U,b+V)$.

\begin{theorem}[Rational contact equals algebraic multiplicity]
\label{r78:thm-contact-multiplicity}
Let $\tau$ be transcendental, $0\ne P\in\Q[X,Y]$, and
$(a,b)\in\Q_{>0}^2$.  If $s=\operatorname{mult}_{(a,b)}P$, then
\begin{equation}
 \operatorname{ord}_{z=0}P(ae^z,be^{\tau z})=s.
 \label{r78:eq-contact-multiplicity}
\end{equation}
\end{theorem}

If $H_s(U,V)=\sum_{j=0}^s h_jU^{s-j}V^j$ is the first nonzero homogeneous Taylor layer, the
coefficient of $z^s$ is
$H_s(a,b\tau)=\sum_jh_ja^{s-j}b^j\tau^j$.  This is a nonzero rational polynomial evaluated
at a transcendental number, so it cannot vanish.  In particular, for bidegree at most $(d,e)$
the contact order is at most $d+e$; equality is attained only by a rational scalar multiple of
$(X-a)^d(Y-b)^e$.  The square of $(d+1)(e+1)$ coefficient slots thus collapses to the diameter
$d+e$ over $\Q$.  Over $\Q(\tau)$ the ordinary distinct-frequency Vandermonde instead permits
order $(d+1)(e+1)-1$.

For not-both-zero $A,B\in\Q[t]$ of degree at most $d$ this gives
\[
 \operatorname{ord}_{t=1}\bigl(B(t)t^\tau-A(t)\bigr)\le d+1.
\]
If $A(1)=B(1)$ and $B(1)\ne0$, the intersection is already transverse because
\begin{equation}
 B(1)\tau+B'(1)-A'(1)\ne0.
 \label{r78:eq-pade-transverse}
\end{equation}
The rational-polynomial transcendence lemma, nonvanishing of the homogeneous leading
coefficient, the derivative identity, and the Pad\'e transversality endpoint are
\statuslean{} in \lean{RationalContactRigidity}.  The general Taylor-order packaging in
Theorem~\ref{r78:thm-contact-multiplicity} remains a self-contained paper proof.

Finally, for a finite support $S\subset\Nzero^2$ and nonzero integer coefficients,
\[
 F(x)=\sum_{(i,j)\in S}c_{ij}x^{i+j\thetaq}
\]
is integral at every candidate.  If $\|F\|_I<1$, it vanishes at all candidates in $I$, while
the accepted exponential-polynomial zero bound gives at most $|S|-1$ positive zeros.  Hence
\begin{equation}
 \#(\Cand\cap I)\le |S|-1.
 \label{r78:eq-muntz-certificate}
\end{equation}
This arithmetic M\"untz certificate is a direct corollary of the existing
\lean{ExponentialPolynomialZeros} machinery, not a new approximation theorem.  The genuinely
open target is to find an integral coefficient vector with norm below one after all
archimedean, denominator, and $p$-adic height costs are charged.  The schematic adelic height
functional proposed in report 8 records this target but is not itself a proved estimate.

\begin{statusbox}{Reports 7 and 8 are independent continuations from the common unified base.
The exact drift/escape package is \statuslean{} in \lean{AdelicDrift}; the finite-difference and
normalization-tax package is \statuslean{} in \lean{PolynomialApproximationTax}; the exact
carry, prefix, mean, variance, joint-probability, and covariance package is \statuslean{} in
\lean{SturmianProbability}; and the
homogeneous contact obstruction plus rational Pad\'e transversality are \statuslean{} in
\lean{RationalContactRigidity}.  The triangle--Haar, profinite Beatty, logarithmic Haar,
Hausdorff-dimension, full contact-order, and area-to-diameter statements above have
self-contained paper proofs.  The effective discrepancy uses the inherited
Baker--Matveev interface, the near-Sidon law uses the external finite-rank unit-equation
theorem, and the sharp Bernstein rate is classical.  The reported $2^{29}$ computation is not
promoted: this unified report retains its reproducible $2^{27}$ trust frontier.}
\end{statusbox}

\section{Ninth report: analytic common divisors and natural boundaries}
\label{sec:report9}

Report 9 is another independent continuation from the common unified base.  Its approximation
section overlaps report 8 and is not repeated.  Its distinct contribution is a three-part
analytic reformulation: cancel the automatic integer zeros of the two sine compositions,
identify the resulting conditional zero lattice and its canonical divisor, and encode the
synchronized denominator exponents by integer power series with a natural boundary.  A final
compact-height count sharpens the earlier qualitative compact-orbit equivalence.

\subsection{Removing the automatic zeros and counting the conditional divisor}
\label{r9:common-divisor}

Put
\[
 F_a(z)=\sin(\pi a^z),\qquad a\in\{2,3\}.
\]
The accepted complex-localization theorem already shows that every common zero of $F_2$ and
$F_3$ is real and simple.  The nonnegative integers are automatic common zeros.  They can be
removed without introducing new ones by filling the removable singularities in
\begin{equation}
 G_a(z)=\Gamma(-z)\sin(\pi a^z).
 \label{r9:eq-deintegralized}
\end{equation}
Indeed, at $n\in\Nzero$ one has
\begin{equation}
 G_a(n)=\frac{\pi(\log a)a^n}{n!}(-1)^{n+1+a^n}\ne0,
 \label{r9:eq-integer-value}
\end{equation}
so $G_2$ and $G_3$ have a common zero exactly when the conjecture has a nonintegral solution.

Under failure, let $\beta$ be the least nonintegral solution.  The exact conditional monoid
classification identifies the de-integralized common zero set as
\[
 \Lambda_\beta=\{n+k\beta:n\in\Nzero,\ k\in\N\}.
\]
Its representation is unique, and elementary lattice counting gives
\begin{equation}
 N_\beta(R):=\#(\Lambda_\beta\cap[0,R])
 =\frac{R^2}{2\beta}+O_\beta(R).
 \label{r9:eq-quadratic-zero-count}
\end{equation}
Consequently its exponent of convergence is exactly two.  With
$E_2(w)=(1-w)\exp(w+w^2/2)$, the genus-two product
\begin{equation}
 P_\beta(z)=\prod_{\lambda\in\Lambda_\beta}E_2(z/\lambda)
 \label{r9:eq-canonical-product}
\end{equation}
converges locally uniformly, has precisely those simple zeros, and divides both $G_2$ and
$G_3$ in the ring of entire functions.  Conversely, any common nonunit entire divisor has a
zero and hence gives a nonintegral solution.  Thus the conjecture is equivalently the absence
of a common nonunit entire divisor.

This is not an ordinary finite-order zero-counting problem: along suitable horizontal lines,
$\log M(r,G_a)\ge c_a a^r$ for infinitely many $r$.  The conditional divisor has only
quadratic zero density inside functions of infinite order.  It is also too large for a
bounded-characteristic function on the right half-plane, since partial summation gives
\begin{equation}
 \sum_{\lambda\in\Lambda_\beta}\frac{\lambda}{1+\lambda^2}=\infty.
 \label{r9:eq-blaschke-divergence}
\end{equation}
This supplies a precise but still open analytic target: a growth-controlled common divisor or
B\'ezout theorem for the pair $G_2,G_3$.  Canonical-product convergence and the
bounded-characteristic/Blaschke implication are classical analytic interfaces, not new Lean
claims.

\subsection{The synchronized word has a natural boundary}
\label{r9:natural-boundary}

Write $\beta=B+\alpha$ with $0<\alpha<1$, and define
\[
 e_k=\lfloor k\alpha\rfloor,\qquad s_k=e_{k+1}-e_k\in\{0,1\}.
\]
For the formal power series
\[
 \mathcal S(X)=\sum_{k\ge0}s_kX^k,\qquad
 \mathcal E(X)=\sum_{k\ge1}e_kX^k,
\]
telescoping gives the exact identity
\begin{equation}
 X\mathcal S(X)=(1-X)\mathcal E(X).
 \label{r9:eq-carry-floor-series}
\end{equation}
Both analytic series have radius one.  The binary Sturmian word $(s_k)$ is not eventually
periodic because $\alpha$ is irrational.  A rational finite-alphabet power series has an
eventually periodic coefficient sequence; hence $\mathcal S$ is not rational.  The classical
P\'olya--Carlson dichotomy now makes the unit circle a natural boundary for $\mathcal S$, and
\eqref{r9:eq-carry-floor-series} transfers that boundary to $\mathcal E$.

The exact reduced denominator exponents have the form
\[
 q_{r,k}=rk+e_k,\qquad
 \mathcal Q_r(X)=\frac{rX}{(1-X)^2}+\mathcal E(X).
\]
Thus every $\mathcal Q_r$ has the same natural boundary, and
\begin{equation}
 \mathcal Q_r-\mathcal Q_s=\frac{(r-s)X}{(1-X)^2}.
 \label{r9:eq-denominator-difference}
\end{equation}
The carry telescope, formal-series identity, denominator decomposition, twice-cleared defect,
and cross-base cancellation are \statuslean{} in \lean{SturmianGeneratingSeries}.  The
natural-boundary conclusion itself invokes the classical P\'olya--Carlson theorem.  It is a
one-slope diagnostic, not a contradiction: every irrational rotation already has this
boundary, so a successful continuation theorem must use both prime supports together.

\subsection{A logarithmic compact-height criterion}
\label{r9:compact-height}

Let
\[
 \mathcal C_{2,3}=\{(2^t,3^t):0\le t<1\}
 \cap\bigl(\Z[1/2]\times\Z[1/3]\bigr),
\]
and for a reduced rational pair define
$H(a/b,c/d)=\max(|a|,b,|c|,d)$.  Write $N_{2,3}(H)$ for the number of points of
$\mathcal C_{2,3}$ of height at most $H$.  The accepted compact-orbit theorem gives the
qualitative equivalence $\mathcal C_{2,3}=\{(1,1)\}$ if and only if the conjecture holds.
Report 9 adds the exact growth dichotomy
\begin{equation}
 N_{2,3}(H)=
 \begin{cases}
  1,&\text{if the conjecture holds},\\
  \asymp_\beta\log H,&\text{if it fails}.
 \end{cases}
 \label{r9:eq-compact-height-dichotomy}
\end{equation}
In the failure branch, the distinct points
$(2^{\{k\alpha\}},3^{\{k\alpha\}})$ exhaust the localized arc, and their reduced numerator
and denominator heights grow between two fixed exponentials in $k$.  Therefore each of
\[
 N_{2,3}(H)=O(1),\qquad N_{2,3}(H)=o(\log H)
\]
is equivalent to the conjecture.  Compact normalization has removed the automatic integer
direction from the inherited $\asymp(\log H)^2$ unbounded count and leaves precisely one
conditional generator.

\begin{statusbox}{Report 9 was audited independently from the same common base as the other
reports.  The formal generating-series algebra is \statuslean{} in
\lean{SturmianGeneratingSeries}.  The Gamma cancellation, quadratic lattice count, compact
height dichotomy, and Blaschke calculation have self-contained paper proofs; canonical-product
and P\'olya--Carlson steps are classical interfaces.  The duplicated Bernstein approximation
material is represented once in Section~\ref{sec:reports78}.  The inherited reported
$2^{29}$ scan is not promoted; the unified trust frontier remains the reproducible $2^{27}$
computation.}
\end{statusbox}

\section{Tenth report: matrix powers, Loewner certificates, and the one-half barrier}
\label{sec:report10}

Report 10 is an independent matrix-theoretic continuation from the common base.  It produces
one exact reformulation and a family of genuine determinant certificates, but it also proves
that polynomial representations of the original matrix carry no information beyond the two
assumed integer outputs.  The optimized determinant family reaches a limiting one-half height
exponent and therefore remains on the same side of the four-exponentials wall.

\subsection{One fixed integral matrix}
\label{r10:fixed-matrix}

Consider
\begin{equation}
 A_0=\begin{pmatrix}0&-6\\1&5\end{pmatrix}.
 \label{r10:eq-matrix}
\end{equation}
Its eigenvalues are $2,3$, and the complementary integral spectral projectors are
$P_2=3I-A_0$ and $P_3=A_0-2I$.  Real functional calculus therefore gives
\begin{equation}
 A_0^t=2^tP_2+3^tP_3
 =(3\,2^t-2\,3^t)I+(3^t-2^t)A_0.
 \label{r10:eq-real-power}
\end{equation}
Consequently
\begin{equation}
 A_0^t\in M_2(\Z)
 \quad\Longleftrightarrow\quad
 2^t,3^t\in\Z.
 \label{r10:eq-integral-power-iff}
\end{equation}
For the nontrivial direction, if $B=A_0^t$ is integral then
$c=B_{21}=3^t-2^t$ and $d=B_{11}=3\,2^t-2\,3^t$ are integers, while
$2^t=d+2c$ and $3^t=d+3c$.  Equivalently,
\[
 \text{Alaoglu--Erd\H{o}s}
 \quad\Longleftrightarrow\quad
 \{t\in\R:A_0^t\in M_2(\Z)\}=\Nzero.
\]
The concrete matrix, all four entry formulas, the exact entry-integrality equivalence, and
this conjectural equivalence are \statuslean{} in \lean{IntegralPowerMatrix}.

This reformulation is exact but intentionally not advertised as extra arithmetic.  The
unimodular matrix
$\bigl(\begin{smallmatrix}-3&-2\\1&1\end{smallmatrix}\bigr)$ diagonalizes $A_0$ over
$\Z$, so the lattice already splits into its $2$- and $3$-eigensummands.  Every polynomial,
tensor, symmetric, or exterior-power construction from $(A_0,A_0^t)$ is a polynomial in the
two integers $(2^t,3^t)$ and their monomials.  It cannot recover the common real time unless it
introduces analytic dependence, a denominator-sensitive inequality, or a new logarithmic
direction.

\subsection{Cross-Loewner and generalized-alternant certificates}
\label{r10:loewner}

For distinct positive nodes $z_1,\ldots,z_N$ and distinct real exponents
$\gamma_1,\ldots,\gamma_N$, strict total positivity of $e^{uv}$ gives
\[
 \det(z_i^{\gamma_j})\ne0,
\]
with positive sign when both lists are increasing.  If $t$ is a nonintegral solution, the
two-block exponent list
\begin{equation}
 0,1,\ldots,r-1,\quad t,t+1,\ldots,t+r-1
 \label{r10:eq-two-block-exponents}
\end{equation}
is distinct.  At $2r$ distinct nodes in the $2,3$-smooth semigroup every entry is integral,
so the alternant is a nonzero integer.

The same certificate admits an exact divided-difference form.  For disjoint $r$-tuples
$X=(x_i)$ and $Y=(y_j)$ and a function $f$ on their union, put
\[
 L_f(X,Y)_{ij}=\frac{f(x_i)-f(y_j)}{x_i-y_j}.
\]
If $M_f(X,Y)$ is the $2r$-row alternant with columns
$1,z,\ldots,z^{r-1},f(z),zf(z),\ldots,z^{r-1}f(z)$, then
\begin{equation}
 \det M_f(X,Y)=(-1)^{r(r+1)/2}
 \left(\prod_{i,j}(x_i-y_j)\right)\det L_f(X,Y).
 \label{r10:eq-sylvester-loewner}
\end{equation}
For $f(u)=u^t$ and smooth integer nodes, the right-hand side is therefore a nonzero integer.
Writing $t=k+\alpha$ with $0<\alpha<1$, the fractional-power kernel
\[
 K_\alpha(u,v)=\frac{u^\alpha-v^\alpha}{u-v}
\]
has the positive Stieltjes representation
\[
 K_\alpha(u,v)=\frac{\sin(\pi\alpha)}\pi
 \int_0^\infty\frac{\lambda^\alpha}{(u+\lambda)(v+\lambda)}\,d\lambda.
\]
Thus its ordered cross minors are strictly positive.  Combining positivity with
\eqref{r10:eq-sylvester-loewner} gives the explicit positive integer
\begin{equation}
 \left(\prod_i x_i^k y_i^k\right)
 \left|\prod_{i,j}(x_i-y_j)\right|
 \det(K_\alpha(x_i,y_j))\in\Z_{>0}.
 \label{r10:eq-positive-loewner-integer}
\end{equation}
The positivity holds for every $0<\alpha<1$; all arithmetic content is in the displayed
normalization.  These finite determinant identities have paper proofs here.  They are not
silently conflated with the already formal generalized-Vandermonde positivity theorem.

\subsection{Quantitative repulsion and the optimized barrier}
\label{r10:one-half}

Let $z_i=e^{u_i}$ with pairwise distinct $u_i\in[T,T+\delta]$, $0<\delta\le1$, and
pairwise distinct exponents $0\le\gamma_j\le B$.  (With repeated nodes or exponents the
determinant vanishes and the upper bound is trivial.)  Put $\Sigma=\sum_j\gamma_j$,
$D=N(N-1)/2$, and $B_*=\max(1,B)$.  Newton divided differences followed by Hadamard's
inequality give the explicit upper bound
\begin{equation}
 \left|\det(z_i^{\gamma_j})\right|
 \le e^{T\Sigma}\delta^D
 \frac{N^{N/2}e^{NB}B_*^D}{\prod_{\ell=0}^{N-1}\ell!}.
 \label{r10:eq-alternant-upper}
\end{equation}
Together with the nonzero-integer lower bound, this excludes a sufficiently tight cluster of
$2r$ smooth nodes.  For the equal two-block list
\eqref{r10:eq-two-block-exponents} it gives an explicit bound of the shape
\[
 \delta\ge C_{r,t}^{-1/[r(2r-1)]}
 R^{-(t+r-1)/(2r-1)}.
\]
The exponent tends to $1/2$ as $r\to\infty$, much smaller than the logarithmic spacings
available from counting smooth numbers.

Allowing unequal block sizes does not improve the limiting ratio.  With $p+q=N$ and exponent
set $0,\ldots,p-1,t,\ldots,t+q-1$, the total exponent is
\begin{equation}
 \Sigma(p,q;t)=\frac{p(p-1)}2+qt+\frac{q(q-1)}2
 =D+q(q+t-N),
 \label{r10:eq-mixed-sum}
\end{equation}
so
\begin{equation}
 \min_{p+q=N}\frac{\Sigma(p,q;t)}D
 =\frac12+O_t(N^{-1}).
 \label{r10:eq-half-barrier}
\end{equation}
Initial consecutive shifts already minimize the height inside each coset.  Translating signed
exponents into the first quadrant restores exactly the denominator height that cancellation
appeared to save.  This is the matrix one-half barrier.

There is a dual exact identity for candidate exponents.  Let $d\ge1$, let
$p,q\in\N$ with $q>0$, and put $\varepsilon=qt-p$ and
$\tau_i=(d-1-i)p+iqt$ for $0\le i<d$.  Then at distinct smooth bases $s_j$,
$0\le j<d$,
\begin{equation}
 \det(s_j^{\tau_i})
 =\left(\prod_j s_j^{(d-1)p}\right)
   \prod_{j<k}(s_k^\varepsilon-s_j^\varepsilon).
 \label{r10:eq-approx-vandermonde}
\end{equation}
Its nonzero-integral lower bound yields an explicit but exponentially weak irrationality
estimate for $|qt-p|$.  Higher dimension organizes the elementary integer
$2^{qt}-2^p$ but does not change its exponential scale.  The affine-exponent determinant,
its real-power specialization, the exact semigroup exponent identity, and
\eqref{r10:eq-approx-vandermonde} are \statuslean{} in \lean{CandidateAlternant}; the
mean-value upper bound and resulting irrationality estimate remain the paper layer.

\begin{statusbox}{Report 10 was audited independently from the common base.  The exact fixed
matrix reformulation is \statuslean{} in \lean{IntegralPowerMatrix}.  Generalized-Vandermonde
nonvanishing and ordered-chamber positivity are separately \statuslean{} in the inherited
determinant modules, and the exact rational-approximation alternant is \statuslean{} in
\lean{CandidateAlternant}.  The Sylvester--Loewner factorization, positive normalized
determinant, explicit alternant bound, and optimized one-half barrier are retained as paper
theorems.  No claim is made that formal representation theory manufactures a third
exponential, or that matrix positivity by itself selects integral exponents.}
\end{statusbox}

\section{Eleventh report: finite-place transversality and local rank-drop certificates}
\label{sec:report11}

Report 11 is an independent $p$-adic and matrix-coefficient continuation from the common
base.  Its real prime-symbol polynomial, rational matrix-coefficient obstruction, and sharp
support bound overlap the already reconciled reports 1--4.  They are represented once by
\lean{FormalPrimeDeterminant}, \lean{RestrictedMatrixCoefficient}, and
\lean{SparsePowerCurve}.  The genuinely distinct layer compares the singular real logarithm
matrix with its finite-place analogues, isolates the first Fermat-quotient obstruction, and
shows why ordinary real rational approximation does not propagate to a local rank drop.

\subsection{The structural primes are locally transverse}
\label{r11:structural-primes}

For the Iwasawa logarithm $\log_p:\Q_p^*\to\Q_p$, normalized by $\log_p p=0$, define
\begin{equation}
 L_p(M,A)=
 \begin{pmatrix}\log_p2&\log_pM\\\log_p3&\log_pA\end{pmatrix},
 \qquad
 D_p(M,A)=\log_pM\log_p3-\log_pA\log_p2.
 \label{r11:eq-local-matrix}
\end{equation}
On positive rationals the kernel is exactly
\begin{equation}
 \log_pu=0\quad\Longleftrightarrow\quad u=p^k
 \text{ for some }k\in\Z.
 \label{r11:eq-iwasawa-kernel}
\end{equation}
Hence a nonintegral counterexample pair $(M,A)=(2^\beta,3^\beta)$ satisfies
\begin{align}
 D_2(M,A)&=\log_2M\log_23\ne0,\\
 D_3(M,A)&=-\log_3A\log_32\ne0.
 \label{r11:eq-structural-nonzero}
\end{align}
The real matrix has rank one, but both finite-place matrices at the primes defining the
problem have rank two.  If, for a prime $p$, the pair $(M,A)$ belonged to the closure of
$\{(2^n,3^n):n\in\Z\}$ in $(\Q_p^*)^2$, continuity of $\log_p$ would instead force
$D_p(M,A)=0$.  Thus a real-to-local closure principle at either structural prime would prove
the conjecture, but no such transfer theorem is known.

This is not merely a change of language.  It reverses the naive local expectation: a
hypothetical counterexample is singular only at the real place among the three places that
are directly visible from its definition.  Any adelic proof must compare a real zero with
proven local nonzeros, not copy the real relation into $\Q_2$ or $\Q_3$.

\subsection{The Fermat-quotient first layer}
\label{r11:fermat}

For an odd prime $p\nmid u$, set
\[
 q_p(u)=\frac{u^{p-1}-1}{p}\pmod p.
\]
The first Iwasawa-logarithm digit is
\begin{equation}
 \frac{\log_pu}{p}\equiv-q_p(u)\pmod p.
 \label{r11:eq-log-fermat}
\end{equation}
Therefore, for $p\ge5$ and $p\nmid6MA$, the determinant
\begin{equation}
 \Delta_p(M,A)=q_p(M)q_p(3)-q_p(A)q_p(2)\pmod p
 \label{r11:eq-fermat-determinant}
\end{equation}
satisfies
\begin{equation}
 \frac{D_p(M,A)}{p^2}\equiv\Delta_p(M,A)\pmod p.
 \label{r11:eq-local-first-layer}
\end{equation}
Thus $\Delta_p\ne0$ certifies the exact first possible valuation
$v_p(D_p)=2$.  Conversely, for every integer-exponent pair
$(M,A)=(2^n,3^n)$,
\[
 q_p(u^n)\equiv n q_p(u)\pmod p
\]
makes $\Delta_p(M,A)=0$.  A local rank drop must therefore pass a generalized Wieferich
congruence and then an infinite tower of higher congruences.  The finite near-miss scan in
report 11 is retained only as a software diagnostic; it is not evidence that every primitive
pair has a bounded detecting prime.

\subsection{Real approximants and the exact adelic wedge}
\label{r11:adelic-wedge}

Let $r/s$ be a rational approximation to $\beta$, $s>0$, and put
\[
 R_2(r,s)=\frac{M^s}{2^r},\qquad
 R_3(r,s)=\frac{A^s}{3^r}.
\]
At the real place the errors are perfectly synchronized:
\[
 \log R_2=(s\beta-r)\log2,\qquad
 \log R_3=(s\beta-r)\log3.
\]
At a finite place set
\[
 E_p(r,s)=
 (s\log_pM-r\log_p2,\ s\log_pA-r\log_p3).
\]
Direct expansion gives the exact wedge identity
\begin{equation}
 E_{p,1}(r,s)\log_p3-E_{p,2}(r,s)\log_p2=sD_p(M,A).
 \label{r11:eq-adelic-wedge}
\end{equation}
At the structural primes the numerator $r$ disappears from one coordinate:
\begin{align}
 \log_2R_2(r,s)&=s\log_2M,&
 v_2(\log_2R_2)&=v_2(s)+v_2(\log_2M),\\
 \log_3R_3(r,s)&=s\log_3A,&
 v_3(\log_3R_3)&=v_3(s)+v_3(\log_3A).
 \label{r11:eq-structural-errors}
\end{align}
Choosing a better continued-fraction numerator improves the real approximation but has no
effect on these two local depths; only divisibility of the denominator $s$ helps.  This exact
anti-transfer explains why an archimedean convergent does not automatically produce a local
near-singularity.

\subsection{Cascade depth and the Kummer norm diagnostic}
\label{r11:cascade-kummer}

The report also records two useful no-go estimates.  First, suppose each of the first $J(T)$
active tropical layers contributes at most $CT^{3/2}$ fixed-base logarithmic factors, each of
valuation at most $C'(\log T)^2$.  Their total exposed valuation is then
\begin{equation}
 O\!\left(J(T)T^{3/2}(\log T)^2\right).
 \label{r11:eq-cascade-depth}
\end{equation}
Against $H_T\asymp T^{5/2}$ this is lower order whenever
$J(T)=o(T/(\log T)^2)$.  This is a conditional bookkeeping proposition for a factorized
layer model, not an upper bound on arbitrary determinant cancellation; a global higher
initial form may behave differently.

Second, for positive real $N$th roots define the algebraic integer
\[
 \Xi_N=(M^{1/N}-1)(3^{1/N}-1)
      -(A^{1/N}-1)(2^{1/N}-1).
\]
The determinant relation cancels the quadratic term and leaves
\begin{equation}
 \Xi_N=
 \frac{\beta\log2\log3\,(\beta-1)(\log2-\log3)}{2N^3}
 +O_{M,A}(N^{-4}).
 \label{r11:eq-kummer-asymptotic}
\end{equation}
Thus the distinguished embedding is polynomially small and nonzero for large $N$.  The field
degree and varying, generally uncontrolled root-of-unity twists (subject to the algebraic
relations among the radicands) grow with $N$, however, so a crude product-formula
bound becomes worse rather than better.  A viable algebraic-integer construction would need
simultaneous smallness in many conjugates or exponential smallness at one place.

\begin{statusbox}{Report 11 was audited independently from the common base.  Its real
prime-symbol and matrix-coefficient layers are already \statuslean{} in
\lean{FormalPrimeDeterminant} and \lean{RestrictedMatrixCoefficient}; its sparse-grid core is
\statuslean{} in \lean{SparsePowerCurve}.  The exact Fermat-quotient divisibility identity,
power law modulo $p$, determinant interface, and its vanishing on $(2^n,3^n)$ are
\statuslean{} in \lean{FermatQuotientDeterminant}.  The passage from that residue determinant
to the Iwasawa logarithm uses the classical congruence \eqref{r11:eq-log-fermat}; the
structural-prime theorem uses the classical Iwasawa-logarithm kernel.  The exact wedge,
conditional cascade-depth, and Kummer expansion remain elementary paper results.  No
$p$-adic Matrix Coefficient or real-to-local transfer theorem is asserted.}
\end{statusbox}

\section{Twelfth report: prime switches and microscopic resonance}
\label{sec:report12}

Report 12 is another independent continuation from the common unified base.  It returns to
the historical origin of the problem: simultaneous changes in the prime exponents of
colossally abundant numbers.  This is a valuable structured test family, but its logical
scope is narrower than the integer-output conjecture.  A switch produces rational values
$p^\varepsilon$ and $q^\varepsilon$; excluding all switches would prove only a special case
of the stronger rational-output statement.  No reduction from arbitrary rational outputs,
or from this switch family, to Conjecture~\ref{conj:AE} is known.

\subsection{Exact switch heights and logarithmic means}
\label{r12:switch-logmean}

For a prime $p$ and an integer $a\ge1$, set
\begin{equation}
 H_{p,a}=p+p^2+\cdots+p^a,
 \qquad
 \varepsilon_{p,a}=\frac{\log(1+H_{p,a}^{-1})}{\log p}.
 \label{r12:eq-switch-height}
\end{equation}
Comparing the two consecutive local factors in
$\sigma(n)/n^{1+\varepsilon}$ gives the exact tie identity
\begin{equation}
 p^{\varepsilon_{p,a}}=
 \frac{p^{a+1}-1}{p^{a+1}-p}
 =1+\frac1{H_{p,a}}.
 \label{r12:eq-switch-output}
\end{equation}
For fixed $a$ this parameter decreases with $p$, and for fixed $p$ it decreases with $a$.
If $p<q$ and $\varepsilon_{p,a}=\varepsilon_{q,b}$, then
\begin{equation}
 H_{p,a}>H_{q,b},\qquad a>b.
 \label{r12:eq-antichain}
\end{equation}
Thus simultaneous switches form an antichain in prime size and local exponent.

The analytic center of the calculation is the consecutive logarithmic mean
\[
 R(H)=\frac1{\log(1+1/H)},
 \qquad
 J(H)=H+\frac12-\frac1{12H}.
\]
For every real $H\ge1$ one has the globally valid two-sided enclosure
\begin{equation}
 J(H)<R(H)<J(H)+\frac1{24H^2}.
 \label{r12:eq-logmean-enclosure}
\end{equation}
The proof differentiates two rational majorants for $\log(1+x)$ on $0<x\le1$; it is an
inequality with an exact remainder, not merely an asymptotic expansion.  The definitions and
the complete enclosure \eqref{r12:eq-logmean-enclosure} are \statuslean{} in
\lean{ConsecutiveLogMeanBounds}.

At a collision, $R(H_{p,a})\log p=R(H_{q,b})\log q$.  Writing
$H=H_{p,a}$ and $K=H_{q,b}$, \eqref{r12:eq-logmean-enclosure} therefore gives
\begin{equation}
 \left|J(H)\log p-J(K)\log q\right|
 <\frac{\log p}{24H^2}+\frac{\log q}{24K^2}.
 \label{r12:eq-second-order-resonance}
\end{equation}
This is a highly structured Dirichlet-scale approximation: the coefficients are geometric
sums, and exploiting that structure rather than merely invoking a general irrationality
measure is essential.

\subsection{A microscopic prescribed-prime window}
\label{r12:microscopic-window}

For fixed $b\ge1$, put
\[
 H_b(t)=t+t^2+\cdots+t^b,
 \qquad \Phi_b(t)=J(H_b(t))\log t.
\]
Since $\Phi_b'(t)>bt^{b-1}\log t$ on $t>1$, there is a unique real
$Q_{p,a,b}>1$ satisfying
\[
 \Phi_b(Q_{p,a,b})=J(H_{p,a})\log p.
\]
If primes $p<q$ switch simultaneously, the mean-value theorem and
\eqref{r12:eq-second-order-resonance} imply
\begin{equation}
 |q-Q_{p,a,b}|<
 \frac{\frac{\log p}{24H_{p,a}^{2}}+
       \frac{\log q}{24H_{q,b}^{2}}}
      {b(q-\frac12)^{b-1}\log(q-\frac12)}
 =O_b(q^{1-3b}).
 \label{r12:eq-microscopic-window}
\end{equation}
The interval has length less than one.  Already for $b=1$ its radius is $O(q^{-2})$.
Consequently modern prime-in-short-interval theorems, whose intervals still have positive
power length, cannot decide the switch: before primality enters, the nearest integer has
already been prescribed.  The report's suggested modernization of the 1944 prime-block
argument is retained only as a schematic transference of known prime-distribution input; its
required uniform error estimates are not reconstructed here.  Likewise, the first-order
formula for the location of $Q_{p,a,b}$ is a useful calibration, but its displayed proof in
the source report omits a bootstrap and is not used as a theorem in this synthesis.

\subsection{The exact $2$--$3$ convergent sieve}
\label{r12:two-three-sieve}

For a possible switch between $2$ at level $a$ and $3$ at level $b$, define
\begin{equation}
 P=2^{a+2}-3,
 \qquad Q=3^{b+1}-2.
 \label{r12:eq-PQ}
\end{equation}
The second-order enclosure, together with the elementary bounds
$3/2<\thetaq<8/5$, yields the exact necessary condition
\begin{equation}
 0<\thetaq-\frac PQ<\frac1{2Q^2}.
 \label{r12:eq-legendre-window}
\end{equation}
After reduction, $P/Q$ must therefore be a continued-fraction convergent of $\thetaq$.
Equivalently, a collision must hit an exponentially shrinking shifted-power target
\[
 0<\thetaq(3^{b+1}-2)-(2^{a+2}-3)
 <\frac1{2(3^{b+1}-2)}.
\]
This converts the transcendental equality into a discrete $S$-unit sieve, but no uniform
argument over all convergents is currently known.

The literal exact-integer verifier in report 12 checks this necessary condition for
$1\le b\le10000$.  It certifies $10000$ unambiguous floors, finds the sole power-shape hit
$(a,b,P,Q)=(5,3,125,79)$, and rejects that hit by the stronger Legendre window; hence there
are no survivors in the tested range.  The verifier source has SHA-256
\begin{center}
\sha{a7aef69a62ce5b07bd8f909966244b3a5e92e23ac9bb2e1b951433d1f6902ce8}.
\end{center}
An independent rerun reproduced every mathematical field.  The source report's JSON digest
is not used as a certificate because its payload includes nondeterministic elapsed time.
This remains \statuscomp{}, not a kernel proof and not a finite test of the original
integer-output conjecture.

\subsection{What prime distribution can and cannot force}
\label{r12:prime-barriers}

For fixed integers $U,V\ge2$, define the universal Fermat-synchronized divisor
\[
 F_{U,V}(n)=
 \prod_{\substack{\ell\ {\rm prime},\ \ell\nmid UV\\ \ell-1\mid n}}\ell.
\]
Fermat's theorem gives
\begin{equation}
 F_{U,V}(n)\mid\gcd(U^n-1,V^n-1),
 \qquad
 \log F_{U,V}(n)\le\tau(n)\log(n+1)=n^{o(1)}.
 \label{r12:eq-fermat-sync}
\end{equation}
Along $n_y=\operatorname{lcm}(1,\ldots,\lfloor y\rfloor)$, the prime number theorem and the
standard divisor estimate still leave $\log F_{U,V}(n_y)=n_y^{o(1)}=o(n_y)$.  This is far
below the exponential common-divisor growth that would contradict the
Bugeaud--Corvaja--Zannier bound for multiplicatively independent $U,V$.  Generic prime
distribution therefore cannot manufacture the missing real-to-finite-place correlation.
It becomes relevant only after another argument has produced a rigid prescribed integer,
such as the nearest integer in \eqref{r12:eq-microscopic-window} or a convergent satisfying
\eqref{r12:eq-legendre-window}.

\begin{statusbox}{Report 12 was audited independently from the same common base as all other
numbered reports.  The exact logarithmic-mean enclosure is \statuslean{} in
\lean{ConsecutiveLogMeanBounds}.  The switch formula, antichain, resonance, microscopic
window, convergent sieve, and Fermat-synchronization barrier are self-contained paper
arguments.  The low-level switch exclusion is \statuscomp{} with the scope and software
qualification stated above.  Modern short-interval results are scale calibrations only, and
the claimed transference to the complete 1944 block argument is not promoted to a theorem.
Nothing in this section reduces the rational switch problem to the original integer-output
conjecture.}
\end{statusbox}

\section{Thirteenth report: exterior powers and boundary realizations}
\label{sec:report13}

Report 13 is independent of report 10 and of every other numbered continuation: it starts
from the same inherited unified base and pursues a different matrix-analysis program.  Its
new content is not another reformulation of the Matrix Coefficient principle.  It asks what
finite-dimensional matrix phenomena a hypothetical counterexample can realize, whether
proper exterior powers improve the determinant budget, and whether a $K^2$-point
multiplicative grid can be compressed to its $O(K)$ boundary data.

Throughout this section assume hypothetically that
\begin{equation}
 M=2^\beta\in\Z_{>0},\qquad A=3^\beta\in\Z_{>0},
 \qquad \beta\notin\Q,
 \label{r13:eq-counterexample}
\end{equation}
and retain $\thetaq=\log_2 3$.  Then
\begin{equation}
 G=\{2^rM^s:r,s\in\Z\}\subset\Q_{>0}
 \label{r13:eq-dense-group}
\end{equation}
is dense in $\R_{>0}$ and every $g\in G$ has rational $\thetaq$-th power.

\subsection{Rational conditioning does not create rational rank}
\label{r13:rational-conditioning}

The singular logarithm matrix factors as
\[
 L=\begin{pmatrix}\log2&\log M\\ \log3&\log A\end{pmatrix}
 =\binom{\log2}{\log3}(1,\beta).
\]
For a nonzero $c\in\R^2$ and $B\ge1$, set
\[
 \mu_c(B)=\min_{0\ne z\in\Z^2,\ \|z\|_\infty\le B}|z^{\mathsf T}c|,
 \qquad
 \eta_L(B)=\min_{\substack{0\ne u,v\in\Z^2\\
   \|u\|_\infty,\|v\|_\infty\le B}}|u^{\mathsf T}Lv|.
\]
The rank-one factorization gives the exact identity
\begin{equation}
 \eta_L(B)=
 \mu_{(\log2,\log3)}(B)\,\mu_{(1,\beta)}(B).
 \label{r13:eq-bilinear-minimum}
\end{equation}
Dirichlet approximation supplies $\eta_L(B)\ll_M B^{-2}$, while effective linear forms in
logarithms give only a computable polynomial lower bound.  Approximate zero coefficients after
bounded rational basis changes are therefore plentiful and entirely compatible with exact
nonvanishing.  The missing conclusion remains rational orthogonality in one of the two
factors, exactly as in \lean{RestrictedMatrixCoefficient}.

\subsection{Projective arithmetic density and a Hadamard-power no-go}
\label{r13:projective-density}

The density in \eqref{r13:eq-dense-group} has a simultaneous finite-matrix strengthening.
Given $n\ge1$, a positive-entry real matrix $B$, and $\epsilon>0$, choose each entry from
$G$ within a sufficiently small tolerance and clear all negative exponents by one common
$\lambda=2^RM^S$ with $R,S\ge0$.  Then
$\lambda,\lambda^{\thetaq}\in\Z_{>0}$, and this produces
\begin{equation}
 C\in M_n(\Z_{>0}),\qquad C^{\circ\thetaq}\in M_n(\Z_{>0}),
 \label{r13:eq-integral-hadamard-pair}
\end{equation}
such that
\begin{equation}
 \|\lambda^{-1}C-B\|_{\max}<\epsilon,
 \qquad
 \|\lambda^{-\thetaq}C^{\circ\thetaq}-B^{\circ\thetaq}\|_{\max}<\epsilon.
 \label{r13:eq-projective-density}
\end{equation}
The choices may be made symmetrically.  More precisely, if $\mathcal U$ and $\mathcal V$ are
separately open and invariant under positive scalar multiplication, then
$B\in\mathcal U$ and $B^{\circ\thetaq}\in\mathcal V$ imply that the integral pair may be
chosen with $C\in\mathcal U$ and $C^{\circ\thetaq}\in\mathcal V$.

This retires a natural positivity route.  For distinct positive $r_1,\ldots,r_4$ and small
$t>0$, Cauchy--Binet gives
\begin{equation}
 \det\bigl((1+tr_ir_j)^\alpha\bigr)_{i,j=1}^4
 =t^6\left(\prod_{k=0}^3\binom\alpha k\right)
   \prod_{i<j}(r_j-r_i)^2+O(t^7).
 \label{r13:eq-hadamard-sign}
\end{equation}
For $1<\alpha<2$ the leading coefficient is negative.  After a small positive diagonal
perturbation, the unpowered matrix is positive definite while its entrywise $\alpha$-power
is indefinite.  Applying \eqref{r13:eq-projective-density} with $\alpha=\thetaq$ yields a
positive-definite integer matrix $C$ for which $C^{\circ\thetaq}$ is both integral and
indefinite.  Thus critical-exponent failure is compatible with a counterexample; it cannot
serve as an integrality contradiction.

Supports placed on a rational near-kernel direction give no hidden collective gain either.
Their alternants factor into an ordinary Vandermonde times a Schur polynomial; for consecutive
support the determinant is exactly a product of pairwise integer binomials.  This is the
homogeneous counterpart of the affine factorization already \statuslean{} in
\lean{CandidateAlternant}, so it is cross-referenced rather than counted as a separate formal
advance.

\subsection{Determinantal divisors and the compound ceiling}
\label{r13:compounds}

Let $E\in M_N(\Z)$ be nonsingular and $1\le k\le N$.  Write $d_k(E)$ for the gcd of its
$k$-minors, let $\sigma_1(E)\ge\cdots\ge\sigma_N(E)>0$ be its singular values, adopt
$v_p(0)=+\infty$, and define
\begin{align}
 \tau_{p,k}(E)&=
 \min_{\substack{|I|=|J|=k\\\sigma:I\overset\sim\longrightarrow J}}
 \sum_{i\in I}v_p(E_{i,\sigma(i)}),
 &Q_k(E)&=\prod_p p^{\tau_{p,k}(E)}.
 \label{r13:eq-compound-content}
\end{align}
Only finitely many primes contribute to $Q_k(E)$.
The exact adelic compound sandwich is
\begin{equation}
 Q_k(E)\mid d_k(E),
 \qquad
 d_k(E)\le\|\textstyle\bigwedge^kE\|_2
 =\prod_{j=1}^k\sigma_j(E).
 \label{r13:eq-compound-sandwich}
\end{equation}
The divisibility uses only termwise prime divisibility of every minor; the inequality views
the minors as the entries of $\bigwedge^kE$.  It is the entire Smith/exterior-power analogue
of the inherited full-determinant comparison.

On the complete triangular evaluation grid $E_T$, put $Q_{T,k}=Q_k(E_T)$ and let
$E^-_{T,k}$ and $E^+_{T,k}$ denote,
respectively, the minimum reverse-pairing and maximum same-order $k$-point assignment
energies, and take $k/N_T\to\rho\in(0,1]$.  The limiting quantile is $\sqrt t$, so
\begin{equation}
 \frac{E^-_{T,k}}{N_TTL_T\log2}\longrightarrow
 \int_0^\rho\sqrt{t(\rho-t)}\,dt=\frac\pi8\rho^2,
 \qquad
 \frac{E^+_{T,k}}{N_TTL_T\log2}\longrightarrow
 \int_{1-\rho}^1t\,dt=\rho-\frac{\rho^2}{2}.
 \label{r13:eq-compound-energies}
\end{equation}
The combinatorial factor $\binom{N_T}{k}k!$ is lower order on the $T^4$ scale.  Under these
specific termwise-divisor and largest-minor/Frobenius bounds one obtains the ceiling
\begin{equation}
 \limsup_{T\to\infty}
 \frac{\log Q_{T,k}}
      {E^+_{T,k}+\log\!\left(\binom{N_T}{k}k!\right)}
 \le R(\rho)=\frac{\pi\rho}{8-4\rho},
 \qquad R(\rho)<R(1)=\frac\pi4\quad(0<\rho<1).
 \label{r13:eq-compound-ceiling}
\end{equation}
Thus every proper compound is worse than the full determinant under the inherited estimates.
This is not a no-go theorem for exterior powers themselves: a sharper singular-value free
energy, cancellation among minors, or extra Smith divisibility could still change the
comparison.

The report also applies a two-variable finite-difference transform to the square-grid
evaluation matrix.  It extracts diagonal weights $a!b!$.  If, for each prime $p$, the
multiset $\{v_p(a!b!):0\le a,b<K\}$ is arranged as
$\omega_{p,1}\le\cdots\le\omega_{p,K^2}$ and
\[
 \mathfrak F_{K,k}=\prod_p p^{\sum_{\ell=1}^k\omega_{p,\ell}},
\]
then the every-compound divisor is
\begin{equation}
 \mathfrak F_{K,k}\mid d_k(\mathcal E_K)\qquad(1\le k\le K^2).
 \label{r13:eq-factorial-compounds}
\end{equation}
At top degree this becomes
\[
 \left(\prod_{r=0}^{K-1}r!\right)^{2K}\mid\det\mathcal E_K,
 \qquad
 \log\left(\prod_{r=0}^{K-1}r!\right)^{2K}
 =K^3\log K-\frac32K^3+O(K^2\log K).
\]
The mechanism is the two-variable/all-compound extension of
\lean{SuperfactorialDeterminant}, \lean{AllLayerUnitBasis}, and
\lean{RowBlockUnitDivisibility}; its contribution
is lower order than the $K^4$ determinant budget and is retained as a paper specialization,
not advertised as a new leading divisor.

\subsection{Exact boundary displacement and transfer determinants}
\label{r13:boundary-displacement}

Let $K>0$, $\mathcal I_K=\{0,\ldots,K-1\}^2$, and set
\begin{equation}
 \mathcal E_K[(i,j),(u,v)]
 =(2^uM^v)^i(3^uA^v)^j.
 \label{r13:eq-grid-matrix}
\end{equation}
Put $d_{2;u,v}=2^uM^v$ and $d_{3;u,v}=3^uA^v$.  Let $D_2,D_3$ be the diagonal matrices
with these respective entries, let
$S_2,S_3$ be the two truncated forward row shifts, and let $P_2,P_3$ project onto their
terminal row boundaries.  Directly shifting a row index gives
\begin{align}
 S_2\mathcal E_K-\mathcal E_KD_2&=-P_2\mathcal E_KD_2,
 &S_3\mathcal E_K-\mathcal E_KD_3&=-P_3\mathcal E_KD_3.
 \label{r13:eq-displacement}
\end{align}
After sorting, irrationality of $\thetaq$ and $\beta$, respectively, makes the row parameters
and column frequencies pairwise distinct; strict total positivity of $e^{xy}$ therefore gives
$\det\mathcal E_K\ne0$.  Since the coordinate multipliers are nonzero, each right side in
\eqref{r13:eq-displacement} has rank exactly $K$.  The commuting
multiplication operators
\[
 T_r=\mathcal E_KD_r\mathcal E_K^{-1}\qquad(r=2,3)
\]
have the boundary-feedback form $T_r=S_r+R_r$ with $\operatorname{rank}R_r=K$.  Their commutativity gives
the exact boundary relation
\begin{equation}
 [S_2,R_3]+[R_2,S_3]+[R_2,R_3]=0.
 \label{r13:eq-boundary-commutator}
\end{equation}

If $J_r:\R^K\to\R^{K^2}$ injects the corresponding boundary and $R_r=J_rC_r$, the matrix
determinant lemma compresses the full spectrum, for $z\ne0$, to
\begin{equation}
 \prod_{0\le u,v<K}(z-d_{r;u,v})
 =z^{K^2}\det\!\left(I_K-C_r(zI-S_r)^{-1}J_r\right),
 \qquad
 (zI-S_r)^{-1}=\sum_{h=0}^{K-1}z^{-h-1}S_r^h.
 \label{r13:eq-transfer-determinant}
\end{equation}
This is an exact $K\times K$ encoding of a $K^2$-point spectrum.  It does not by itself
provide an arithmetic gain: the primitive denominators and local Smith factors of the
boundary feedback remain unknown.

Because the $d_{2;u,v}$ are distinct, there is also a unique
$\widehat Q_K\in\Q[X]$ of degree below $K^2$ with
$\widehat Q_K(d_{2;u,v})=d_{3;u,v}$.  Simultaneous diagonalization gives the exact paper
identity
\begin{equation}
 T_3=\widehat Q_K(T_2).
 \label{r13:eq-spectral-interpolation}
\end{equation}
Together with \eqref{r13:eq-boundary-commutator}, this concentrates the unresolved arithmetic
in two rank-$K$ boundary feedbacks.

\subsection{Kraus--Loewner positivity and its denominator tax}
\label{r13:kraus}

For $f(t)=t^{\thetaq}$ and $a,x,y>0$, let $K_a(x,y)=f[x,a,y]$ be the second divided
difference.  The operator-convex integral representation gives
\begin{equation}
 K_a(x,y)=c_{\thetaq}\int_0^\infty
 \frac{s^{\thetaq}}{(x+s)(a+s)(y+s)}\,ds,
 \qquad c_{\thetaq}=\frac{\sin(\pi(\thetaq-1))}{\pi}>0.
 \label{r13:eq-kraus-kernel}
\end{equation}
Let $0<x_1<\cdots<x_n$ and $0<y_1<\cdots<y_n$.  Andr\'eief's identity and the two Cauchy determinants
give a strictly positive multiple integral for $\det(K_a(x_i,y_j))$.  Thus the cross-kernel
is strictly totally positive.

When $a,x_i,y_j$ are pairwise distinct positive integers and all their $\thetaq$-th powers
are integers, each entry is rational and
\begin{equation}
 \left(\prod_i(x_i-a)\right)
 \left(\prod_j(a-y_j)\right)
 \left(\prod_{i,j}(x_i-y_j)\right)
 \det(K_a(x_i,y_j))\in\Z\setminus\{0\}
 \label{r13:eq-kraus-clearing}
\end{equation}
But scaling all nodes by $R>0$ multiplies this cleared expression by
$R^{n^2+n\thetaq}$, a positive power.  (It remains an integer after scaling whenever the
scaled nodes satisfy the same integrality hypotheses.)  The derivative decay is therefore overwhelmed by the
naive denominator; a useful argument would need primitive denominator cancellation or new
numerator divisibility.

\begin{statusbox}{Report 13 was audited independently from the common base, not as a revision
of report 10.  Its projective density and Hadamard-power no-go, compound sandwich and
triangular ceiling, every-compound finite-difference specialization, transfer determinant,
rational-conditioning factorization, and Kraus--Loewner package have self-contained paper
proofs.  The exact shift/displacement identities are \statuslean{} in
\lean{BoundaryDisplacement}; its rank-$K$ endpoints explicitly assume $K>0$, nonsingularity,
and nonzero coordinate multipliers, all supplied for the displayed grid by the paper argument
above.  The
transfer determinant and spectral interpolation remain the paper layer.  The rational-direction
alternant and factorial-basis mechanisms overlap existing
accepted modules and are not double-counted.  The compound ceiling applies only to the stated
termwise-divisor/Frobenius estimates; the boundary compression and Kraus positivity are exact
interfaces, not claimed contradictions.}
\end{statusbox}

\section{Fourteenth report: structural-prime residuals and dilation resultants}
\label{sec:report14}

Report 14 is another independent continuation from the common unified base.  Its first
finite-place observations overlap report 11: the real logarithm matrix has rank one, whereas
the Iwasawa logarithm matrices at the structural primes $2$ and $3$ have rank two; the first
good-prime shadow is the Fermat-quotient determinant already isolated in
\lean{FermatQuotientDeterminant}.  The genuinely new part begins when the square-grid
interpolant is placed in its integral coefficient lattice and its two boundary residuals are
normalized at a structural prime.

\subsection{Local rank is thick, but it is not a product-formula input}
\label{r14:local-rank}

Under a hypothetical nonintegral solution write
\[
 M=2^\beta\in\N,\qquad A=3^\beta\in\N,\qquad \beta\notin\Q.
\]
For Iwasawa logarithms the positive-rational kernel is
\[
 \log_p r=0\quad\Longleftrightarrow\quad r\in p^\Z
 \qquad(r\in\Q_{>0}).
\]
Consequently
\[
 \det\begin{pmatrix}\log_2 2&\log_2M\\ \log_2 3&\log_2A\end{pmatrix}
   =-\log_2M\log_23\ne0,
\quad
 \det\begin{pmatrix}\log_3 2&\log_3M\\ \log_3 3&\log_3A\end{pmatrix}
   =\log_32\log_3A\ne0.
\]
After principal-unit normalization, the two-parameter orbit therefore has open closure in a
two-dimensional $p$-adic torus at both structural primes.  For a full-rank logarithm matrix
$B_p$ there is a constant $c_p$ such that, for $z\ne0$,
\[
 \bigl|\min_i v_p((B_pz)_i)-\min_i v_p(z_i)\bigr|\le c_p.
\]
Thus for $K\ge2$, one fixed pair of distinct parameters in $\{0,\ldots,K-1\}^2$ cannot have
depths $r_p$ at a finite set $S$ of full-rank primes unless
\begin{equation}
 \sum_{p\in S}r_p\log p
 \le \log(K-1)+\sum_{p\in S}c_p\log p.
 \label{r14:eq-collision-budget}
\end{equation}
This is a genuine local packing ceiling, not a contradiction.  The determinants of local
logarithms are different transcendental periods at different places; they are not absolute
values of one algebraic number and cannot be inserted directly into a product formula.

At a good prime $\ell\nmid6MA$, after clearing denominators an integer relation
$a\log_\ell2+b\log_\ell M=0$ first implies
$2^aM^b\in\ell^\Z$.  Its $\ell$-valuation is zero, hence it equals one; multiplicative
independence then gives $a=b=0$.  This kernel step is essential before applying
Baker--Brumer to a putative local shadow exponent.  It is explicitly included here because
omitting it would turn multiplicative independence into an invalid assertion about
$p$-adic logarithms.

\subsection{The interpolation denominator is a cokernel order}
\label{r14:denominator-lattice}

Let $K\ge2$, $R=K^2$, $n=K-1$, and enumerate the nodes and values
\[
 x_{ij}=2^i3^j,\qquad y_{ij}=M^iA^j\qquad(0\le i,j<K).
\]
Let $E_K=(x_{ij}^{\,r})_{(i,j),\,0\le r<R}$ and let $Q_K\in\Q[X]$ be the unique
interpolant with $Q_K(x_{ij})=y_{ij}$.  Define
\[
 \delta_K=\min\{d>0:dQ_K\in\Z[X]\},\qquad P_K=\delta_KQ_K,
\]
and put
\[
 d_{ij}^{(K)}=\prod_{(a,b)\ne(i,j)}(x_{ij}-x_{ab}),
 \qquad D_K^*=\operatorname{lcm}_{i,j}|d_{ij}^{(K)}|.
\]
The Lagrange basis and Smith normal form give three exact descriptions:
\begin{equation}
 D_K^*=\min\{D>0:DE_K^{-1}\in M_R(\Z)\}=s_R,
 \qquad
 \delta_K=\operatorname{ord}_{\Z^R/E_K\Z^R}[y],
 \label{r14:eq-denominator-cokernel}
\end{equation}
where $s_R$ is the largest invariant factor of $E_K$.  If
$UE_KV=\operatorname{diag}(s_1,\ldots,s_R)$, then more explicitly
\[
 \delta_K=\operatorname{lcm}_{1\le r\le R}
 \frac{s_r}{\gcd(s_r,(Uy)_r)}.
\]
Thus all node-only denominator growth is separated from the order of the specialized value
class.  Any unexpectedly small denominator must be an arithmetic statement about that class,
not merely a small Vandermonde determinant.

The exact endpoint valuation calculation finds the maximal dyadic barycentric denominator
only in the top row and a maximizer at $(n,0)$; triadically it lies only in the rightmost
column with a maximizer at $(0,n)$.  This yields two favorable, nonexhaustive branches:
\begin{align}
 M\text{ odd},\ A\text{ even}:&\quad
   v_2([X^{R-1}]P_K)=0,\\
 3\mid M,\ 3\nmid A:&\quad
   v_3([X^{R-1}]P_K)=0.
 \label{r14:eq-favorable-branches}
\end{align}
The restrictions in \eqref{r14:eq-favorable-branches} are part of the theorem.  They do not
cover every hypothetical counterexample.

\subsection{Exact scaling, roots, and an honest algebraic resultant}
\label{r14:scaled-resultants}

Let $F_K,G_K\in\Z[X]$ be the two degree-$n$ residuals obtained from the dyadic and triadic
boundary-dilation factorizations after multiplication by $\delta_K$.  In the first branch of
\eqref{r14:eq-favorable-branches}, report 14 proves
\begin{equation}
 \mathcal F_K(Z)=2^{-n^2}F_K(2^nZ)\in\Z_2[Z],
 \qquad
 \overline{\mathcal F_K}(Z)\doteq1-(Z-1)^n,
 \label{r14:eq-dyadic-scaling}
\end{equation}
with unit leading coefficient and
$v_2(\mathcal F_K(0))=v_2(M-1)$.  In the second branch,
\begin{equation}
 \mathcal G_K(Z)=3^{-n^2}G_K(3^nZ)\in\Z_3[Z],
 \qquad
 \overline{\mathcal G_K}(Z)\doteq R_n(0)-\overline A R_n(Z),
 \quad R_n(Z)=\prod_{i=1}^n(Z-2^i).
 \label{r14:eq-triadic-scaling}
\end{equation}
Consequently
\[
 v_2([X^r]F_K)\ge n(n-r),\qquad
 v_3([X^r]G_K)\ge n(n-r),
\]
and all roots have structural-prime depth at least $n$.  The total excess depth is exactly
$v_2(M-1)$ or $v_3(A-1)$, respectively.  When $n$ is odd the reductions are sufficiently
separable to determine the exceptional root classes.

The categorical payoff is a single algebraic integer rather than unrelated local periods:
\begin{align}
 \mathcal R_{2,K}&=\operatorname{Res}(F_K(X),F_K(3X)),\\
 \mathcal R_{3,K}&=\operatorname{Res}(G_K(X),G_K(2X)).
\end{align}
If the relevant resultant is nonzero, then
\begin{align}
 v_2(\mathcal R_{2,K})&\ge n^3+n+v_2(M-1),\\
 v_3(\mathcal R_{3,K})&\ge n^3+v_3(A-1).
 \label{r14:eq-resultant-depths}
\end{align}
For odd $n$ it is nonzero and equality holds.  The finite algebraic hinge is the identity
\begin{equation}
 P(3X)-P(X)=2X\,\mathcal D_3(P),\qquad
 \overline{\mathcal D_3(P)}=\overline P'\quad\text{in }\mathbb F_2[X],
 \label{r14:eq-dilation-quotient}
\end{equation}
where $\mathcal D_3(P)\in\Z[X]$ is an explicit geometric-sum quotient.  Equation
\eqref{r14:eq-dilation-quotient} is \statuslean{} in
\lean{DilationDifferenceQuotient}.

The resultants are the right kind of adelic objects, but their known archimedean heights are
far larger than the local divisor in \eqref{r14:eq-resultant-depths}.  No product-formula
contradiction follows.  The next missing theorem would have to control primitive resultant
height or force compatibility between the two favorable branches; neither is supplied by
local rank or by the Smith denominator alone.

\begin{statusbox}{Report 14 was audited independently from the common base.  Structural-prime
rank and Fermat-quotient statements overlapping report 11 are cross-referenced rather than
double-counted.  The integral dilation quotient and its mod-$2$ derivative reduction are
\statuslean{} in \lean{DilationDifferenceQuotient}.  The denominator-lattice, endpoint,
scaled-root, and resultant theorems are self-contained paper arguments, explicitly restricted
to $M$ odd and $A$ even or to $3\mid M$ and $3\nmid A$ where stated.  The report's larger
symbolic audit is retained only as an unbundled diagnostic: the displayed source does not
reconstruct every claimed interpolation, Newton polygon, and resultant.  No conclusion here
proves the conjecture.}
\end{statusbox}

\section{Fifteenth report: synchronized information and decoder complexity}
\label{sec:report15}

Report 15 is independent of report 14 and of every earlier continuation.  It treats a
hypothetical semigroup as a finite information source.  Several of its dynamical statements
overlap the Sturmian probability and natural-boundary packages of reports 7 and 9; those are
cross-referenced.  Its genuinely new contribution is a sharp fixed-rank sum-channel
calculation, an exact two-view decoder, quantitative degree obstructions for algebraic
decoders, and a finite Fourier formulation of tied determinant cancellation.

\subsection{One magnitude channel and a rank-two valuation profile}
\label{r15:magnitude-channel}

For $0\le n,k<L$ set
\[
 X_{n,k}=2^nM^k=2^{n+\beta k},\qquad
 Y_{n,k}=3^nA^k=3^{n+\beta k}.
\]
The binary digit length of $X_{n,k}$ and ternary digit length of $Y_{n,k}$ are the same
integer, $n+\lfloor\beta k\rfloor+1$.  Thus the two apparent size observations are one channel,
not two independent constraints.  Its entropy is $\log L+O_\beta(1)$, while the source has
entropy $2\log L$.  Prefix-code redundancies are bounded quasiperiodic functions of the same
rotation phase, so their rate vanishes.

The synchronized carry word has zero entropy rate and logarithmic block entropy; its exact
prefix mean, variance, and two-shift covariance are already \statuslean{} in
\lean{SturmianProbability}, while the carry/floor generating-series identities are
\statuslean{} in \lean{SturmianGeneratingSeries}.  P\'olya--Carlson supplies a natural
boundary only after the classical radius-one and non-eventual-periodicity hypotheses.  That
one-variable obstruction does not create a two-prime functional equation.

Fixed integer valuation observations have entropy growth equal to the rank of the displayed
linear forms times $\log L+O(1)$.  For the two output boxes the normalized profile is a
rank-two representable polymatroid.  Consequently every universal homogeneous linear rank or
Shannon inequality is automatically satisfied at leading order.  Any contradiction must see
heights, growing moduli, or a subleading arithmetic term.

\subsection{Fixed-rank sums are asymptotically unordered-pair codes}
\label{r15:sunit-entropy}

Let $E$ be an $N$-element subset of a fixed finitely generated multiplicative group of
positive algebraic numbers, of rank $r$.
The Evertse--Schlickewei--Schmidt unit-equation theorem
\cite{EvertseSchlickeweiSchmidt2002} implies that the nondegenerate
solutions of $x_1+x_2=x_3+x_4$ contribute only $O_r(N)$ beyond the two trivial permutations.
Hence
\begin{align}
 \mathcal E_+(E)&=2N^2-N+O_r(N),\\
 |E+E|&=\frac{N^2}{2}+O_r(N),\\
 H_2(Z+Z')&=2\log N-\log2+O_r(N^{-1}),\\
 H(Z+Z')&=2\log N-\log2+O_r((\log N)/N),
 \label{r15:eq-sunit-entropy}
\end{align}
for independent uniform $Z,Z'\in E$.  Thus the sum determines the unordered input pair with
error $O_r(N^{-1})$.  This is a strong exact-shape theorem, but it is compatible with the
hypothetical rank-two boxes: their additive channel is already almost maximally informative.

\subsection{Exact synchronized decoding and degree growth}
\label{r15:decoders}

For any real $p>1$, strict convexity gives a finite injectivity statement with no number
theory in its proof:
\begin{equation}
 a+b=c+d,\quad a^p+b^p=c^p+d^p,\quad a,b,c,d>0
 \quad\Longrightarrow\quad
 \{a,b\}=\{c,d\}.
 \label{r15:eq-exact-decoder}
\end{equation}
Indeed, with $s=a+b$ the function $x\mapsto x^p+(s-x)^p$ is strictly decreasing on
$[0,s/2]$.  This complete unordered-pair decoder is \statuslean{} in
\lean{SynchronizedSumDecoder}; the formal theorem is
\lean{eq_or_swap_of_sum_eq_rpow_sum_eq}.

Exact decodability does not mean bounded algebraic complexity.  Suppose $L\ge2$ and
$R=P/Q\in\R(z)$ has $\deg P,\deg Q\le d$, no pole at $2^n+1$, and
\begin{equation}
 R(2^n+1)=3^n+1\qquad(0\le n<L).
 \label{r15:eq-rational-decoder}
\end{equation}
After cross multiplication and translation by one, the left side is a nonzero exponential
polynomial whose bases lie in
\[
 \{2^j:0\le j\le d\}\cup\{3\cdot2^j:0\le j\le d\}.
\]
These $2d+2$ bases are distinct.  The Chebyshev zero bound therefore gives
\begin{equation}
 L\le2d+1,\qquad d\ge\left\lceil\frac{L-1}{2}\right\rceil.
 \label{r15:eq-rational-decoder-degree}
\end{equation}
The cross-multiplied inequality $L\le2d+1$ is \statuslean{} as
\lean{rational_sum_decoder_card_le} in \lean{RationalSumDecoder}; the formal hypotheses are
$\deg P,\deg Q\le d$, $Q\ne0$, and the cross-multiplied identities
$P(2^n+1)=(3^n+1)Q(2^n+1)$ for all $0\le n<L$.  The rational-function formulation follows
from its no-pole premise.
More generally, a nonzero $F\in\R[z,w]$ of total degree at most $d$ satisfying
$F(2^n+1,3^n+1)=0$ for $0\le n<L$ gives
\begin{equation}
 L\le\binom{d+2}{2}-1=\frac{d(d+3)}2.
 \label{r15:eq-algebraic-decoder-degree}
\end{equation}
The scope is the displayed exponential slice.  A completion criterion may use these bounds
only when its synchronized family contains $(2^n,1)$ for arbitrarily long initial ranges;
they do not prove degree growth on every unbounded collection of semigroup pairs.

\subsection{Fourier separation identifies, but does not solve, tied cancellation}
\label{r15:fourier-separation}

Let $B\in\{\pm1\}$ be a balanced parity bit and $Z$ a random element of a finite abelian
group $G$.  The signed measure
\[
 \nu(g)=\mathbb E[B\mathbf1_{Z=g}]
\]
has Fourier coefficients $\widehat\nu(\chi)=\mathbb E[B\overline{\chi(Z)}]$.  Finite Fourier
inversion and Parseval give
\begin{align}
 B\text{ independent of }Z
 &\Longleftrightarrow \widehat\nu(\chi)=0\quad(\chi\in\widehat G),\\
 \operatorname{TV}(\mathcal L(B,Z),\mathcal L(B)\otimes\mathcal L(Z))
 &\le\frac12\left(\sum_{\chi\in\widehat G}|\widehat\nu(\chi)|^2\right)^{1/2}.
 \label{r15:eq-fourier-separation}
\end{align}
If $B$ is determined by $Z$, at least one character correlation has magnitude
$|G|^{-1/2}$.  Applied to tied determinant terms, this says exactly what would force
noncancellation: a legal twist family spanning all residue characters while preserving the
archimedean estimate and the tied fiber.  No current determinant supplies that family, so
\eqref{r15:eq-fourier-separation} is a precise target rather than a completed argument.

\begin{statusbox}{Report 15 was audited independently from the common base.  Its exact
ordinary/power-sum decoder is \statuslean{} in \lean{SynchronizedSumDecoder}.  The rational
decoder degree bound is \statuslean{} in \lean{RationalSumDecoder}.  The algebraic decoder
degree bound and finite Fourier lemma are unconditional paper proofs;
the fixed-rank entropy estimates use the classical unit-equation theorem.  Sturmian and
P\'olya--Carlson material overlapping reports 7 and 9 is not double-counted.  The report's
uniform-decoder criterion is stated only for families containing the exponential slice
$(2^n,1)$ through arbitrarily large initial ranges.  Character-twisted determinant
separation remains an open construction.  None of the entropy identities by itself proves
the conjecture.}
\end{statusbox}

\section{Sixteenth report: transfer spectra and mixed Mahler resonance}
\label{sec:report16}

Report 16 is another independent sibling of reports 1--15: it starts from the same unified
base and asks whether classical special functions can amplify the multiplicative rank of the
two known outputs.  Its generic spectrum theorem strengthens the existing $(2,3)$
classification, while its gamma, Hurwitz, polylogarithmic, Barnes, double-sine, and
$q$-digamma packages diagnose several exact but rank-preserving reformulations.  The report's
recorded metadata searches all failed, so no ``no proof found'' conclusion is imported here;
only mathematical claims with an internal proof or an identified classical input are used.

\subsection{The algebraic-base spectrum}
\label{r16:spectrum}

Let $a,b,c>0$ be algebraic, assume that the nonnegative monomials in $a,b$ are unique, and
let $x\notin\Q$ with $a^x,b^x$ algebraic.  Write
\begin{equation}
 \operatorname{Rel}_{a,b}(c)\quad\Longleftrightarrow\quad
 \exists n>0\ \exists i,j,k,l\in\Nzero:\
 a^k b^l c^n=a^i b^j.
 \label{r16:eq-pair-relation}
\end{equation}
Then
\begin{equation}
 \boxed{c^x\in\Qbar\quad\Longleftrightarrow\quad
        \operatorname{Rel}_{a,b}(c).}
 \label{r16:eq-spectrum}
\end{equation}
Indeed, the forward implication follows from the algebraic Six Exponentials theorem unless
the three-base nonnegative monomial map has a collision; the reverse implication follows by
raising \eqref{r16:eq-pair-relation} to $x$ and extracting a positive real $n$th root.  This is
\statuslean{} as
\lean{AlgebraicBaseSpectrum.real_rpow_isAlgebraic_iff_hasPositivePairPowerRelation}; the exact
collision equivalence is \lean{AlgebraicBaseSpectrum.tripleMonomial_injective_iff}.

For a hypothetical nonintegral solution and $(a,b)=(2,3)$, this recovers the exact statement
that the positive algebraic bases with algebraic $x$th power are precisely the divisible hull
of $\langle2,3\rangle$.  Thus no positive algebraic third base can simultaneously have an
algebraic output and supply an independent logarithmic direction.  Gamma quotients, radicals,
norms, and structural-prime constructions do not evade this theorem once their positivity and
algebraicity are established.

\subsection{Entire cancellation and rank-preserving transfer formulas}
\label{r16:entire-transfer}

For an integer $b\ge2$, the apparent meromorphic quotient
\begin{equation}
 Q_b(z):=\frac{\Gamma(-z)}{\Gamma(1-b^z)}
 \label{r16:eq-entire-quotient}
\end{equation}
extends to an entire function: at every $k\in\Nzero$ the two simple poles cancel with a
nonzero ratio.  Its other zeros are the nonintegral points for which $b^z$ is a positive
integer.  Since the two vertical phase lattices for bases $2$ and $3$ meet only on the real
axis,
\begin{equation}
 \text{Alaoglu--Erd\H{o}s}\quad\Longleftrightarrow\quad
 Z(Q_2)\cap Z(Q_3)=\varnothing.
 \label{r16:eq-entire-coprime}
\end{equation}
At a nonintegral real zero with $b^x=B\in\N$ one has
\begin{equation}
 Q_b'(x)=\Gamma(-x)(-1)^B B!\log b.
 \label{r16:eq-entire-derivative}
\end{equation}
These paper identities remove all known integral solutions exactly, but they do not create a
two-dimensional common zero lattice; a common Barnes divisor would therefore be spurious.

Gauss multiplication, Hurwitz distribution, and the Cartier identity for
$F_x(z)=\operatorname{Li}_{-x}(z)$ all yield the same multiplier $q^x$.  Under a
counterexample their normalized values are algebraic exactly for $3$-smooth $q$, by
\eqref{r16:eq-spectrum}.  The dichotomy
\begin{equation}
 \operatorname{Li}_{-x}(z)\text{ is D-finite over }\C(z)
 \quad\Longleftrightarrow\quad x\in\Nzero
 \label{r16:eq-holonomic}
\end{equation}
is classical: nonintegral order produces infinitely many independent fractional-power
monodromy germs.  Likewise, after defining $a_0=0$, the sequence $a_n=n^x$ is $q$-regular
exactly for integral $x$.  Integer Cartier eigenvalues at $2$ and $3$, however, do not imply
either finite holonomic rank or finite $q$-kernel rank; that missing implication is another
form of the original wall.

\subsection{A finite Mahler identity and the corrected resonance theorem}
\label{r16:mahler}

For a commutative ring and finite cutoffs define
\begin{equation}
 H_{A,B}(z):=\sum_{0\le a<A}\sum_{0\le b<B}
 M^aN^b z^{2^a3^b}.
 \label{r16:eq-mahler-trunc}
\end{equation}
The exact rectangular inclusion--exclusion identity
\begin{align}
 H_{A+1,B+1}(z)
 &=z+M H_{A,B+1}(z^2)+N H_{A+1,B}(z^3)\\
 &\qquad-MN H_{A,B}(z^6)
 \label{r16:eq-mahler-finite}
\end{align}
is \statuslean{} as \lean{mixedMahlerTrunc_succ_succ} in
\lean{MixedMahlerFinite}.  It is a finite algebra theorem only; convergence, natural
boundaries, Mellin transforms, and asymptotics are not hidden in its statement.

For positive integers $M,N$, the limiting integer-coefficient series
\begin{equation}
 H_{M,N}(z)=\sum_{a,b\ge0}M^aN^b z^{2^a3^b}
 \label{r16:eq-mahler-series}
\end{equation}
has radius one and a non-eventually-periodic zero pattern, so the classical
P\'olya--Carlson/Skolem--Mahler--Lech argument
\cite{Carlson1921,Lech1953} makes the unit circle a natural boundary.  This
happens for both integral and hypothetical nonintegral exponents and is therefore only a
diagnostic.

For $M,N\ge2$, put
\begin{equation}
 \alpha=\frac{\log M}{\log2},\qquad
 \beta=\frac{\log N}{\log3}.
\end{equation}
Termwise Mellin integration initially gives
\begin{equation}
 \int_0^\infty H_{M,N}(e^{-t})t^{s-1}\,dt
 =\frac{\Gamma(s)}{(1-M2^{-s})(1-N3^{-s})}
 \qquad(\Re s>\max\{\alpha,\beta\}).
 \label{r16:eq-mellin}
\end{equation}
The two geometric pole lattices meet exactly when $\alpha=\beta$, and then only at the real
point $s=x$.  Consequently resonance gives
\begin{equation}
 H_{M,N}(e^{-t})
 =t^{-x}\left[
 \frac{\Gamma(x)}{\log2\log3}\log\frac1t+B+P_2(\log(1/t))+P_3(\log(1/t))
 \right]+O(t^{-x+\eta}),
 \label{r16:eq-resonance}
\end{equation}
where $P_2,P_3$ are absolutely convergent analytic periodic functions with periods
$\log2,\log3$.  The restriction $M,N\ge2$ is essential: $M=1$ or $N=1$ introduces an
unrelated Gamma-pole collision at $s=0$.  The Barnes kernel
\begin{equation}
 \frac1{(1-e^{-v\log2})(1-e^{-v\log3})}
 =\sum_{a,b\ge0}e^{-v(a\log2+b\log3)}
 \label{r16:eq-barnes-kernel}
\end{equation}
is valid as a geometric series for $\Re v>0$; the corresponding Barnes double-zeta integral
starts in $\Re w>2$, $\Re a>0$ and is then continued.  The double pole detects equality of
the two exponents, not their integrality.

\subsection{Double-sine, $q$-digamma, and the surviving target}
\label{r16:qgamma}

The reciprocal double-sine normalization with periods $\log2,\log3$ packages the two outputs
as adjacent shift quotients at $z_x=i x\log2\log3/(2\pi)$.  The shift compatibility is built
into the function and supplies no contradiction.  Separately, Borwein's Pad\'e theorem gives
the unconditional irrationality of
\begin{equation}
 L_b(B)=\sum_{k\ge0}\frac1{Bb^k-1}
 =-\frac{\psi_{1/b}(x)+\log(1-b^{-1})}{\log b}
 \qquad(B=b^x\in\N,\ b\ge2).
 \label{r16:eq-qdigamma}
\end{equation}
This one-base irrationality \cite{Borwein1992} also holds for every genuine integral
exponent.  Finally,
${}_1F_0(-x;;-z)=(1+z)^x$ (continued analytically at the displayed arguments) makes the
rank-amplification target explicit: the known values at $1+z=2,3$ do not algebraize the value
at $1+z=5$.  A cross-base value theorem doing so would contradict
\eqref{r16:eq-spectrum}, but no such theorem is presently available.

\begin{statusbox}{Report 16 was audited independently from the common base.  The generic
algebraic-base spectrum is \statuslean{} in \lean{AlgebraicBaseSpectrum}; the finite mixed
Mahler identity is \statuslean{} in \lean{MixedMahlerFinite}.  The entire quotient, corrected
$M,N\ge2$ Mellin resonance, and transfer identities are self-contained paper arguments.  The
holonomic dichotomy, P\'olya--Carlson boundary, Barnes continuation, and Borwein irrationality
use named classical inputs.  Double-sine rigidity, Cartier-to-holonomy, and all proposed
rank-amplification statements remain open.  Nothing in the section proves the conjecture.}
\end{statusbox}

\section{Seventeenth report: weighted Smith profiles and rational Loewner clusters}
\label{sec:report17}

Report 17 is independent of reports 1--16.  It replaces the single determinant by the full
sequence of determinantal divisors and compares that arithmetic profile with every exterior
singular-value level.  It also develops a rational cross-Loewner ensemble and proves that
ordinary probability and entropy are already saturated by the hypothetical semigroup.  The
new certificates are exact sufficient criteria, not verified contradictions.  The basic
exterior/Frobenius identity overlaps report 13, the fractional Loewner integral and rational
positivity overlap report 10, and source-recovery entropy overlaps report 15.  The genuinely
new refinements are the weighted all-minor inequality, the $p^t$-level modular divisor, the
full geometric-Vandermonde Smith form, and the level-by-level Loewner singular hierarchy.

\subsection{From one determinant to every exterior level}
\label{r17:smith-profile}

For a full-rank $A\in\Z^{n\times n}$ let $\delta_k(A)$ be the gcd of its $k$-minors and let
$\sigma_1\ge\cdots\ge\sigma_n$ be its singular values.  Cauchy--Binet in exterior form gives
\begin{equation}
 \sum_{|I|=|J|=k}\det(A_{I,J})^2
 =e_k(\sigma_1^2,\ldots,\sigma_n^2).
 \label{r17:eq-exterior-minors}
\end{equation}
If all $k$-minors are nonzero, then
\begin{equation}
 e_k(\sigma(A)^2)\ge\binom nk^2\delta_k(A)^2.
 \label{r17:eq-smith-mass}
\end{equation}
Thus, if $A=B+R$, $\operatorname{rank}B\le r<k$, $\|R\|_2\le\varepsilon$, and
$\|A\|_2\le S$, then
\begin{equation}
 \boxed{\sqrt{\binom nk}\,\delta_k(A)\le S^r\varepsilon^{k-r}.}
 \label{r17:eq-bulk-smith}
\end{equation}
Every intermediate $k$ is a possible contradiction level.

For a diagonal gauge $K=D^{-1}AE^{-1}$ with positive diagonal entries $\rho_i,\kappa_j$,
the exact weighted identity is
\begin{equation}
 e_k(\sigma(K)^2)
 =\sum_{I,J}\frac{\det(A_{I,J})^2}{\rho_I^2\kappa_J^2}
 \ge\delta_k(A)^2 e_k(\rho^{-2})e_k(\kappa^{-2}).
 \label{r17:eq-weighted-smith}
\end{equation}
This retains Smith mass after removing analytically irrelevant row and column scales, unlike
a determinant-only normalization.

\subsection{Multilevel residue collapse and a geometric Smith model}
\label{r17:modular-smith}

For a $d$-variable integer monomial-evaluation matrix, let
$0\le\nu_{p,1}\le\cdots\le\nu_{p,n}$ be the valuations of its Smith factors and put
$r_{p,t}=\#\{i:\nu_{p,i}<t\}$.  The exact Ferrers-level identity and residue-class bound are
\begin{equation}
 v_p(\delta_k)=\sum_{t\ge1}(k-r_{p,t})_+,
 \qquad r_{p,t}\le p^{td}.
 \label{r17:eq-level-decomposition}
\end{equation}
Consequently every such matrix satisfies
\begin{equation}
 \boxed{D_{k,d}:=\prod_p p^{\sum_{t\ge1}(k-p^{td})_+}\mid\delta_k.}
 \label{r17:eq-universal-divisor}
\end{equation}
For fixed $d$,
\begin{equation}
 \log D_{k,d}=\left(\frac d{d+1}+o(1)\right)k^{1+1/d};
 \label{r17:eq-divisor-asymptotic}
\end{equation}
the bivariate constant is $2/3$.  This paper theorem applies to arbitrary supports, but its
universal size is lower order in the natural global determinant budgets.

The one-variable geometric Vandermonde
\begin{equation}
 V_n(q)=(q^{ij})_{0\le i,j<n},\qquad q\ge2,
\end{equation}
has the exact Smith factors
\begin{equation}
 s_i(q)=q^{\binom i2}\prod_{r=1}^i(q^r-1),\qquad 0\le i<n.
 \label{r17:eq-q-smith}
\end{equation}
The proof changes from monomials to the Newton basis
$\prod_{r<i}(X-q^r)$, leaving a diagonal matrix times an upper unitriangular Gaussian-binomial
matrix.  Hence
\begin{equation}
 \log\delta_k(V_n(q))=\frac{\log q}{3}k^3+O_q(k^2).
 \label{r17:eq-q-smith-asymptotic}
\end{equation}
This large one-base Smith mass occurs for every ordinary integer $q$ and is therefore a
calibration and no-go, not evidence of exceptional two-base structure.  The complete
determinant-one two-sided reduction, the integral $q$-Pascal identity, the divisibility chain,
and positivity for $q\ge2$ are \statuslean{} in \lean{GeometricVandermondeSmith}; the
formal endpoint is
\lean{geometricVandermonde_isUnimodularEquivalent_qSmithDiagonal}.

\subsection{Rational cross-Loewner clusters}
\label{r17:loewner}

Write a hypothetical exponent as $x=m+\alpha$ with $0<\alpha<1$.  For
$u,v$ in the dense group $\{2^a3^b:a,b\in\Z\}$, both $u^\alpha,v^\alpha$ are rational, so the
cross-Loewner kernel
\begin{equation}
 L_\alpha(u,v)=\frac{u^\alpha-v^\alpha}{u-v}
 =\frac{\sin(\pi\alpha)}{\pi}
   \int_0^\infty\frac{t^\alpha}{(u+t)(v+t)}\,dt
 \label{r17:eq-loewner}
\end{equation}
is rational at distinct group nodes.  Andreief's identity and the Cauchy determinant give,
for disjoint increasing node sets,
\begin{align}
 \det(L_\alpha(u_i,v_j))
 &=\frac{c_\alpha^n}{n!}\,\Delta(u)\Delta(v)\\
 &\quad\times\int_{(0,\infty)^n}
 \frac{\prod t_\ell^\alpha\Delta(t)^2}
      {\prod_{i,\ell}(u_i+t_\ell)(v_i+t_\ell)}\,dt,
 \label{r17:eq-loewner-andreief}
\end{align}
which is strictly positive.  For affine clusters
$u_i=1+\varepsilon\xi_i$, $v_j=1+\varepsilon\eta_j$, Selberg's integral yields
\begin{equation}
 \det(L_\alpha(u_i,v_j))
 =\kappa_{n,\alpha}\Delta(\xi)\Delta(\eta)
   \varepsilon^{n(n-1)}(1+O(\varepsilon)),
 \label{r17:eq-loewner-cluster}
\end{equation}
with an explicit positive gamma-product constant.  Orthogonal projection in the Stieltjes
feature space gives the sharper individual hierarchy
\begin{equation}
 \sigma_j=\Theta(\varepsilon^{2(j-1)})\qquad(1\le j\le n)
 \label{r17:eq-loewner-singular}
\end{equation}
for fixed nondegenerate affine clusters; the determinant constant is the corresponding
Selberg integral \cite{Selberg1944}.

If $A=DKE$ clears a rational cluster to integers, combining
\eqref{r17:eq-weighted-smith} and \eqref{r17:eq-loewner-singular} forces
\begin{equation}
 \delta_k(A)
 \left(\frac{e_k(D^{-2})e_k(E^{-2})}{\binom nk}\right)^{1/2}
 \le C_{n,k,\alpha}\varepsilon^{k(k-1)}.
 \label{r17:eq-loewner-certificate}
\end{equation}
A finite violation would prove the conjecture.  Naive numerator/denominator clearing is
exponentially more expensive than this polynomial target; reduced Smith content or genuine
row--column cancellation is still missing.

\subsection{Positivity and entropy no-go theorems}
\label{r17:no-go}

The most natural $3$-smooth Gram matrices are infinitely divisible: if
\begin{equation}
 K_{ij}=2^{\langle a_i,a_j\rangle}3^{\langle b_i,b_j\rangle},
\end{equation}
then $K^{\circ t}\succeq0$ for every $t\ge0$, and under a counterexample
$K^{\circ x}$ is integral.  Thus positivity of Hadamard powers cannot distinguish integral
from nonintegral $x$ without an additional arithmetic positivity principle.

Probability is equally tight.  If $N,K$ are independent uniform variables on
$\{0,\ldots,H\}$ and
\begin{equation}
 X_H=2^N M^K,\qquad Y_H=3^N A^K,
\end{equation}
then one external-prime valuation and the structural valuation recover $(N,K)$ from either
output separately.  Hence
\begin{equation}
 \mathsf H(X_H)=\mathsf H(Y_H)=\mathsf H(X_H,Y_H)=2\log(H+1),
 \qquad I(X_H;Y_H)=2\log(H+1).
 \label{r17:eq-entropy-saturation}
\end{equation}
The hypothetical semigroup saturates, rather than violates, ordinary anti-concentration and
Shannon inequalities.  Random minors and modular scouting remain useful discovery tools only
when followed by an exact Smith or interval certificate.

\begin{statusbox}{Report 17 was audited independently from the common base.  The exact
geometric-Vandermonde Smith form is \statuslean{} in \lean{GeometricVandermondeSmith}.  The
exterior-minor identities, weighted Smith inequalities, multilevel modular divisor,
Loewner--Andreief/Selberg formulas, singular hierarchy, and entropy saturation remain
self-contained paper arguments with the stated classical linear-algebra and integral inputs.
The universal modular divisor is a baseline, not a closing estimate; the geometric Smith form
is a one-base no-go; and the Loewner certificate still requires reduced denominator or Smith
information not supplied by height clearing.  None of these results proves the conjecture.}
\end{statusbox}

\section{Eighteenth report: logarithmic spectral rank and approximation capacity}
\label{sec:report18}

Report 18 is an independent asymptotic reading of the common exact conditional structure.
It regards the candidate set as the spectrum of a diagonal operator, resolves the associated
heat trace by Mellin inversion, records a second mixed support equation, and compares the
best polynomial approximation rate on a fixed multiplicative window with the number of
conditional integer nodes.  These constructions distinguish a formal rank-one semigroup from
a formal rank-two semigroup; by themselves they do not use the additional arithmetic fact
that the synchronized $\thetaq$th powers are also integers.

\subsection{Spectral zeta and logarithmic Weyl rank}
\label{r18:spectral}

On $\ell^2(\Sol)$ define $De_x=2^x e_x$ and $Ee_x=3^x e_x$.  Positive functional calculus
gives $E=D^{\thetaq}$.  For $\Re s>0$, the exact conditional structure gives
\begin{equation}
 Z_D(s)=\operatorname{Tr}(D^{-s})=
 \begin{cases}
 (1-2^{-s})^{-1},&\Sol=\Nzero,\\[0.3em]
 \bigl((1-2^{-s})(1-M^{-s})\bigr)^{-1},
   &\Sol=\Nzero+\Nzero\beta,
 \end{cases}
 \label{r18:eq-spectral-zeta}
\end{equation}
and $Z_E(s)=Z_D(\thetaq s)$.  Thus the origin is a simple pole in rank one and a double pole
in rank two.  Equivalently, for
\begin{equation}
 N_D(\Lambda)=\#\{x\in\Sol:2^x\le\Lambda\},
\end{equation}
the exact conditional count, together with the paper-level Baker--discrepancy refinement,
gives
\begin{align}
 N_D(\Lambda)&=\frac{\log\Lambda}{\log2}+O(1)
 &&\text{in rank one},\label{r18:eq-weyl-one}\\
 N_D(\Lambda)&=\frac{(\log\Lambda)^2}{2\log2\log M}
 +\left(\frac1{2\log2}+\frac1{2\log M}\right)\log\Lambda
 +O_M((\log\Lambda)^{1-\delta})
 &&\text{in rank two}.\label{r18:eq-weyl-two}
\end{align}
Consequently the limit
\begin{equation}
 \dim_{\log}D:=\lim_{\Lambda\to\infty}
 \frac{\log N_D(\Lambda)}{\log\log\Lambda}
\end{equation}
is respectively $1$ or $2$.  This is an exact reformulation of the monoid dichotomy, not an
independent method for deciding which case holds.

\subsection{Complete heat-trace expansion}
\label{r18:heat}

For a formal independent pair $(2,M)$ put $a=\log2$, $b=\log M$ and
\begin{equation}
 H_{a,b}(t)=\sum_{n,k\ge0}\exp(-t e^{an+bk}).
\end{equation}
Mellin transformation gives
\begin{equation}
 \int_0^\infty H_{a,b}(t)t^{s-1}\,dt
 =\frac{\Gamma(s)}{(1-e^{-as})(1-e^{-bs})}
 \qquad(\Re s>0).
 \label{r18:eq-heat-mellin}
\end{equation}
Since $2$ and $M$ are multiplicatively independent, Baker's theorem supplies polynomial
lower bounds for the near-collisions of the two imaginary pole lattices.  The exponential
vertical decay of $\Gamma$ then makes the residue series absolutely convergent.  With
$L=\log(1/t)$, contour shifting yields, for every $R\ge0$,
\begin{align}
 H_{a,b}(t)
 &=Q_{a,b}(L)+P_a(L)+P_b(L)\\
 &\quad+\sum_{r=1}^{R}
 \frac{(-1)^r t^r}{r!(1-e^{ar})(1-e^{br})}
 +O_{a,b,R}(t^{R+1/2}),
 \label{r18:eq-heat-expansion}
\end{align}
where
\begin{align}
 Q_{a,b}(L)
 &=\frac{L^2}{2ab}
 +\left(\frac1{2a}+\frac1{2b}-\frac{\gamma}{ab}\right)L+C_{a,b},\\
 P_a(L)
 &=\sum_{\ell\ne0}
 \frac{\Gamma(2\pi i\ell/a)}{a(1-e^{-2\pi i\ell b/a})}
 e^{2\pi i\ell L/a},
\end{align}
and $P_b$ is the symmetric series.  The two nonzero frequency sets are disjoint; hence the
bounded remainder contains a genuine quasiperiodic component.  In rank one the leading term
is $L/a$ with one periodic frequency lattice.  These paper formulas are mathematically
sharper than the rough count, but the output heat trace is only the rescaling
$Z_E(s)=Z_D(\thetaq s)$ and supplies no second independent constraint.

\subsection{The finite two-scale support equation}
\label{r18:support}

For natural scales $p,q$ and a commutative ring, set
\begin{equation}
 F_{A,B}(z)=\sum_{0\le a<A}\sum_{0\le b<B}z^{p^a q^b}.
\end{equation}
The exact finite identity
\begin{align}
 F_{A+1,B+1}(z)
 &=z+F_{A,B+1}(z^p)+F_{A+1,B}(z^q)-F_{A,B}(z^{pq})
 \label{r18:eq-support-finite}
\end{align}
is \statuslean{} as \lean{twoScaleSupportTrunc_succ_succ} in
\lean{TwoScaleSupportFinite}.  It is coefficientwise finite and has no analytic hypotheses.

For multiplicatively independent integers $2,M$, passage to the convergent support series
\begin{equation}
 F_{2,M}(z)=\sum_{n,k\ge0}z^{2^nM^k},\qquad |z|<1,
\end{equation}
gives
\begin{equation}
 F_{2,M}(z)-F_{2,M}(z^2)-F_{2,M}(z^M)+F_{2,M}(z^{2M})=z.
 \label{r18:eq-support-infinite}
\end{equation}
Its $0$--$1$ coefficient sequence has density-zero infinite support and is not eventually
periodic, so P\'olya--Carlson gives the unit circle as a natural boundary.  Rank-one lacunary
series already have the same boundary; only the radial singularity order sees the rank.
Existing Mahler theorems for two separate independent functional equations do not apply to
this one mixed inclusion--exclusion equation.

\subsection{A sharp fixed-window node-budget barrier}
\label{r18:approximation}

Fix $h>0$.  On $[1,e^h]$ the nearest singularity of $u^{\thetaq}$ maps, under the affine
normalization to $[-1,1]$, to the Bernstein ellipse parameter
\begin{equation}
 \rho(h)=\coth(h/4).
 \label{r18:eq-ellipse}
\end{equation}
Therefore the least polynomial degree needed to approximate $x^{\thetaq}$ uniformly within
$1/2$ on $[X,e^hX]$ obeys
\begin{equation}
 d_{\min}(X,h)\ge
 \left(\frac{\thetaq}{\log\rho(h)}+o(1)\right)\log X.
 \label{r18:eq-degree-cost}
\end{equation}
Under a counterexample the number of conditional nodes in the same window is
\begin{equation}
 K_X(h)=\frac{h}{\beta(\log2)^2}\log X+o(\log X).
 \label{r18:eq-window-nodes}
\end{equation}
The elementary maximization
\begin{equation}
 \sup_{h>0}h\log\coth(h/4)=2\operatorname{arsinh}(1)^2
 =1.5536387997\ldots
 \label{r18:eq-capacity}
\end{equation}
occurs at $h=2\operatorname{arsinh}(1)$.  Using the kernel-verified lower bound $\beta>12$
therefore gives
\begin{equation}
 \limsup_{X\to\infty}\frac{K_X(h)}{d_{\min}(X,h)}
 \le\frac{h\log\coth(h/4)}{\beta\log2\log3}
 <0.170020.
 \label{r18:eq-node-barrier}
\end{equation}
Thus a fixed-window strategy that first approximates by a degree-$d$ polynomial to sub-unit
accuracy and then demands at least $d+2$ integer nodes is starved by a factor greater than
$5.88$, even before coefficient height and denominator taxes.  This does not exclude rational,
multivariate, semigroup-adapted, or growing-window approximants.

The report also compares the conditional $(\log X)^2$ set with recent irrational-power
additive-relation bounds.  Unequal-length relations such as
$(2m)^{\thetaq}=3m^{\thetaq}$ already occur at the scale allowed by those theorems, while no
extra equal-length energy is forced.  Because the relevant preprints and constants are
time-sensitive, this synthesis retains only the mathematical threshold diagnosis and does
not elevate a dated literature search to a theorem.

\begin{statusbox}{Report 18 was audited independently from the common base.  The finite
two-scale support identity is \statuslean{} in \lean{TwoScaleSupportFinite}.  The zeta/Weyl
reformulations use the kernel-verified exact monoid and leading counting asymptotic, with the
displayed power-saving remainder kept at paper level.  The
complete heat expansion and sharp fixed-window capacity theorem are paper arguments using,
respectively, classical Mellin--Baker input and classical best-polynomial-approximation
asymptotics.  The natural boundary is classical after the finite support algebra.  None of
the spectral, heat, Mahler, approximation, or additive signatures eliminates a formal
rank-two integer semigroup; the missing step is still the arithmetic coupling
$M^{\thetaq}\in\N$.}
\end{statusbox}

\section{Nineteenth report: arithmetic approximation and geometric Newton interpolation}
\label{sec:report19}

Report 19 is an independent approximation-theoretic sibling of the preceding reports.  Its
fixed-window optimization converges on the same capacity constant as report 18, but starts
from the integer-valued-polynomial lattice on the $3$-smooth nodes rather than from a spectral
count.  Its genuinely complementary finite content is an exact interpolation certificate,
the cofactor formula for discrete minimax error, and a geometric Newton expansion with an
explicit first extrapolation error.  These results isolate the arithmetic unit barrier; they
do not manufacture the missing integer-valued approximant.

\subsection{Subunit locking and the optimal multiplicative block}
\label{r19:locking}

Put $\mathcal B=\{2^a3^b:a,b\in\Nzero\}$.  Under a hypothetical solution $\beta$,
$q^\beta\in\Z$ for every $q\in\mathcal B$.  If $E\subset\mathcal B$ is finite and a
degree-$d$ polynomial $P$ is integer-valued on $E$, then
\begin{equation}
 \max_{q\in E}|q^\beta-P(q)|<1
 \quad\Longrightarrow\quad q^\beta=P(q)\quad(q\in E).
 \label{r19:eq-unit-lock}
\end{equation}
The generalized power system $1,t,\ldots,t^d,t^\beta$ is Chebyshev on $(0,\infty)$, so a
nonintegral $\beta$ makes \eqref{r19:eq-unit-lock} impossible as soon as $|E|\ge d+2$.
The same argument gives an exact unit floor for every integer-coefficient M\"untz polynomial
with distinct exponents in $\Nzero+\Nzero\beta$.

For $c>1$, the $3$-smooth nodes in $[X,cX]$ have cardinality
\begin{equation}
 \frac{\log c}{\log2\log3}\log X+o(\log X),
\end{equation}
whereas Bernstein approximation to $t^\beta$ on that block requires degree at least
$\beta\log X/\log\rho(c)+o(\log X)$, with
\begin{equation}
 \rho(c)=\frac{\sqrt c+1}{\sqrt c-1}.
\end{equation}
The exact optimization is
\begin{equation}
 \max_{c>1}\log c\log\rho(c)=2\log^2(1+\sqrt2),
 \qquad c_*=3+2\sqrt2.
 \label{r19:eq-block-optimum}
\end{equation}
Consequently this particular method could have a positive degree-versus-node margin only for
\begin{equation}
 M=2^\beta<\exp\!\left(\frac{2\log^2(1+\sqrt2)}{\log3}\right)
 =4.113124586\ldots.
\end{equation}
Among nonpowers of two only $M=3$ lies in that formal range, and the accepted finite corpus
already excludes it (indeed it proves $M\ge4096$ for a counterexample).  Thus the exact
constant is a sharp no-go for ordinary fixed-block polynomial approximation, consistent with
the report-18 capacity calculation.

\subsection{Dyadic certificates and the discrete minimax unit barrier}
\label{r19:certificates}

Let $0\le n_0<\cdots<n_d$, and let $I_{\mathbf n}\in\Q[T]$ be the degree-at-most-$d$
interpolant with $I_{\mathbf n}(2^{n_j})=3^{n_j}$.  For an integer $M$ outside the node set,
write $I_{\mathbf n}(M)=a/b$ in lowest terms, $b>0$.  The exact certificate is
\begin{equation}
 \left|M^{\thetaq}-\frac ab\right|<\frac1b
 \quad\Longrightarrow\quad M^{\thetaq}\notin\Z.
 \label{r19:eq-dyadic-certificate}
\end{equation}
If the target were integral, its nonzero difference from $a/b$ would have size at least
$1/b$; equality with the interpolant is excluded by the $d+2$-zero Chebyshev bound.  The
quadratic interpolants
\begin{equation}
 I_{0,1,2}(T)=\frac{T^2+3T-1}{3},\qquad
 I_{1,2,3}(T)=\frac{T^2+6T-4}{4}
\end{equation}
give short exact exclusions for $M=3,5,6,7$.  They are conceptually transparent certificates,
not an improvement on the much larger kernel-verified finite range.

More generally, let $A$ be an $(N+1)\times N$ generalized evaluation matrix, $f$ the target
column, and $C_i$ the signed maximal cofactors of $A$.  With
\begin{equation}
 D=\det[A\ f],\qquad S=\sum_{i=0}^N|C_i|,
\end{equation}
strict total positivity gives the exact real minimax value
\begin{equation}
 \min_{c\in\R^N}\|f-Ac\|_\infty=\frac{|D|}{S}.
 \label{r19:eq-minimax}
\end{equation}
When $A$ and $f$ are integral, clearing this canonical denominator produces $z\in\Z^N$ with
\begin{equation}
 Sf-Az=D\,(\operatorname{sgn}C_i)_{i=0}^N.
 \label{r19:eq-arithmetic-alternant}
\end{equation}
Thus the small real residual $D/S$ becomes the nonzero integer residual $D$.  Smith content
or a different integer-valued basis must save part of $S$; ordinary coefficient clearing
cannot cross the unit threshold.

\subsection{The geometric Newton shadow}
\label{r19:qnewton}

For integers $q\ge2$, $r\ge1$, define
\begin{equation}
 B_k^{(q)}(X)=\prod_{i=0}^{k-1}\frac{X-q^i}{q^k-q^i},\qquad
 a_k(q,r)=\prod_{i=0}^{k-1}(r-q^i).
\end{equation}
At $X=q^n$, the first factor is the Gaussian binomial
$\genfrac{[}{]}{0pt}{}{n}{k}_q\in\Z$, and the exact
Newton identity is
\begin{equation}
 r^n=\sum_{k=0}^n \genfrac{[}{]}{0pt}{}{n}{k}_q
      \prod_{i=0}^{k-1}(r-q^i).
 \label{r19:eq-qnewton}
\end{equation}
Hence the truncation through $k=d$ is integer-valued at every geometric node, interpolates
$r^n$ for $n\le d$, and has first omitted error
\begin{equation}
 r^{d+1}-P_d^{q,r}(q^{d+1})=\prod_{i=0}^{d}(r-q^i).
 \label{r19:eq-first-extrapolation}
\end{equation}
Unless $r$ is a nonnegative integral power of $q$, the Newton series never terminates and this
extrapolation product grows rather than shrinks.  This complements the exact geometric Smith
form of report 17: one-axis integer-valued interpolation is abundant, but it supplies no
cross-axis smallness.

\begin{statusbox}{Report 19 was audited independently from the common base.  The unit-locking,
optimal-block, dyadic-certificate, minimax-cofactor, and arithmetic-alternant statements are
self-contained paper arguments using the stated classical Chebyshev/Bernstein inputs.  The
fixed-window constant overlaps report 18 and is already incompatible with the kernel-verified
counterexample range.  The geometric Newton identity and its exact truncation/remainder are
\statuslean{} in \lean{GeometricNewtonInterpolation}, with endpoints
\lean{qNewtonExpansion_eq_pow} and \lean{qNewtonTruncValue_first_omitted}; all broader
Bhargava-ordering, Hermite--Pad\'e, and
integer-Chebyshev proposals remain research interfaces.  No dated literature-search claim is
promoted, and none of these statements proves the conjecture.}
\end{statusbox}

\section{Twentieth report: operator coherence and local-to-global arithmetic}
\label{sec:report20}

Report 20 is an independent operator-theoretic sibling from the same unified base.  It asks
which exact operator extensions would add arithmetic information beyond the two scalar values
$2^x$ and $3^x$.  Bilateral shifts, two-mode thermodynamics, semisimple functional calculus,
and bounded Sturmian cocycles are proved blind; global or finite valuation lattices, one
affine coherence, one rational first jet, or one exact finite realization would close the
problem.  These are sufficient criteria and no-go theorems, not constructions supplied by the
original hypotheses.

\subsection{Universal two-mode dynamics}
\label{r20:local-universality}

On $\ell^2(\Q_{>0})$ let $U_a e_q=e_{aq}$.  If $a,b$ are multiplicatively independent,
the pair $(U_a,U_b)$ is a countable multiple of the regular representation of $\Z^2$:
decompose $\Q_{>0}$ into cosets of $\langle a,b\rangle$.  For the substantive case $x>0$,
the pairs attached to $(2,3)$ and $(2^x,3^x)$ are therefore unitarily equivalent.  On the
two-mode Fock space,
\begin{equation}
 H_{2^x,3^x}=xH_{2,3},
\end{equation}
so heat traces, local KMS data, and Weyl laws are merely reparametrized in time.  This
reconciles with report 18: logarithmic spectral rank detects one versus two generators, but
cannot quantize a positive rescaling inside the rank-two sector.

The same blindness holds for every rational semisimple matrix $T$ with
$(T-2I)(T-3I)=0$.  Writing $M=2^x$ and $A=3^x$,
\begin{equation}
 T^x=(A-M)T+(3M-2A)I,
 \label{r20:eq-semisimple}
\end{equation}
which is rational, or integral when the displayed data are integral.  Coupling the distinct
spectral points probes only their already-known secant.

\subsection{Global and finite arithmetic covariance}
\label{r20:arithmetic-covariance}

On $\ell^2(\N)$ let $He_n=(\log n)e_n$.  For $c>0$ the following are equivalent:
\begin{equation}
 \exists\text{ an isometry }T\text{ with }HT=cTH;
 \quad n^c\in\N\ (n\in\N);
 \quad c\in\N.
 \label{r20:eq-global-covariance}
\end{equation}
Spectral simplicity gives the middle condition, while the forward difference of order
$\lfloor c\rfloor+1$ is a nonzero integer tending to zero unless $c$ is integral.  Conversely
$T e_n=e_{n^c}$ works for positive integral $c$.  This is much stronger than the original
hypotheses, whose canonical partial intertwiner is supported only on the $3$-smooth sector.

A finite prime lattice already suffices.  Let $P=\{p_1,\ldots,p_d\}$ contain $2$, assume
$p_j^c\in\langle P\rangle$ for every $j$, and suppose $2^c$ is algebraic.  The valuation
columns form $V\in M_d(\Z)$ with
\begin{equation}
 V^{\mathsf T}(\log p_i)_i=c(\log p_i)_i.
\end{equation}
Thus $c$ is an algebraic integer; Gelfond--Schneider and algebraicity of $2^c$ make it
rational, hence integral.  If the lattice map is onto, $V$ is unimodular and positive $c$
must equal $1$.  Finiteness and exact valuation support are the arithmetic input; an abstract
Hilbert-space extension does not provide them.

\subsection{Affine rulers and the first-jet criterion}
\label{r20:affine-jet}

On $\ell^2(\Nzero)$ let $Se_k=e_{k+1}$ and $V_a e_k=e_{ak}$.  For positive integers
$a,b,d$,
\begin{equation}
 D_{a,b,d}=V_aS^d-S^{bd}V_a
\end{equation}
vanishes exactly when $a=b$; otherwise every basis vector is sent to the difference of two
distinct basis vectors, so $D_{a,b,d}$ is noncompact and
$\|D_{a,b,d}\|_{\mathrm{ess}}\ge\sqrt2$.  Hence preserving the single relation
$V_2S=S^2V_2$ after $V_2\mapsto V_{2^x}$, even modulo compacts, forces $2^x=2$ and $x=1$.
Translation acts as an external integral ruler absent from the multiplicative subsystem.

At one positive repeated spectral point the terminating functional calculus gives, for
$b>0$ and $N^2=0$,
\begin{equation}
 (bI+N)^x=b^xI+xb^{x-1}N.
 \label{r20:eq-jordan}
\end{equation}
If $b^x$ and one nonzero first-jet coefficient $xb^{x-1}$ are rational, then
$x=b\,(xb^{x-1})/b^x$ is rational.  Under the original integral-output hypotheses this
forces $x\in\Z$.  Distinct-node Loewner entries are secants and automatic; a self-extension
is the first finite-dimensional datum that sees the missing derivative.

\subsection{Sturmian similarity and exact finite realization}
\label{r20:koopman}

Writing $x=n+\alpha$, define the rotation jump
$w(t)=\lfloor t+\alpha\rfloor$ on $[0,1)$.  The rational two-valued cocycles
\begin{equation}
 c_b(t)=b^{w(t)-\alpha}\qquad(b=2,3)
\end{equation}
are bounded coboundaries because $w(t)-\alpha=t-\{t+\alpha\}$.  Their weighted Koopman
operators are therefore similar to the rotation itself; spectra, entropy, Lyapunov growth,
and determinant-like invariants are blind.  The surviving projection $P=M_w$ has trace
$\tau(P)=\alpha$.  An exact finite-dimensional trace model would make $\alpha$ rational.
Likewise, finite Hankel rank of
$w_k=\lfloor(k+1)\alpha\rfloor-\lfloor k\alpha\rfloor$ would make this binary sequence
eventually periodic and its mean rational.  Either exact realization closes the conjecture;
ordinary finite-dimensional approximation does not.  The deterministic carry identities
overlap the accepted \lean{SturmianProbability} and \lean{SturmianGeneratingSeries} modules.

\begin{statusbox}{Report 20 was audited independently from the common base.  Its bilateral
universality, time-homothety, global and finite covariance criteria, affine essential-norm
rigidity, semisimple blindness, Jordan first-jet criterion, and finite trace/Hankel criteria
are self-contained paper arguments with the named standard operator inputs.  The maximal
algebraic support is already covered more generally by \lean{AlgebraicBaseSpectrum}, and the
Sturmian/heat cores overlap earlier accepted modules.  The rational scalar and matrix-entry
first-jet recovery is \statuslean{} in \lean{JordanFirstJet}; the theorem
\lean{integer_of_two_rpow_integer_of_jordanFirstJet_rational} formalizes the arithmetic
closing step once the square-zero holomorphic functional-calculus formula is supplied.  No
global intertwiner, finite prime completion, affine
coherence, or finite realization is derived from the two original outputs.  None of these
criteria proves the conjecture.}
\end{statusbox}

\section{Varying the antecedent, the domain, and the codomain}
\label{sec:variants}

The conjecture fixes three choices that look incidental: the values are asked to be
\emph{integers}, the exponent is asked to be \emph{real}, and the codomain is $\Z$ rather than
a larger ring.  This section records what each choice is worth.  The short answer is that the
conjecture predicts the same thing in every reading; what changes is only which wall stands
guard.

\subsection{The antecedent ladder}

Write the three antecedents
\[
  (\mathbb A_2)\ 2^x,3^x\in\Qbar\Rightarrow x\in\Q,
  \qquad
  (\Q_2)\ 2^x,3^x\in\Q\Rightarrow x\in\Z,
  \qquad
  (\Z_2)\ 2^x,3^x\in\Z\Rightarrow x\in\Z .
\]
The three-base theorem decomposes into two fully decoupled mechanisms, and only one of them
ever sees the difference between $\Z$, $\Q$ and $\Qbar$.

\emph{Mechanism one: transcendence.}  If $2^x,3^x,5^x\in\Qbar$ and $x\notin\Q$ then $\{1,x\}$
is $\Q$-linearly independent, $\{\log2,\log3,\log5\}$ is $\Q$-linearly independent by unique
factorization, and the six products exponentiate to $2,3,5,2^x,3^x,5^x$, all algebraic,
contradicting six exponentials.  This step is blind to integrality: it uses only algebraicity.
It is the entire transcendence content of the three-base theorem.

\emph{Mechanism two: descent from $\Q$ to $\Z$.}  If the values are rational, write $x=p/q$ in
lowest terms.  Then $2^{p/q}\in\Q$ is a rational root of the monic $X^q-2^p$, hence an integer
$m$, and comparing $2$-adic valuations in $2^p=m^q$ gives $p=q\,\vpp2m$, so $q\mid p$ and
$q=1$.  One rational value at one base kills the denominator, and the same works for values in
any subring of $\Q$.  The dividing line is simply whether the value ring sits inside $\Q$ or
contains irrational algebraic numbers.

\begin{statusbox}{\statusknown{} for mechanism one; \statuslean{} for mechanism two, which is
\lean{integer_of_rational_of_two_rpow_integer} in \lean{IntegerExponent}.  The implications
$(\mathbb A_2)\Rightarrow(\Q_2)\Rightarrow(\Z_2)$ follow, the first by mechanism two and the
second trivially.}
\end{statusbox}

The $\Qbar$-conclusion is sharp.  Every rational $x=p/q$ genuinely occurs, since
$2^{p/q},3^{p/q},5^{p/q}$ are algebraic of degree exactly $q$: the binomial $X^q-2^p$ is
irreducible by Capelli's criterion for $\gcd(p,q)=1$, because $2^p$ is positive and is not an
$\ell$-th power for any $\ell\mid q$.  So the $\Qbar$-antecedent solution set is exactly $\Q$,
and no strengthening of six exponentials can improve the conclusion.  The same degree
computation gives a graded refinement: if $2^x,3^x,5^x$ are algebraic of degree at most $d$,
then $x\in\Q$ with denominator at most $d$.  The kernel-verified exact radical degree of
\lean{RadicalDegree} is the two-base conditional shadow of exactly this statement.

\begin{remark}[Where primality is used]
Mechanism one needs only that the bases be multiplicatively independent, but mechanism two
needs a base that is not a perfect power.  The bases $4,9,25$ are multiplicatively independent,
yet $x=1/2$ gives $4^x=2$, $9^x=3$, $25^x=5$, all integers.  For general multiplicatively
independent integer bases the correct statement is: rational values force $x\in\Q$ with a
denominator $q$ such that every base is a perfect $q$-th power.  This is exactly the phenomenon
that the power-primitive normalization of Section~\ref{sec:conditional} exists to quotient
away.
\end{remark}

The rational form $(\Q_2)$ deserves separate billing: it is the conjecture as Alaoglu and
Erd\H{o}s posed it, and it is the form the colossally abundant application consumes.  Its
solution set is exactly $\Z$, negative exponents included, which restores the $x\leftrightarrow
-x$ symmetry that the $\Z$ form breaks by truncating to $\Nzero$.

\begin{proposition}[The exact cost of the rational form]
\label{va:prop-rational-cost}
$(\Q_2)$ is equivalent to $(\Z_2)$ together with a denominator-support lemma: if $2^x=a/b$ and
$3^x=c/d$ in lowest terms, then $b$ is a power of $2$ and $d$ is a power of $3$.
\end{proposition}

\begin{proof}
Given the support lemma, $y=x+s$ for large $s$ has $2^y=2^sa/b\in\Z$ and $3^y=3^sc/d\in\Z$,
with $y$ nonintegral exactly when $x$ is, so $(\Z_2)$ applies.  Conversely $(\Q_2)$ makes the
support lemma vacuous, every solution being integral.
\end{proof}

\begin{statusbox}{\statusnew{} --- absent the support lemma no reduction is known in either
direction.  It is the same kind of statement as the kernel-verified rational gateway of
\lean{RationalThirdBase}, and the modules \lean{ThreeDenominatorNormalization} and
\lean{RationalPowerIndex} operate inside exactly this gap without stating the equivalence.}
\end{statusbox}

\subsection{Complexifying the exponent is free}

On the principal branch the restriction to real exponents is cosmetic.

\begin{theorem}[Every complex solution is real]
\label{va:thm-complex-real}
Let $x\in\C$ and let $2^x=\exp(x\log2)$, $3^x=\exp(x\log3)$ be the principal determinations.
If $2^x,3^x\in\Z$ then $x\in\R$.  The same holds with $\Z$ replaced by $\Q$.
\end{theorem}

\begin{proof}
Write $x=s+it$.  The modulus of $2^x$ is $2^s$, which reproduces the real hypothesis, and its
argument is $t\log2$; an integer value is real, so $t\log2\in\pi\Z$, and likewise
$t\log3\in\pi\Z$.  Write $t\log2=\pi a$ and $t\log3=\pi b$ with $a,b\in\Z$.  If $t\ne0$ then
$a\ne0$, and dividing gives $\thetaq=b/a\in\Q$, contradicting irrationality of $\thetaq$.
\end{proof}

\begin{statusbox}{\statuslean{} for the integer-codomain version ---
\lean{im_eq_zero_of_complexTwoBaseIntegralSolution} and
\lean{twoBaseIntegralSolution_re_of_complex} in \lean{ComplexSolutions}.  The rational-codomain
realness statement follows by the same paper argument but is not exported there.  Only
irrationality of $\thetaq$ enters; no transcendence input is used.}
\end{statusbox}

Geometrically the sets $\{x:2^x\in\Z\}$ and $\{x:3^x\in\Z\}$ lie on horizontal gratings with
spacings $\pi/\log2$ and $\pi/\log3$, and the two spacings are incommensurable \emph{precisely
because} $\thetaq$ is irrational, so the gratings meet only on the real axis.  The conjecture
may therefore be stated over $\C$ at no mathematical cost.

\subsection{The algebraic antecedent over the complex numbers}

Here the reduction to $\R$ genuinely fails --- algebraic values need not have argument in
$\{0,\pi\}$ --- but something cleaner is true.  Since $\Qbar$ is closed under conjugation and
division, $z\in\Qbar\setminus\{0\}$ has $|z|=\sqrt{z\bar z}\in\Qbar$.  So the complex solution
set splits as a direct sum
\[
  \{x\in\C:2^x,3^x\in\Qbar\}
  =\{s\in\R:2^s,3^s\in\Qbar\}+i\,\{t\in\R:2^{it},3^{it}\in\Qbar\} .
\]
The real slice is $(\mathbb A_2)$; the imaginary slice is the unit-circle problem.  Both are
four-exponentials instances --- genuinely different instances, with rows $\{1,s\}$ and
$\{1,it\}$ against the same columns $\{\log2,\log3\}$ --- so the complex two-base algebraic
problem is exactly two decoupled instances of the same conjecture, and under it the conclusion
is still $x\in\Q$, still sharp.

For three bases everything closes unconditionally: six exponentials is a statement over $\C$,
the real slice gives $s\in\Q$, and the imaginary slice is the standard corollary that
$2^{it},3^{it},5^{it}$ cannot all be algebraic for real $t\ne0$.  The three-base algebraic
solution set over $\C$ is therefore exactly $\Q$; complexification adds nothing.  One
unconditional crumb survives for two bases: nonreal \emph{algebraic} $x$ are excluded, because
Gelfond--Schneider holds for complex algebraic irrational exponents, so the transcendence
theorem of Section~\ref{sec:corpus} extends to $\C$ verbatim.

\subsection{Multivalued branches widen the wall by one prime}

Demand only that \emph{some} determination of each power be an integer:
$e^{xL_2}=m$ and $e^{xL_3}=m'$ with $L_2=\log2+2\pi iJ$ and $L_3=\log3+2\pi iK$.  Writing
$xL_2=\log|m|+i\pi a$ and $xL_3=\log|m'|+i\pi b$ and eliminating $x$, the imaginary part gives
a $\Z$-linear relation among logarithms, hence by unique factorization the multiplicative
constraint $|m|^{2K}3^a=|m'|^{2J}2^b$, while the real part gives
\begin{equation}
  \log|m|\log3-\log|m'|\log2=2\pi^2\,(aK-bJ).
  \label{eq:va-branch}
\end{equation}
If $aK-bJ=0$ the left side vanishes, which is the log-product wall of
Section~\ref{sec:logproduct}: either $|m|=|m'|=1$, forcing $x=0$, or $(|m|,|m'|)$ is a real
counterexample pair.  If $aK-bJ\ne0$ one needs the log-product form to equal a nonzero integer
multiple of $2\pi^2$.  Since $i\pi=\log(-1)$, this is exactly the same wall over the prime set
extended by the ``prime'' $-1$: the bilinear lattice gains the directions
$\pi^2=-(\log(-1))^2$ and $i\pi\log p$.

\begin{statusbox}{\statusnew{} --- the independence criterion of
Section~\ref{sec:logproduct}, extended to this larger lattice, covers the multivalued reading
as well, and it remains a consequence of Schanuel's conjecture.  The four exponentials
conjecture masters the multivalued reading wholesale, because the branch logarithms stay
$\Q$-independent: it gives $x\in\Q$, and then the modulus argument forces $x\in\Z$ with all
phases trivial.  The predicted answer is unchanged; only the wall widens.}
\end{statusbox}

\subsection{Complexifying the codomain: the Gaussian version is equivalent}

Since extending the domain is free, the honest place to look for new content is the codomain.

\begin{theorem}[Gaussian equivalence]
\label{va:thm-gaussian}
The statement ``$2^x,3^x\in\Z[i]\Rightarrow x\in\Z$'' is equivalent to
Conjecture~\ref{conj:AE}.
\end{theorem}

\begin{proof}
One direction is trivial since $\Z\subset\Z[i]$.  Conversely assume the conjecture and let
$2^x,3^x\in\Z[i]$ with $x=s+it$.  Taking norms, $N(2^x)=2^{2s}$ and $N(3^x)=3^{2s}$ are
positive integers, which is the real hypothesis at exponent $2s$, so $2s\in\Z$.  Now $3$ is
inert in $\Z[i]$, so a Gaussian integer of norm $3^{2s}$ is a unit times $3^{k}$ with
$9^{k}=3^{2s}$, forcing $s=k\in\Z$: the inert prime upgrades the half-integer ambiguity to
integrality.  Then $2^{2s}=4^{s}$ shows that $2^x$ is divisible only by the ramified prime
$1+i$, so $2^x=u(1+i)^{2s}=2^{s}i^{s}u$.  In both bases the values are $2^{s}$ and $3^{s}$
times fourth roots of unity, the only Gaussian integers of modulus one, so the argument
conditions read $t\log2,\ t\log3\in(\pi/2)\Z$, and irrationality of $\thetaq$ kills $t$ as in
Theorem~\ref{va:thm-complex-real}.
\end{proof}

\begin{statusbox}{\statusnew{} --- a self-contained equivalence whose only inputs are Gaussian
factorization, irrationality of $\thetaq$, and the real conjecture as a hypothesis.  It is a
low-cost Lean target on top of \lean{ComplexSolutions}.}
\end{statusbox}

Two boundary markers make the picture sharp.  Replacing $\Z[i]$ by $\Q(i)$ breaks the argument
at exactly one point: unit-modulus elements of $\Q(i)$ need not be roots of unity, as
$(3+4i)/5$ shows, so the imaginary slice reopens as a genuine unit-circle four-exponentials
instance.  It is integrality up to units, not rationality, that closes the argument direction.
Replacing $\Z[i]$ by the integers of an imaginary quadratic field in which $2$ or $3$
\emph{splits} --- for instance $\Q(\sqrt{-2})$, where $3$ splits and $1+\sqrt{-2}$ has modulus
$\sqrt3$ --- destroys the root-of-unity conclusion, and eliminating $t$ produces relations
mixing $\thetaq\pi$ with logarithms of complex algebraic numbers: the product-of-logarithms
wall again, with more rooms.  These split-prime variants are recorded as open rather than
pursued; they meet the same wall with worse constants.

\subsection{Summary}

\begin{center}
\begin{tabular}{p{0.34\textwidth}p{0.58\textwidth}}
\toprule
Reading & Status \\
\midrule
Exponent in $\C$, principal branch, values in $\Z$ & reduces to the real integer formulation,
\statuslean \\
Exponent in $\C$, principal branch, values in $\Q$ & reduces to the real rational-output
formulation at paper level; equivalence to integer AE still needs the denominator-support
lemma of Proposition~\ref{va:prop-rational-cost} \\
Values in $\Qbar$, exponent real & two-base: four exponentials; conclusion $\Q$, sharp \\
Values in $\Qbar$, exponent in $\C$ & splits into real and unit-circle slices \\
Three bases, values in $\Qbar$, exponent in $\C$ & solution set exactly $\Q$, \statusknown \\
Multivalued branches & wall widened by $\log(-1)$, \statusnew \\
Values in $\Z[i]$ & unconditionally equivalent, \statusnew \\
Values in $\Q(i)$ or split-prime rings & open \\
\bottomrule
\end{tabular}
\end{center}

In every reading the answer the conjecture predicts is the same.

\appendix

\section{Status matrix for major statements}
\label{app:status}

The status tags have the following meanings.  \statuslean{} marks a statement that the Lean
kernel accepts at the current repository state, with no \lean{sorry}, no extra axioms, and no
\lean{native_decide}.  \statusnew{} marks a paper proof given in this report that has not yet
been formalized.  \statusdraft{} marks a statement for which Lean source exists in the
repository but has not yet been accepted by the kernel.  \statuscomp{} marks a finite
software computation with its local trust boundary---for example CPython/\lean{mpmath}
interval arithmetic or C/MPFR/GMP directed rounding---which is not a kernel certificate.
\statusknown{} marks classical or open background.  No entry below asserts the
full Alaoglu--Erd\H{o}s conjecture; it remains open.

No row of the table carries \statusdraft{} any longer.  Every module that earlier revisions
of this report listed as a pending draft has since been accepted by the kernel and placed on
the aggregate import surface.  The tag still applies to the separate, unassembled $16384$
finite-certificate attempt described in Section~\ref{sec:finitecheck}, which is not a row of
this major-statement table and supplies no theorem.

{\small
\begin{longtable}{>{\raggedright\arraybackslash}p{0.37\textwidth}
                  >{\raggedright\arraybackslash}p{0.22\textwidth}
                  >{\raggedright\arraybackslash}p{0.33\textwidth}}
\toprule
Statement & Status & Location \\
\midrule
\endfirsthead
\toprule
Statement & Status & Location \\
\midrule
\endhead
\multicolumn{3}{@{}l}{\emph{Classical frame}}\\
$2^x,3^x,5^x\in\Z\Rightarrow x\in\Z$ & \statuslean & \lean{IntegerExponent} \\
Gelfond--Schneider theorem & \statuslean & \lean{GelfondSchneider} \\
Four exponentials conjecture (hypothesis only) & \statusknown &
\cite{Waldschmidt2023,TheoremDB2026} \\
Full Alaoglu--Erd\H{o}s conjecture & Open &
Section~\ref{sec:radical}, Section~\ref{sec:program} \\
\midrule
\multicolumn{3}{@{}l}{\emph{Transcendence and reduction}}\\
Nonintegral solution irrational and transcendental & \statuslean &
\lean{TwoBaseIntegerExponent}, \lean{Transcendence} \\
$\thetaq=\log_2 3$ transcendental & \statuslean & \lean{Transcendence} \\
Four-exponentials conjecture $\Rightarrow$ Conjecture~\ref{conj:AE} & \statuslean &
\lean{FourExponentialsReduction} \\
Prime-log-product independence $\Rightarrow$ Conjecture~\ref{conj:AE} & \statuslean &
Section~\ref{sec:logproduct}; \lean{LogProductIndependence} \\
Prime-log-product independence $\Rightarrow$ the rational-output form
\eqref{eq:lp-rational-form} & \statuslean &
Section~\ref{sec:logproduct}; \lean{LogProductIndependence} \\
Two-prime case of the independence criterion, and its sharpness & \statuslean &
Section~\ref{lp:twoprime} \\
\midrule
\multicolumn{3}{@{}l}{\emph{Zero-free regions and finite searches}}\\
$2^x<100$ and $2^x<256\Rightarrow x\in\Z$ & \statuslean &
\lean{FiniteCheck}, \lean{FiniteCheck256} \\
$2^x<4096\Rightarrow x\in\Z$\ \ (so $x\ge12$) & \statuslean & \lean{FiniteCheck4096} \\
No nonintegral candidate with $2^x<2^{20}$ & \statuscomp &
Section~\ref{sec:finitecheck}, Appendix~\ref{app:interval} \\
No nonintegral candidate with $2^x<2^{27}$ & \statuscomp &
Section~\ref{sec:finite27} \\
Reported scan through $2^{29}$; canonical logs and replay artifacts absent from the merged
report directory & \statuscomp & Section~\ref{mc:finite29} \\
\midrule
\multicolumn{3}{@{}l}{\emph{Exact conditional structure}}\\
Rank of candidate exponent vectors $\le2$ & \statuslean & \lean{MultiplicativeRank} \\
Cross-valuation identity and conditional box growth & \statuslean &
\lean{CandidateStructure} \\
Common odd core $2^iw^j$ (existence) & \statuslean & \lean{OddCore} \\
Counting $\clog_2N\cdot\clog_3N$; linear dyadic blocks & \statuslean & \lean{OddCore} \\
Primitive (canonical) odd core; unique coordinates & \statuslean &
\lean{PrimitiveOddCore} \\
Unique least solution; affine primitive output pair & \statuslean & \lean{LeastSolution} \\
Simultaneous primitive output normalization; gcd one; size dichotomy & \statuslean &
\lean{OutputNormalization} \\
Power-of-three denominator normalization & \statuslean &
\lean{ThreeDenominatorNormalization} \\
Rational-power index and exact solution monoid $\Nzero+\Nzero\beta$ & \statuslean &
\lean{RationalPowerIndex}, \lean{PrimitiveGenerator} \\
Canonical primitive generator; exact candidate monoid $\{2^nM^k\}$ & \statuslean &
\lean{PrimitiveGenerator} \\
Denominator-controlled rational third output and determinant margin & \statuslean &
\lean{RationalThirdBase} \\
Unrestricted rational-third-output determinant grid & \statuslean &
\lean{RationalSixExponentials}, \lean{RationalBaseThird} \\
$3$-smooth classification of further bases & \statuslean & \lean{ThreeSmooth} \\
Both outputs $3$-smooth $\Rightarrow x\in\Z$ (unconditional exclusion) & \statuslean &
\lean{SmoothOutputs} \\
Sharp form: a $3$-smooth output forces a partner with a prime factor $\ge5$ & \statuslean &
\lean{SmoothOutputs} \\
Primitive output cores are never exactly $3$ and $2$; exact core rationality
classification & \statuslean & \lean{CoreIndependence} \\
\midrule
\multicolumn{3}{@{}l}{\emph{Algebraic auxiliary bases and output loci}}\\
Six exponentials for three positive rational bases and rational outputs & \statuslean &
\lean{RationalSixExponentials} \\
Field-norm determinant bound over rings of integers & \statuslean &
\lean{AlgebraicIntegerDeterminant} \\
One number field and one natural denominator for a finite algebraic family & \statuslean &
\lean{AlgebraicOutputBridge}, \lean{FiniteAlgebraicOutputBridge} \\
Six exponentials for three positive algebraic real bases and algebraic outputs &
\statuslean & \lean{AlgebraicSixExponentials} \\
Algebraic output at a natural third base & \statuslean & \lean{AlgebraicThirdBase} \\
Exact rational-base spectrum: $q^x\in\Q$ iff $x\in\Z$ or $q$ is a $2,3$-unit & \statuslean &
\lean{RationalBaseThird} \\
Exact algebraic-output form of the same classification & \statuslean &
\lean{AlgebraicRationalBase} \\
Divisible-hull criterion: an algebraic base with algebraic power has a positive
$2,3$-unit power & \statuslean & \lean{AlgebraicBaseUnitPower} \\
Pointwise algebraic natural-base classification and iterated-output transcendence &
\statuslean & \lean{AlgebraicConsequences} \\
Exact algebraic output monomial locus is a coordinate face or the origin & \statuslean &
\lean{AlgebraicOutputLocus} \\
Algebraic output ratios lie on one external $p$-adic line (lattice rank $\le1$) &
\statuslean & \lean{AlgebraicOutputLocus} \\
Common algebraic-output exponent plane: $y\in\Q+\Q x$ & \statuslean &
\lean{AlgebraicExponentLocus} \\
Exact rational-function locus $P/Q\in\Q+\Q x$ iff $P=(q+rX)Q$; canonical
$\Nzero$-cone comparison & \statuslean &
{\scriptsize\lean{RationalFunctionAlgebraicLocus}} \\
General two-base saturated spectrum $\GammaSat(a,b)=\{a^rb^s:r,s\in\Q\}$ & \statusnew &
algebraic spectrum (this report) \\
Strong spectrum: $x\log\alpha\notin\LL$ for $\alpha\in\Apos\setminus\GammaSat$ &
\statusnew & algebraic spectrum (this report) \\
Every $\alpha\in\Apos\setminus\GammaSat$ is a one-output detector for
Conjecture~\ref{conj:AE} & \statusnew & algebraic spectrum (this report) \\
\midrule
\multicolumn{3}{@{}l}{\emph{One-base equivalent formulations}}\\
Conjecture~\ref{conj:AE} $\iff$ every solution has $5^x\in\Qbar$ & \statuslean &
\lean{AlgebraicConsequences} \\
Conjecture~\ref{conj:AE} $\iff$ every solution has $2^{x^2},3^{x^2}\in\Qbar$ &
\statuslean & \lean{AlgebraicConsequences} \\
Conjecture~\ref{conj:AE} $\iff$ every solution has $6^{x^2}\in\Qbar$ & \statuslean &
\lean{AlgebraicConsequences} \\
Conjecture~\ref{conj:AE} $\iff$ the algebraic locus among $(2^i/3^j)^{x^2}$ has rank two &
\statuslean & \lean{AlgebraicOutputLocus} \\
Conjecture~\ref{conj:AE} $\iff$ the algebraic real-power output locus has rank three &
\statuslean & \lean{AlgebraicRankThreeEquivalence} \\
Conjecture~\ref{conj:AE} $\iff$ no squared-monomial algebraic coordinate face &
\statuslean & \lean{AlgebraicOutputLocus} \\
Conjecture~\ref{conj:AE} $\iff$ $2^y,3^y\in\Qbar$ at every natural power $y=x^n$ &
\statuslean & \lean{AlgebraicExponentLocus} \\
Conjecture~\ref{conj:AE} $\iff$ both outputs are algebraic after evaluation of any fixed
$P\in\Q[X]$ of degree at least two & \statuslean &
{\scriptsize\lean{RationalFunctionAlgebraicLocus}} \\
\midrule
\multicolumn{3}{@{}l}{\emph{Canonical radical: valuation, degree, and hull}}\\
Odd-prime valuation gcd; outer-power exponents of the canonical radical & \statuslean &
\lean{CanonicalRadicalValuation} \\
Exact rational outer-power criterion; the triple gcd is not replaceable by either
pairwise gcd & \statuslean & \lean{CanonicalRadicalSharpness} \\
Norm and trace of the pure-radical generator; the base-swapped relation blocks
monomial injectivity & \statuslean & \lean{CanonicalRadicalFieldNorm} \\
Algebraic integrality of the radical iff trivial reduced denominator & \statuslean &
\lean{CanonicalRadicalThirdOutput} \\
Algebraicity of the radical's $\beta$-th power iff $2,3$-smooth numerator & \statuslean &
{\scriptsize\lean{CanonicalRadicalAlgebraicOutput}} \\
External numerator prime excludes the positive $2,3$-unit divisible hull &
\statuslean & \lean{CanonicalRadicalDivisibleHull} \\
Mixed radical monomials retain an external prime; transcendental outputs & \statuslean &
\lean{CanonicalRadicalMixedOutput} \\
Exact rational-power indices and radical degrees of the simultaneous radicals; product
degree $de$ & \statuslean & \lean{SimultaneousRadicalDegrees} \\
Localized-radical conjecture target itself & Open & Section~\ref{sec:radical} \\
\midrule
\multicolumn{3}{@{}l}{\emph{Valuation lattice, kernel, and towers}}\\
Rational-output lattice $\Lambda_x$ of the prime-valuation matrix $V_x$ & \statusnew &
valuation lattice (this report) \\
Smith classification: $\Lambda_x/\Z^2\cong\Z/d_1\Z\oplus\Z/d_2\Z$ and the finite
rational-output defect & \statusnew & valuation lattice (this report) \\
Second-iterate kernel $K_x=\{(r,s)\in\Q^2:M^rA^s\in\GammaSat\}$; $\dim_\Q K_x\le1$ &
\statuslean & \lean{SecondIterateKernel} \\
Kernel invariance $K_x=K_y$ across nonintegral solutions & \statuslean &
\lean{SecondIterateKernel} \\
M\"obius criterion: $K\ne\{0\}$ iff $x\in\PGL_2(\Q)\cdot\thetaq$, iff
$1,\thetaq,x,x\thetaq$ are $\Q$-dependent & \statuslean &
\lean{KernelDichotomy} \\
Multiplicative-dependence form: $K\ne\{0\}$ iff $M,A,2,3$ are multiplicatively
dependent & \statuslean & \lean{KernelDichotomy} \\
Four-type classification of counterexamples by external valuation rank & \statusnew &
second-iterate kernel (this report) \\
No nontrivial two-step return inside $\GammaSat$ & \statusnew &
depth-three tower (this report) \\
Depth-three tower rigidity: some $\alpha^{x^j}$, $j\le3$, is transcendental, with exact
first depth & \statusnew & depth-three tower (this report) \\
Universal three-level detector: Conjecture~\ref{conj:AE} $\iff$
$\alpha^x,\alpha^{x^2},\alpha^{x^3}\in\Qbar$ for one fixed $\alpha\ne1$ & \statusnew &
depth-three tower (this report) \\
Base-$2$ cubic-tower equivalence: Conjecture~\ref{conj:AE} $\iff$
$2^{x^2},2^{x^3}\in\Qbar$ & \statuslean & \lean{TowerRigidity} \\
Base-$2$ depth-two dichotomy: $2^{x^2}\in\Qbar$ iff $2^x$ is $3$-smooth, and then
$2^{x^2}\in\N$ while $2^{x^3}$ is transcendental & \statuslean & \lean{TowerRigidity} \\
Exact transition $2^y,2^{xy}\in\Qbar$ iff $y\in\Q+\Q\thetaq$; at most one positive depth
starts an adjacent algebraic pair; two fixed distinct depths give an AE equivalence & \statuslean &
\lean{AlgebraicTowerTransition} \\
Exact base-$3$ transition locus; no positive base-$2$/base-$3$ pair can coexist; across both
principal towers at most one positive transition depth, with higher rational-unit directions
excluded in the dependent-output branch & \statuslean & \lean{CrossBaseTowerTransition} \\
Symmetric tower constants: at least one of $\log E_3$, $\log E_2$ lies outside $\LL$ &
\statusnew & symmetric tower constants (this report) \\
Exact conditional classification of $E_3=3^{\thetaq}$ and $E_2=2^{1/\thetaq}$ &
\statuslean & \lean{SymmetricTowerConstants} \\
$E_3$ and $E_2$ are never both algebraic under a nonintegral solution & \statuslean &
\lean{SymmetricTowerConstants} \\
\midrule
\multicolumn{3}{@{}l}{\emph{Counting, growth, and equidistribution}}\\
Exact block counts and cumulative finite sum & \statuslean & \lean{ExactCounting} \\
Cumulative $T^2/(2\beta)+O(T)$ asymptotic; subquadratic-growth equivalence & \statuslean &
\lean{QuadraticAsymptotic} \\
Vanishing nonzero block Fourier modes & \statuslean & \lean{BlockWeyl} \\
Exact block averages for finite trigonometric polynomials & \statuslean &
\lean{BlockEquidistribution} \\
Full continuous-test/interval blockwise equidistribution & \statusnew &
Section~\ref{sec:conditional} \\
Density dichotomy, gap equivalences, infinitude & \statuslean &
\lean{FractionalPartDensity} \\
Nonintegral candidate refinement; density zero & \statuslean & \lean{NonintegerDensity} \\
Shrinking-target Diophantine bounds & \statuslean & \lean{ShrinkingTarget} \\
Convergent-denominator growth $q_{n+1}<2m^{q_n}$ & \statuslean & \lean{DenominatorGrowth} \\
\midrule
\multicolumn{3}{@{}l}{\emph{Difference and progression obstructions}}\\
Affine descent, $r+8k\le x$ & \statuslean & \lean{ArithmeticRigidity} \\
Three-term AP obstruction under $h^5\le n$ & \statuslean & \lean{ZeroDensity} \\
Sharp three-term and four-term AP obstructions ($r=2,3$) & \statuslean &
\lean{SharperProgressions}, \lean{ThirdDifference} \\
Uniform arbitrary-order AP obstruction ($r\ge2$) & \statuslean & \lean{HigherDifference} \\
Effective rational three-term criterion $h^{400}\le n^{83}$ & \statuslean &
\lean{SharpProgression} \\
\midrule
\multicolumn{3}{@{}l}{\emph{Radical and localized reformulations}}\\
Localized-radical and symmetric reformulations & \statuslean &
\lean{RadicalReformulation}, Appendix~\ref{app:symmetric} \\
Primitive localized-radical reduction & \statuslean & \lean{PrimitiveRadical} \\
Exact radical degree and binomial irreducibility & \statuslean & \lean{RadicalDegree} \\
Candidate-form localization equivalence & \statuslean & \lean{Localization} \\
\midrule
\multicolumn{3}{@{}l}{\emph{New results of this report}}\\
Restricted logarithm-matrix coefficient principle $\Longleftrightarrow$
Conjecture~\ref{conj:AE}; rational anisotropy of a counterexample & \statuslean &
\lean{RestrictedMatrixCoefficient}; Section~\ref{mc:restricted} \\
Sparse irrational-power-curve zero bound, nonzero integral evaluation determinant, and
the exact $R+1$-term positive-node interpolant & \statuslean &
\lean{SparsePowerCurve}; Section~\ref{mc:sparse} \\
Exact rational rectangular orbit-collision count, irrational grid cardinality, and the
one-monomial-saving rational-orbit square-grid certificate & \statuslean &
\lean{RationalOrbitGrid}; Section~\ref{fc:support-spectrum} \\
Formal prime-symbol determinant and its specialization; formal vanishing exactly in the
integral branch; sparse Structural Rank equivalence & \statuslean &
\lean{FormalPrimeDeterminant}; Section~\ref{mc:formal-prime} \\
BK degree threshold $R-1$ and the factor-four comparison on the $2N$ grid & \statusnew &
Section~\ref{mc:bk} \\
Boundary-dilation defect factorization, unique degree-$K-1$ residuals, and the degree-$2K-2$
commutator & \statuslean & \lean{BoundaryDilationDefect}; Section~\ref{mc:boundary} \\
Unbounded $o((\log H)^2)$ and synchronized compact $o(\log H)$ point-counting targets &
\statusnew & Section~\ref{mc:counting-monodromy} \\
Transcendence of $e^{2\pi iq\beta}$ for $q\in\Q^\times$, deduced from classical six
exponentials & \statusnew & Section~\ref{mc:counting-monodromy}; classical input
\statusknown \\
Exponential-polynomial zero count; generalized Vandermonde nonvanishing at ordered
nodes & \statuslean & \lean{ExponentialPolynomialZeros} \\
Generalized Vandermonde positivity (ordered chamber) & \statuslean &
\lean{OrderedChamberPositivity} \\
Arbitrary-node determinant integrality and divided-difference repulsion bound &
\statuslean & \lean{ArbitraryNodeDeterminant} \\
Universal nonnegative-node superfactorial interpolation divisor, exact consecutive-node witness, and
strengthened candidate repulsion/packing & \statuslean &
\lean{SuperfactorialDeterminant}; final repulsion endpoints retain the named
mean-value hypothesis \\
Normalized geometric signed-progression determinant, primitive quotient coprimality, and exact
denominator-tax decomposition with bounded remainder & \statuslean &
\lean{SignedProgressionDeterminant}; Section~\ref{fc:signed-progressions} \\
Normalized Fej\'er kernel, exact finite weighted Fourier score, and exact integer-phase hit
bound & \statuslean & \lean{FejerKernelCertificate}; Section~\ref{pc:fejer} \\
Uniform metric near-integrality and Haar random-phase extinction & \statusnew &
Section~\ref{pc:metric-phase}; finite Fej\'er bridge separately kernel-checked \\
Exact denominator-exponent drift, common centered defect, and structural-prime $p$-adic escape
& \statuslean & \lean{AdelicDrift}; Section~\ref{r78:adelic-drift} \\
Sturmian carry telescope, exact prefix mean/variance, joint-carry probability, and covariance
kernel & \statuslean & \lean{SturmianProbability}; Section~\ref{r78:sturmian} \\
Dense denominator clearing, finite-difference approximation floor, and all-degree normalized
tax obstruction & \statuslean & \lean{PolynomialApproximationTax};
Section~\ref{r78:approximation-tax} \\
Rational homogeneous contact nonvanishing and first-order Pad\'e transversality & \statuslean &
\lean{RationalContactRigidity}; Section~\ref{r78:contact} \\
Carry/floor and denominator formal-series identities, including the exact common defect &
\statuslean & \lean{SturmianGeneratingSeries}; Section~\ref{r9:natural-boundary} \\
De-integralized common-zero lattice and canonical divisor; compact $\Theta(\log H)$ height
dichotomy & \statusnew & Section~\ref{r9:common-divisor}, Section~\ref{r9:compact-height} \\
P\'olya--Carlson natural boundary for the carry, floor, and reduced-denominator series &
\statusknown & Section~\ref{r9:natural-boundary}; exact series algebra separately
kernel-checked \\
Fixed $2\times2$ integral matrix: integral real power iff both original outputs are integral &
\statuslean & \lean{IntegralPowerMatrix}; Section~\ref{r10:fixed-matrix} \\
Affine-exponent and rational-approximation alternant factorizations & \statuslean &
\lean{CandidateAlternant}; Section~\ref{r10:one-half} \\
Cross-Loewner positive integer certificate and optimized two-coset one-half barrier &
\statusnew & Section~\ref{r10:loewner}, Section~\ref{r10:one-half} \\
Fermat-quotient exact identity, power law, first-log determinant interface, and vanishing on
$(2^n,3^n)$ & \statuslean & \lean{FermatQuotientDeterminant}; Section~\ref{r11:fermat} \\
Structural-prime Iwasawa transversality, exact adelic wedge, conditional cascade depth, and
Kummer third-order diagnostic & \statusnew & Section~\ref{sec:report11}; classical Iwasawa
kernel where stated \\
Exact consecutive-logarithmic-mean enclosure & \statuslean &
\lean{ConsecutiveLogMeanBounds}; Section~\ref{r12:switch-logmean} \\
Prime-switch antichain, microscopic resonance window, and $2$--$3$ convergent sieve &
\statusnew & Section~\ref{sec:report12}; rational-output switch family only \\
Exact exclusion of simultaneous $2$--$3$ switches through $b=10000$ & \statuscomp &
Section~\ref{r12:two-three-sieve}; Python big-integer/rational verifier \\
Exact two-coordinate grid displacement identities and rank-$K$ boundary defects &
\statuslean & \lean{BoundaryDisplacement}; rank under $K>0$, nonsingularity, and nonzero
multipliers; Section~\ref{r13:boundary-displacement} \\
Exact rational bilinear minimum, projective arithmetic density, and integral Hadamard-power
no-go & \statusnew & Section~\ref{r13:rational-conditioning},
Section~\ref{r13:projective-density} \\
Determinantal-divisor compound sandwich and triangular method ceiling & \statusnew &
Section~\ref{r13:compounds}; ceiling applies to the stated termwise/Frobenius bounds \\
Two-variable finite-difference divisor for every compound of the square-grid matrix &
\statusnew & Section~\ref{r13:compounds}; top gain remains lower order \\
Boundary transfer determinant and Kraus--Loewner positive integral/denominator tax &
\statusnew & Section~\ref{r13:boundary-displacement}, Section~\ref{r13:kraus} \\
Integral dilation-difference quotient and its mod-$2$ derivative reduction & \statuslean &
\lean{DilationDifferenceQuotient}; Section~\ref{r14:scaled-resultants} \\
Structural-prime denominator lattice, favorable-branch residual scaling, root depths, and
odd-degree dilation-resultant valuations & \statusnew & Section~\ref{sec:report14}; only the
two displayed residue branches \\
Exact unordered-pair decoding from the ordinary and synchronized $p$-power sums & \statuslean &
\lean{SynchronizedSumDecoder}; Section~\ref{r15:decoders} \\
Rational synchronized-sum decoder degree bound $L\le2d+1$ & \statuslean &
\lean{RationalSumDecoder}; Section~\ref{r15:decoders} \\
Fixed-rank sum entropy, algebraic decoder degree growth, and finite Fourier parity--residue
separation & \statusnew & Section~\ref{sec:report15}; unit-equation input classical where
stated \\
Generic algebraic-base spectrum for an arbitrary independent positive algebraic pair &
\statuslean & \lean{AlgebraicBaseSpectrum}; Section~\ref{r16:spectrum} \\
Finite rectangular mixed-Mahler inclusion--exclusion identity & \statuslean &
\lean{MixedMahlerFinite}; Section~\ref{r16:mahler}; no infinite convergence claim \\
Entire gamma quotient, corrected mixed-Mahler double-pole resonance, and transfer-spectrum
diagnosis & \statusnew & Section~\ref{sec:report16}; holonomic, P\'olya--Carlson, Barnes, and
$q$-digamma inputs classical where stated \\
Exact determinant-one Smith reduction of the geometric Vandermonde, $q$-Pascal identity,
factor formula, divisibility order, and positivity & \statuslean &
\lean{GeometricVandermondeSmith}; Section~\ref{r17:modular-smith} \\
Exterior-minor probability identity, weighted Smith-tail inequalities, and multilevel modular
divisor & \statusnew & Section~\ref{r17:smith-profile}, Section~\ref{r17:modular-smith} \\
Rational cross-Loewner Andreief--Selberg determinant, quadratic singular hierarchy, and
weighted compound certificate & \statusnew & Section~\ref{r17:loewner}; reduced denominator
or Smith gain remains open \\
Infinitely divisible $3$-smooth Gram no-go and exact semigroup entropy saturation &
\statusnew & Section~\ref{r17:no-go}; no closing arithmetic inequality \\
Finite ordinary two-scale support inclusion--exclusion identity & \statuslean &
\lean{TwoScaleSupportFinite}; Section~\ref{r18:support}; no infinite convergence claim \\
Spectral-zeta/logarithmic-Weyl reformulations and complete rank-two heat expansion &
\statusnew & Section~\ref{r18:spectral}, Section~\ref{r18:heat}; exact counting core already
kernel-checked, Mellin--Baker expansion paper-level \\
Sharp fixed-window polynomial node-budget barrier and capacity constant & \statusnew &
Section~\ref{r18:approximation}; no claim about rational or growing-window approximants \\
Exact finite geometric Newton expansion, interpolation, omitted tail, and first extrapolation
product & \statuslean & \lean{GeometricNewtonInterpolation}; Section~\ref{r19:qnewton} \\
Arithmetic unit locking, optimal-block no-go, dyadic certificates, and cofactor minimax unit
barrier & \statusnew & Section~\ref{sec:report19}; no improvement to the accepted finite scan \\
Rational scalar/matrix-entry first-jet recovery and base-two integral-exponent closure &
\statuslean & \lean{JordanFirstJet}; Section~\ref{r20:affine-jet}; matrix real-power formula
remains paper-level \\
Operator local-universality no-go and global/finite/affine/trace/Hankel closing criteria &
\statusnew & Section~\ref{sec:report20}; criteria are not consequences of the two outputs \\
Arbitrary-node repulsion theorem and its packing corollary & \statuslean &
\lean{ArbitraryNodeRepulsion}, conditional on its named divided-difference mean-value
hypothesis \\
Newton divided-difference factorization of the interpolation determinant & \statuslean &
\lean{ArbitraryNodeRepulsion}, unconditional \\
Semigroup generalized-Vandermonde hierarchy (integrality half) & \statuslean &
\lean{MixedExponentSemigroup}, \lean{SemigroupDeterminant} \\
Falling-factorial coefficient bound; dyadic-strip divided-difference bound & \statuslean &
\lean{FallingFactorialBound} \\
Semigroup gap theorem in quantitative form & \statuslean & \lean{SemigroupGapBound} \\
Leibniz determinant bound from separated row/column bounds & \statuslean &
\lean{SemigroupGapBound}, unconditional \\
Window-span counting bridge for mixed exponents & \statuslean &
\lean{SemigroupWeightBound} \\
Explicit dyadic $O((\log X)^2)$ bound from that hierarchy & \statuslean &
\lean{SemigroupGapBound}, conditional on the named hypothesis
\lean{NewtonFactorization} \\
One-logarithm deficit analysis & \statusnew & Section~\ref{sec:deficit} \\
Min-plus ceiling: term-wise $p$-adic content never exceeds a correct archimedean bound &
\statuslean & \lean{MinPlusCeiling} \\
Exact termwise-divisor ceilings ($4/5$, $7/10$, $3/10$, $\pi/4$, $0.8710$) & \statusnew &
Section~\ref{sec:minplus} \\
A unique tropical minimizer forbids cancellation & \statuslean & \lean{MinPlusCeiling} \\
Tropical quotient factorization, all-depth divisibility, and tied-fiber first-layer
criterion (no cancellation lower bound) & \statuslean & \lean{TropicalInitialForm} \\
Block-preserving tied sums factor into block determinants and consecutive-power
Vandermondes; Beatty minimizers, floor gap, quotient divisibility, and exact expansion &
\statuslean & \lean{BlockFiberFactorization}, \lean{TropicalBlockApplication},
\lean{BeattyTropicalFiber} \\
Exact all-layer row-partition expansion and cutoff-$K$ support truncation
(no bound on surviving cancellation) & \statuslean & \lean{RowBlockLaplaceExpansion} \\
Prefix-crossing lower bound and stars-and-bars count of row partitions below a Beatty
excess cutoff (no valuation/cancellation bound) & \statuslean &
\lean{AssignmentExcessGeometry} \\
Exact Vandermonde collision energy, elementary power-difference valuation, and finite
variable-block cube budget & \statuslean & \lean{BlockVandermondeValuationBudget} \\
Whole-Vandermonde dyadic and triadic unit fixed divisors under monic column changes
(no bound on the residual cascade) & \statuslean & \lean{AllLayerUnitBasis},
\lean{UnitOrderingGain} \\
Uniform scaled-block minor divisors and additive whole-determinant gain
$p^{\tau+U}\mid\det$ (no residual-cancellation bound) & \statuslean &
\lean{RowBlockUnitDivisibility} \\
Column-permuted raw mixed-matrix bridge and additive dyadic/triadic divisors under explicit
finite fiber certificates & \statuslean & \lean{MixedDeterminantUnitBridge} \\
Genuine global first-$N$ mixed prefix, automatic consecutive fibers, exact dyadic/triadic
cutoff expansions, and structural-output divisors (no residual-cancellation upper bound) &
\statuslean & \lean{GlobalMixedPrefixCascade}; strictly decreasing depths, with $A$ odd for
the dyadic structural divisor and $3\nmid M$ for the triadic one \\
Exact synthetic cascade computations, including counterexamples to first-layer equality and
minimum-effective-layer equality & \statusexactcomp & Section~\ref{sec:cascade-experiments} \\
Compact synchronized-orbit reformulation & \statuslean & \lean{CompactOrbit} \\
Block-growth reformulations (bounded, $o(T)$, $\liminf$, $\equiv1$) & \statuslean &
\lean{BlockGrowthReformulations} \\
Rational-function rigidity, multiplication and quotient criteria, closure
reformulations & \statusnew & Section~\ref{sec:rigidity} \\
Perfect-power content, matched affine decompositions, primitive-output statistics &
\statusnew & Section~\ref{sec:powers} \\
Candidate gcd/lcm lattice; prime-support rigidity of the whole branch & \statuslean &
\lean{CandidateLattice} \\
Ambient root saturation index $d$ & \statusnew & Section~\ref{sec:lattice} \\
Operator formulation on the prime-logarithm space & \statuslean &
\lean{PrimeLogSpaceOperator} \\
Every complex solution is real (principal branch) & \statuslean & \lean{ComplexSolutions} \\
Gaussian-integer version is equivalent to the conjecture & \statusnew & Section~\ref{sec:variants} \\
Multivalued branches: the log-product wall widened by $\log(-1)$ & \statusnew & Section~\ref{sec:variants} \\
Exact cost of the rational antecedent (denominator-support lemma) & \statusnew & Section~\ref{sec:variants} \\
\bottomrule
\end{longtable}
}

\section{Formalization inventory}
\label{app:inventory}

All modules live under
\lean{Analysis/ExponentialIdentities/Lean/ExponentialIdentities/}, in namespace
\lean{LeanProofs}.  Every theorem module listed below is reachable from the aggregate
import surface; \lean{AuditNewModules} is built separately as an assumption-audit target.
The merged tree builds with no errors.  None of the theorem modules contains
\lean{sorry}, an extra axiom, or \lean{native_decide}; certificate tables are replayed by
the kernel.

\subsection{Kernel-accepted modules}

{\footnotesize
\begin{longtable}{>{\raggedright\arraybackslash}p{0.385\textwidth}
                  >{\raggedright\arraybackslash}p{0.535\textwidth}}
\toprule
Module & Contents \\
\midrule
\endfirsthead
\toprule
Module & Contents \\
\midrule
\endhead
\lean{IntegerExponent} (+3 modules) & three-base theorem, interpolation determinants \\
\lean{TwoBaseIntegerExponent} & conjecture statement, candidates, gaps, $2^x<10$ \\
\lean{.../FiniteCheck}, \lean{FiniteCheck256} & zero-free regions $2^x<100$, $2^x<256$ \\
\lean{.../FiniteCheck4096} & zero-free region $2^x<4096$ (kernel \lean{decide} table) \\
\lean{.../ZeroDensity} & AP obstruction, Roth-number counting, $o(N)$ \\
\lean{.../NonintegerDensity} & nonintegral candidate refinement, density zero \\
\lean{.../MultiplicativeRank} & rank $\le2$, $\thetaq$ irrational, $(\clog_2N)^2$ \\
\lean{.../CandidateStructure} & cross-valuation identity, conditional box growth \\
\lean{.../OddCore} & common odd core, $\clog_2N\,\clog_3N$, dyadic blocks \\
\lean{.../PrimitiveOddCore} & gcd-normalized canonical core, unique coordinates \\
\lean{.../LeastSolution} & unique positive least solution, exact affine primitivity \\
\lean{.../ThreeDenominatorNormalization} & reduced denominator supported at $3$ \\
\lean{.../RationalPowerIndex} & least rational-power index, exact solution monoid \\
\lean{.../PrimitiveGenerator} & canonical primitive generator and unique coordinates \\
\lean{.../ExactCounting} & exact unit/dyadic block counts and cumulative finite sum \\
\lean{.../QuadraticAsymptotic} & sharp quadratic limit; subquadratic-growth equivalence \\
\lean{.../BlockWeyl} & exact-block geometric sums and vanishing nonzero Fourier modes \\
\lean{.../BlockEquidistribution} & finite trigonometric-polynomial block averages \\
\lean{.../OutputNormalization} & simultaneous primitive outputs, gcd one, size dichotomy \\
\lean{.../RationalThirdBase} & denominator divisibility and determinant margin \\
\lean{.../RadicalReformulation} & localized-radical equivalences, base-swapped form \\
\lean{.../PrimitiveRadical} & canonical primitive-power reduction of radical target \\
\lean{.../RadicalDegree} & irreducible binomial, exact minimal polynomial and degree \\
\lean{.../SharperProgressions} & sharp three-candidate curvature obstruction \\
\lean{.../ThirdDifference} & exact third difference, four-candidate obstruction \\
\lean{.../HigherDifference} & uniform arbitrary-order forward-difference obstruction \\
\lean{.../PolynomialApproximationTax} & exact dense denominator clearing, integer resolution,
finite-difference approximation floor, and the all-degree normalized tax obstruction \\
\lean{.../SharpProgression} & effective rational three-term criterion $h^{400}\le n^{83}$ \\
\lean{.../ArithmeticRigidity} & descent, $r+8k\le x$ \\
\lean{.../ThreeSmooth} & $3$-smooth classification of further bases \\
\lean{.../ShrinkingTarget} & fractional-part shrinking targets, Diophantine lower bounds \\
\lean{.../DenominatorGrowth} & convergent-denominator growth cap $q_{n+1}<2m^{q_n}$ \\
\lean{.../Transcendence} & Gelfond--Schneider applications: solutions and $\log_23$ \\
\lean{.../FractionalPartDensity} & density dichotomy, gap equivalences, infinitude \\
\lean{.../Localization} & candidate-form localization equivalence \\
\lean{.../FourExponentialsReduction} & four exponentials $\Rightarrow$ the conjecture \\
\lean{GelfondSchneider/...} & Gelfond--Schneider theorem (port of Karatarakis) \\
\midrule
\multicolumn{2}{@{}l}{\emph{Rational and algebraic auxiliary bases}}\\
\lean{.../RationalSixExponentials} & six exponentials for three positive rational bases and
rational outputs \\
\lean{.../RationalBaseThird} & specialization to bases $2,3$; exact $2,3$-unit
classification of positive rational third bases \\
\lean{.../AlgebraicOutputBridge} & a number field, an embedding, and one natural denominator
clearing an algebraic value to an algebraic integer \\
\lean{.../FiniteAlgebraicOutputBridge} & one common number field and one denominator for a
finite family of algebraic values \\
\lean{.../AlgebraicIntegerDeterminant} & field-norm lower bound for a nonzero
algebraic-integer determinant; house and factorial bounds on the $n^{11}$ scale \\
\lean{.../AlgebraicSixExponentials} & field-norm six-exponentials determinant for three
positive algebraic real bases and three algebraic outputs \\
\lean{.../AlgebraicThirdBase} & three-base determinant for algebraic output at a natural
third base, first over a number field and then by clearing denominators \\
\lean{.../AlgebraicRationalBase} & algebraic outputs at positive rational bases; the
algebraic locus is exactly the $2,3$-units unless the exponent is integral \\
\lean{.../AlgebraicBaseUnitPower} & positive $2,3$-unit powers and the multiplicative rank
of an algebraic auxiliary base \\
\lean{.../AlgebraicBaseSpectrum} & arbitrary independent positive algebraic pair; exact
positive-power/divisible-hull classification of positive algebraic bases with algebraic
$x$th output \\
\lean{.../AlgebraicConsequences} & pointwise algebraic natural-base classification;
iterated-output transcendence; the $5^x$, $6^{x^2}$ and both-iterate equivalences \\
\lean{.../AlgebraicOutputLocus} & exact algebraic output monomial locus (coordinate face or
origin); external $p$-adic line for output ratios; coordinate-face and rank-two
equivalences \\
\lean{.../AlgebraicRankThreeEquivalence} & rank three in the algebraic real-power output
locus forces integrality; the resulting equivalence \\
\lean{.../AlgebraicExponentLocus} & the common algebraic-output exponent plane
$\Q+\Q x$, and its natural-power equivalence \\
\lean{.../RationalFunctionAlgebraicLocus} & exact rational-function affine identity at a
counterexample; nonaffine rational-function and degree-$\ge2$ polynomial detectors; canonical
rational affine plane versus nonnegative integral cone \\
\midrule
\multicolumn{2}{@{}l}{\emph{Canonical radical arithmetic}}\\
\lean{.../CanonicalRadicalValuation} & odd-prime valuation gcd; divisors of the gcd are
honest outer-power exponents \\
\lean{.../CanonicalRadicalSharpness} & exact criterion for a reduced $c/3^a$ to be a
rational outer power; the triple gcd resists both pairwise relaxations \\
\lean{.../CanonicalRadicalFieldNorm} & norm and trace of the pure-radical generator; the
base-swapped relation blocks monomial injectivity \\
\lean{.../CanonicalRadicalThirdOutput} & algebraic integrality of the normalized radical is
equivalent to a trivial reduced denominator \\
\lean{.../CanonicalRadicalAlgebraicOutput} & algebraicity of the radical's $\beta$-th power
is equivalent to $2,3$-smoothness of the normalized numerator \\
\lean{.../CanonicalRadicalDivisibleHull} & an external numerator prime excludes the whole
positive $2,3$-unit divisible hull \\
\lean{.../CanonicalRadicalMixedOutput} & genuinely mixed radical monomials keep an external
prime, so all their outputs are transcendental \\
\lean{.../SimultaneousRadicalDegrees} & the displayed exponents are the least
rational-power indices; exact radical degrees and the product degree $de$ \\
\lean{.../CoreIndependence} & the two primitive output cores are never exactly $3$ and $2$;
exact rationality classification of the simultaneous cores \\
\midrule
\multicolumn{2}{@{}l}{\emph{Smoothness and log-products}}\\
\lean{.../SmoothOutputs} & both outputs $3$-smooth forces an integral exponent; sharp form
supplying a partner with a prime factor $\ge5$ \\
\lean{.../LogProductIndependence} & prime-log-product independence as a sufficient condition
distinct from four exponentials, in both the integer-output and rational-output forms \\
\midrule
\multicolumn{2}{@{}l}{\emph{Determinant hierarchy and reformulations}}\\
\lean{.../ArbitraryNodeDeterminant} & interpolation determinant at arbitrary ordered nodes;
integrality and the divided-difference repulsion bound \\
\lean{.../SuperfactorialDeterminant} & universal nonnegative-node superfactorial divisor and exact
consecutive-node witness; strengthened candidate determinant, repulsion, and packing bounds \\
\lean{.../SignedProgressionDeterminant} & geometric Vandermonde factorization; primitive
invariant quotient and coprimality; exact finite logarithmic denominator-tax identity \\
\lean{.../GeometricVandermondeSmith} & integral $q$-Pascal and Newton-basis matrices;
determinant-one two-sided reduction to the exact positive Smith diagonal for $q\ge2$ \\
\lean{.../GeometricNewtonInterpolation} & exact finite $q$-Newton power expansion;
interpolation through the truncation degree, omitted-tail identity, and first extrapolation
product over arbitrary integer parameters \\
\lean{.../FejerKernelCertificate} & normalized finite Fej\'er kernel; exact weighted Fourier
score; exact phase and integer-phase hit-count certificates \\
\lean{.../ExponentialPolynomialZeros} & zero-count theorem for real exponential polynomials
and nonvanishing of the generalized Vandermonde determinant \\
\lean{.../SparsePowerCurve} & sparse zero bound on an irrational power curve; nonzero
integral evaluation determinant; exact-support positive-node interpolant \\
\lean{.../RationalOrbitGrid} & exact rational rectangle collision count; irrational grid
cardinality and sparse lower bound; canonical one-monomial-saving rational-orbit certificate \\
\lean{.../RationalContactRigidity} & rational polynomial evaluation at a transcendental
direction; homogeneous leading-contact nonvanishing and Pad\'e transversality \\
\lean{.../RestrictedMatrixCoefficient} & concrete logarithm matrix, rank-one coefficient
factorization, rational anisotropy, and exact equivalence with the conjecture \\
\lean{.../IntegralPowerMatrix} & explicit real powers of a fixed split integral matrix and
their exact entry-integrality/conjecture equivalence \\
\lean{.../JordanFirstJet} & rationality recovery from a scalar value and its nonzero first
jet; rational matrix-entry extraction; base-two integral-exponent closure conditional on the
formal first-jet matrix \\
\lean{.../CandidateAlternant} & affine real-exponent Vandermonde factorization and the exact
rational-approximation candidate alternant \\
\lean{.../SynchronizedSumDecoder} & strict monotonicity on fixed-sum pairs and exact
unordered-pair recovery from ordinary and synchronized real-power sums \\
\lean{.../RationalSumDecoder} & exact $L\le2d+1$ bound for a degree-$d$ rational decoder on
the synchronized slice $(2^n+1,3^n+1)$ \\
\lean{.../FormalPrimeDeterminant} & multivariate prime-symbol determinant, specialization at
prime logarithms, exact formal-zero classification, and Structural Rank equivalence \\
\lean{.../FermatQuotientDeterminant} & exact Fermat quotient and power law modulo $p$; first
local-log determinant interface and integral-power determinant vanishing \\
\lean{.../ConsecutiveLogMeanBounds} & exact two-sided enclosure for the reciprocal
consecutive logarithmic mean, with an explicit $1/(24H^2)$ remainder \\
\lean{.../BoundaryDisplacement} & exact two-coordinate evaluation-matrix displacement
identities, terminal-boundary projector ranks, and rank-$K$ displacement consequences under
explicit $K>0$, nonsingularity, and nonzero-multiplier hypotheses \\
\lean{.../BoundaryDilationDefect} & rectangular node products, dilation-defect divisibility,
unique square-grid residuals, and the degree-$2K-2$ boundary commutator \\
\lean{.../DilationDifferenceQuotient} & integral factorization of $P(cX)-P(X)$; dyadic
$2X$ quotient and its formal-derivative reduction modulo two \\
\lean{.../OrderedChamberPositivity} & positivity of generalized real-power Vandermonde
determinants and of the semigroup determinant on ordered chambers \\
\lean{.../MixedExponentSemigroup} & pair-indexed exponent $a+b\thetaq$, injectivity, and
integrality of $m^{a+b\thetaq}$ \\
\lean{.../MixedMahlerFinite} & exact finite rectangular inclusion--exclusion identity behind
the mixed $(2,3)$-Mahler equation; no convergence or natural-boundary assertion \\
\lean{.../TwoScaleSupportFinite} & exact finite rectangular inclusion--exclusion identity
behind the ordinary two-scale support equation; no convergence assertion \\
\lean{.../SemigroupDeterminant} & integrality half of the semigroup generalized-Vandermonde
hierarchy, with $\lvert\det\rvert\ge1$ \\
\lean{.../FallingFactorialBound} & falling-factorial coefficient bound and dyadic-strip
divided-difference bound \\
\lean{.../SemigroupWeightBound} & mixed-exponent family construction and the
window-summation counting bridge \\
\lean{.../BlockGrowthReformulations} & boundedness, $o(T)$, $\liminf$, and $\equiv1$
reformulations of the conjecture \\
\lean{.../CompactOrbit} & compact synchronized-orbit reformulation with shared floor
exponent, in both directions \\
\lean{.../AdelicDrift} & exact integer denominator exponents, common centered defect,
structural-prime normalized orbit norms, and $p$-adic escape \\
\lean{.../SturmianProbability} & carry/floor telescope, exact prefix law, mean, variance,
joint-carry intervals, probability, and covariance kernel \\
\lean{.../SturmianGeneratingSeries} & exact carry/floor and denominator formal power-series
identities, twice-cleared defect, and cross-base cancellation \\
\lean{.../ArbitraryNodeRepulsion} & divided differences, the Newton factorization of the
interpolation determinant, and the repulsion and packing bounds conditional on the named
mean-value hypothesis \\
\lean{.../SemigroupGapBound} & Leibniz determinant bound, semigroup span inequality, and the
explicit dyadic $10^4(\log X)^2$ block bound conditional on the named Newton factorization \\
\lean{.../TropicalInitialForm} & common-power factorization of a Leibniz determinant;
all-depth quotient divisibility and the modulo-$p$ tied-fiber criterion, without a
cancellation lower bound \\
\lean{.../BlockFiberFactorization} & block-preserving signed Leibniz sums as products of
block determinants, including variable-size and consecutive-power Vandermonde forms \\
\lean{.../TropicalBlockApplication} & first-extra-prime divisibility iff the product of
block determinants, or its consecutive-power Vandermonde specialization, vanishes modulo
$p$; application hypotheses remain explicit \\
\lean{.../BeattyTropicalFiber} & exact repeated-weight minimizers for rank-one assignment
costs; Beatty floor gap; normalized quotient divisibility; exact block-determinant expansion \\
\lean{.../RowBlockLaplaceExpansion} & exact grouping of all assignment layers by row
partitions; normalized and cutoff-$K$ expansions into signed products of square unit minors \\
\lean{.../AssignmentExcessGeometry} & exact prefix-crossing energy; Beatty excess lower
bound; stars-and-bars binomial count for row partitions below any fixed cutoff \\
\lean{.../BlockVandermondeValuationBudget} & exact collision-depth formula for finite
Vandermondes; elementary difference-of-powers valuation and variable-block cubic budgets \\
\lean{.../AllLayerUnitBasis} & complete-determinant divisors from monic unitriangular or
determinant-one column changes; strengthened positive dyadic-odd and triadic-unit $p$-ordering weights, without a residual-cascade
bound \\
\lean{.../UnitOrderingGain} & exact unit block-weight identities and finite square/cube
aggregate bounds; no triangular-lattice asymptotic or residual-cascade estimate \\
\lean{.../RowBlockUnitDivisibility} & row-scaled consecutive-power matrices; a common
$p$-ordering divisor for every assignment-minor product; determinant-one fiberwise basis
change; additive divisor $p^{\tau+U}$ for the full rank-one power determinant \\
\lean{.../MixedDeterminantUnitBridge} & finite sorted-box mixed-exponent prefix; exact
column-permuted identification with the rank-one scaled model; dyadic and triadic additive
divisors under explicit block-fiber certificates \\
\lean{.../GlobalMixedPrefixCascade} & global initial-segment certificate, automatic
lexicographic/colexicographic fiber equivalences, exact truncated cascade identities, and
structural-output determinant divisors under the stated decreasing-depth and residue
hypotheses \\
\midrule
\multicolumn{2}{@{}l}{\emph{Second iterates, towers, and the candidate lattice}}\\
\lean{.../CandidateLattice} & coordinatewise gcd and lcm for $2^s w^t$; unconditional closure
of the candidate set under gcd and lcm; prime-support rigidity of the branch \\
\lean{.../KernelDichotomy} & multiplicative dependence of $M,A,2,3$ is equivalent to a
rational linear relation among $1,\thetaq,x,x\thetaq$ and to the M\"obius orbit \\
\lean{.../SecondIterateKernel} & the second-iterate kernel as a subgroup of $\Z^2$; no
same-sign vectors, rank at most one, smooth coordinate cases, and solution-independence \\
\lean{.../TowerRigidity} & base-two second iterate as a natural number; the depth-two
algebraic dichotomy; the depth-three transcendence and the cubic-tower equivalence \\
\lean{.../AlgebraicTowerTransition} & exact transition locus
$2^y,2^{xy}\in\Qbar\iff y\in\Q+\Q\thetaq$ under a counterexample; global uniqueness of a
positive transition depth, the all-depth three-level exclusion, and the fixed-distinct-depth
AE equivalence \\
\lean{.../CrossBaseTowerTransition} & exact scaled base-three transition locus; cross-base
mutual exclusion, combined uniqueness across both principal towers, and higher smooth-direction
exclusions in the dependent-output branch \\
\lean{.../SymmetricTowerConstants} & $E_3=3^{\thetaq}$ and $E_2=2^{1/\thetaq}$, their
conditional rational powers, and their mutual exclusion \\
\lean{.../PrimeLogSpaceOperator} & the $\thetaq$-pullback of the prime-logarithm span is
exactly $\Q\log2$ if and only if the rational-output conjecture holds \\
\lean{.../ComplexSolutions} & on the principal branch every complex solution is real, by
incommensurability of the two argument gratings \\
\lean{.../MinPlusCeiling} & the $\{2,3\}$-part of an integer divides it, hence term-wise
$p$-adic content can never exceed a correct archimedean bound \\
\lean{.../AuditNewModules} & \texttt{\#print axioms} audit of the headline algebraic-locus,
rational-function, tower, radical-degree, mixed-output, determinant, tropical, and block-fiber
theorems \\
\bottomrule
\end{longtable}
}

\subsection{Former theorem drafts and the remaining finite-certificate draft}
\label{st:no-drafts}

No theorem module in the nine-item list from earlier revisions remains pending.  Those
revisions listed the following Lean sources as not yet accepted:
\begin{center}
\lean{ArbitraryNodeDeterminant}, \lean{ExponentialPolynomialZeros},
\lean{MixedExponentSemigroup},\\
\lean{SemigroupDeterminant}, \lean{FallingFactorialBound}, \lean{SemigroupWeightBound},\\
\lean{BlockGrowthReformulations}, \lean{CompactOrbit}, \lean{LogProductIndependence}.
\end{center}
All nine have since been accepted by the kernel, placed on the aggregate import surface, and
moved into the table above.  Statements that those modules actually contain are therefore
\statuslean{}.

Separately, the attempted $2^x<16384$ finite-certificate assembly described in
Section~\ref{sec:finitecheck} remains a genuine \statusdraft{}: its individual arithmetic
rows replay, but its covering theorem does not compile end to end.  It is not on the aggregate
import surface and raises no kernel-verified bound.  The accepted finite frontier remains
$2^x<4096$.

A second round has since added eight further modules: \lean{CandidateLattice},
\lean{KernelDichotomy}, \lean{SecondIterateKernel}, \lean{TowerRigidity},
\lean{SymmetricTowerConstants}, \lean{PrimeLogSpaceOperator}, \lean{ArbitraryNodeRepulsion}
and \lean{SemigroupGapBound}, all on the aggregate import surface and all audited with
\texttt{\#print axioms}.

The current wave adds the following modules:
\begin{center}
\lean{AlgebraicExponentLocus}, \lean{AlgebraicOutputLocus},\\
\lean{RationalFunctionAlgebraicLocus}, \lean{AlgebraicTowerTransition},\\
\lean{CrossBaseTowerTransition}, \lean{BlockVandermondeValuationBudget},
\lean{AllLayerUnitBasis}, \lean{UnitOrderingGain},\\
\lean{CanonicalRadicalMixedOutput}, \lean{SimultaneousRadicalDegrees},\\
\lean{TropicalInitialForm}, \lean{BlockFiberFactorization},
\lean{TropicalBlockApplication}, \lean{BeattyTropicalFiber},\\
\lean{RowBlockLaplaceExpansion}, \lean{AssignmentExcessGeometry},
\lean{RowBlockUnitDivisibility}, \lean{MixedDeterminantUnitBridge},\\
\lean{RestrictedMatrixCoefficient}, \lean{FormalPrimeDeterminant},
\lean{SparsePowerCurve}, \lean{BoundaryDilationDefect},
\lean{SuperfactorialDeterminant}, \lean{RationalOrbitGrid},
\lean{SignedProgressionDeterminant}, \lean{FejerKernelCertificate},\\
\lean{AdelicDrift}, \lean{SturmianProbability},
\lean{SturmianGeneratingSeries}, \lean{PolynomialApproximationTax},\\
\lean{RationalContactRigidity}, \lean{IntegralPowerMatrix},
\lean{CandidateAlternant}, \lean{FermatQuotientDeterminant},
\lean{ConsecutiveLogMeanBounds}, \lean{BoundaryDisplacement},\\
\lean{GlobalMixedPrefixCascade}, \lean{DilationDifferenceQuotient},
\lean{SynchronizedSumDecoder}, \lean{RationalSumDecoder},\\
\lean{AlgebraicBaseSpectrum}, \lean{MixedMahlerFinite},
\lean{GeometricVandermondeSmith}, \lean{TwoScaleSupportFinite},
\lean{GeometricNewtonInterpolation}, \lean{JordanFirstJet}.
\end{center}
Their headline APIs are likewise on the aggregate surface
and covered by the axiom-audit module.

Two cautions survive both promotions, and the status matrix reflects them.  First, a module
being accepted certifies exactly the theorems it states.  The analytic halves of the repulsion
and gap theorems are now formal, but each carries its missing analytic input as an explicit,
named \emph{hypothesis} --- \lean{RpowDividedDifferenceMeanValue} and
\lean{NewtonFactorization} --- rather than as an axiom, so those two conclusions are
conditional in a way the reader can see in the theorem statement; everything upstream and
downstream of them is unconditional.  \lean{ExponentialPolynomialZeros} supplies the
nonvanishing engine, while the separate accepted module \lean{OrderedChamberPositivity}
supplies the ordered-chamber sign proved on paper in Appendix~\ref{app:gvproof}.  Second, the
interval scan through $2^{20}$ and the combined
directed-rounding scan through $2^{27}$ are \statuscomp{}, not kernel certificates; the
kernel-verified exponent bound is still
$x\ge12$ from \lean{FiniteCheck4096}.

Each further headline theorem should be added to the aggregate import surface and checked
with \texttt{\#print axioms} before its status is promoted.  Numerical constants in the
tables of Appendix~\ref{app:constants} are explanatory only; the formal statements should use
exact expressions or safe rational constants such as $1/24$ and $10000$.

\section{A symmetric formulation}
\label{app:symmetric}

Everything above has a base-swapped version.  Put $\eta=\log_32=1/\thetaq$.  Failure of the
conjecture is also equivalent to the existence of an integer $v>1$ with $3\nmid v$ and
\[
v^{\eta}\in\Ztwo.
\]
\begin{statusbox}{\statuslean{} ---
\lean{not_alaogluErdosConjecture_iff_exists_threeFree_localized_radical} and
\lean{alaogluErdosConjecture_iff_threeFreeLocalizedRadicalExclusion} in
\lean{RadicalReformulation}.}
\end{statusbox}
For the primitive pair $(M,A)$, factoring $A$ into a power of $3$ times a $3$-coprime part
produces a symmetric primitive core and rational-power index; the two normalizations give the
same least exponent $\beta$, and their compatibility condition is exactly
$\log M\log3=\log A\log2$.  The symmetry is useful in formalization --- some divisibility
arguments are easier on one side --- and confirms that the localized-radical reduction of
Section~\ref{sec:radical} is not an artifact of privileging the base $2$.

\section{Detailed proof of the generalized Vandermonde sign}
\label{app:gvproof}

This appendix expands the generalized Vandermonde positivity lemma of
Section~\ref{sec:arbnode} to make its analytic input fully transparent.

Let $0\le\lambda_0<\cdots<\lambda_r$ and define $f_j(x)=x^{\lambda_j}$ on $(0,\infty)$.  For
$k\le r$,
\[
 f_j^{(k)}(x)=\fall{\lambda_j}{k}x^{\lambda_j-k}.
\]
The $k$th initial Wronskian is
\begin{align*}
 W_k(x)
 &=\det\bigl(f_j^{(i)}(x)\bigr)_{0\le i,j\le k}\\
 &=x^{\sum_{j=0}^{k}\lambda_j-k(k+1)/2}
   \det\bigl(\fall{\lambda_j}{i}\bigr)_{0\le i,j\le k}.
\end{align*}
The polynomial $\fall{z}{i}$ is monic of degree $i$.  Replacing the monomial basis
$1,z,\ldots,z^k$ by the falling-factorial basis is unit lower-triangular.  Therefore
\[
 \det\bigl(\fall{\lambda_j}{i}\bigr)_{i,j=0}^{k}
 =\prod_{0\le i<j\le k}(\lambda_j-\lambda_i)>0.
\]
Hence every initial Wronskian is positive.

One standard characterization of extended complete Chebyshev systems \cite{KarlinStudden} now
implies that for $x_0<\cdots<x_r$,
\[
 \det(f_j(x_i))_{i,j=0}^{r}>0.
\]
To avoid using that characterization as a black box, proceed inductively.  Suppose a
nontrivial combination
\[
 F(x)=\sum_{j=0}^{r}c_jx^{\lambda_j}
\]
has $r+1$ distinct positive zeros.  Put $x=e^t$ and factor the lowest exponential:
\[
 e^{-\lambda_0t}F(e^t)
 =c_0+\sum_{j=1}^{r}c_je^{(\lambda_j-\lambda_0)t}.
\]
Rolle's theorem gives at least $r$ zeros of its derivative, an exponential polynomial with at
most $r$ terms.  Induction gives a contradiction unless all coefficients vanish.  Thus the
evaluation determinant is nonzero.  Its sign cannot change on the connected chamber
$x_0<\cdots<x_r$.  Let the nodes coalesce as $x_i=x+i\varepsilon$ and divide by the positive
ordinary Vandermonde.  As $\varepsilon\downarrow0$, the quotient tends to
$W_r(x)/\prod_{k=0}^{r}k!>0$.  Hence the original determinant is positive everywhere in the
ordered chamber.

\begin{statusbox}{\statusnew{} and \statuslean{} --- paper proof in this report;
kernel-accepted zero-count/nonvanishing in \lean{ExponentialPolynomialZeros} and the positive
ordered-chamber determinant in \lean{OrderedChamberPositivity}.}
\end{statusbox}

\section{Numerical constants}
\label{app:constants}

For reference,
\[
 \thetaq=1.58496250072115618145\ldots,
 \qquad
 \thetaq(\thetaq-1)=0.92714362797110482756\ldots.
\]
The three-point span constant is
\[
 \left(\frac{8}{\thetaq(\thetaq-1)}\right)^{1/3}
 =2.051072395661874\ldots.
\]
For the ordinary-power order $r$, the raw span constant and exponent are
\[
 C_r=\left(\frac{r!}{|\fall{\thetaq}{r}|}\right)^{2/(r(r+1))},
 \qquad
 \alpha_r=\frac{2(r-\thetaq)}{r(r+1)}.
\]
The continuous critical point for $\alpha_r$ solves
\[
 -r^2+2\thetaq r+\thetaq=0,
\]
so
\[
 r_*=\thetaq+\sqrt{\thetaq^2+\thetaq}=3.6090841940\ldots,
\]
confirming that $r=4$ is the optimal integer order.

For the semigroup hierarchy of Section~\ref{sec:semigroup}, the first exponents of
$\LambdaTheta=\{a+b\thetaq:a,b\in\Nzero\}$ and their pair labels are:

\begin{center}
\small
\begin{tabular}{rcl}
\toprule
Index & Pair $(a,b)$ & Exponent $a+b\thetaq$ \\
\midrule
0 & $(0,0)$ & $0$ \\
1 & $(1,0)$ & $1$ \\
2 & $(0,1)$ & $\thetaq$ \\
3 & $(2,0)$ & $2$ \\
4 & $(1,1)$ & $1+\thetaq$ \\
5 & $(3,0)$ & $3$ \\
6 & $(0,2)$ & $2\thetaq$ \\
7 & $(2,1)$ & $2+\thetaq$ \\
8 & $(4,0)$ & $4$ \\
9 & $(1,2)$ & $1+2\thetaq$ \\
10 & $(3,1)$ & $3+\thetaq$ \\
11 & $(0,3)$ & $3\thetaq$ \\
12 & $(5,0)$ & $5$ \\
\bottomrule
\end{tabular}
\end{center}

These decimal values are explanatory.  They are not part of any machine-checked statement,
and formal versions of the corresponding theorems should carry exact expressions instead.

\section{Interval computation: script and outputs}
\label{app:interval}

The script below verifies one range.  Set \lean{START} to $1$ or $2^{19}$ and \lean{LIMIT} to
$2^{19}$ or $2^{20}$ to obtain the two halves reported afterwards.

\begin{lstlisting}[style=pythonstyle,caption={Interval verification},label={ap:lst-interval}]
import math, time, hashlib
import mpmath as mp

START = 1
LIMIT = 1 << 19
DPS = 40

mp.iv.dps = DPS
ln2 = mp.iv.log(2)
ln3 = mp.iv.log(3)
theta_float = math.log(3.0) / math.log(2.0)

checked = 0
failures = []
min_lower = (float("inf"), None, None)
min_upper = (float("inf"), None, None)
manifest = hashlib.sha256()
start_time = time.time()

for m in range(START, LIMIT):
    if (m & (m - 1)) == 0:      # integral exponent: m is a power of two
        continue

    # Only proposes a neighboring integer.  The interval tests below certify it.
    n = math.floor(math.exp(theta_float * math.log(m)))

    lhs = mp.iv.log(m) * ln3
    lower_diff = lhs - mp.iv.log(n) * ln2
    upper_diff = mp.iv.log(n + 1) * ln2 - lhs

    lower_lb = float(lower_diff.a)
    upper_lb = float(upper_diff.a)
    if not (lower_lb > 0.0 and upper_lb > 0.0):
        failures.append((m, n, str(lower_diff), str(upper_diff)))
        break

    min_lower = min(min_lower, (lower_lb, m, n))
    min_upper = min(min_upper, (upper_lb, m, n))
    manifest.update(f"{m}:{n};".encode())
    checked += 1

print("checked_nonpowers=", checked)
print("failures=", len(failures))
print("min_lower_log_margin=", min_lower)
print("min_upper_log_margin=", min_upper)
print("sha256_m_n_manifest=", manifest.hexdigest())
print("elapsed_seconds=", time.time() - start_time)
assert not failures
\end{lstlisting}

The two 40-digit runs produced:

\begin{lstlisting}[style=pythonstyle,numbers=none,caption={Lower half output.}]
LIMIT=524288
DPS=40
checked_nonpowers=524268
failures=0
min_lower_log_margin=(1.4147089634746718e-15, 189427, 231498127)
min_upper_log_margin=(3.0841010256409534e-15, 306982, 497586055)
sha256_m_n_manifest=02b589461fa8c1721bc251575b98b1e3d58001f54cb7faa2cd6c59edb8c819a5
elapsed_seconds=18.505
\end{lstlisting}

\begin{lstlisting}[style=pythonstyle,numbers=none,caption={Upper half output.}]
LIMIT=1048576
DPS=40
checked_nonpowers=524287
failures=0
min_lower_log_margin=(9.490232037370244e-16, 958460, 3023921230)
min_upper_log_margin=(2.3019781003173366e-15, 556061, 1275861767)
sha256_m_n_manifest=fccccf5de8673379d37ff0fc840ab091d9ff7a2352ff64a93091f01c21bf4cb4
elapsed_seconds=18.904
\end{lstlisting}

The two SHA-256 manifests are
\[
\text{lower: }\sha{02b589461fa8c1721bc251575b98b1e3d58001f54cb7faa2cd6c59edb8c819a5},
\]
\[
\text{upper: }\sha{fccccf5de8673379d37ff0fc840ab091d9ff7a2352ff64a93091f01c21bf4cb4}.
\]
Elapsed times are machine-dependent sanity checks; the counts, zero-failure results, extremal
tuples, and hashes are the reproducibility data.  The hashes cover the ASCII stream
\lean{m:n;} of proposed-and-certified pairs and make accidental changes to the candidate
sequence detectable across reruns.

\begin{statusbox}{\statuscomp{} --- interval arithmetic in \lean{mpmath} \cite{mpmath} at 40
digits.  This is a software trust boundary, not a kernel certificate.  The kernel-verified
zero-free region remains $2^x<4096$ in \lean{FiniteCheck4096}.}
\end{statusbox}

\section{Suggested experimental determinant families}
\label{app:experiments}

Before developing a full $p$-adic theory, several finite experiments can test whether the
required amplification exists.

\subsection{Rectangular row and column boxes}

Choose rows $(n,k)$ with $0\le n<N_0$, $0\le k<K_0$ and columns $(a,b)$ with
$0\le a<A_0$, $0\le b<B_0$.  The entry is
\[
 (2^a3^b)^n(M^aA^b)^k.
\]
For complete rectangles, the matrix is a rearranged Kronecker product only when the row and
column cardinalities match in a compatible way.  Its determinant or maximal minors can be
factored symbolically for small boxes.  Record valuations at $2$ and $3$ as polynomials in
$a_0=\vpp{2}{M}$ and $b_0=\vpp{3}{A}$.

\subsection{Triangular boxes adapted to height}

Use
\[
 n+k\beta\le T,
 \qquad
 a+b\thetaq\le L.
\]
These are the natural row and column shapes from the exact counting theorem and the
mixed-exponent semigroup.  Search for square minors with minimal archimedean weight and
maximal forced valuation.  The linear boundary slopes are irrational, so floor-sum
fluctuations may prevent cancellation patterns present in rectangles.

\subsection{Finite-difference row operations}

Order rows by $n$ for fixed $k$, or by $k$ for fixed $n$, and apply repeated difference
operators.  Differences introduce factors such as
\[
 2^a3^b-1,
 \qquad
 M^aA^b-1.
\]
These may have systematic $2$- or $3$-adic valuations for carefully selected column pairs.
Lifting-the-exponent lemmas can then provide exact lower bounds.  The risk is that the
archimedean size of the same factors grows too quickly; experiments should track both sides.

\subsection{Confluent columns}

Instead of only distinct mixed monomials, consider derivatives with respect to an auxiliary
continuous exponent parameter before specializing.  Confluent generalized Vandermonde
matrices introduce logarithmic factors, which are not integral, but common-denominator or
linear-form constructions may restore arithmetic control.  This is closer to Laurent's
interpolation determinant method \cite{Waldschmidt2000} and may reduce total exponent weight.

\section{Dependency map and Lean-oriented proof sketches}
\label{app:sketches}

The conditional arithmetic criteria of this report were written with formalization in mind:
each one is a short argument over an already kernel-verified base.  This appendix records the
shortest route from the existing corpus to each of them, and then gives skeleton statements
for the four criteria whose formalization is least obvious.

\begin{statusbox}{Design note.  The snippets below are schematic ASCII transliterations and
are not claimed to compile unchanged.  The mathematics they target is \statusnew{}; none of
it has been accepted by the kernel, and nothing here may be read as a claim of verification.}
\end{statusbox}

\subsection{Dependency map}
\label{a3:depmap}

Each row names one new result, the principal existing input it is built on, and the extra
mathematics a formalization would have to supply.  All of the new results are conditional:
they describe the structure forced on $\Sol$ and $\Cand$ by a hypothetical counterexample to
Conjecture~\ref{conj:AE}, and they are vacuous if the conjecture holds.

\begin{center}
\small
\begin{tabular}{@{}p{0.26\textwidth}p{0.33\textwidth}p{0.31\textwidth}@{}}
\toprule
New result & Principal existing input & Additional mathematics \\
\midrule
rational-function rigidity &
canonical primitive generator; transcendence of $\beta$ &
injectivity of polynomial evaluation at a transcendental element \\
\addlinespace
product and power criteria &
unique $(n,k)$ coordinates &
quadratic coefficient comparison \\
\addlinespace
quotient criterion &
unique $(n,k)$ coordinates &
denominator clearing and coefficient comparison \\
\addlinespace
three-smooth output test &
transcendence of $\thetaq$ &
unique factorization in $\N$ \\
\addlinespace
product transcendence &
three-smooth-base theorem &
classical six exponentials \\
\addlinespace
common power degree &
unique $(n,k)$ coordinates &
positive root uniqueness; divisibility \\
\addlinespace
primitive statistics &
total lattice count &
M\"obius inversion \\
\addlinespace
gcd and lcm lattice &
exact candidate monoid &
prime-valuation min and max \\
\addlinespace
ambient root saturation &
primitive valuation gcd of $w$ &
B\'ezout and linear cone arithmetic \\
\addlinespace
optimized differences &
arbitrary-order difference theorem &
gamma identity and Stirling estimate \\
\bottomrule
\end{tabular}
\end{center}

Every entry in the middle column is \statuslean{} in the corpus catalogued in
Section~\ref{sec:corpus}.  In the right-hand column, only the six exponentials theorem is
external background (\statusknown); the remaining items are ordinary mathlib-level algebra
and analysis, which is what makes these criteria cheap to formalize once the base is in
place.

\paragraph{Transliteration conventions.}
The listings write Unicode in plain ASCII: \lean{R}, \lean{Q}, \lean{Z} and \lean{N} are the
real, rational, integer and natural types; arrows are \lean{->} and \lean{<->}; membership is
\lean{in} and divisibility is \lean{dvd}; anonymous constructors use angle brackets; the two
logical connectives use their two-character ASCII spellings; and the focusing dot is written
as a period.  Coercions are written as explicit casts such as \lean{Int.cast} rather than as
the coercion arrow.

\subsection{Multiplication skeleton}
\label{a3:mult-skeleton}

A direct statement can use the repository predicate \lean{TwoBaseIntegralSolution}, which is
the defining conjunction $2^x\in\Z$ and $3^x\in\Z$.

\begin{lstlisting}[style=leanstyle,caption={Product criterion: skeleton.},label={a3:lst-prod}]
theorem product_solution_iff_integer_factor
    {x y : R}
    (hx : TwoBaseIntegralSolution x)
    (hy : TwoBaseIntegralSolution y) :
    TwoBaseIntegralSolution (x * y) <->
      x in Set.range (Int.cast : Z -> R) \/
      y in Set.range (Int.cast : Z -> R) := by
  classical
  constructor
  . intro hxy
    by_cases hAE : AlaogluErdosConjecture
    . left
      exact hAE x hx
    . obtain <w, d, a, c, beta, ..., hcoord, ...> :=
        exists_canonical_primitiveGenerator_of_not_alaogluErdosConjecture hAE
      obtain <nx, kx, hxcoord> := (hcoord x).mp hx
      obtain <ny, ky, hycoord> := (hcoord y).mp hy
      obtain <np, kp, hpcoord> := (hcoord (x * y)).mp hxy
      -- If kx and ky are both positive, expanding the product gives a
      -- nonzero quadratic polynomial over Q vanishing at beta, which
      -- contradicts transcendence of beta.  Hence kx = 0 or ky = 0.
      ...
  . rintro (hxint | hyint)
    . exact solution_nsmul_left hxint hy
    . exact solution_nsmul_right hyint hx
\end{lstlisting}

In the finished theorem the integer range can be narrowed to $\Nzero$, because every solution
is nonnegative by \lean{nonneg_of_two_rpow_integer}.  The algebraic step should be isolated,
since three of the four criteria reduce to it.  Note that the coefficients must be rational
--- with arbitrary real $a,b,c$ the statement is false --- and that is exactly how the
applications supply them, as differences of the integer coordinates:

\begin{lstlisting}[style=leanstyle,caption={The shared coefficient lemma.},label={a3:lst-quad}]
lemma transcendental_quadratic_coeff_eq_zero
    {beta : R} {a b c : Q}
    (hbeta : Transcendental Q beta)
    (h : (a : R) * beta ^ 2 + (b : R) * beta + (c : R) = 0) :
    a = 0 /\ b = 0 /\ c = 0 := ...
\end{lstlisting}

Mathlib's polynomial evaluation API is probably the cleanest route: build the polynomial
$aX^{2}+bX+c$, rewrite \lean{h} as an evaluation, and apply the definition of transcendence
to conclude that the polynomial is zero.

\subsection{Output smoothness skeleton}
\label{a3:smooth-skeleton}

Define the predicate
\[
  \operatorname{ThreeSmooth}(m)
  :\Longleftrightarrow
  \exists u,v\in\Nzero,\ m=2^{u}3^{v}.
\]
A practical implementation should avoid a partial ``natural output'' function and instead
carry explicit witnesses $M,A\in\N$ for the two integer powers:

\begin{lstlisting}[style=leanstyle,caption={Output smoothness: skeleton.},label={a3:lst-smooth}]
theorem integer_of_outputs_threeSmooth
    {x : R} {M A : N}
    (hM : (2 : R) ^ x = (M : R))
    (hA : (3 : R) ^ x = (A : R))
    (hMs : ThreeSmooth M) (hAs : ThreeSmooth A) :
    x in Set.range (Int.cast : Z -> R) := by
  obtain <a, b, rfl> := hMs
  obtain <c, d, rfl> := hAs
  -- Taking logarithms base two: x = a + b * theta and x * theta = c + d * theta.
  have hpoly :
      (b : R) * logThreeDivLogTwo ^ 2
        + ((a : R) - (d : R)) * logThreeDivLogTwo - (c : R) = 0 := by
    ...
  have htheta := transcendental_logThreeDivLogTwo
  -- The coefficients are integers, so all three vanish:
  -- b = 0, a = d, c = 0, whence 2 ^ x = 2 ^ a and x = a.
  ...
\end{lstlisting}

The converse direction is immediate: for integral $x=a$ the outputs are $2^{a}$ and $3^{a}$.
The interesting content is that three-smoothness of \emph{both} outputs is already enough to
force integrality, with no appeal to the conditional structure theory.

\subsection{Common-power skeleton}
\label{a3:power-skeleton}

The formal theorem is most reusable when it is stated for explicit output witnesses and for a
fixed canonical generator, so that the coordinate comparison is available as a hypothesis
rather than reconstructed inside the proof.

\begin{lstlisting}[style=leanstyle,caption={Common powers: skeleton.},label={a3:lst-power}]
theorem simultaneous_perfectPower_iff_divides_coordinates
    (hfail : not AlaogluErdosConjecture)
    {x beta : R} {n k r : N}
    (hr : 0 < r)
    (hbeta : CanonicalPrimitiveGenerator beta)
    (hxcoord : x = (n : R) + (k : R) * beta) :
    (Exists fun U : N => (U : R) ^ r = (2 : R) ^ x) /\
    (Exists fun V : N => (V : R) ^ r = (3 : R) ^ x) <->
      r dvd n /\ r dvd k := by
  -- Convert the two perfect powers into a solution at x / r,
  -- apply uniqueness of coordinates, and compare.
  ...
\end{lstlisting}

The positive-root step can be discharged either by strict monotonicity of $t\mapsto t^{r}$ on
the nonnegative reals or by injectivity of the logarithm; the first is usually shorter in
mathlib.

\subsection{Ambient saturation skeleton}
\label{a3:saturation-skeleton}

The only genuinely arithmetic lemma in the saturation argument is the reconstruction of a
root of a power of a primitive number:

\begin{lstlisting}[style=leanstyle,caption={Primitive-core root lemma.},label={a3:lst-root}]
lemma primitive_core_root
    {w u q r : N}
    (hw : oddPrimeValuationGCD w = 1)
    (hr : 0 < r)
    (hpow : u ^ r = w ^ q) :
    Exists fun t : N => q = r * t /\ u = w ^ t := by
  -- Compare every coordinate of the prime factorizations.
  -- A Bezout combination of the positive valuations of w gives r | q;
  -- factorization extensionality then gives u = w ^ t.
  ...
\end{lstlisting}

The proof in this report only needs the case $q=dk$, but the general lemma is worth stating
once: the same reconstruction is used implicitly in \lean{PrimitiveOddCore} and in
\lean{OutputNormalization}, and a shared version would remove that duplication.

\subsection{Suggested theorem order for a pull request}
\label{a3:pr-order}

A low-conflict sequence, each step reviewable on its own, is:

\begin{enumerate}
\item add a small \lean{ThreeSmooth} API together with the output-smoothness theorem;
\item add coordinate comparison lemmas for affine, quadratic and quotient expressions;
\item prove the product, power and quotient criteria;
\item prove the common-power and affine-descent criteria;
\item add the candidate divisibility and support results;
\item add primitive-core root reconstruction and ambient saturation;
\item only then add the asymptotic and six-exponentials extensions.
\end{enumerate}

This ordering keeps the first several commits elementary, and it defers everything that
depends on external transcendence input to the end, where it can be reviewed against the
status discipline of Appendix~\ref{app:status} rather than mixed into routine algebra.

\section{Directed-rounding verifier}
\label{app:scanner}

Appendix~\ref{app:interval} gives the interval-arithmetic script behind the $2^{20}$
verification.  This appendix gives the second, independent verifier: a C program that uses
MPFR directed rounding rather than an interval library, and that extends the scanned range
sixty-fourfold, to $2^{26}$.

\begin{statusbox}{\statuscomp{} --- MPFR $4.2.2$ at $192$ bits with directed rounding modes.
Like the $2^{20}$ scan this is a software trust boundary, not a kernel certificate.  The
kernel-verified zero-free region remains $2^{x}<4096$ in \lean{FiniteCheck4096}.}
\end{statusbox}

\subsection{What the verifier certifies}
\label{a3:verifier-contract}

For a non-power of two $m$, proving $m^{\thetaq}\notin\N$ amounts to exhibiting $n\in\N$ with
$n<m^{\thetaq}<n+1$, and, as in Section~\ref{sec:finitecheck}, that pair of strict
inequalities is certified in logarithmic form so that no exponentiation is ever performed.
At precision $p=192$ the program computes directed enclosures
\[
  L_m^{-}\le\log m\le L_m^{+},
  \qquad
  L_2^{-}\le\log2\le L_2^{+},
  \qquad
  L_3^{-}\le\log3\le L_3^{+},
\]
together with $L_n^{+}\ge\log n$ and $L_{n+1}^{-}\le\log(n+1)$, and then rounds
\emph{downward} in
\[
  \delta_{-}=L_m^{-}L_3^{-}-L_n^{+}L_2^{+},
  \qquad
  \delta_{+}=L_{n+1}^{-}L_2^{-}-L_m^{+}L_3^{+}.
\]
Positivity of $\delta_{-}$ proves the lower inequality and positivity of $\delta_{+}$ proves
the upper one.  Every input on the $n$ side is below $2^{42}$ and is therefore exactly
representable at $192$ bits.

An ordinary \lean{long double} evaluation proposes $n\approx\lfloor m^{\thetaq}\rfloor$ and
the offsets $0,-1,1,-2,2,-3,3$ are tried in that order.  The proposal has no logical role: a
proposal is accepted only when both directed lower bounds are strictly positive, and a
failure to locate any admissible $n$ is reported as a failure rather than silently
reclassified.  Locator error can therefore produce a false alarm but not a false proof.
Powers of two are skipped rather than tested, since $\left(2^{q}\right)^{\thetaq}=3^{q}$ is
the known integral candidate; an exact integral value at a non-power of two would satisfy
neither strict inequality for any $n$ and would be reported.

\subsection{Source identity and environment}
\label{a3:verifier-env}

The recorded run used MPFR $4.2.2$ linked as \lean{libmpfr.so.6}, GCC $14.2.0$ with
\lean{-O3 -fopenmp}, GNU OpenMP with five worker threads, Linux on x86-64, and $192$-bit
precision with the directed modes \lean{MPFR_RNDD} and \lean{MPFR_RNDU}.  The SHA-256 digest
of the exact executed source reproduced below is
\begin{center}
\sha{f888c0d89be9366f01e2d26f339dcbe446556d39d4c5dbcf8f579ec0d0052e1d}
\end{center}
The execution image contained the MPFR shared library but not its development header, so the
source carries a guarded minimal declaration of the MPFR~4.x ABI; on a system with
\lean{mpfr.h} present the standard header branch is taken instead.  That ABI detail is part
of the software trust boundary and is disclosed here rather than hidden.

\subsection{Source}
\label{a3:verifier-source}

The listing below is the executed source verbatim.  Its language is C99; it is typeset in
this report's generic code style, which has no C-specific highlighting.

\begin{lstlisting}[style=pythonstyle,caption={Directed-rounding verifier (C99, MPFR, OpenMP).}]
#define _GNU_SOURCE
#include <stdio.h>
#include <stdlib.h>
#include <stdint.h>
#include <inttypes.h>
#include <math.h>
#include <omp.h>

#if defined(__has_include)
#  if __has_include(<mpfr.h>)
#    include <mpfr.h>
#    define HAVE_SYSTEM_MPFR_HEADER 1
#  endif
#endif

#ifndef HAVE_SYSTEM_MPFR_HEADER
/* Minimal declarations for the MPFR 4.x ABI.  The execution environment had
   libmpfr.so.6 but not the development header package.  A normal build with
   libmpfr-dev installed takes the branch above instead. */
typedef unsigned long mp_limb_t;
typedef long mpfr_prec_t;
typedef long mpfr_exp_t;
typedef int mpfr_sign_t;
typedef struct {
  mpfr_prec_t _mpfr_prec;
  mpfr_sign_t _mpfr_sign;
  mpfr_exp_t _mpfr_exp;
  mp_limb_t *_mpfr_d;
} __mpfr_struct;
typedef __mpfr_struct mpfr_t[1];
typedef __mpfr_struct *mpfr_ptr;
typedef const __mpfr_struct *mpfr_srcptr;
typedef enum {
  MPFR_RNDN = 0, MPFR_RNDZ = 1, MPFR_RNDU = 2,
  MPFR_RNDD = 3, MPFR_RNDA = 4, MPFR_RNDF = 5,
  MPFR_RNDNA = -1
} mpfr_rnd_t;
extern void mpfr_init2(mpfr_ptr, mpfr_prec_t);
extern void mpfr_clear(mpfr_ptr);
extern int mpfr_set_ui(mpfr_ptr, unsigned long, mpfr_rnd_t);
extern int mpfr_log(mpfr_ptr, mpfr_srcptr, mpfr_rnd_t);
extern int mpfr_mul(mpfr_ptr, mpfr_srcptr, mpfr_srcptr, mpfr_rnd_t);
extern int mpfr_sub(mpfr_ptr, mpfr_srcptr, mpfr_srcptr, mpfr_rnd_t);
extern int mpfr_cmp_ui(mpfr_srcptr, unsigned long);
extern double mpfr_get_d(mpfr_srcptr, mpfr_rnd_t);
extern const char *mpfr_get_version(void);
#endif

static inline int is_power_of_two_u64(uint64_t x) {
  return x && ((x & (x - 1)) == 0);
}

typedef struct {
  uint64_t checked;
  uint64_t failures;
  double min_lower;
  double min_upper;
  uint64_t min_lower_m, min_lower_n;
  uint64_t min_upper_m, min_upper_n;
  uint64_t first_fail_m, first_fail_guess;
} Stats;

int main(int argc, char **argv) {
  uint64_t start = 1;
  uint64_t limit = (1ULL << 20);
  int bits = 192;
  if (argc > 1) start = strtoull(argv[1], NULL, 0);
  if (argc > 2) limit = strtoull(argv[2], NULL, 0);
  if (argc > 3) bits = atoi(argv[3]);
  if (start < 1 || limit <= start) {
    fprintf(stderr, "bad range\n");
    return 2;
  }

  int nt = omp_get_max_threads();
  Stats *stats = calloc((size_t)nt, sizeof(Stats));
  if (!stats) return 3;
  for (int i = 0; i < nt; ++i) {
    stats[i].min_lower = INFINITY;
    stats[i].min_upper = INFINITY;
  }

  const long double theta = logl(3.0L) / logl(2.0L);
  double t0 = omp_get_wtime();

  mpfr_t two, three, ln2_lo, ln2_hi, ln3_lo, ln3_hi;
  mpfr_init2(two, bits); mpfr_init2(three, bits);
  mpfr_init2(ln2_lo, bits); mpfr_init2(ln2_hi, bits);
  mpfr_init2(ln3_lo, bits); mpfr_init2(ln3_hi, bits);
  mpfr_set_ui(two, 2, MPFR_RNDN);
  mpfr_set_ui(three, 3, MPFR_RNDN);
  mpfr_log(ln2_lo, two, MPFR_RNDD);
  mpfr_log(ln2_hi, two, MPFR_RNDU);
  mpfr_log(ln3_lo, three, MPFR_RNDD);
  mpfr_log(ln3_hi, three, MPFR_RNDU);

#pragma omp parallel
  {
    int tid = omp_get_thread_num();
    Stats *st = &stats[tid];
    mpfr_t xm, xn, xnp1;
    mpfr_t logm_lo, logm_hi, logn_hi, lognp1_lo;
    mpfr_t lhs_lo, lhs_hi, rhs_lo, rhs_hi, dl, du;
    mpfr_init2(xm, bits); mpfr_init2(xn, bits); mpfr_init2(xnp1, bits);
    mpfr_init2(logm_lo, bits); mpfr_init2(logm_hi, bits);
    mpfr_init2(logn_hi, bits); mpfr_init2(lognp1_lo, bits);
    mpfr_init2(lhs_lo, bits); mpfr_init2(lhs_hi, bits);
    mpfr_init2(rhs_lo, bits); mpfr_init2(rhs_hi, bits);
    mpfr_init2(dl, bits); mpfr_init2(du, bits);

#pragma omp for schedule(static)
    for (uint64_t m = start; m < limit; ++m) {
      if (is_power_of_two_u64(m)) continue;

      /* This is only a locator.  The MPFR inequalities below are the proof. */
      uint64_t guess = (uint64_t)floorl(expl(theta * logl((long double)m)));

      mpfr_set_ui(xm, (unsigned long)m, MPFR_RNDN);
      mpfr_log(logm_lo, xm, MPFR_RNDD);
      mpfr_log(logm_hi, xm, MPFR_RNDU);
      mpfr_mul(lhs_lo, logm_lo, ln3_lo, MPFR_RNDD);
      mpfr_mul(lhs_hi, logm_hi, ln3_hi, MPFR_RNDU);

      int ok = 0;
      uint64_t certn = 0;
      double dl_d = 0.0, du_d = 0.0;
      static const int offsets[7] = {0, -1, 1, -2, 2, -3, 3};
      for (int oi = 0; oi < 7; ++oi) {
        int off = offsets[oi];
        if (off < 0 && guess < (uint64_t)(-off) + 1) continue;
        uint64_t n = (uint64_t)((int64_t)guess + off);
        if (n < 1) continue;

        mpfr_set_ui(xn, (unsigned long)n, MPFR_RNDN);
        mpfr_set_ui(xnp1, (unsigned long)(n + 1), MPFR_RNDN);

        /* Directed lower bound for log(m)log(3)-log(n)log(2). */
        mpfr_log(logn_hi, xn, MPFR_RNDU);
        mpfr_mul(rhs_hi, logn_hi, ln2_hi, MPFR_RNDU);
        mpfr_sub(dl, lhs_lo, rhs_hi, MPFR_RNDD);
        if (mpfr_cmp_ui(dl, 0) <= 0) continue;

        /* Directed lower bound for log(n+1)log(2)-log(m)log(3). */
        mpfr_log(lognp1_lo, xnp1, MPFR_RNDD);
        mpfr_mul(rhs_lo, lognp1_lo, ln2_lo, MPFR_RNDD);
        mpfr_sub(du, rhs_lo, lhs_hi, MPFR_RNDD);
        if (mpfr_cmp_ui(du, 0) <= 0) continue;

        ok = 1;
        certn = n;
        dl_d = mpfr_get_d(dl, MPFR_RNDD);
        du_d = mpfr_get_d(du, MPFR_RNDD);
        break;
      }

      if (!ok) {
        ++st->failures;
        if (!st->first_fail_m || m < st->first_fail_m) {
          st->first_fail_m = m;
          st->first_fail_guess = guess;
        }
        continue;
      }

      ++st->checked;
      if (dl_d < st->min_lower) {
        st->min_lower = dl_d;
        st->min_lower_m = m;
        st->min_lower_n = certn;
      }
      if (du_d < st->min_upper) {
        st->min_upper = du_d;
        st->min_upper_m = m;
        st->min_upper_n = certn;
      }
    }

    mpfr_clear(xm); mpfr_clear(xn); mpfr_clear(xnp1);
    mpfr_clear(logm_lo); mpfr_clear(logm_hi);
    mpfr_clear(logn_hi); mpfr_clear(lognp1_lo);
    mpfr_clear(lhs_lo); mpfr_clear(lhs_hi);
    mpfr_clear(rhs_lo); mpfr_clear(rhs_hi);
    mpfr_clear(dl); mpfr_clear(du);
  }

  Stats total = {0};
  total.min_lower = INFINITY;
  total.min_upper = INFINITY;
  for (int i = 0; i < nt; ++i) {
    Stats *s = &stats[i];
    total.checked += s->checked;
    total.failures += s->failures;
    if (s->min_lower < total.min_lower) {
      total.min_lower = s->min_lower;
      total.min_lower_m = s->min_lower_m;
      total.min_lower_n = s->min_lower_n;
    }
    if (s->min_upper < total.min_upper) {
      total.min_upper = s->min_upper;
      total.min_upper_m = s->min_upper_m;
      total.min_upper_n = s->min_upper_n;
    }
    if (s->first_fail_m &&
        (!total.first_fail_m || s->first_fail_m < total.first_fail_m)) {
      total.first_fail_m = s->first_fail_m;
      total.first_fail_guess = s->first_fail_guess;
    }
  }

  double elapsed = omp_get_wtime() - t0;
  printf("mpfr_version=%s\n", mpfr_get_version());
  printf("range=[%" PRIu64 ",%" PRIu64 ")\n", start, limit);
  printf("precision_bits=%d\nthreads=%d\n", bits, nt);
  printf("checked_nonpowers=%" PRIu64 "\n", total.checked);
  printf("failures=%" PRIu64 "\n", total.failures);
  printf("min_lower_log_margin=%.17g at (m,n)=(%" PRIu64 ",%" PRIu64 ")\n",
         total.min_lower, total.min_lower_m, total.min_lower_n);
  printf("min_upper_log_margin=%.17g at (m,n)=(%" PRIu64 ",%" PRIu64 ")\n",
         total.min_upper, total.min_upper_m, total.min_upper_n);
  if (total.failures) {
    printf("first_failure=(m,guess)=(%" PRIu64 ",%" PRIu64 ")\n",
           total.first_fail_m, total.first_fail_guess);
  }
  printf("elapsed_seconds=%.6f\n", elapsed);

  mpfr_clear(two); mpfr_clear(three);
  mpfr_clear(ln2_lo); mpfr_clear(ln2_hi);
  mpfr_clear(ln3_lo); mpfr_clear(ln3_hi);
  free(stats);
  return total.failures ? 1 : 0;
}
\end{lstlisting}

\subsection{Build and invocation}
\label{a3:verifier-build}

On a system with the MPFR development headers installed, the whole build and a single
full-range run are:

\begin{lstlisting}[style=pythonstyle,numbers=none,caption={Build and invocation (POSIX shell).}]
gcc -O3 -fopenmp mpfr_scan.c -lmpfr -lm -o mpfr_scan
OMP_NUM_THREADS=5 ./mpfr_scan 1 67108864 192
\end{lstlisting}

The recorded audit instead used sixteen shorter invocations over half-open subranges, so that
each chunk produces an independently inspectable log.  The exit status is nonzero exactly
when some input failed to be certified, which makes the chunked run scriptable without
parsing the output.

\section{Chunk audit and reproducibility}
\label{app:chunkaudit}

\begin{statusbox}{\statuscomp{} --- the numbers in this appendix are the output of the
verifier of Appendix~\ref{app:scanner}.  They are reproducible data at a software trust
level, not kernel certificates; the kernel-verified exponent bound remains $x\ge12$.}
\end{statusbox}

There are $2^{26}-1=67{,}108{,}863$ positive integers below $2^{26}$, of which exactly $26$
are powers of two, namely $2^{0}$ through $2^{25}$.  The verifier therefore tests
$67{,}108{,}863-26=67{,}108{,}837$ inputs.  The range was split at multiples of $2^{22}$ into
one initial chunk $[1,2^{22})$ and fifteen chunks of width $2^{22}$.  The canonical chunk log
holds the full standard output of every invocation and has SHA-256 digest
\begin{center}
\sha{dcaa7f3b7e5be44acaee9015a2045da14ee42d8fa2d96b9619cdd44c072a6fd7}
\end{center}

Table~\ref{a3:tab:chunkaudit} reports, for each chunk, the exact scanned integer interval, the
number of non-powers certified, and the smallest directed lower and upper logarithmic margins
recorded in that chunk.  Every chunk returned a failure count of zero.

{\small
\begin{longtable}{@{}r p{0.26\textwidth} r r r@{}}
\caption{Directed-rounding audit through $2^{26}$.}\label{a3:tab:chunkaudit}\\
\toprule
Chunk & Integer range & Checked & Smallest lower gap & Smallest upper gap \\
\midrule
\endfirsthead
\toprule
Chunk & Integer range & Checked & Smallest lower gap & Smallest upper gap \\
\midrule
\endhead
0 & 1--4{,}194{,}303 & 4{,}194{,}281 & $1.1539\times10^{-17}$ & $3.0076\times10^{-17}$ \\
1 & 4{,}194{,}304--8{,}388{,}607 & 4{,}194{,}303 & $4.6319\times10^{-20}$ &
$1.3515\times10^{-19}$ \\
2 & 8{,}388{,}608--12{,}582{,}911 & 4{,}194{,}303 & $1.2811\times10^{-18}$ &
$1.3515\times10^{-19}$ \\
3 & 12{,}582{,}912--16{,}777{,}215 & 4{,}194{,}304 & $4.6319\times10^{-20}$ &
$2.9082\times10^{-19}$ \\
4 & 16{,}777{,}216--20{,}971{,}519 & 4{,}194{,}303 & $1.1898\times10^{-18}$ &
$1.3515\times10^{-19}$ \\
5 & 20{,}971{,}520--25{,}165{,}823 & 4{,}194{,}304 & $5.4212\times10^{-19}$ &
$1.5111\times10^{-19}$ \\
6 & 25{,}165{,}824--29{,}360{,}127 & 4{,}194{,}304 & $4.6319\times10^{-20}$ &
$1.7669\times10^{-19}$ \\
7 & 29{,}360{,}128--33{,}554{,}431 & 4{,}194{,}304 & $1.4439\times10^{-19}$ &
$5.4132\times10^{-21}$ \\
8 & 33{,}554{,}432--37{,}748{,}735 & 4{,}194{,}303 & $4.2839\times10^{-19}$ &
$1.3515\times10^{-19}$ \\
9 & 37{,}748{,}736--41{,}943{,}039 & 4{,}194{,}304 & $3.9292\times10^{-19}$ &
$1.0460\times10^{-19}$ \\
10 & 41{,}943{,}040--46{,}137{,}343 & 4{,}194{,}304 & $1.9177\times10^{-19}$ &
$1.5111\times10^{-19}$ \\
11 & 46{,}137{,}344--50{,}331{,}647 & 4{,}194{,}304 & $1.0117\times10^{-19}$ &
$3.4145\times10^{-19}$ \\
12 & 50{,}331{,}648--54{,}525{,}951 & 4{,}194{,}304 & $5.6927\times10^{-20}$ &
$1.2322\times10^{-19}$ \\
13 & 54{,}525{,}952--58{,}720{,}255 & 4{,}194{,}304 & $3.6447\times10^{-20}$ &
$1.1345\times10^{-19}$ \\
14 & 58{,}720{,}256--62{,}914{,}559 & 4{,}194{,}304 & $6.7779\times10^{-20}$ &
$5.4132\times10^{-21}$ \\
15 & 62{,}914{,}560--67{,}108{,}863 & 4{,}194{,}304 & $5.0923\times10^{-20}$ &
$3.9197\times10^{-20}$ \\
\midrule
\multicolumn{2}{@{}r}{\textbf{total}} & \textbf{67{,}108{,}837} & & \\
\bottomrule
\end{longtable}
}

\subsection{Extremal records}
\label{a3:records}

The global lower-gap record is the chunk~13 value $3.6447\times10^{-20}$, attained at
\[
  (m,n)=(58{,}570{,}438,\ 2{,}048{,}703{,}434{,}093),
\]
where a separate $100$-digit evaluation gives
$m^{\thetaq}-n=1.07725158750919017548\ldots\times10^{-7}$.  The global upper-gap record is
the chunk~7 value $5.4132\times10^{-21}$, attained at
\[
  (m,n)=(29{,}411{,}704,\ 687{,}581{,}893{,}443),
\]
where $n+1-m^{\thetaq}=5.36974926710988874953\ldots\times10^{-9}$.  The same upper margin
reappears in chunk~14 at the scaled pair $(58{,}823{,}408,\ 2{,}062{,}745{,}680{,}331)$; this
is not a coincidence but the identity $(2m)^{\thetaq}=3\,m^{\thetaq}$, which multiplies the
input by two and its $\thetaq$-power by three while leaving the logarithmic margin unchanged.

Both extremal inequalities were reevaluated independently at $100$ decimal digits.  That
secondary check is not an independent proof system --- it shares the same floating-point
philosophy --- but it does guard against transcription and aggregation mistakes, and the
certified margins sit many orders of magnitude above the nominal $192$-bit rounding scale of
the logarithmic operations.

The $2^{20}$ run of Appendix~\ref{app:interval} and the present run are independent
implementations, one using interval arithmetic in \lean{mpmath} and the other MPFR directed
rounding, and they agree on the overlapping range $1\le m<2^{20}$: no failures.  Their
recorded extremal metadata are not directly comparable, because the chunk boundaries differ.

\subsection{Reproducibility checklist}
\label{a3:checklist}

To reproduce the finite claim independently:

\begin{enumerate}
\item obtain MPFR~4.2.x and GMP on a $64$-bit platform;
\item save the source of Appendix~\ref{app:scanner} and verify its SHA-256 digest;
\item compile with optimization and OpenMP, or delete the OpenMP pragmas for a serial run;
\item run the sixteen half-open intervals of Table~\ref{a3:tab:chunkaudit} at $192$ bits;
\item verify that every invocation reports \lean{failures=0};
\item sum the reported \lean{checked_nonpowers} and obtain $67{,}108{,}837$;
\item compare the two global extremal records above;
\item optionally rerun at $256$ or $320$ bits, and on a second MPFR build;
\item for a theorem-grade result, translate the intervals into exact rational certificates and
replay them in Lean rather than trusting the executable.
\end{enumerate}

The finite statement is deterministic.  Thread count changes execution time and the reduction
order that produces the recorded minimum metadata, but it cannot change whether an individual
directed inequality is certified.  Step nine is the one that would change the status of the
result rather than its range, and it is the reason the scan is presented here as data for
certificate generation --- in the sense of Section~\ref{sec:program} --- and not as an
extension of the kernel-verified zero-free region.

\section{Classical transcendence inputs and their exact roles}\label{app:classical}

\begin{statusbox}{\statusknown{} --- this appendix states external theorems, not new ones.
Exactly one of them, Gelfond--Schneider, is formalized in this repository
(Theorem~\ref{co:thm-gs}); the remaining ones enter only through paper proofs, which are
tagged \statusnew{} where they appear.  The corpus does not introduce any of them as an
axiom: where a conjectural or unformalized statement is needed in Lean it is defined as a
proposition and the implication is proved conditionally on it.  Deleting six exponentials,
five exponentials, Baker's theorems, or Roy's theorem from the report would change which
paper results stand, and would change none of the kernel-verified results catalogued in
Appendix~\ref{app:status}.}
\end{statusbox}

The purpose of this appendix is auditability.  A reader who wants to know how much of the
report is elementary consequence of the repository hypotheses, and how much rests on major
transcendence theory, should be able to answer that question by looking in one place.  For
each classical input below we give the statement in the form actually used, the exact
places where the report uses it, and the exact reason it stops short of
Conjecture~\ref{conj:AE}.  The last reason is the same in every case, seen from a different
angle: every input listed here constrains a \emph{linear} configuration of logarithms or an
\emph{algebraic} exponent, while the hypothetical counterexample is a \emph{bilinear}
relation between logarithms with a transcendental exponent.

\begin{description}[leftmargin=1.2em,style=nextline,itemsep=0.9em]

\item[Gelfond--Schneider (Hilbert's seventh problem)]
\emph{Statement.}  Let $\alpha,\beta\in\Qbar$ with $\alpha\ne0,1$ and $\beta\notin\Q$.  Then
every value of $\alpha^{\beta}$ is transcendental \cite{Hua1982,Schneider1957,Baker1975}.
\statusknown{}

\emph{Role here.}  This is the one classical transcendence theorem the corpus proves rather
than cites; the Lean proof follows the classical exposition \cite{Hua1982} and is recorded as
Theorem~\ref{co:thm-gs}.  It is used twice.  First, a nonintegral solution $x$ is shown
irrational by the elementary rational-exponent argument, and then must be transcendental
because $2^{x}$ is an algebraic number; this is the algebraic--integral dichotomy of
Corollary~\ref{co:cor-transcendental}, and it is what makes the conditional primitive
generator $\beta$ and the odd-core exponent $\log_2 w$ transcendental throughout
Section~\ref{sec:conditional} and Section~\ref{sec:radical}.  Second, $\thetaq=\log_2 3$ is
irrational, hence transcendental, because $2^{\thetaq}=3$.

\emph{Why it does not settle Conjecture~\ref{conj:AE}.}  The theorem speaks only about
algebraic irrational exponents.  Once the dichotomy has removed algebraic exponents, the
remaining hypothetical exponent is transcendental, and Gelfond--Schneider has nothing further
to say about $2^{x}$ for transcendental $x$.  The theorem closes the easy half of the problem
and defines the hard half.

\item[Six exponentials]
\emph{Statement.}  Let $u_1,u_2,u_3\in\C$ be linearly independent over $\Q$ and let
$v_1,v_2\in\C$ be linearly independent over $\Q$.  Then at least one of the six numbers
\[
  e^{u_iv_j}\qquad(1\le i\le3,\ 1\le j\le2)
\]
is transcendental \cite{Waldschmidt2000,Waldschmidt2005,WaldschmidtSix}.  \statusknown{}

\emph{Role here.}  It is the analytic engine behind every exact spectrum statement in the
report.  The saturated algebraic-base spectrum --- for multiplicatively independent
$a,b\in\Apos$ with $a^{x},b^{x}\in\Qbar$ and $x$ irrational, $\alpha^{x}\in\Qbar$ if and only
if $\alpha\in\GammaSat(a,b)$ --- is the case with rows $\log a,\log b,\log\alpha$ and columns
$1,x$.  The conditional classification of the symmetric constant $E_3=3^{\thetaq}$ is the
case with rows $1,x,\thetaq$ and columns $\log2,\log3$, and the classification of
$E_2=2^{\eta}$ is its mirror image under $\thetaq\mapsto\eta=1/\thetaq$.  The ordinary form
of the unconditional $E_3/E_2$ dichotomy uses the same configuration.  The kernel-verified
$3$-smooth classification of auxiliary bases and the three-base theorem
$2^{x},3^{x},5^{x}\in\Z\Rightarrow x\in\Z$ are the formalized integer-level shadows of this
theorem; the Lean proofs proceed by interpolation determinants and do not cite it.

\emph{Why it does not settle Conjecture~\ref{conj:AE}.}  Six exponentials needs three
independent directions on one side of the configuration.  The two-base problem supplies
exactly two logarithmic directions, $\log2$ and $\log3$, and every auxiliary base built from
$2$ and $3$ has logarithm in their rational span, so it creates no third column.  The
remaining $2\times2$ configuration of Section~\ref{sec:fourexp} is precisely the four
exponentials conjecture, which is open.

\item[Five exponentials]
\emph{Statement.}  Let $u_1,u_2$ be linearly independent over $\Q$, let $v_1,v_2$ be linearly
independent over $\Q$, and let $\gamma\in\Qbar$ with $\gamma\ne0$.  Then at least one of the
five numbers
\[
  e^{u_1v_1},\quad e^{u_1v_2},\quad e^{u_2v_1},\quad e^{u_2v_2},\quad e^{\gamma u_2/u_1}
\]
is transcendental \cite{Waldschmidt2005,Roy1992}.  \statusknown{}

\emph{Role here.}  It supplies the two forbidden exponential rays of the hypothetical world:
for a nonintegral solution $x$ and any nonzero algebraic $\gamma$, both $e^{\gamma x}$ and
$e^{\gamma\thetaq}$ are transcendental.  For the first, apply the theorem to the independent
pairs $(1,x)$ and $(\log2,\log3)$, whose four rectangular exponentials are the algebraic
numbers $2,3,M,A$, leaving the fifth transcendental; the second is the same argument with the
roles of the two logarithms exchanged.  Formalizing it would also upgrade the
denominator-controlled auxiliary-base theorem in \lean{RationalThirdBase} from a controlled
denominator to an unrestricted rational output, which is why it appears as a research
priority in Section~\ref{sec:program}.

\emph{Why it does not settle Conjecture~\ref{conj:AE}.}  The theorem does not remove the
fifth exponential, so it does not prove four exponentials.  More concretely, the natural
auxiliary quantities in this problem are $p^{x}=e^{x\log p}$ for a prime $p$, whose
coefficient $\log p$ is transcendental rather than an admissible algebraic $\gamma$.  The
forbidden rays therefore never convert a prime divisor of the integer outputs into a
contradiction, as Section~\ref{sec:attempts} records.

\item[Baker's theorem on linear forms in logarithms]
\emph{Statement.}  Let $\gamma_1,\dots,\gamma_n$ be nonzero algebraic numbers of degree at
most $d$ and height at most $H$, and let $\alpha_1,\dots,\alpha_n$ be algebraic numbers of
height at most $B$, not all zero.  If
\[
  \Lambda=\alpha_1\log\gamma_1+\cdots+\alpha_n\log\gamma_n\ne0,
\]
then $|\Lambda|>B^{-C}$ for an effectively computable $C>0$ depending only on
$n,d,H$ and the degree of the $\alpha_i$ \cite{Baker1975,Waldschmidt2000}.  \statusknown{}

\emph{Role here.}  The effective form is used as a measuring stick rather than as a step in
any proof.  In the rational-approximation analysis of Section~\ref{sec:attempts} it fixes the
scale of what is available: Baker-type lower bounds are polynomial in the relevant heights,
whereas the integer gaps produced by the hypotheses only assert $|\varepsilon|\gtrsim2^{-p}$,
exponentially weaker, so the two scales are consistent and no contradiction arises.  It is
also the theorem that effective irrationality measures for $\thetaq$ would refine.

\emph{Why it does not settle Conjecture~\ref{conj:AE}.}  Baker's theory bounds forms whose
coefficients are \emph{algebraic}.  The decisive relation under failure is
$\log M\,\log3=\log A\,\log2$, and its radical form is
$d\,\log w\,\log3=(\log A-a\log3)\log2$: each side is a product of two logarithms, so the
coefficient of every logarithm is itself a logarithm.  Expanded over prime factorizations
$m=\prod_p p^{m_p}$ and $A=\prod_p p^{a_p}$, failure is the nontrivial vanishing of
\[
  D=\sum_{p}\bigl(m_p\log3-a_p\log2\bigr)\log p,
\]
an integer-coefficient linear form in the \emph{products} $\log p\cdot\log2$ and
$\log p\cdot\log3$.  No lower-bound technology exists for such forms, and even the
$\Q$-linear independence of $\{\log p\cdot\log q\}$ is open; Section~\ref{sec:logproduct}
turns exactly that independence into a sufficient condition.

\item[Baker's homogeneous theorem]
\emph{Statement.}  Let $\gamma_1,\dots,\gamma_n$ be nonzero algebraic numbers.  If
$\log\gamma_1,\dots,\log\gamma_n$ are linearly independent over $\Q$, then
$1,\log\gamma_1,\dots,\log\gamma_n$ are linearly independent over $\Qbar$.  The homogeneous
consequence used here is the shorter one: logarithms of algebraic numbers that are linearly
independent over $\Q$ are linearly independent over $\Qbar$
\cite{Baker1975,Waldschmidt2000}.  \statusknown{}

\emph{Role here.}  It is the bridge from a rational-independence hypothesis to the
$\Qbar$-independence hypothesis that Roy's theorem requires.  In the strong spectrum, the
three logarithms $\log a,\log b,\log\alpha$ are first shown $\Q$-independent by an elementary
saturation argument, and Baker's homogeneous theorem upgrades them to $\Qbar$-independence;
the same upgrade produces the logarithmic independence used in the strong forms of the $E_3$
and $E_2$ classifications.  The ordinary spectrum does not need it, since rational
independence already suffices for six exponentials.

\emph{Why it does not settle Conjecture~\ref{conj:AE}.}  It is a statement about linear
relations among logarithms.  No independence statement about the family $\{\log p\}$, over
$\Q$ or over $\Qbar$, implies anything about the family of products $\{\log p\cdot\log q\}$,
and it is the products that carry the obstruction.

\item[Roy's strong and sharp six exponentials theorems]
\emph{Statement.}  Let $\LL$ be the $\Qbar$-vector space spanned by $1$ together with all
logarithms of nonzero algebraic numbers.  If $u_1,u_2,u_3$ are linearly independent over
$\Qbar$ and $v_1,v_2$ are linearly independent over $\Qbar$, then at least one of the six
products $u_iv_j$ lies outside $\LL$.  The sharp form replaces the independence hypotheses by
a rank condition on a matrix whose entries are linear forms in logarithms of algebraic
numbers \cite{Roy1992,Waldschmidt2005}.  \statusknown{}

\emph{Role here.}  It produces the strengthened conclusions only.  Concretely: the strong
auxiliary-base spectrum, which asserts $x\log\alpha\notin\LL$ for
$\alpha\notin\GammaSat(a,b)$ rather than merely $\alpha^{x}\notin\Qbar$; the strong form of
the symmetric-constant dichotomy, which places $\thetaq\log3$ and $\eta\log2$ outside $\LL$;
and the log-strong transcendence clauses in the second-iterate kernel and depth-three tower
theorems.

\emph{Why it does not settle Conjecture~\ref{conj:AE}.}  Every ordinary transcendence
conclusion in the report survives the removal of this theorem; only the ``lies outside
$\LL$'' clauses disappear.  Structurally it inherits the six exponentials dimension count and
still demands three independent directions on one side, so it does not reach the $2\times2$
configuration where the conjecture actually lives.

\item[Lindemann--Weierstrass]
\emph{Statement.}  If $\alpha_1,\dots,\alpha_n$ are distinct algebraic numbers, then
$e^{\alpha_1},\dots,e^{\alpha_n}$ are linearly independent over $\Qbar$.  The special case
used here is Hermite--Lindemann: for algebraic $\alpha\ne0$ the value $e^{\alpha}$ is
transcendental, equivalently every nonzero value of $\log\gamma$ is transcendental for
algebraic $\gamma\ne0,1$ \cite{Baker1975}.  \statusknown{}

\emph{Role here.}  It is the standing background fact that each $\log p$, and in particular
$\log2$ and $\log3$, is transcendental.  That fact is what makes the coefficients in the
log-product form $D$ above transcendental rather than algebraic, and therefore what places
the wall outside Baker's reach; it is also the precise content of the audit warning against
treating $\log2$ and $\log3$ as algebraic coefficients (Remark~\ref{a4:warn-bilinear}).

\emph{Why it does not settle Conjecture~\ref{conj:AE}.}  Lindemann--Weierstrass constrains
exponentials of \emph{algebraic} arguments.  Under failure the exponent $x$ is transcendental
and the arguments $x\log2$ and $x\log3$ are transcendental, so none of the four corner
exponentials $2,3,2^{x},3^{x}$ is within its scope, and it gives no information whatever
about products of two logarithms.

\end{description}

\subsection{Logical dependency graph}

The table records, for each family of results in the report, the classical input it actually
consumes.  The rows whose input column reports no classical theorem are elementary or
kernel-verified, and depend on no transcendence theory beyond the Gelfond--Schneider
dichotomy that the corpus already formalizes.

\begin{longtable}{p{0.45\textwidth}p{0.47\textwidth}}
\caption{Classical transcendence input consumed by each family of results.}%
\label{a4:tab:deps}\\
\toprule
Result & Required classical input \\
\midrule
\endfirsthead
\toprule
Result & Required classical input \\
\midrule
\endhead
Nonintegral solution is transcendental &
Rational-exponent lemma and Gelfond--Schneider; both kernel-verified \\
Reduction to real four exponentials &
None beyond $\Q$-independence of $1,x$ and of $\log2,\log3$; kernel-verified \\
$3$-smooth auxiliary-base classification and the three-base theorem &
None; interpolation determinants, kernel-verified \\
Exact algebraic-base spectrum for an independent pair &
Six exponentials \\
Strong $\LL$ spectrum &
Exact row independence, Baker's homogeneous theorem, and Roy \\
Valuation lattice and Smith normal form of rational outputs &
Exact spectrum and unique factorization \\
Second-iterate kernel: dimension and the four branch types &
Exact spectrum and elementary prime valuations \\
Kernel invariance across all nonintegral solutions &
Exact spectrum applied twice \\
M\"obius orbit criterion &
Kernel definition and real logarithms \\
Depth-three tower theorem &
Exact spectrum and kernel dimension; Roy only for the strong form \\
Base-$2$ tower equivalence &
Natural-base algebraicity classification and smooth branch algebra \\
Symmetric constants $E_3$ and $E_2$ &
Six exponentials; Baker's homogeneous theorem and Roy for the strong form \\
Two forbidden exponential rays $e^{\gamma x}$, $e^{\gamma\thetaq}$ &
Five exponentials \\
Transcendence of the coefficients $\log p$ in the log-product form &
Lindemann--Weierstrass, in its Hermite--Lindemann case \\
\bottomrule
\end{longtable}

\section{Common false shortcuts and audit warnings}\label{app:falseshortcuts}

\begin{statusbox}{\statusknown{} --- this appendix records negative facts and standard
cautions.  Nothing in it is a new result, and nothing in it changes the status of any
statement elsewhere in the report.}
\end{statusbox}

The problem attracts several plausible but invalid one-line arguments.  Recording them is
part of the audit trail: each is a specific way in which a proposed proof can be circular,
can silently assume the conclusion, or can mistake numerical evidence for exclusion.  Future
paper and formal work should be checked against this list before a short argument is
believed.

\begin{remark}[Fractional-part error]
\label{a4:warn-fract}
Writing $x=n+\alpha$ with $n\in\Z$ and $0<\alpha<1$, and assuming $2^{x}\in\Z$, does
\emph{not} imply $2^{\alpha}\in\Z$.  It implies only $2^{\alpha}\in\Ztwo$, and an irrational
$\alpha$ is not an ordinary radical exponent $1/q$.  The integrality of the output does not
descend to the fractional part.
\end{remark}

\begin{remark}[Rational-root circularity]
\label{a4:warn-circular}
From $M^{\thetaq}\in\Z$ one cannot conclude that $\thetaq$ is rational, nor that $M$ is a
perfect power.  Gelfond--Schneider constrains algebraic irrational exponents, whereas
$\thetaq$ is transcendental.  An argument that quietly treats $\thetaq$ as a rational
exponent has assumed the conclusion.
\end{remark}

\begin{remark}[Baker is linear, not bilinear]
\label{a4:warn-bilinear}
The relation $\log M\,\log3=\log A\,\log2$ is not a Baker linear form.  Treating $\log2$ and
$\log3$ as algebraic coefficients is invalid: by Hermite--Lindemann both are transcendental.
The same objection applies to the radical form $d\,\log w\,\log3=(\log A-a\log3)\log2$ and to
the prime-expanded log-product form of Section~\ref{sec:logproduct}.
\end{remark}

\begin{remark}[No modular real exponent]
\label{a4:warn-modular}
The expression $2^{x}\bmod p$ has no canonical meaning for a general real $x$.  A $p$-adic
exponential requires a $p$-adic exponent and a unit in its convergence domain; the real
equality $2^{x}\in\Z$ supplies neither.  Congruence arguments must therefore be stated for
the integer outputs $M$ and $A$, never for the real exponent.
\end{remark}

\begin{remark}[Candidate quotients need not be candidates]
\label{a4:warn-quotient}
If $m_1,m_2\in\Cand$ and $m_1\mid m_2$, it does not follow that $m_2/m_1\in\Cand$, because
the corresponding quotient of the outputs need not be integral.  Multiplication is safe;
division requires a separate valuation check.  This is the reason $\Cand$ is treated as a
multiplicative monoid rather than as a divisor-closed set.
\end{remark}

\begin{remark}[Density alone is compatible]
\label{a4:warn-density}
The orbit $\{kx\}$ is dense modulo one for irrational $x$, but the associated rational points
have exponential height.  A sequence of $O(\log H)$ points up to height $H$ is compatible
with generic transcendental-curve counting, so density by itself yields no contradiction
without a denominator-sensitive refinement.
\end{remark}

\begin{remark}[Numerical tolerance is not exclusion]
\label{a4:warn-tolerance}
Computing $M^{\thetaq}$ in floating point and comparing its distance to the nearest integer
against a fixed epsilon cannot certify a zero-free region, however wide.  Outward directed
interval bounds or an exact certificate are required, which is why the finite scans in this
report are tagged \statuscomp{} and the kernel-verified exponent bound remains $x\ge12$.
\end{remark}

\begin{remark}[Algebraicity of a power is branch-sensitive]
\label{a4:warn-branch}
From $z^{d}\in\Qbar$ with an integer $d>0$ one may conclude $z\in\Qbar$, but from
$z^{x}\in\Qbar$ with transcendental $x$ no such algebraic closure argument is available.  The
finite-degree argument uses the polynomial $X^{d}-z^{d}$, which has no analogue when the
exponent is not an integer.
\end{remark}

Every warning above was re-derived independently during the preparation of the present
unified report: each was proposed at some point as a short route to the conjecture, and each
failed on exactly the ground stated here.



\section{Logarithmic-energy convergence for projected lattice triangles}
\label{app:energy-convergence}

We record the technical point needed in the triangular determinant limit.  Let
$\lambda_1,\lambda_2>0$ and assume their ratio has finite irrationality type:
there are $c,\tau>0$ such that
\[
 |r\lambda_1+s\lambda_2|\ge c(1+|r|+|s|)^{-\tau}
\]
for all nonzero $(r,s)\in\Z^2$.  Let $t_{1,N}<\cdots<t_{N,N}$ be the first $N$
values of $a\lambda_1+b\lambda_2$ with $a,b\in\Nzero$, and set
$\widetilde t_{i,N}=t_{i,N}/\sqrt N$.  Then the empirical measures converge to
\[
 d\mu(t)=\frac{2t}{U^2}\,dt,\qquad U=\sqrt{2\lambda_1\lambda_2},
\]
and
\begin{equation}\label{eq:energy-limit-app}
 \frac{2}{N^2}\sum_{i<j}\log|\widetilde t_{j,N}-\widetilde t_{i,N}|
 \longrightarrow\iint\log|x-y|\,d\mu(x)d\mu(y).
\end{equation}

The weak convergence follows from a Riemann-sum lattice count in the triangle
$\lambda_1a+\lambda_2b\le U\sqrt N$.  For the singular kernel, truncate
$\log|x-y|$ below at $-M$.  Weak convergence handles the truncated kernel.
Uniform integrability of the negative tail follows by counting difference
vectors $(r,s)$ in thin strips around
$r\lambda_1+s\lambda_2=0$; the finite-type lower bound prevents an individual
difference from being smaller than a fixed power of $N$, while the strip count
shows that the proportion below $\varepsilon$ is $O(\varepsilon)$ up to a
power-saving discrepancy term.  Integrating the distribution function of
$-\log|x-y|$ and then sending $M\to\infty$ proves
\eqref{eq:energy-limit-app}.

For $(\lambda_1,\lambda_2)=(\log2,\log3)$, the required finite type follows from
linear forms in logarithms.  For $(1,\beta)$, it follows from
$\beta=\log A/\log2$ and the multiplicative independence of $A$ and $2$.

\section{Integral evaluations used in the determinant constants}

\subsection{The beta-two logarithmic energy}

For the probability density $f(x)=2x$ on $[0,1]$,
\begin{align*}
 \iint_0^1 4xy\log|x-y|\,dx\,dy
 &=8\int_0^1\int_0^xxy\log(x-y)\,dy\,dx\\
 &=8\int_0^1x^3\left(\frac12\log x-\frac34\right)dx\\
 &=-\frac74.
\end{align*}
Scaling by $U$ adds $\log U$, proving
Equation~\eqref{eq:beta2-energy}.

\subsection{The simplex assignment constant}

A coordinate of a point uniformly distributed on the standard two-simplex has
quantile $q(s)=1-\sqrt{1-s}$.  Therefore
\begin{align*}
 J&=\int_0^1q(s)q(1-s)\,ds\\
 &=\int_0^1\left(1-\sqrt s-\sqrt{1-s}+\sqrt{s(1-s)}\right)ds\\
 &=1-\frac43+\frac\pi8
 =\frac\pi8-\frac13.
\end{align*}
Its comonotone second-moment constant is
$K=\int_0^1q(s)^2ds=1/6$.

\subsection{The balanced-rectangle energy}

Let $Z=(X+Y)/2$ with independent uniform $X,Y\in[0,1]$.  The density of
$S=X+Y$ is triangular, and the density of the difference of two independent
copies of $S$ is
\[
 h(t)=
 \begin{cases}
  \frac23-t^2+\frac12t^3,&0\le t\le1,\\
  \frac16(2-t)^3,&1\le t\le2.
 \end{cases}
\]
A direct integration gives
\[
 \E\log|S-S'|=\frac43\log2-\frac{25}{12}.
\]
Since $Z-Z'=(S-S')/2$,
\[
 \E\log|Z-Z'|=\frac13\log2-\frac{25}{12}.
\]
Also $\E Z^2=7/24$.

\section{General two-prime version}

The structural and determinant arguments extend verbatim to distinct primes
$p,q$.  If a nonintegral $\beta$ satisfies $p^\beta,q^\beta\in\Z$, choose the
least such $\beta$.  Six Exponentials gives the same affine-rank collapse, and
valuation minimality yields the exact monoid $\Nzero+\Nzero\beta$.  The
triangular determinant then gives
\[
 g\!\left(\sqrt{\beta\log p\log q}\right)\ge0.
\]
This form may be useful for comparing the strength of the method across base
pairs and for testing variational improvements.

\section{Finite determinant certificates}

The asymptotic argument can be converted into finite certificates.  For chosen
finite node sets, a certificate consists of:
\begin{enumerate}
\item the exact integer or rational matrix entries;
\item sorted rational intervals for $u_i,v_i$ and their differences;
\item an upper bound from Equation~\eqref{eq:HCIZ-bound};
\item for each relevant prime, dual potentials for the assignment problem in
      Lemma~\ref{lem:tropical}, proving the claimed minimum;
\item a product-formula comparison showing the upper bound is smaller than the
      certified divisor.
\end{enumerate}
Such certificates are compact enough for independent checking and eventual
Lean verification.  They also avoid trusting floating-point determinant
computations, which are extremely ill-conditioned in the near-singular regime.


\section{Source of the finite verifier}
\label{app:finite-verifier}

The following is the complete source used to establish
Theorem~\ref{thm:finite-range}.  Its only non-rigorous arithmetic---the
binary64 call used to suggest a starting integer---cannot affect the result:
the subsequent directed-interval sign search is monotone and independently
certifies the final pair of inequalities.

\begin{lstlisting}[style=proofpython]
#!/usr/bin/env python3
"""Rigorous finite verification for the Alaoglu--Erdos problem.

For every integer A in [2, LIMIT], this proves that A^(log(3)/log(2))
is an integer only when A is a power of two.  All decisive comparisons use
Decimal directed rounding plus the positive atanh series for logarithms.
The ordinary binary64 approximation is used only to choose an initial nearby
integer; a monotone certified search repairs any inaccurate guess.
"""
from __future__ import annotations

from dataclasses import dataclass
from decimal import Decimal, Context, ROUND_FLOOR, ROUND_CEILING, localcontext
import json
import math
import sys
import time

LIMIT = int(sys.argv[1]) if len(sys.argv) > 1 else 100_000
PREC = 70
TERMS = 28

DOWN = Context(prec=PREC, rounding=ROUND_FLOOR)
UP = Context(prec=PREC, rounding=ROUND_CEILING)

@dataclass(frozen=True)
class Interval:
    lo: Decimal
    hi: Decimal


def _atanh_log_ratio(p: int, q: int, terms: int = TERMS) -> Interval:
    """Enclose 2*atanh(p/q), for 0 <= p/q <= 1/3."""
    assert 0 <= p <= q // 3 + (1 if 3 * p == q else 0)
    if p == 0:
        z = Decimal(0)
        return Interval(z, z)

    with localcontext(DOWN):
        z_lo = Decimal(p) / Decimal(q)
        z2_lo = z_lo * z_lo
        power = z_lo
        total = Decimal(0)
        for j in range(terms):
            total = total + power / Decimal(2 * j + 1)
            power = power * z2_lo
        lo = Decimal(2) * total

    with localcontext(UP):
        z_hi = Decimal(p) / Decimal(q)
        z2_hi = z_hi * z_hi
        power = z_hi
        total = Decimal(0)
        for j in range(terms):
            total = total + power / Decimal(2 * j + 1)
            power = power * z2_hi
        partial = Decimal(2) * total
        # power is now z^(2*terms+1).  Since all omitted terms are positive,
        # replace every odd denominator by the first one and sum geometrically.
        rem = Decimal(2) * power / (
            Decimal(2 * terms + 1) * (Decimal(1) - z2_hi)
        )
        hi = partial + rem
    return Interval(lo, hi)


LOG2 = _atanh_log_ratio(1, 3)


def log_integer(n: int, terms: int = TERMS) -> Interval:
    """Rigorous interval for log(n), n >= 1, by binary range reduction."""
    assert n >= 1
    if n == 1:
        z = Decimal(0)
        return Interval(z, z)
    k = n.bit_length() - 1
    two_k = 1 << k
    # n/2^k in [1,2), and z=(y-1)/(y+1) lies in [0,1/3).
    local = _atanh_log_ratio(n - two_k, n + two_k, terms)
    with localcontext(DOWN):
        lo = Decimal(k) * LOG2.lo + local.lo
    with localcontext(UP):
        hi = Decimal(k) * LOG2.hi + local.hi
    return Interval(lo, hi)


LOG3 = log_integer(3)


def determinant_log_interval(log_a: Interval, b: int) -> Interval:
    """Enclose F(A,b)=log(b)log(2)-log(A)log(3)."""
    log_b = log_integer(b)
    with localcontext(DOWN):
        lo = log_b.lo * LOG2.lo - log_a.hi * LOG3.hi
    with localcontext(UP):
        hi = log_b.hi * LOG2.hi - log_a.lo * LOG3.lo
    return Interval(lo, hi)


def certified_sign(log_a: Interval, b: int) -> tuple[int, Interval]:
    f = determinant_log_interval(log_a, b)
    if f.hi < 0:
        return -1, f
    if f.lo > 0:
        return +1, f
    # At this finite precision an interval meeting zero is not silently accepted.
    raise ArithmeticError(f"inconclusive interval at A-log={log_a}, b={b}: {f}")


def is_power_of_two(n: int) -> bool:
    return n > 0 and (n & (n - 1)) == 0


def main() -> None:
    start = time.monotonic()
    rho_approx = math.log(3.0) / math.log(2.0)
    min_abs_endpoint = Decimal('Infinity')
    max_repairs = 0
    checked = 0

    for a in range(2, LIMIT + 1):
        if is_power_of_two(a):
            continue
        checked += 1
        log_a = log_integer(a)
        guess = max(1, int(math.floor(math.exp(rho_approx * math.log(float(a))))))
        repairs = 0

        # Find consecutive integers b,b+1 on opposite sides of the unique real root.
        while True:
            s0, f0 = certified_sign(log_a, guess)
            s1, f1 = certified_sign(log_a, guess + 1)
            min_abs_endpoint = min(
                min_abs_endpoint,
                abs(f0.lo), abs(f0.hi), abs(f1.lo), abs(f1.hi)
            )
            if s0 < 0 < s1:
                break
            if s1 < 0:
                guess += 1
            elif s0 > 0:
                guess -= 1
                assert guess >= 1
            else:
                raise AssertionError((a, guess, s0, s1))
            repairs += 1
            if repairs > 100:
                raise AssertionError(f"unexpectedly bad initial approximation at A={a}")
        max_repairs = max(max_repairs, repairs)

    elapsed = time.monotonic() - start
    result = {
        "limit": LIMIT,
        "checked_nonpowers": checked,
        "precision_decimal_digits": PREC,
        "atanh_terms": TERMS,
        "max_initial_guess_repairs": max_repairs,
        "minimum_absolute_interval_endpoint": str(min_abs_endpoint),
        "elapsed_seconds": elapsed,
        "conclusion": (
            f"For every integer 2 <= A <= {LIMIT}, "
            "A^(log(3)/log(2)) is integral only if A is a power of two."
        ),
    }
    print(json.dumps(result, indent=2, sort_keys=True))


if __name__ == '__main__':
    main()
\end{lstlisting}

\section{External literature reconnaissance}
\label{app:literature}

This appendix merges two literature-and-strategy surveys of the same corpus.  It records what
the surrounding literature does and does not supply, which imports are immediately usable, and
which are techniques that would need adaptation.  Two claims made in the source surveys are
retracted here; both are flagged where they occur.

Citations are given in full in the text rather than through the bibliography, because the
purpose of the appendix is provenance rather than argument: nothing below is used as a step in
any proof in this report.

\subsection{Effective measures for the two logarithms}

The sharpest linear-independence measure for the relevant triple is Qiang Wu, \emph{On the
linear independence measure of logarithms of rational numbers}, Math.\ Comp.\ \textbf{72}
(2003), 901--911, which gives $\mu(1,\log2,\log3)\le7.6155$, improving Rhin's $7.616$.  For the
single logarithm, Wu and Wang, J.\ Number Theory \textbf{142} (2014), 264--273, give
$\mu(\log3)\le5.1163051$, improving Salikhov's $5.125$.  For two-logarithm forms,
Laurent--Mignotte--Nesterenko, J.\ Number Theory \textbf{55} (1995), 285--321, and Laurent,
Acta Arith.\ \textbf{133} (2008), 325--348, give explicit
$\log|b_1\log2+b_2\log3|\ge-C(\log B)^2$ with $C$ of the order of a few tens.  On the
non-archimedean side, Bugeaud--Laurent (1996) and Bugeaud (1999) are the two-logarithm
workhorses, and K.~C.~Chim, J.\ Number Theory (2025), improves the $(\log B)^2$ dependence to
$\log B$ for two $p$-adic logarithms.

\begin{warning}[Retracted claim]
The first survey proposed that Wu's measure would yield a \emph{uniform} zero-free criterion
replacing the per-value certificates of Section~\ref{sec:finitecheck}.  That is wrong, and the
error is instructive.  Near-integrality of $U^{\thetaq}$ is not a rational approximation to
$\thetaq$; it is the vanishing of the bilinear form $\log3\log U-\log2\log V$, which is exactly
the wall of Section~\ref{sec:barrier}.  No linear-in-logarithms measure touches it.  The finite
verification stays per-value.
\end{warning}

What the measure does give unconditionally is an effective irrationality exponent
$\mu_{\mathrm{eff}}(\thetaq)\le8.62$, hence a polynomial recurrence $q_{n+1}\ll q_n^{7.62}$ for
the convergent denominators of $\thetaq$.  That is the strongest known bound of its kind and
supersedes the exponential estimate; the kernel-verified module \lean{DenominatorGrowth} proves
only the weaker exponential form, which remains the strongest \statuslean{} statement here.

\subsection{Interpolation determinants with non-archimedean amplification}

The closest existing realization of the amplification idea audited in
Section~\ref{sec:minplus} is Bombieri, \emph{Effective Diophantine approximation on
$\mathbb G_m$}, Ann.\ Scuola Norm.\ Sup.\ Pisa \textbf{20} (1993), 61--89, and Bombieri--Cohen,
\emph{ibid.} \textbf{24} (1997), 205--225, with the elementary account in LMS Lecture Notes
\textbf{303} (2003), 41--62.  These combine archimedean and non-archimedean interpolation
determinants through the equivariant Thue--Siegel principle and are the only non-Baker route to
effective results on the multiplicative group.  Laurent's interpolation-determinant papers
(Acta Arith.\ \textbf{66} (1994), 181--199, and \textbf{133} (2008)) supply the determinant
method in the form this report already uses.

This is a technique needing adaptation, not a plug-in.  Its relevance is sharpened by
Corollary~\ref{mp:cor-inert}: any amplification must obtain its archimedean bound from genuine
cancellation, and Bombieri--Cohen is where that is done systematically.

\subsection{Small-value estimates and the four-exponentials frontier}

Roy, \emph{An arithmetic criterion for the values of the exponential function}, Acta Arith.\
\textbf{97} (2001), 183--194, constructs a small-value criterion on
$\mathbb G_a\times\mathbb G_m$ and proves it \emph{equivalent} to Schanuel's conjecture.  His
subsequent estimates --- Acta Arith.\ \textbf{135} (2008), 357--393; Int.\ J.\ Number Theory
\textbf{6} (2010), 919--956; Mathematika \textbf{59} (2013), 333--363 --- and Nguyen--Roy, Int.\
J.\ Number Theory \textbf{12} (2016), 1273--1293, are the deepest partial progress toward the
four-exponentials and Schanuel circle.  Roy's Strong Six Exponentials Theorem, J.\ Number
Theory \textbf{41} (1992), 22--47, is the object whose formalization would most directly
strengthen the algebraic-base spectrum of Section~\ref{sec:spectrum}.

Butler, \emph{Some cases of Wilkie's conjecture}, Bull.\ London Math.\ Soc.\ \textbf{44}
(2012), 642--660, shows that certain cases of Wilkie's conjecture are equivalent to real forms
of the three- and four-exponentials conjectures.  That is a direct bridge between the counting
route below and the transcendence core, and it was not previously cited in this corpus.

\subsection{Counting rational points on the transcendental arc}

A counterexample produces $\asymp T$ \emph{rational} points
\[
  P_k=\left(\frac{M^k}{2^{r_k}},\ \frac{A^k}{3^{r_k}}\right)
   =\bigl(2^{\fract{k\beta}},3^{\fract{k\beta}}\bigr)
\]
of logarithmic height at most $T$, lying simultaneously in a \emph{fixed} rank-two
multiplicative subgroup $\Gamma\subset\mathbb G_m^2(\Qbar)$ and on the \emph{fixed}
transcendental Pfaffian arc $\mathcal C=\{(u,u^{\thetaq}):1\le u<2\}$.  This unlikely-
intersection formulation is the most promising framing the surveys contribute, and it is not
used elsewhere in this report.

The counting technology is Binyamini--Novikov, Ann.\ of Math.\ \textbf{186} (2017), 237--275,
and Binyamini--Novikov--Zak, Ann.\ of Math.\ \textbf{199} (2024), 795--821, which prove
Wilkie's conjecture for restricted elementary and Pfaffian structures with $(\log H)^\kappa$
counting; and, on the non-archimedean side, Binyamini--Cluckers--Novikov (Duke Math.\ J.),
Cluckers--Comte--Loeser, Forum Math.\ Pi \textbf{3} (2015), and \emph{Counting rational points
on transcendental curves in valued fields}, arXiv:2506.19411 (2025).  Boxall--Jones--Schmidt,
J.\ Eur.\ Math.\ Soc.\ \textbf{24} (2022), 1567--1592, combine archimedean counting with
$p$-adic methods.

The honest verdict is negative in a familiar way.  Polylogarithmic counting gives
$(\log H)^\kappa$ where a counterexample supplies $\asymp\log H$ points, so the framing
reproduces the one-logarithm deficit of Section~\ref{sec:deficit} in counting form: a third
independent route arriving at the same shortfall.  Moreover no cited theorem is sensitive to
\emph{which} primes divide the denominators, which is precisely what the compact-arc
reformulation of Section~\ref{sec:reform} needs.  The statement required --- that a compact
transcendental arc of nonzero curvature carries $o(H)$ points of height at most $H$ whose first
coordinate has $2$-power denominator and second $3$-power denominator --- is not implied by any
of them and would be a new theorem.

\subsection{Multiplicative dependence, \texorpdfstring{$S$}{S}-units, and rigidity}

Bugeaud--Corvaja--Zannier, Math.\ Z.\ \textbf{243} (2003), 79--84, prove that for
multiplicatively independent $a,b\ge2$ and every $\varepsilon>0$ one has
$\log\gcd(a^n-1,b^n-1)\le\varepsilon n$ for large $n$, and that this is sharp.  Together with
the Corvaja--Zannier subspace-theorem machinery and Evertse--Gy\H{o}ry's effective $S$-unit
theory, this confirms that $S$-unit methods deliver finiteness, not nonexistence, and cannot
supply the exponential-in-$q$ divisibility the rational-approximation route would need.

For the $\times2\times3$ side, Bourgain--Lindenstrauss--Michel--Venkatesh, Ergodic Theory
Dynam.\ Systems \textbf{29} (2009), 1705--1722, give the only \emph{effective} Furstenberg-type
density result, with a density gap of order $(\log\log L)^{-\kappa}$.  Shmerkin and Wu, Ann.\
of Math.\ \textbf{189} (2019), resolved Furstenberg's intersection conjecture.  None of these
can force visits to targets shrinking like $M^{-k}$; the target-width mismatch is fatal for a
direct application, exactly as Section~\ref{sec:reform} anticipates.  The effective
Bourgain--Lindenstrauss--Michel--Venkatesh rate is the one usable quantitative handle and
confirms, rather than defeats, the shrinking-target obstruction.

\subsection{Kummer theory and the radical reformulation}

Perucca--Sgobba--Tronto, Int.\ J.\ Number Theory \textbf{17} (2021), 1091--1110, and
\emph{Kummer theory for number fields via entanglement groups}, Manuscripta Math.\ \textbf{169}
(2022), 251--270, compute degrees of Kummer extensions exactly through entanglement groups, and
Chan--Pajaziti--Perissinotto--Perucca, arXiv:2508.19211 (2025), complete Kneser's theorem on
additive relations among radicals.  These make the algebraic half of the radical reformulation
of Section~\ref{sec:radical} rigorous and formalizable.  They do not see the special value
$\thetaq$: the mixed archimedean--radical invariant that the reformulation needs is supplied by
no Kummer-theoretic result found.

\subsection{Why the generic-power model theory excludes this exponent}

Bays, Kirby and Wilkie, J.\ London Math.\ Soc.\ (2010), prove a Schanuel property for raising
to an \emph{exponentially transcendental} power: for such $\lambda$ and multiplicatively
independent $\bar y$ one has $\operatorname{td}(\bar y,\bar y^{\lambda})\ge n$.  The result does
not apply here, and the reason is worth recording precisely, because it explains a pattern
rather than merely reporting a failure.

The exponent $\thetaq$ is exponentially \emph{algebraic}.  The pair
$(z,w)=(\log2,\thetaq)$ is a nonsingular solution of the system
\[
  e^{z}-2=0,\qquad e^{zw}-3=0,
\]
whose Jacobian determinant is $6\log2\ne0$.  So $\thetaq$ lies in exactly the exceptional set
that the generic-power theorems are stated to avoid.  The same holds for every version of the
problem in this report, since all of them involve the same exponent.  Generic o-minimal and
model-theoretic machinery should therefore be treated as a structural guide here, not as a
black box expected to apply.

\subsection{Prioritized directions, and what the ceiling changes}

The second survey graded the located literature by how directly it attaches to the obstruction.
Its ranking, compressed, was: highest priority to the rank-two multiplicative group combined
with specialized o-minimal or algebraic covering and $S$-unit rigidity; to a Laurent-style
interpolation-determinant redesign of the semigroup determinant of
Section~\ref{sec:semigroup}; and to tropical or $p$-adic determinant amplification with modern
$p$-adic logarithmic-form control.  High priority to the Corvaja--Zannier recurrence and
power-sum machinery, once an integral quotient has been manufactured, and to the
Mahler--Ridout--Corvaja--Zannier exponential-nearness criterion, which supplies a sharp
conditional target: construct one non-Pisot algebraic power sequence exponentially close to the
integers.  Medium priority to the explicit logarithmic-form technology of Matveev, Laurent,
Gouillon and Mignotte--Voutier, excellent for finite verification and auxiliary lemmas but
unlikely to cross the bilinear barrier alone, and to integer-valued o-minimal function theory,
methodologically close but not fitted to a sparse multiplicative grid.  Low priority, as a
direct attack, to Skolem--Mahler--Lech, generic-power model theory and general Zilber--Pink.

Every item in that survey was recorded as entirely absent from the corpus at the time, and the
present report's bibliography still does not carry most of them; they are listed here so that
the gap is visible rather than implicit.

One entry of the ranking has since been settled negatively.  The third highest-priority
direction, tropical or $p$-adic determinant amplification, is retired at the term-wise level by
Corollary~\ref{mp:cor-inert}: automatic divisibility can never contradict a correct archimedean
bound, whatever the primes, nodes or columns.  What survives of that direction is the
cancellation problem localized by Proposition~\ref{mp:prop-unique}, which is a different and
harder target than the one the ranking had in view.

\subsection{Provenance of the conjecture}

Alaoglu and Erd\H{o}s proved that the ratio of consecutive superabundant numbers is prime, and,
using the three-prime case that Siegel had claimed, that the ratio of consecutive colossally
abundant numbers is a prime or a semiprime.  Erd\H{o}s and Nicolas later showed that exactly
one, two, or four integers can attain the maximum in the colossally abundant construction, so a
positive answer to the rational-power form of the conjecture bounds this to at most two and
makes the ratio of consecutive colossally abundant numbers prime.  Lagarias, Amer.\ Math.\
Monthly \textbf{109} (2002), 534--543, records the Robin/Riemann reformulation of the
surrounding circle.

\begin{warning}[Unverified figure]
The first survey quoted a frequently repeated count of colossally abundant numbers below
$10^{18}$, attributed to Briggs, Experiment.\ Math.\ \textbf{15} (2006), 251--256.  That figure
could not be verified against the accessible text and is not repeated here.  Any use of it
should re-derive it.
\end{warning}

\subsection{The formalization landscape}

Karatarakis and Wiedijk formalized the Gelfond--Schneider theorem in Lean~4, including a
formalized Siegel lemma for algebraic integers and the auxiliary-function machinery; that
development is the basis of the port used by this corpus and cited in
Appendix~\ref{app:classical}.  Mathlib's Lindemann--Weierstrass development supplies analytic
scaffolding only.  No formalization of Baker's theorem, of the six exponentials theorem, or of
linear forms in logarithms was found in any proof assistant, so the transcendence inputs
isolated behind interfaces in this report would be greenfield work.

\subsection{Summary of what the literature does not supply}

Three gaps are unanimous across the surveyed material.  No lower-bound technology handles
linear or bilinear forms whose \emph{coefficients} are themselves logarithms, which is the wall
of Section~\ref{sec:barrier}.  No counting theorem is sensitive to the arithmetic of
denominators, which is what the compact-arc reformulation needs.  And no unconditional
algebraic-independence result reaches products of two logarithms, which is what the criterion
of Section~\ref{sec:logproduct} would need.  The conjecture remains open, and every route
surveyed inherits the same rank-two-to-rank-one collapse.

\begin{statusbox}{\statusknown{} throughout, with the two retractions marked as such.  Nothing
in this appendix is used as a step in any argument of this report; it records provenance and
the state of the surrounding literature.}
\end{statusbox}

\begin{thebibliography}{99}

\bibitem{AE1944}
L.~Alaoglu and P.~Erd\H{o}s,
\newblock \emph{On highly composite and similar numbers},
\newblock Transactions of the American Mathematical Society \textbf{56} (1944), 448--469.
\newblock DOI:
\href{https://doi.org/10.1090/S0002-9947-1944-0011087-2}{10.1090/S0002-9947-1944-0011087-2}.

\bibitem{Ram1915}
S.~Ramanujan,
\newblock \emph{Highly composite numbers},
\newblock Proceedings of the London Mathematical Society (2) \textbf{14} (1915), 347--409;
annotated manuscript in Ramanujan Journal \textbf{1} (1997), 119--153.

\bibitem{Baker1975}
A.~Baker,
\newblock \emph{Transcendental Number Theory},
\newblock Cambridge University Press, 1975.

\bibitem{Hua1982}
L.-K.~Hua,
\newblock \emph{Introduction to Number Theory},
\newblock Springer, 1982, Chapter~17.9.

\bibitem{Butler2017}
L.~A.~Butler,
\newblock \emph{A Diophantine approach to the three and four exponentials conjectures},
\newblock Ramanujan Journal \textbf{42}(1) (2017), 199--221.
\newblock DOI: \href{https://doi.org/10.1007/s11139-015-9728-2}{10.1007/s11139-015-9728-2}.

\bibitem{KarlinStudden}
S.~Karlin and W.~J.~Studden,
\newblock \emph{Tchebycheff Systems: With Applications in Analysis and Statistics},
\newblock Interscience, 1966.

\bibitem{KuipersNiederreiter}
L.~Kuipers and H.~Niederreiter,
\newblock \emph{Uniform Distribution of Sequences},
\newblock Wiley, 1974; Dover reprint, 2006.

\bibitem{Roy1992}
D.~Roy,
\newblock \emph{Matrices whose coefficients are linear forms in logarithms},
\newblock Journal of Number Theory \textbf{41}(1) (1992), 22--47.
\newblock DOI:
\href{https://doi.org/10.1016/0022-314X(92)90081-Y}{10.1016/0022-314X(92)90081-Y}.

\bibitem{Schneider1957}
Th.~Schneider,
\newblock \emph{Einf\"uhrung in die transzendenten Zahlen},
\newblock Springer, 1957.

\bibitem{Waldschmidt2000}
M.~Waldschmidt,
\newblock \emph{Diophantine Approximation on Linear Algebraic Groups: Transcendence
Properties of the Exponential Function in Several Variables},
\newblock Grundlehren der mathematischen Wissenschaften 326, Springer, 2000.
\newblock DOI: \href{https://doi.org/10.1007/978-3-662-11569-5}{10.1007/978-3-662-11569-5}.

\bibitem{Waldschmidt2005}
M.~Waldschmidt,
\newblock \emph{Variations on the Six Exponentials Theorem},
\newblock in \emph{Algebra and Number Theory}, pp.~338--355, 2005.
\newblock DOI:
\href{https://doi.org/10.1007/978-93-86279-23-1_23}{10.1007/978-93-86279-23-1\_23}.

\bibitem{WaldschmidtSix}
M.~Waldschmidt,
\newblock \emph{Six exponentials theorem --- irrationality},
\newblock expository note.

\bibitem{Waldschmidt2023}
M.~Waldschmidt,
\newblock \emph{The Four Exponentials Problem and Schanuel's Conjecture},
\newblock in \emph{Mathematics Going Forward}, Lecture Notes in Mathematics, 2023,
pp.~579--592.
\newblock DOI:
\href{https://doi.org/10.1007/978-3-031-12244-6_39}{10.1007/978-3-031-12244-6\_39}.

\bibitem{TheoremDB2026}
TheoremDB contributors,
\newblock \emph{Four exponentials conjecture}, problem P4, reviewed July 31, 2026.
\newblock \url{https://theoremdb.org/statements/P4/}.

\bibitem{mpmath}
F.~Johansson et al.,
\newblock \emph{mpmath: a Python library for arbitrary-precision floating-point arithmetic},
\newblock version 1.3.0.
\newblock \url{https://mpmath.org/}.

\bibitem{ProveIt}
The ProveIt project,
\newblock \emph{ProveIt repository},
\newblock \url{https://github.com/VladimirReshetnikov/ProveIt},
\newblock pinned in this report at commit
\commit{5d3982ceb12644fdd3499569d9db5f62f9a31dc6}.

\bibitem{CombinedReport}
The ProveIt project,
\newblock \emph{The Alaoglu--Erd\H{o}s Integer-Exponent Conjecture: Machine-Checked Partial
Results, Exact Conditional Structure, and Research Directions},
\newblock source directory
\lean{Analysis/ExponentialIdentities/Article/alaoglu-erdos-combined-report/},
\newblock file \lean{alaoglu-erdos-combined-report.tex}, August 20, 2026.

\bibitem{FurtherReport}
The ProveIt project,
\newblock \emph{The Alaoglu--Erd\H{o}s Integer-Exponent Conjecture: Further Determinant
Methods},
\newblock source directory
\lean{Analysis/ExponentialIdentities/Article/alaoglu-erdos-further-report/},
\newblock file \lean{alaoglu-erdos-further-report.tex}, August 20, 2026.

% --- imported with the continuation part ---

\bibitem{BakerWustholz1993}
A.~Baker and G.~W\"ustholz,
\newblock \emph{Logarithmic forms and group varieties},
\newblock Journal f\"ur die reine und angewandte Mathematik \textbf{442} (1993), 19--62.
\newblock DOI: \href{https://doi.org/10.1515/crll.1993.442.19}{10.1515/crll.1993.442.19}.

\bibitem{Matveev2000}
E.~M.~Matveev,
\newblock \emph{An explicit lower bound for a homogeneous rational linear form in the logarithms of algebraic numbers. II},
\newblock Izvestiya: Mathematics \textbf{64}(6) (2000), 1217--1269.
\newblock DOI: \href{https://doi.org/10.1070/IM2000v064n06ABEH000314}{10.1070/IM2000v064n06ABEH000314}.

\bibitem{BugeaudLaurent1996}
Y.~Bugeaud and M.~Laurent,
\newblock \emph{Minoration effective de la distance $p$-adique entre puissances de nombres alg\'ebriques},
\newblock Journal of Number Theory \textbf{61}(2) (1996), 311--342.
\newblock DOI: \href{https://doi.org/10.1006/jnth.1996.0152}{10.1006/jnth.1996.0152}.

\bibitem{Bugeaud2002}
Y.~Bugeaud,
\newblock \emph{Linear forms in two $m$-adic logarithms and applications to Diophantine problems},
\newblock Compositio Mathematica \textbf{132}(2) (2002), 137--158.
\newblock DOI: \href{https://doi.org/10.1023/A:1015825809661}{10.1023/A:1015825809661}.

\bibitem{Yu1998}
K.~Yu,
\newblock \emph{$p$-adic logarithmic forms and group varieties I},
\newblock Journal f\"ur die reine und angewandte Mathematik \textbf{502} (1998), 29--92.
\newblock DOI: \href{https://doi.org/10.1515/crll.1998.090}{10.1515/crll.1998.090}.

\bibitem{Wu2003}
Q.~Wu,
\newblock \emph{On the linear independence measure of logarithms of rational numbers},
\newblock Mathematics of Computation \textbf{72}(242) (2003), 901--911.
\newblock DOI: \href{https://doi.org/10.1090/S0025-5718-02-01442-4}{10.1090/S0025-5718-02-01442-4}.

\bibitem{MorseHedlund1940}
M.~Morse and G.~A.~Hedlund,
\newblock \emph{Symbolic dynamics II. Sturmian trajectories},
\newblock American Journal of Mathematics \textbf{62}(1) (1940), 1--42.
\newblock DOI: \href{https://doi.org/10.2307/2371431}{10.2307/2371431}.

\bibitem{Lothaire2002}
M.~Lothaire,
\newblock \emph{Algebraic Combinatorics on Words},
\newblock Encyclopedia of Mathematics and its Applications 90, Cambridge University Press, 2002.

\bibitem{Cobham1972}
A.~Cobham,
\newblock \emph{Uniform tag sequences},
\newblock Mathematical Systems Theory \textbf{6} (1972), 164--192.
\newblock DOI: \href{https://doi.org/10.1007/BF01706087}{10.1007/BF01706087}.

\bibitem{AlloucheShallit2003}
J.-P.~Allouche and J.~Shallit,
\newblock \emph{Automatic Sequences: Theory, Applications, Generalizations},
\newblock Cambridge University Press, 2003.

\bibitem{HarishChandra1957}
Harish-Chandra,
\newblock \emph{Differential operators on a semisimple Lie algebra},
\newblock American Journal of Mathematics \textbf{79}(1) (1957), 87--120.
\newblock DOI: \href{https://doi.org/10.2307/2372387}{10.2307/2372387}.

\bibitem{ItzyksonZuber1980}
C.~Itzykson and J.-B.~Zuber,
\newblock \emph{The planar approximation. II},
\newblock Journal of Mathematical Physics \textbf{21}(3) (1980), 411--421.
\newblock DOI: \href{https://doi.org/10.1063/1.524438}{10.1063/1.524438}.

\bibitem{McSwiggen2018}
C.~McSwiggen,
\newblock \emph{The Harish--Chandra integral: an introduction with examples},
\newblock L'Enseignement Math\'ematique \textbf{67}(3--4) (2021), 229--299;
preprint arXiv:1806.11155.
\newblock DOI: \href{https://doi.org/10.4171/LEM/1017}{10.4171/LEM/1017}.

\bibitem{PilaWilkie2006}
J.~Pila and A.~J.~Wilkie,
\newblock \emph{The rational points of a definable set},
\newblock Duke Mathematical Journal \textbf{133}(3) (2006), 591--616.
\newblock DOI: \href{https://doi.org/10.1215/S0012-7094-06-13336-7}{10.1215/S0012-7094-06-13336-7}.

% --- imported with the second continuation ---

\bibitem{BombieriGubler}
E.~Bombieri and W.~Gubler,
\newblock \emph{Heights in Diophantine Geometry},
\newblock Cambridge University Press, 2006.

\bibitem{HardyLittlewoodPolya}
G.~H.~Hardy, J.~E.~Littlewood, and G.~P\'olya,
\newblock \emph{Inequalities}, second edition,
\newblock Cambridge University Press, 1952.

\bibitem{LaurentMignotteNesterenko}
M.~Laurent, M.~Mignotte, and Yu.~Nesterenko,
\newblock Formes lin\'eaires en deux logarithmes et d\'eterminants
 d'interpolation,
\newblock \emph{Journal of Number Theory} 55 (1995), 285--321.

\bibitem{LangTranscendental}
S.~Lang,
\newblock \emph{Introduction to Transcendental Numbers},
\newblock Addison--Wesley, 1966.

\bibitem{DasguptaKakde2026}
S.~Dasgupta and M.~Kakde,
\newblock \emph{Ranks of matrices of logarithms of algebraic numbers II: The Matrix
Coefficient Conjecture},
\newblock Advances in Mathematics \textbf{487} (2026), Article 110753.
\newblock DOI:
\href{https://doi.org/10.1016/j.aim.2025.110753}{10.1016/j.aim.2025.110753};
arXiv:2408.08178.

\bibitem{Massold2025}
H.~Massold,
\newblock \emph{Transcendence degrees of fields generated by exponentials of products},
\newblock arXiv:2506.01123v1, June 2025.

\bibitem{EvertseSchlickeweiSchmidt2002}
J.-H.~Evertse, H.~P.~Schlickewei, and W.~M.~Schmidt,
\newblock \emph{Linear equations in variables which lie in a multiplicative group},
\newblock Annals of Mathematics \textbf{155} (2002), 807--836.

\bibitem{Carlson1921}
F.~Carlson,
\newblock \emph{"Uber Potenzreihen mit ganzzahligen Koeffizienten},
\newblock Mathematische Zeitschrift \textbf{9} (1921), 1--13.

\bibitem{Lech1953}
C.~Lech,
\newblock \emph{A note on recurring series},
\newblock Arkiv f\"or Matematik \textbf{2} (1953), 417--421.

\bibitem{Borwein1992}
P.~B.~Borwein,
\newblock \emph{On the irrationality of certain series},
\newblock Mathematical Proceedings of the Cambridge Philosophical Society
\textbf{112} (1992), 141--146.

\bibitem{Selberg1944}
A.~Selberg,
\newblock \emph{Bemerkninger om et multipelt integral},
\newblock Norsk Matematisk Tidsskrift \textbf{26} (1944), 71--78.

\end{thebibliography}
\end{document}
